\documentclass[11pt,a4paper]{article}

\usepackage[T1]{fontenc}
\usepackage[utf8]{inputenc}
\usepackage{lmodern}
\usepackage[margin=2.4cm]{geometry}
\usepackage{amsmath,amssymb,amsthm}
\usepackage{booktabs}
\usepackage{tabularx}
\usepackage{longtable}
\usepackage{array}
\usepackage{float}
\usepackage{listings}
\usepackage{xcolor}
\usepackage{algorithm}
\usepackage[noend]{algpseudocode}
\usepackage{enumitem}
\usepackage{mdframed}
\usepackage{titlesec}
\usepackage{needspace}
\usepackage[hidelinks,bookmarksnumbered]{hyperref}

% --------------------------------------------------------------------------
% Notation macros. Section files may use only these.
% --------------------------------------------------------------------------
\newcommand{\Fr}{\ensuremath{\mathbb{F}_r}}
\newcommand{\Fp}{\ensuremath{\mathbb{F}_p}}
\newcommand{\Fq}{\ensuremath{\mathbb{F}_q}}
\newcommand{\GO}{\ensuremath{\mathbb{G}_1}}
\newcommand{\GT}{\ensuremath{\mathbb{G}_2}}
\newcommand{\id}{\ensuremath{\mathcal{O}}}
\newcommand{\pair}[2]{\ensuremath{e\!\left(#1,\,#2\right)}}
\newcommand{\mptr}[1]{\texttt{#1}}
\newcommand{\yul}[1]{\texttt{#1}}
\newcommand{\opcode}[1]{\textsc{#1}}
\newcommand{\revert}[1]{\texttt{#1()}}
% mathtools provides \coloneqq; define it here so the document needs only the
% packages listed above.
\providecommand{\coloneqq}{\mathrel{\vcentcolon=}}
\providecommand{\vcentcolon}{:}

% --------------------------------------------------------------------------
% Listings
% --------------------------------------------------------------------------
\definecolor{codebg}{rgb}{0.975,0.975,0.97}
\definecolor{codecomment}{rgb}{0.30,0.45,0.32}
\definecolor{codekw}{rgb}{0.10,0.20,0.60}
\definecolor{codestr}{rgb}{0.60,0.20,0.10}
\definecolor{coderule}{rgb}{0.85,0.85,0.85}

\lstdefinestyle{solidity}{
  backgroundcolor=\color{codebg},
  basicstyle=\ttfamily\scriptsize,
  commentstyle=\color{codecomment}\itshape,
  keywordstyle=\color{codekw}\bfseries,
  stringstyle=\color{codestr},
  breaklines=true,
  breakatwhitespace=false,
  columns=fullflexible,
  keepspaces=true,
  showstringspaces=false,
  frame=single,
  rulecolor=\color{coderule},
  framesep=4pt,
  xleftmargin=6pt,
  xrightmargin=2pt,
  numbers=none,
  captionpos=b,
  abovecaptionskip=4pt,
  belowcaptionskip=2pt,
  morecomment=[l]{//},
  morecomment=[s]{/*}{*/},
  morestring=[b]",
  morekeywords={function,let,if,for,switch,case,default,leave,revert,return,
    mstore,mload,mstore8,mcopy,calldataload,calldatacopy,calldatasize,
    staticcall,returndatasize,extcodesize,extcodehash,extcodecopy,keccak256,
    addmod,mulmod,and,or,xor,not,shl,shr,byte,eq,lt,gt,iszero,add,sub,mul,div,
    mod,pop,assembly,contract,constructor,external,internal,view,pure,returns,
    uint256,bytes32,bytes,address,immutable,constant,require,error,pragma,
    solidity,memory,calldata,public,true,false}
}
\lstset{style=solidity}
\renewcommand{\lstlistingname}{Listing}

% --------------------------------------------------------------------------
% Finding call-outs
% --------------------------------------------------------------------------
\newcounter{findingctr}
\newenvironment{finding}[2][]{%
  \needspace{3\baselineskip}%
  \refstepcounter{findingctr}%
  \begin{mdframed}[linewidth=0.8pt,linecolor=black!35,backgroundcolor=black!3,
                   innerleftmargin=8pt,innerrightmargin=8pt,
                   innertopmargin=6pt,innerbottommargin=6pt,skipabove=8pt,skipbelow=8pt]%
  \noindent\textbf{Finding \thefindingctr\ifx\relax#1\relax\else\ (#1)\fi: #2}\par\vspace{2pt}%
}{%
  \end{mdframed}%
}

\theoremstyle{definition}
\newtheorem{remark}{Remark}[section]

\titleformat{\section}{\Large\bfseries}{\thesection}{0.7em}{}
\setlist[itemize]{itemsep=1pt,topsep=3pt,parsep=0pt}
\setlist[enumerate]{itemsep=1pt,topsep=3pt,parsep=0pt}

% Many sections carry several wide tables. Loosen LaTeX's float placement so
% they are set promptly instead of accumulating into a "too many unprocessed
% floats" abort.
\setcounter{topnumber}{5}
\setcounter{bottomnumber}{5}
\setcounter{totalnumber}{10}
\renewcommand{\textfraction}{0.05}
\renewcommand{\topfraction}{0.95}
\renewcommand{\bottomfraction}{0.95}
\renewcommand{\floatpagefraction}{0.70}

\setcounter{tocdepth}{2}
\setcounter{secnumdepth}{3}
\sloppy
\emergencystretch=3em

\title{\textbf{A Specification of the Moonlight Wrap Halo2 BLS12-381 KZG Verifier}\\[4pt]
\large \texttt{fixtures/moonlight-wrap/Halo2Verifier.sol} and
\texttt{fixtures/moonlight-wrap/Halo2VerifyingKey.sol}}
\author{Extracted and assessed from the rendered artifacts}
\date{\today}

\begin{document}
\maketitle
\begin{abstract}
\noindent
This document is a line-faithful specification of the two Solidity contracts that
make up the Moonlight \emph{wrap} on-chain verifier: a circuit-specialised Halo2
KZG verifier over BLS12-381 (\texttt{Halo2Verifier.sol}, 4{,}130 lines) and the
data-only verifying-key contract it is pinned to
(\texttt{Halo2VerifyingKey.sol}, 695 lines). Every guard, constant, loop bound,
memory address, precompile call and revert condition present in the artifacts is
described here, together with the algebraic relations the code computes. The
specification is interleaved with an assessment of the design and implementation:
call-outs marked \emph{Finding} record observations a reviewer should act on or
consciously accept. Empirical results from an adversarial replay suite executed
against these exact artifacts are reported in
Section~\ref{sec:testing}.
\end{abstract}
\tableofcontents
\newpage


%% ==== 00-intro.tex ======================================================
\section{Scope, artifacts and reading guide}
\label{sec:intro:scope}

\subsection{The two artifacts}

This document specifies exactly two files:

\begin{center}
\begin{tabularx}{\textwidth}{@{}l r r X@{}}
\toprule
File & Lines & Deployed runtime & Role \\
\midrule
\texttt{Halo2Verifier.sol} & 4{,}130 & 21{,}880 B & Circuit-specialised Halo2/KZG verifier \\
\texttt{Halo2VerifyingKey.sol} & 695 & 17{,}025 B & Data-only contract holding the verifying key \\
\bottomrule
\end{tabularx}
\end{center}

Runtime sizes are measured with the pinned compiler
(\texttt{solc 0.8.30+commit.73712a01}, \texttt{--via-ir --optimize
--optimize-runs 1 --evm-version cancun --no-cbor-metadata}), which is the
configuration the artifact header pins. Both are below the EIP-170 limit of
24{,}576 bytes, with 2{,}696 bytes of headroom on the verifier.

Neither file is hand-written. Both are rendered by the Rust generator in this
repository from one \texttt{VerifyingKey<Fq, KZGCommitmentScheme<Bls12>>}. The
proof layout, the verifying-key payload, the quotient identity program, the
memory map and the gas bounds are all specialised to a single circuit --- the
Moonlight \emph{wrap} decider. Nothing in either contract is generic: changing
the circuit changes the bytecode.

\subsection{What the verifier proves}

\texttt{verifyProof(bytes calldata proof, uint256[] calldata instances)}
returns \texttt{true} if and only if \texttt{proof} is a valid Halo2 proof for
the pinned constraint system on the public inputs \texttt{instances}. It never
returns \texttt{false}: every rejection reverts with one of seven typed errors.

The proof system is Halo2 in the midnight-proofs flavour, instantiated with
KZG over BLS12-381 and a Keccak256 Fiat--Shamir transcript. The verifier
reduces the whole argument to a single BLS12-381 pairing check performed by
the EIP-2537 precompile at address \texttt{0x0f}:
\begin{equation}
\pair{\mathsf{C} - v\cdot G - (-x_3)\cdot\pi}{[1]_2}
\;=\;\pair{\pi}{[s]_2},
\label{eq:intro:kzg}
\end{equation}
where $\mathsf{C}$ is a single commitment fused from every opened commitment
in the proof, $v$ is the corresponding batched claimed evaluation, $x_3$ is the
Fiat--Shamir evaluation point and $\pi$ is the KZG opening proof. Because this
render carries a public IVC accumulator, a second pairing equation is folded
into the same check with a verifier-local randomiser $\alpha$
(Section~\ref{sec:fin:batch}).

\subsection{Protocol shape of this render}
\label{sec:intro:shape}

Every count below is a generated constant baked into the bytecode; the verifier
rejects any calldata that does not match it exactly.

\begin{center}
\begin{tabularx}{\textwidth}{@{}l l X@{}}
\toprule
Parameter & Value & Where pinned \\
\midrule
Domain $k$, $n = 2^k$ & $20$, $n = 1{,}048{,}576$ & VK word 2; cross-checked at \texttt{Halo2Verifier.sol:1401} \\
Public instances & $19$ & VK word 1; \texttt{Halo2Verifier.sol:1400} \\
Accumulator offset & $11$ & VK word 8; \texttt{Halo2Verifier.sol:1403} \\
Accumulator limbs & $7\times 56$ bits & VK words 9--10; \texttt{Halo2Verifier.sol:1404--1405} \\
Blinding rows & $9$ (rotation $-10$) & implied by \texttt{OMEGA\_INV\_TO\_L}; see \S\ref{sec:lag:convert} \\
Proof length & \texttt{0x1e60} $= 7{,}776$ B & \texttt{Halo2Verifier.sol:1417} \\
Calldata size & \texttt{0x2144} $= 8{,}516$ B & \texttt{Halo2Verifier.sol:1423} \\
Advice commitments & $15$ & parser loop, \texttt{0x0780} bytes \\
Lookup multiplicities & $2$ & parser loop, \texttt{0x0100} bytes \\
Permutation $Z$ & $6$ & parser loop, \texttt{0x0300} bytes \\
Lookup helper $+$ $Z$ & $2 + 2$ & two unrolled per-lookup blocks \\
Trashcan commitments & $1$ & parser loop, \texttt{0x0080} bytes \\
Quotient limbs & $4$ & parser loop, \texttt{0x0200} bytes \\
Polynomial evaluations & $102$ & parser loop, \texttt{0x0cc0} bytes \\
Point sets & $5$ (sizes $3,3,2,2,1$) & PCS blocks \\
$q$ evaluations & $5$ & parser loop, \texttt{0x00a0} bytes \\
Final MSM terms & $78$ & \texttt{Halo2Verifier.sol:3999} \\
Simple selectors & $10$ & selector zeroing loop, \texttt{0x0140} bytes \\
\bottomrule
\end{tabularx}
\end{center}

The proof stream therefore carries $34$ group elements
($15+2+6+4+1+4$ commitments, plus \texttt{f\_com} and $\pi$) and $107$ field
elements ($102$ evaluations plus $5$ $q$-evaluations).

\subsection{Execution outline}
\label{sec:intro:outline}

\begin{algorithm}[H]
\caption{\texttt{verifyProof} at a glance. Section references point at the
full specification of each step.}
\label{alg:intro:outline}
\begin{algorithmic}[1]
\State \textbf{ABI guard}: two head words and the solc-side decode
  (\S\ref{sec:abi:headguard})
\State \textbf{Memory guard}: assert \texttt{mload(0x40) $\le$ TRANSCRIPT\_MPTR}
  (\S\ref{sec:abi:mf2})
\State \textbf{VK load}: re-check \texttt{extcodesize}/\texttt{extcodehash},
  \texttt{extcodecopy} payload to \mptr{VK\_MPTR}, cross-check six header words
  (\S\ref{sec:vkload:pin})
\State \textbf{Envelope}: proof length, instance count, exact
  \texttt{calldatasize} (\S\ref{sec:vkload:abi})
\State \textbf{Accumulator}: decode the trailing $8$ instance words into two
  \GO{} points, validate the packing and route both through G1MSM
  (\S\ref{sec:acc:validate})
\State \textbf{Transcript}: absorb $\mathrm{vk\_digest}$, the committed-instance
  identity, the instance count, $19$ instances, then alternate
  commitment blocks and challenge squeezes
  $\theta,\beta,\gamma,\tau,y,x,x_1,x_2,x_3,x_4$ (\S\ref{sec:tp:schedule})
\State \textbf{Domain}: $x^n$, one batched inversion for
  $L_{\mathrm{last}},L_{\mathrm{blind}},L_0,\dots,L_{18}$ and $(x^n-1)^{-1}$,
  and the instance evaluation (\S\ref{sec:lag:top})
\State \textbf{Quotient}: interpret the VK-resident identity bytecode to
  rebuild $\nu_y(x)$; store $-\nu_y(x)$ and one scalar per simple selector
  (\S\ref{sec:qvms:why}--\S\ref{sec:qeps:lineval})
\State \textbf{PCS}: build the rotation points, the $x_1$ and $x_4$ power
  tables, the five point sets, and the batched value $v$
  (\S\ref{sec:pcs:sets})
\State \textbf{Fused MSM}: one $78$-pair G1MSM producing $\mathsf{C}$
  (\S\ref{sec:fin:msm})
\State \textbf{Pairing}: stage $\pi$ and $\mathsf{C} - vG + x_3\pi$, fold in the
  randomised accumulator equation, call \texttt{0x0f}
  (\S\ref{sec:fin:pairing})
\State \Return \texttt{true}
\end{algorithmic}
\end{algorithm}

\subsection{Notation}
\label{sec:intro:notation}

\begin{center}
\begin{tabularx}{\textwidth}{@{}l X@{}}
\toprule
Symbol & Meaning \\
\midrule
$r$ & BLS12-381 scalar-field order,
\texttt{0x73eda753299d7d483339d80809a1d80553bda402fffe5bfeffffffff00000001} \\
$p$ & BLS12-381 base-field order,
\texttt{0x1a0111ea...b9feffffffffaaab} ($381$ bits) \\
$n,\ k$ & Domain size $n = 2^k$ with $k = 20$ \\
$\omega$ & Primitive $n$-th root of unity in \Fr{} \\
$\theta,\beta,\gamma,\tau,y,x,x_1,x_2,x_3,x_4$ & Fiat--Shamir challenges, in
squeeze order \\
$\mathsf{A}_i,\mathsf{M}_i,\mathsf{Z}_i,\mathsf{Q}_i$ & Advice, lookup
multiplicity, permutation product and quotient limb commitments \\
$\mathsf{F},\ \pi$ & \texttt{f\_com} and the KZG opening proof \\
$\mathbf{e}=(e_0,\dots,e_{101})$ & The alleged polynomial evaluations at $x$ \\
$\nu_y(x)$ & The $y$-batched constraint numerator reconstructed by the
quotient VM \\
$L_0,L_{\mathrm{last}},L_{\mathrm{blind}}$ & Lagrange values at $x$ \\
\bottomrule
\end{tabularx}
\end{center}

Memory addresses are written as they appear in the artifact, e.g.\
\mptr{VK\_MPTR} for a named constant and \texttt{0xb280} for a bare literal.
A significant fraction of the generated body uses bare literals; where that
matters for readability it is called out in the assessment.

\subsection{How to read this document}

Sections~\ref{sec:const:top}--\ref{sec:fin:return} follow the artifact
top to bottom. Each begins with the structure of the region, gives the formulas
and algorithms it implements, quotes the code verbatim with line references,
and then assesses the design. Assessment call-outs are numbered
\emph{Findings} and carry a severity; they are consolidated in
Section~\ref{sec:assess:summary}. Section~\ref{sec:testing} reports what an
adversarial replay suite executed against these exact artifacts actually
observed; Section~\ref{sec:quality} critiques the artifact as something a human
must read; and Section~\ref{sec:robust} asks the separate question of whether
the checks it performs are the right checks, expressed the right way.

\begin{remark}
Line numbers in this document refer to the artifacts as committed at
\texttt{fixtures/moonlight-wrap/}. The companion render under
\texttt{deployments/sepolia/moonlight-wrap/} is an older, \texttt{\^{}0.8.24}
version of the same circuit; it shares the proof layout but predates the typed
error taxonomy and part of the quotient program encoding, so line numbers do
not transfer.
\end{remark}


%% ==== 01-constants.tex ==================================================
\section{Contract header, typed error taxonomy, and the generated constant/memory map}
\label{sec:const:top}

This section specifies \texttt{Halo2Verifier.sol} lines 1--350 of the
\texttt{fixtures/moonlight-wrap} artifact: the license and pragma header, the
NatSpec contract documentation, the eight declared custom errors and their rendered
selector constants, the verifying-key pinning constants, the ABI calldata
cursors, the complete generated memory map, the precompile gas bounds, the
build identity, and the field-arithmetic constants. Nothing in this range
executes; every declaration in it is either compile-time metadata or a literal
substituted into assembly elsewhere in the file. The section therefore does two
things: it states each constant exactly, and it derives the runtime
\emph{extent} of every memory region those constants name, so that later
sections can be read against a closed address map.

All numeric claims below were recomputed from the fixture. Where a claim
depends on code outside lines 1--350 (region extents are only observable at
their use sites) the defining line is cited explicitly and the behaviour itself
is specified in the section that owns it.

\subsection{Artifact identity and derived circuit shape}
\label{sec:const:shape}

The constants in this range are not free parameters; they are a projection of
one \texttt{VerifyingKey} and one proof shape. The following table records the
circuit parameters that the constants in lines 1--350 determine uniquely. Each
value is \emph{derived} from the fixture (by differencing adjacent generated
addresses or reading a pinned literal), not copied from the generator, and each
derivation is stated in the sections that follow.

\begin{table}[htbp]
\centering
\begin{tabular}{@{}llll@{}}
\toprule
Quantity & Symbol & Value & Derivation \\
\midrule
Public instances                & $m$        & $19$    & literal, line 1420 \\
Proof byte length               & --         & $\texttt{0x1e60} = 7776$ & literal, line 1419 \\
Exact \opcode{CALLDATASIZE}     & --         & $\texttt{0x2144} = 8516$ & line 1426 \\
Domain size exponent            & $k$        & $20$    & literal, lines 1401, 1800; $n = 2^{20}$ \\
Proof evaluations               & $|\mathbf{e}|$ & $102$ & loop bound \texttt{0x0cc0}, line 1699 \\
Advice commitments (all phases) & --         & $15$    & loop bound \texttt{0x0780}, line 1553 \\
Lookups                         & --         & $2$     & $(\texttt{0xa9c0}-\texttt{0xa8c0})/128$ \\
Permutation $Z$ chunks          & --         & $6$     & $(\texttt{0xacc0}-\texttt{0xa9c0})/128$ \\
Lookup helper chunks            & --         & $2$     & $(\texttt{0xadc0}-\texttt{0xacc0})/128$ \\
Trashcans                       & --         & $1$     & $(\texttt{0xaf40}-\texttt{0xaec0})/128$ \\
Quotient limbs                  & --         & $4$     & $(\texttt{0xb140}-\texttt{0xaf40})/128$ \\
Proof G1 commitments (total)    & --         & $32$    & $(\texttt{0xb140}-\texttt{0xa140})/128$ \\
Fixed (VK) G1 commitments       & --         & $45$    & $(\texttt{0x7900}-\texttt{0x6280})/128$ \\
Simple selectors                & --         & $10$    & line 3351 offset \texttt{0x120}; line 3066 bound \\
User-phase challenges           & --         & $0$     & $\mptr{CHALLENGE\_MPTR} = \mptr{THETA\_MPTR}$ \\
Last rotation                   & $\mathrm{rot}_{\mathrm{last}}$ & $-10$ & $-(\texttt{0x03a0}/32 - m) = -(29-19)$, line 1817 \\
Fused final-MSM terms           & --         & $78$    & $\texttt{0x30c0}/\texttt{0xa0}$, line 3999 \\
Powers of $x_1$ materialised    & --         & $43$    & loop bound \texttt{0x2a}, line 3450 \\
Accumulator encoding            & --         & \texttt{point\_pair} & lines 1270--1318 \\
Accumulator instance offset     & --         & $11$ words & line 1257 ($\texttt{0x0160}$), line 1403 \\
Accumulator limbs / limb bits   & --         & $7 / 56$ & lines 1248--1249, 1404--1405 \\
\bottomrule
\end{tabular}
\caption{Circuit shape recovered from the generated constants of
\texttt{Halo2Verifier.sol} lines 1--350.}
\label{tab:const:shape}
\end{table}

\subsection{Pragma pinning and the EIP-170 code-size budget}
\label{sec:const:pragma}

\begin{lstlisting}[style=solidity,caption={\texttt{Halo2Verifier.sol} lines 1--13: license and pinned pragma},label={lst:const:pragma}]
// SPDX-License-Identifier: CC0-1.0
// Pinned, not floating. Two properties of this artifact are compiler- and
// optimiser-dependent, and neither is visible in the source:
//   1. The generated layout writes absolute addresses from TRANSCRIPT_MPTR
//      upward. That is only safe while solc's stack-spill reservation stays
//      below it -- measured 0x8c0 on 0.8.24 and 0x8e0 on 0.8.26+, so it is not
//      a constant this file controls. verifyProof now asserts the separation.
//   2. Runtime size depends on --optimize-runs. Measured: 0.8.24 at runs=1
//      emits 29,567 bytes and 0.8.30 at runs=100000 emits 29,836 -- both over
//      the EIP-170 24,576-byte limit, so neither can be deployed. Only the
//      pinned (version, runs) pair is known to produce a deployable contract.
// A floating `^0.8.24` advertises compatibility this contract does not have.
pragma solidity 0.8.30;
\end{lstlisting}

\paragraph{What the pin buys.}
The contract body is one \yul{assembly ("memory-safe")} block that writes
absolute memory addresses starting at $\texttt{0x1000}$
(\mptr{TRANSCRIPT\_MPTR}) and never consults the free-memory pointer. The
\yul{memory-safe} annotation is what enables solc's via-IR stack-to-memory
mover --- without it the block does not compile (stack too deep) --- and that
mover reserves spill slots upward from $\texttt{0x80}$, recording the top of the
reservation in the runtime's \yul{mstore(0x40, ...)} prologue. The size of that
reservation is a property of the compiler release and the optimiser schedule,
not of this source file. Pinning the exact patch version removes one degree of
freedom; the runtime guard at line 674 removes the rest by failing closed:

\[
  \yul{mload(0x40)} > \mptr{TRANSCRIPT\_MPTR} \;\Longrightarrow\;
  \revert{MemoryLayoutViolated}
\]

The same guard is repeated in the creation frame at line 363. Both are
specified in detail in the \verb|verifyProof| prologue section; here it matters
only that the pragma pin and the guard are two halves of one argument, and that
the header comment says so.

\paragraph{The EIP-170 problem.}
EIP-170 caps deployed runtime code at $24{,}576$ bytes. The header records two
measured $(\text{version}, \texttt{--optimize-runs})$ pairs that produce
$29{,}567$ and $29{,}836$ runtime bytes respectively --- both undeployable. This
is the direct reason the verifying key is a \emph{separate contract}: the VK
payload alone is $17{,}024$ bytes (\S\ref{sec:const:vkpin}), and no arrangement
puts payload and verifier logic in one $24$\,KiB runtime. It is also why the
quotient identity program is stored as \emph{data} inside the VK runtime
(\texttt{0x11cf} bytes at \texttt{0x50a0} in the memory image, line 2016) and
interpreted by a compact VM, rather than unrolled into bytecode.

\begin{finding}[const-1]{Pragma pins the compiler but cannot pin the optimiser setting --- Low}
\texttt{Halo2Verifier.sol}:1--13. The header states that ``only the pinned
(version, runs) pair is known to produce a deployable contract'', but a Solidity
\texttt{pragma} can express the version and \emph{not} the
\texttt{--optimize-runs} value. The two measured pairs the comment names are
both over EIP-170; the deployable pair is not named anywhere in the artifact.
The value lives in the generator
(\texttt{src/evm.rs}:72, \texttt{DEFAULT\_OPTIMIZE\_RUNS = 200}; and
\texttt{src/evm.rs}:78,
\texttt{PINNED\_SOLC\_VERSION = "0.8.30+commit.73712a01"}), which an integrator
who receives only the \texttt{.sol} file does not have. Note also that the
pragma is an equality pin (\texttt{pragma solidity 0.8.30;}) and so does at
least prevent a silent minor-version bump, but it cannot distinguish
\texttt{0.8.30+commit.73712a01} from another build of the same version.

Impact is confined to reproducibility and deployment, not to soundness: a
mis-optimised build either fails to deploy (EIP-170) or trips the
\revert{MemoryLayoutViolated} guard, and in both cases fails closed. The
mitigation is documentation, not code: the deployment record must publish the
$(\text{solc version}+\text{commit}, \text{runs})$ pair alongside
$\yul{BUILD\_ID}$ (\S\ref{sec:const:buildid}). Note that $\yul{BUILD\_ID}$
does \emph{not} cover the compiler settings --- it is computed by the Rust
generator before solc runs --- so it cannot be used to detect this class of
drift.
\end{finding}

\subsection{Contract-level documentation contract}
\label{sec:const:natspec}

Lines 15--46 are NatSpec. They are not executable, but three claims in them are
load-bearing for the rest of this document and are restated here as
specification, each verified against the code:

\begin{enumerate}
  \item \textbf{Not a generic verifier} (lines 22--24). The proof layout, VK
        payload, quotient identity program, memory layout, and optional
        quotient evaluator are generated for one
        $\mathsf{VerifyingKey}\langle \Fq, \mathsf{KZG}(\mathsf{Bls12})\rangle$.
        Confirmed: every offset in Table~\ref{tab:const:map} is a literal.
  \item \textbf{$\Fp$ does not fit a word} (lines 30--32). $p$ is $381$ bits, so
        each $\Fp$ coordinate is carried EIP-2537 padded as $16$ zero bytes
        followed by $48$ coordinate bytes. A $\GO$ point is therefore $4$ words
        ($128$ bytes) and a $\GT$ point $8$ words ($256$ bytes). Every $\GO$
        slot in the map is $\texttt{0x80}$ bytes wide and every $\GT$ slot
        $\texttt{0x100}$; see \mptr{G2\_BASE\_MPTR} and
        \mptr{NEG\_S\_G2\_BASE\_MPTR}, which are exactly $\texttt{0x100}$ apart.
  \item \textbf{Calldata carries uncompressed padded $\GO$} (lines 33--39). The
        prover's native compressed $48$-byte stream is repacked off chain; the
        verifier hashes the uncompressed $128$-byte form verbatim into the
        transcript, matching
        \texttt{Hashable<Keccak256> for G1Projective::to\_input}. This is why
        \mptr{PROOF\_LEN\_CPTR} pins $\texttt{0x1e60}$ rather than a compressed
        length, and why there is no on-chain decompression routine to budget
        memory for.
\end{enumerate}

Lines 44--46 state the deployment precondition: compile with Solidity
$\geq 0.8.24$ and deploy only on chains supporting \opcode{MCOPY} and EIP-2537.
The constructor enforces both by execution (\S\ref{sec:const:gas}).

\subsection{Typed failure taxonomy}
\label{sec:const:errors}

\begin{lstlisting}[style=solidity,caption={\texttt{Halo2Verifier.sol} lines 48--83: taxonomy rationale and the eight declared errors (NatSpec bodies elided)},label={lst:const:errors}]
    // ----------------------------------------------------------------------
    // Typed failure taxonomy (P4/L-3, docs/audit/HALO2_VERIFIER_REVIEW).
    // verifyProof is success-or-revert; these errors let integrators and
    // incident responders distinguish malformed calldata from a swapped VK,
    // a non-canonical scalar, a failed precompile, or a rejected proof.
    // Constructor smoke probes intentionally keep bare reverts -- they report
    // a chain-capability failure, and the deployment transaction identifies
    // itself. The one constructor exception is the memory-layout guard
    // (MF-2), which reports a BUILD fault and is typed on both paths.
    // ----------------------------------------------------------------------
    // ... (each declaration below is preceded in the source by a
    // two-to-five-line /// @notice block, elided here; lines 58--83)
    error BadCalldataShape();
    error VkMismatch();
    error NonCanonicalScalar();
    error BadPointEncoding();
    error PrecompileFailed();
    error ProofRejected();
    error QuotientProgramInvalid();
    error MemoryLayoutViolated();
\end{lstlisting}

\subsubsection{Selector constants}

Lines 320--327 render the four-byte selectors as Yul-readable
\yul{uint256} constants (the rationale comment occupies lines 313--319). The generated \yul{fail(sel)} helper (lines 690--693)
writes \yul{shl(224, sel)} to scratch \texttt{0x00} and executes
\yul{revert(0x00, 0x04)}. The two guards that run before \yul{fail} is in scope
(lines 636--641 and 674--677) inline the same two statements, and
\yul{q\_program\_fail} (lines 1890--1893) inlines them a third time with a
literal selector; the source explains both exceptions at lines 637--638 and
669--673, and 1885--1889.

\begin{lstlisting}[style=solidity,caption={\texttt{Halo2Verifier.sol} lines 313--327: selector constants},label={lst:const:selectors}]
    // ----------------------------------------------------------------------
    // Typed-error selectors (P4/L-3): bytes4(keccak256("Name()")) of the
    // errors declared on the contract, as Yul-readable constants. The
    // `fail(sel)` helper in AssemblyHelpers.yul writes the selector to
    // scratch 0x00 and reverts with 4 bytes. Pinned by
    // `p4_error_selectors_match_declared_errors` in src/lowering/tests.rs.
    // ----------------------------------------------------------------------
    uint256 internal constant ERR_BAD_CALLDATA_SHAPE      = 0x1b99e37c;
    uint256 internal constant ERR_VK_MISMATCH             = 0xa447d73e;
    uint256 internal constant ERR_NON_CANONICAL_SCALAR    = 0x77530042;
    uint256 internal constant ERR_BAD_POINT_ENCODING      = 0xf27905ec;
    uint256 internal constant ERR_PRECOMPILE_FAILED       = 0x84e81692;
    uint256 internal constant ERR_PROOF_REJECTED          = 0xc3b0d8cd;
    uint256 internal constant ERR_QUOTIENT_PROGRAM_INVALID = 0x3cc81b89;
    uint256 internal constant ERR_MEMORY_LAYOUT_VIOLATED   = 0xc9888d23;
\end{lstlisting}

All eight values were recomputed as
$\mathrm{keccak256}(\text{signature})[0..4]$ and match bit for bit. The
rightmost column lists every site in the fixture that raises the error, so a
reader can see the taxonomy's actual discrimination power rather than its
intended one.

\begin{table}[htbp]
\centering
\begin{tabularx}{\textwidth}{@{}llX@{}}
\toprule
Error & Selector (verified) & Raised at (\texttt{Halo2Verifier.sol}) \\
\midrule
\revert{BadCalldataShape}     & \texttt{0x1b99e37c} & 639 (ABI heads), 1430 (proof length, instance count, exact \opcode{CALLDATASIZE}), 1785 (proof cursor did not land on \mptr{NUM\_INSTANCE\_CPTR}) \\
\revert{VkMismatch}           & \texttt{0xa447d73e} & 1389 (VK \opcode{EXTCODESIZE}/\opcode{EXTCODEHASH}), 1406 (VK header cross-check) \\
\revert{NonCanonicalScalar}   & \texttt{0x77530042} & 708, 709 (\yul{scalar\_inv} domain), 1523 (instance word), 1706 (proof evaluation), 1759, 1789 \\
\revert{BadPointEncoding}     & \texttt{0xf27905ec} & 780, 781 (EIP-2537 zero padding), 786, 789 (coordinate $\ge p$), 1452 \\
\revert{PrecompileFailed}     & \texttt{0x84e81692} & 719, 720 (modexp), 982 (pairing), 1451, 1881, 4113 \\
\revert{ProofRejected}        & \texttt{0xc3b0d8cd} & 962 (\yul{ec\_pairing} entered with \yul{success} already clear), 988 (pairing returned $\neq 1$), 1882 (Lagrange denominator rejected), 4125 (terminal guard) \\
\revert{QuotientProgramInvalid} & \texttt{0x3cc81b89} & 1891 only, inside \yul{q\_program\_fail()} (lines 1890--1893), reached from $26$ call sites in the quotient VM (see finding const-3) \\
\revert{MemoryLayoutViolated} & \texttt{0xc9888d23} & 364 (creation frame), 675 (runtime frame) \\
\bottomrule
\end{tabularx}
\caption{Declared errors, their verified selectors, and every raise site in the
fixture.}
\label{tab:const:errors}
\end{table}

\paragraph{Assessment.}
The taxonomy is well chosen for its stated purpose. The important split is
between \revert{PrecompileFailed} (the chain could not execute the operation)
and \revert{ProofRejected} / \revert{BadPointEncoding} (the operation executed
and the input lost). That distinction is preserved end to end: the accumulator
helper at lines 1246--1352 threads a separate \yul{precompile\_failed} return
value (set at lines 1290 and 1346) precisely so a \opcode{STATICCALL} that could
not run is not reported as a malformed public input, and the caller discriminates
at lines 1448--1453; the Lagrange block does the same with
\yul{lagrange\_precompile\_failed} (hoisted at line 1798, set inside
\yul{batch\_invert} at lines 875 and 926, discriminated at lines 1877--1883). The cost is one extra stack slot and
one extra \yul{if} per call site --- cheap, and it is the difference between
``your chain is misconfigured'' and ``your prover is broken'' during an
incident.

Writing the selector at memory \texttt{0x00} is safe: \texttt{0x00}--\texttt{0x40}
is Solidity's designated scratch space, and the generated layout starts at
\texttt{0x1000}, so \yul{fail} can never clobber live verifier state. Every
\yul{fail} call is also terminal (\opcode{REVERT}), so no aliasing question
arises.

\begin{finding}[const-2]{Constructor VK check reverts \texttt{Error(string)}, outside the typed taxonomy --- Observation}
\texttt{Halo2Verifier.sol}:580--584. The comment at lines 53--56 states that the
\emph{only} constructor exception to bare reverts is the memory-layout guard.
In fact the constructor's verifying-key identity check is
\yul{require(..., "invalid vk")}, which encodes as
\texttt{Error(string)} (\texttt{0x08c379a0}) rather than \revert{VkMismatch}.
The runtime path for the same condition (line 1389) does use
\revert{VkMismatch}. No security impact: the check itself is correct and
deployment-time only, and a string revert is strictly more informative than a
bare one. It is a documentation/consistency gap, noted because the header
enumerates the exceptions and this one is missing.
\end{finding}

\begin{finding}[const-3]{\texttt{ERR\_QUOTIENT\_PROGRAM\_INVALID} has no reader; the selector is duplicated as a literal --- Observation}
\texttt{Halo2Verifier.sol}:326 declares
\yul{ERR\_QUOTIENT\_PROGRAM\_INVALID = 0x3cc81b89}, and no code in the fixture
references that identifier. The $26$ structural guards in the quotient VM (an
operand-window guard at 2216, 2253, 2261, 2308, 2344, 2357, 2364, 2384, 2385,
2411, 2426, 2427; stack overflow at 2218, 2332; stack underflow/balance at
2229, 2464, 2580, 2772, 3105; opcode dispatch at 3051, 3091; selector-index and
gap clamps at 3066, 3067, 3068; and program-termination checks at 3098, 3099)
all call \yul{q\_program\_fail()}, which hardcodes the literal:
\begin{quote}\ttfamily
function q\_program\_fail() \{ mstore(0x00, shl(224, 0x3cc81b89)) revert(0x00, 0x04) \}
\end{quote}
(lines 1890--1893). The reason is legitimate and documented at lines 1885--1889:
the quotient VM renders into both this contract and the standalone
\texttt{Halo2QuotientEvaluator}, where the \yul{internal constant} is not in
scope. Both copies are pinned by the same generator test
(\texttt{src/lowering/tests.rs}:2900--2952, which asserts the literal
\texttt{shl(224, 0x3cc81b89)} appears in
\texttt{templates/partials/quotient\_numerator/QuotientHelpers.yul}). So the
duplication is checked, not accidental. It is recorded here only because the
selector block advertises itself as the single source of truth and has one
exception.
\end{finding}

\subsection{Verifying-key pinning constants}
\label{sec:const:vkpin}

\begin{lstlisting}[style=solidity,caption={\texttt{Halo2Verifier.sol} lines 86--95: VK pinning},label={lst:const:vkpin}]
    /// @notice Verifying-key contract address authorized for this verifier.
    /// @dev The runtime length and codehash are pinned by generated constants and checked at construction time.
    address public immutable AUTHORIZED_VK;
    // Expected VK runtime metadata. The deployed VK runtime is
    // INVALID || payload, hence EXPECTED_VK_LENGTH is one byte longer than
    // EXPECTED_VK_PAYLOAD_LENGTH.
    uint256 internal constant EXPECTED_VK_PAYLOAD_LENGTH = 17024;
    uint256 internal constant EXPECTED_VK_LENGTH = 17025;
    uint256 internal constant EXPECTED_VK_CODEHASH_WORD = 0xe68d89362065c8b7107055774f1e045b69bc3bffde4704541c5bc5c91c94cf52;
    bytes32 internal constant EXPECTED_VK_CODEHASH = bytes32(EXPECTED_VK_CODEHASH_WORD);
\end{lstlisting}

The relations are exact:
\[
  \yul{EXPECTED\_VK\_LENGTH} = \yul{EXPECTED\_VK\_PAYLOAD\_LENGTH} + 1,
  \qquad
  \yul{EXPECTED\_VK\_PAYLOAD\_LENGTH} = \texttt{0x4280} = 17024 .
\]
The extra byte is a leading \opcode{INVALID} (\texttt{0xfe}) that makes the VK
runtime non-callable; the loader skips it with
\yul{extcodecopy(vk, VK\_MPTR, 0x01, EXPECTED\_VK\_PAYLOAD\_LENGTH)} (line 1393).
The payload therefore occupies memory $[\texttt{0x3680}, \texttt{0x7900})$,
which is exactly $\mptr{VK\_MPTR} + \texttt{0x4280} = \mptr{CHALLENGE\_MPTR}$.
The two constant blocks (VK length, and the map at lines 128--144) agree by
construction, and this is the strongest internal consistency check available in
lines 1--350.

\yul{EXPECTED\_VK\_CODEHASH} is a \yul{bytes32} restatement of
\yul{EXPECTED\_VK\_CODEHASH\_WORD}; the Solidity-level constructor check (line
582) uses the former and the Yul loader (line 1388) the latter, so both are
live. The VK is validated \emph{twice}: once at construction
(\opcode{EXTCODESIZE} plus \opcode{EXTCODEHASH}, lines 580--584) and again at the
top of every \verb|verifyProof| call (lines 1386--1389), followed by a header
cross-check (lines 1400--1406) that pins \emph{six} loaded header words against
generated literals and folds them into \yul{success}, enforced at line 1406:
\[
  \begin{array}{llll}
    \yul{num\_instances} = 19 & \text{(1400)} &
    \yul{acc\_offset} = 11 & \text{(1403)} \\
    k = 20 & \text{(1401)} &
    \yul{num\_acc\_limbs} = 7 & \text{(1404)} \\
    \yul{has\_accumulator} = 1 & \text{(1402)} &
    \yul{num\_acc\_limb\_bits} = 56 & \text{(1405)}
  \end{array}
\]
The $k = 20$ check is the one with the widest blast radius: the Lagrange block
hard-codes \yul{let k := 20} (line 1800) and squares $x$ that many times to form
$x^{n}$ with $n = 2^{20}$, so a VK header carrying a different $k$ would
silently evaluate the wrong domain. Note that these six are the \emph{only}
header words cross-checked. $\mptr{N\_INV\_MPTR}$, $\mptr{OMEGA\_MPTR}$,
$\mptr{OMEGA\_INV\_MPTR}$, $\mptr{OMEGA\_INV\_TO\_L\_MPTR}$,
$\mptr{G1\_BASE\_MPTR}$, $\mptr{G2\_BASE\_MPTR}$ and
$\mptr{NEG\_S\_G2\_BASE\_MPTR}$ are consumed exactly as loaded, with no on-chain
consistency test (for instance nothing verifies
$\omega \cdot \omega^{-1} = 1$ or that $\omega$ has order $n$). That is sound
only because the entire payload is pinned by \opcode{EXTCODEHASH} one line
earlier: the cross-check exists to catch \emph{generator} drift, not a
substituted VK.

\paragraph{Assessment (clean).}
Re-checking \opcode{EXTCODEHASH} on every proof rather than only at
construction costs $100$ gas warm and defends against the one attack the
immutable alone does not: a VK address that is a proxy, or an address whose code
was placed by \opcode{CREATE2} and later \opcode{SELFDESTRUCT}ed and re-created
(pre-Cancun semantics) or re-deployed at the same address. Because the codehash
is checked immediately before \opcode{EXTCODECOPY} in the same frame, there is
no window in which the code can change between check and use. This region is
sound.

\subsection{ABI calldata cursors}
\label{sec:const:abi}

\begin{lstlisting}[style=solidity,caption={\texttt{Halo2Verifier.sol} lines 97--103: calldata cursors},label={lst:const:abi}]
    // Solidity ABI calldata cursors. The generated verifier accepts exactly
    // verifyProof(bytes proof, uint256[] instances), then parses the `proof`
    // bytes itself in the same order as the Rust verifier transcript.
    uint256 internal constant    PROOF_LEN_CPTR = 0x44;
    uint256 internal constant        PROOF_CPTR = 0x64;
    uint256 internal constant NUM_INSTANCE_CPTR = 0x1ec4;
    uint256 internal constant     INSTANCE_CPTR = 0x1ee4;
\end{lstlisting}

The ABI signature is
\verb|verifyProof(bytes calldata proof, uint256[] calldata instances)|. The
generated calldata frame is fully rigid:

\begin{table}[htbp]
\centering
\begin{tabular}{@{}llll@{}}
\toprule
Offset & Constant & Contents & Pinned by \\
\midrule
\texttt{0x00} & --                     & 4-byte selector & dispatcher \\
\texttt{0x04} & --                     & head of \texttt{proof} $= \texttt{0x40}$ & line 636 \\
\texttt{0x24} & --                     & head of \texttt{instances} $= \texttt{0x1ec0}$ & line 636 \\
\texttt{0x44} & \mptr{PROOF\_LEN\_CPTR} & \texttt{proof.length} $= \texttt{0x1e60}$ & line 1419 \\
\texttt{0x64} & \mptr{PROOF\_CPTR}      & first proof byte & line 1545 \\
\texttt{0x1ec4} & \mptr{NUM\_INSTANCE\_CPTR} & \texttt{instances.length} $= 19$ & line 1420 \\
\texttt{0x1ee4} & \mptr{INSTANCE\_CPTR}  & first instance word & line 1510 \\
\texttt{0x2144} & --                     & end of calldata & line 1426 \\
\bottomrule
\end{tabular}
\caption{Generated calldata frame. Every field is a compile-time constant.}
\label{tab:const:abi}
\end{table}

The four cursor constants satisfy three identities that must hold for the frame
to be self-consistent, and all three do:
\begin{align}
  \mptr{PROOF\_CPTR} &= \mptr{PROOF\_LEN\_CPTR} + \texttt{0x20}
    & \texttt{0x64} &= \texttt{0x44} + \texttt{0x20} \\
  \mptr{NUM\_INSTANCE\_CPTR} &= \mptr{PROOF\_CPTR} + \texttt{0x1e60}
    & \texttt{0x1ec4} &= \texttt{0x64} + \texttt{0x1e60} \\
  \mptr{INSTANCE\_CPTR} &= \mptr{NUM\_INSTANCE\_CPTR} + \texttt{0x20}
    & \texttt{0x1ee4} &= \texttt{0x1ec4} + \texttt{0x20}
\end{align}
and the exact size check closes the frame:
$\opcode{CALLDATASIZE} = \mptr{INSTANCE\_CPTR} + 19\cdot\texttt{0x20}
 = \texttt{0x1ee4} + \texttt{0x260} = \texttt{0x2144}$.
The ABI head for \texttt{instances} is checked against
$\mptr{NUM\_INSTANCE\_CPTR} - \texttt{0x04} = \texttt{0x1ec0}$, i.e.\ relative to
the post-selector origin, as the ABI requires.

Finally, after the transcript parser has consumed the whole proof stream it
asserts that its running cursor landed exactly on \mptr{NUM\_INSTANCE\_CPTR}
(line 1785), which is an independent confirmation that the parser's generated
step sizes sum to \texttt{0x1e60}.

\paragraph{Assessment.}
Pinning \opcode{CALLDATASIZE} exactly (rather than $\ge$) is the right call for
a verifier whose parser is a sequence of fixed offsets: it removes any
possibility of trailing-byte smuggling and makes the parser's cursor arithmetic
checkable in one comparison. The cost is documented at lines 606--611 and is
real: ERC-2771 forwarders, multicall wrappers and paymaster contexts append
calldata and therefore cannot call this contract directly. That is an
integration constraint, not a defect, and the NatSpec states it prominently.
This region is clean.

\subsection{The generated memory map}
\label{sec:const:map}

\subsubsection{Design: absolute addressing and phase reuse}

The verifier abandons Solidity's free-memory pointer entirely. Every address is
a generated literal, and the region below $\texttt{0x1000}$ is left untouched so
that Solidity's scratch words (\texttt{0x00}--\texttt{0x40}), the free-memory
pointer (\texttt{0x40}), the zero slot (\texttt{0x60}) and solc's via-IR spill
reservation (observed up to \texttt{0x8e0}) all remain valid. What this buys is
that no address in the hot loops has to be computed, loaded, or bounds-checked
at runtime; what it costs is that correctness of the layout is a
\emph{generation-time} property, defended on chain only by the single
free-memory-pointer guard of \S\ref{sec:const:pragma}.

The generator models the layout as an arena of regions, each tagged with a
lifetime that is either \texttt{Permanent}, a single \texttt{MemoryPhase}, or a
\texttt{PhaseSpan}. Two regions may share an address if and only if their
lifetimes do not intersect. The declared phase order
(\texttt{MemoryPhase::IN\_TEMPLATE\_ORDER},
\texttt{src/lowering/layout/memory.rs}:168--185) has sixteen entries:

\begin{center}
\texttt{VkConstructorPayload} $\to$ \texttt{ConstructorSmoke} $\to$
\texttt{AccumulatorMsm} $\to$ \texttt{Transcript} $\to$
\texttt{LagrangeBatchInvert} $\to$ \texttt{QuotientVm} $\to$
\texttt{LinearizationTrace} $\to$ \texttt{PcsQEvalSourceTable} $\to$
\texttt{PcsQComTrace} $\to$ \texttt{ScalarInv} $\to$ \texttt{PcsFinalMsm}
$\to$ \texttt{PcsPairing} $\to$ \texttt{AccumulatorPairingBatch} $\to$
\texttt{FinalPairing} $\to$ \texttt{VerifierReturn} $\to$
\texttt{QuotientReturn}
\end{center}

Each rendered section carries a \yul{// \_\_phase:} marker comment naming its
phase, so the template's execution order is checkable against this array on
every render. Two caveats a reader should carry forward. First, the source
comment at lines 110--111 gives the chain as
``\texttt{Transcript} $\to$ \texttt{QuotientReturn} $\to$ \texttt{PcsPairing}
$\to$ \texttt{VerifierReturn}'', which is \emph{not} the declared order ---
\texttt{QuotientReturn} sorts last, after \texttt{VerifierReturn} --- and the
comment omits \texttt{AccumulatorPairingBatch}, which also lives at
\texttt{0x1000} in this render. Second, and more importantly, the arena's
notion of ``both live'' is the \emph{declared} lifetime, not a computed
liveness: \yul{MemoryLifetime::intersects} returns false for any two distinct
single-phase lifetimes (\texttt{src/lowering/layout/memory.rs}:230--249),
whatever their byte ranges. Regions whose value really must survive across
phases have to be declared \texttt{PhaseSpan} by hand for the check to bite.

The comment at lines 104--118 states this contract to the reader in one
sentence: ``seeing the same address on several constants here means \emph{reused
in sequence}, never \emph{written twice at once}''.

\begin{lstlisting}[style=solidity,caption={\texttt{Halo2Verifier.sol} lines 104--120: the deliberate phase-reuse base},label={lst:const:reuse}]
    // First general-purpose memory words reserved by the generated verifier.
    //
    // TRANSCRIPT_MPTR, RETURN_MPTR and QUOTIENT_RETURN_MPTR below share one
    // base address, and so do the constructor smoke-test and PCS pairing
    // scratch windows that the assembly addresses inline. That is deliberate,
    // not a collision: each is a distinct `MemoryPhase` whose lifetime ends
    // before the next begins (Transcript -> QuotientReturn -> PcsPairing ->
    // VerifierReturn, with ConstructorSmoke confined to deployment). The
    // generator allocates them from one arena that rejects any two regions
    // sharing an address while both are live, so the reuse is checked at
    // generation time rather than argued in a comment.
    //
    // Practical consequence for a reader: seeing the same address on several
    // constants here means "reused in sequence", never "written twice at
    // once". RETURN_MPTR is a single word set to 1 on success.
    uint256 internal constant    TRANSCRIPT_MPTR = 0x1000;
    uint256 internal constant        RETURN_MPTR = 0x1000;
\end{lstlisting}

\subsubsection{The \texttt{0x1000} band}
\label{sec:const:band1000}

Six regions share the base $\texttt{0x1000}$. Three are named by constants
(\mptr{TRANSCRIPT\_MPTR}, line 119; \mptr{RETURN\_MPTR}, line 120;
\mptr{QUOTIENT\_RETURN\_MPTR}, line 230) and three are addressed by bare
literals in assembly. Their extents, and the phase in which each is live:

\begin{table}[htbp]
\centering
\begin{tabularx}{\textwidth}{@{}lllX@{}}
\toprule
Region & Extent & Phase & Evidence \\
\midrule
constructor smoke scratch & $[\texttt{0x1000},\texttt{0x1320})$ & \texttt{ConstructorSmoke} & \yul{scratch := 0x1000} line 369; zeroing loops over \texttt{0x300} bytes at lines 411 and 537; the pairing probes write their result word at \texttt{scratch+0x300} (lines 525, 532), one word beyond the region's registered \texttt{0x300} length --- see finding const-12 \\
transcript buffer & $[\texttt{0x1000},\,\cdot\,)$, ceiling \texttt{0x3580} & \texttt{Transcript} & \yul{buf\_len := TRANSCRIPT\_MPTR} line 735; \yul{squeeze\_to} rehashes and resets \yul{buf\_len} to $\mptr{TRANSCRIPT\_MPTR} + 32$, lines 801--809, so the high-water mark is the longest inter-squeeze absorb run, every length of which is a literal \\
\mptr{QUOTIENT\_RETURN\_MPTR} & $[\texttt{0x1000},\texttt{0x1180})$ & \texttt{QuotientReturn} & planner size $(2+10)$ words (flag, value, $10$ selector accumulators); \emph{not written in this single-contract render} \\
PCS pairing scratch & $[\texttt{0x1000},\texttt{0x1120})$ obs. & \texttt{PcsPairing} & \yul{\_\_phase:pcs\_pairing} block, lines 4008--4033; touches \texttt{0x1000}, \texttt{0x1080}, \texttt{0x1100}. Planner reserves $32$ words ($\texttt{PCS\_STATIC\_WORKING\_WORDS}$) \\
accumulator pairing batch & $[\texttt{0x1000},\texttt{0x1240})$ & \texttt{AccumulatorPairingBatch} & \yul{batch\_ptr := 0x1000} line 4052; \yul{keccak256(batch\_ptr, 0x0240)} line 4066; MSM/ADD staging reuses $[\texttt{0x1000},\texttt{0x1100})$, lines 4072--4096 \\
final pairing scratch & $[\texttt{0x1240},\texttt{0x1540})$ touched, $[\texttt{0x1240},\texttt{0x1560})$ reserved & \texttt{FinalPairing} & \yul{let scratch := 0x1240} line 969; the \texttt{0x300}-byte two-pair input is at \texttt{scratch} and the \texttt{0x20} result is written back \emph{over} \texttt{scratch}, not after it (line 980), so only $24$ words are touched; the planner reserves $25$ ($\texttt{PAIRING\_STATIC\_WORKING\_WORDS}$) \\
\mptr{RETURN\_MPTR} & $[\texttt{0x1000},\texttt{0x1020})$ & \texttt{VerifierReturn} & \yul{mstore(RETURN\_MPTR, 1); return(RETURN\_MPTR, 0x20)} lines 4127--4128 \\
\bottomrule
\end{tabularx}
\caption{The \texttt{0x1000} band: six regions based at \texttt{0x1000} plus
the final-pairing frame immediately above them, with the phase that owns each.}
\label{tab:const:band1000}
\end{table}

The ordering argument is verifiable from source order in the fixture: the
transcript is complete at line 1785 (cursor assertion); the PCS pairing block
runs at lines 4011--4033, the accumulator pairing batch at 4050--4097, and the
final pairing at 4112--4114; \mptr{RETURN\_MPTR} is written only on the terminal
success path at line 4127. Each of these blocks carries an explicit
\yul{\_\_phase:} marker in the source. \mptr{QUOTIENT\_RETURN\_MPTR} is rendered
for the split-evaluator configuration and is not written in this
single-contract render; the address is nonetheless reserved so that a
split render places its return frame here.

Note the important detail in the accumulator batch: the randomizer is
$\alpha = \mathrm{keccak256}(\text{tag} \Vert \yul{vk\_digest} \Vert
\mathsf{PairingRhs} \Vert \mathsf{PairingLhs} \Vert \mathsf{AccRhs}
\Vert \mathsf{AccLhs}) \bmod r$ over exactly \texttt{0x240} bytes at
\texttt{0x1000} (line 4066; the tag is the literal
the $21$ ASCII bytes \texttt{"pairing-batch-acc-kzg"} written at line 4057 as
the word \texttt{0x70616972696e672d62617463682d6163632d6b7a67} followed by $8$
zero bytes, i.e.\ occupying bytes $3$--$23$ of the word), with a zero draw
replaced by $1$ at line
4067. The same bytes at \texttt{0x1000} are then overwritten as G1MSM input
(line 4072). The hash is consumed before the overwrite, so the reuse within one
phase is also ordered, not simultaneous.

\begin{finding}[const-4]{Six regions alias \texttt{0x1000}; separation rests entirely on generation-time analysis --- Observation}
\texttt{Halo2Verifier.sol}:104--120, 230. The aliasing is deliberate,
documented, and correct for this artifact: no two of the regions in
Table~\ref{tab:const:band1000} are live simultaneously, and the transcript ---
the longest-lived of them --- is provably finished before the first pairing
scratch write. There is no runtime check that any of these regions stays inside
its budget; in particular \yul{common\_*} advances \yul{buf\_len} by generated
constants with no upper-bound comparison against $\texttt{0x3580}$ (the next
allocated region, \S\ref{sec:const:vkimage}).

This is sound in the deployed artifact for a specific reason worth stating:
\opcode{CALLDATASIZE} is pinned (\S\ref{sec:const:abi}), every absorb length is
a generated literal, the parse loops have literal bounds, and the transcript
resets to \texttt{0x1000} at each squeeze --- so \yul{buf\_len} follows a single
deterministic trajectory independent of proof content, and its maximum is a
build-time property, not an input-dependent one. No proof can steer it.

What deserves qualification is the header comment's justification, that ``the
generator allocates them from one arena that rejects any two regions sharing an
address while both are live, so the reuse is checked at generation time rather
than argued in a comment.'' The arena's test is
\yul{MemoryLifetime::intersects}, and for two regions declared
\texttt{Phase(a)} and \texttt{Phase(b)} it returns \yul{false} whenever
$a \neq b$; the arena compares byte ranges only for pairs whose lifetimes
\emph{do} intersect, so two distinct single-phase regions are never range-checked
against each other at all, and no real liveness is computed. Everything
therefore rests on each region's \emph{declared} lifetime being correct. The
generator itself records one place where this is not enough:
\texttt{src/lowering/layout/mod.rs}:111--120 notes that the final-pairing
scratch and the accumulator pairing-batch hash frame ``carry different
\texttt{MemoryPhase}s, and \texttt{MemoryLifetime::intersects} treats distinct
phases as never co-live, so the planner cannot catch that overlap -- it has to
be avoided by construction here.'' The reuse is thus checked at generation time
in part and argued by construction in part; the comment states only the first
half. The observation is recorded because none of this is visible from lines
1--350 alone, and because it does not survive a change to the absorb schedule
or to a phase boundary.
\end{finding}

\begin{finding}[const-12]{The constructor smoke probe writes one word past its planner-registered region --- Observation}
\texttt{Halo2Verifier.sol}:369, 525, 532. The constructor scratch is
\yul{let scratch := 0x1000} (line 369) and the generator registers it as
\yul{constructor\_smoke\_scratch} with length
$\texttt{PAIRING\_TWO\_PAIR\_BYTES} = 2\,(\texttt{0x80} + \texttt{0x100})
= \texttt{0x300}$ (\texttt{src/lowering/layout/memory.rs}:739--744,
\texttt{src/lowering/layout/mod.rs}:94--97), i.e.\
$[\texttt{0x1000}, \texttt{0x1300})$. Both pairing probes, however, direct the
precompile's $\texttt{0x20}$-byte output to \yul{add(scratch, 0x300)}:
\begin{quote}\ttfamily
if iszero(staticcall(PAIRING\_GAS\_2PAIR, 0x0f, scratch, 0x0300, add(scratch, 0x300), 0x20)) \{ revert(0, 0) \}
\end{quote}
(lines 525 and 532; the result word is then read at lines 527 and 534). The
touched extent is therefore $[\texttt{0x1000}, \texttt{0x1320})$ --- one word
outside the range the arena knows about, so that word is excluded from the
overlap analysis the header comment at lines 104--114 cites as the reason the
aliasing is safe.

No impact in this artifact. The write happens only in the creation frame, where
the sole other live region is the $\texttt{0x30c0}$-byte G1MSM probe buffer at
\texttt{0xb140} (line 549), and constructor memory is discarded. It is recorded
because it is a concrete instance of the general point in finding const-4: the
arena checks declared region extents, and a template that writes outside its
declared extent is not something the check can see. Contrast the runtime
final-pairing frame, which \emph{is} registered with the extra return word
($\texttt{PAIRING\_STATIC\_WORKING\_WORDS} = \texttt{0x300}/32 + 1 = 25$) and
then does not use it, because \yul{ec\_pairing} writes its result back over the
input at \texttt{scratch} (line 980).
\end{finding}

\subsubsection{The verifying-key image}
\label{sec:const:vkimage}

\begin{lstlisting}[style=solidity,caption={\texttt{Halo2Verifier.sol} lines 122--144: VK header slots},label={lst:const:vkhdr}]
    // ----------------------------------------------------------------------
    // Verifying-key memory map. The VK header lives at VK_MPTR, followed
    // by the quotient VM payload and commitments. After the full VK
    // runtime comes the challenge slots (challenge_mptr..) and the
    // per-stage scratch (theta_mptr..).
    // ----------------------------------------------------------------------
    uint256 internal constant                VK_MPTR = 0x3680;
    uint256 internal constant         VK_DIGEST_MPTR = 0x3680;
    uint256 internal constant     NUM_INSTANCES_MPTR = 0x36a0;
    uint256 internal constant                 K_MPTR = 0x36c0;
    uint256 internal constant             N_INV_MPTR = 0x36e0;
    uint256 internal constant             OMEGA_MPTR = 0x3700;
    uint256 internal constant         OMEGA_INV_MPTR = 0x3720;
    uint256 internal constant    OMEGA_INV_TO_L_MPTR = 0x3740;
    uint256 internal constant   HAS_ACCUMULATOR_MPTR = 0x3760;
    uint256 internal constant        ACC_OFFSET_MPTR = 0x3780;
    uint256 internal constant     NUM_ACC_LIMBS_MPTR = 0x37a0;
    uint256 internal constant NUM_ACC_LIMB_BITS_MPTR = 0x37c0;
    uint256 internal constant            G1_BASE_MPTR = 0x37e0;
    uint256 internal constant            G2_BASE_MPTR = 0x3860;
    uint256 internal constant      NEG_S_G2_BASE_MPTR = 0x3960;

    uint256 internal constant CHALLENGE_MPTR = 0x7900;
\end{lstlisting}

The header is $31$ words: eleven scalars, one padded $\GO$ point
($\mptr{G1\_BASE\_MPTR}$, $4$ words), and two padded $\GT$ points
($\mptr{G2\_BASE\_MPTR}$ and $\mptr{NEG\_S\_G2\_BASE\_MPTR}$, $8$ words each).
Header bytes: $31 \cdot 32 = \texttt{0x3e0}$, ending at $\texttt{0x3a60}$.

Everything after the header is addressed by bare literals in the assembly, not
by named constants. The full VK image decomposes as:

\begin{table}[htbp]
\centering
\begin{tabular}{@{}lllll@{}}
\toprule
Sub-region & Start & End & Size & Cited at \\
\midrule
VK header ($31$ words)          & \texttt{0x3680} & \texttt{0x3a60} & \texttt{0x03e0} & lines 128--142 \\
Quotient $\Fr$ constant pool ($178$ words) & \texttt{0x3a60} & \texttt{0x50a0} & \texttt{0x1640} & \yul{q\_const\_mptr}, line 2013 \\
Quotient program bytes          & \texttt{0x50a0} & \texttt{0x626f} & \texttt{0x11cf} & \yul{q\_program\_mptr} line 2016, \yul{q\_end} line 2170 \\
Word-alignment padding          & \texttt{0x626f} & \texttt{0x6280} & $17$ bytes & implied \\
Fixed VK commitments ($45\times\GO$) & \texttt{0x6280} & \texttt{0x7900} & \texttt{0x1680} & lines 3848--3917 \\
\midrule
Total                           & \texttt{0x3680} & \texttt{0x7900} & \texttt{0x4280} & $= 17024$ \\
\bottomrule
\end{tabular}
\caption{Decomposition of the $17{,}024$-byte VK payload in memory. The sizes
sum \emph{exactly} to \yul{EXPECTED\_VK\_PAYLOAD\_LENGTH}.}
\label{tab:const:vkimage}
\end{table}

This closure is a strong check: it independently confirms the $178$-word
constant pool, the $\texttt{0x11cf}$-byte program, and the $45$ fixed
commitments, none of which is named by a constant.

Immediately below \mptr{VK\_MPTR} sits the scalar-inversion modexp scratch, a
$\texttt{0x100}$-byte window at $\texttt{0x3580}$ (\yul{let p := 0x3580}, line
710). Only the first $\texttt{0xc0}$ bytes are the EIP-198 frame; the remaining
$\texttt{0x40}$ is deliberate padding kept below \mptr{VK\_MPTR} so no later
allocation can claim it. Its comment (lines 695--699) explains the placement: it
reuses the dead transcript buffer rather than a fixed post-VK address that would
collide with live PCS scratch when the VK shrinks.

\subsubsection{Challenge window and the theta slots}
\label{sec:const:theta}

\begin{lstlisting}[style=solidity,caption={\texttt{Halo2Verifier.sol} lines 146--187: challenge and theta-relative slots (two comment blocks and the four 4-word G1 slot declarations elided at the marked points)},label={lst:const:theta}]
    // Challenge layout. Squeeze order in midnight-proofs:
    //   user_phase challenges (variable count)
    //   theta -> beta, gamma -> trash_challenge -> y -> x ->
    //   x1, x2 -> x3 -> x4
    uint256 internal constant            THETA_MPTR = 0x7900;
    uint256 internal constant             BETA_MPTR = 0x7920;
    uint256 internal constant            GAMMA_MPTR = 0x7940;
    uint256 internal constant TRASH_CHALLENGE_MPTR = 0x7960;
    uint256 internal constant                Y_MPTR = 0x7980;
    uint256 internal constant                X_MPTR = 0x79a0;
    uint256 internal constant               X1_MPTR = 0x79c0;
    uint256 internal constant               X2_MPTR = 0x79e0;
    uint256 internal constant               X3_MPTR = 0x7a00;
    uint256 internal constant               X4_MPTR = 0x7a20;
    // ... F_COM_MPTR, PI_MPTR, ACC_LHS_MPTR, ACC_RHS_MPTR: 4-word G1 slots
    uint256 internal constant              X_N_MPTR = 0x7c40;
    uint256 internal constant  X_N_MINUS_1_INV_MPTR = 0x7c60;
    uint256 internal constant           L_LAST_MPTR = 0x7c80;
    uint256 internal constant          L_BLIND_MPTR = 0x7ca0;
    uint256 internal constant              L_0_MPTR = 0x7cc0;
    uint256 internal constant     INSTANCE_EVAL_MPTR = 0x7ce0;
    // ... 5-line comment on LINEARIZATION_EVAL_MPTR (lines 176-180)
    uint256 internal constant LINEARIZATION_EVAL_MPTR = 0x7d00;
    uint256 internal constant         QUOTIENT_MPTR = 0x7d20;   // 4 words
    uint256 internal constant            F_EVAL_MPTR = 0x7dc0;
    uint256 internal constant                 V_MPTR = 0x7de0;
    uint256 internal constant         FINAL_COM_MPTR = 0x7e00;   // 4 words
    uint256 internal constant      PAIRING_LHS_MPTR = 0x7e80;   // 4 words
    uint256 internal constant      PAIRING_RHS_MPTR = 0x7f00;   // 4 words
\end{lstlisting}

(The four $\GO$ slots elided above are
$\mptr{F\_COM\_MPTR} = \texttt{0x7a40}$, $\mptr{PI\_MPTR} = \texttt{0x7ac0}$,
$\mptr{ACC\_LHS\_MPTR} = \texttt{0x7b40}$,
$\mptr{ACC\_RHS\_MPTR} = \texttt{0x7bc0}$; lines 162--167.)

Every one of these is $\mptr{THETA\_MPTR} + 32w$ for a fixed word index $w$, and
the word indices are a fixed historical schedule. Writing $T = \texttt{0x7900}$:

\begin{table}[htbp]
\centering
\begin{tabular}{@{}llll@{}}
\toprule
$w$ & Address & Constant & Width \\
\midrule
$0$--$9$   & \texttt{0x7900}--\texttt{0x7a20} & $\theta,\beta,\gamma,\tau,y,x,x_1,x_2,x_3,x_4$ & $1$ word each \\
$10$       & \texttt{0x7a40} & \mptr{F\_COM\_MPTR}            & $4$ \\
$14$       & \texttt{0x7ac0} & \mptr{PI\_MPTR}                & $4$ \\
$18$       & \texttt{0x7b40} & \mptr{ACC\_LHS\_MPTR}          & $4$ \\
$22$       & \texttt{0x7bc0} & \mptr{ACC\_RHS\_MPTR}          & $4$ \\
$26$       & \texttt{0x7c40} & \mptr{X\_N\_MPTR}              & $1$ \\
$27$       & \texttt{0x7c60} & \mptr{X\_N\_MINUS\_1\_INV\_MPTR} & $1$ \\
$28$--$30$ & \texttt{0x7c80}--\texttt{0x7cc0} & $L_{\mathrm{last}}, L_{\mathrm{blind}}, L_0$ & $1$ each \\
$31$       & \texttt{0x7ce0} & \mptr{INSTANCE\_EVAL\_MPTR}    & $1$ \\
$32$       & \texttt{0x7d00} & \mptr{LINEARIZATION\_EVAL\_MPTR} & $1$ \\
$33$--$36$ & \texttt{0x7d20} & \mptr{QUOTIENT\_MPTR}          & $4$ \\
$37$       & \texttt{0x7da0} & \emph{unnamed padding}         & $1$ \\
$38$       & \texttt{0x7dc0} & \mptr{F\_EVAL\_MPTR}           & $1$ \\
$39$       & \texttt{0x7de0} & \mptr{V\_MPTR}                 & $1$ \\
$40$--$43$ & \texttt{0x7e00} & \mptr{FINAL\_COM\_MPTR}        & $4$ \\
$44$--$47$ & \texttt{0x7e80} & \mptr{PAIRING\_LHS\_MPTR}      & $4$ \\
$48$--$51$ & \texttt{0x7f00} & \mptr{PAIRING\_RHS\_MPTR}      & $4$ \\
$52$       & \texttt{0x7f80} & \mptr{ROT\_POINTS\_MPTR}       & window \\
\bottomrule
\end{tabular}
\caption{Theta-relative word schedule. Word $37$ is a hole (see
\S\ref{sec:const:gaps}).}
\label{tab:const:theta}
\end{table}

\paragraph{\mptr{LINEARIZATION\_EVAL\_MPTR}.}
The comment at lines 176--180 is a substantive claim about the protocol, not a
naming note, and it should be read carefully. The slot holds the \emph{expected
opening scalar} of the linearised commitment: the negated $y$-batched identity
numerator reconstructed from the alleged evaluations at $x$. It is
\emph{deliberately not} $h(x)$ --- the proof carries no $h(x)$ scalar the
verifier could trust, and the commitment side is rebuilt from the four quotient
limb commitments instead. The scalar is stored at line 3359 and consumed by the
PCS $q$-eval source table (line 3505). This is specified fully in the quotient
and linearisation section.

\begin{finding}[const-5]{\mptr{CHALLENGE\_MPTR} and \mptr{THETA\_MPTR} are the same address --- Observation}
\texttt{Halo2Verifier.sol}:144, 150. Both are $\texttt{0x7900}$. This is not a
bug: the user-phase challenge window is sized by the number of user challenges,
which is zero for this circuit, so the window is empty and the named challenge
slots begin immediately. The reason it is worth stating is the failure mode it
resembles. The transcript parser writes user challenge $j$ to
$\mptr{CHALLENGE\_MPTR} + 32j$, and the very next squeeze targets
\mptr{THETA\_MPTR} unconditionally; if a future circuit had user challenges and
an under-sized window, theta would silently overwrite a user challenge while the
transcript still matched the prover --- a liveness break at best and a soundness
break at worst. The generator asserts
$\mptr{CHALLENGE\_MPTR} + 32\cdot(\text{user challenges}) \le \mptr{THETA\_MPTR}$
(\texttt{src/lowering/layout/memory.rs}:835--841, whose assertion message names
exactly this failure: ``theta's squeeze would silently overwrite a user
challenge''). In
\emph{this} artifact the
equality is safe and is corroborated independently: the VK payload ends exactly
at $\texttt{0x3680} + \texttt{0x4280} = \texttt{0x7900}$, leaving no room for a
challenge window.
\end{finding}

\subsubsection{PCS fixed windows}
\label{sec:const:pcswindows}

\begin{lstlisting}[style=solidity,caption={\texttt{Halo2Verifier.sol} lines 189--240: PCS windows, staging bases and phase scratch},label={lst:const:pcswin}]
    // Multi-prepare scratch (sized at codegen time).
    uint256 internal constant       ROT_POINTS_MPTR = 0x7f80;
    uint256 internal constant       X1_POWERS_MPTR = 0x8300;
    // Per-set q_com points are folded straight into the final MSM rather than
    // materialized, so no q_com window is rendered. The generator still plans
    // and bounds-checks one (`q_com_words`, currently 0); an emitter that
    // starts materializing q_com must add the constant back here.
    uint256 internal constant      Q_EVAL_SET_MPTR = 0x8b20;
    // ... PCS_FINAL_MSM_MPTR comment
    uint256 internal constant     PCS_FINAL_MSM_MPTR = 0xb280;
    // Q_EVAL_CPTR is set at runtime once the verifier reaches the q_evals
    // block of the proof; we keep it as a memory slot for symmetry.
    uint256 internal constant         Q_EVAL_CPTR_MPTR = 0x9220;
    // ... G1_IDENTITY_MPTR comment
    uint256 internal constant       G1_IDENTITY_MPTR = 0x9320;
    // ... REVERSED_EVALS_MPTR comment (Optimisation H3)
    uint256 internal constant     REVERSED_EVALS_MPTR = 0x9480;
    // SELECTOR_ACC_MPTR holds one accumulator per simple selector, live from
    // the quotient VM through the final MSM.
    uint256 internal constant      SELECTOR_ACC_MPTR = 0xb140;
    uint256 internal constant   QUOTIENT_RETURN_MPTR = 0x1000;
    // Shares SELECTOR_ACC_MPTR's base by construction: the Lagrange batch
    // inversion runs to completion before the quotient VM writes the first
    // selector accumulator, so the two never overlap in time. Pinned by
    // `selector_accumulators_are_live_from_quotient_to_final_msm`.
    uint256 internal constant  BATCH_INV_SCRATCH_MPTR = 0xb140;
    // Lagrange batch-inversion input run: denominators, in-place inverses,
    // then Lagrange values, consumed and distilled into the named theta
    // slots by the Lagrange block. Planner-registered phase scratch.
    uint256 internal constant    LAGRANGE_DENOMS_MPTR = 0xb580;
    uint256 internal constant        TRACE_U256_MPTR = 0xe340;
\end{lstlisting}

Five windows sit above the theta slots, each with a fixed \emph{capacity} and a
smaller \emph{used} extent in this render:

\begin{table}[htbp]
\centering
\begin{tabular}{@{}lllll@{}}
\toprule
Window & Base & Capacity & Used & Evidence for the used extent \\
\midrule
\mptr{ROT\_POINTS\_MPTR}  & \texttt{0x7f80} & $28$ w & $4$ w  & max offset \texttt{0x60}, lines 3424--3439 \\
\mptr{X1\_POWERS\_MPTR}   & \texttt{0x8300} & $65$ w & $43$ w & loop bound \texttt{0x2a} plus slot $0$, lines 3447--3453 \\
(\yul{q\_com}, unrendered) & \texttt{0x8b20} & $0$ w  & $0$ w  & lines 192--195 \\
\mptr{Q\_EVAL\_SET\_MPTR} & \texttt{0x8b20} & $56$ w & $11$ w & max offset \texttt{0x140}, line 3637 \\
\mptr{Q\_EVAL\_CPTR\_MPTR}& \texttt{0x9220} & $1$ w  & $1$ w  & written line 1753, read line 3648 \\
\mptr{G1\_IDENTITY\_MPTR} & \texttt{0x9320} & $4$ w  & $0$ w  & never referenced in this render \\
\mptr{REVERSED\_EVALS\_MPTR} & \texttt{0x9480} & $102$ w & $102$ w & \S\ref{sec:const:evals} \\
\bottomrule
\end{tabular}
\caption{PCS fixed windows. ``w'' = 32-byte words.}
\label{tab:const:pcswin}
\end{table}

\paragraph{The zero-width \yul{q\_com} window (lines 192--195).}
The window exists in the planner but renders no constant, because per-set
$q_{\mathrm{com}}$ points are folded straight into the fused final MSM rather
than materialised. Consequently $\mptr{Q\_EVAL\_SET\_MPTR}$ starts at the same
address the $q_{\mathrm{com}}$ window would have. The comment correctly warns
that an emitter which starts materialising $q_{\mathrm{com}}$ must reintroduce
the constant --- otherwise the two windows would genuinely overlap, since both
are \texttt{Permanent} (never phase-scoped) and the arena's phase-disjointness
argument does not apply to them.

\paragraph{\mptr{REVERSED\_EVALS\_MPTR} and the eval count.}
\label{sec:const:evals}
The window holds $102$ decoded proof evaluations
$\mathbf{e} = (e_0,\dots,e_{101})$, spilled one per iteration by the transcript
loop at lines 1697--1715 so that downstream references become $3$-gas
\yul{mload}s instead of calldata reads. Three independent facts pin the count at
$102$:
\begin{enumerate}
  \item the spill loop's literal calldata bound is \texttt{0x0cc0} (line 1699)
        and its step is \texttt{0x20} (line 1713), so it executes
        $\texttt{0x0cc0}/32 = 102$ times and writes $102$ words starting at
        \mptr{REVERSED\_EVALS\_MPTR} (lines 1708--1709);
  \item $\mptr{ADVICE\_COMMS\_MPTR\_BASE} - \mptr{REVERSED\_EVALS\_MPTR}
        = \texttt{0xa140} - \texttt{0x9480} = \texttt{0xcc0} = 102 \cdot 32$
        (the commitment bases are allocated immediately after this window);
  \item the quotient VM's operand bound is
        $\yul{q\_ptr} - \texttt{0x3680} \le \texttt{0x6aa0}$ (lines 2216, 2253,
        2261, 2308, 2344, 2357, 2364, 2411, 2426, 2427), i.e.\
        $\yul{q\_ptr} \le \texttt{0xa120}$, and
        $\texttt{0xa120} = \mptr{REVERSED\_EVALS\_MPTR} + 101\cdot 32$ is
        exactly the last evaluation word.
\end{enumerate}
The VM's operand window is therefore the closed byte range
$[\texttt{0x3680}, \texttt{0xa120}]$ --- the guard is written on the raw literal
\texttt{0x3680}, not on \mptr{VK\_MPTR}, and the subtraction is unsigned, so an
operand below \texttt{0x3680} wraps to a huge value and is rejected by the same
comparison. The strided variants use \texttt{0x69e0} instead (lines 2384--2385)
because those operands are the \emph{base} of a seven-word run read at offsets
up to $6\cdot 32 = \texttt{0xc0}$, and
$\texttt{0x69e0} + \texttt{0xc0} = \texttt{0x6aa0}$, so the effective ceiling is
the same $\texttt{0xa120}$.

That window is a single unsigned comparison, which is its point --- but it is
worth being precise about what it admits. It is \emph{not} ``the VK image plus
the decoded evaluations'': it is every byte between them as well, namely the
challenge/theta band $[\texttt{0x7900},\texttt{0x7f80})$ and the PCS fixed
windows $[\texttt{0x7f80},\texttt{0x9480})$. The bound is therefore loose by
$\texttt{0x1b80}$ bytes relative to the two regions a VM operand is intended to
name. This costs nothing here, because the program bytes are part of the
codehash-pinned VK payload and no proof input reaches \yul{q\_pc}, but it means
the guard bounds \emph{where} an operand can point, not \emph{what} it can
name.

\begin{finding}[const-6]{\mptr{G1\_IDENTITY\_MPTR} and \mptr{TRACE\_U256\_MPTR} have no reader in this render --- Observation}
\texttt{Halo2Verifier.sol}:218 and 240. Neither identifier appears anywhere else
in the fixture. Both are intentional:
\mptr{G1\_IDENTITY\_MPTR} reserves four never-written words so that PCS emitters
staging an identity commitment can \opcode{MCOPY} the all-zero EIP-2537 identity
encoding from a fixed address, and this render's PCS plan never stages one;
\mptr{TRACE\_U256\_MPTR} is the one-word log scratch used only by trace-enabled
builds, and this is a production render. Solidity \yul{internal constant}s that
are never referenced contribute no runtime bytecode, so the cost is zero and the
EIP-170 budget of \S\ref{sec:const:pragma} is unaffected. Recorded so a reader
auditing the map does not go looking for the missing accesses.
\end{finding}

\subsubsection{Per-category commitment bases}
\label{sec:const:comms}

\begin{lstlisting}[style=solidity,caption={\texttt{Halo2Verifier.sol} lines 242--263: commitment category bases and the cumulative-offset rule},label={lst:const:comms}]
    // ----------------------------------------------------------------------
    // Per-category bases for EIP-2537 padded G1 commitments. The proof
    // calldata carries 128-byte uncompressed/padded G1s after the off-chain
    // proof shim repacks midnight-proofs' native compressed stream; this
    // region stores the 4-word slots used by PCS / quotient-fold sections.
    //
    // Cumulative offsets (in words from `comms_mptr_base`):
    //   ADVICE_COMMS_MPTR_BASE          + 0
    //   LOOKUP_M_COMMS_MPTR_BASE        + 4*total_advices
    //   PERM_Z_COMMS_MPTR_BASE          + 4*total_advices + 4*num_lookups
    //   LOOKUP_HELPER_COMMS_MPTR_BASE   + ... + 4*num_permutation_zs
    //   LOOKUP_Z_COMMS_MPTR_BASE        + ... + 4*lookup_helper_chunks_total
    //   TRASHCAN_COMMS_MPTR_BASE        + ... + 4*num_lookups
    //   QUOTIENT_LIMB_COMMS_MPTR_BASE   + ... + 4*num_trashcans
    // ----------------------------------------------------------------------
    uint256 internal constant         ADVICE_COMMS_MPTR_BASE = 0xa140;
    uint256 internal constant       LOOKUP_M_COMMS_MPTR_BASE = 0xa8c0;
    uint256 internal constant         PERM_Z_COMMS_MPTR_BASE = 0xa9c0;
    uint256 internal constant  LOOKUP_HELPER_COMMS_MPTR_BASE = 0xacc0;
    uint256 internal constant       LOOKUP_Z_COMMS_MPTR_BASE = 0xadc0;
    uint256 internal constant     TRASHCAN_COMMS_MPTR_BASE = 0xaec0;
    uint256 internal constant QUOTIENT_LIMB_COMMS_MPTR_BASE = 0xaf40;
\end{lstlisting}

Let $B = \mptr{ADVICE\_COMMS\_MPTR\_BASE}$ and let $W = 4$ be the word width of
a padded $\GO$ point ($G = 32W = 128$ bytes). Write
$a$ for the total advice commitment count over all user phases, $\ell$ for the
number of lookups, $z$ for the number of permutation $Z$ chunks, $c$ for the
total lookup helper chunks, and $t$ for the number of trashcans. The rendered
bases are the prefix sums of the read schedule:

\begin{algorithm}[htbp]
\caption{Commitment-base derivation (generator-side; rendered as literals)}
\label{alg:const:comms}
\begin{algorithmic}[1]
\State $B_{\mathrm{advice}} \gets B$
\State $B_{\mathrm{lookup\,m}} \gets B_{\mathrm{advice}} + G\,a$
\State $B_{\mathrm{perm\,z}} \gets B_{\mathrm{lookup\,m}} + G\,\ell$
\State $B_{\mathrm{lookup\,helper}} \gets B_{\mathrm{perm\,z}} + G\,z$
\State $B_{\mathrm{lookup\,z}} \gets B_{\mathrm{lookup\,helper}} + G\,c$
\State $B_{\mathrm{trashcan}} \gets B_{\mathrm{lookup\,z}} + G\,\ell$
\State $B_{\mathrm{quotient\,limb}} \gets B_{\mathrm{trashcan}} + G\,t$
\State \Return the seven bases
\end{algorithmic}
\end{algorithm}

In closed form, with all quantities in bytes:
\begin{align}
  B_{\mathrm{lookup\,m}}        &= B + 128\,a \\
  B_{\mathrm{perm\,z}}          &= B + 128(a+\ell) \\
  B_{\mathrm{lookup\,helper}}   &= B + 128(a+\ell+z) \\
  B_{\mathrm{lookup\,z}}        &= B + 128(a+\ell+z+c) \\
  B_{\mathrm{trashcan}}         &= B + 128(a+2\ell+z+c) \\
  B_{\mathrm{quotient\,limb}}   &= B + 128(a+2\ell+z+c+t)
\end{align}

Note that $\ell$ appears \emph{twice}: once for the lookup $m$ commitments and
once for the lookup $Z$ commitments. The comment's cumulative list reflects this
(``$+\,4\cdot\texttt{num\_lookups}$'' appears at both the
\yul{PERM\_Z\_COMMS\_MPTR\_BASE} and \yul{TRASHCAN\_COMMS\_MPTR\_BASE} steps).
Inverting the rendered addresses gives the counts, and the inversion is
over-determined --- $\ell$ is recovered twice and both agree:

\begin{table}[htbp]
\centering
\begin{tabular}{@{}lllrl@{}}
\toprule
Category & Base & Extent & Points & Count recovered \\
\midrule
Advice          & \texttt{0xa140} & $[\texttt{0xa140},\texttt{0xa8c0})$ & $15$ & $a = 15$ \\
Lookup $m$      & \texttt{0xa8c0} & $[\texttt{0xa8c0},\texttt{0xa9c0})$ & $2$  & $\ell = 2$ \\
Permutation $Z$ & \texttt{0xa9c0} & $[\texttt{0xa9c0},\texttt{0xacc0})$ & $6$  & $z = 6$ \\
Lookup helper   & \texttt{0xacc0} & $[\texttt{0xacc0},\texttt{0xadc0})$ & $2$  & $c = 2$ \\
Lookup $Z$      & \texttt{0xadc0} & $[\texttt{0xadc0},\texttt{0xaec0})$ & $2$  & $\ell = 2$ (agrees) \\
Trashcan        & \texttt{0xaec0} & $[\texttt{0xaec0},\texttt{0xaf40})$ & $1$  & $t = 1$ \\
Quotient limbs  & \texttt{0xaf40} & $[\texttt{0xaf40},\texttt{0xb140})$ & $4$  & $4$ limbs \\
\midrule
Total           & \texttt{0xa140} & $[\texttt{0xa140},\texttt{0xb140})$ & $32$ & $\texttt{0x1000}$ bytes \\
\bottomrule
\end{tabular}
\caption{Commitment region, instantiated. The upper bound
$\texttt{0xb140} = \mptr{SELECTOR\_ACC\_MPTR}$ closes the region exactly.}
\label{tab:const:comms}
\end{table}

The quotient-limb count is corroborated by the final-MSM staging block, which
copies exactly four limb commitments at offsets
$\texttt{0x0}, \texttt{0x80}, \texttt{0x100}, \texttt{0x180}$ from
\mptr{QUOTIENT\_LIMB\_COMMS\_MPTR\_BASE} (lines 3920--3929).

\paragraph{Assessment.}
Storing commitments by category rather than by proof order is what makes the PCS
and quotient-fold emitters addressable: a fold over ``all lookup $Z$
commitments'' becomes a literal base plus a stride, with no indirection. The
cost is that the transcript parser must scatter into seven regions instead of
appending to one, which the generator absorbs at codegen time. The trade is
clearly worthwhile here, and the region is closed --- there is no slack and no
overlap between adjacent categories (Table~\ref{tab:const:comms}).

\subsubsection{Selector accumulators, batch-inversion scratch, and PCS staging}
\label{sec:const:b140}

Two constants share $\texttt{0xb140}$, and two further regions alias the same
base without being named:

\begin{table}[htbp]
\centering
\begin{tabularx}{\textwidth}{@{}lllX@{}}
\toprule
Region & Extent & Phase & Evidence \\
\midrule
constructor G1MSM smoke & $[\texttt{0xb140},\texttt{0xe200})$ & \texttt{ConstructorSmoke} & \yul{msm\_scratch := 0xb140}, length \texttt{0x30c0}, lines 549--555 \\
accumulator RHS MSM     & $[\texttt{0xb140},\texttt{0xb1e0})$ & \texttt{AccumulatorMsm} & \yul{acc\_scratch := 0xb140} line 1269; one \texttt{0xa0} pair staged twice in sequence (LHS lines 1282--1283, RHS lines 1314--1315) \\
\mptr{BATCH\_INV\_SCRATCH\_MPTR} & $[\texttt{0xb140},\texttt{0xb580})$ & \texttt{LagrangeBatchInvert} & \S\ref{sec:const:batchinv} \\
\mptr{SELECTOR\_ACC\_MPTR} & $[\texttt{0xb140},\texttt{0xb280})$ & \texttt{QuotientVm}--\texttt{PcsFinalMsm} & $10$ words; max offset \texttt{0x120}, lines 3315--3352 \\
\bottomrule
\end{tabularx}
\caption{The four uses of memory base \texttt{0xb140}.}
\label{tab:const:b140}
\end{table}

Above $\texttt{0xb280}$ the picture is the same, but none of the regions there
is named except \mptr{PCS\_FINAL\_MSM\_MPTR}:

\begin{table}[htbp]
\centering
\begin{tabularx}{\textwidth}{@{}lllX@{}}
\toprule
Region & Extent & Phase & Evidence \\
\midrule
quotient VM state       & $[\texttt{0xb280},\texttt{0xb8e0})$ & \texttt{QuotientVm} & \yul{mstore(0xb280, 0)} line 2021; highest observed access \yul{mload(add(0xb2c0, 0x0600))} $= \texttt{0xb8c0}$, line 3316 \\
quotient VM stack       & $[\texttt{0xb8e0},\texttt{0xbfa0})$ & \texttt{QuotientVm} & \yul{q\_sp := 0xb8e0} line 2172; overflow guard \yul{iszero(lt(q\_sp, 0xbfa0))} at 2218, 2332; underflow/balance guards on \yul{eq(q\_sp, 0xb8e0)} at 2229, 2464, 2580, 2772, 3105 \\
PCS q-eval source table & $[\texttt{0xb280},\texttt{0xb7e0})$ & \texttt{PcsQEvalSourceTable} & \yul{mstore(0xb280, 0x9980)} \dots \yul{mstore(0xb7c0, ...)}, lines 3463--3505 \\
\mptr{PCS\_FINAL\_MSM\_MPTR} & $[\texttt{0xb280},\texttt{0xe340})$ & \texttt{PcsFinalMsm} & $78 \times \texttt{0xa0} = \texttt{0x30c0}$, line 3999 \\
\bottomrule
\end{tabularx}
\caption{The four uses of memory base \texttt{0xb280}.}
\label{tab:const:b280}
\end{table}

The \mptr{SELECTOR\_ACC\_MPTR} region is the one long-lived allocation here: it
is zeroed at line 2027, written by the quotient VM through line 3087, rescaled
by the selector-tail block at lines 3315--3353, and read a last time by the
fused final MSM at lines 3932--3950 (e.g.\ line 3932,
\yul{mload(add(SELECTOR\_ACC\_MPTR, 0x0))}). Its span
$[\texttt{0xb140},\texttt{0xb280})$ ends exactly where
\mptr{PCS\_FINAL\_MSM\_MPTR} begins, so the two are disjoint in \emph{address},
not merely in time --- which is necessary, since their lifetimes overlap.
Ten simple selectors at $32$ bytes each is $\texttt{0x140}$ bytes, and
$\texttt{0xb140} + \texttt{0x140} = \texttt{0xb280}$. This is the tightest
adjacency in the map and it is exact.

\paragraph{The batch-inversion scratch is exactly sized, with zero slack.}
\label{sec:const:batchinv}
The Lagrange block (lines 1816--1826) builds its denominator run at
\mptr{LAGRANGE\_DENOMS\_MPTR}:
\begin{quote}\ttfamily
let mptr := LAGRANGE\_DENOMS\_MPTR \\
let mptr\_end := add(mptr, 0x03a0) \\
\dots \\
mstore(mptr\_end, x\_n\_minus\_1) \\
success, lagrange\_precompile\_failed := batch\_invert(success, LAGRANGE\_DENOMS\_MPTR, add(mptr\_end, 0x20), BATCH\_INV\_SCRATCH\_MPTR, r)
\end{quote}
so the run length is
$n = \texttt{0x03a0}/32 + 1 = 29 + 1 = 30$ words, occupying
$[\texttt{0xb580}, \texttt{0xb940})$. Reading \yul{batch\_invert} (lines
836--951): the forward pass stores one prefix product per element
\emph{excluding} the first and last, i.e.\ $n-2$ words at
\yul{scratch\_mptr}, and then writes the $6$-word EIP-198 modexp frame
immediately after them. The scratch requirement is therefore
\[
  \mathrm{scratch}(n) = 32(n-2) + 192 .
\]
For $n = 30$ this is $32\cdot 28 + 192 = 1088 = \texttt{0x440}$, and
\[
  \mptr{LAGRANGE\_DENOMS\_MPTR} - \mptr{BATCH\_INV\_SCRATCH\_MPTR}
  = \texttt{0xb580} - \texttt{0xb140} = \texttt{0x440}
\]
--- the scratch window is sized to the byte with no slack. Note also that
$n = 30 = m + |\mathrm{rot}_{\mathrm{last}}| + 1$ with $m = 19$ recovers
$\mathrm{rot}_{\mathrm{last}} = -10$.

\begin{finding}[const-7]{Batch-inversion scratch abuts the denominator run with zero slack and no runtime bound --- Informational}
\texttt{Halo2Verifier.sol}:229--239 (constants), 836--951 (\yul{batch\_invert}),
1816--1826 (call site). \yul{batch\_invert} takes \yul{scratch\_mptr} as a bare
pointer and never compares its advancing \yul{gp\_mptr} against any limit. The
scratch window $[\texttt{0xb140},\texttt{0xb580})$ is exactly
$32(n-2)+192$ bytes for the $n = 30$ this render passes, and
\mptr{LAGRANGE\_DENOMS\_MPTR} starts at the very next byte. If the run length
and the window ever disagreed by one word, the forward pass's modexp frame would
land on \texttt{denominator[0]} and corrupt it \emph{silently}: the modexp still
succeeds, \yul{ret} stays $1$, and the backward pass returns wrong inverses
without reverting --- a wrong-answer failure, not a revert.

This is \textbf{not reachable from any input} in the deployed artifact. $n$ is
determined by the literal loop bound \texttt{0x03a0} (line 1817) and the pinned
instance count (line 1420); no proof or instance value can change it. The risk is
codegen drift, and the generator does defend against it: it recomputes the
template's run length independently and asserts both equality with its own count
and $\mathrm{scratch} \ge 32(n-2)+192$
(\texttt{src/lowering/layout/memory.rs}:698--731, whose own comment describes
exactly this failure mode). Recorded at Informational
severity because the on-chain artifact has no defence of its own for a
failure mode whose signature is a wrong answer rather than a revert, and because
a reader of the constants block cannot see that the two addresses are related at
all.
\end{finding}

\begin{finding}[const-8]{Six live memory regions are addressed by bare literals and appear in no constant --- Observation}
The constant block of lines 104--240 does not name: the scalar-inversion modexp
scratch at \texttt{0x3580} (line 710); the quotient constant pool at
\texttt{0x3a60} (line 2013); the quotient program bytes at \texttt{0x50a0}
(line 2016); the fixed VK commitments at \texttt{0x6280} (lowest source address
in the final-MSM staging block, line 3860; first use line 3848); the quotient VM
state base at \texttt{0xb280} (line 2021); and the quotient VM operand stack at
\texttt{0xb8e0} with its ceiling \texttt{0xbfa0} (lines 2172, 2218). The VM's
own operand-window guard is written on the bare literal \texttt{0x3680} rather
than \mptr{VK\_MPTR} for the same reason \yul{q\_program\_fail} inlines its
selector (finding const-3): the VM renders into a second contract where the
constants are out of scope. Four of the six also \emph{overlap} named regions in address and
are separated only by phase ordering --- in particular the quotient VM state and
stack together span $[\texttt{0xb280},\texttt{0xbfa0})$, which contains all of
\mptr{LAGRANGE\_DENOMS\_MPTR}'s $[\texttt{0xb580},\texttt{0xb940})$.

The ordering is correct and checkable from source order: \yul{batch\_invert}
returns at line 1826 and the Lagrange block distils its results into the named
theta slots at lines 1830--1870, all before the quotient VM's first write at
line 2019. But this is exactly the class of relationship the constants block
exists to make visible, and here it is not visible. The comment at lines
198--203 even asserts the opposite convention for the final-MSM block (``Every
generated address in that block is written as an offset from this base'').
No security impact in this artifact; a reviewability and drift-resistance
observation.
\end{finding}

\begin{finding}[const-9]{Four final-MSM stores use absolute literals, contradicting the stated convention --- Observation}
\texttt{Halo2Verifier.sol}:198--203 documents \mptr{PCS\_FINAL\_MSM\_MPTR} with
``Every generated address in that block is written as an offset from this base
so a reader can see which term a store belongs to; solc folds the addition away,
so the symbolic form costs no runtime bytes.'' Four \opcode{MCOPY} destinations
in that block are absolute literals instead:
\texttt{0xcc20} (line 3920), \texttt{0xccc0} (3923), \texttt{0xcd60} (3926),
\texttt{0xce00} (3929) --- the four quotient-limb commitment terms. The
addresses are correct:
$\texttt{0xcc20}-\texttt{0xb280} = \texttt{0x19a0}$ and the matching scalar goes
to \yul{add(PCS\_FINAL\_MSM\_MPTR, 0x1a20)} $= \texttt{0x19a0}+\texttt{0x80}$,
and likewise for the other three. So this is a cosmetic inconsistency with a
documented convention, not a mis-addressed store. It matters only because a
reader auditing the staging buffer against the stated convention will find four
stores that cannot be attributed to a term by inspection.
\end{finding}

\subsubsection{Complete memory map, sorted by address}
\label{sec:const:fullmap}

Table~\ref{tab:const:map} lists every region the artifact uses, in address
order, whether or not it is named by a constant. ``$\phi$'' marks a
phase-scoped region (may be aliased by a region of a different phase);
``P'' marks a \texttt{Permanent} region (must be disjoint from everything).

\begin{table}[htbp]
\centering
\small
\begin{tabularx}{\textwidth}{@{}llllX@{}}
\toprule
Start & End & Life & Constant / name & Contents \\
\midrule
\texttt{0x0000} & \texttt{0x0040} & --- & (Solidity scratch) & \yul{fail(sel)} writes the 4-byte selector at \texttt{0x00} \\
\texttt{0x0040} & \texttt{0x0060} & --- & (free-memory pointer) & read by the layout guard, lines 363, 674 \\
\texttt{0x0060} & \texttt{0x0080} & --- & (zero slot) & untouched \\
\texttt{0x0080} & \texttt{0x1000} & --- & (solc via-IR spill) & reserved; guard asserts \yul{mload(0x40)} $\le \texttt{0x1000}$ \\
\midrule
\texttt{0x1000} & \texttt{0x1320} & $\phi$ & constructor smoke & \opcode{MCOPY}, modexp, G1ADD, G1MSM, pairing probes \\
\texttt{0x1000} & \texttt{0x1240} & $\phi$ & accumulator pairing batch & tag, \yul{vk\_digest}, four $\GO$; hashed as \texttt{0x240} bytes \\
\texttt{0x1000} & $\cdot$         & $\phi$ & \mptr{TRANSCRIPT\_MPTR} & Keccak absorb buffer; resets to base on squeeze \\
\texttt{0x1000} & \texttt{0x1180} & $\phi$ & \mptr{QUOTIENT\_RETURN\_MPTR} & split-evaluator return frame; unwritten in this render \\
\texttt{0x1000} & \texttt{0x1120} & $\phi$ & PCS pairing scratch & observed extent; $32$ words reserved \\
\texttt{0x1000} & \texttt{0x1020} & $\phi$ & \mptr{RETURN\_MPTR} & one word $=1$ on success \\
\texttt{0x1240} & \texttt{0x1540} & $\phi$ & final pairing scratch & $24$ words: two $(\GO,\GT)$ pairs; the result word overwrites the first ($25$ words reserved) \\
\texttt{0x3580} & \texttt{0x3680} & $\phi$ & scalar-inv scratch & \texttt{0xc0} modexp frame, \texttt{0x40} padding \\
\midrule
\texttt{0x3680} & \texttt{0x3a60} & P & \mptr{VK\_MPTR} header & $31$ words (Table~\ref{tab:const:vkimage}) \\
\texttt{0x3680} & \texttt{0x36a0} & P & \mptr{VK\_DIGEST\_MPTR} & transcript seed \\
\texttt{0x36a0} & \texttt{0x36c0} & P & \mptr{NUM\_INSTANCES\_MPTR} & $=19$ \\
\texttt{0x36c0} & \texttt{0x36e0} & P & \mptr{K\_MPTR} & $\log_2 n$ \\
\texttt{0x36e0} & \texttt{0x3700} & P & \mptr{N\_INV\_MPTR} & $n^{-1}$ \\
\texttt{0x3700} & \texttt{0x3720} & P & \mptr{OMEGA\_MPTR} & $\omega$ \\
\texttt{0x3720} & \texttt{0x3740} & P & \mptr{OMEGA\_INV\_MPTR} & $\omega^{-1}$ \\
\texttt{0x3740} & \texttt{0x3760} & P & \mptr{OMEGA\_INV\_TO\_L\_MPTR} & $\omega^{-|\mathrm{rot}_{\mathrm{last}}|}$ \\
\texttt{0x3760} & \texttt{0x3780} & P & \mptr{HAS\_ACCUMULATOR\_MPTR} & $=1$ \\
\texttt{0x3780} & \texttt{0x37a0} & P & \mptr{ACC\_OFFSET\_MPTR} & $=11$ \\
\texttt{0x37a0} & \texttt{0x37c0} & P & \mptr{NUM\_ACC\_LIMBS\_MPTR} & $=7$ \\
\texttt{0x37c0} & \texttt{0x37e0} & P & \mptr{NUM\_ACC\_LIMB\_BITS\_MPTR} & $=56$ \\
\texttt{0x37e0} & \texttt{0x3860} & P & \mptr{G1\_BASE\_MPTR} & padded $\GO$ generator \\
\texttt{0x3860} & \texttt{0x3960} & P & \mptr{G2\_BASE\_MPTR} & padded $\GT$ generator \\
\texttt{0x3960} & \texttt{0x3a60} & P & \mptr{NEG\_S\_G2\_BASE\_MPTR} & padded $-[s]\GT$ \\
\texttt{0x3a60} & \texttt{0x50a0} & P & (quotient constants) & $178$ $\Fr$ words \\
\texttt{0x50a0} & \texttt{0x626f} & P & (quotient program) & \texttt{0x11cf} packed bytes \\
\texttt{0x626f} & \texttt{0x6280} & P & (alignment padding) & $17$ bytes \\
\texttt{0x6280} & \texttt{0x7900} & P & (fixed commitments) & $45 \times \GO$ \\
\midrule
\texttt{0x7900} & \texttt{0x7900} & P & \mptr{CHALLENGE\_MPTR} & empty (no user challenges) \\
\texttt{0x7900} & \texttt{0x7f80} & P & \mptr{THETA\_MPTR} band & $52$ words (Table~\ref{tab:const:theta}) \\
\texttt{0x7da0} & \texttt{0x7dc0} & P & \emph{(hole)} & unnamed, never written \\
\texttt{0x7f80} & \texttt{0x8300} & P & \mptr{ROT\_POINTS\_MPTR} & cap $28$ w, $4$ used \\
\texttt{0x8300} & \texttt{0x8b20} & P & \mptr{X1\_POWERS\_MPTR} & cap $65$ w, $43$ used \\
\texttt{0x8b20} & \texttt{0x9220} & P & \mptr{Q\_EVAL\_SET\_MPTR} & cap $56$ w, $11$ used \\
\texttt{0x9220} & \texttt{0x9240} & P & \mptr{Q\_EVAL\_CPTR\_MPTR} & runtime calldata cursor \\
\texttt{0x9240} & \texttt{0x9320} & P & \emph{(hole)} & $7$ words padding \\
\texttt{0x9320} & \texttt{0x93a0} & P & \mptr{G1\_IDENTITY\_MPTR} & zero $\GO$; never read here \\
\texttt{0x93a0} & \texttt{0x9480} & P & \emph{(hole)} & $7$ words padding \\
\texttt{0x9480} & \texttt{0xa140} & P & \mptr{REVERSED\_EVALS\_MPTR} & $\mathbf{e} = (e_0,\dots,e_{101})$ \\
\midrule
\texttt{0xa140} & \texttt{0xa8c0} & P & \mptr{ADVICE\_COMMS\_MPTR\_BASE} & $15 \times \GO$ \\
\texttt{0xa8c0} & \texttt{0xa9c0} & P & \mptr{LOOKUP\_M\_COMMS\_MPTR\_BASE} & $2 \times \GO$ \\
\texttt{0xa9c0} & \texttt{0xacc0} & P & \mptr{PERM\_Z\_COMMS\_MPTR\_BASE} & $6 \times \GO$ \\
\texttt{0xacc0} & \texttt{0xadc0} & P & \mptr{LOOKUP\_HELPER\_COMMS\_MPTR\_BASE} & $2 \times \GO$ \\
\texttt{0xadc0} & \texttt{0xaec0} & P & \mptr{LOOKUP\_Z\_COMMS\_MPTR\_BASE} & $2 \times \GO$ \\
\texttt{0xaec0} & \texttt{0xaf40} & P & \mptr{TRASHCAN\_COMMS\_MPTR\_BASE} & $1 \times \GO$ \\
\texttt{0xaf40} & \texttt{0xb140} & P & \mptr{QUOTIENT\_LIMB\_COMMS\_MPTR\_BASE} & $4 \times \GO$ \\
\midrule
\texttt{0xb140} & \texttt{0xe200} & $\phi$ & constructor G1MSM smoke & \texttt{0x30c0} zero bytes \\
\texttt{0xb140} & \texttt{0xb1e0} & $\phi$ & accumulator RHS MSM & one $(\GO,\Fr)$ pair \\
\texttt{0xb140} & \texttt{0xb580} & $\phi$ & \mptr{BATCH\_INV\_SCRATCH\_MPTR} & $28$ prefix products $+$ modexp frame \\
\texttt{0xb140} & \texttt{0xb280} & $\phi$ & \mptr{SELECTOR\_ACC\_MPTR} & $10$ selector accumulators \\
\texttt{0xb280} & \texttt{0xb8e0} & $\phi$ & (quotient VM state) & $51$ words: Horner accumulator at \texttt{0xb280}, $y$-power table from \texttt{0xb2c0} \\
\texttt{0xb280} & \texttt{0xb7e0} & $\phi$ & (PCS q-eval source table) & $43$ source addresses \\
\texttt{0xb280} & \texttt{0xe340} & $\phi$ & \mptr{PCS\_FINAL\_MSM\_MPTR} & $78 \times (\GO,\Fr)$ pairs \\
\texttt{0xb580} & \texttt{0xb940} & $\phi$ & \mptr{LAGRANGE\_DENOMS\_MPTR} & $30$-word inversion run \\
\texttt{0xb8e0} & \texttt{0xbfa0} & $\phi$ & (quotient VM stack) & $54$ words, guarded both ends \\
\texttt{0xe340} & \texttt{0xe360} & P & \mptr{TRACE\_U256\_MPTR} & trace log word; unused here \\
\bottomrule
\end{tabularx}
\caption{Complete memory map of the fixture, sorted by start address.
Regions marked $\phi$ are phase-scoped and may share addresses with
$\phi$ regions of other phases; regions marked P are permanent and must be
globally disjoint.}
\label{tab:const:map}
\end{table}

\subsubsection{Provable gaps and overlaps}
\label{sec:const:gaps}

\paragraph{Gaps (allocated but never written).}
The following holes are provable from the fixture alone:
\begin{enumerate}
  \item \textbf{Theta word $37$}, $[\texttt{0x7da0},\texttt{0x7dc0})$, one word.
        \mptr{QUOTIENT\_MPTR} is $4$ words ending at \texttt{0x7da0}, and
        \mptr{F\_EVAL\_MPTR} is \texttt{0x7dc0}. No constant names
        \texttt{0x7da0} and no store targets it.
  \item \textbf{Q-eval cursor padding}, $[\texttt{0x9240},\texttt{0x9320})$,
        $7$ words. \mptr{Q\_EVAL\_CPTR\_MPTR} is one word.
  \item \textbf{G1-identity padding}, $[\texttt{0x93a0},\texttt{0x9480})$,
        $7$ words.
  \item \textbf{Rotation-point slack}, $[\texttt{0x8000},\texttt{0x8300})$,
        $24$ of $28$ words unused.
  \item \textbf{$x_1$-power slack}, $[\texttt{0x8860},\texttt{0x8b20})$,
        $22$ of $65$ words unused.
  \item \textbf{Q-eval-set slack}, $[\texttt{0x8c80},\texttt{0x9220})$,
        $45$ of $56$ words unused.
  \item \textbf{VK alignment padding}, $[\texttt{0x626f},\texttt{0x6280})$,
        $17$ bytes, rounding the quotient program end up to a word boundary.
  \item \textbf{Constructor-smoke tail}, $[\texttt{0xe200},\texttt{0xe340})$:
        the constructor's \texttt{0x30c0}-byte probe starts at \texttt{0xb140}
        and ends at \texttt{0xe200}, whereas \mptr{PCS\_FINAL\_MSM\_MPTR}'s
        equally sized buffer starts $\texttt{0x140}$ higher and ends at
        \texttt{0xe340}. The last $\texttt{0x140}$ bytes are therefore never
        touched during deployment. Harmless --- constructor memory is discarded
        and the probe's purpose is input-\emph{length} coverage --- but it means
        the probe does not pre-expand the exact runtime byte range.
\end{enumerate}
The three ``slack'' entries (4--6) are capacity windows fixed by the generator's
historical layout rather than by this circuit's needs; they cost nothing at
runtime (EVM memory expansion is charged on the high-water mark, which is
\texttt{0xe340} regardless) and they keep addresses stable across circuits of
different sizes.

\paragraph{Overlaps.}
Every address collision in Table~\ref{tab:const:map} was checked. There are
exactly two aliasing bases, \texttt{0x1000} and \texttt{0xb140}/\texttt{0xb280},
and in each case the colliding regions carry disjoint phases whose order is
verifiable from source order in the fixture:
\begin{itemize}
  \item \texttt{0x1000}: Table~\ref{tab:const:band1000}. The transcript is
        finished at line 1785; the PCS pairing scratch is first written at line
        4012; the accumulator batch frame at line 4057;
        \mptr{RETURN\_MPTR} at line 4127.
  \item \texttt{0xb140}: Table~\ref{tab:const:b140}. The constructor probe runs
        only in the creation frame. \yul{validate\_public\_accumulator}
        (defined line 1246) is called at line 1447, before the transcript block
        begins at line 1455. \yul{batch\_invert} returns at
        line 1826 and the Lagrange results are distilled into theta slots at
        lines 1869--1874, before the selector-accumulator zeroing loop at line
        2027.
  \item \texttt{0xb280}: Table~\ref{tab:const:b280}. The quotient VM completes
        by line 3359 (last selector write at line 3087, last read of the VM
        Horner word \texttt{0xb280} at line 3358); the q-eval source table is
        built at lines 3463--3505; the
        final MSM buffer is filled at lines 3836--3996 and consumed at line
        3999.
\end{itemize}
The only genuinely simultaneous pair --- \mptr{SELECTOR\_ACC\_MPTR}, live from
the quotient VM through the final MSM, and \mptr{PCS\_FINAL\_MSM\_MPTR} --- is
separated in \emph{address}, exactly and with no slack, at \texttt{0xb280}.

\paragraph{Conclusion.} No unintended overlap is provable from the fixture, and
none was found. The map is sound. What carries the soundness is generation-time
analysis plus execution ordering, not any runtime check; findings const-4,
const-7 and const-8 record where that is worth a reader's attention.

\subsection{Precompile gas bounds}
\label{sec:const:gas}

\begin{lstlisting}[style=solidity,caption={\texttt{Halo2Verifier.sol} lines 290--299: gas bound constants (rationale comment at 265--289 elided)},label={lst:const:gas}]
    uint256 internal constant          G1ADD_GAS = 375;
    uint256 internal constant   G1MSM_GAS_1PAIR = 12000;
    uint256 internal constant PAIRING_GAS_2PAIR = 102900;
    uint256 internal constant        MODEXP_GAS = 4080;
    // Exact cost of the deployment-time worst-case G1MSM smoke probe.
    uint256 internal constant G1MSM_GAS_SMOKE = 525096;
    // Worst-case accumulator RHS MSM: carried RHS point plus every generated
    // fixed-base tail scalar nonzero. Zero tail scalars are omitted at
    // runtime, which only lowers the actual cost below this bound.
    uint256 internal constant ACC_RHS_MSM_GAS = 12000;
\end{lstlisting}

\paragraph{The design.}
Every generated \opcode{STATICCALL} to a precompile forwards an exact scheduled
cost instead of \yul{gas()}. The reason is stated at lines 269--275: a failing
EIP-2537 or modexp call consumes \emph{all} gas supplied to the frame. Under
EIP-150's 63/64 rule, forwarding \yul{gas()} would let one malformed proof point
burn essentially the whole transaction budget; forwarding the scheduled cost
caps the loss at that cost. Since the schedule is the spec-guaranteed price ---
it does not depend on the input values, only on their sizes, which are all
compile-time constants here --- the bounds are sufficient by construction on any
conformant chain. On success the unused portion is returned, so the bound costs
nothing when the proof is valid.

\paragraph{Derivations.}
All six values were checked against the published schedules.

\begin{table}[htbp]
\centering
\begin{tabularx}{\textwidth}{@{}lllX@{}}
\toprule
Constant & Value & Schedule & Derivation \\
\midrule
\yul{G1ADD\_GAS}         & $375$    & EIP-2537 & flat price of \texttt{0x0b} G1ADD \\
\yul{G1MSM\_GAS\_1PAIR}  & $12000$  & EIP-2537 & $\lfloor k\cdot 12000\cdot d(k)/1000\rfloor$ at $k=1$, $d(1)=1000$ \\
\yul{PAIRING\_GAS\_2PAIR}& $102900$ & EIP-2537 & $32600k + 37700$ at $k=2$: $65200+37700$ \\
\yul{MODEXP\_GAS}        & $4080$   & EIP-7883 & $16 \times 255$; see below \\
\yul{G1MSM\_GAS\_SMOKE}  & $525096$ & EIP-2537 & $\lfloor 78\cdot 12000\cdot 561/1000\rfloor = \lfloor 936000\cdot 0.561\rfloor$ \\
\yul{ACC\_RHS\_MSM\_GAS} & $12000$  & EIP-2537 & $k=1$; see below \\
\bottomrule
\end{tabularx}
\caption{Gas bound constants and their schedule derivations.}
\label{tab:const:gas}
\end{table}

$\yul{G1MSM\_GAS\_SMOKE}$ is consistent to the gas: $78 \cdot 12000 = 936000$
and $936000 \cdot 561 / 1000 = 525096$ exactly, matching the $78$-pair input
length \texttt{0x30c0} used by both the constructor probe (line 555) and the
runtime fused MSM, which forwards the same number as a literal (line 3999).

\paragraph{The modexp bound.}
The frame is the Fermat inversion $x^{r-2} \bmod r$ with $32$-byte base,
exponent and modulus (lines 713--719). Under EIP-2565 the price is
$\max(200,\, \lfloor \mathrm{mc}\cdot \mathrm{ic}/3\rfloor)$ with
multiplication complexity $\mathrm{mc} = \lceil 32/8\rceil^2 = 16$; under
EIP-7883 (Osaka/Fusaka) the price is $\max(500,\, \mathrm{mc}\cdot\mathrm{ic})$
with the same $\mathrm{mc}=16$ for $\le 32$-byte operands --- the $/3$ divisor
is removed. Taking $\mathrm{ic} = 255$, the worst case over all $32$-byte
exponents:
\[
  \text{EIP-2565: } \left\lfloor \tfrac{16\cdot 255}{3}\right\rfloor = 1360,
  \qquad
  \text{EIP-7883: } 16\cdot 255 = 4080 .
\]
The larger is rendered, which is the only safe choice: forwarding the EIP-7883
bound on a pre-Osaka chain is free on success, whereas forwarding the EIP-2565
bound on a repriced chain would out-of-gas and revert \emph{every} proof.

\begin{finding}[const-10]{\yul{MODEXP\_GAS} is a $16$-gas over-approximation of this frame's exact price --- Informational}
\texttt{Halo2Verifier.sol}:277--282, 293. The comment states that EIP-2565
``prices this frame at 1360''. For \emph{this} frame the exponent is $r-2$ with
$r$ the $255$-bit BLS12-381 scalar modulus, so
$\mathrm{ic} = \mathrm{bitlen}(r-2) - 1 = 254$ and the exact prices are
$\lfloor 16\cdot 254/3\rfloor = 1354$ (EIP-2565) and $16\cdot 254 = 4064$
(EIP-7883). The rendered $4080$ corresponds to $\mathrm{ic}=255$, i.e.\ the
worst case over all $32$-byte exponents rather than this specific one.

The direction is the safe one: $4080 > 4064$, so the forwarded bound never
under-funds the call, and the $16$-gas surplus is refunded on success. Recorded
only so a reader reconciling the constant with the EIP does not conclude the
arithmetic is wrong. The constructor probes the same bound against the same
frame (line 405), so a chain priced above $4080$ fails at deployment rather than
on the first proof.
\end{finding}

\begin{finding}[const-11]{\yul{ACC\_RHS\_MSM\_GAS} is the one-pair price but funds a runtime-computed length --- Observation}
\texttt{Halo2Verifier.sol}:296--299 and 1322--1345. The accumulator RHS MSM
forwards the constant $12000$ while passing a \emph{runtime} input length
\yul{acc\_msm\_len := sub(acc\_pair\_ptr, acc\_scratch)}. In this render the
block that appends pairs (lines 1306--1318) is unconditional and appends exactly
one $\texttt{0xa0}$-byte pair, so $\yul{acc\_msm\_len} \equiv \texttt{0xa0}$ and
$12000$ is the exact EIP-2537 price for $k=1$. The artifact is correct.

The observation is about the coupling, not this artifact. The comment describes
the constant as ``the compile-time worst case (every tail scalar nonzero)'',
which is accurate here only because the tail is empty; a render with even one
fixed-base tail term would need $k=2$, priced at
$\lfloor 2\cdot 12000\cdot d(2)/1000\rfloor = 22776$ for $d(2)=949$, and the
under-funded call would revert \revert{PrecompileFailed} on every proof. That
failure would be caught at deployment by the constructor's \texttt{0x30c0}-byte
G1MSM probe only if the probe's bound were the binding one, which it is not ---
the probe uses \yul{G1MSM\_GAS\_SMOKE}. The safety of this constant rests
entirely on the generator emitting it from the same shape that emits the
staging block.
\end{finding}

\paragraph{Constructor coverage (clean).}
Lines 284--288 make a liveness claim that the code honours: the constructor
smoke probes forward the same bounds for \emph{every} precompile the runtime
calls. Verified --- \yul{MODEXP\_GAS} at line 405, \yul{G1ADD\_GAS} at 418 and
439, \yul{G1MSM\_GAS\_1PAIR} at 473, \yul{G1MSM\_GAS\_SMOKE} at 555, and
\yul{PAIRING\_GAS\_2PAIR} at 525, 532 and 565. A chain whose schedule was
repriced above any of these bounds therefore fails at deployment rather than
bricking silently. The modexp probe additionally verifies semantics rather than
mere success: it computes $2^{r-2} \bmod r$ and checks
$2 \cdot \text{result} \equiv 1 \pmod r$ (line 407), which a stub returning
zeros or echoing its input cannot pass. Line 500 is a deliberate
\emph{negative} probe (\yul{if staticcall(200000, 0x0c, ...) \{ revert \}}),
asserting that a malformed input is rejected. This is a well-constructed
deployment gate.

\subsection{Build identity}
\label{sec:const:buildid}

\begin{lstlisting}[style=solidity,caption={\texttt{Halo2Verifier.sol} lines 301--311: \yul{BUILD\_ID}},label={lst:const:buildid}]
    /// @notice Build identity for this generated artifact (P10/L-8).
    /// @dev keccak256 over: the domain tag "halo2-solidity-verifier-build-v1",
    ///      the u64-length-prefixed generator feature profile, the vk_digest,
    ///      the expected VK runtime codehash (zero when the VK is embedded),
    ///      the SRS fingerprint keccak("halo2-solidity-verifier-srs-v1" || n
    ///      || G2 || s_g2 || [tau]G1), and an optional 32-byte deployment
    ///      provenance tag (0x00 marker when absent, 0x01 || tag when set).
    ///      The deployment record must publish these preimage components so
    ///      third parties can recompute the id; see
    ///      docs/reference/DEPLOYMENT_AND_INCIDENT_RESPONSE.md.
    bytes32 public constant BUILD_ID = 0x76eca8e8f19f0bb8c7d8d0b7757bf44687a635676d3a55e5ca8efb47b4481999;
\end{lstlisting}

The preimage is the concatenation, in this exact order:
\begin{align*}
  \yul{BUILD\_ID} = \mathrm{keccak256}\Big(&
    \texttt{"halo2-solidity-verifier-build-v1"} \\
  &\Vert\; \mathrm{be}_{64}(|F|) \;\Vert\; F \\
  &\Vert\; \mathrm{be}_{256}(\yul{vk\_digest}) \\
  &\Vert\; \mathrm{be}_{256}(\yul{EXPECTED\_VK\_CODEHASH}) \\
  &\Vert\; \sigma_{\mathrm{srs}} \\
  &\Vert\; \begin{cases}
      \texttt{0x00} & \text{no provenance tag} \\
      \texttt{0x01} \Vert \mathrm{tag}_{32} & \text{otherwise}
    \end{cases}\Big)
\end{align*}
where $F$ is the generator's feature-profile string (a build-time environment
value), the codehash field is zero when the VK is embedded rather than pinned
by address, and
\[
  \sigma_{\mathrm{srs}} = \mathrm{keccak256}\big(
    \texttt{"halo2-solidity-verifier-srs-v1"} \Vert \mathrm{be}_{64}(n)
    \Vert \GT \Vert [s]\GT \Vert [\tau]\GO \big),
\]
with each group element serialised as its EIP-2537 padded $256$-bit words.

\paragraph{Assessment.}
The constant is a public getter, so it costs a few dozen runtime bytes ---
non-trivial against the EIP-170 budget of \S\ref{sec:const:pragma}, and
deliberately spent. What it buys is that an incident responder can bind a
deployed address to a specific generator input set (VK, SRS, feature profile,
optional deployment tag) with one \yul{eth\_call}, without decompiling.

Two limits are worth stating precisely because the docstring does not:
\begin{itemize}
  \item The identity is computed by the Rust generator \emph{before} solc runs.
        It therefore covers the generated \emph{source}, not the compiled
        bytecode, and cannot detect the compiler/optimiser drift of
        finding const-1. Bytecode identity must be established separately, by
        \opcode{EXTCODEHASH} against a published value.
  \item The preimage is only recomputable by a third party if all components are
        published --- in particular the feature profile string $F$, which exists
        nowhere in the artifact. The docstring makes publishing them a
        requirement of the deployment record, which is the right place for it,
        but the on-chain constant is unverifiable without that record.
\end{itemize}

\subsection{Field constants and packed accumulator sentinels}
\label{sec:const:field}

\begin{lstlisting}[style=solidity,caption={\texttt{Halo2Verifier.sol} lines 329--342: field constants and coordinate sentinels},label={lst:const:field}]
    // BLS12-381 scalar-field modulus, used for transcript challenges and all
    // Halo2 verifier arithmetic.
    uint256 internal constant FR_MODULUS        = 0x73eda753299d7d483339d80809a1d80553bda402fffe5bfeffffffff00000001;

    // BLS12-381 Fp modulus minus one, split like an EIP-2537 coordinate:
    // high word = 16 zero bytes || top 16 coordinate bytes, low word =
    // bottom 32 coordinate bytes.
    uint256 internal constant BLS_P_HI             = 0x000000000000000000000000000000001a0111ea397fe69a4b1ba7b6434bacd7;
    uint256 internal constant BLS_P_MINUS_ONE_LO   = 0x64774b84f38512bf6730d2a0f6b0f6241eabfffeb153ffffb9feffffffffaaaa;
    // Packed public-accumulator sentinels for the shifted coordinate codec.
    // The `_WITH_ID_FLAG` variant is used only for the first x-coordinate word.
    uint256 internal constant BLS_P_MINUS_ONE_PACKED_0 = 0x00000000f38512bf6730d2a0f6b0f6241eabfffeb153ffffb9feffffffffaaaa;
    uint256 internal constant BLS_P_MINUS_ONE_PACKED_0_WITH_ID_FLAG = 0x00000000f38512bf6730d2a0f6b0f6241eabfffeb153ffffbafeffffffffaaaa;
    uint256 internal constant BLS_P_MINUS_ONE_PACKED_1 = 0x0000000000000000000000001a0111ea397fe69a4b1ba7b6434bacd764774b84;
\end{lstlisting}

\subsubsection{\texorpdfstring{$\Fr$}{Fr} modulus}

$r = \texttt{0x73eda753299d7d483339d80809a1d80553bda402fffe5bfeffffffff00000001}$
is the BLS12-381 scalar modulus, $255$ bits wide. It is the modulus for every
\yul{addmod}/\yul{mulmod} in the file, the modulus of the Fermat inversion
$x^{r-2}$, and the reduction modulus for transcript squeezes. Because
$r < 2^{255} < 2^{256}$, a $256$-bit calldata word can encode a non-canonical
representative; the verifier therefore range-checks every scalar it reads
(\revert{NonCanonicalScalar} sites in Table~\ref{tab:const:errors}) rather than
relying on \yul{mulmod}'s implicit reduction. That is the correct choice: silent
reduction would make two distinct calldata encodings verify identically, which
is a proof-malleability surface even when it is not a soundness break.

\subsubsection{\texorpdfstring{$\Fp$}{Fp} bounds in EIP-2537 split form}

$p$ is the $381$-bit base-field modulus. The two constants are the EIP-2537
two-word split of $p - 1$:
\[
  \yul{BLS\_P\_HI} = \left\lfloor \frac{p}{2^{256}} \right\rfloor
    = \texttt{0x1a0111ea397fe69a4b1ba7b6434bacd7},
  \qquad
  \yul{BLS\_P\_MINUS\_ONE\_LO} = (p-1) \bmod 2^{256} .
\]
Both were recomputed and match. The name asymmetry is correct and worth noting:
the high word of $p$ and of $p-1$ are equal (the decrement does not borrow out
of the low word, since $p \equiv \texttt{0xaab} \pmod{2^{12}}$), so one constant
serves both roles.

These two words implement the canonical-range test for a padded coordinate
$(hi, lo)$ as a lexicographic comparison:
\[
  hi < \yul{BLS\_P\_HI}
  \;\lor\;
  \big(hi = \yul{BLS\_P\_HI} \;\land\; lo \le \yul{BLS\_P\_MINUS\_ONE\_LO}\big)
  \iff hi\cdot 2^{256} + lo \le p-1 .
\]
This is exactly the predicate at lines 785 and 788 (rejecting on its negation
with \revert{BadPointEncoding}) and at line 1126. Using $p-1$ with a
non-strict $\le$ rather than $p$ with a strict $<$ avoids needing a separate
constant for $p$ and makes both branches of the comparison use the same two
words.

Note that the range check is only half of the encoding check: lines 780--784
separately require the top $16$ bytes of each padded word to be zero ---
\yul{if shr(128, x\_hi\_word) \{ fail(ERR\_BAD\_POINT\_ENCODING) \}}, and likewise
for $y$ --- which is the EIP-2537 padding requirement; the surviving low $128$
bits are then masked out with
\yul{and(x\_hi\_word, 0xffffffffffffffffffffffffffffffff)} (lines 783--784) to
form the $hi$ operand of the range test above. Both halves are present, and both
raise \revert{BadPointEncoding}.

\subsubsection{Packed accumulator sentinels}

The public-accumulator codec packs a $381$-bit $\Fp$ coordinate into two
$256$-bit words differently from the EIP-2537 split: it uses a
$224$/$157$ split rather than $256$/$125$. Verified:
\[
  \yul{BLS\_P\_MINUS\_ONE\_PACKED\_0} = (p-1) \bmod 2^{224},
  \qquad
  \yul{BLS\_P\_MINUS\_ONE\_PACKED\_1} = \left\lfloor \frac{p-1}{2^{224}} \right\rfloor,
\]
with $\yul{PACKED\_0}$ occupying $224$ bits and $\yul{PACKED\_1}$ $157$ bits,
recomposing to $p-1$ exactly. The identity-flag variant differs from
$\yul{PACKED\_0}$ by exactly one bit:
\[
  \yul{BLS\_P\_MINUS\_ONE\_PACKED\_0\_WITH\_ID\_FLAG}
   = \yul{BLS\_P\_MINUS\_ONE\_PACKED\_0} + 2^{56} .
\]
$2^{56}$ is the limb radix of the accumulator encoding
($\yul{num\_acc\_limb\_bits} = 56$, line 1405), so the flag is a $+1$ in limb
$1$ --- i.e.\ the sentinel is the canonical ``encoded $p-1$'' coordinate with one
distinguished limb incremented, which keeps it inside the same $224$-bit word
and outside the value range any real coordinate limb can take.

The sentinels are used at lines 1053--1055 and 1066--1076 to recognise the canonical
identity encoding of a carried accumulator point:
\begin{itemize}
  \item Line 1054 tests the \emph{decoded} EIP-2537 pair against
        $(\yul{BLS\_P\_HI}, \yul{BLS\_P\_MINUS\_ONE\_LO})$, i.e.\ ``this
        coordinate is $p-1$''.
  \item \yul{is\_acc\_encoded\_identity} (lines 1066--1076) tests the four
        \emph{raw calldata} words directly:
        the $x$ pair must equal
        $(\yul{PACKED\_0\_WITH\_ID\_FLAG}, \yul{PACKED\_1})$ and the $y$ pair
        must equal $(\yul{PACKED\_0}, \yul{PACKED\_1})$.
\end{itemize}
The asymmetry --- flag on $x$, no flag on $y$ --- is what makes the identity
encoding a single distinguished point rather than a family, and it is why only
the first $x$-coordinate word needs the variant constant, as the comment says.

\paragraph{Assessment (clean).}
The design here is careful in a way that is easy to miss. EIP-2537 reserves
affine $(0,0)$ for the point at infinity, so a decoded infinity is a legitimate
group element; but the accumulator codec maps encoded $p-1$ to decoded zero.
Without a distinguished encoding, a malformed public input decoding to $(0,0)$
would be indistinguishable from a legitimately-carried identity. The code
resolves this by requiring the exact packed sentinel for the identity case
(lines 1066--1076), and by \emph{rejecting} a decoded zero in the non-identity
branch (lines 1213--1214, \yul{decoded\_zero} $\Rightarrow$ not \yul{ok}). Both
directions are checked. Furthermore every decoded accumulator point --- identity
or not, scalar $0$, $1$ or otherwise --- is routed through EIP-2537 G1MSM (lines
1282--1290 and 1314--1345) specifically so the precompile performs on-curve and
subgroup validation; the comment at lines 1277--1280 states the reason
explicitly (``skipping it would let a malformed non-identity point hide behind
scalar 0''). This region is sound.

\subsection{Section assessment}
\label{sec:const:assessment}

Lines 1--350 are declaration-only, so there is no control flow to attack here.
The security-relevant content is entirely in whether the declared values are the
values the rest of the file needs. On that question the region verifies well:

\begin{itemize}
  \item All eight error selectors are correct
        (Table~\ref{tab:const:errors}, recomputed).
  \item All six gas bounds match the EIP-2537 / EIP-7883 schedules
        (Table~\ref{tab:const:gas}); the one deviation, \yul{MODEXP\_GAS}, errs
        by $16$ gas in the safe direction (const-10).
  \item All five field constants are exact functions of $r$ and $p$
        (\S\ref{sec:const:field}, recomputed).
  \item The calldata frame is internally consistent and exactly closed
        (\S\ref{sec:const:abi}).
  \item The VK image decomposes to exactly \yul{EXPECTED\_VK\_PAYLOAD\_LENGTH}
        (Table~\ref{tab:const:vkimage}), and the commitment-base
        cumulative-offset rule inverts consistently, recovering the lookup count
        twice with agreement (Table~\ref{tab:const:comms}).
  \item No unintended memory overlap exists; every aliasing base is
        phase-disjoint with an ordering verifiable from source order
        (\S\ref{sec:const:gaps}). This is an execution-order argument made here
        against the fixture, not a property the on-chain artifact or the
        generator's arena check establishes on its own
        (const-4, const-12).
\end{itemize}

The findings raised are, in descending order of what a reviewer should act on:
const-1 (the pragma cannot pin the optimiser setting, so the deployment record
must); const-7 (the zero-slack batch-inversion scratch, whose failure mode is a
wrong answer rather than a revert, defended only at generation time); const-11
(a gas constant coupled to a runtime-computed length, correct here but by
construction rather than by check); const-4 and const-8 (aliasing and unnamed
regions, which are correct but not visible from the constants block);
const-10 (a harmless $16$-gas over-approximation); const-12 (a constructor
probe writing one word outside its planner-registered region, which is the
concrete instance of const-4's general point); and const-2, const-3, const-6 and
const-9 (documentation and dead-declaration inconsistencies with no runtime
effect). None is a soundness defect in this artifact.


%% ==== 02-deployment.tex =================================================
\section{Deployment-time runtime capability probes and the constructor}
\label{sec:deploy}

\subsection{Scope}
\label{sec:deploy:scope}

This section specifies \texttt{Halo2Verifier.sol} lines 319--623: the typed-error
selector constants (320--327, closing the banner comment that runs 313--319), the
field constants consumed by --- or merely adjacent to --- the deployment path
(329--342), the whole of \yul{require\_eip2537\_precompiles} (344--569), and the
constructor (572--586; line 571 holds only whitespace). Lines 588--623 are the
NatSpec block of \yul{verifyProof}; that function's declaration begins at line
624 and its body is specified elsewhere, referenced here only by line number.

Two conventions apply to every listing below. First, leading indentation is
trimmed uniformly: by 12 columns for the twelve excerpts taken from inside the
\yul{assembly} block of \yul{require\_eip2537\_precompiles}, and by 4 columns for
the constructor excerpt of \S\ref{sec:deploy:ctor}, which sits at Solidity's
contract-member level. The code is otherwise byte-for-byte verbatim; every
excerpt in this section was diffed character by character against the cited
range. Second,
$r$ denotes \yul{FR\_MODULUS}, the BLS12-381 scalar modulus declared at line 331:
\[
  r = \texttt{0x73eda753299d7d483339d80809a1d80553bda402fffe5bfeffffffff00000001},
\]
and $p$ denotes the BLS12-381 base-field modulus. Nothing in the deployment path
reads $p$: it appears in this range only through five declarations that the
point-decoding and public-accumulator logic (specified elsewhere) consumes ---
\yul{BLS\_P\_HI} and \yul{BLS\_P\_MINUS\_ONE\_LO} (lines 336--337, the
EIP-2537-style hi/lo split of $p-1$), and the three shifted-codec sentinels
\yul{BLS\_P\_MINUS\_ONE\_PACKED\_0}, \yul{BLS\_P\_MINUS\_ONE\_PACKED\_0\_WITH\_ID\_FLAG}
and \yul{BLS\_P\_MINUS\_ONE\_PACKED\_1} (lines 340--342). Of the constants
declared at 329--342 only \yul{FR\_MODULUS} is read by any code in this range,
and only by the modexp probe P1.

\subsection{What this range is for}
\label{sec:deploy:purpose}

The generated verifier is not portable. Its correctness and its liveness both
rest on properties of the host chain that no Solidity type system expresses:

\begin{enumerate}
  \item the \opcode{MCOPY} opcode (EIP-5656, Cancun) must exist, because the
        proof-time point and scratch staging uses it;
  \item the \opcode{MODEXP} precompile at \texttt{0x05} must be priced at or
        below \yul{MODEXP\_GAS} $= 4080$, because every call site forwards
        exactly that bound rather than \yul{gas()};
  \item the EIP-2537 precompiles \texttt{0x0b} (G1ADD), \texttt{0x0c} (G1MSM)
        and \texttt{0x0f} (PAIRING\_CHECK) must exist, must be priced at or below
        the rendered bounds, and must implement the specified \emph{rejection}
        semantics --- G1MSM in particular is the verifier's only subgroup
        validator for absorbed proof commitments, and PAIRING\_CHECK is the sole
        accept gate.
\end{enumerate}

A chain that fails any of these produces a contract that either rejects every
proof (a permanent brick) or, in the PAIRING\_CHECK case, accepts every proof.
The range specified here converts all of those outcomes into a failed deployment
transaction. It is pure fail-fast machinery: nothing in it is reachable after
construction.

\subsection{The typed-error selector constants}
\label{sec:deploy:selectors}

Lines 320--327 declare the eight typed-error selectors as Yul-readable
\yul{uint256} constants, under a banner comment (313--319) stating that each is
$\mathrm{bytes4}(\mathrm{keccak256}(\text{``Name()''}))$ of the corresponding
\yul{error} declared at lines 60--83, and that the pairing is pinned by the
generator test \yul{p4\_error\_selectors\_match\_declared\_errors} in
\texttt{src/lowering/tests.rs}. Only one of the eight, the last, is reachable
from the range specified here; the other seven are consumed by
\yul{verifyProof}, specified elsewhere. They are tabulated in full because the
declarations themselves are in range and because the correctness of a selector
constant is exactly the kind of claim a generated artifact should not be taken on
trust for. I recomputed all eight.

\begin{table}[htbp]
\centering
\caption{Typed-error selectors declared at lines 320--327, recomputed from the error signatures at lines 60--83.}
\label{tab:deploy:selectors}
\small
\begin{tabularx}{\linewidth}{lllX}
\toprule
Constant (line) & Error (line) & Selector & Recomputed \\
\midrule
\yul{ERR\_BAD\_CALLDATA\_SHAPE} (320)       & \revert{BadCalldataShape} (60)      & \texttt{0x1b99e37c} & matches \\
\yul{ERR\_VK\_MISMATCH} (321)               & \revert{VkMismatch} (63)            & \texttt{0xa447d73e} & matches \\
\yul{ERR\_NON\_CANONICAL\_SCALAR} (322)     & \revert{NonCanonicalScalar} (66)    & \texttt{0x77530042} & matches \\
\yul{ERR\_BAD\_POINT\_ENCODING} (323)       & \revert{BadPointEncoding} (69)      & \texttt{0xf27905ec} & matches \\
\yul{ERR\_PRECOMPILE\_FAILED} (324)         & \revert{PrecompileFailed} (71)      & \texttt{0x84e81692} & matches \\
\yul{ERR\_PROOF\_REJECTED} (325)            & \revert{ProofRejected} (73)         & \texttt{0xc3b0d8cd} & matches \\
\yul{ERR\_QUOTIENT\_PROGRAM\_INVALID} (326) & \revert{QuotientProgramInvalid} (77)& \texttt{0x3cc81b89} & matches \\
\yul{ERR\_MEMORY\_LAYOUT\_VIOLATED} (327)   & \revert{MemoryLayoutViolated} (83)  & \texttt{0xc9888d23} & matches \\
\bottomrule
\end{tabularx}
\end{table}

All eight are correct. Note the consequence for the one that matters here: a
transposed selector would not be a soundness bug --- a revert is a revert
whatever four bytes it carries --- but it would silently misroute incident
response, since \revert{MemoryLayoutViolated} is the signal that distinguishes a
build fault from a chain fault (\S\ref{sec:deploy:reverts}).

\subsection{Constants consumed by the deployment path}
\label{sec:deploy:constants}

\begin{table}[htbp]
\centering
\caption{Constants read by lines 344--586, with their declaration sites.}
\label{tab:deploy:constants}
\begin{tabularx}{\linewidth}{llrX}
\toprule
Constant & Line & Value & Role in this range \\
\midrule
\yul{ERR\_MEMORY\_LAYOUT\_VIOLATED} & 327 & \texttt{0xc9888d23} & selector of \revert{MemoryLayoutViolated}, the one typed revert in the probe suite \\
\yul{FR\_MODULUS} & 331 & $r$ & modexp probe modulus and exponent base \\
\yul{G1ADD\_GAS} & 290 & 375 & forwarded to both \texttt{0x0b} probes \\
\yul{G1MSM\_GAS\_1PAIR} & 291 & 12000 & forwarded to the single-pair \texttt{0x0c} probe \\
\yul{PAIRING\_GAS\_2PAIR} & 292 & 102900 & forwarded to all three \texttt{0x0f} probes \\
\yul{MODEXP\_GAS} & 293 & 4080 & forwarded to the \texttt{0x05} probe \\
\yul{G1MSM\_GAS\_SMOKE} & 295 & 525096 & forwarded to the 78-pair worst-case \texttt{0x0c} probe \\
\yul{EXPECTED\_VK\_LENGTH} & 93 & 17025 & constructor VK runtime-length pin \\
\yul{EXPECTED\_VK\_CODEHASH} & 95 & \texttt{0xe68d\ldots cf52} & constructor VK codehash pin \\
\bottomrule
\end{tabularx}
\end{table}

The three EIP-2537 bounds are exactly the schedule of that EIP, so they are
sufficient by construction on a conformant chain:
\[
  \mathrm{gas}_{\mathrm{G1ADD}} = 375,\qquad
  \mathrm{gas}_{\mathrm{G1MSM}}(k) = \left\lfloor\frac{12000\,k\,d_k}{1000}\right\rfloor,\qquad
  \mathrm{gas}_{\mathrm{PAIR}}(k) = 32600\,k + 37700 .
\]
For $k=1$ the discount is $d_1 = 1000$, giving $12000$; for $k=2$ the pairing
formula gives $65200 + 37700 = 102900$. For the worst-case MSM of this artifact,
$k = 78$ and $d_{78} = 561$, giving $78 \cdot 12000 \cdot 561 / 1000 = 525096$
exactly, with no rounding --- which is precisely \yul{G1MSM\_GAS\_SMOKE}.
\yul{MODEXP\_GAS} is deliberately the \emph{larger} of the two live modexp
schedules (EIP-2565 prices this frame at 1360, EIP-7883 at 4080), because
over-forwarding is free on success --- unused gas is refunded to the caller ---
whereas under-forwarding reverts every proof.

\subsection{\texorpdfstring{\yul{require\_eip2537\_precompiles}}{require\_eip2537\_precompiles}: overview}
\label{sec:deploy:overview}

The function is declared \yul{private view} at line 351 and its entire body is a
single \yul{assembly ("memory-safe")} block (line 353) tagged with the phase
marker \yul{// \_\_phase:constructor\_smoke} (line 352). It takes no arguments,
returns nothing, and its only observable effect is reverting. It performs one
opcode check and nine \opcode{STATICCALL}s, in the fixed order of
Algorithm~\ref{alg:deploy:suite}.

\begin{algorithm}[htbp]
\caption{\yul{require\_eip2537\_precompiles}, lines 351--569}
\label{alg:deploy:suite}
\begin{algorithmic}[1]
\Require creation frame; free-memory pointer not yet raised past \texttt{0x1000}
\If{$\mathrm{mload}(\texttt{0x40}) > \texttt{0x1000}$} \textbf{revert} \revert{MemoryLayoutViolated} \EndIf
\State $\mathit{scratch} \gets \texttt{0x1000}$
\State \textbf{P0} \opcode{MCOPY} round trip of the word \texttt{0x1234}
\State \textbf{P1} \texttt{0x05} modexp: $2^{\,r-2} \bmod r$, assert $2 \cdot \text{result} \equiv 1 \pmod r$
\State zero-fill $\mathit{scratch}[\texttt{0x000},\texttt{0x300})$ \Comment{identity encoding for G1 and G2}
\State \textbf{P2} \texttt{0x0b} G1ADD$(\id,\id) = \id$
\State \textbf{P3} \texttt{0x0b} G1ADD$(G,G) = 2G$
\State \textbf{P4} \texttt{0x0c} G1MSM$([2] \cdot G) = 2G$
\State \textbf{P5} \texttt{0x0c} G1MSM$([1] \cdot Q)$ \textbf{must fail}, $Q \in E(\Fp) \setminus \GO$
\State \textbf{P6} \texttt{0x0f} $\pair{G}{G_2}\cdot\pair{-G}{G_2} = 1$
\State \textbf{P7} \texttt{0x0f} $\pair{G}{G_2}\cdot\pair{G}{G_2} \neq 1$
\State zero-fill $\mathit{scratch}[\texttt{0x000},\texttt{0x300})$ again
\State $\mathit{msm\_scratch} \gets \texttt{0xb140}$; zero-fill
       $[\texttt{0xb140}, \texttt{0xe200})$ \Comment{390 iterations}
\State \textbf{P8} \texttt{0x0c} G1MSM over 78 identity/zero terms $= \id$, output to $\mathit{scratch}$
\State \textbf{P9} \texttt{0x0f} $\pair{\id}{\id}\cdot\pair{\id}{\id} = 1$
\end{algorithmic}
\end{algorithm}

\subsection{The free-memory-pointer guard}
\label{sec:deploy:fmp}

\begin{lstlisting}[style=solidity,caption={\texttt{Halo2Verifier.sol} lines 354--369: creation-frame memory guard},label={lst:deploy:fmp}]
// Same free-memory-pointer guard as verifyProof. This body runs in
// the *creation* frame, which the generator's memoryguard test does
// not inspect (it parses the runtime prologue only).
//
// MF-2: typed like the runtime guard, and for a stronger reason.
// The runtime guard turns a bad recompile into a revert on every
// proof; this one turns it into a failed DEPLOYMENT, which is
// where a build fault belongs. The probes below keep bare reverts
// (a chain-capability failure, not a build fault).
if gt(mload(0x40), 0x1000) {
    mstore(0x00, shl(224, ERR_MEMORY_LAYOUT_VIOLATED))
    revert(0x00, 0x04)
}

// Scratch is reused for every runtime-prerequisite probe.
let scratch := 0x1000
\end{lstlisting}

The guard is a strict inequality: $\mathrm{mload}(\texttt{0x40}) = \texttt{0x1000}$
passes. That is the correct boundary, because Solidity's reserved region is the
half-open interval $[\texttt{0x00}, \mathrm{FMP})$ and the lowest address this
function writes is exactly \texttt{0x1000}. The error selector is written at
\texttt{0x00}, Solidity's legal scratch space, and four bytes are returned, so
the failure decodes as \revert{MemoryLayoutViolated}.

Three facts make this guard load-bearing rather than decorative:

\begin{itemize}
  \item It is what justifies the \yul{("memory-safe")} annotation on line 353.
        Solidity's memory-safety contract permits an assembly block to use
        memory located above the free-memory pointer as it stands on entry to
        the block, without updating that pointer. Every write in this function
        is at \texttt{0x1000} or above (the probe scratch) or at \texttt{0xb140}
        or above (the MSM scratch), so once the guard establishes
        $\mathrm{FMP} \le \texttt{0x1000}$ the annotation is sound. Without the
        guard the annotation would be an unchecked assertion about a
        compiler-internal layout.
  \item The value \texttt{0x1000} rendered here is the generator's
        \yul{constructor\_smoke\_scratch\_mptr}, allocated by the same memory
        arena that allocates \mptr{TRANSCRIPT\_MPTR}. It coincides numerically
        with \mptr{TRANSCRIPT\_MPTR} $= \texttt{0x1000}$ (line 119), and with
        \mptr{RETURN\_MPTR} and \mptr{QUOTIENT\_RETURN\_MPTR} at the same
        address (lines 120 and 230), but is a distinct arena allocation with a
        disjoint lifetime.
  \item The runtime guard at line 674 is written as
        \yul{gt(mload(0x40), TRANSCRIPT\_MPTR)} using the named constant,
        whereas this creation-frame guard hardcodes the literal \texttt{0x1000}.
        Both are generator-rendered from their respective arena pointers, so
        they cannot drift within a single render.
\end{itemize}

The header comment at lines 2--11 records why the pragma is pinned to
\texttt{0.8.30}: solc's stack-spill reservation was measured at \texttt{0x8c0} on
0.8.24 and \texttt{0x8e0} on 0.8.26+, both comfortably below \texttt{0x1000},
but that margin is not a property this file controls. The guard converts a
future toolchain that widens the reservation past \texttt{0x1000} from a silently
mis-compiled artifact into a deployment that cannot complete.

\begin{finding}[deploy-1]{Creation-frame memory guard is correct and is the sole justification for the \texttt{memory-safe} annotation --- clean}
\textbf{Severity: Informational.}
\texttt{Halo2Verifier.sol:363--366}. The guard is sound: strict \yul{gt} against
the lowest address the function writes, typed revert, executed before any
absolute-address store. It is also the reason the \yul{("memory-safe")}
annotation on line 353 is legitimate rather than assumed. No defect. The one
cosmetic wrinkle is that the NatSpec for \revert{MemoryLayoutViolated} (lines
78--82, with the declaration at 83) says the fault occurs when the free-memory
pointer starts ``at or above
\mptr{TRANSCRIPT\_MPTR}'', which reads as a non-strict comparison, while both
guards (363 and 674) are strict. The code is right and the prose is loose by one
byte; equality is genuinely safe.
\end{finding}

\subsection{P0: the \texorpdfstring{\opcode{MCOPY}}{MCOPY} probe}
\label{sec:deploy:mcopy}

\begin{lstlisting}[style=solidity,caption={\texttt{Halo2Verifier.sol} lines 371--376: \opcode{MCOPY} availability probe},label={lst:deploy:mcopy}]
// MCOPY must be available because the verifier uses it for
// proof-time point/scratch staging. Execute the opcode here so a
// non-Cancun fork fails during deployment instead of later proofs.
mstore(scratch, 0x1234)
mcopy(add(scratch, 0x20), scratch, 0x20)
if iszero(eq(mload(add(scratch, 0x20)), 0x1234)) { revert(0, 0) }
\end{lstlisting}

One word is stored at \texttt{0x1000}, copied to \texttt{0x1020} by
\opcode{MCOPY}, and read back. The observable behaviour splits three ways:

\begin{enumerate}
  \item On a Cancun-or-later chain the opcode exists and the equality holds; the
        explicit check is redundant with the EVM specification.
  \item On a pre-Cancun chain \texttt{0x5e} is an undefined opcode, so the
        creation frame halts exceptionally, consumes all gas, and the deployment
        fails. The equality check never runs.
  \item The check is only reachable and only informative against a chain that
        \emph{defines} \texttt{0x5e} but implements it incorrectly (a no-op, or
        a different operation). That is the residual case it buys.
\end{enumerate}

The cost of case 2 is a diagnostic one: the deployment fails with
``out of gas'' rather than anything that names \opcode{MCOPY}. That is inherent
to probing an undefined opcode and cannot be improved from inside the contract.

Note that the compiler side of this requirement is enforced separately: solc
emits \texttt{0x5e} for \yul{mcopy} only when compiled with an EVM version of
\texttt{cancun} or later, and refuses otherwise. The probe therefore tests the
\emph{chain}, not the build.

\subsection{P1: the modexp known-answer probe}
\label{sec:deploy:modexp}

\begin{lstlisting}[style=solidity,caption={\texttt{Halo2Verifier.sol} lines 399--407: \texttt{0x05} known-answer probe (leading 21-line rationale comment at 378--398 elided)},label={lst:deploy:modexp}]
mstore(add(scratch, 0x00), 0x20)        // base len
mstore(add(scratch, 0x20), 0x20)        // exp len
mstore(add(scratch, 0x40), 0x20)        // mod len
mstore(add(scratch, 0x60), 2)
mstore(add(scratch, 0x80), sub(FR_MODULUS, 2))
mstore(add(scratch, 0xa0), FR_MODULUS)
if iszero(staticcall(MODEXP_GAS, 0x05, scratch, 0xc0, scratch, 0x20)) { revert(0, 0) }
if iszero(eq(returndatasize(), 0x20)) { revert(0, 0) }
if iszero(eq(mulmod(mload(scratch), 2, FR_MODULUS), 1)) { revert(0, 0) }
\end{lstlisting}

\subsubsection{Frame shape}
The EIP-198 input frame is six words at $\mathit{scratch} + \{0,\texttt{0x20},
\ldots,\texttt{0xa0}\}$: three 32-byte length headers all equal to
\texttt{0x20}, then the base $2$, the exponent $r-2$, and the modulus $r$. Total
input length \texttt{0xc0} $= 192$ bytes; requested output length \texttt{0x20}.
This is bit-for-bit the frame the runtime builds at lines 719, 873 and 924, so
the probe exercises the exact base/exponent/modulus widths that determine the
EIP-2565 and EIP-7883 gas formulas, not a cheaper stand-in.

\subsubsection{The known answer, and why the check is a \texorpdfstring{\yul{mulmod}}{mulmod}}
Since $r$ is prime and $2 \not\equiv 0 \pmod r$, Fermat's little theorem gives
\[
  2^{\,r-2} \equiv 2^{-1} \pmod r ,
\]
which is exactly the inversion the runtime performs in \yul{scalar\_inv} and in
the Lagrange batch inversion. The verification step is
\[
  \mathrm{mulmod}\bigl(\text{result},\,2,\,r\bigr) \stackrel{!}{=} 1 ,
\]
rather than a comparison against a rendered literal $2^{-1} \bmod r$. Three
things this buys:

\begin{itemize}
  \item \textbf{Self-containment.} The generator does not have to render, and a
        reader does not have to trust, a 32-byte magic constant; the assertion is
        checkable by inspection from the definition of a modular inverse.
  \item \textbf{Stub rejection.} A precompile that returns all zeros gives
        $0 \cdot 2 \bmod r = 0 \neq 1$. A precompile that echoes its input gives
        $2 \cdot 2 = 4 \neq 1$ (because $r > 4$). A precompile that returns the
        modulus gives $0$. All three fail.
  \item \textbf{No false negative.} In $\Fr$ the condition $2x \equiv 1$ has the
        unique solution $x = 2^{-1}$, because $\Fr$ is a field. The probe is
        therefore not a weaker checksum than a rendered-constant comparison; it
        cannot reject a correct answer.
\end{itemize}

One precision the code's own comment does not make, and which a reader should
not overstate in the other direction: the check is on the 256-bit \emph{word}
loaded from $\mathit{scratch}$, not on a field element, and
$\mathrm{mulmod}(x, 2, r) = 1$ has exactly \emph{two} solutions in
$[0, 2^{256})$ --- namely
\[
  2^{-1} \bmod r = \texttt{0x39f6d3a9\ldots80000001}
  \quad\text{and}\quad
  (2^{-1} \bmod r) + r = \texttt{0xade47afc\ldots80000002},
\]
since $2^{-1} + 2r > 2^{256} > 2^{-1} + r$. The second word is unreachable from
any EIP-198-conformant implementation, which must return exactly
\yul{mod\_len} bytes holding a value strictly below the modulus, and the probe's
purpose is to reject stubs rather than to certify a non-reducing precompile.
The known-answer property is therefore ``accepts the correct answer and one
non-canonical representation of it'', not ``accepts exactly one word''. This
does not weaken the probe: a precompile that returned $2^{-1} + r$ would also
pass every runtime \yul{scalar\_inv} use, because every consumer feeds the
result straight back into \yul{mulmod}/\yul{addmod} modulo $r$.

\subsubsection{Success shape}
Two assertions guard the call itself. Line 405 checks the \opcode{STATICCALL}
success flag; line 406 checks $\yul{returndatasize()} = \texttt{0x20}$. The
second is not redundant: a \opcode{STATICCALL} to an account with no code
succeeds and copies nothing, so without the size check a chain missing
\texttt{0x05} would fall through to the value check --- which would then read
whatever the input frame left at $\mathit{scratch}$, namely the base-length
header \texttt{0x20} $= 32$, and $32 \cdot 2 = 64 \neq 1$ would still revert.
The size check makes that reasoning unnecessary and makes the failure
unambiguous.

\subsubsection{Assessment}
This probe is the one addition that closes a genuine gap flagged in the code's
own \texttt{MF-1} note (lines 381--389): \texttt{0x05} is mandatory for the
Lagrange batch inversion and for every \yul{scalar\_inv}, two live schedules
price the same frame at 1360 and 4080, and forwarding a bound below the chain's
price does not degrade gracefully --- the fixed forward simply OOGs and
\emph{every} proof reverts \revert{PrecompileFailed}. Because the probe forwards
\yul{MODEXP\_GAS} rather than \yul{gas()}, a chain that repriced modexp above
4080 fails at deployment. That is exactly the right place for that failure.

\subsection{Identity zero-fill and the EIP-2537 encoding}
\label{sec:deploy:identity}

\begin{lstlisting}[style=solidity,caption={\texttt{Halo2Verifier.sol} lines 409--413: identity zero-fill},label={lst:deploy:zerofill}]
// Start the EIP-2537 probes with the identity encoding for G1/G2:
// all-zero padded words. This also clears the modexp frame above.
for { let off := 0 } lt(off, 0x0300) { off := add(off, 0x20) } {
    mstore(add(scratch, off), 0)
}
\end{lstlisting}

The three-line loop body (411--413) is emitted a second time, byte for byte, at
lines 537--539, under a different one-line lead comment (536: ``Restore the
identity encoding for the probes below''); the two comment lines 409--410 have no
counterpart there. Nothing else in the range repeats.

The loop runs $\texttt{0x0300}/\texttt{0x20} = 24$ iterations, zeroing
$[\texttt{0x1000}, \texttt{0x1300})$. \texttt{0x300} $= 768$ bytes is the
two-pair PAIRING\_CHECK frame size, i.e.\ the largest EIP-2537 input any probe
in the low scratch region uses; it strictly covers the \texttt{0x100}-byte G1ADD
frame, the \texttt{0xa0}-byte one-pair G1MSM frame, and the \texttt{0xc0}-byte
modexp frame left over from P1.

Under EIP-2537 an $\Fp$ coordinate is a 64-byte field: 16 zero bytes of padding
followed by the 48-byte big-endian coordinate. A $\GO$ point is therefore
$2 \times 64 = 128$ bytes (\texttt{0x80}, four words: \texttt{x\_hi},
\texttt{x\_lo}, \texttt{y\_hi}, \texttt{y\_lo}), and a $\GT$ point is
$4 \times 64 = 256$ bytes (\texttt{0x100}, eight words:
$x_{c_0}, x_{c_1}, y_{c_0}, y_{c_1}$, each as a hi/lo pair). The point at
infinity in both groups is encoded as all zeros. Hence a zero-filled buffer is a
\emph{valid} EIP-2537 identity input, which is what makes this loop double as
both a frame reset and a probe vector.

\subsection{P2: G1ADD identity}
\label{sec:deploy:g1add-id}

\begin{lstlisting}[style=solidity,caption={\texttt{Halo2Verifier.sol} lines 415--422: \texttt{0x0b} identity probe},label={lst:deploy:g1add-id}]
// G1ADD(identity, identity) -> identity, 128-byte return.
// This catches chains where the precompile is missing or returns a
// non-standard success shape.
if iszero(staticcall(G1ADD_GAS, 0x0b, scratch, 0x0100, scratch, 0x80)) { revert(0, 0) }
if iszero(eq(returndatasize(), 0x80)) { revert(0, 0) }
if or(or(mload(scratch), mload(add(scratch, 0x20))), or(mload(add(scratch, 0x40)), mload(add(scratch, 0x60)))) {
    revert(0, 0)
}
\end{lstlisting}

Input: two $\GO$ points at \texttt{0x1000} and \texttt{0x1080}, both all-zero,
total \texttt{0x100} $= 256$ bytes. Output: \texttt{0x80} bytes written back over
the first input point. Gas forwarded: \yul{G1ADD\_GAS} $= 375$.

The success shape is asserted twice --- flag, then
$\yul{returndatasize()} = \texttt{0x80}$ --- and the value assertion is that all
four returned words are zero, expressed as a four-way \yul{or} that must be
falsy. Here the returndatasize assertion at line 419 is strictly
\emph{load-bearing}, not belt-and-braces: the output buffer was pre-zeroed by
the loop of \S\ref{sec:deploy:identity}, so a chain with no \texttt{0x0b}
precompile would succeed, copy zero bytes, leave the zeros in place, and satisfy
the value check. Line 419 is the only thing that distinguishes ``returned the
identity'' from ``returned nothing''.

\subsection{P3: G1ADD known answer}
\label{sec:deploy:g1add-kat}

\begin{lstlisting}[style=solidity,caption={\texttt{Halo2Verifier.sol} lines 434--450: \texttt{0x0b} known-answer probe $G+G=2G$ (10-line rationale comment at 424--433 elided)},label={lst:deploy:g1add-kat}]
mstore(add(scratch, 0x00), 0x0000000000000000000000000000000017f1d3a73197d7942695638c4fa9ac0f)
mstore(add(scratch, 0x20), 0xc3688c4f9774b905a14e3a3f171bac586c55e83ff97a1aeffb3af00adb22c6bb)
mstore(add(scratch, 0x40), 0x0000000000000000000000000000000008b3f481e3aaa0f1a09e30ed741d8ae4)
mstore(add(scratch, 0x60), 0xfcf5e095d5d00af600db18cb2c04b3edd03cc744a2888ae40caa232946c5e7e1)
mcopy(add(scratch, 0x80), scratch, 0x80)
if iszero(staticcall(G1ADD_GAS, 0x0b, scratch, 0x0100, scratch, 0x80)) { revert(0, 0) }
if iszero(eq(returndatasize(), 0x80)) { revert(0, 0) }
if iszero(and(
    and(
        eq(mload(add(scratch, 0x00)), 0x000000000000000000000000000000000572cbea904d67468808c8eb50a9450c),
        eq(mload(add(scratch, 0x20)), 0x9721db309128012543902d0ac358a62ae28f75bb8f1c7c42c39a8c5529bf0f4e)
    ),
    and(
        eq(mload(add(scratch, 0x40)), 0x00000000000000000000000000000000166a9d8cabc673a322fda673779d8e38),
        eq(mload(add(scratch, 0x60)), 0x22ba3ecb8670e461f73bb9021d5fd76a4c56d9d4cd16bd1bba86881979749d28)
    )
)) { revert(0, 0) }
\end{lstlisting}

The four literals at lines 434--437 are the canonical BLS12-381 $\GO$ generator
$G = (x_G, y_G)$ split into hi/lo halves; \yul{mcopy} then duplicates them into
the second operand slot, so the call computes $G + G$. The expected result at
lines 443--448 is $2G$.

I recomputed all four constants independently from the curve parameters:
$G$ satisfies $y_G^2 = x_G^3 + 4$ over $\Fp$, $[r]G = \id$ (so $G$ is in the
$r$-order subgroup), and the doubling formula applied to $G$ reproduces the
expected words at lines 443--448 exactly. The known answer is correct.

\begin{align*}
  x_G    &= \texttt{0x17f1d3a7\ldots db22c6bb}, &  y_G   &= \texttt{0x08b3f481\ldots46c5e7e1},\\
  x_{2G} &= \texttt{0x0572cbea\ldots29bf0f4e}, &  y_{2G} &= \texttt{0x166a9d8c\ldots79749d28}.
\end{align*}

The comment at 424--433 states the design reason plainly and correctly: every
preceding probe uses $\id$, which is precisely the input an implementation gets
right \emph{without doing curve arithmetic}. A stub that returns its zero-filled
input, or zeros unconditionally, passes P2. P3 is the first vector whose answer a
stub cannot guess.

\subsection{P4: G1MSM known answer}
\label{sec:deploy:g1msm-kat}

\begin{lstlisting}[style=solidity,caption={\texttt{Halo2Verifier.sol} lines 467--484: \texttt{0x0c} known-answer probe $[2]G = 2G$},label={lst:deploy:g1msm-kat}]
// (a) G1MSM known answer: [2]*G == 2G.
mstore(add(scratch, 0x00), 0x0000000000000000000000000000000017f1d3a73197d7942695638c4fa9ac0f)
mstore(add(scratch, 0x20), 0xc3688c4f9774b905a14e3a3f171bac586c55e83ff97a1aeffb3af00adb22c6bb)
mstore(add(scratch, 0x40), 0x0000000000000000000000000000000008b3f481e3aaa0f1a09e30ed741d8ae4)
mstore(add(scratch, 0x60), 0xfcf5e095d5d00af600db18cb2c04b3edd03cc744a2888ae40caa232946c5e7e1)
mstore(add(scratch, 0x80), 2)
if iszero(staticcall(G1MSM_GAS_1PAIR, 0x0c, scratch, 0xa0, scratch, 0x80)) { revert(0, 0) }
if iszero(eq(returndatasize(), 0x80)) { revert(0, 0) }
if iszero(and(
    and(
        eq(mload(add(scratch, 0x00)), 0x000000000000000000000000000000000572cbea904d67468808c8eb50a9450c),
        eq(mload(add(scratch, 0x20)), 0x9721db309128012543902d0ac358a62ae28f75bb8f1c7c42c39a8c5529bf0f4e)
    ),
    and(
        eq(mload(add(scratch, 0x40)), 0x00000000000000000000000000000000166a9d8cabc673a322fda673779d8e38),
        eq(mload(add(scratch, 0x60)), 0x22ba3ecb8670e461f73bb9021d5fd76a4c56d9d4cd16bd1bba86881979749d28)
    )
)) { revert(0, 0) }
\end{lstlisting}

One MSM pair is $128 + 32 = 160$ bytes (\texttt{0xa0}): the padded $\GO$ point
followed by a 32-byte big-endian scalar. The probe writes $G$ and the scalar $2$
and asserts the result equals $2G$, with the same two-step success-shape check.
The forwarded bound \yul{G1MSM\_GAS\_1PAIR} $= 12000$ is the exact one-pair
schedule and is also the bound used at every one-pair runtime site (lines 1288,
4015, 4026, 4075, 4089). \yul{ACC\_RHS\_MSM\_GAS} (declared at line 299, value
$12000$) is numerically the same bound and is forwarded by the multi-line
\yul{staticcall} at lines 1337--1344; that one call passes a runtime-computed
input length \yul{acc\_msm\_len} rather than a literal \texttt{0xa0}, but in this
render exactly one $(\GO, \text{scalar})$ pair is ever appended (lines
1305--1317), so \yul{acc\_msm\_len} is always \texttt{0xa0} and $12000$ is the
correct schedule for it. Both facts are outside this range and covered
elsewhere; they are cited here only to justify the coverage claim of
\S\ref{sec:deploy:absent}.

\subsection{P5: the G1MSM subgroup-rejection probe}
\label{sec:deploy:g1msm-neg}

\begin{lstlisting}[style=solidity,caption={\texttt{Halo2Verifier.sol} lines 486--500: \texttt{0x0c} negative probe},label={lst:deploy:g1msm-neg}]
// (b) G1MSM negative probe. (4, y) satisfies y^2 = x^3 + 4 over Fp
// but is NOT in the r-order subgroup (checked off-chain: r*P != O).
// EIP-2537 requires G1MSM to reject it. This is the one property
// the verifier's deferred-validation strategy depends on and the
// one property no other probe exercises.
//
// Gas is bounded on purpose: a precompile that rejects its input
// consumes everything forwarded to it, so an unbounded `gas()` here
// would burn 63/64 of the deployment gas before the probes below.
mstore(add(scratch, 0x00), 0x0000000000000000000000000000000000000000000000000000000000000000)
mstore(add(scratch, 0x20), 0x0000000000000000000000000000000000000000000000000000000000000004)
mstore(add(scratch, 0x40), 0x000000000000000000000000000000000a989badd40d6212b33cffc3f3763e9b)
mstore(add(scratch, 0x60), 0xc760f988c9926b26da9dd85e928483446346b8ed00e1de5d5ea93e354abe706c)
mstore(add(scratch, 0x80), 1)
if staticcall(200000, 0x0c, scratch, 0xa0, scratch, 0x80) { revert(0, 0) }
\end{lstlisting}

This is the single most security-relevant probe in the range and the only one
with inverted polarity: the guard reverts when the \opcode{STATICCALL}
\emph{succeeds}.

The vector is $Q = (4,\,y_Q)$ with
$y_Q = \texttt{0x0a989bad\ldots4abe706c}$. I verified both claims in the comment
independently: $y_Q^2 \equiv 4^3 + 4 \pmod p$ holds, so $Q \in E(\Fp)$; and
$[r]Q \neq \id$, so $Q \notin \GO$. EIP-2537 requires G1MSM (unlike G1ADD) to
perform a subgroup check on every input point and to fail on a curve point
outside the $r$-order subgroup.

Why this matters for \emph{this} verifier: the transcript absorption path
(\yul{common\_uncompressed\_g1}) performs no curve or subgroup check on absorbed
proof commitments; validation is deferred to the fact that every such point is
eventually fed to G1MSM. If the host chain's \texttt{0x0c} accepted non-subgroup
points, that deferral would silently evaporate. No other probe in the suite tests
rejection behaviour of \texttt{0x0c}.

Two implementation details deserve comment.

\paragraph{The bounded 200{,}000 gas forward.} A rejecting EIP-2537 precompile
raises an exceptional failure and consumes \emph{all} gas forwarded to it. An
unbounded \yul{gas()} would therefore burn $63/64$ of the remaining deployment
budget on this one probe, leaving nothing for P6--P9. The literal $200000$ caps
that loss. The cost is that 200{,}000 gas is consumed unconditionally on every
successful deployment; it is not refunded, because the call is meant to fail.

\paragraph{What a passing P5 does and does not establish.} P5 is a single
negative vector, so it certifies exactly one implication: this chain's
\texttt{0x0c} refuses \emph{this} point at \emph{this} MSM shape. It does not
establish that a general subgroup check runs. An implementation that rejected
$Q$ for an unrelated reason --- a coordinate-range bug, a rejection of $x = 4$
as a special case --- would pass. The vector does close off the most likely
wrong reason: $Q$ is a genuine point of $E(\Fp)$, so ``rejected because
off-curve'' is not available to a conformant implementation. See also
Finding~deploy-8 on the shape restriction.

\paragraph{The OOG/rejection ambiguity, and why it does not bite.} A failing
\opcode{STATICCALL} returns $0$ whether the precompile rejected the input or
merely ran out of the 200{,}000 gas forwarded. In isolation P5 could therefore
pass on a chain that repriced a one-pair G1MSM above 200{,}000 gas. It cannot in
context: P4 (line 473) forwards only \yul{G1MSM\_GAS\_1PAIR} $= 12000$ for a
same-shape one-pair call and runs \emph{first}, so any such chain has already
failed at line 473. The ordering is what makes the composition sound.

\begin{finding}[deploy-2]{G1MSM negative probe is correct and its gas bound is safe only because P4 precedes it}
\textbf{Severity: Observation.}
\texttt{Halo2Verifier.sol:500}. The literal $200000$ is not derived from any
schedule constant, and a failing \opcode{STATICCALL} cannot distinguish
``precompile rejected the point'' from ``precompile ran out of gas''. The probe
is nonetheless sound, because the one-pair positive probe at line 473 forwards
the exact one-pair schedule (12000) and executes earlier, so a chain whose
one-pair G1MSM costs more than 200{,}000 gas has already failed. This is a
correct but \emph{ordering-dependent} argument that is not stated in the code
comment; a reordering of the probe battery in a future generator revision would
silently weaken P5 to a vacuous check. Recording the dependency in the comment
(or deriving the bound from \yul{G1MSM\_GAS\_1PAIR}) would make the invariant
local. Also note the unconditional 200{,}000 gas burned per deployment.
\end{finding}

\subsection{P6 and P7: the pairing known answers}
\label{sec:deploy:pairing-kat}

\begin{lstlisting}[style=solidity,caption={\texttt{Halo2Verifier.sol} lines 502--522: two-pair pairing frame layout},label={lst:deploy:pairing-layout}]
// (c)+(d) Pairing known answers. Lay out [G1 | G2 | G1' | G2] once:
// with G1' = -G the product is 1, with G1' = +G it is not. G2 is
// written literally because the VK payload is not loaded during
// construction.
mstore(add(scratch, 0x000), 0x0000000000000000000000000000000017f1d3a73197d7942695638c4fa9ac0f)
mstore(add(scratch, 0x020), 0xc3688c4f9774b905a14e3a3f171bac586c55e83ff97a1aeffb3af00adb22c6bb)
mstore(add(scratch, 0x040), 0x0000000000000000000000000000000008b3f481e3aaa0f1a09e30ed741d8ae4)
mstore(add(scratch, 0x060), 0xfcf5e095d5d00af600db18cb2c04b3edd03cc744a2888ae40caa232946c5e7e1)
mstore(add(scratch, 0x080), 0x00000000000000000000000000000000024aa2b2f08f0a91260805272dc51051)
mstore(add(scratch, 0x0a0), 0xc6e47ad4fa403b02b4510b647ae3d1770bac0326a805bbefd48056c8c121bdb8)
mstore(add(scratch, 0x0c0), 0x0000000000000000000000000000000013e02b6052719f607dacd3a088274f65)
mstore(add(scratch, 0x0e0), 0x596bd0d09920b61ab5da61bbdc7f5049334cf11213945d57e5ac7d055d042b7e)
mstore(add(scratch, 0x100), 0x000000000000000000000000000000000ce5d527727d6e118cc9cdc6da2e351a)
mstore(add(scratch, 0x120), 0xadfd9baa8cbdd3a76d429a695160d12c923ac9cc3baca289e193548608b82801)
mstore(add(scratch, 0x140), 0x000000000000000000000000000000000606c4a02ea734cc32acd2b02bc28b99)
mstore(add(scratch, 0x160), 0xcb3e287e85a763af267492ab572e99ab3f370d275cec1da1aaa9075ff05f79be)
mstore(add(scratch, 0x180), 0x0000000000000000000000000000000017f1d3a73197d7942695638c4fa9ac0f)
mstore(add(scratch, 0x1a0), 0xc3688c4f9774b905a14e3a3f171bac586c55e83ff97a1aeffb3af00adb22c6bb)
mstore(add(scratch, 0x1c0), 0x00000000000000000000000000000000114d1d6855d545a8aa7d76c8cf2e21f2)
mstore(add(scratch, 0x1e0), 0x67816aef1db507c96655b9d5caac42364e6f38ba0ecb751bad54dcd6b939c2ca)
mcopy(add(scratch, 0x200), add(scratch, 0x80), 0x100)
\end{lstlisting}

\begin{table}[htbp]
\centering
\caption{Two-pair PAIRING\_CHECK frame built at lines 506--522, relative to \texttt{scratch} $=$ \texttt{0x1000}.}
\label{tab:deploy:pairing-frame}
\begin{tabular}{llll}
\toprule
Offset range & Size & Content & Written by \\
\midrule
\texttt{0x000}--\texttt{0x080} & \texttt{0x80} & $\GO$ point $G$              & lines 506--509 \\
\texttt{0x080}--\texttt{0x180} & \texttt{0x100} & $\GT$ generator $G_2$        & lines 510--517 \\
\texttt{0x180}--\texttt{0x200} & \texttt{0x80} & $\GO$ point $-G$ (P6) / $G$ (P7) & lines 518--521, 530--531 \\
\texttt{0x200}--\texttt{0x300} & \texttt{0x100} & copy of $G_2$                & line 522 (\yul{mcopy}) \\
\midrule
\texttt{0x300}--\texttt{0x320} & \texttt{0x20} & output word                   & precompile \\
\bottomrule
\end{tabular}
\end{table}

Each PAIRING\_CHECK pair is $128 + 256 = 384$ bytes; two pairs are
$768 = \texttt{0x300}$, which is the input length passed at lines 525 and 532.
The $\GT$ words at \texttt{0x080}--\texttt{0x180} are the canonical BLS12-381
generator $G_2$ in EIP-2537 order $(x_{c_0}, x_{c_1}, y_{c_0}, y_{c_1})$; they
are written as literals rather than read from the verifying key because the VK
payload is not loaded during construction (the runtime loads it via
\yul{extcodecopy} at line 1393). Line 522 duplicates $G_2$ into the second pair
with a single \yul{mcopy}, which also means P6/P7 exercise \opcode{MCOPY} a
second time.

I verified that the word pair at lines 520--521 is exactly $p - y_G$, i.e.\ the
$y$-coordinate of $-G$.

\begin{lstlisting}[style=solidity,caption={\texttt{Halo2Verifier.sol} lines 524--534: the two pairing assertions},label={lst:deploy:pairing-kat}]
// (c) e(G, G2) * e(-G, G2) == 1.
if iszero(staticcall(PAIRING_GAS_2PAIR, 0x0f, scratch, 0x0300, add(scratch, 0x300), 0x20)) { revert(0, 0) }
if iszero(eq(returndatasize(), 0x20)) { revert(0, 0) }
if iszero(eq(mload(add(scratch, 0x300)), 1)) { revert(0, 0) }

// (d) e(G, G2) * e(G, G2) != 1. Flip the second G1 back to +G.
mstore(add(scratch, 0x1c0), 0x0000000000000000000000000000000008b3f481e3aaa0f1a09e30ed741d8ae4)
mstore(add(scratch, 0x1e0), 0xfcf5e095d5d00af600db18cb2c04b3edd03cc744a2888ae40caa232946c5e7e1)
if iszero(staticcall(PAIRING_GAS_2PAIR, 0x0f, scratch, 0x0300, add(scratch, 0x300), 0x20)) { revert(0, 0) }
if iszero(eq(returndatasize(), 0x20)) { revert(0, 0) }
if iszero(iszero(mload(add(scratch, 0x300)))) { revert(0, 0) }
\end{lstlisting}

The two assertions are:
\begin{align*}
  \textbf{P6:}\quad & \pair{G}{G_2} \cdot \pair{-G}{G_2} = \pair{G}{G_2} \cdot \pair{G}{G_2}^{-1} = 1 &&\Rightarrow \text{output word} = 1,\\
  \textbf{P7:}\quad & \pair{G}{G_2} \cdot \pair{G}{G_2} = \pair{G}{G_2}^{2} \neq 1 &&\Rightarrow \text{output word} = 0 .
\end{align*}
P7 holds because $\pair{G}{G_2}$ has order $r$ in $\GT$ and $r$ is odd, so
squaring cannot reach the identity.

Note the output buffer choice: the result is written to
$\mathit{scratch} + \texttt{0x300} = \texttt{0x1300}$, \emph{outside} the input
region, so the frame survives and P7 needs only to rewrite the two
$y$-coordinate words at \texttt{0x1c0}/\texttt{0x1e0}. This is the reason P6 and
P7 share one 22-line layout instead of two. The word at \texttt{0x1300} is
never read again after line 534: the restore loop at 537--539 covers only
$[\texttt{0x1000}, \texttt{0x1300})$, P8 writes its output at \texttt{0x1000},
and P9 both reads its input from and writes its output to \texttt{0x1000}. The
stale $0$ left there by P7 is therefore harmless.

Failure-closure of P7 deserves an explicit argument, because its expected value
is $0$ and an untouched buffer also reads $0$. Two independent mechanisms make
it non-vacuous. First, line 533 asserts
$\yul{returndatasize()} = \texttt{0x20}$, so a precompile that returns nothing is
rejected regardless of the buffer contents. Second, even without that check the
buffer at \texttt{0x1300} still holds the value $1$ written by P6, so a
no-op precompile would leave $1$ and line 534 would revert. The probe fails
closed both ways.

\begin{finding}[deploy-3]{The rejection-side probes for \texttt{0x0c} and \texttt{0x0f} are present and correctly composed --- clean}
\textbf{Severity: Informational.}
\texttt{Halo2Verifier.sol:486--534}. The suite tests both precompiles that decide
acceptance in \emph{both} directions: P4/P6 confirm they accept and compute
correctly, P5/P7 confirm they reject. That closes the two failure modes that
would be catastrophic rather than merely bricking --- a \texttt{0x0c} that skips
subgroup checks (silently disabling the verifier's deferred point validation)
and a \texttt{0x0f} that always returns $1$ (accepting every proof). All four
known answers were recomputed independently against the BLS12-381 parameters and
match. The composition order (accept-probe before reject-probe, for each
precompile) is what removes the OOG/rejection ambiguity in the reject-probes. No
defect.
\end{finding}

\subsection{P8: the worst-case G1MSM smoke bound}
\label{sec:deploy:g1msm-smoke}

\begin{lstlisting}[style=solidity,caption={\texttt{Halo2Verifier.sol} lines 541--559: full-size MSM probe},label={lst:deploy:g1msm-smoke}]
// Worst-case generated G1MSM with all identity/zero terms ->
// identity, 128-byte return. This exercises the largest MSM input
// LENGTH rendered by this verifier instead of only a one-pair
// smoke call, proving the target chain's precompile accepts the
// full-size input. It runs in the creation frame at its own
// scratch base, so it does not (and cannot) pre-expand the
// runtime call frame's memory -- constructor memory is discarded;
// only the input size coverage carries over.
let msm_scratch := 0xb140
for { let off := 0 } lt(off, 0x30c0) { off := add(off, 0x20) } {
    mstore(add(msm_scratch, off), 0)
}
// The production verifier uses G1MSM both for commitments and as
// the subgroup validator for absorbed proof points.
if iszero(staticcall(G1MSM_GAS_SMOKE, 0x0c, msm_scratch, 0x30c0, scratch, 0x80)) { revert(0, 0) }
if iszero(eq(returndatasize(), 0x80)) { revert(0, 0) }
if or(or(mload(scratch), mload(add(scratch, 0x20))), or(mload(add(scratch, 0x40)), mload(add(scratch, 0x60)))) {
    revert(0, 0)
}
\end{lstlisting}

\subsubsection{The bound}
The input length \texttt{0x30c0} $= 12480$ bytes is
\[
  12480 = 78 \times 160 = 78 \times \texttt{0xa0},
\]
i.e.\ exactly 78 MSM pairs. This is not an arbitrary large number: it is
bit-for-bit the input length of the fused final MSM the runtime issues at line
3999, \yul{staticcall(525096, 0x0c, PCS\_FINAL\_MSM\_MPTR, 0x30c0, FINAL\_COM\_MPTR, 0x80)}.
The gas bound matches too, since
$\lfloor 78 \cdot 12000 \cdot 561 / 1000 \rfloor = 525096 = \yul{G1MSM\_GAS\_SMOKE}$.
The zero-fill loop runs $\texttt{0x30c0}/\texttt{0x20} = 390$ iterations over
$[\texttt{0xb140}, \texttt{0xe200})$.

\subsubsection{What it proves and what it does not}
It proves that the target chain's \texttt{0x0c} accepts an input of the largest
\emph{length} this artifact will ever submit, at the exact gas the runtime will
forward. That rules out an implementation with a lower input-size ceiling, and
rules out a discount table that prices $k=78$ above 525{,}096. The gas half of
that claim is exact rather than probabilistic: EIP-2537 prices G1MSM purely from
the pair count, so an all-identity input costs exactly what 78 arbitrary points
would.

The \emph{arithmetic} half does not carry over. Every one of the 78 terms is
$([0], \id)$, which is precisely the input class the code's own comment at
lines 426--433 identifies as the one an implementation answers correctly
``without doing any curve arithmetic'' --- and, in particular, without running a
subgroup check on anything. P8 therefore validates the input length and the
price of the $k = 78$ path; it validates neither its arithmetic nor its
input validation. See Finding~deploy-8.

It does not, and cannot, pre-warm memory for the runtime: the constructor's
memory is discarded when the creation frame ends. The comment at lines 545--548
says this explicitly and correctly.

The scratch base is \texttt{0xb140}, which is \mptr{SELECTOR\_ACC\_MPTR}
(line 229) and \mptr{BATCH\_INV\_SCRATCH\_MPTR} (line 235) in the runtime layout
--- a different address from \mptr{PCS\_FINAL\_MSM\_MPTR} $= \texttt{0xb280}$
(line 204), which is where the runtime builds the 78-pair frame this probe
imitates. The reuse is harmless because it is a different memory phase in a
different frame; the generator's arena is what enforces that no two live regions
collide, as the comment block at lines 104--118 records. High-water memory for the whole
constructor is therefore $\texttt{0xb140} + \texttt{0x30c0} = \texttt{0xe200}
= 57856$ bytes $= 1808$ words.

\subsubsection{Implicit dependency on the identity result}
The MSM output is written to $\mathit{scratch} = \texttt{0x1000}$, overwriting
the first \texttt{0x80} bytes of the region that the zero-fill loop at 537--539
prepared for P9. Lines 557--559 assert that all four output words are zero. That
assertion does double duty: it validates P8, \emph{and} it re-establishes the
precondition on the only sub-range P8 disturbed, $[\texttt{0x1000},
\texttt{0x1080})$. The remaining $[\texttt{0x1080}, \texttt{0x1300})$ is
untouched between line 539 and line 565, so the full \texttt{0x300}-byte
all-identity pairing frame P9 submits is zero on every path that reaches it.
The dependency is real but is discharged by a check that is already present ---
note that if line 557's polarity were ever inverted, P9 would silently degrade
into a pairing check over whatever $\GO$ point P8 returned.

\subsection{P9: the identity pairing probe}
\label{sec:deploy:pairing-id}

\begin{lstlisting}[style=solidity,caption={\texttt{Halo2Verifier.sol} lines 561--567: two-pair identity pairing probe},label={lst:deploy:pairing-id}]
// PAIRING_CHECK([(identity_g1, identity_g2), (identity_g1, identity_g2)])
// -> true, 32-byte return. This matches the runtime two-pair KZG
// pairing input size and catches absent pairing precompiles,
// short return data, and obviously incompatible semantics.
if iszero(staticcall(PAIRING_GAS_2PAIR, 0x0f, scratch, 0x0300, scratch, 0x20)) { revert(0, 0) }
if iszero(eq(returndatasize(), 0x20)) { revert(0, 0) }
if iszero(eq(mload(scratch), 1)) { revert(0, 0) }
\end{lstlisting}

Input \texttt{0x300} bytes of zeros at \texttt{0x1000}, i.e.\ two
$(\id_{\GO}, \id_{\GT})$ pairs; the empty product is $1$, so the expected output
word is $1$. The output overwrites the head of its own input, which is safe
because the precompile has already consumed the input. This is the same
two-pair frame shape and the same forwarded bound (\yul{PAIRING\_GAS\_2PAIR})
the runtime uses at line 980.

Note the asymmetry with P8: here the value assertion is
$\mathrm{mload}(\mathit{scratch}) = 1$ against a pre-zeroed buffer, so an absent
precompile fails on the value check alone; line 566 is redundant defence rather
than load-bearing. The reverse held for P2 and P8, where the expected value was
zero and line 419 / line 556 were the only discriminators. The suite is
consistent in always asserting both, which is the right policy: it removes the
need to reason case-by-case about which check carries the weight.

\subsection{Probe summary}
\label{sec:deploy:summary}

\begin{table}[htbp]
\centering
\caption{All ten deployment probes, lines 371--567. \texttt{rds} abbreviates \yul{returndatasize()}. \emph{scratch} $=$ \texttt{0x1000}.}
\label{tab:deploy:probes}
\small
\begin{tabularx}{\linewidth}{llllllX}
\toprule
Id & Lines & Target & Gas fwd. & In (ptr, len) & Out (ptr, len) & Assertions \\
\midrule
P0 & 374--376 & \opcode{MCOPY} & --- & --- & --- & word round-trips as \texttt{0x1234} \\
P1 & 399--407 & \texttt{0x05} & 4080 & \texttt{0x1000}, \texttt{0xc0} & \texttt{0x1000}, \texttt{0x20} & ok; \texttt{rds}=\texttt{0x20}; $2x \equiv 1$ \\
P2 & 418--422 & \texttt{0x0b} & 375 & \texttt{0x1000}, \texttt{0x100} & \texttt{0x1000}, \texttt{0x80} & ok; \texttt{rds}=\texttt{0x80}; 4 words zero \\
P3 & 434--450 & \texttt{0x0b} & 375 & \texttt{0x1000}, \texttt{0x100} & \texttt{0x1000}, \texttt{0x80} & ok; \texttt{rds}=\texttt{0x80}; $=2G$ \\
P4 & 468--484 & \texttt{0x0c} & 12000 & \texttt{0x1000}, \texttt{0xa0} & \texttt{0x1000}, \texttt{0x80} & ok; \texttt{rds}=\texttt{0x80}; $=2G$ \\
P5 & 495--500 & \texttt{0x0c} & 200000 & \texttt{0x1000}, \texttt{0xa0} & \texttt{0x1000}, \texttt{0x80} & \textbf{must fail} \\
P6 & 506--527 & \texttt{0x0f} & 102900 & \texttt{0x1000}, \texttt{0x300} & \texttt{0x1300}, \texttt{0x20} & ok; \texttt{rds}=\texttt{0x20}; $=1$ \\
P7 & 530--534 & \texttt{0x0f} & 102900 & \texttt{0x1000}, \texttt{0x300} & \texttt{0x1300}, \texttt{0x20} & ok; \texttt{rds}=\texttt{0x20}; $=0$ \\
P8 & 549--559 & \texttt{0x0c} & 525096 & \texttt{0xb140}, \texttt{0x30c0} & \texttt{0x1000}, \texttt{0x80} & ok; \texttt{rds}=\texttt{0x80}; 4 words zero \\
P9 & 565--567 & \texttt{0x0f} & 102900 & \texttt{0x1000}, \texttt{0x300} & \texttt{0x1000}, \texttt{0x20} & ok; \texttt{rds}=\texttt{0x20}; $=1$ \\
\bottomrule
\end{tabularx}
\end{table}

\subsection{Precompiles that are \emph{not} probed}
\label{sec:deploy:absent}

There is no G2ADD (\texttt{0x0d}) probe, no G2MSM (\texttt{0x0e}) probe, no
map-to-curve (\texttt{0x10}, \texttt{0x11}) probe, and no probe for any of the
pre-Byzantium precompiles.

This is correct and complete rather than an omission. A scan of every
\yul{staticcall} in \texttt{Halo2Verifier.sol} shows exactly four distinct
callee addresses across the whole contract:

\begin{table}[htbp]
\centering
\caption{Every precompile the artifact calls, and its probe coverage.}
\label{tab:deploy:coverage}
\small
\begin{tabularx}{\linewidth}{lllX}
\toprule
Address & Runtime call sites & Bounds used at runtime & Probe coverage \\
\midrule
\texttt{0x05} modexp & 719, 873, 924 & \yul{MODEXP\_GAS} & P1 (same bound, same frame) \\
\texttt{0x0b} G1ADD  & 4020, 4030, 4080, 4094 & \yul{G1ADD\_GAS} & P2, P3 \\
\texttt{0x0c} G1MSM  & 1288, 1337, 3999, 4015, 4026, 4075, 4089 & \yul{G1MSM\_GAS\_1PAIR}, \yul{ACC\_RHS\_MSM\_GAS} ($=12000$), literal $525096$ & P4, P5 (1-pair), P8 (78-pair, identity terms only) \\
\texttt{0x0f} PAIRING & 980 & \yul{PAIRING\_GAS\_2PAIR} & P6, P7, P9 \\
\bottomrule
\end{tabularx}
\end{table}

The NatSpec claim at lines 346--350 --- that ``the probes forward the same exact
gas bounds the runtime uses, for every precompile it calls'' --- therefore holds
for this fixture: every distinct forwarded bound in the runtime ($375$, $4080$,
$12000$, $102900$, $525096$) appears as a forwarded bound in the probe suite.
The converse does not hold, and does not need to: P5's $200000$ (line 500) is a
probe-only bound with no runtime counterpart, for the reason given in
\S\ref{sec:deploy:g1msm-neg}. Probing \texttt{0x0d} or \texttt{0x0e} would only
add deployment gas for capabilities the artifact never exercises.

The discount factor $d_{78} = 561$ used above is not merely back-derived from the
rendered constant: it is entry $78$ of the generator's \yul{G1\_MSM\_DISCOUNT}
table (\texttt{src/lowering/layout/mod.rs}, lines 179--189), whose
\yul{g1msm\_gas} helper computes $\lfloor k \cdot 12000 \cdot d_k / 1000\rfloor$
at lines 192--200. Reading $d_{78} = 561$ from that table and evaluating the
formula reproduces $525096$ exactly.

\begin{finding}[deploy-4]{The 78-pair runtime MSM hardcodes its gas bound while the matching probe uses the named constant}
\textbf{Severity: Observation.}
\texttt{Halo2Verifier.sol:555} forwards \yul{G1MSM\_GAS\_SMOKE}; the single
runtime call site with that shape, \texttt{Halo2Verifier.sol:3999} (outside this
range, covered elsewhere), forwards the bare literal \texttt{525096}. The two
values are equal in this fixture --- I confirmed both are
$\lfloor 78 \cdot 12000 \cdot 561/1000\rfloor$ --- so the probe is faithful as
rendered. The observation is a coupling one: the constant declared at line 295
carries the comment ``Exact cost of the deployment-time worst-case G1MSM smoke
probe'', yet it is also, silently, the runtime's bound. A generator change that
adjusted \yul{G1MSM\_GAS\_SMOKE} without touching the emitter that produces line
3999 would leave the deployment probe validating a bound the runtime does not
use, and the artifact would deploy cleanly on a chain where every proof OOGs.
The runtime site is not undocumented --- line 3998 carries the comment
``exact EIP-2537 G1MSM cost for 78 pair(s)'' --- but a comment is not a
coupling. Emitting line 3999 through the same named constant would make it
mechanical.
\end{finding}

\begin{finding}[deploy-8]{The subgroup-rejection probe covers only the one-pair MSM shape, while the runtime's deferred point validation happens in the 78-pair MSM}
\textbf{Severity: Observation.}
\texttt{Halo2Verifier.sol:495--500} (P5) submits its non-subgroup vector as a
\emph{one-pair} G1MSM: input length \texttt{0xa0}, one point, scalar $1$. But the
call site on which the verifier's deferred point-validation strategy actually
rests for absorbed proof commitments is the fused final MSM at
\texttt{Halo2Verifier.sol:3999} (outside this range, covered elsewhere), whose
input is \texttt{0x30c0} bytes $= 78$ pairs assembled from the decoded proof
points and challenge scalars by the store sequence at lines 3836--3996. The probe that exercises \emph{that} shape, P8 at
lines 549--559, feeds 78 all-zero terms --- the identity --- which is exactly the
input class the contract's own comment at lines 426--433 calls out as the one an
implementation answers correctly without doing curve arithmetic, and hence
without running a subgroup check on anything.

The suite therefore establishes two propositions that do not compose into the one
it needs: (i) \texttt{0x0c} rejects a non-subgroup point when $k = 1$, and
(ii) \texttt{0x0c} accepts a 12{,}480-byte input at 525{,}096 gas. A host
implementation whose single-scalar fast path validates subgroup membership but
whose multi-scalar (Pippenger-style) path does not --- a plausible shape for such
a bug, since the two paths are usually separate code --- passes the entire
battery, and the verifier then accepts absorbed commitments outside $\GO$ at
proof time. This is a \emph{coverage} observation about the probe suite, not a
defect in the verifier's own logic: the contract is correct against a conformant
EIP-2537 chain, and P5 does close the case that motivated it. Closing the gap
costs one additional \opcode{STATICCALL} --- cheapest as a two-pair negative
vector (\texttt{0x140} input, one non-subgroup point, one identity), which
exercises the multi-scalar path for roughly 21{,}000 gas rather than repeating
the full 78-pair frame.
\end{finding}

\subsection{Revert policy: bare versus typed}
\label{sec:deploy:reverts}

The deployment path uses three distinct failure encodings, and the split is
deliberate. The contract's own taxonomy comment (lines 48--57, banner-delimited)
states the rule at lines 53--56:
``Constructor smoke probes intentionally keep bare reverts --- they report a
chain-capability failure, and the deployment transaction identifies itself. The
one constructor exception is the memory-layout guard (MF-2), which reports a
BUILD fault and is typed on both paths.''

\begin{table}[htbp]
\centering
\caption{Revert encodings on the deployment path.}
\label{tab:deploy:reverts}
\small
\begin{tabularx}{\linewidth}{llrX}
\toprule
Encoding & Sites & Count & Rationale as implemented \\
\midrule
\revert{MemoryLayoutViolated} (4-byte selector) & 363--366 & 1 & build fault: the artifact was compiled with a toolchain whose spill reservation collides with the generated layout \\
bare \yul{revert(0, 0)} & 376; 405--407; 418--422; 439--450; 473--484; 500; 525--527; 532--534; 555--559; 565--567 & 26 & chain-capability fault: the bytecode is fine, the host chain is not \\
\yul{require(..., "invalid vk")} $\to$ \texttt{Error(string)} & 580--584 & 1 & dependency-pin fault \\
\bottomrule
\end{tabularx}
\end{table}

\subsubsection{Assessment of the bare reverts}
\textbf{What it buys.} Bare reverts are the smallest possible failure encoding:
two \yul{PUSH1 0} and a \yul{REVERT}, versus a \yul{PUSH4}, \yul{SHL},
\yul{MSTORE}, and a longer \yul{REVERT} for a typed error. Across 26 sites the
saving is on the order of a hundred creation-code bytes. That matters more than
usual here: the header comment at lines 8--11 records that this artifact's
\emph{runtime} is already at or over the EIP-170 24{,}576-byte limit for
non-pinned (compiler, \texttt{--optimize-runs}) pairs, and the probe function is
\yul{private} and called only from the constructor, so solc emits it into the
creation object only. Creation-object bytes still count against EIP-3860's
49{,}152-byte initcode limit and against initcode gas.

\textbf{What it costs.} A deployment that fails inside the probe suite returns
zero bytes of revert data. The operator learns only that construction failed;
they cannot tell whether the chain lacks \opcode{MCOPY}, repriced modexp,
lacks \texttt{0x0f}, or --- most alarmingly --- has a \texttt{0x0c} that accepts
non-subgroup points, without re-running the transaction under a tracer. Since
all 26 sites are in one function called once, the deployment transaction does
identify the \emph{function}, but not the \emph{probe}.

\textbf{Verdict.} The choice is defensible and is correctly and consistently
applied. A cheap middle ground exists and is worth noting: returning a one-byte
probe index (\yul{mstore8(0x00, i)} then \yul{revert(0x00, 0x01)}) would cost
roughly two extra bytes per site and would name the failing probe exactly, which
is a materially better incident signal for the price. That is an improvement, not
a defect.

\subsubsection{Assessment of the typed exception}
Typing only the memory-layout guard is the right discrimination. Its failure mode
is categorically different from the others: it means \emph{this bytecode} is
wrong, not \emph{this chain} is wrong, and the remedy is to rebuild from the
pinned toolchain rather than to deploy elsewhere. Making that one signal
unambiguous while leaving the capability probes bare correctly matches encoding
cost to diagnostic value. The runtime guard at line 674 carries the identical
selector, so a build fault produces the same error whether it is caught at
deployment or at proof time.

\begin{finding}[deploy-5]{Probe failures carry no revert data identifying which capability is missing}
\textbf{Severity: Informational.}
\texttt{Halo2Verifier.sol:376--567}, 26 bare \yul{revert(0, 0)} sites. This is a
documented, deliberate trade (lines 53--56, 361--362) and is not a
vulnerability: every probe fails closed and deployment aborts. It is an
incident-response cost. The failure that matters most --- P5, the G1MSM subgroup
rejection probe at line 500 --- is exactly the one an operator would most want
named in the revert data, because it is the only probe whose failure indicates a
chain on which the verifier's deferred point-validation strategy is unsound
rather than merely a chain on which it will not run. A one-byte probe index
would close this at negligible cost.
\end{finding}

\subsection{Deployment cost}
\label{sec:deploy:gas}

The probe suite's cost is dominated by forwarded precompile gas:

\[
  \underbrace{4080}_{\text{P1}}
  + \underbrace{2 \times 375}_{\text{P2,P3}}
  + \underbrace{12000}_{\text{P4}}
  + \underbrace{200000}_{\text{P5}}
  + \underbrace{3 \times 102900}_{\text{P6,P7,P9}}
  + \underbrace{525096}_{\text{P8}}
  = 1{,}050{,}626 .
\]

On a conformant EIP-7883 chain every one of these is charged in full: the eight
succeeding calls each cost exactly their scheduled price (which is what is
forwarded), and P5 burns its entire 200{,}000-gas forward because a rejecting
EIP-2537 precompile consumes all gas supplied. On a pre-Osaka EIP-2565 chain P1
charges 1360 instead of 4080, giving $1{,}047{,}906$. Adding nine
\opcode{STATICCALL} base costs ($900$), the memory expansion to 1808 words
\[
  3 \cdot 1808 + \left\lfloor \frac{1808^2}{512} \right\rfloor = 5424 + 6384 = 11808 ,
\]
and roughly 438 zero-fill loop iterations ($24 + 24 + 390$), the suite costs on
the order of $1.06$--$1.08$ million gas. The \opcode{STATICCALL} base figure
assumes all four callee addresses are warm; if a client does not seed the
EIP-2537 addresses \texttt{0x0b}, \texttt{0x0c} and \texttt{0x0f} into
EIP-2929's \yul{accessed\_addresses} at transaction start the way the
pre-Byzantium precompiles are, first touch of each costs $2600$ rather than
$100$, adding at most $7500$ --- immaterial against the forwarded total, which is
why the figure is quoted as a range. Against a full deployment --- creation plus a
$\sim$24.5\,KB code deposit at 200 gas per byte, i.e.\ around 5 million gas ---
that is a $\sim$20\% surcharge paid once. Given what it detects, the trade is
clearly favourable; the one line item a deployer should recognise as pure,
non-refundable overhead is P5's 200{,}000 gas.

The forwarded amounts also interact with EIP-150's $63/64$ rule: the largest
single forward, P8's 525{,}096, requires at least
$\lceil 525096 \cdot 64/63 \rceil = 533{,}431$ gas remaining at that point. With
the suite running first in the constructor, before the code deposit is charged,
that is never the binding constraint.

\subsection{The constructor}
\label{sec:deploy:ctor}

\begin{lstlisting}[style=solidity,caption={\texttt{Halo2Verifier.sol} lines 572--586: the constructor},label={lst:deploy:ctor}]
/// @notice Create a verifier pinned to a generated verifying key.
/// @dev Checks MCOPY/EIP-2537 availability and verifies the VK runtime before storing its address.
/// @param authorizedVk Address of the generated `Halo2VerifyingKey` runtime.
constructor(address authorizedVk) {
    // Embedded quotient path: only the external VK runtime needs to be
    // pinned, but the runtime opcode/precompile prerequisites are still
    // mandatory.
    require_eip2537_precompiles();
    require(
        authorizedVk.code.length == EXPECTED_VK_LENGTH
            && authorizedVk.codehash == EXPECTED_VK_CODEHASH,
        "invalid vk"
    );
    AUTHORIZED_VK = authorizedVk;
}
\end{lstlisting}

The constructor performs exactly three actions, in order:

\begin{enumerate}
  \item calls \yul{require\_eip2537\_precompiles()} (line 579), so chain
        capability is settled before anything else;
  \item pins the verifying-key runtime by \emph{both} its length and its
        codehash (lines 580--584);
  \item assigns the validated address to the \yul{public immutable}
        \yul{AUTHORIZED\_VK} (line 585).
\end{enumerate}

The ordering has a cost consequence worth stating: the roughly $1.05$ million
gas of \S\ref{sec:deploy:gas} is spent \emph{before} the cheap VK pin is
evaluated, so a deployment that merely passes the wrong \yul{authorizedVk}
address still burns the whole probe suite. Reversing the two would make the
common operator error cheap to discover, at the price of letting a wrong-chain
deployment be diagnosed as a VK error. Neither ordering is wrong; the rendered
one prioritises the diagnosis that is harder to reproduce off-chain.

There is no other constructor logic: no owner, no pause flag, no event, no
second dependency. This artifact takes the embedded-quotient generator path, so
the split-evaluator constructor variant --- which additionally pins
\yul{EXPECTED\_QUOTIENT\_LENGTH} and \yul{EXPECTED\_QUOTIENT\_CODEHASH} --- is
not rendered, and no such constants appear anywhere in the file.

\subsubsection{The VK pin, verified}
\begin{align*}
  \yul{EXPECTED\_VK\_PAYLOAD\_LENGTH} &= 17024 = \texttt{0x4280} &&\text{(line 92)},\\
  \yul{EXPECTED\_VK\_LENGTH}          &= 17025 = \texttt{0x4281} &&\text{(line 93)},\\
  \yul{EXPECTED\_VK\_CODEHASH}        &= \texttt{0xe68d89362065c8b7107055774f1e045b69bc3bffde4704541c5bc5c91c94cf52} &&\text{(lines 94--95)}.
\end{align*}
The off-by-one is the leading \opcode{INVALID} byte: the deployed
\texttt{Halo2VerifyingKey} runtime is $\texttt{0xfe} \,\|\, \text{payload}$, so
the runtime is one byte longer than the payload the verifier later
\yul{extcodecopy}s from offset \texttt{0x01} (line 1393).

I verified both pins against the companion fixture rather than taking them on
trust. \texttt{Halo2VerifyingKey.sol} builds its runtime from a single
\yul{mstore8(runtime, 0xfe)} plus 532 contiguous 32-byte \yul{mstore} literals at
offsets \texttt{0x0000}--\texttt{0x4260}, and returns
\yul{return(runtime, 0x4281)} (line 692 of that file). Reconstructing those
$1 + 532 \cdot 32 = 17025$ bytes and hashing them gives
\[
  \mathrm{keccak256}(\text{VK runtime}) = \texttt{0xe68d8936\ldots1c94cf52},
\]
matching \yul{EXPECTED\_VK\_CODEHASH} exactly, and the length matches
\yul{EXPECTED\_VK\_LENGTH} exactly. Both constants in this range are correct for
the fixture VK. (The VK contract's own header notes that its runtime is pure
returned data, hence compiler-independent, which is why this reconstruction is
valid without recompiling.)

\subsubsection{Soundness of the pin}
Three properties make this construction sound despite its brevity:

\begin{itemize}
  \item \textbf{The codehash subsumes everything else.} \yul{.codehash} compiles
        to \opcode{EXTCODEHASH}, which is $\mathrm{keccak256}$ of the full
        deployed runtime. Fixing it fixes every byte: the \opcode{INVALID}
        prefix, the \yul{vk\_digest}, $n$, $\omega$, $\omega^{-1}$, the $\GO$ and
        $\GT$ base points, the quotient program, and every fixed and permutation
        commitment. The length check at line 581 is strictly implied by the
        codehash check and is redundant --- but it is a two-gas-class
        \opcode{EXTCODESIZE} and it also short-circuits the common operator error
        of passing an EOA or the zero address, for which \yul{code.length} is
        $0$. Keeping both is cheap belt-and-braces, not confusion.
  \item \textbf{No zero-address check is needed.} \yul{address(0).code.length}
        is $0 \neq 17025$, so the guard already rejects it.
  \item \textbf{An attacker-deployed VK is harmless.} Nothing restricts
        \emph{who} deployed \yul{authorizedVk}. It does not need to: any contract
        whose codehash equals \yul{EXPECTED\_VK\_CODEHASH} is byte-identical to
        the intended verifying key.
\end{itemize}

The one scenario the constructor alone cannot cover is a dependency whose code
changes after construction --- historically reachable by deploying the VK through
\opcode{CREATE2} with a self-destructing constructor and redeploying different
code at the same address. That is why the runtime repeats the check: lines
1386--1389 (outside this range, covered elsewhere) assert
\yul{eq(extcodesize(vk), EXPECTED\_VK\_LENGTH)} and
\yul{eq(extcodehash(vk), EXPECTED\_VK\_CODEHASH\_WORD)} before every
\yul{extcodecopy} of the payload, and fail with the typed
\revert{VkMismatch}. The constructor pin is therefore fail-fast, and the
security property is carried by the runtime re-check. That layering is correct.

\begin{finding}[deploy-6]{Constructor VK pin uses a string \texttt{require} rather than the declared \revert{VkMismatch} error}
\textbf{Severity: Informational.}
\texttt{Halo2Verifier.sol:580--584}. The contract declares a typed failure
taxonomy at lines 48--83 and states its policy explicitly at lines 53--56: probes
stay bare, the memory-layout guard is typed. The constructor's VK pin falls into
neither category --- it emits \texttt{Error(string)} with the payload
\texttt{"invalid vk"}, even though \revert{VkMismatch} (selector
\texttt{0xa447d73e}, line 321) exists and is exactly what the runtime path uses
for the identical condition at line 1389. This is a consistency defect, not a
security one: the deployment fails closed either way, and a failed deployment is
observed by the deployer, not by an integrator. It also costs creation-code bytes
--- a 10-character string plus \texttt{Error(string)} ABI framing, versus a
4-byte selector --- on an artifact whose size budget is documented as tight
(lines 8--11). Using \revert{VkMismatch} here would be smaller, would match the
runtime, and would let deployment tooling decode the same error on both paths.
\end{finding}

\begin{finding}[deploy-7]{No post-deployment way to re-run the capability probes}
\textbf{Severity: Observation.}
\texttt{Halo2Verifier.sol:351} declares \yul{require\_eip2537\_precompiles} as
\yul{private}, and line 579 is its only call site. The gas-bound comment block at
lines 284--289 states the liveness caveat itself: ``if a future fork reprices
these precompiles above the bounds below, this verifier must be regenerated and
redeployed.'' Once deployed, however, there is no entry point that lets an
operator or an automated monitor evaluate that condition --- for example, on a
testnet running a fork candidate --- without submitting a full proof. Exposing
the same body as an \yul{external view} function would make repricing detectable
in advance for the cost of a few hundred runtime bytes. Mitigating context, which
is why this is an Observation rather than a finding of substance: an operator can
already achieve equivalent (in fact broader) coverage by \yul{eth\_call}-ing
\yul{verifyProof} with a known-good canary proof against a fork candidate, which
exercises the real call sites at the real bounds. The runtime-size budget (lines
8--11) is also a legitimate reason not to spend bytes on this.
\end{finding}

\subsection{What the constructor does not do}
\label{sec:deploy:notdone}

For completeness, and because absence of a check is as much a specification fact
as its presence:

\begin{itemize}
  \item It does not validate \yul{BUILD\_ID} (declared at line 311, documented
        by the NatSpec at 301--310) against anything. That
        constant is a published identifier for off-chain deployment records; it
        is not, and is not intended to be, self-checking on chain.
  \item It does not check that byte $0$ of the VK runtime is \opcode{INVALID}.
        It does not need to: the codehash pin fixes that byte along with all
        17{,}024 others.
  \item It does not check \yul{EXPECTED\_VK\_PAYLOAD\_LENGTH} (line 92); that
        constant is consumed only by the runtime \yul{extcodecopy} at line 1393.
  \item It does not emit an event and does not store anything besides
        \yul{AUTHORIZED\_VK}, which is \yul{immutable} and therefore baked into
        the runtime code rather than into storage. The verifier consequently has
        no mutable state at all, and no privileged role: there is no pause, no
        owner, and no way to repoint the VK. The NatSpec at lines 594--597 makes
        replaceability, pausing, and chain/domain binding explicit
        \emph{requirements on an application wrapper}, documented in
        \texttt{docs/reference/DEPLOYMENT\_AND\_INCIDENT\_RESPONSE.md}. That
        division of responsibility is a design choice this section records
        rather than evaluates; the verifier itself being immutable and
        stateless is what makes the deployment-time pins meaningful.
\end{itemize}

\subsection{Section verdict}
\label{sec:deploy:verdict}

Lines 319--623 are clean. The probe suite is unusually thorough for generated
code: it covers every precompile the artifact actually calls and no others, at
the exact gas bounds the runtime forwards and the exact frame shapes the runtime
builds; it tests the two acceptance-deciding precompiles in both the accept and
the reject direction; and every one of its eight known-answer constants is
correct, which I confirmed by recomputation against the BLS12-381 parameters.
The memory guard is correctly strict and is what makes the \yul{memory-safe}
annotation sound. The constructor's VK pins are correct, which I confirmed by
reconstructing and hashing the companion fixture's runtime.

The five substantive items raised above are Informational or Observation: the
one-pair-only shape of the subgroup-rejection coverage (deploy-8), no revert data
on probe failure (deploy-5), the ordering dependency that makes the negative
G1MSM probe non-vacuous (deploy-2), the literal-versus-constant coupling on the
78-pair MSM bound (deploy-4), and the string \texttt{require} in the constructor
where a declared typed error exists (deploy-6). None of them changes the
accept/reject behaviour of a deployed verifier on a conformant chain, and nothing
in this range weakens proof soundness. deploy-8 is the only one that bounds what
a successful deployment lets an operator conclude about a \emph{non}-conformant
chain, and it is a limit of the probe suite's coverage rather than an error in
what it does check.


%% ==== 03-vk-contract.tex ================================================
\section{The \texttt{Halo2VerifyingKey} payload contract}
\label{sec:vk:contract}

This section specifies the whole of
\texttt{fixtures/moonlight-wrap/Halo2VerifyingKey.sol} (696 lines; the file has
no terminating newline, so \texttt{wc -l} reports 695). The contract has no
state variables, no functions, no fallback and no receive function. Its entire
observable behaviour is: at construction, assemble a fixed 17{,}025-byte string
in memory and return it as the deployed runtime code. Everything else in the
file is data.

The companion verifier, \texttt{Halo2Verifier.sol}, treats this runtime as an
immutable, content-addressed blob: it pins the byte length and the
\opcode{EXTCODEHASH} in generated constants
(\texttt{Halo2Verifier.sol}:92--95) and copies the payload into memory with
\opcode{EXTCODECOPY} before every proof
(\texttt{Halo2Verifier.sol}:1393). The verifier-side loader, the header
re-checks and the quotient interpreter are specified elsewhere in this
document; this section cites them only where the payload's meaning depends on
them.

\subsection{Contract role: a constructor-as-data-blob}
\label{sec:vk:pattern}

\begin{lstlisting}[style=solidity,caption={\texttt{Halo2VerifyingKey.sol} lines 50--67: the constructor prologue},label={lst:vk:prologue}]
contract Halo2VerifyingKey {
    /// @notice Deploy the verifying-key payload as this contract's runtime bytecode.
    /// @dev The constructor writes an INVALID byte followed by generated words into memory and returns that prefixed runtime.
    /// @dev The transient construction buffer starts at `0x80`, preserving Solidity's reserved memory words.
    constructor() {
        assembly {
            // Runtime layout:
            //   byte 0      : INVALID, so the payload cannot be executed
            //   byte 1..end : generated VK payload copied by Halo2Verifier
            //
            // `runtime` includes the INVALID prefix; `payload` points to word
            // zero of the verifier-key data described in the contract NatSpec.
            let runtime := 0x80
            let payload := add(runtime, 0x01)
            mstore8(runtime, 0xfe)
            // Header, base-point, and quotient-program words generated from
            // VkPayloadLayout. The inline names on each mstore identify the
            // exact slot in the rendered source.
\end{lstlisting}

The pattern is the standard EVM ``data contract'' (sometimes called SSTORE2 /
\texttt{write} contracts): a constructor that writes a byte string to memory and
returns it, so that the string becomes the deployed runtime code. The important
consequence for a reader is that the Solidity \emph{runtime} object solc would
otherwise emit --- the dispatcher, and the appended CBOR metadata --- never
reaches the chain. The explicit \yul{return} at line 693 pre-empts solc's own
constructor epilogue, so the deployed code is exactly the 17{,}025 bytes
assembled at \texttt{0x80}. That is what makes the claim in the file header
(lines 3--5, ``this contract's runtime is pure returned data, so its codehash is
compiler-independent'') true: no solc version marker, optimiser setting or
metadata hash can perturb the deployed bytes.

\paragraph{Verification.} We reconstructed the runtime independently by parsing
the 532 \yul{mstore} immediates out of the source, prefixing \texttt{0xfe}, and
hashing:
\[
\mathrm{keccak256}\bigl(\texttt{0xfe} \,\|\, \text{payload}\bigr)
= \texttt{0xe68d8936\ldots1c94cf52},
\]
which is byte-for-byte the
\yul{EXPECTED\_VK\_CODEHASH\_WORD} literal at
\texttt{Halo2Verifier.sol}:94. This is a complete, source-only reproduction of
the pin: it confirms both that the deployed runtime carries no compiler
metadata and that the verifier's pinned hash matches this exact fixture.

\paragraph{The two implicit compiler-inserted guards.} The source declares no
guards of its own, but two are inserted by solc and belong in a faithful
specification. First, \texttt{constructor()} is not \texttt{payable}, so solc
emits a \yul{callvalue} check that reverts the deployment if any wei is sent.
(Ether force-sent afterwards --- via a \opcode{SELFDESTRUCT} beneficiary or a
block reward --- would be permanently locked, since byte 0 of the runtime is
\opcode{INVALID} and no code path can move it. This is inherent to every data
contract and is noted only for completeness.) Second, line 6 pins
\texttt{pragma solidity 0.8.30} exactly (no caret), so the artifact is
reproducible from one toolchain. The file's own header comment (lines 3--5) is
careful about what that pin is for: because the runtime is returned data, this
contract's codehash does not depend on the compiler at all --- the pin exists so
that the VK and the verifier are provably built together, not because this
contract needs it.

\paragraph{What the pattern buys.} Reading a word of the VK costs
\opcode{EXTCODECOPY} once per proof (\texttt{100} warm-account gas --- the
\texttt{2600} cold-access charge was already paid by the
\opcode{EXTCODESIZE} at \texttt{Halo2Verifier.sol}:1387 --- plus
$3\lceil \ell/32\rceil$ copy gas plus memory expansion, $\ell = 17024$), i.e.\
$3 \times 532 = \texttt{1{,}596}$ gas of copy cost for the whole 532-word
table, versus
$532 \times 2100 \approx \texttt{1.1M}$ gas for cold \opcode{SLOAD}s or
$\approx \texttt{21kB}$ of \opcode{PUSH32} immediates inlined into the verifier
runtime. It also physically separates the circuit-dependent data from the
circuit-dependent \emph{code}, so the verifier's own runtime stays under
EIP-170 (Section~\ref{sec:vk:size}).

\paragraph{What it costs.} Three things. (i) The payload is a bare byte string:
it carries no magic number, no version tag, no section lengths and no internal
checksum, so every offset into it is a literal duplicated on the verifier side
(Section~\ref{sec:vk:layout}). (ii) Deployment is a two-step ceremony --- the VK
must exist before the verifier's constructor can pin it --- with a
correspondingly larger operational surface. (iii) The blob is subject to
EIP-170's 24{,}576-byte runtime limit, which caps circuit size
(Section~\ref{sec:vk:size}).

\subsection{The \opcode{INVALID} prefix and the runtime layout}
\label{sec:vk:invalid}

The deployed runtime is
\[
\text{runtime} \;=\; \underbrace{\texttt{0xfe}}_{\opcode{INVALID}}
\;\|\; \underbrace{\text{payload}}_{17024 \text{ bytes}},
\qquad |\text{runtime}| = 17025 = \texttt{0x4281}.
\]
Line 62 sets \yul{runtime := 0x80}; line 63 sets
\yul{payload := add(runtime, 0x01)}; line 64 writes the single prefix byte with
\yul{mstore8(runtime, 0xfe)}. Because \yul{payload} is \texttt{runtime + 1}, the
first payload \yul{mstore} covers \texttt{0x81}--\texttt{0xa0} and cannot
clobber the prefix byte at \texttt{0x80}; the write order (prefix first, then
payload) is therefore immaterial here, but it would not be if
\yul{payload} were word-aligned with \yul{runtime}.

\paragraph{Why the prefix exists.} A data contract's runtime is still
\emph{code}. Any \opcode{CALL}, \opcode{STATICCALL} or \opcode{DELEGATECALL}
to this address begins execution at PC 0 and interprets the payload as an
opcode stream. Byte 0 being \texttt{0xfe} (\opcode{INVALID}) makes every such
entry abort unconditionally, consuming all forwarded gas, with no return data
and no state change. Three concrete hazards are removed:

\begin{enumerate}
\item \textbf{Accidental success.} Without the prefix, byte 0 is the first byte
of \yul{vk\_digest}, which is generator-determined. If it happened to be
\texttt{0x00} (\opcode{STOP}) --- a $1/256$ event across regenerations --- a
plain call to the VK address would \emph{succeed} with empty return data. An
integration that guards on the boolean result of a low-level call would then
silently accept the VK contract as a valid callee.
\item \textbf{Executable payload bytes.} The 17{,}024 payload bytes of this
fixture contain 137 occurrences of \texttt{0xff} (\opcode{SELFDESTRUCT}), 29 of
\texttt{0xf4} (\opcode{DELEGATECALL}), 32 of \texttt{0xf1} (\opcode{CALL}), 28
of \texttt{0xf5} (\opcode{CREATE2}), 32 of \texttt{0x55} (\opcode{SSTORE}) and
56 occurrences of \texttt{0x5b} (\opcode{JUMPDEST}), of which 33 survive EVM
jumpdest analysis as genuine jump targets (the other 23 fall inside
\opcode{PUSH} immediate data). Reaching any of them requires
a specific stack shape that the payload itself would have to produce, and here
byte 0 is \texttt{0x56} (\opcode{JUMP}), which underflows an empty stack
immediately --- but a generator cannot be asked to re-argue that for every
circuit it emits. The prefix makes the question moot for all payloads.
\item \textbf{Undeployability under EIP-3541.} Since London, a contract whose
runtime code begins with \texttt{0xEF} is rejected at deployment. Without the
prefix, a \yul{vk\_digest} whose leading byte is \texttt{0xEF} --- again a
$1/256$ event --- would produce a VK that simply cannot be deployed. Pinning
byte 0 to \texttt{0xfe} eliminates the class. (The payload's \emph{interior}
contains 37 \texttt{0xEF} bytes; only byte 0 is constrained by EIP-3541.)
\end{enumerate}

\paragraph{Cost of the prefix.} One byte, and one unit of arithmetic in the
verifier: every offset the verifier uses is payload-relative, so the loader must
copy from source offset \texttt{0x01} and the pinned length constants come in
two flavours, \yul{EXPECTED\_VK\_LENGTH} $=$ 17025 (runtime, used for
\opcode{EXTCODESIZE} and \opcode{EXTCODEHASH}) and
\yul{EXPECTED\_VK\_PAYLOAD\_LENGTH} $=$ 17024 (used for the
\opcode{EXTCODECOPY} length). Both literals appear at
\texttt{Halo2Verifier.sol}:92--93 and are used at 1387 and 1393 respectively.
Conflating them is a real hazard in three distinct ways, and the generated
code avoids all three. Using \yul{EXPECTED\_VK\_PAYLOAD\_LENGTH} in the
\opcode{EXTCODESIZE} comparison at line 1387 would make every proof revert
\revert{VkMismatch} (fails closed, but bricks the verifier). Using
\yul{EXPECTED\_VK\_LENGTH} as the \opcode{EXTCODECOPY} length would copy
17{,}025 bytes to $\mptr{VK\_MPTR}$: 17{,}024 real payload bytes plus one
zero byte spilled at
$\mptr{VK\_MPTR} + \texttt{0x4280} = \texttt{0x7900} = \mptr{THETA\_MPTR}$
(\opcode{EXTCODECOPY} zero-fills past the end of the code), silently writing
outside the VK window. Copying from source offset \texttt{0x00} rather than
\texttt{0x01} would shift the entire payload by one byte and place the
\opcode{INVALID} prefix in the top byte of \yul{vk\_digest}. The generated code
keeps the two constants distinct and comments the distinction at
\texttt{Halo2Verifier.sol}:1383--1385.

\begin{finding}[VK-1]{The \opcode{INVALID} prefix is complete for entry-point
protection}
\textbf{Severity:} Informational (clean region).\\
Every EVM entry into a contract's code starts at PC 0, including
\opcode{DELEGATECALL} from another contract. A single \opcode{INVALID} at byte 0
is therefore a total barrier: there is no \opcode{JUMP} target, calldata shape
or gas amount that can begin execution at a later byte of this runtime. The
construction is sound, and the fixture implements it exactly
(\texttt{Halo2VerifyingKey.sol}:62--64, 693).
\end{finding}

\subsection{Payload layout}
\label{sec:vk:layout}

The payload is a flat array of 532 big-endian 32-byte words written by 532
consecutive \yul{mstore}s at strictly increasing offsets
\texttt{0x0000}, \texttt{0x0020}, \dots, \texttt{0x4260}. We verified
mechanically that the offset sequence is contiguous with no gaps, no
duplicates and no overlaps, and that every immediate is exactly 64 hex digits
wide. Consequently every byte of the returned region
$[\texttt{0x80}, \texttt{0x4301})$ is explicitly written: byte \texttt{0x80} by
the \yul{mstore8}, and \texttt{0x81}--\texttt{0x4300} by the 532 word stores
($\texttt{0x81} + \texttt{0x4280} = \texttt{0x4301}$). No uninitialised or
stale memory is returned.

Table~\ref{tab:vk:sections} gives the section map, in payload-relative byte
offsets and in the absolute memory addresses the verifier sees after
\yul{extcodecopy(vk, VK\_MPTR, 0x01, 17024)} with
$\mptr{VK\_MPTR} = \texttt{0x3680}$.

\begin{table}[h]
\centering
\begin{tabular}{@{}llrllr@{}}
\toprule
Section & Payload bytes & Words & Memory range & Source lines & Items \\
\midrule
Header scalars      & \texttt{0x0000}--\texttt{0x015f} & 11  & \texttt{0x3680}--\texttt{0x37df} & 68--78   & 11 \\
\texttt{G1\_BASE}       & \texttt{0x0160}--\texttt{0x01df} & 4   & \texttt{0x37e0}--\texttt{0x385f} & 79--82   & 1 pt \\
\texttt{G2\_BASE}       & \texttt{0x01e0}--\texttt{0x02df} & 8   & \texttt{0x3860}--\texttt{0x395f} & 83--90   & 1 pt \\
\texttt{NEG\_S\_G2\_BASE} & \texttt{0x02e0}--\texttt{0x03df} & 8 & \texttt{0x3960}--\texttt{0x3a5f} & 91--98   & 1 pt \\
\texttt{quotient\_const}   & \texttt{0x03e0}--\texttt{0x1a1f} & 178 & \texttt{0x3a60}--\texttt{0x509f} & 99--276  & 178 \\
\texttt{quotient\_program} & \texttt{0x1a20}--\texttt{0x2bff} & 143 & \texttt{0x50a0}--\texttt{0x627f} & 277--419 & 4559 B \\
\texttt{fixed\_comms}      & \texttt{0x2c00}--\texttt{0x397f} & 108 & \texttt{0x6280}--\texttt{0x6fff} & 420--581 & 27 pts \\
\texttt{permutation\_comms}& \texttt{0x3980}--\texttt{0x427f} & 72  & \texttt{0x7000}--\texttt{0x78ff} & 582--689 & 18 pts \\
\midrule
\textbf{Total}      & \texttt{0x0000}--\texttt{0x427f} & \textbf{532} & \texttt{0x3680}--\texttt{0x78ff} & 68--689 & \\
\bottomrule
\end{tabular}
\caption{Payload section map. ``Memory range'' is the address seen by the
verifier after the \opcode{EXTCODECOPY} at \texttt{Halo2Verifier.sol}:1393.}
\label{tab:vk:sections}
\end{table}

Three of these absolute addresses are hard-coded on the verifier side and match
exactly:
\begin{align*}
\yul{q\_const\_mptr}   &= \texttt{0x3a60} = \mptr{VK\_MPTR} + \texttt{0x03e0}
   &&\text{(\texttt{Halo2Verifier.sol}:2013)},\\
\yul{q\_program\_mptr} &= \texttt{0x50a0} = \mptr{VK\_MPTR} + \texttt{0x1a20}
   &&\text{(\texttt{Halo2Verifier.sol}:2016)},\\
\texttt{0x6280} &= \mptr{VK\_MPTR} + \texttt{0x2c00},\quad
\texttt{0x7000} = \mptr{VK\_MPTR} + \texttt{0x3980}
   &&\text{(\texttt{Halo2Verifier.sol}:3860, 3882)}.
\end{align*}
The payload also abuts the next verifier region exactly:
$\mptr{VK\_MPTR} + \texttt{0x4280} = \texttt{0x7900} = \mptr{CHALLENGE\_MPTR}$
(\texttt{Halo2Verifier.sol}:144). There is neither a gap nor an overlap.

The contract's own NatSpec (lines 22--39) states the layout parametrically. It
checks out for this fixture: with $Q_{\text{payload}} = 178 + 143 = 321$ and
$N_{\text{fixed}} = 27$,
\[
31 + Q_{\text{payload}} = 352 \;\Rightarrow\; \texttt{0x2c00},
\qquad
31 + Q_{\text{payload}} + 4 N_{\text{fixed}} = 460 \;\Rightarrow\; \texttt{0x3980},
\]
which are exactly the rendered \texttt{fixed\_comms[0].x\_hi} and
\texttt{permutation\_comms[0].x\_hi} offsets.

\begin{finding}[VK-2]{The payload is not self-describing; the codehash pin
protects the VK half only}
\textbf{Severity:} Low.\\
Nothing inside the blob identifies where the header ends, how many constants
follow, or how long the packed program is. The verifier carries every boundary
as an independent literal --- \texttt{0x3a60}, \texttt{0x50a0},
\texttt{0x11cf}, \texttt{0x6280}, \texttt{0x7000}, 17024, 17025
(\texttt{Halo2Verifier.sol}:92--93, 2013, 2016, 2170, 3860, 3882) --- and the
VK carries the matching \yul{mstore} offsets.\\
\emph{The pin is one-sided.} \yul{EXPECTED\_VK\_CODEHASH} covers all
17{,}025 runtime bytes of the \emph{VK}, so editing the VK is caught at
verifier construction (\texttt{Halo2Verifier.sol}:580--584) and on every proof
(1386--1389). It does \emph{not} cover the verifier's own literals: editing
\yul{q\_const\_mptr}, \yul{q\_program\_mptr}, \yul{q\_end}'s \texttt{0x11cf},
or either commitment base address leaves the deployed VK, and hence its
codehash, unchanged, so such a verifier deploys successfully against the
unmodified VK and then mis-parses the blob. Of the seven boundary literals,
only 17024 and 17025 are checked against anything on chain (against
\opcode{EXTCODESIZE}); the five section offsets
\texttt{0x3a60}, \texttt{0x50a0}, \texttt{0x11cf}, \texttt{0x6280} and
\texttt{0x7000} are compared with nothing at all.\\
\emph{What partially back-stops a verifier/VK shape mismatch:} the six header
re-checks at
\texttt{Halo2Verifier.sol}:1400--1406 (\texttt{num\_instances}, \texttt{k},
\texttt{has\_accumulator}, \texttt{acc\_offset}, \texttt{num\_acc\_limbs},
\texttt{num\_acc\_limb\_bits}), and the absorption of \yul{vk\_digest} into
the transcript (line 1481), which makes a shape mismatch overwhelmingly likely
to produce a failing --- not a forgeable --- proof. Neither mechanism covers
the section offsets, and neither covers the circuit-shape parameters the
contract's own NatSpec (lines 48--49) says are baked into the verifier rather
than into the blob: the per-lookup chunk counts, the trashcan structure and
\texttt{num\_simple\_selectors}. Those live only in the verifier's generated
code, which no on-chain digest pins.\\
\emph{So the residual hazard is operational rather than cryptographic, but it
is real:} a maintainer who regenerates one file and hand-patches the other ---
in either direction --- gets a pair that deploys. Correctness rests on the two
halves being emitted by one generator run from one \texttt{VkPayloadLayout}
(\texttt{src/lowering/render/models.rs}:278--311 rejects an offset mismatch at
generation time, but only at generation time). Recommendation: never edit
either file by hand; treat the pair as one generated artifact, and consider
having the verifier's constructor assert its own section literals against a
generated self-describing header word.
\end{finding}

\subsection{Header words 0--10}
\label{sec:vk:header}

\begin{lstlisting}[style=solidity,caption={\texttt{Halo2VerifyingKey.sol} lines 68--78: the eleven header scalars},label={lst:vk:header}]
            mstore(add(payload, 0x0000), 0x56c0824fcff237dd8dc7b15f527346d9e1647d191815acb142500b0293e84f66) // vk_digest
            mstore(add(payload, 0x0020), 0x0000000000000000000000000000000000000000000000000000000000000013) // num_instances
            mstore(add(payload, 0x0040), 0x0000000000000000000000000000000000000000000000000000000000000014) // k
            mstore(add(payload, 0x0060), 0x73eda0144f284aae5b6554d46c21576b363d4ec725be2bff1a400fff00001001) // n_inv
            mstore(add(payload, 0x0080), 0x03e1c54bcb947035a57a6e07cb98de4a2f69e02d265e09d9fece7e0e39898d4b) // omega
            mstore(add(payload, 0x00a0), 0x6c39442eade0092768ac033fa6f608750624a1bb17dbc026ef97c3573a28fc8c) // omega_inv
            mstore(add(payload, 0x00c0), 0x2a0ccbaa0613f093f2bb6e97859513f0b613d8587eaa92db9e5604b8d6b68d45) // omega_inv_to_l
            mstore(add(payload, 0x00e0), 0x0000000000000000000000000000000000000000000000000000000000000001) // has_accumulator
            mstore(add(payload, 0x0100), 0x000000000000000000000000000000000000000000000000000000000000000b) // acc_offset
            mstore(add(payload, 0x0120), 0x0000000000000000000000000000000000000000000000000000000000000007) // num_acc_limbs
            mstore(add(payload, 0x0140), 0x0000000000000000000000000000000000000000000000000000000000000038) // num_acc_limb_bits
\end{lstlisting}

\begin{table}[h]
\centering
\begin{tabular}{@{}rlllp{5.1cm}@{}}
\toprule
W & Offset & Name & Value & Meaning \\
\midrule
0 & \texttt{0x0000} & \texttt{vk\_digest} & \texttt{0x56c0824f\ldots93e84f66} & $\Fq$ transcript representation of the constraint system \\
1 & \texttt{0x0020} & \texttt{num\_instances} & \texttt{0x13} $= 19$ & number of public-input words \\
2 & \texttt{0x0040} & \texttt{k} & \texttt{0x14} $= 20$ & $\log_2 n$, so $n = 2^{20} = 1{,}048{,}576$ \\
3 & \texttt{0x0060} & \texttt{n\_inv} & \texttt{0x73eda014\ldots00001001} & $n^{-1} \bmod r$ \\
4 & \texttt{0x0080} & \texttt{omega} & \texttt{0x03e1c54b\ldots39898d4b} & $\omega$, primitive $2^{20}$-th root of unity \\
5 & \texttt{0x00a0} & \texttt{omega\_inv} & \texttt{0x6c39442e\ldots3a28fc8c} & $\omega^{-1}$ \\
6 & \texttt{0x00c0} & \texttt{omega\_inv\_to\_l} & \texttt{0x2a0ccbaa\ldots d6b68d45} & $\omega^{-10}$ \\
7 & \texttt{0x00e0} & \texttt{has\_accumulator} & $1$ & the instance vector carries a KZG accumulator \\
8 & \texttt{0x0100} & \texttt{acc\_offset} & \texttt{0x0b} $= 11$ & first instance index of the accumulator \\
9 & \texttt{0x0120} & \texttt{num\_acc\_limbs} & $7$ & limbs per $\Fp$ coordinate \\
10 & \texttt{0x0140} & \texttt{num\_acc\_limb\_bits} & \texttt{0x38} $= 56$ & bits per limb \\
\bottomrule
\end{tabular}
\caption{Header scalars, with the concrete values of this fixture
(lines 68--78).}
\label{tab:vk:header}
\end{table}

\paragraph{Domain constants: independent verification.} All four domain scalars
are internally consistent, computed over $r = \texttt{0x73eda753\ldots00000001}$
with $n = 2^{20}$:
\begin{align}
n \cdot \texttt{n\_inv} &\equiv 1 \pmod r, \label{eq:vk:ninv}\\
\omega^{n} &\equiv 1 \pmod r, \qquad \omega^{n/2} \not\equiv 1 \pmod r,
   \label{eq:vk:omega}\\
\omega \cdot \omega^{-1} &\equiv 1 \pmod r, \label{eq:vk:ominv}\\
\texttt{omega\_inv\_to\_l} &= (\omega^{-1})^{10} \bmod r. \label{eq:vk:oitl}
\end{align}
Equation~\eqref{eq:vk:omega} establishes that $\omega$ has exact order $n$, not
merely order dividing $n$ (which is the property the Lagrange evaluations
actually need: a non-primitive $\omega$ would collapse $L_0$ and
$L_{\mathrm{last}}$). Equation~\eqref{eq:vk:oitl} pins
$|\mathrm{rotation\_last}| = 10$, i.e.\ $9$ blinding rows plus one. The
verifier's Lagrange pass confirms the reading: the denominator loop at
\texttt{Halo2Verifier.sol}:1818--1823 seeds
\yul{pow\_of\_omega} with \yul{mload(OMEGA\_INV\_TO\_L\_MPTR)} $= \omega^{-10}$
and multiplies by $\omega$ on each of its $\texttt{0x03a0}/\texttt{0x20} = 29$
iterations, so the powers walk \emph{up} from $\omega^{-10}$ toward
$\omega^{18}$, not down. Slot $0$ of that run is $L_{\mathrm{last}}$, and
$L_{\mathrm{blind}}$ is the sum of slots $1$--$9$
(\texttt{Halo2Verifier.sol}:1841--1847, whose bound
$\texttt{0x0140}/\texttt{0x20} - 1 = 9$ is exactly the blinding-row count
implied by $\omega^{-10}$).

\paragraph{Accumulator parameters: cross-check against the verifier.} With
$b = 56$ bits per limb and $n_\ell = 7$ limbs, the verifier packs
$\lfloor 256/b \rfloor = 4$ limbs per instance word, so an $\Fp$ coordinate
occupies
\[
\texttt{coord\_words} = \left\lceil \frac{n_\ell}{4} \right\rceil = 2
\quad\text{words} \qquad (\text{\texttt{Halo2Verifier.sol}:1254}),
\]
and one G1 point occupies $2 \times 2 = 4$ instance words. The accumulator is an
already-collapsed point pair (LHS point, RHS point, both with implicit unit
scalars; \texttt{Halo2Verifier.sol}:1264--1300), hence exactly 8 instance
words. Therefore
\[
\texttt{acc\_offset} + 8 = 11 + 8 = 19 = \texttt{num\_instances},
\]
i.e.\ the accumulator occupies the final 8 of the 19 public inputs and ends
exactly at the end of the instance vector. Independently, the verifier computes
the accumulator's calldata cursor as
$\mptr{INSTANCE\_CPTR} + \texttt{0x0160}$
(\texttt{Halo2Verifier.sol}:1257), and $\texttt{0x0160} = 352 = 11 \times 32$,
confirming \texttt{acc\_offset} $= 11$. Note also
$n_\ell \cdot b = 7 \times 56 = 392 \ge 381 = \lceil \log_2 p \rceil$, so the
limb decomposition is wide enough to represent any BLS12-381 base-field
coordinate.

\paragraph{Six of the eleven header words are re-checked against literals.} The
verifier re-loads and compares \texttt{num\_instances}, \texttt{k},
\texttt{has\_accumulator}, \texttt{acc\_offset}, \texttt{num\_acc\_limbs} and
\texttt{num\_acc\_limb\_bits} against the generated integers 19, 20, 1, 11, 7
and 56 (\texttt{Halo2Verifier.sol}:1400--1406), reverting
\revert{VkMismatch} on any difference. The other five scalars
(\texttt{vk\_digest}, \texttt{n\_inv}, $\omega$, $\omega^{-1}$,
$\omega^{-10}$) and all 20 base-point words are \emph{not} re-checked against
literals. That is sound rather than an omission: the codehash pin already
covers all 17{,}024 payload bytes, so the six literal comparisons are strictly
redundant belt-and-braces --- valuable only as a louder failure mode if someone
deploys a VK for a different circuit shape.

\subsection{Base points: \texttt{G1\_BASE}, \texttt{G2\_BASE},
\texttt{NEG\_S\_G2\_BASE} (words 11--30)}
\label{sec:vk:bases}

\begin{lstlisting}[style=solidity,caption={\texttt{Halo2VerifyingKey.sol} lines 79--98: the three EIP-2537 base points},label={lst:vk:bases}]
            mstore(add(payload, 0x0160), 0x0000000000000000000000000000000017f1d3a73197d7942695638c4fa9ac0f) // g1_x_hi
            mstore(add(payload, 0x0180), 0xc3688c4f9774b905a14e3a3f171bac586c55e83ff97a1aeffb3af00adb22c6bb) // g1_x_lo
            mstore(add(payload, 0x01a0), 0x0000000000000000000000000000000008b3f481e3aaa0f1a09e30ed741d8ae4) // g1_y_hi
            mstore(add(payload, 0x01c0), 0xfcf5e095d5d00af600db18cb2c04b3edd03cc744a2888ae40caa232946c5e7e1) // g1_y_lo
            mstore(add(payload, 0x01e0), 0x00000000000000000000000000000000024aa2b2f08f0a91260805272dc51051) // g2_x_c0_hi
            mstore(add(payload, 0x0200), 0xc6e47ad4fa403b02b4510b647ae3d1770bac0326a805bbefd48056c8c121bdb8) // g2_x_c0_lo
            mstore(add(payload, 0x0220), 0x0000000000000000000000000000000013e02b6052719f607dacd3a088274f65) // g2_x_c1_hi
            mstore(add(payload, 0x0240), 0x596bd0d09920b61ab5da61bbdc7f5049334cf11213945d57e5ac7d055d042b7e) // g2_x_c1_lo
            mstore(add(payload, 0x0260), 0x000000000000000000000000000000000ce5d527727d6e118cc9cdc6da2e351a) // g2_y_c0_hi
            mstore(add(payload, 0x0280), 0xadfd9baa8cbdd3a76d429a695160d12c923ac9cc3baca289e193548608b82801) // g2_y_c0_lo
            mstore(add(payload, 0x02a0), 0x000000000000000000000000000000000606c4a02ea734cc32acd2b02bc28b99) // g2_y_c1_hi
            mstore(add(payload, 0x02c0), 0xcb3e287e85a763af267492ab572e99ab3f370d275cec1da1aaa9075ff05f79be) // g2_y_c1_lo
            mstore(add(payload, 0x02e0), 0x000000000000000000000000000000000632aaf712568f19c297802268a7ad9d) // neg_s_g2_x_c0_hi
            mstore(add(payload, 0x0300), 0xceea6ef7ab6f75a7d26781c8e90c7432bc5e99dcc219ba64010f3052123983ab) // neg_s_g2_x_c0_lo
            mstore(add(payload, 0x0320), 0x00000000000000000000000000000000191ff4920e077a2f8cb3969ba8f05bc2) // neg_s_g2_x_c1_hi
            mstore(add(payload, 0x0340), 0xaa9da8c95d640b1e051be7cf344ee7f01996df2568bf0e7ccd9eb70978820045) // neg_s_g2_x_c1_lo
            mstore(add(payload, 0x0360), 0x0000000000000000000000000000000005f434ebf45460a864ad5b17497c7903) // neg_s_g2_y_c0_hi
            mstore(add(payload, 0x0380), 0x71820c70c83aa186029536d22dff54373251152c28bc43269f95281eba1b012e) // neg_s_g2_y_c0_lo
            mstore(add(payload, 0x03a0), 0x0000000000000000000000000000000004d1c747141bcac15e77e3e1d3853254) // neg_s_g2_y_c1_hi
            mstore(add(payload, 0x03c0), 0xc8687afdad35345a04f79d9c2759007f6640676eb44aee7011ce5ad80744bb23) // neg_s_g2_y_c1_lo
\end{lstlisting}

\paragraph{Encoding.} Each BLS12-381 base-field element is 381 bits and is
stored in the EIP-2537 precompile ABI form: 64 bytes consisting of 16 zero
bytes followed by the 48-byte big-endian value. In 32-byte words that is a
\emph{hi} word whose top 128 bits are zero and whose low 128 bits are the top 16
bytes of the coordinate, and a \emph{lo} word holding the low 32 bytes:
\[
c = \bigl(\texttt{hi} \bmod 2^{128}\bigr)\cdot 2^{256} + \texttt{lo},
\qquad \texttt{hi} < 2^{128}.
\]
A $\GO$ point is thus 4 words (\texttt{x\_hi}, \texttt{x\_lo}, \texttt{y\_hi},
\texttt{y\_lo}) and a $\GT$ point is 8 words in the order
$x_{c_0}, x_{c_1}, y_{c_0}, y_{c_1}$, each split hi/lo. This is exactly the
byte layout the \texttt{0x0c} (G1MSM) and \texttt{0x0f} (pairing) precompiles
expect, which is why the verifier can \yul{mcopy} these slots straight into
precompile input buffers without re-serialising
(\texttt{Halo2Verifier.sol}:971, 973, 4012).

\paragraph{Independent verification.} We decoded all three points and checked
them against the BLS12-381 parameters
$p = \texttt{0x1a0111ea}\ldots\texttt{ffffaaab}$ and
$r = \texttt{0x73eda753}\ldots\texttt{00000001}$:

\begin{table}[h]
\centering
\begin{tabularx}{\textwidth}{@{}lccccX@{}}
\toprule
Point & hi words $< 2^{128}$ & coords $< p$ & on curve & $[r]P = \id$ & Identification \\
\midrule
\texttt{G1\_BASE}       & yes & yes & yes & yes & the standard BLS12-381 $\GO$ generator \\
\texttt{G2\_BASE}       & yes & yes & yes & yes & the standard BLS12-381 $\GT$ generator \\
\texttt{NEG\_S\_G2\_BASE} & yes & yes & yes & yes & structured-reference-string element, not independently identifiable \\
\bottomrule
\end{tabularx}
\caption{Validation of the three embedded base points (lines 79--98). ``on
curve'' means $y^2 = x^3 + 4$ over $\Fp$ (resp.\ $y^2 = x^3 + 4(1+u)$ over
$\Fp[u]/(u^2+1)$ for $\GT$).}
\label{tab:vk:bases}
\end{table}

The G1 and G2 base points are bit-for-bit the canonical BLS12-381 generators.
Write $g_2 \in \GT$ for the generator held in \texttt{G2\_BASE}.
\texttt{NEG\_S\_G2\_BASE} is intended to be $[-s]g_2$ for the KZG trapdoor $s$,
used so that the final check can be written as a single product of pairings,
$\pair{A}{g_2} \cdot \pair{B}{[-s]g_2} = 1$
(\texttt{Halo2Verifier.sol}:4045--4046, 4100). We confirmed it is a valid,
prime-order-subgroup $\GT$ point with canonical coordinates.

\begin{finding}[VK-3]{\texttt{NEG\_S\_G2\_BASE} is unattested trusted-setup data}
\textbf{Severity:} Informational.\\
\emph{Location:} \texttt{Halo2VerifyingKey.sol}:91--98.\\
Nothing on-chain --- and nothing in either generated contract --- ties
\texttt{NEG\_S\_G2\_BASE} to a particular powers-of-tau ceremony. It is a
well-formed $\GT$ point in the prime-order subgroup (verified above), but
soundness of the whole KZG argument requires that no party knows the $s$ with
$[-s]g_2$ equal to this value. If the generator were fed an SRS with a known
trapdoor, the deployed verifier would accept forged openings while every
on-chain check --- codehash pin, curve check, subgroup check, pairing --- still
passed.\\
This is inherent to KZG and cannot be fixed inside the verifier; it must be
discharged by attesting the SRS out of band and reviewing that
\texttt{NEG\_S\_G2\_BASE} equals the ceremony's $[-s]g_2$ before deployment. A
reviewer should treat these 8 words as the single most consequential
unverifiable constant in the payload.
\end{finding}

\begin{finding}[VK-4]{The embedded points are never re-validated on chain, and
do not need to be}
\textbf{Severity:} Informational (clean region).\\
The verifier performs no curve or subgroup check on \texttt{G1\_BASE},
\texttt{G2\_BASE} or \texttt{NEG\_S\_G2\_BASE}; it copies them straight into
precompile buffers. This is sound: they are constants covered by
\yul{EXPECTED\_VK\_CODEHASH}, and EIP-2537 precompiles perform their own
subgroup validation on every input point, so a malformed base point would make
the pairing precompile fail rather than admit an invalid proof. Proof-supplied
and instance-supplied points, by contrast, \emph{are} forced through G1MSM
precisely so the precompile validates them
(\texttt{Halo2Verifier.sol}:1277--1288).
\end{finding}

\subsection{The quotient constant table (words 31--208)}
\label{sec:vk:qconst}

178 words at payload \texttt{0x03e0}--\texttt{0x1a1f}, source lines 99--276,
memory \texttt{0x3a60}--\texttt{0x509f}. Every entry is one canonical $\Fr$
element. The quotient interpreter addresses them with a single unsigned byte
operand \yul{qconst} via \yul{mload(add(q\_const\_mptr, shl(5, qconst)))}
(e.g.\ \texttt{Halo2Verifier.sol}:2243--2245).

\begin{lstlisting}[style=solidity,caption={\texttt{Halo2VerifyingKey.sol} lines 99--113: head of the constant table: small integers, powers of two, negated shifts and opaque elements},label={lst:vk:qconst}]
            mstore(add(payload, 0x03e0), 0x0000000000000000000000000000000000000000000000000000000000000001) // quotient_const
            mstore(add(payload, 0x0400), 0x00bbe1fbe9ef1e2d62490b03a82bf9ef10f5e9b2323033669cf6c50481f63e05) // quotient_const
            mstore(add(payload, 0x0420), 0x0000000000000000000000000000000000000000000000000100000000000000) // quotient_const
            mstore(add(payload, 0x0440), 0x0000000000000000000000000000000000010000000000000000000000000000) // quotient_const
            mstore(add(payload, 0x0460), 0x0000000000000000000000000000000000000000000000000000000400000000) // quotient_const
            mstore(add(payload, 0x0480), 0x0000000000000000000000000000000000000000040000000000000000000000) // quotient_const
            mstore(add(payload, 0x04a0), 0x0000000000000000000000000000000000000000000000000000000000001000) // quotient_const
            mstore(add(payload, 0x04c0), 0x0000000000000000000000000000000000000000000000100000000000000000) // quotient_const
            mstore(add(payload, 0x04e0), 0x73eda753299d7d483339d80809a1d80553bda402fffe5bfeffffffff00000000) // quotient_const
            mstore(add(payload, 0x0500), 0x73eda753299d7d483339d80809a1d80553bda402fffe5bfefeffffff00000001) // quotient_const
            mstore(add(payload, 0x0520), 0x73eda753299d7d483339d80809a1d80553bca402fffe5bfeffffffff00000001) // quotient_const
            mstore(add(payload, 0x0540), 0x73eda753299d7d483339d80809a1d80553bda402fffe5bfefffffffb00000001) // quotient_const
            mstore(add(payload, 0x0560), 0x73eda753299d7d483339d80809a1d80553bda402fbfe5bfeffffffff00000001) // quotient_const
            mstore(add(payload, 0x0580), 0x73eda753299d7d483339d80809a1d80553bda402fffe5bfefffffffefffff001) // quotient_const
            mstore(add(payload, 0x05a0), 0x73eda753299d7d483339d80809a1d80553bda402fffe5beeffffffff00000001) // quotient_const
\end{lstlisting}

\noindent The remaining 163 entries (lines 114--276, offsets
\texttt{0x05c0}--\texttt{0x1a00}) are elided; they have exactly the same
one-\yul{mstore}-per-word form and are classified statistically below rather
than reproduced.

\paragraph{Canonicality.} We checked all 178 values: every one satisfies
$0 < v < r$, and there are no duplicates and no zeros. Canonicality is not
enforced anywhere at runtime (nor does it need to be --- the values only ever
reach \yul{addmod}/\yul{mulmod}, which reduce their arguments), but its holding
here confirms the generator reduces before emitting.

\paragraph{Classification.} Writing $-x$ for $(r - x) \bmod r$, the 178
constants fall into five structural classes
(Table~\ref{tab:vk:qconstclass}).

\begin{table}[h]
\centering
\begin{tabular}{@{}lrp{7.4cm}@{}}
\toprule
Class & Count & Description \\
\midrule
$c \cdot 2^e$, $c \in \{1,3\}$   & 32 & positive limb weights; see below \\
$-(1 \cdot 2^e)$ & 19 & negated limb weights (all have $c = 1$) \\
$v < 2^{136}$, no shift structure & 34 & bit-lengths 64 and 129--136 \\
$-v$ with $v < 2^{136}$           & 29 & bit-lengths 76, 129--130 and 132--135 \\
full-width $\Fr$ elements          & 64 & bit-lengths 137, 225, and 248--255 \\
\bottomrule
\end{tabular}
\caption{Structural classification of the 178 \texttt{quotient\_const} words
(lines 99--276).}
\label{tab:vk:qconstclass}
\end{table}

The shift family is highly regular. The positive entries are
$1 \cdot 2^e$ for
\[
e \in \{0,1,12,13,34,35,56,57,68,69,90,91,112,113,168,169,224,225\}
\quad (18 \text{ entries}),
\]
and $3 \cdot 2^e$ for the same set minus $\{12,34,56,168\}$ (14 entries). The
negated entries are $-(1 \cdot 2^e)$ for the same 18 exponents plus $e = 134$
(19 entries). The exponent multiset decomposes as
$\{0,56,112,168,224\} \cup \{12,34,68,90\}$, with each exponent also appearing
shifted by one. The first set is exactly the radix-$2^{56}$ limb weights
implied by \texttt{num\_acc\_limb\_bits} $= 56$ --- $2^{0}, 2^{56}, 2^{112},
2^{168}, 2^{224}$ are the weights of the four limbs packed into one instance
word and of the carry into the next; the second set $\{12,34\}$ and its
$+56$ translates $\{68,90\}$ come from a finer sub-limb decomposition inside
the circuit. The $c = 3$ entries and the $e{+}1$ variants supply the
$3x$ and $2x$ multiples a range-check or non-native multiplication gadget
needs alongside each weight.

\paragraph{The $r$-adjacent constants.} 52 of the 178 words begin with the
leading seven hex digits of $r$ (\texttt{0x73eda75}), which is what makes the
table look, at a glance, as though it were full of near-modulus values. It is
not: 48 of those 52 are exactly the negations already catalogued (the 19 negated
shifts and the 29 negations of sub-$2^{136}$ values), and the remaining 4 are
ordinary full-width elements whose distance from $r$ is about $2^{179}$ ---
close enough to share a prefix, far from being ``adjacent'' to $r$. The six
genuinely closest to $r$ are:

\begin{table}[h]
\centering
\begin{tabular}{@{}llll@{}}
\toprule
Offset & Line & Final 8 bytes & $r - v$ \\
\midrule
\texttt{0x04e0} & 107 & \texttt{ffffffff00000000} & $1$ \\
\texttt{0x10c0} & 202 & \texttt{fffffffeffffffff} & $2$ \\
\texttt{0x0580} & 112 & \texttt{fffffffefffff001} & $2^{12}$ \\
\texttt{0x1400} & 228 & \texttt{fffffffeffffe001} & $2^{13}$ \\
\texttt{0x0540} & 110 & \texttt{fffffffb00000001} & $2^{34}$ \\
\texttt{0x13c0} & 226 & \texttt{fffffff700000001} & $2^{35}$ \\
\midrule
\multicolumn{2}{@{}l}{$r$ itself} & \texttt{ffffffff00000001} & $0$ \\
\bottomrule
\end{tabular}
\caption{The six \texttt{quotient\_const} entries nearest $r$. All six --- and
$r$ --- share the identical leading 48 hex digits
\texttt{0x73eda753299d7d483339d80809a1d80553bda402fffe5bfe} and differ only in
the final 8 bytes shown. The entry at \texttt{0x04e0} is $r - 1 \equiv -1$.}
\label{tab:vk:radjacent}
\end{table}

A reviewer should read \texttt{0x73eda753\ldots00000000} as ``$-1$'', not as ``a
suspiciously large field element''. The trailing bytes of $r$ itself
(\texttt{\ldots ffffffff00000001}) are what make small negations look like
near-$r$ values with sparse low words.

\paragraph{One constant is identifiable.} The word at payload \texttt{0x08a0}
(line 137) is \texttt{0x0000}\ldots\texttt{0000d201000000010000}, i.e.\ the
integer \texttt{0xd201000000010000} zero-extended to 32 bytes.
This is $|x|$ for the BLS12-381 curve parameter
$x = -\texttt{0xd201000000010000}$, and, by the BLS12 identity $p \equiv x
\pmod r$, it is simultaneously
\[
\texttt{0xd201000000010000} \;=\; (-p) \bmod r,
\]
which we verified numerically. That is precisely the reduction constant a
non-native $\Fp$-arithmetic gadget needs in order to assert
``$\sum_i \ell_i 2^{56 i} \equiv 0 \pmod p$'' inside $\Fr$: one adds
$k \cdot ((-p) \bmod r)$. Its presence corroborates the reading of the shift
family as base-field limb weights for the self-emulated accumulator arithmetic.
It is, however, the only entry we could identify. The remaining 126 opaque
values --- all 64 full-width elements and 62 of the 63 sub-$2^{136}$ values
(positive and negated) --- are circuit-specific coefficients, plausibly gate
and round constants, and are not independently checkable from either contract.

\begin{finding}[VK-5]{The constant index is a byte, but the table has only 178
entries}
\textbf{Severity:} Observation.\\
\emph{Location:} \texttt{Halo2VerifyingKey.sol}:99--276 (the table);
\texttt{Halo2Verifier.sol}:2243--2245, 2315--2317, 2347, 2369, 2396, 2414,
2432 (the reads).\\
The interpreter decodes \yul{qconst} as one unsigned byte and computes
\yul{mload(add(q\_const\_mptr, shl(5, qconst)))} with no bound check. Indices
$0..177$ hit the table; indices $178..255$ would read words from the packed
\texttt{quotient\_program} region that immediately follows at payload
\texttt{0x1a20}, silently substituting bytecode for a field constant rather
than reverting. All seven read sites listed above are unguarded.\\
\emph{The omission is an asymmetry, not a uniform policy.} The same interpreter
\emph{does} bound-check its other operand class: every \texttt{u16} memory
pointer decoded from the program is validated with a single unsigned
range test before use ---
\yul{if gt(sub(q\_ptr, 0x3680), 0x6aa0) \{ q\_program\_fail() \}} at
\texttt{Halo2Verifier.sol}:2216, and the pairwise-block variant
\yul{if gt(sub(q\_lhs\_base, 0x3680), 0x69e0) \{ q\_program\_fail() \}} at
2384--2385. Constant indices get no equivalent, even though the identical
one-instruction idiom would serve.\\
\emph{Why this is not exploitable as deployed:} the index bytes live inside the
same pinned blob as the table, so no calldata can influence them, and the
whole 17{,}024 bytes are covered by \yul{EXPECTED\_VK\_CODEHASH}. Establishing
that no emitted index exceeds 177 requires decoding the packed program and is
the responsibility of the quotient-interpreter section
(\texttt{Halo2Verifier.sol}:2200--3098); it is a generator invariant, not a
runtime guard. Adding a one-time \yul{lt(qconst, 178)} assertion would be
cheap, but its absence is a defence-in-depth gap, not a vulnerability.
\end{finding}

\subsection{The packed quotient program (words 209--351)}
\label{sec:vk:qprogram}

143 words at payload \texttt{0x1a20}--\texttt{0x2bff}, source lines 277--419,
memory \texttt{0x50a0}--\texttt{0x627f}. Unlike every other section, this one is
a \emph{byte} stream that happens to be transported in 32-byte words: the
generator's \texttt{PackedProgramCodec} concatenates variable-width
instructions (1-byte opcodes; 1-, 2-, 3-, 4-, 5- and 13-byte operand groups)
and pads the tail to a word boundary.

\begin{lstlisting}[style=solidity,caption={\texttt{Halo2VerifyingKey.sol} lines 277--282: the first six packed program words},label={lst:vk:qprogram-head}]
            mstore(add(payload, 0x1a20), 0x0594e0109a0005958008060d000b0200011b000021020103070002000094a002) // quotient_program
            mstore(add(payload, 0x1a40), 0x94c00394e00495000595200095400295600395800695a00795c00895e0099600) // quotient_program
            mstore(add(payload, 0x1a60), 0x0a96200b96400c96600d96800e96a00496c00596e006970007972094a0009540) // quotient_program
            mstore(add(payload, 0x1a80), 0x0295600395800496c00596e006970007972094c00295400395600495800596c0) // quotient_program
            mstore(add(payload, 0x1aa0), 0x0696e00797000f972094e00395400495600595800696c00796e00f9700109720) // quotient_program
            mstore(add(payload, 0x1ac0), 0x95000495400595600695800796c00f96e0109700119720952005954006956007) // quotient_program
\end{lstlisting}

\noindent Lines 283--414 (offsets \texttt{0x1ae0}--\texttt{0x2b40}) are
elided; they are 132 further \yul{mstore}s of the same shape.

\begin{lstlisting}[style=solidity,caption={\texttt{Halo2VerifyingKey.sol} lines 415--419: the last five packed program words, showing the 17 zero pad bytes},label={lst:vk:qprogram-tail}]
            mstore(add(payload, 0x2b60), 0x9f96e097205597009700a09700972057972097200d000b0800010594a01194a0) // quotient_program
            mstore(add(payload, 0x2b80), 0x1194a005950008060d000b0900000594c01194c01194c005952008060d000b09) // quotient_program
            mstore(add(payload, 0x2ba0), 0x00010594e01194e01194e00595a008060d000b0900011b00021b000305958008) // quotient_program
            mstore(add(payload, 0x2bc0), 0x109aa00594a01194a01195000daf060594c01194c01195200db0060594e01194) // quotient_program
            mstore(add(payload, 0x2be0), 0xe01195a00db1060d000b090001191f0000000000000000000000000000000000) // quotient_program
\end{lstlisting}

\paragraph{Byte length.} The interpreter sets
\[
\yul{q\_pc} := \yul{q\_program\_mptr} = \texttt{0x50a0},
\qquad
\yul{q\_end} := \yul{q\_program\_mptr} + \texttt{0x11cf}
\]
(\texttt{Halo2Verifier.sol}:2168, 2170) and runs
\yul{for \{ \} lt(q\_pc, q\_end) \{ \}} (line 2200), terminating with
\yul{if iszero(eq(q\_pc, q\_end)) \{ q\_program\_fail() \}} (line 3098), which
reverts \revert{QuotientProgramInvalid}. Hence:
\begin{align*}
\text{program bytes} &= \texttt{0x11cf} = 4559,\\
\text{transported bytes} &= 143 \times 32 = 4576 = \texttt{0x11e0},\\
\text{pad bytes} &= 4576 - 4559 = 17.
\end{align*}
This is directly visible in the last word (line 419): the payload word at
\texttt{0x2be0} ends \texttt{\ldots 191f} followed by exactly 34 hex zeros, i.e.\
15 program bytes and 17 zero bytes. The last real program byte is at payload
offset $\texttt{0x2bee}$, and
$\texttt{0x2bef} = \texttt{0x1a20} + \texttt{0x11cf}$, so \yul{q\_end} is
exactly one past it. The trailing zeros are a codec artifact of word packing,
not an instruction: they are never executed, because the loop stops at
\yul{q\_end}.

\paragraph{The pad bytes are not a read-ahead guard.} The interpreter decodes
operands by loading a whole word and shifting --- e.g.\
\yul{shr(240, mload(q\_pc))} for a 2-byte pointer
(\texttt{Halo2Verifier.sol}:2214), and
\yul{let q\_pair\_word := mload(q\_pc)} followed by a 13-byte coefficient walk
(lines 2380--2387). The opcode fetch itself is word-wide too:
\yul{let q\_op := byte(0, mload(q\_pc))} at line 2205. Every such \yul{mload}
reads 32 bytes regardless of how few are consumed. The final loop iteration
has $\yul{q\_pc} = \yul{q\_end} - 1$ at most, i.e.\ payload \texttt{0x2bee}, so
its opcode fetch reads through payload \texttt{0x2c0d} ---
inside \texttt{fixed\_comms[0].x\_hi}. The over-read therefore happens on
every proof, not just in principle. This is harmless --- only the
high-order bytes of the loaded word are used, and \yul{q\_pc} is bounded by
\yul{q\_end} --- but a reader should not mistake the 17 pad bytes for a
32-byte read-ahead margin. They are 15 bytes short of one.

\begin{finding}[VK-6]{Program read-ahead can extend past the program section
into \texttt{fixed\_comms}}
\textbf{Severity:} Observation.\\
\emph{Location:} \texttt{Halo2VerifyingKey.sol}:419 (17 pad bytes);
\texttt{Halo2Verifier.sol}:2205 (word-wide opcode fetch), 2214, 2323, 2340,
2380--2387, 2407, 2421 (word-wide operand loads).\\
The packed program's tail padding (17 bytes) is smaller than the 32-byte width
of an \yul{mload}, so loads issued for the final instructions read into the
\texttt{fixed\_comms} region that follows at payload \texttt{0x2c00}. This is
not merely possible, it is certain: the loop's own opcode fetch is
\yul{let q\_op := byte(0, mload(q\_pc))}
(\texttt{Halo2Verifier.sol}:2205), which reads 32 bytes to extract one, so
every iteration with $\yul{q\_pc} > \yul{q\_end} - 32$ --- and in particular
the last, at payload \texttt{0x2bee} --- crosses the section boundary.\\
\emph{Why this is sound here:} the read is a memory read inside the already
\opcode{EXTCODECOPY}-populated VK window, never out of bounds of allocated
memory; only the top bytes of the loaded word are extracted by the shifts; and
\yul{q\_pc} advancement is driven by the decoded instruction widths, with the
terminal equality check at line 3098 catching any mis-tracking. No behaviour
depends on the padding value. Worth recording because a future change to the
codec (e.g.\ moving \texttt{fixed\_comms} before the program, or reducing the
final instruction's operand width) interacts with this.
\end{finding}

\begin{finding}[VK-7]{Program control flow cannot be influenced by calldata}
\textbf{Severity:} Informational (clean region).\\
The bytecode lives entirely inside the codehash-pinned VK runtime and is copied
to memory by \opcode{EXTCODECOPY} before any calldata is absorbed
(\texttt{Halo2Verifier.sol}:1393 runs before the transcript is opened with
\yul{vk\_digest} at 1481). The
interpreter reads opcodes and operands only from
$[\yul{q\_program\_mptr}, \yul{q\_end})$ and never from calldata. Combined with
the exact-termination check at line 3098, the program is a fixed, verifiable
straight-line schedule; a proof cannot redirect it.
\end{finding}

\subsection{Fixed-column and permutation commitments}
\label{sec:vk:comms}

\begin{lstlisting}[style=solidity,caption={\texttt{Halo2VerifyingKey.sol} lines 420--425: the first fixed-column commitment slot},label={lst:vk:fixed0}]
            // Fixed-column commitment 0, stored as one
            // EIP-2537 padded uncompressed G1 slot.
            mstore(add(payload, 0x2c00), 0x0000000000000000000000000000000016e98a681dca730dcbe651aec5493976) // fixed_comms[0].x_hi
            mstore(add(payload, 0x2c20), 0x11944be61932e7d19a63786871335939b0eefa7e32507b68b4ad6fd89a5c1028) // fixed_comms[0].x_lo
            mstore(add(payload, 0x2c40), 0x0000000000000000000000000000000010e6ae8d3becb251e60db7d788be7029) // fixed_comms[0].y_hi
            mstore(add(payload, 0x2c60), 0x8dbcdfdbb394da2ca9725e07485d6b630569a66ec867aa0d605573ae82d550df) // fixed_comms[0].y_lo
\end{lstlisting}

\noindent Lines 426--581 hold \texttt{fixed\_comms[1..26]} and lines 582--683
hold \texttt{permutation\_comms[0..16]}, all in the identical six-line
(two comment lines, four \yul{mstore}s) shape; they are elided.

\begin{lstlisting}[style=solidity,caption={\texttt{Halo2VerifyingKey.sol} lines 684--689: the last permutation commitment slot},label={lst:vk:permlast}]
            // Permutation commitment 17, also a 4-word
            // EIP-2537 padded uncompressed G1 slot.
            mstore(add(payload, 0x4200), 0x0000000000000000000000000000000002ecf71916494d46153fdd513ead1917) // permutation_comms[17].x_hi
            mstore(add(payload, 0x4220), 0x06401b48839713727c920f4591fe09ac3e8f4140ef1cd9e39b0d0bbe0b5b87d4) // permutation_comms[17].x_lo
            mstore(add(payload, 0x4240), 0x000000000000000000000000000000000e7bab8cc3c714f06826c8f3695dc63c) // permutation_comms[17].y_hi
            mstore(add(payload, 0x4260), 0xe7e931ff0dbe009cbb84fef0b99124905d90cdd6f5356a7ee7673548f6f0692c) // permutation_comms[17].y_lo
\end{lstlisting}

\begin{table}[h]
\centering
\begin{tabular}{@{}lrllll@{}}
\toprule
Group & Points & Words & Payload range & Memory range & Source lines \\
\midrule
\texttt{fixed\_comms[0..26]}        & 27 & 108 & \texttt{0x2c00}--\texttt{0x397f} & \texttt{0x6280}--\texttt{0x6fff} & 420--581 \\
\texttt{permutation\_comms[0..17]}  & 18 & 72  & \texttt{0x3980}--\texttt{0x427f} & \texttt{0x7000}--\texttt{0x78ff} & 582--689 \\
\bottomrule
\end{tabular}
\caption{Commitment sections. Each point is 4 words in the EIP-2537 padded
uncompressed layout \texttt{x\_hi, x\_lo, y\_hi, y\_lo}.}
\label{tab:vk:comms}
\end{table}

The $i$-th fixed commitment starts at payload
$\texttt{0x2c00} + \texttt{0x80}\,i$ and the $j$-th permutation commitment at
$\texttt{0x3980} + \texttt{0x80}\,j$; both are rendered explicitly, one
\yul{mstore} per word, each preceded by a two-line comment naming the index.
There is no loop and no compression. The counts 27 and 18 appear nowhere in
the payload as data --- they exist only as the number of rendered slots ---
so their meaning has to be read off the verifier's consumption of the slots:
all 27 \texttt{fixed\_comms} base addresses (\texttt{0x6280}, \texttt{0x6300},
\dots, \texttt{0x6f80}) and all 18 \texttt{permutation\_comms} base addresses
(\texttt{0x7000}, \dots, \texttt{0x7880}) occur as literal \yul{mcopy} sources
in \texttt{Halo2Verifier.sol}, so no slot is dead. Reading 27 as the circuit's
post-compression fixed-column count and 18 as the width of the copy-constraint
argument is consistent with the generator but is not checkable from either
contract.

\paragraph{Independent verification.} We decoded all 45 points and checked, for
each:
(i) the \texttt{hi} words are $< 2^{128}$, i.e.\ the EIP-2537 zero padding is
intact;
(ii) both coordinates are $< p$, i.e.\ canonically reduced;
(iii) $y^2 = x^3 + 4$ over $\Fp$;
(iv) $[r]P = \id$, i.e.\ membership in the prime-order subgroup.
All 45 points pass all four checks. None is the point at infinity, and no two
share an $x$-coordinate. As with the base points, the verifier does not re-check
these at runtime, which is sound for the reason given in Finding VK-4: they are
codehash-pinned, and the EIP-2537 precompiles that consume them
(\texttt{Halo2Verifier.sol}:3860, 3882 copy them into the final MSM buffer)
validate every input point themselves.

\subsection{The return, and the length cross-check}
\label{sec:vk:return}

\begin{lstlisting}[style=solidity,caption={\texttt{Halo2VerifyingKey.sol} lines 690--696: the terminating return},label={lst:vk:return}]

            // Return exactly the INVALID prefix plus the generated payload. The
            // linked verifier pins this byte length and the resulting codehash.
            return(runtime, 0x4281)
        }
    }
}
\end{lstlisting}

\paragraph{Arithmetic cross-check against
\texttt{EXPECTED\_VK\_PAYLOAD\_LENGTH}.} The assignment asks for this
explicitly; the result is \emph{consistent}. Counting rendered \yul{mstore}
statements by their trailing comment tag:

\begin{table}[h]
\centering
\begin{tabular}{@{}lrrr@{}}
\toprule
Section & Words & $\times 32$ & Cumulative bytes \\
\midrule
Header scalars (11) $+$ three base points ($4+8+8$) & 31 & 992 & 992 \\
\texttt{quotient\_const}                        & 178 & 5{,}696 & 6{,}688 \\
\texttt{quotient\_program}                      & 143 & 4{,}576 & 11{,}264 \\
\texttt{fixed\_comms} ($27 \times 4$)           & 108 & 3{,}456 & 14{,}720 \\
\texttt{permutation\_comms} ($18 \times 4$)     & 72  & 2{,}304 & 17{,}024 \\
\midrule
\textbf{Payload total}                          & \textbf{532} & \textbf{17{,}024} & $= \texttt{0x4280}$ \\
\opcode{INVALID} prefix                          & --- & 1 & \\
\textbf{Runtime total}                          & & \textbf{17{,}025} & $= \texttt{0x4281}$ \\
\bottomrule
\end{tabular}
\caption{Payload length derived by summing the rendered sections.}
\label{tab:vk:lengths}
\end{table}

Formally,
\begin{equation}
\ell_{\text{payload}}
 = 32\bigl(31 + 178 + 143 + 4 \cdot 27 + 4 \cdot 18\bigr)
 = 32 \cdot 532 = 17024,
\label{eq:vk:len}
\end{equation}
\begin{equation}
\ell_{\text{runtime}} = 1 + \ell_{\text{payload}} = 17025 = \texttt{0x4281}.
\label{eq:vk:runlen}
\end{equation}
Three independent corroborations:
\begin{enumerate}
\item The last \yul{mstore} offset is \texttt{0x4260} (line 689), and
$\texttt{0x4260} + \texttt{0x20} = \texttt{0x4280} = 17024$, matching
Equation~\eqref{eq:vk:len}.
\item The literal in the \yul{return} at line 693 is \texttt{0x4281} $= 17025$,
matching Equation~\eqref{eq:vk:runlen}, and equal to
\yul{EXPECTED\_VK\_LENGTH} at \texttt{Halo2Verifier.sol}:93.
\item $\texttt{17024} = \yul{EXPECTED\_VK\_PAYLOAD\_LENGTH}$ at
\texttt{Halo2Verifier.sol}:92, and
$\mptr{VK\_MPTR} + 17024 = \texttt{0x3680} + \texttt{0x4280} = \texttt{0x7900}
= \mptr{CHALLENGE\_MPTR}$, so the copied payload tiles the verifier's memory map
exactly.
\end{enumerate}

\textbf{Result of the cross-check: the declared payload length agrees with the
rendered contract in all three ways, and the reconstructed runtime hashes to the
verifier's pinned \yul{EXPECTED\_VK\_CODEHASH}.} There is no discrepancy to
report.

\subsection{Tamper-evidence}
\label{sec:vk:tamper}

The blob itself carries no integrity mechanism. Its tamper-evidence is entirely
external, and it is layered:

\begin{enumerate}
\item \textbf{Deploy-time pin.} The verifier's constructor requires
\yul{authorizedVk.code.length == EXPECTED\_VK\_LENGTH \&\& authorizedVk.codehash
== EXPECTED\_VK\_CODEHASH} before storing the address in the immutable
\yul{AUTHORIZED\_VK} (\texttt{Halo2Verifier.sol}:580--584). A verifier can never
be constructed against a VK whose runtime differs by one bit.
\item \textbf{Per-proof re-pin.} \yul{verifyProof} re-reads
\opcode{EXTCODESIZE} and \opcode{EXTCODEHASH} and compares both again on every
call, reverting \revert{VkMismatch} otherwise
(\texttt{Halo2Verifier.sol}:1386--1389). The generated comment gives the
rationale: the constructor check catches deployment mistakes, the fresh check
hardens against forks or same-transaction states where code at the authorized
address could differ from what was pinned. Since \yul{AUTHORIZED\_VK} is
\texttt{immutable} and code at an address is itself immutable post-Cancun
(\opcode{SELFDESTRUCT} no longer removes code outside the creating
transaction), the second check is defence in depth --- but it costs one cold
account access and closes the same-transaction \opcode{CREATE2} redeployment
window that existed pre-Cancun.
\item \textbf{Shape re-assertion.} Six header words are compared against
generated integers (\texttt{Halo2Verifier.sol}:1400--1406), giving a
\revert{VkMismatch} rather than an obscure downstream failure if a
structurally different VK were somehow pinned.
\item \textbf{Transcript binding.} \yul{vk\_digest} (payload word 0) is the
first value absorbed into the Fiat--Shamir transcript
(\texttt{Halo2Verifier.sol}:1481), so the challenges --- and therefore the proof
--- are bound to the constraint system this VK describes. A proof produced for a
different circuit cannot be replayed against this VK even if an attacker could
somehow swap the blob.
\end{enumerate}

\begin{finding}[VK-8]{The blob is tamper-evident to the strength of
\opcode{EXTCODEHASH}}
\textbf{Severity:} Informational (clean region).\\
Any modification of any of the 17{,}025 runtime bytes changes
\opcode{EXTCODEHASH} and is rejected both at verifier construction
(\texttt{Halo2Verifier.sol}:580--584) and on every proof
(\texttt{Halo2Verifier.sol}:1386--1389). We reproduced the pinned hash from the
source of this fixture, so the pin is known to correspond to exactly this
payload. The residual trust is not in the blob's integrity but in the
\emph{content} of two of its parts: \texttt{NEG\_S\_G2\_BASE}
(Finding VK-3) and the correctness of the generator that produced the
constants, program and commitments.
\end{finding}

\subsection{The \texttt{0x80} construction buffer}
\label{sec:vk:memory}

\yul{runtime := 0x80} places the transient buffer at the first byte of
Solidity's allocatable memory, immediately above the reserved words: scratch
space \texttt{0x00}--\texttt{0x3f}, the free-memory pointer at \texttt{0x40},
and the zero slot at \texttt{0x60}. The address is a generator constant,
\texttt{VK\_CONSTRUCTOR\_PAYLOAD\_START} $=$
\texttt{SOLIDITY\_ALLOCATABLE\_MEMORY\_START} $=$ \texttt{0x80}
(\texttt{src/lowering/layout/mod.rs}:25, 59); the only check that runs at
generation time is
\texttt{src/lowering/layout/memory.rs}:629--635, which rejects drift of the
rendered pointer. (\texttt{mod.rs}:818 asserts the same value but is a unit
test, not a generation-time guard.) The fixture renders it as the literal
\texttt{0x80}.

\paragraph{Written extent.} The stores cover
\[
[\texttt{0x80},\ \texttt{0x80} + \texttt{0x4281}) = [\texttt{0x80},\
\texttt{0x4301}),
\]
so the EVM's memory word count after construction is
$\lceil \texttt{0x4301}/32 \rceil = 537$ words, i.e.\ a memory-expansion cost of
$3 \cdot 537 + \lfloor 537^2/512 \rfloor = 1611 + 563 = 2174$ gas --- entirely
in the linear regime, negligible next to the 532 \opcode{MSTORE}s and the
$200 \times 17025 \approx 3.4$M gas of code-deposit cost.

\paragraph{Alignment.} Because \yul{payload = 0x81}, every payload
\yul{mstore} is at an odd address. EVM memory is byte-addressed, so this is
legal and carries no extra cost; consecutive stores at \yul{payload+0x00} and
\yul{payload+0x20} tile exactly with no gap and no overlap.

\begin{finding}[VK-9]{The \texttt{0x80} buffer is memory-safe as written}
\textbf{Severity:} Informational (clean region).\\
\emph{Location:} \texttt{Halo2VerifyingKey.sol}:54--55, 62--64, 693.\\
Four properties together make the unmanaged buffer safe:
(i) the constructor body is a \emph{single} assembly block that is terminal ---
it ends in \yul{return}, so no Solidity code runs afterwards and no
Solidity-managed allocation can ever be handed an address inside
$[\texttt{0x80}, \texttt{0x4301})$;
(ii) the block never reads or writes the free-memory pointer at \texttt{0x40},
nor the scratch words at \texttt{0x00}--\texttt{0x3f} or the zero slot at
\texttt{0x60}, so Solidity's reserved region is preserved (which matters only
for hypothetical future code, but is the documented contract);
(iii) the constructor takes no arguments and the contract has no state
variables, so nothing has allocated memory before the block begins --- the
free-memory pointer is still at its initial \texttt{0x80};
(iv) every byte of the returned region is explicitly written, so no
uninitialised memory is deposited as code.\\
The one structural caveat worth recording: property (i) is what makes the
absent free-memory-pointer bump harmless. If a future revision added any
Solidity statement after the assembly block, or lifted this body into a
non-terminal helper, the buffer would silently overlap the next allocation.
The analogous risk on the verifier side is guarded at runtime
(\texttt{Halo2Verifier.sol}:674--677 reverts \revert{MemoryLayoutViolated} if
\yul{mload(0x40)} has reached the generated layout); the VK constructor has no
such guard, relying instead on the generator-side pointer assertion. Given that
this constructor is 100\% generated and terminal, that is a reasonable trade.
\end{finding}

\subsection{Size budget}
\label{sec:vk:size}

\begin{table}[h]
\centering
\begin{tabular}{@{}lrrr@{}}
\toprule
Quantity & This fixture & Limit & Headroom \\
\midrule
Deployed runtime (EIP-170)   & 17{,}025 B & 24{,}576 B & 7{,}551 B (30.7\%) \\
Creation code (EIP-3860)     & $\approx 21$ kB (estimated) & 49{,}152 B & $\approx 28$ kB \\
\bottomrule
\end{tabular}
\caption{Code-size budget. The creation-code figure is an estimate: 532 stores
at roughly 39 bytes each (\opcode{PUSH32} immediate, \opcode{PUSH2} offset,
\opcode{DUP}, \opcode{ADD}, \opcode{MSTORE}) plus preamble and solc metadata.}
\label{tab:vk:size}
\end{table}

7{,}551 bytes of EIP-170 headroom is 236 words: room for about 59 more G1
commitments, or 236 more quotient constants, or 7{,}551 more packed program
bytes --- but not all three. The generator enforces the bound at generation
time rather than letting deployment fail
(\texttt{src/lowering/render/models.rs}:271--332, with
\texttt{EIP\_170\_MAX\_RUNTIME\_BYTES} $= \texttt{0x6000}$ and an explicit
error message telling the operator to reduce the circuit's
constant/commitment count).

\begin{finding}[VK-10]{EIP-170 headroom is finite and the pattern has no
overflow path}
\textbf{Severity:} Low.\\
\emph{Location:} \texttt{Halo2VerifyingKey.sol}:693 (\texttt{0x4281});
\texttt{src/lowering/render/models.rs}:271--332 (the guard).\\
The single-data-contract pattern places a hard ceiling on circuit size: a
circuit needing more than 24{,}575 payload bytes cannot be expressed at all,
and the failure appears at code-generation time with no fallback (there is no
sharding across multiple VK contracts, and no SSTORE2-style chunking). At
17{,}025/24{,}576 this fixture uses 69.3\% of the budget. This is a
maintainability and roadmap risk rather than a security defect --- the guard
fails closed and loudly --- but it should be tracked, since growth in the
quotient program or the number of fixed columns consumes it directly.\\
We found no equivalent generator-side guard for the EIP-3860 initcode limit;
the estimated 21 kB creation code leaves a wide margin, and the initcode limit
is 2$\times$ the runtime limit, so a payload that fits EIP-170 will comfortably
fit EIP-3860 under this rendering strategy. No action needed beyond noting the
asymmetry.
\end{finding}

\subsection{Fixture versus template}
\label{sec:vk:template}

We compared the fixture against its generator template,
\texttt{templates/contracts/Halo2VerifyingKey.sol}. The two agree in every
respect: the same SPDX header, the same pinned
\texttt{pragma solidity 0.8.30}, the same NatSpec layout block, the same
\yul{runtime}/\yul{payload}/\yul{mstore8(runtime, 0xfe)} prologue, the same
per-slot comment convention (\texttt{fixed\_comms[i].x\_hi} and friends), and a
\yul{return} whose length is rendered as
$1 + 32(|\text{constants}| + 4|\text{fixed}| + 4|\text{perm}|)$, which
evaluates to \texttt{0x4281} for this fixture. The template's
\texttt{constructor\_payload\_mptr} placeholder renders to the fixture's
\texttt{0x80}, matching
\texttt{VK\_CONSTRUCTOR\_PAYLOAD\_START} $=$
\texttt{SOLIDITY\_ALLOCATABLE\_MEMORY\_START} $=$ \texttt{0x80}
(\texttt{src/lowering/layout/mod.rs}:25, 59). No divergence between fixture and
template was found, so nothing in this section rests on template intent over
rendered fact.

\subsection{Summary of this section's findings}
\label{sec:vk:summary}

\begin{table}[h]
\centering
\begin{tabular}{@{}llp{8.0cm}@{}}
\toprule
ID & Severity & Subject \\
\midrule
VK-1  & Informational & \opcode{INVALID} prefix is a complete entry-point barrier (clean) \\
VK-2  & Low           & Payload is not self-describing; the codehash pin covers the VK half only \\
VK-3  & Informational & \texttt{NEG\_S\_G2\_BASE} is unattested trusted-setup data \\
VK-4  & Informational & Embedded points are not re-validated, and need not be (clean) \\
VK-5  & Observation   & Byte-wide constant index exceeds the 178-entry table's range \\
VK-6  & Observation   & Program operand read-ahead can reach into \texttt{fixed\_comms} \\
VK-7  & Informational & Program control flow is calldata-independent (clean) \\
VK-8  & Informational & Blob is tamper-evident to \opcode{EXTCODEHASH} strength (clean) \\
VK-9  & Informational & \texttt{0x80} construction buffer is memory-safe (clean) \\
VK-10 & Low           & Finite EIP-170 headroom, no overflow path \\
\bottomrule
\end{tabular}
\caption{Findings raised in Section~\ref{sec:vk:contract}.}
\label{tab:vk:findings}
\end{table}

The contract is, in the main, clean: it is 696 lines of which 532 are pure
data statements, its one non-trivial behaviour (assemble, prefix, return) is
correct and memory-safe, its declared length is consistent with its contents
three ways over, and its pinned codehash reproduces exactly from source. The
security-relevant weight of this file lies not in its logic but in the
provenance of the values it carries --- above all the eight
\texttt{NEG\_S\_G2\_BASE} words.


%% ==== 04-abi-entry.tex ==================================================
\section{\texorpdfstring{\yul{verifyProof} entry: ABI-shape guard, memory-safety guard, and low-level helpers}{verifyProof entry: ABI-shape guard, memory-safety guard, and low-level helpers}}
\label{sec:abi:top}

This section specifies \texttt{Halo2Verifier.sol} lines 588--722: the documented
external contract of \yul{verifyProof}, the two guards that run before any
generated memory is touched, and the first two Yul helpers defined inside the
main assembly block (\yul{fail} and \yul{scalar\_inv}). The streaming Keccak
transcript helpers that begin at line 724 are covered elsewhere.

\subsection{External interface and calldata map}
\label{sec:abi:iface}

The contract exposes exactly one verification entry point,
\[
  \yul{verifyProof(bytes\ calldata\ proof,\ uint256[]\ calldata\ instances)}
  \;\longrightarrow\; \yul{bool},
\]
declared \yul{external view} at lines 624--627. Its ABI selector is
$\mathrm{bytes4}(\mathrm{keccak256}(\text{\texttt{"verifyProof(bytes,uint256[])"}}))
= \texttt{0x1e8e1e13}$.

Neither Solidity parameter is ever read in the function body. The whole verifier
is a hand-rolled calldata parser that indexes calldata through the generated
absolute cursor constants of lines 100--103. No \emph{soundness} property of
this contract therefore depends on solc's ABI decoder: every shape property the
parser needs is asserted explicitly, either in the guard of
Section~\ref{sec:abi:headguard} or in the envelope block at lines 1419--1430
(covered elsewhere).

The decoder is nevertheless emitted and it runs \emph{before} the function body,
so it determines the observable failure mode for a whole class of malformed
calldata. Writing $h_0 = \opcode{CALLDATALOAD}(\texttt{0x04})$ and $h_1 =
\opcode{CALLDATALOAD}(\texttt{0x24})$ for the two ABI head words (the same
notation used in Section~\ref{sec:abi:headguard}), compiling this signature and
the line-635 block with \yul{solc --optimize --via-ir} shows the dispatcher
enforcing, in this order and each with a bare \yul{revert(0x00, 0x00)}:
\opcode{CALLDATASIZE} $\ge
\texttt{0x44}$; $h_0 \le 2^{64}-1$; $h_0 + \texttt{0x24} \le
\opcode{CALLDATASIZE}$; $\yul{proof.length} \le 2^{64}-1$; $h_0 +
\texttt{0x24} + \yul{proof.length} \le \opcode{CALLDATASIZE}$; and the four
analogous bounds for $h_1$ and $32\cdot\yul{instances.length}$. Legacy codegen
emits the same checks. Consequences: the decoder is redundant with (and strictly
weaker than) the generated guards, so it changes nothing about which calldata is
accepted; but every input it rejects reverts with \emph{empty} return data, not
with a typed selector (Finding~abi:F8).

\begin{table}[htbp]
\centering
\caption{Generated calldata cursors (\texttt{Halo2Verifier.sol} lines 100--103) and the
resulting fixed calldata map. ``Pinned at'' gives the line that forces the value.}
\label{tab:abi:calldata}
\begin{tabularx}{\textwidth}{@{}llXll@{}}
\toprule
Offset & Width & Content & Value & Pinned at \\
\midrule
\texttt{0x00}   & 4 bytes        & ABI selector                                  & \texttt{0x1e8e1e13} & dispatcher \\
\texttt{0x04}   & \texttt{0x20}  & head word 0: offset of \yul{proof} tail       & \texttt{0x40}       & line 636 \\
\texttt{0x24}   & \texttt{0x20}  & head word 1: offset of \yul{instances} tail   & \texttt{0x1ec0}     & line 636 \\
\texttt{0x44}   & \texttt{0x20}  & \yul{proof.length} (\mptr{PROOF\_LEN\_CPTR})  & \texttt{0x1e60}     & line 1419 \\
\texttt{0x64}   & \texttt{0x1e60}& proof payload (\mptr{PROOF\_CPTR})            & ---                 & line 1785 \\
\texttt{0x1ec4} & \texttt{0x20}  & \yul{instances.length} (\mptr{NUM\_INSTANCE\_CPTR}) & $19$          & line 1420 \\
\texttt{0x1ee4} & \texttt{0x260} & 19 instance words (\mptr{INSTANCE\_CPTR})     & ---                 & line 1511 \\
\midrule
\multicolumn{3}{@{}l}{Total \opcode{CALLDATASIZE}}
  & \texttt{0x2144} & line 1426 \\
\bottomrule
\end{tabularx}
\end{table}

The map is internally consistent and fully determined:
\begin{align}
  \mptr{PROOF\_LEN\_CPTR} &= \texttt{0x04} + \texttt{0x40} = \texttt{0x44},\\
  \mptr{PROOF\_CPTR}      &= \mptr{PROOF\_LEN\_CPTR} + \texttt{0x20} = \texttt{0x64},\\
  \mptr{NUM\_INSTANCE\_CPTR} &= \mptr{PROOF\_CPTR} + \texttt{0x1e60} = \texttt{0x1ec4},\\
  \mptr{INSTANCE\_CPTR}   &= \mptr{NUM\_INSTANCE\_CPTR} + \texttt{0x20} = \texttt{0x1ee4},\\
  \opcode{CALLDATASIZE}   &= \mptr{INSTANCE\_CPTR} + 19\cdot\texttt{0x20} = \texttt{0x1ee4} + \texttt{0x260} = \texttt{0x2144}.
\end{align}
The proof payload length \texttt{0x1e60} $= 7776 = 243\cdot 32$ is an exact
multiple of the word size, so the canonical ABI encoding inserts no tail padding
between the proof bytes and the \yul{instances} length word. This is what makes
the constant $\mptr{NUM\_INSTANCE\_CPTR}=\mptr{PROOF\_CPTR}+\yul{proof.length}$
correct without a rounding term; a future proof layout whose byte length is not
word-aligned would need $\lceil\cdot\rceil_{32}$ here.

\subsection{The documented NatSpec contract}
\label{sec:abi:natspec}

Lines 588--623 state four obligations. They are normative for integrators, so
they are reproduced here as claims together with where the code discharges them.

\paragraph{(N1) Scope.} Lines 589--597: the function proves only that
\yul{proof} verifies for \yul{instances} under the pinned VK and protocol.
Semantic binding of the instances (state roots, program identifiers, IVC
outputs, chain/domain separation, authorization) is explicitly \emph{out of
scope} and delegated to an application wrapper; wrapper obligations
(replaceable verifier address, wrapper-held pause, chain-id/address/anti-replay
binding) are declared requirements of
\texttt{docs/reference/DEPLOYMENT\_AND\_INCIDENT\_RESPONSE.md}. Nothing in the
assigned range binds a chain id, a contract address, or a nonce; the contract is
correctly described as a raw verifier ABI.

\paragraph{(N2) Success-or-revert.} Lines 598--605: accepted proofs return
\yul{true}; the function \emph{never} returns \yul{false}. The comment
enumerates seven typed errors (\revert{BadCalldataShape}, \revert{VkMismatch},
\revert{NonCanonicalScalar}, \revert{BadPointEncoding},
\revert{PrecompileFailed}, \revert{ProofRejected},
\revert{QuotientProgramInvalid}). That list is incomplete as a statement about
call-time behaviour: the eighth declared error, \revert{MemoryLayoutViolated}
(line 83), is raised at line 675 inside \yul{verifyProof} itself, and bare
\yul{revert(0x00, 0x00)}s from solc's decoder precede the body entirely
(Section~\ref{sec:abi:iface}, Finding~abi:F8). The omission is
defensible -- line 44--48 classifies \revert{MemoryLayoutViolated} as a BUILD
fault, not a proof fault -- but an integrator writing an exhaustive
\yul{catch} from this comment alone would miss two cases.
The success-or-revert half of the claim is verifiable from the fixture: the
only \opcode{RETURN} reachable from the main assembly block is line 4128,
\yul{return(RETURN\_MPTR, 0x20)}, preceded by \yul{mstore(RETURN\_MPTR, 1)} at
line 4127; every other exit is a \opcode{REVERT}. The Solidity epilogue that
would return the default-initialised \yul{false} is unreachable because the
assembly block is terminal (Section~\ref{sec:abi:memsafe}).

\paragraph{(N3) Exact calldata size.} Lines 606--611: calldata must be exactly
selector $\|$ proof bytes $\|$ instance words; \opcode{CALLDATASIZE} is pinned
and any trailing byte reverts \revert{BadCalldataShape}. \textbf{The pin itself
is not in this range}: the guard at line 636 constrains only the two head words.
Exact size equality is asserted at line 1426,
\yul{eq(calldatasize(), add(INSTANCE\_CPTR, 0x0260))}, i.e.\ against
\texttt{0x2144}, and folded into \yul{success} which is enforced at line 1430.
The claim is accurate for a \emph{suffix}: solc's decoder bounds the tails with
$\le$, not $=$, so appended bytes pass the decoder and are caught at 1426 with
the typed selector. It is not accurate for \emph{truncated} or otherwise
ill-formed calldata, which the decoder rejects first with empty revert data.

\paragraph{(N4) Absolute memory addressing.} Lines 612--620: the verifier uses
absolute Yul addresses instead of the Solidity free-memory pointer, generated
scratch starts at \mptr{TRANSCRIPT\_MPTR}, and the assembly block asserts the
separation on entry rather than assuming it. Discharged by line 674
(Section~\ref{sec:abi:mf2}).

\paragraph{(N5) Pre-repacked proof encoding.} Lines 621--623, easy to overlook
because they are \yul{@param}/\yul{@return} tags rather than \yul{@dev} prose,
carry a hard integration requirement: \yul{proof} is ``Solidity-facing proof
bytes, with G1 elements repacked into EIP-2537 padded uncompressed form'' and
\yul{instances} are ``canonical BLS12-381 scalar-field words''. The native
midnight-proofs serialization is therefore \emph{not} accepted; an off-chain
shim must rewrite it first (the fixture names this optimisation H3 at lines
1689--1690). Nothing in the assigned range checks the repacking -- the encoding
is validated later, per element, by the G1 codec and the scalar canonicality
checks, and any mismatch surfaces as \revert{BadPointEncoding} or
\revert{NonCanonicalScalar} rather than as a distinguishable ``wrong
serialization'' error. \yul{@return} restates (N2).

\begin{finding}[abi:F1]{Success-or-revert makes \yul{if (!verifyProof(...))} dead code}
\textbf{Severity: Informational.} \texttt{Halo2Verifier.sol}:598--605, 4113--4128.
The function is total on its accept path only: the boolean return channel carries
the single value \yul{true}. An integrator writing the idiomatic
\yul{if (!verifier.verifyProof(p, i)) \{ revert Invalid(); \}} gets a branch that
can never be taken -- not a vulnerability by itself, but it means the
integrator's own rejection path (custom event, refund, fallback) is never
exercised and never tested. The bubbling revert also propagates the verifier's
selector rather than the wrapper's. The hazard is explicitly documented at lines
602--605, and the correct integration patterns are
\yul{require(verifier.verifyProof(...))} (equivalent, since the false branch is
unreachable) or a \yul{try}/\yul{catch} that decodes the revert data -- which
must tolerate \emph{empty} revert data as well as a four-byte selector, see
Finding~abi:F8. No code change is implied; this is recorded so a reviewer
confirms the wrapper under audit does not rely on a \yul{false} return.
\end{finding}

\begin{finding}[abi:F2]{Exact-calldatasize policy excludes calldata-appending relayers}
\textbf{Severity: Observation.} \texttt{Halo2Verifier.sol}:606--611 (policy),
1424--1427 (enforcement). Because \opcode{CALLDATASIZE} must equal
\texttt{0x2144} exactly, ERC-2771 forwarders -- which append a 20-byte
\yul{\_msgSender()} suffix, giving \texttt{0x2158} -- are rejected with
\revert{BadCalldataShape}, as are multicall wrappers and paymaster contexts that
append context bytes. This is a deliberate, sound trade: the alternative
(accepting a suffix) would let a relayer append arbitrary bytes that the
hand-rolled parser never reads, weakening the ``the caller saw exactly what the
verifier parsed'' property and creating an unhashed, unauthenticated calldata
channel. The cost is purely integrational and is documented: such traffic must
be routed through an application wrapper that reassembles exact calldata. The
failure is loud, deterministic, and input-independent, so it cannot be
discovered late.
\end{finding}

\begin{finding}[abi:F8]{The typed-error contract is not total: solc's decoder reverts with empty data}
\textbf{Severity: Low.} \texttt{Halo2Verifier.sol}:598--611 (the claim),
624--627 (the signature), 635--642 (the first body statement).
The NatSpec at lines 598--605 tells integrators that ``every rejection reverts
with one of the typed errors declared above'' and lines 606--611 attribute
mis-shaped calldata specifically to \revert{BadCalldataShape}. Neither statement
covers the code that actually runs first. Because \yul{verifyProof} is declared
with two \yul{calldata} array parameters, solc emits an ABI decoder in the
dispatcher, ahead of the function body, that enforces
$\opcode{CALLDATASIZE} \ge \texttt{0x44}$, $h_0,h_1 \le 2^{64}-1$, both length
words in bounds, and both tails in bounds -- each with a bare
\yul{revert(0x00, 0x00)}. This was confirmed by compiling the signature and the
line-635 block with \yul{solc --optimize --via-ir} and reading the emitted
dispatcher; legacy codegen emits the same checks. A second untyped path exists in
the constructor (bare reverts in the precompile probes, lines 405--407 onward),
and a third typed-but-undocumented one at line 675
(\revert{MemoryLayoutViolated}, absent from the NatSpec's list of seven).
Impact is confined to error reporting -- every one of these paths is a revert, so
no invalid proof is accepted -- but an integrator who writes
\yul{catch (bytes memory e)} and switches on \yul{bytes4(e)} from the
documentation alone will index an empty \yul{bytes} for a large, easily reached
class of malformed calldata. The documentation, not the code, is what should
change: the NatSpec should state that malformed ABI envelopes may revert without
return data.
\end{finding}

\subsection{The head-word guard}
\label{sec:abi:headguard}

The first executable statement of the function is a small, self-contained
assembly block that validates the two ABI head words before any generated memory
work, before the VK \opcode{EXTCODECOPY} of line 1393, and before the transcript
parser starts.

\begin{lstlisting}[style=solidity,caption={\texttt{Halo2Verifier.sol} lines 624--642: function head and the ABI-shape guard},label={lst:abi:headguard}]
    function verifyProof(
        bytes calldata proof,
        uint256[] calldata instances
    ) external view returns (bool) {
        // Cheap ABI-shape guard before any generated memory work:
        //   - proof head must point at the bytes payload;
        //   - instances head must point at the generated instance array.
        //
        // The verifier below is a hand-rolled calldata parser. Failing here
        // keeps malformed dynamic-argument layouts from being interpreted as a
        // valid Midfall proof stream.
        assembly ("memory-safe") {
            if iszero(and(eq(calldataload(0x04), 0x40), eq(calldataload(0x24), sub(NUM_INSTANCE_CPTR, 0x04)))) {
                // BadCalldataShape() -- fail() is not in scope in this early
                // guard block, so write the selector inline.
                mstore(0x00, shl(224, ERR_BAD_CALLDATA_SHAPE))
                revert(0x00, 0x04)
            }
        }
\end{lstlisting}

\paragraph{Predicate.} Writing $h_0 = \opcode{CALLDATALOAD}(\texttt{0x04})$ and
$h_1 = \opcode{CALLDATALOAD}(\texttt{0x24})$, execution continues iff
\begin{equation}
  \label{eq:abi:headpred}
  \bigl(h_0 = \texttt{0x40}\bigr) \;\wedge\;
  \bigl(h_1 = \mptr{NUM\_INSTANCE\_CPTR} - \texttt{0x04} = \texttt{0x1ec0}\bigr),
\end{equation}
and otherwise the call reverts with the four bytes
$\texttt{0x1b99e37c} = \revert{BadCalldataShape}$. ABI head words are byte
offsets measured from the end of the selector, so \eqref{eq:abi:headpred} says
precisely: the \yul{proof} tail begins at absolute calldata offset
$\texttt{0x04}+\texttt{0x40}=\mptr{PROOF\_LEN\_CPTR}$, and the \yul{instances}
tail begins at $\texttt{0x04}+\texttt{0x1ec0}=\mptr{NUM\_INSTANCE\_CPTR}$.

\paragraph{Behaviour on short calldata.} Considered in isolation the guard is
total: \opcode{CALLDATALOAD} zero-extends past \opcode{CALLDATASIZE}, so a call
carrying only the selector would read $h_0=h_1=0$ and fail
\eqref{eq:abi:headpred}. The guard therefore needs no separate length
precondition to be memory- or control-safe. In the deployed artifact that path is
unreachable, however: solc's decoder already requires $\opcode{CALLDATASIZE} \ge
\texttt{0x44}$ and both tails in bounds, and reverts with empty data before the
body runs (Section~\ref{sec:abi:iface}). Truncated calldata therefore fails
closed, but \emph{untyped}; the guard's own zero-extension argument is
defence in depth against a future build in which the decoder is absent or
weakened.

\paragraph{What this buys.} The guard compiles to two \opcode{CALLDATALOAD}s and
a handful of stack ops (two constant pushes, two \opcode{EQ}, one
\opcode{AND}, one \opcode{ISZERO}, and the branch) -- the
\yul{sub(NUM\_INSTANCE\_CPTR, 0x04)} is constant-folded, so no runtime
arithmetic on attacker data occurs. It runs \emph{before} the 17\,024-byte VK
\opcode{EXTCODECOPY} (line 1393) and its \yul{EXTCODEHASH} comparison. Rejecting
a mis-shaped envelope here avoids that copy and its memory expansion. More
importantly it is a semantic guard, not an optimisation: because the body never
dereferences the Solidity \yul{proof}/\yul{instances} slices, an ABI encoding
that placed the two tails anywhere other than the generated offsets would be
silently reinterpreted as a Midfall proof stream at the hard-coded cursors. The
guard turns that class of confusion into a typed revert.

\paragraph{What it costs, and what it does not check.} The guard pins offsets
only. It does \emph{not} check \yul{proof.length}, \yul{instances.length}, or
\opcode{CALLDATASIZE}; those are asserted later at lines 1419, 1420 and 1426
respectively, and the parser's own terminal position check at line 1785
(\yul{eq(proof\_cptr, NUM\_INSTANCE\_CPTR)}) closes the loop a third time. The
region is therefore layered rather than complete on its own, which the comment
at lines 628--634 states accurately.

\begin{finding}[abi:F3]{Head-word guard: clean}
\textbf{Severity: Observation (no defect).} \texttt{Halo2Verifier.sol}:635--642.
The block is sound for three independent reasons worth recording. (i) The
\yul{and(...)} at line 636 combines two \opcode{EQ} results, each of which is
definitionally $0$ or $1$, so the bitwise \opcode{AND} is a correct logical
conjunction -- this is the one shape of \yul{and} that is safe, and the repo
lints for the unsafe shape (commit 06537c3a). (ii) The revert payload is built
with \yul{shl(224, ERR\_BAD\_CALLDATA\_SHAPE)} into scratch at \texttt{0x00} and
returned as exactly four bytes, matching Solidity's custom-error ABI; the
constant is a genuine four-byte value so no high bits are lost in the shift.
(iii) The block's only memory effect is a write inside $[\texttt{0x00},
\texttt{0x40})$, which is unconditionally permitted for a \yul{memory-safe}
block, so this annotation needs none of the reasoning of
Section~\ref{sec:abi:memsafe}. The one caveat is scope, not correctness: the
guard is the first statement of the \emph{body}, but solc's ABI decoder runs
before the body and rejects a strict subset of the same malformed inputs with an
untyped revert, so this guard is not the first line of defence it reads as.
\end{finding}

\subsection{The \yul{memory-safe} annotation on the main block (MF-12)}
\label{sec:abi:memsafe}

Between the two guards, line 646 loads the pinned VK address into a local,
\yul{address vk = AUTHORIZED\_VK;}. The main assembly block then opens at line
647, again annotated \yul{("memory-safe")}, and runs to line 4129. The
in-fixture comment at lines 648--673 argues the annotation, and the assignment
requires that argument to be assessed rigorously rather than repeated.

The comment at lines 643--645 states the role of that local: this is a
non-embedded render, so the VK is pinned by address \emph{and} codehash, the Yul
loader rechecks the runtime before every proof, and it copies the
\yul{INVALID}-prefixed payload into \mptr{VK\_MPTR} (lines 1386--1393, covered
elsewhere). Reading the immutable in Solidity rather than in Yul is what makes
\yul{vk} available to the assembly block at all.

\paragraph{What the annotation does.} \yul{memory-safe} is a promise \emph{to}
the compiler, and in exchange solc enables the stack-to-memory mover, which
relocates stack slots that would otherwise trigger ``stack too deep'' into a
reserved memory region growing upward from \texttt{0x80} and raises the initial
value of the free-memory pointer past it. For a verifier body of this size the
mover is effectively not optional; the file header's discussion of a measured
spill reservation (lines 5--7) presupposes it. This specification did not
recompile the fixture to confirm that the body fails to compile without the
annotation.

\paragraph{The promise, and the block's actual memory effects.} Solidity's
memory-safety contract permits an annotated block to touch only four kinds of
memory: (1) memory it allocated itself through the free-memory pointer;
(2) memory Solidity allocated, within the bounds of a referenced object;
(3) the scratch space $[\texttt{0x00},\texttt{0x40})$; and (4) temporary memory
located \emph{after} the value of the free-memory pointer as read at the start
of the block, used without updating that pointer. Enumerating every memory write
in the main block of the fixture:
\begin{itemize}
  \item $[\texttt{0x00},\texttt{0x04})$ -- the three inline selector stores at
        lines 675, 691 (inside \yul{fail}) and 1891. All lie in scratch:
        category (3).
  \item $[\texttt{0x1000},\dots)$ -- every other write. A mechanical scan of
        every \opcode{MSTORE}, \opcode{MSTORE8}, \opcode{MCOPY},
        \opcode{CALLDATACOPY}, \opcode{RETURNDATACOPY} and \opcode{EXTCODECOPY}
        destination, and of every \opcode{STATICCALL} output pointer, over lines
        647--4129 finds exactly three constant destinations below
        \texttt{0x1000} -- the three scratch stores of the previous bullet -- and
        nothing at all in $[\texttt{0x40},\texttt{0x1000})$. Every remaining
        destination is either a \texttt{*\_MPTR} constant, an offset from one, or
        a bare hex literal. The generated body contains $124$ distinct bare
        literals used as memory operands -- for example \texttt{0x3580} (line
        710, the \yul{scalar\_inv} frame), \texttt{0x1240} (line 969, the
        \yul{ec\_pairing} staging area), \texttt{0xb8e0} (line 2172, the
        quotient-VM stack base), and \texttt{0x1000}, \texttt{0x1080} and
        \texttt{0x1100} (lines 4012--4033, the PCS pairing staging area) -- and
        every one of them lies in $[\texttt{0x1000}, \texttt{0xce00}]$. The
        lowest address written outside scratch is thus \mptr{TRANSCRIPT\_MPTR}
        $= \texttt{0x1000}$.
  \item The free-memory pointer word at \texttt{0x40} is \emph{read} (line 674)
        and never written.
\end{itemize}
Category (4) therefore applies to the second bullet \emph{provided} the
free-memory pointer at block entry is $\le \texttt{0x1000}$ -- which is exactly
what the guard of Section~\ref{sec:abi:mf2} enforces, as the block's first
statement. Under that guard the annotation is defensible by the letter of the
contract, not merely in spirit: the block writes scratch and memory at or above
the entry free-memory pointer, and nothing else.

\paragraph{The three stated side conditions, weighed.} The comment at lines
657--665 claims three properties must hold together. They are not equally
load-bearing:
\begin{description}
  \item[Terminality.] Necessary, and verified. The block never updates
        \texttt{0x40} after writing above it, so if control returned to Solidity
        the compiler would be free to allocate over the generated arena. It does
        not return: every path ends in \opcode{REVERT}, or in
        \yul{return(RETURN\_MPTR, 0x20)} at line 4128, and the function's closing
        braces at lines 4129--4130 admit no further statement.
  \item[Pinned pragma.] Line 13 is \yul{pragma solidity 0.8.30;}, an exact pin,
        not a caret range. Its primary justification is not memory safety but
        deployability: the header at lines 8--11 records that 0.8.24 at
        \texttt{runs=1} emits 29\,567 runtime bytes and 0.8.30 at
        \texttt{runs=100000} emits 29\,836, both above the EIP-170 limit, so only
        the pinned (version, runs) pair yields a deployable artifact. For memory
        safety the pin is defence in depth, because the runtime guard measures the
        real reservation rather than assuming it.
  \item[Fail-closed guard.] The genuinely load-bearing condition; see
        Section~\ref{sec:abi:mf2}.
\end{description}

\paragraph{Assessment of the ``the annotation is a lie'' comment.} The comment at
lines 657--659 is more pessimistic than the code warrants. ``Writes memory it
never allocated through the free-memory pointer'' is precisely what category (4)
sanctions; the annotation is not a lie under the published contract once the
guard establishes the ordering. The residual imprecision is narrower: category
(4) is written for memory used \emph{temporarily} and returned unobserved to
Solidity, whereas this block never returns at all, so the ``temporary'' qualifier
is vacuous here rather than violated. Erring conservative in a comment is not a
defect, and the conservatism keeps a future maintainer from reusing the body in a
non-terminal context; the assessment is recorded only so a reviewer does not
mistake the comment for an admission of an unfixed problem.

\begin{finding}[abi:F4]{The \yul{memory-safe} annotation is sound, but conditionally so}
\textbf{Severity: Informational.} \texttt{Halo2Verifier.sol}:647--677, 13, 4125--4130.
The annotation on the main block is justified, and the justification is
mechanised on chain by the free-memory-pointer guard at line 674 rather than
argued only in prose or in a generator-side test. The condition set is real and
should be re-checked by any reviewer who touches the file: the block must remain
terminal, the guard must remain the first statement of the block, and no write
below \mptr{TRANSCRIPT\_MPTR} other than scratch $[\texttt{0x00},\texttt{0x40})$
may be introduced. Violating any of these turns a compiler assumption into
silent memory corruption of Solidity-owned data, which is the worst available
failure mode for a verifier. No violation exists in this fixture.
\end{finding}

\subsection{The free-memory-pointer guard (MF-2)}
\label{sec:abi:mf2}

\begin{lstlisting}[style=solidity,caption={\texttt{Halo2Verifier.sol} lines 647--677: main assembly block opening and the MF-2 guard (comment retained verbatim; it is the fixture's own argument)},label={lst:abi:mf2}]
        assembly ("memory-safe") {
            // The `memory-safe` annotation above is what enables solc's
            // stack-to-memory mover, which reserves spill slots upward from
            // 0x80. The generated layout below writes absolute addresses from
            // TRANSCRIPT_MPTR upward and never consults the free-memory
            // pointer, so the two regions must not meet. The size of that
            // reservation is compiler-version and optimiser dependent, so
            // assert the invariant in the deployed bytecode instead of relying
            // on a generator-side test the integrator never runs. ~6 gas.
            //
            // MF-12: by the letter of Solidity's memory-safety contract this
            // annotation is a lie -- the block writes memory it never
            // allocated through the free-memory pointer. Three properties
            // make it safe HERE, and all three must hold together: this
            // block is terminal (no Solidity executes after it), the pragma
            // is pinned so codegen cannot shift underneath it, and the guard
            // below fails closed if the spill reservation ever reaches the
            // generated layout. Lifting this body into a non-terminal
            // context, or unpinning the pragma, invalidates the annotation.
            //
            // MF-2: this is the only on-chain guard against a recompile that
            // silently moves the spill region, and the failure it catches is
            // permanent (no input can verify). `fail()` is not in scope this
            // early, so write the MemoryLayoutViolated() selector inline
            // rather than reverting bare -- an empty revert here is
            // indistinguishable from every other empty revert, which is
            // exactly the wrong signal for a build fault.
            if gt(mload(0x40), TRANSCRIPT_MPTR) {
                mstore(0x00, shl(224, ERR_MEMORY_LAYOUT_VIOLATED))
                revert(0x00, 0x04)
            }
\end{lstlisting}

\paragraph{Invariant.} Let $\mathit{fmp}$ be the word at \texttt{0x40} on entry
to the block. Execution continues iff
\begin{equation}
  \label{eq:abi:mf2}
  \mathit{fmp} \;\le\; \mptr{TRANSCRIPT\_MPTR} \;=\; \texttt{0x1000},
\end{equation}
otherwise the call reverts with
$\texttt{0xc9888d23} = \revert{MemoryLayoutViolated}$. Since Solidity's
dispatcher initialises \texttt{0x40} to $\texttt{0x80} + (\text{spill
reservation})$ and the only code executed before this point is the calldata
decoder (which allocates no memory: both parameters are \yul{calldata} slices,
so nothing is copied), the head-word guard, and an immutable read,
$\mathit{fmp}$ here \emph{is} the top of solc's reserved prefix. \eqref{eq:abi:mf2} is thus
exactly the separation property the layout needs, stated against the smallest
generated address, and it is the tightest such statement: any write in the block
is at $\ge\texttt{0x1000}\ge\mathit{fmp}$ or in scratch.

\paragraph{Headroom in this build.} The file header (lines 5--7) records
\texttt{0x8c0} on solc 0.8.24 and \texttt{0x8e0} on 0.8.26 and later. The header
does not say whether these are the reservation \emph{size} (entry
$\mathit{fmp} = \texttt{0x80} + \texttt{0x8e0} = \texttt{0x960}$) or the
resulting pointer \emph{value} ($\mathit{fmp} = \texttt{0x8e0}$), and this
specification did not recompile the fixture to settle it. Either reading leaves
$\texttt{0x6a0}$--$\texttt{0x720}$ bytes, i.e.\ 53--57 further 32-byte spill
slots, before the guard begins to fire. The margin is comfortable but is not a
constant the source controls, which is precisely the comment's point.

\paragraph{Failure semantics.} The guard is input-independent: if it fires it
fires for every call, so the contract is bricked rather than unsound. That is the
correct direction. Emitting the typed \revert{MemoryLayoutViolated} instead of a
bare \yul{revert(0, 0)} is a real operational improvement, since a bare revert is
indistinguishable from an out-of-gas or a precompile fault and would point
incident responders at the chain instead of at the build. Note that
\yul{fail(sel)} is defined \emph{after} this guard (line 690), hence the inline
duplication of the two-instruction revert sequence; the same duplication appears
in the constructor's copy of the guard at lines 363--366 and in the head-word
guard at lines 639--640.

\paragraph{Boundary case.} The comparison is strict (\opcode{GT}), so
$\mathit{fmp} = \texttt{0x1000}$ is accepted. This is correct: at that value
Solidity considers $[\texttt{0x1000},\infty)$ unallocated, and the block writes
there without returning control, so no Solidity object can be clobbered.
Tightening it to \opcode{LT} would reject a legal configuration.

\paragraph{Cost.} The comment's ``\textasciitilde 6 gas'' undercounts. On the
accept path the guard executes \opcode{PUSH} \texttt{0x1000} (3), \opcode{PUSH}
\texttt{0x40} (3), \opcode{MLOAD} (3), \opcode{GT} (3), \opcode{ISZERO} (3),
\opcode{PUSH} target (3) and \opcode{JUMPI} (10): about 28 gas, roughly
$4.5\times$ the stated figure, which appears to count only the \opcode{MLOAD}
and \opcode{GT}. The discrepancy is immaterial -- for a function costing
hundreds of thousands of gas the guard is still free -- but the comment should
not be quoted as measured.

\begin{finding}[abi:F5]{MF-2 guard: clean, and correctly scoped}
\textbf{Severity: Observation (no defect).} \texttt{Halo2Verifier.sol}:674--677.
The guard checks the one quantity that matters (the top of solc's reserved
prefix) against the one constant that matters (the bottom of the generated
arena), at the only moment where both are meaningful, and fails closed. Two
scoping details are worth confirming rather than assuming. First, the earlier
\yul{memory-safe} block at lines 635--642 runs before the guard, but it writes
only scratch, so it needs no guard of its own. Second, the guard bounds the
generated layout only from below; nothing on chain bounds it from above, but the
upper end is a fixed generated constant and its only consequence is
memory-expansion gas, not correctness. The constructor carries an equivalent
guard at line 363 against the same \texttt{0x1000} threshold, so a build fault of
this kind is caught at deployment as well as at call time.
\end{finding}

\subsection{\yul{fail(sel)}}
\label{sec:abi:fail}

\begin{lstlisting}[style=solidity,caption={\texttt{Halo2Verifier.sol} lines 687--693: the typed-revert helper (indentation reproduced from the fixture)},label={lst:abi:fail}]
                // Revert with a 4-byte custom-error selector (P4/L-3). Writing at
            // 0x00 is Solidity's legal scratch space and never touches the
            // generated layout, which starts at TRANSCRIPT_MPTR.
            function fail(sel) {
                mstore(0x00, shl(224, sel))
                revert(0x00, 0x04)
            }
\end{lstlisting}

\yul{fail(sel)} writes \yul{sel} left-aligned into the top four bytes of the
scratch word at \texttt{0x00} and reverts returning exactly those four bytes,
which is Solidity's custom-error encoding for a zero-argument error. It never
returns to its caller. Every typed rejection inside the main assembly block
funnels through it except four sites that predate its definition or deliberately
bypass it: the constructor's copy of the MF-2 guard (lines 363--366), the
head-word guard (lines 639--640), the runtime MF-2 guard (lines 675--676), and
\yul{q\_program\_fail} at line 1891, which duplicates the body so the quotient
VM can render into a second, standalone assembly block. Note also that the
constructor's precompile probes (lines 405--407 and following) deliberately keep
bare \yul{revert(0, 0)}: per the comment at lines 361--362 those signal a
chain-capability failure rather than a build fault, so they are outside
\yul{fail}'s typed-error contract by design.

The eight selector constants at lines 320--327 were recomputed independently for
this specification and all eight match
$\mathrm{bytes4}(\mathrm{keccak256}(\text{signature}))$ of the corresponding
\yul{error} declaration at lines 60--83.

\begin{table}[htbp]
\centering
\caption{Typed-error selectors (\texttt{Halo2Verifier.sol} lines 320--327), each verified
against the declaration at lines 60--83.}
\label{tab:abi:selectors}
\begin{tabular}{@{}llll@{}}
\toprule
Yul constant & Declaration & Selector & Verified \\
\midrule
\yul{ERR\_BAD\_CALLDATA\_SHAPE}       & \yul{BadCalldataShape()}       & \texttt{0x1b99e37c} & yes \\
\yul{ERR\_VK\_MISMATCH}               & \yul{VkMismatch()}             & \texttt{0xa447d73e} & yes \\
\yul{ERR\_NON\_CANONICAL\_SCALAR}     & \yul{NonCanonicalScalar()}     & \texttt{0x77530042} & yes \\
\yul{ERR\_BAD\_POINT\_ENCODING}       & \yul{BadPointEncoding()}       & \texttt{0xf27905ec} & yes \\
\yul{ERR\_PRECOMPILE\_FAILED}         & \yul{PrecompileFailed()}       & \texttt{0x84e81692} & yes \\
\yul{ERR\_PROOF\_REJECTED}            & \yul{ProofRejected()}          & \texttt{0xc3b0d8cd} & yes \\
\yul{ERR\_QUOTIENT\_PROGRAM\_INVALID} & \yul{QuotientProgramInvalid()} & \texttt{0x3cc81b89} & yes \\
\yul{ERR\_MEMORY\_LAYOUT\_VIOLATED}   & \yul{MemoryLayoutViolated()}   & \texttt{0xc9888d23} & yes \\
\bottomrule
\end{tabular}
\end{table}

Two soundness notes. The shift \yul{shl(224, sel)} discards any bits of
\yul{sel} above bit 31; all eight constants are genuine four-byte values, so no
selector is truncated, but the helper would silently mis-encode a wider argument.
And because the write lands in $[\texttt{0x00},\texttt{0x04})\subset$ scratch, it
cannot disturb the generated arena at $\ge\texttt{0x1000}$ -- which matters
because \yul{fail} is called from deep inside routines that hold live data in
that arena. The stray indentation of the first comment line (line 687) is a
template-seam artifact with no semantic effect.

\subsection{\yul{scalar\_inv(x)}: inversion in $\Fr$ via \opcode{MODEXP}}
\label{sec:abi:scalarinv}

\begin{lstlisting}[style=solidity,caption={\texttt{Halo2Verifier.sol} lines 695--722: Fermat inversion in $\mathbb{F}_r$},label={lst:abi:scalarinv}]
            // Inverse of a Fr scalar via modexp(x, r-2, r). The verifier
            // calls this only after transcript absorption is complete, so it
            // reuses the dead transcript buffer just below VK_MPTR instead of
            // a fixed post-VK address that can collide with live PCS scratch
            // when the VK payload becomes smaller.
            function scalar_inv(x) -> inv {
                // Zero has no multiplicative inverse in Fr; callers rely on a
                // revert here rather than a bogus modexp result. Check the
                // full canonical range, not just the literal word 0: for any
                // x congruent to 0 mod r (x = r, say) modexp returns 0, which
                // downstream mulmod chains would silently absorb. Every
                // current call site feeds addmod/mulmod output, so this only
                // guards against a future emitter passing a raw scalar.
                if iszero(lt(x, FR_MODULUS)) { fail(ERR_NON_CANONICAL_SCALAR) }
                if iszero(x) { fail(ERR_NON_CANONICAL_SCALAR) }
                let p := 0x3580
                // EIP-198 modexp frame:
                //   [base_len, exp_len, mod_len, base, exponent, modulus]
                mstore(add(p, 0x00), 0x20)        // base len
                mstore(add(p, 0x20), 0x20)        // exp len
                mstore(add(p, 0x40), 0x20)        // mod len
                mstore(add(p, 0x60), x)
                mstore(add(p, 0x80), sub(FR_MODULUS, 2))
                mstore(add(p, 0xa0), FR_MODULUS)
                if iszero(staticcall(MODEXP_GAS, 0x05, p, 0xc0, p, 0x20)) { fail(ERR_PRECOMPILE_FAILED) }
                if iszero(eq(returndatasize(), 0x20)) { fail(ERR_PRECOMPILE_FAILED) }
                inv := mload(p)
            }
\end{lstlisting}

\paragraph{Specification.} For $x \in \Fr$ the routine returns
\begin{equation}
  \yul{scalar\_inv}(x) \;=\; x^{\,r-2} \bmod r \;=\; x^{-1} \in \Fr,
  \qquad r = \texttt{0x73eda753299d7d483339d80809a1d80553bda402fffe5bfeffffffff00000001},
\end{equation}
by Fermat's little theorem, since $r$ is prime. It reverts rather than returning
on any of the following: $x \ge r$; $x = 0$; the \opcode{STATICCALL} to
precompile \texttt{0x05} returning failure; or that call returning a payload
whose length is not \texttt{0x20}.

\begin{algorithm}[htbp]
\caption{\yul{scalar\_inv}(x) (\texttt{Halo2Verifier.sol} lines 700--722)}
\label{alg:abi:scalarinv}
\begin{algorithmic}[1]
\Require a word $x$
\Ensure $x^{-1} \bmod r$, or a revert
\If{$\neg (x < r)$} \State \textbf{fail}(\revert{NonCanonicalScalar}) \EndIf
\If{$x = 0$}       \State \textbf{fail}(\revert{NonCanonicalScalar}) \EndIf
\State $p \gets \texttt{0x3580}$
\State $\mathrm{mem}[p+\texttt{0x00}] \gets \texttt{0x20}$ \Comment{base length}
\State $\mathrm{mem}[p+\texttt{0x20}] \gets \texttt{0x20}$ \Comment{exponent length}
\State $\mathrm{mem}[p+\texttt{0x40}] \gets \texttt{0x20}$ \Comment{modulus length}
\State $\mathrm{mem}[p+\texttt{0x60}] \gets x$
\State $\mathrm{mem}[p+\texttt{0x80}] \gets r-2$
\State $\mathrm{mem}[p+\texttt{0xa0}] \gets r$
\State $ok \gets \opcode{STATICCALL}(\yul{MODEXP\_GAS}, \texttt{0x05}, p, \texttt{0xc0}, p, \texttt{0x20})$
\If{$\neg ok$} \State \textbf{fail}(\revert{PrecompileFailed}) \EndIf
\If{$\opcode{RETURNDATASIZE} \ne \texttt{0x20}$} \State \textbf{fail}(\revert{PrecompileFailed}) \EndIf
\State \Return $\mathrm{mem}[p]$
\end{algorithmic}
\end{algorithm}

\paragraph{Input validation.} Two separate rejections are emitted, both with
\revert{NonCanonicalScalar}. The range check $x < r$ is the stronger one and is
not redundant with the zero check: \opcode{MODEXP} computes $x^{r-2} \bmod r$
for \emph{any} input word, so $x = r$ (or any multiple of $r$ below $2^{256}$)
would return $0$, and a $0$ propagating into the downstream \yul{mulmod} chains
of the PCS interpolation would silently annihilate the result instead of
reverting. Ordering the checks range-first, zero-second is immaterial since both
revert with the same selector. As the comment states, every call site in this
fixture (lines 3673, 3717, 3753, 3782, 3802) passes the output of an
\yul{addmod}/\yul{mulmod} chain already reduced mod $r$, so both checks are
currently defence against future emitter changes rather than against proof data.
That is the right place to put such a check: at the primitive, not at each call
site.

\paragraph{Frame layout and memory reuse.} The EIP-198 frame occupies
$[\texttt{0x3580}, \texttt{0x3640})$, i.e.\ \texttt{0xc0} bytes. The output word
is written back over the frame's first word at $[\texttt{0x3580},
\texttt{0x35a0})$ and read from there.

\begin{table}[htbp]
\centering
\caption{Memory neighbourhood of the \yul{scalar\_inv} frame.}
\label{tab:abi:frame}
\begin{tabularx}{\textwidth}{@{}lllX@{}}
\toprule
Range & Size & Owner & Note \\
\midrule
$[\texttt{0x1000},\texttt{0x3580})$ & \texttt{0x2580} & transcript buffer (\mptr{TRANSCRIPT\_MPTR}) & dead when \yul{scalar\_inv} runs; later reused by the PCS pairing and \yul{ec\_pairing} staging areas and by \mptr{RETURN\_MPTR} \\
$[\texttt{0x3580},\texttt{0x3640})$ & \texttt{0xc0}   & \yul{scalar\_inv} modexp frame & written here \\
$[\texttt{0x3640},\texttt{0x3680})$ & \texttt{0x40}   & tail of the same \texttt{0x100} scratch region & allocated, never written \\
$[\texttt{0x3680},\texttt{0x7900})$ & \texttt{0x4280} & VK payload (\mptr{VK\_MPTR}) & live; never touched \\
\bottomrule
\end{tabularx}
\end{table}

Four facts make the reuse sound, and all four are checkable from the fixture
alone.
\begin{enumerate}
  \item \textbf{Lifetime.} Transcript absorption ends at line 1773 (the final
        \yul{common\_uncompressed\_g1} for $\pi$); the last squeeze is
        \yul{squeeze\_to(buf\_len, X4\_MPTR)} at line 1767. Every
        \yul{scalar\_inv} call site is at line 3673 or later, inside the PCS
        $f$-eval interpolation. The buffer is therefore provably dead before the
        frame is written.
  \item \textbf{High-water mark.} Even ignoring lifetime, the transcript never
        reaches \texttt{0x3580} in this render. \yul{squeeze\_to} (line 801)
        reseeds the buffer to $\mptr{TRANSCRIPT\_MPTR}+\texttt{0x20}$ after every
        digest, so only the longest inter-squeeze segment matters. There are
        nine such segments, delimited by the ten \yul{squeeze\_to} calls at
        lines 1568, 1586, 1587, 1646, 1663, 1686, 1720, 1721, 1733 and 1767; all
        nine were enumerated, and every segment other than the two named below
        ends at or under \texttt{0x1520}. The first
        segment absorbs the VK digest (\texttt{0x20}), the identity
        \yul{committed\_pi} as 128 zero bytes (\texttt{0x80}), the instance-count
        word (\texttt{0x20}), 19 instance words (\texttt{0x260}) and
        \texttt{0x780} bytes of phase-1 advice commitments, ending at
        \texttt{0x1aa0}; the evaluation segment (line 1699) absorbs
        \texttt{0xcc0} bytes over the \texttt{0x20} reseed, ending at
        \texttt{0x1ce0}. The maximum is \texttt{0x1ce0}, some \texttt{0x18a0}
        bytes below the frame.
  \item \textbf{No other occupant.} Enumerating the generated pointer constants
        at lines 119--263, the \emph{only} ones below \mptr{VK\_MPTR} are
        \mptr{TRANSCRIPT\_MPTR}, \mptr{RETURN\_MPTR} and
        \mptr{QUOTIENT\_RETURN\_MPTR}, all three equal to \texttt{0x1000};
        every other region starts at \texttt{0x3680} or above. The only other
        traffic in $[\texttt{0x1000},\texttt{0x3680})$ is the
        \yul{ec\_pairing} staging area at \texttt{0x1240} (line 969, extent
        \texttt{0x300}) and the PCS pairing staging area at
        \texttt{0x1000}--\texttt{0x1180} (lines 4012--4033), both of which run
        strictly after the last \yul{scalar\_inv} call at line 3802 and neither
        of which reaches \texttt{0x3580}.
  \item \textbf{Disjointness from the VK.} The frame ends at \texttt{0x3640},
        strictly below $\mptr{VK\_MPTR}=\texttt{0x3680}$, so the 17\,024 bytes
        copied by \yul{extcodecopy(vk, VK\_MPTR, 0x01, EXPECTED\_VK\_PAYLOAD\_LENGTH)}
        at line 1393 -- which are live for the whole call -- are untouched.
\end{enumerate}
The \texttt{0x40} gap at $[\texttt{0x3640},\texttt{0x3680})$ is not slack and is
not unowned: the generator registers a \texttt{0x100}-byte region
(\texttt{MODEXP\_SCRATCH\_BYTES} $= 8$ words,
\texttt{src/lowering/layout/mod.rs}:92--93) at $\mptr{VK\_MPTR}-\texttt{0x100}$
while the frame itself is only \texttt{0xc0} bytes
(\texttt{MODEXP\_FRAME\_BYTES} $= 6$ words, ibid.\ lines 86--87), so the tail
stays owned by \texttt{scalar\_inv\_scratch} and cannot be claimed by a later
allocation (\texttt{src/lowering/layout/memory.rs}:786--795). The comment's rationale for
anchoring the frame \emph{below} the VK rather than after it is also correct and
worth preserving: a post-VK fixed address would move relative to live PCS scratch
whenever the VK payload shrinks, whereas $\mptr{VK\_MPTR}-\texttt{0x100}$ tracks
the VK automatically.

\paragraph{Gas bound.} \yul{MODEXP\_GAS} is the fixed literal $4080$ (line 293),
forwarded as the \opcode{STATICCALL} gas limit rather than \opcode{GAS}. The
frame is the smallest legal one
($\text{base\_len}=\text{exp\_len}=\text{mod\_len}=32$, exponent $r-2$ of bit
length $255$). The fixture's own comment (line 384) names the two live schedules
that price it as ``EIP-2565: 1360, EIP-7883: 4080''. Those figures should be read
as conservative round-ups rather than as the schedules' outputs: both EIPs price
a $32$-byte exponent by the \emph{adjusted} exponent length
$\mathrm{bitlen}(r-2)-1 = 254$, not by $255$, and the multiplication complexity
for $32$-byte operands is $\lceil 32/8\rceil^2 = 16$, giving
\begin{equation}
  \label{eq:abi:modexpgas}
  \text{EIP-2565: } \max\!\left(200, \left\lfloor \tfrac{16\cdot 254}{3}
  \right\rfloor\right) = 1354,
  \qquad
  \text{EIP-7883: } \max(500,\, 16\cdot 254) = 4064,
\end{equation}
against the quoted $1360 = 16\cdot 255/3$ and $4080 = 16\cdot 255$. The error is
in the safe direction: the pinned constant exceeds the price under both
schedules, so the call succeeds under either, with $16$ gas of slack under the
more expensive one. This specification did \emph{not} verify either schedule
against a primary source, and the arithmetic above is therefore evidence that
the constant is conservative, not a proof that it is sufficient on any given
chain; the actual guarantee is the deployment-time known-answer probe at lines
399--407, which exercises this exact frame at this exact gas bound and fails the
constructor if the chain prices it higher. A capped forward is the right choice
here regardless: it makes the cost of a \yul{scalar\_inv} independent of the
caller's gas and prevents a mispriced or repriced precompile from consuming the
whole remaining budget.

\paragraph{Return-data handling.} Precompile \texttt{0x05} returns exactly
$\text{mod\_len}=32$ bytes on success. The explicit
\yul{eq(returndatasize(), 0x20)} check at line 720 is necessary and not
cosmetic: \opcode{STATICCALL} copies only $\min(\texttt{0x20},
\opcode{RETURNDATASIZE})$ bytes into $p$ and does \emph{not} zero-pad the
remainder of the output area, so a short return leaves some or all of the frame's
stale first word (the literal \texttt{0x20} base-length) in place. In the
degenerate case $\opcode{RETURNDATASIZE}=0$ -- what a chain without the
precompile, or a stub that succeeds and returns nothing, would produce --
\yul{inv} would be read as exactly $32$: a plausible-looking but wrong scalar,
silently accepted by every downstream \yul{mulmod}. Checking the size rather
than the value closes that. The same pattern is used by
\yul{batch\_invert} at lines 873 and 924.

\begin{finding}[abi:F6]{Fixed \opcode{MODEXP} gas cap conflates caller under-gassing with a chain fault}
\textbf{Severity: Informational.} \texttt{Halo2Verifier.sol}:719, 293.
The \opcode{STATICCALL} forwards a fixed $4080$ gas. Under EIP-150 the callee
receives $\min\!\left(4080, \left\lfloor \tfrac{63 G}{64} \right\rfloor\right)$
where $G$ is the gas remaining at the call, so a
caller whose $G$ dips below $\lceil 4080 \cdot 64/63 \rceil = 4145$ gas at this point
starves the precompile, which then out-of-gases and returns $0$, and the verifier
reverts \revert{PrecompileFailed}. The revert selector then reports a chain or
precompile fault when the actual cause was the caller's gas limit. In practice
the surrounding verification consumes far more than that before reaching this
line, so the window is narrow, but a responder triaging a \revert{PrecompileFailed}
report should confirm the caller's gas limit before concluding the node or the
gas schedule is at fault. The genuine repricing case -- a chain that prices the
frame above $4080$ -- is caught at deployment by the known-answer probe at lines
399--407, which uses the same frame shape, exponent width and gas bound
(computing $2^{r-2} = 2^{-1}$) and checks
$\mathrm{mulmod}(\text{result}, 2, r) = 1$ at line 407, after a
\opcode{RETURNDATASIZE} check at line 406. Unlike the runtime helper the probe
reverts bare, by design.
\end{finding}

\begin{finding}[abi:F7]{Generated memory bases are rendered as bare literals, hiding the relations that make them correct}
\textbf{Severity: Low.} \texttt{Halo2Verifier.sol}:710; also 363, 969, 2172,
4012--4033.
\yul{let p := 0x3580} is correct in this fixture -- $\texttt{0x3580} =
\mptr{VK\_MPTR} - \texttt{0x100}$, a relation the generator asserts explicitly
(\texttt{src/lowering/layout/memory.rs}:1236--1243, ``scalar inversion scratch
drifted'') -- but nothing in the deployed source records that relation. A reader
of \texttt{Halo2Verifier.sol} alone sees an unexplained constant, and a hand-edit
of \mptr{VK\_MPTR} that missed this literal would place the \texttt{0xc0}-byte
frame either inside the VK payload (corrupting it, since the frame is written
after the \opcode{EXTCODECOPY} at line 1393) or lower in the transcript arena
(harmless).

The literal at line 710 is not exceptional, which is the wider point: the main
assembly block uses $124$ distinct bare hex literals as memory operands, among
them \texttt{0x1240} (line 969, \yul{ec\_pairing} staging),
\texttt{0xb8e0} (line 2172, quotient-VM stack base) and
\texttt{0x1000}/\texttt{0x1080}/\texttt{0x1100} (lines 4012--4033, PCS pairing
staging), alongside absolute VK-payload source addresses such as
\texttt{0x6280} in the final-MSM assembly at lines 3860 onward. The constructor's
MF-2 guard compares against the literal \texttt{0x1000} (line 363) rather than
\mptr{TRANSCRIPT\_MPTR}, which is the same issue applied to a \emph{guard
threshold} and so is arguably the more concerning instance: a change to
\mptr{TRANSCRIPT\_MPTR} would silently desynchronise the deployment-time guard
from the runtime one at line 674, which does use the named constant.
Rendering these as named constants (\yul{SCALAR\_INV\_SCRATCH\_MPTR},
\yul{EC\_PAIRING\_SCRATCH\_MPTR}, and so on) would carry the invariants into the
artifact. This is a maintainability and review-surface issue, not an exploitable
defect in the deployed fixture: every literal was checked against the layout in
Section~\ref{sec:abi:scalarinv} and Table~\ref{tab:abi:frame} and all are
consistent.
\end{finding}

\subsection{Summary of guards in this range}
\label{sec:abi:summary}

\begin{table}[htbp]
\centering
\caption{Every revert raised by code in \texttt{Halo2Verifier.sol} lines 588--722.
Untyped reverts raised \emph{before} this range, by solc's ABI decoder, are not
listed; see Finding~abi:F8.}
\label{tab:abi:guards}
\begin{tabularx}{\textwidth}{@{}llXl@{}}
\toprule
Line & Condition to revert & Meaning & Error \\
\midrule
636--640 & $h_0 \ne \texttt{0x40}$ or $h_1 \ne \texttt{0x1ec0}$ & ABI head words do not point at the generated tails & \revert{BadCalldataShape} \\
674--676 & $\mathrm{mem}[\texttt{0x40}] > \texttt{0x1000}$ & solc's spill reservation reaches the generated arena (build fault) & \revert{MemoryLayoutViolated} \\
708 & $\neg(x < r)$ in \yul{scalar\_inv} & non-canonical scalar handed to inversion & \revert{NonCanonicalScalar} \\
709 & $x = 0$ in \yul{scalar\_inv} & zero has no inverse in $\Fr$ & \revert{NonCanonicalScalar} \\
719 & \opcode{STATICCALL} to \texttt{0x05} returned failure & modexp precompile missing, reverted, or out of gas & \revert{PrecompileFailed} \\
720 & $\opcode{RETURNDATASIZE} \ne \texttt{0x20}$ & modexp returned a short or oversized payload & \revert{PrecompileFailed} \\
\bottomrule
\end{tabularx}
\end{table}

Checks a reader might expect in this range and which are deliberately
\emph{absent} here, with the line that performs them instead: \yul{proof.length}
equality (line 1419), \yul{instances.length} equality (line 1420),
\opcode{CALLDATASIZE} equality (line 1426), the VK address/codehash re-check and
\opcode{EXTCODECOPY} (lines 1380--1406), and the parser's terminal-cursor check
(line 1785). All are covered in other sections.

Two further classes of rejection are outside the range in the other direction,
i.e.\ they fire \emph{before} line 636 and are therefore invisible in
Table~\ref{tab:abi:guards}: solc's generated ABI decoder (bare
\yul{revert(0x00, 0x00)} on \opcode{CALLDATASIZE} below \texttt{0x44} or on
either dynamic tail out of bounds) and the dispatcher's own selector and
\opcode{CALLVALUE} checks, likewise untyped. Finding~abi:F8 records the
documentation gap this creates.


%% ==== 05-transcript.tex =================================================
\section{The streaming Keccak256 transcript}
\label{sec:tsh:main}

This section specifies the Fiat--Shamir layer of the generated verifier: the four
Yul helpers \yul{transcript\_init}, \yul{common\_word},
\yul{common\_uncompressed\_g1} and \yul{squeeze\_to}, which occupy lines
724--809 of \texttt{Halo2Verifier.sol}. The assigned range 700--836 additionally
contains \yul{scalar\_inv} (lines 700--722), the \texttt{EC primitives} banner
comment (lines 811--816), and the doc comment plus signature line of
\yul{batch\_invert} (lines 818--836); those belong to the field-arithmetic
section and are covered there. Because the audit contract requires every
precompile call inside the assigned range to be catalogued, and because
\yul{scalar\_inv} places its scratch frame inside the transcript arena,
Section~\ref{sec:tsh:boundary} records its numeric facts and discharges the
adjacency argument; its arithmetic is not respecified here.

Every helper in this block is a pure Yul function. Yul function bodies do not
capture variables of the enclosing block, so the local
\yul{let r := FR\_MODULUS} at line 805 inside \yul{squeeze\_to} is a fresh
binding and does not shadow the outer \yul{r} declared at line 1355; the two are
equal by construction.

\subsection{State, constants and the buffer invariant}
\label{sec:tsh:state}

The transcript is not a sponge object. It is a \emph{byte buffer} in EVM memory
plus a single \texttt{uint256} cursor carried on the Yul stack. Write

\[
  T \;=\; \mptr{TRANSCRIPT\_MPTR} \;=\; \texttt{0x1000},
  \qquad
  L \;=\; \yul{buf\_len} \;\in\; [\,T,\;\infty),
\]
\[
  B \;=\; \mathrm{memory}[\,T,\;L\,) \;\in\; \{0,1\}^{8(L-T)} .
\]

$B$ is the \emph{live} transcript buffer; every byte at address $\ge L$ is stale
and is never read. All constants used by this block are fixed at generation time
and are listed in Table~\ref{tab:tsh:constants}.

\begin{table}[h]
\centering
\begin{tabularx}{\textwidth}{@{}l l X@{}}
\toprule
Constant & Value & Role \\
\midrule
\mptr{TRANSCRIPT\_MPTR} & \texttt{0x1000} & Base of the transcript buffer (line 119). \\
\yul{FR\_MODULUS} & \texttt{0x73eda753299d7d48\allowbreak 3339d80809a1d805\allowbreak 53bda402fffe5bfe\allowbreak ffffffff00000001} & $r$, the BLS12-381 scalar modulus (line 331). \\
\yul{BLS\_P\_HI} & \texttt{0x0000000000000000\allowbreak 0000000000000000\allowbreak 1a0111ea397fe69a\allowbreak 4b1ba7b6434bacd7} & $\lfloor (p-1)/2^{256}\rfloor = \lfloor p/2^{256}\rfloor$ (line 336). \\
\yul{BLS\_P\_MINUS\_ONE\_LO} & \texttt{0x64774b84f38512bf\allowbreak 6730d2a0f6b0f624\allowbreak 1eabfffeb153ffff\allowbreak b9feffffffffaaaa} & $(p-1) \bmod 2^{256}$ (line 337). \\
\yul{ERR\_BAD\_POINT\_ENCODING} & \texttt{0xf27905ec} & \revert{BadPointEncoding} selector (line 323). \\
\bottomrule
\end{tabularx}
\caption{Constants consumed by the transcript block. Line numbers refer to
\texttt{Halo2Verifier.sol}.}
\label{tab:tsh:constants}
\end{table}

Note the asymmetric naming: \yul{BLS\_P\_HI} is the high word of \emph{both} $p$
and $p-1$ (they differ only in the low limb, since $p \bmod 2^{256} \neq 0$),
while \yul{BLS\_P\_MINUS\_ONE\_LO} is the low word of $p-1$. A reader must not
read \yul{BLS\_P\_HI} as "the high word of $p$ only"; the pair
$(\yul{BLS\_P\_HI}, \yul{BLS\_P\_MINUS\_ONE\_LO})$ is exactly the 512-bit
big-endian decomposition of $p-1$. This matters for the exactness argument in
Section~\ref{sec:tsh:g1}.

\subsection{The Fiat--Shamir construction}
\label{sec:tsh:fs}

Let $a_1, a_2, \dots, a_m$ be the byte strings absorbed between successive
squeezes, i.e. $a_i$ is the concatenation of everything appended after the
$(i-1)$-th squeeze and before the $i$-th. The block implements the chained-hash
transcript

\begin{align}
  h_1     &= \mathrm{Keccak256}\bigl(a_1\bigr), \label{eq:tsh:h1}\\
  h_i     &= \mathrm{Keccak256}\bigl(h_{i-1} \,\|\, a_i\bigr), && i \ge 2, \label{eq:tsh:hi}\\
  c_i     &= \mathrm{int}_{\mathrm{BE}}(h_i) \bmod r \;\in\; \Fr. \label{eq:tsh:ci}
\end{align}

The 32-byte digest $h_{i-1}$ -- not the reduced challenge $c_{i-1}$ -- is what is
carried forward into the next hash. That distinction is load-bearing twice: it
preserves the full 256 bits of chain entropy, and it is what the native
implementation does (\texttt{midfall/proofs/src/transcript/implementors.rs},
\yul{impl TranscriptHash for Keccak256}: \yul{squeeze} finalizes a clone,
then re-initialises the state with \yul{new\_hasher.update(\&out)} where
\yul{out} is the raw digest). Reseeding with $c_{i-1}$ instead would produce a
verifier that rejects every honestly generated proof.

Concretely $m = 10$ for this fixture, and
$(c_1,\dots,c_{10}) = (\theta,\beta,\gamma,\tau,y,x,x_1,x_2,x_3,x_4)$; see
Section~\ref{sec:tsh:schedule}. Note that $x_3$ is additionally truncated to its
low 128 bits \emph{after} the squeeze (line 1743); that masking happens at the
call site and is specified in the PCS section, not here --
\yul{squeeze\_to} itself always stores the full $c_i$.

\subsection{\texttt{transcript\_init}}
\label{sec:tsh:init}

\begin{lstlisting}[style=solidity,caption={\texttt{Halo2Verifier.sol} lines 733--743: \yul{transcript\_init} and \yul{common\_word}},label={lst:tsh:init}]
            function transcript_init() -> buf_len {
                // Empty transcript buffer starts exactly at TRANSCRIPT_MPTR.
                buf_len := TRANSCRIPT_MPTR
            }

            // Append one 32-byte big-endian field/transcript word at the
            // current end of the transcript buffer.
            function common_word(buf_len, word) -> ret {
                mstore(buf_len, word)
                ret := add(buf_len, 32)
            }
\end{lstlisting}

\yul{transcript\_init} sets $L := T$, i.e. $B := \varepsilon$ (the empty
string). It performs no memory write, zeroes nothing, and cannot fail. There is
no hash-state initialisation because there is no hash state: the first
\yul{keccak256} is taken over whatever has been appended by then. This mirrors
\yul{Keccak256::new()} on the native side, which likewise absorbs no domain
separator (contrast the Blake2b transcript in the same Rust file, which is keyed
with \texttt{"Domain separator for transcript"} and tags every absorb and squeeze
with a one-byte prefix). See Finding~TSH-1.

Because nothing is zeroed, the emptiness of $B$ is a \emph{cursor} property, not
a memory property: bytes left over at addresses $\ge T$ from a previous phase
would be invisible as long as the cursor discipline holds. The transcript is
\emph{not} the first thing in the call frame to touch the generated layout --
the VK \opcode{EXTCODECOPY} at line 1393 fills \mptr{VK\_MPTR} $=$
\texttt{0x3680}, and \yul{validate\_public\_accumulator} (called at line 1447,
before \yul{transcript\_init} at line 1474) writes \mptr{ACC\_LHS\_MPTR} $=$
\texttt{0x7b40}, \mptr{ACC\_RHS\_MPTR} $=$ \texttt{0x7bc0} and its MSM scratch
at \texttt{0xb140}. What matters is the weaker and verifiable fact that
\emph{nothing writes below \texttt{0x3680} before line 1474}: every pre-transcript
write in this contract lands at $\ge$ \texttt{0x3680}, so the window
$[T, L_{\max})$ is untouched when \yul{transcript\_init} runs and $B$ genuinely
starts empty on every call. \mptr{RETURN\_MPTR} and
\mptr{QUOTIENT\_RETURN\_MPTR} do share the address \texttt{0x1000}, but they
belong to strictly later memory phases (documented at lines 106--118; the
\yul{mcopy(0x1000, G1\_BASE\_MPTR, 0x80)} at line 4012 is the first reuse), so
they cannot precede the transcript.

\subsection{\texttt{common\_word}}
\label{sec:tsh:word}

\yul{common\_word} appends one 32-byte big-endian word:

\[
  B \;\leftarrow\; B \,\|\, \mathrm{BE}_{32}(\yul{word}),
  \qquad
  L \;\leftarrow\; L + 32 .
\]

It is a two-instruction function: one \opcode{MSTORE} and one \opcode{ADD}. It
performs \emph{no} validation whatsoever -- in particular it does not check
$\yul{word} < r$. Canonicality is a caller obligation, and every one of the
\emph{five} \yul{common\_word} call sites in the contract (lines 1481, 1508,
1521, 1712, 1760 -- there are no others) discharges it.
Table~\ref{tab:tsh:wordcallers} enumerates them exhaustively.

\begin{table}[h]
\centering
\begin{tabularx}{\textwidth}{@{}l l X@{}}
\toprule
Call site & Word & Canonicality argument \\
\midrule
1481 & \yul{mload(VK\_DIGEST\_MPTR)} & Not range-checked at the call site. The word is the pinned VK digest (\texttt{0x56c0824f...}, \texttt{Halo2VerifyingKey.sol}:68), which is $< r$; the VK runtime is pinned by length and codehash (\yul{EXPECTED\_VK\_CODEHASH}, line 95), so the value cannot vary per call. \\
1508 & literal \texttt{19} & Generated constant; the calldata instance-array length is separately checked $= 19$ at line 1420 and the VK header at line 1400. \\
1521 & public instance & \yul{success := and(success, lt(inst\_be, r))} at line 1518, enforced at line 1523 with \revert{NonCanonicalScalar}. \\
1712, 1760 & proof evaluation / \texttt{q\_eval} & \yul{if iszero(lt(eval, r)) \{ fail(ERR\_NON\_CANONICAL\_SCALAR) \}} at lines 1706 and 1759 respectively, i.e.\ immediate revert before the absorb. \\
\bottomrule
\end{tabularx}
\caption{All five \yul{common\_word} call sites and the range check that guards
each. Call-site lines are outside this section's range and are specified in the
transcript-schedule section.}
\label{tab:tsh:wordcallers}
\end{table}

\paragraph{One absorb bypasses both helpers.} The buffer is not written
exclusively by \yul{common\_word} and \yul{common\_uncompressed\_g1}. The
\yul{committed\_pi} identity is appended by four raw \yul{mstore(.., 0)} at
lines 1496--1499 with the cursor advanced by hand at line 1500
(\yul{buf\_len := add(buf\_len, 0x80)}). That is safe because the value is a
compile-time constant of exactly 128 zero bytes -- there is nothing to
range-check -- but a reader auditing "what can enter the transcript" must not
stop at the two helpers.

The 32-byte big-endian form is exactly the native
\yul{Hashable<Keccak256> for midnight\_curves::Fq::to\_input}, which takes
\yul{to\_repr()} (canonical little-endian) and reverses it. Since $T$ is
32-byte aligned and every absorb is 32 or 128 bytes, $L$ is always 32-byte
aligned and the \opcode{MSTORE} never straddles a previously written word
boundary.

\subsection{\texttt{common\_uncompressed\_g1}}
\label{sec:tsh:g1}

\subsubsection{Wire format}

A $\GO$ point is transmitted to the transcript as four consecutive calldata
words (128 bytes), the EIP-2537 padded uncompressed encoding:

\[
  \underbrace{\texttt{x\_hi}}_{32}\;\|\;
  \underbrace{\texttt{x\_lo}}_{32}\;\|\;
  \underbrace{\texttt{y\_hi}}_{32}\;\|\;
  \underbrace{\texttt{y\_lo}}_{32},
\]
where each coordinate occupies a 64-byte pair consisting of 16 zero pad bytes
followed by the 48 big-endian bytes of an $\Fp$ element:
\[
  X \;=\; (\texttt{x\_hi} \bmod 2^{128})\cdot 2^{256} + \texttt{x\_lo},
  \qquad
  Y \;=\; (\texttt{y\_hi} \bmod 2^{128})\cdot 2^{256} + \texttt{y\_lo}.
\]
The identity is encoded as 128 zero bytes (the EIP-2537 $(0,0)$ convention, not
the BLS-spec \texttt{0x40}-leading-byte form). This matches the patched
\yul{Hashable<Keccak256> for midnight\_curves::G1Projective::to\_input} in
midnight-proofs, whose doc comment describes the same 128-byte layout; the
48-byte compressed ZCash encoding is still what travels on the \emph{wire}
(\yul{to\_bytes}), but the Solidity-facing proof shim repacks it, as documented
at line 621.

\subsubsection{Code}

\begin{lstlisting}[style=solidity,caption={\texttt{Halo2Verifier.sol} lines 775--796: \yul{common\_uncompressed\_g1} (function body, verbatim)},label={lst:tsh:g1}]
            function common_uncompressed_g1(buf_len, cptr) -> ret {
                let x_hi_word := calldataload(cptr)
                let x_lo := calldataload(add(cptr, 0x20))
                let y_hi_word := calldataload(add(cptr, 0x40))
                let y_lo := calldataload(add(cptr, 0x60))
                if shr(128, x_hi_word) { fail(ERR_BAD_POINT_ENCODING) }
                if shr(128, y_hi_word) { fail(ERR_BAD_POINT_ENCODING) }

                let x_hi := and(x_hi_word, 0xffffffffffffffffffffffffffffffff)
                let y_hi := and(y_hi_word, 0xffffffffffffffffffffffffffffffff)
                if iszero(or(lt(x_hi, BLS_P_HI), and(eq(x_hi, BLS_P_HI), iszero(gt(x_lo, BLS_P_MINUS_ONE_LO))))) {
                    fail(ERR_BAD_POINT_ENCODING)
                }
                if iszero(or(lt(y_hi, BLS_P_HI), and(eq(y_hi, BLS_P_HI), iszero(gt(y_lo, BLS_P_MINUS_ONE_LO))))) {
                    fail(ERR_BAD_POINT_ENCODING)
                }

                // Memcpy the 4 calldata words (128 bytes) verbatim
                // into the keccak buffer.
                calldatacopy(buf_len, cptr, 0x80)
                ret := add(buf_len, 0x80)
            }
\end{lstlisting}

\subsubsection{Guards, exhaustively}

\begin{table}[h]
\centering
\begin{tabularx}{\textwidth}{@{}l l X l@{}}
\toprule
\# & Line & Reject condition & Revert \\
\midrule
G1 & 780 & $\lfloor \texttt{x\_hi\_word} / 2^{128} \rfloor \neq 0$, i.e.\ any of the top 16 bytes of the $X$ high word is non-zero & \revert{BadPointEncoding} \\
G2 & 781 & $\lfloor \texttt{y\_hi\_word} / 2^{128} \rfloor \neq 0$ & \revert{BadPointEncoding} \\
G3 & 785--787 & $\neg\bigl(\,\texttt{x\_hi} < \yul{BLS\_P\_HI} \;\vee\; (\texttt{x\_hi} = \yul{BLS\_P\_HI} \wedge \texttt{x\_lo} \le \yul{BLS\_P\_MINUS\_ONE\_LO})\,\bigr)$ & \revert{BadPointEncoding} \\
G4 & 788--790 & same predicate on $(\texttt{y\_hi}, \texttt{y\_lo})$ & \revert{BadPointEncoding} \\
\bottomrule
\end{tabularx}
\caption{Every guard inside \yul{common\_uncompressed\_g1}. All four share one
selector, so the four failure modes are externally indistinguishable.}
\label{tab:tsh:g1guards}
\end{table}

\paragraph{G1/G2 -- the 16-zero-byte padding check.} \yul{shr(128, w)} is
non-zero iff any bit of the top 128 bits of $w$ is set, i.e.\ iff any of the top
16 pad bytes is non-zero. This is the full EIP-2537 padding requirement, not a
weaker top-byte test. Only the two \texttt{\_hi} words are checked; the
\texttt{\_lo} words carry 32 significant bytes each and have no pad.

\paragraph{G3/G4 -- exact $\Fp$ range.} After G1/G2 the masks at lines 783--784
are redundant (\texttt{x\_hi\_word} already has zero high bits), but they make
the comparison independent of guard order. In Yul, \yul{iszero(gt(a,b))} is
$a \le b$, and both \yul{eq} and \yul{iszero} return canonical booleans
$\{0,1\}$, so the bitwise \yul{and}/\yul{or} coincide with logical
conjunction/disjunction here -- this is \emph{not} an instance of the
bitwise-\yul{and} trap that the repository lints for. The accept predicate is
therefore the lexicographic 512-bit comparison
\[
  (\texttt{x\_hi}, \texttt{x\_lo}) \;\le_{\mathrm{lex}}\;
  \bigl(\lfloor (p-1)/2^{256}\rfloor,\; (p-1)\bmod 2^{256}\bigr)
  \;\Longleftrightarrow\;
  X \le p-1 \;\Longleftrightarrow\; X < p,
\]
and identically for $Y$. The equivalence is exact -- no off-by-one -- precisely
because $\lfloor (p-1)/2^{256}\rfloor = \lfloor p/2^{256}\rfloor$, so
\yul{BLS\_P\_HI} may be reused as the high word of both bounds. Rejecting $X = p$
matters: $p$ and $0$ are the same field element but distinct 128-byte strings,
and hashing both would give one point two transcripts.

\paragraph{Absorb.} On success the function copies the four calldata words
verbatim:
\[
  B \;\leftarrow\; B \,\|\, \mathrm{calldata}[\texttt{cptr},\, \texttt{cptr}+128),
  \qquad L \leftarrow L + 128 .
\]
The bytes hashed are byte-for-byte the bytes validated -- there is no
re-encoding step between check and absorb that could drift. This is why the
generator rejects rather than normalises the pad bytes: masking them to zero
before hashing (which the native doc comment loosely describes as "masking")
would let two distinct calldata encodings map to one transcript, as the fixture
comment at lines 761--764 states. Rejecting is strictly stronger than masking,
and the fixture takes the stronger option.

\subsubsection{What is deliberately \emph{not} checked}

\yul{common\_uncompressed\_g1} performs \textbf{no on-curve check} and
\textbf{no $\GO$ subgroup check}. It also does not reject the identity
$(0,0)$, does not reject $Y = 0$, and does not reject a point whose $X$ is a
non-residue. Two further absent checks are local to the helper and must be
named: it never compares \yul{cptr} against \yul{calldatasize()} or against the
proof-payload extent, and it never compares \yul{buf\_len} against any upper
bound before the \yul{calldatacopy}. Both are caller obligations discharged
whole-program (Section~\ref{sec:tsh:assess} and Finding~TSH-5); neither is
checkable from the helper alone. The generated comment states the justification
for the missing curve check (lines 766--774, reproduced verbatim):

\begin{lstlisting}[style=solidity,caption={\texttt{Halo2Verifier.sol} lines 766--774: the deferral justification, verbatim},label={lst:tsh:defer}]
            // This helper does not run an independent curve/subgroup
            // check. Instead, ProtocolPlan::validate rejects generated
            // plans where an absorbed proof commitment would not later be
            // consumed by an EIP-2537 G1MSM or pairing path, and those
            // precompiles perform the curve/subgroup validation.
            //
            // The point's uncompressed form remains in calldata; the
            // call site is responsible for `calldatacopy`-ing it into
            // memory afterwards if it needs the on-curve coordinates.
\end{lstlisting}

The cost model is the point: an in-EVM subgroup check for BLS12-381 costs a
scalar multiplication by the cofactor (or an endomorphism-based test) in
384-bit arithmetic, which is unimplementable at reasonable gas without the
precompiles; EIP-2537's \texttt{G1MSM} and \texttt{PAIRING} are specified to
reject inputs that are off-curve or outside the prime-order subgroup, so a
second check would be pure duplication for every commitment that reaches them.
See Finding~TSH-3 for the residual obligation this creates.

\subsection{\texttt{squeeze\_to}}
\label{sec:tsh:squeeze}

\begin{lstlisting}[style=solidity,caption={\texttt{Halo2Verifier.sol} lines 798--809: \yul{squeeze\_to}},label={lst:tsh:squeeze}]
            // One Keccak finalization + reseed. Returns the new buffer
            // cursor (= TRANSCRIPT_MPTR + 32) and stores the squeezed Fq at
            // `mptr`.
            function squeeze_to(buf_len, mptr) -> ret {
                let h0 := keccak256(TRANSCRIPT_MPTR, sub(buf_len, TRANSCRIPT_MPTR))
                // Reseed: write the 32-byte digest at start of buffer.
                mstore(TRANSCRIPT_MPTR, h0)
                let r := FR_MODULUS
                // Sample Fq as uint256(keccak_digest_be) mod r.
                mstore(mptr, mod(h0, r))
                ret := add(TRANSCRIPT_MPTR, 32)
            }
\end{lstlisting}

In three steps:
\begin{enumerate}
  \item $h \;=\; \mathrm{Keccak256}\bigl(\mathrm{memory}[T, L)\bigr)$ --
        the hash covers exactly the live buffer, length $L - T$.
  \item \emph{Reseed}: $\mathrm{memory}[T, T+32) \leftarrow h$, and the returned
        cursor is $T + 32$. The digest becomes the new buffer, i.e.\
        $B \leftarrow h$; every byte previously at $[T+32, L)$ is now stale and
        will be overwritten by the next absorbs (or, if the next segment is
        shorter, simply left outside the next hash range -- the hash always runs
        from $T$ to the \emph{current} cursor, so stale tails can never leak into
        a later digest).
  \item \emph{Sample}: $\mathrm{memory}[\yul{mptr}] \leftarrow h \bmod r$.
        The stored challenge is reduced; the chained state $h$ is not.
\end{enumerate}

\begin{algorithm}[h]
\caption{\yul{squeeze\_to}$(L, \yul{mptr})$}
\label{alg:tsh:squeeze}
\begin{algorithmic}[1]
\Require $L \ge T$, $L - T$ a multiple of 32
\State $h \gets \mathrm{Keccak256}(\mathrm{mem}[T, L))$
\State $\mathrm{mem}[T] \gets h$ \Comment{reseed: digest is the new buffer prefix}
\State $\mathrm{mem}[\yul{mptr}] \gets h \bmod r$ \Comment{challenge $c \in \Fr$}
\State \Return $T + 32$
\end{algorithmic}
\end{algorithm}

The function cannot revert. Two degenerate inputs are worth naming and
dismissing: $L = T$ would hash the empty string (Keccak of $\varepsilon$), and
$L < T$ would make \yul{sub} underflow to a near-$2^{256}$ length and
out-of-gas. Neither is reachable -- $L$ is initialised to $T$ by
\yul{transcript\_init} and only ever increases, and the first squeeze at line
1568 is preceded by 2720 bytes of absorbs (Section~\ref{sec:tsh:schedule}), so
every hashed length in this contract is $\ge 32$.

\paragraph{Consecutive squeezes.} The schedule squeezes twice in a row at lines
1586--1587 ($\beta$, then $\gamma$) and at lines 1720--1721 ($x_1$, then $x_2$).
The second call then hashes exactly the 32 reseeded bytes, so
\[
  \gamma = \mathrm{Keccak256}(h_\beta) \bmod r,
  \qquad
  x_2 = \mathrm{Keccak256}(h_{x_1}) \bmod r .
\]
$\gamma$ is a deterministic function of $\beta$'s digest with nothing new
absorbed in between. This is not a defect: it is exactly what the native
\yul{squeeze} does when called twice, and in the random-oracle model
$(c_i, c_{i+1})$ from a chained oracle are indistinguishable from independent.
It does mean an implementer must not "optimise" the second squeeze away or fold
two challenges out of one digest.

\subsection{The absorb/squeeze schedule and the buffer extent}
\label{sec:tsh:schedule}

The call sites are outside this section's range (they are specified in the
transcript/proof-parser section, lines 1474--1785), but the schedule is needed
here to bound the buffer and to argue memory disjointness.
Table~\ref{tab:tsh:schedule} gives it exactly, with all loop bounds resolved to
item counts.

\begin{table}[h]
\centering
\begin{tabularx}{\textwidth}{@{}r X r r l@{}}
\toprule
$i$ & Absorbed since previous squeeze ($a_i$) & bytes & words & squeezed to \\
\midrule
1 & \yul{vk\_digest} (32) $+$ \yul{committed\_pi} identity as 128 zero bytes $+$ instance count word \texttt{19} (32) $+$ 19 instances (608) $+$ 15 phase-1 advice $\mathsf{A}_j$ (1920) & 2720 & 85 & $\theta$ (1568) \\
2 & reseed (32) $+$ 2 lookup multiplicity $\GO$ (256) & 288 & 9 & $\beta$ (1586) \\
3 & reseed (32) & 32 & 1 & $\gamma$ (1587) \\
4 & reseed (32) $+$ 6 permutation $\mathsf{Z}$ (768) $+$ 2 lookup helpers (256) $+$ 2 lookup accumulators (256) & 1312 & 41 & $\tau$ (1646) \\
5 & reseed (32) $+$ 1 trashcan $\GO$ (128) & 160 & 5 & $y$ (1663) \\
6 & reseed (32) $+$ 4 quotient limbs $\mathsf{Q}_j$ (512) & 544 & 17 & $x$ (1686) \\
7 & reseed (32) $+$ 102 evaluations $\mathbf{e}$ (3264) & 3296 & 103 & $x_1$ (1720) \\
8 & reseed (32) & 32 & 1 & $x_2$ (1721) \\
9 & reseed (32) $+$ $\mathsf{F}$ (128) & 160 & 5 & $x_3$ (1733) \\
10 & reseed (32) $+$ 5 \texttt{q\_evals} (160) & 192 & 6 & $x_4$ (1767) \\
\midrule
-- & reseed (32) $+$ $\pi$ (128), never hashed & 160 & 5 & -- \\
\bottomrule
\end{tabularx}
\caption{The complete absorb/squeeze schedule of this fixture. "reseed" is the
32-byte digest written by the preceding \yul{squeeze\_to}. Line numbers are the
\yul{squeeze\_to} call sites.}
\label{tab:tsh:schedule}
\end{table}

\paragraph{Buffer extent.} The peak cursor is reached in segment 7:
\[
  L_{\max} = T + 3296 = \texttt{0x1000} + \texttt{0xce0} = \texttt{0x1ce0},
\]
with segment 1 a close second at \texttt{0x1aa0}. The transcript window
$[\texttt{0x1000}, \texttt{0x1ce0})$ is disjoint from every other live region:
\yul{scalar\_inv}'s EIP-198 frame sits at \texttt{0x3580}--\texttt{0x3640}
(line 710, six words from \texttt{0x3580}), the VK payload begins at
\mptr{VK\_MPTR} $=$ \texttt{0x3680}, and every remaining generated region --
challenge slots from \mptr{THETA\_MPTR} $=$ \texttt{0x7900}, the PCS commitment
slots \mptr{F\_COM\_MPTR} $=$ \texttt{0x7a40} and \mptr{PI\_MPTR} $=$
\texttt{0x7ac0}, the accumulator and pairing slots at
\texttt{0x7b40}--\texttt{0x7f00}, the evaluation spill
\mptr{REVERSED\_EVALS\_MPTR} $=$ \texttt{0x9480}, and the commitment bases from
\texttt{0xa140} -- lies at $\ge$ \texttt{0x7900}. Below the transcript, the
entry guard at line 674 (\yul{if gt(mload(0x40), TRANSCRIPT\_MPTR)}) fails
closed with \revert{MemoryLayoutViolated} if solc's spill reservation ever
reaches \texttt{0x1000}. See Finding~TSH-5.

\paragraph{Cost.} Total hashed input is $8736$ bytes $= 273$ words across 10
squeezes, so the \opcode{KECCAK256} work is
$10 \cdot 30 + 273 \cdot 6 = 1938$ gas -- negligible. (That is the opcode
schedule only; the one-off memory expansion to \texttt{0x1ce0} and the
\opcode{MSTORE}/\opcode{CALLDATACOPY} of the absorbs are charged separately and
are likewise small next to the EIP-2537 precompiles.) The streaming design is
what buys this: a naive "re-hash the whole transcript at each challenge"
implementation would be $\Theta(m \cdot n)$ in absorbed bytes rather than
$\Theta(n + 32m)$, and would additionally need the entire transcript kept live
rather than a window that collapses to 32 bytes at each squeeze.

\subsection{Range boundary: \texttt{scalar\_inv} and the dead transcript window}
\label{sec:tsh:boundary}

\yul{scalar\_inv} (lines 700--722) is the only precompile call inside the
assigned range 700--836. Its arithmetic belongs to the field-arithmetic
section; the facts recorded here are the ones this section depends on, plus the
call catalogue the audit contract requires.

\begin{table}[h]
\centering
\begin{tabularx}{\textwidth}{@{}l l X@{}}
\toprule
Property & Value & Source \\
\midrule
Guard 1 & \yul{if iszero(lt(x, FR\_MODULUS))} $\Rightarrow$ \revert{NonCanonicalScalar} & line 708 \\
Guard 2 & \yul{if iszero(x)} $\Rightarrow$ \revert{NonCanonicalScalar} & line 709 \\
Frame base & \yul{let p := 0x3580} (literal, not a named constant) & line 710 \\
Frame layout & \texttt{0x20},\,\texttt{0x20},\,\texttt{0x20},\,$x$,\,$r-2$,\,$r$ at \texttt{+0x00,0x20,0x40,0x60,0x80,0xa0} & lines 713--718 \\
Precompile & \opcode{STATICCALL} address \texttt{0x05} (\texttt{MODEXP}) & line 719 \\
Gas bound & \yul{MODEXP\_GAS} $= 4080$ & line 293 \\
Argument window & in \texttt{[0x3580, 0x3640)} ($\texttt{0xc0}$ bytes), out \texttt{[0x3580, 0x35a0)} ($\texttt{0x20}$ bytes) & line 719 \\
Call failure & \yul{iszero(staticcall(...))} $\Rightarrow$ \revert{PrecompileFailed} & line 719 \\
Return-size check & \yul{iszero(eq(returndatasize(), 0x20))} $\Rightarrow$ \revert{PrecompileFailed} & line 720 \\
\bottomrule
\end{tabularx}
\caption{The single precompile call inside the assigned range. Full treatment is
in the field-arithmetic section.}
\label{tab:tsh:scalarinv}
\end{table}

The adjacency argument this section needs is that the frame
$[\texttt{0x3580}, \texttt{0x3640})$ lies inside the transcript \emph{arena}
$[T, \mptr{VK\_MPTR}) = [\texttt{0x1000}, \texttt{0x3680})$ but strictly above
the live transcript window $[\texttt{0x1000}, \texttt{0x1ce0})$, and that it
stops $\texttt{0x40}$ bytes short of \mptr{VK\_MPTR}. The comment at lines
695--699 asserts that \yul{scalar\_inv} "calls this only after transcript
absorption is complete" and "reuses the dead transcript buffer just below
\yul{VK\_MPTR}". A comment is not a proof, so it was checked against the
code: the five call sites are lines 3673, 3717, 3753, 3782 and 3802, all of them
after the transcript block ends at line 1789. The claim holds for this fixture,
but it is a whole-program ordering property, not a local one -- \yul{scalar\_inv}
carries no guard that would stop an emitter from calling it mid-transcript, and
doing so would silently corrupt bytes \texttt{0x3580}--\texttt{0x3640} of a
transcript segment that had grown past \texttt{0x2580} bytes.

\subsection{Assessment}
\label{sec:tsh:assess}

\subsubsection{What is clean, and why}

\begin{itemize}
  \item \textbf{Reseed uses the digest, not the challenge.} Line 804 stores
        $h$; line 807 stores $h \bmod r$ at a different address. The chain
        keeps 256 bits of entropy and matches the native
        \yul{TranscriptHash} implementation exactly.
  \item \textbf{Absorb is verbatim.} \yul{common\_uncompressed\_g1} hashes the
        same calldata bytes it validated (line 794), and \yul{common\_word}
        hashes the same word the caller range-checked. There is no
        canonicalisation step between validation and hashing, so no
        validate/hash divergence is possible.
  \item \textbf{The $\Fp$ range check is exact.} G3/G4 accept exactly
        $[0, p)$, with no off-by-one, and use canonical-boolean operands so the
        bitwise \yul{and}/\yul{or} are safe.
  \item \textbf{The padding check is complete.} \yul{shr(128, w)} tests all
        16 pad bytes, not a prefix of them.
  \item \textbf{All calldata reads are in bounds.} \yul{calldataload} past
        \yul{calldatasize} silently returns zeros, which would turn a truncated
        proof into a stream of identity points. That cannot happen here:
        \yul{calldatasize()} is pinned to
        $\mptr{INSTANCE\_CPTR} + \texttt{0x260} = \texttt{0x2144}$ at line 1426,
        the proof length is pinned to \texttt{0x1e60} at line 1419, the ABI heads
        are pinned at line 636, and the parser must land exactly on
        \mptr{NUM\_INSTANCE\_CPTR} at line 1785. (Those checks are specified in
        the calldata-shape section.)
  \item \textbf{Fixture matches template.} Lines 724--809 are byte-identical to
        \texttt{templates/partials/verifier/AssemblyHelpers.yul} lines 38--123;
        this block contains no template substitution at all. There is nothing to
        report under the "fixture wins" rule.
\end{itemize}

\subsubsection{Findings}

\begin{finding}[TSH-1]{No domain separation and no per-item type tags in the Keccak transcript}
\textbf{Severity: Informational.}
\yul{transcript\_init} (line 733) installs no domain separator, and neither
\yul{common\_word} (740) nor \yul{common\_uncompressed\_g1} (775) prefixes a
type or length tag; absorbs are raw concatenation. The same protocol's Blake2b
transcript does the opposite -- it keys the state with
\texttt{"Domain separator for transcript"} and tags absorbs and squeezes with
\yul{BLAKE2B\_PREFIX\_COMMON} / \yul{BLAKE2B\_PREFIX\_CHALLENGE} -- so the
omission is specific to the Keccak variant, and the Solidity verifier is
faithful to it.

\emph{Absorb ambiguity.} There is none \emph{within} this verifier. Every
absorbed item has a fixed length (32 or 128 bytes), the number of items in every
segment is a compile-time constant baked into the loop bounds
(Table~\ref{tab:tsh:schedule}), and the parser is pinned to consume exactly the
generated layout (lines 1419, 1420, 1426, 1785). The map
$(\text{instances}, \text{proof}) \mapsto (a_1,\dots,a_{10})$ is therefore
injective, which is the property tags would otherwise have to supply.
\emph{Cross-instance} separation is supplied by the first absorbed word: the VK
digest (line 1481), which commits to the constraint system before any proof byte
is read, followed by the absorbed instance count. What is genuinely absent is
\emph{cross-protocol} separation: no string identifies "this is a Halo2 KZG
BLS12-381 transcript", so an unrelated protocol that happened to Keccak the same
byte string would produce the same challenges. That is a theoretical concern
only, and it cannot be fixed in the Solidity verifier alone -- any tag added here
would have to be added to midnight-proofs first or every honest proof would fail.

\emph{Length extension.} Not a concern, twice over. Keccak256 is a sponge with
capacity and is not length-extendable, and independently, each squeeze is a
\emph{fresh} \yul{keccak256} over a bounded memory range rather than a
continuation of a partially absorbed state.
\end{finding}

\begin{finding}[TSH-2]{\texttt{mod r} challenge sampling is measurably biased}
\textbf{Severity: Informational.}
Line 807 samples $c = \mathrm{int}_{\mathrm{BE}}(h) \bmod r$ from a single
256-bit digest. Since $r$ is a 255-bit prime with
$r/2^{256} \approx 0.452845$, we have $\lfloor 2^{256}/r \rfloor = 2$ and
$t := 2^{256} \bmod r = 2^{256} - 2r \approx 0.2083\,r$. Assuming $h$ is
uniform on $\{0,1\}^{256}$:
\[
  \Pr[c = v] =
  \begin{cases}
    3 \cdot 2^{-256}, & v < t \quad (\approx 20.83\% \text{ of } \Fr),\\[2pt]
    2 \cdot 2^{-256}, & v \ge t \quad (\approx 79.17\% \text{ of } \Fr).
  \end{cases}
\]
The statistical distance from uniform on $\Fr$ is
\[
  \Delta = \tfrac12\Bigl[\,t\bigl|\tfrac{3}{2^{256}}-\tfrac1r\bigr| +
  (r-t)\bigl|\tfrac{2}{2^{256}}-\tfrac1r\bigr|\Bigr] \approx 0.0747,
\]
which is enormous by the usual "hash at least 64 bits wider than the modulus"
rule of thumb (that rule would demand $\Delta < 2^{-64}$). The
\emph{security}-relevant quantity is milder: the largest per-value probability
is $3\cdot 2^{-256} = 1.3585/r$, so the min-entropy is $254.42$ bits against
$\log_2 r = 254.86$, and any single-challenge soundness bound degrades by a
factor of at most $1.3585$ ($0.442$ bits). Across all ten squeezed challenges the
worst-case degradation compounds to $1.3585^{10} \approx 21.4$, i.e.\ at most
$4.4$ bits of the roughly $254$-bit challenge space. No concrete soundness bound
for this circuit is derivable from the assigned range, so none is asserted here;
the point is only that the loss is a small additive number of bits, not a
structural break.

\emph{Cost/benefit.} \yul{mod} is one opcode (5 gas); rejection sampling or a
512-bit wide reduction would remove the bias, but both would change the sampled
value and therefore break agreement with the native
\yul{sample\_keccak\_digest\_be\_mod\_r}
(\texttt{midfall/proofs/src/transcript/implementors.rs}), which computes
\yul{BigUint::from\_bytes\_be(\&hash\_output) \% modulus} -- bit-for-bit the
same operation. The verifier has no freedom here: this is a property of the
protocol, not of the Solidity rendering, and it must be fixed (if ever) on the
Rust side first.
\end{finding}

\begin{finding}[TSH-3]{Curve and subgroup validation is deferred out of the contract}
\textbf{Severity: Informational.}
\yul{common\_uncompressed\_g1} (lines 775--796) validates encoding and $\Fp$
range only. A caller can therefore absorb 128 bytes that decode to a
well-formed-but-not-on-curve $(X,Y)$, and the transcript will accept it. The
stated mitigation (lines 766--770) is a \emph{generator-side} invariant:
\yul{ProtocolPlan::validate} rejects plans in which an absorbed commitment is not
later consumed by an EIP-2537 \texttt{G1MSM} or pairing input, and those
precompiles perform the curve and subgroup checks.

A generated comment is not a proof, so the invariant was checked directly
against this fixture rather than accepted. All 34 absorbed $\GO$ points are
spilled into seven contiguous commitment regions plus two singleton slots, and
every one of those is read back into an EIP-2537 precompile argument:
\texttt{0xa140} (15 advice, line 3963), \texttt{0xa8c0} (2 multiplicities,
3838), \texttt{0xa9c0} (6 permutation $\mathsf{Z}$, 3985), \texttt{0xacc0}
(2 lookup helpers, 3840), \texttt{0xadc0} (2 lookup accumulators, 3959),
\texttt{0xaec0} (1 trashcan, 3846) and \mptr{QUOTIENT\_LIMB\_COMMS\_MPTR\_BASE}
$=$ \texttt{0xaf40} (4 limbs, 3920--3929) all \yul{mcopy} into
\mptr{PCS\_FINAL\_MSM\_MPTR} $=$ \texttt{0xb280}, which is submitted as a
78-pair \texttt{G1MSM} at line 3999 (\opcode{STATICCALL} \texttt{0x0c},
\texttt{0x30c0} input bytes, return size checked $= \texttt{0x80}$ at 4000);
\mptr{F\_COM\_MPTR} joins the same buffer at 3995; and \mptr{PI\_MPTR} becomes
\mptr{PAIRING\_LHS\_MPTR} at 4011, i.e.\ a pairing argument. The trashcan
commitment was the case worth checking, since \mptr{TRASHCAN\_COMMS\_MPTR\_BASE}
is written only at lines 1650--1657 and never named again -- it is read by its
raw address \texttt{0xaec0} at line 3846 -- and it does reach the MSM. So the
invariant holds \emph{for this fixture}; six of the seven regions are likewise
read only by literal address, which is why the check is easy to get wrong.

What remains is a process finding, not a live defect: the property is not
locally checkable from the deployed bytecode, the contract asserts nothing about
it, and re-verifying it required tracing eight memory bases through 2000 lines of
generated code. Should a future plan absorb a commitment that is never consumed,
the omission would be silent at the transcript layer. Note also that an
off-curve point that \emph{is} consumed produces a precompile rejection, i.e.\
the failure mode is fail-closed (\revert{PrecompileFailed} /
\revert{ProofRejected}), never a silent accept.
\end{finding}

\begin{finding}[TSH-4]{The identity point is accepted for every absorbed commitment}
\textbf{Severity: Observation.}
$(X,Y) = (0,0)$ -- 128 zero bytes -- passes G1--G4 unmodified, so any proof
commitment may be the point at infinity. This is intentional and matches the
native encoding, where \yul{G1Projective::to\_input} emits 128 zero bytes for the
identity; indeed the verifier itself absorbs exactly that string for
\yul{committed\_pi} at lines 1496--1500. Two consequences worth recording: (i)
the byte string absorbed for \yul{committed\_pi} is indistinguishable from an
advice commitment that happens to be the identity, though their transcript
\emph{positions} differ, so no ambiguity arises
(Finding~TSH-1); and (ii) a commitment set to the identity is a commitment to the
zero polynomial, which the KZG opening relations handle without special-casing.
No action; recorded because a reviewer expecting an "is-identity" rejection will
not find one.
\end{finding}

\begin{finding}[TSH-5]{Buffer writes are unbounded locally and bounded only by the fixed schedule}
\textbf{Severity: Observation.}
Neither \yul{common\_word} (line 741) nor \yul{common\_uncompressed\_g1}
(line 794) compares \yul{buf\_len} against any upper limit before writing. A
schedule that absorbed more than
$\mptr{VK\_MPTR} - \mptr{TRANSCRIPT\_MPTR} = \texttt{0x2680}$ bytes
$= 9856$ bytes in one segment would silently overwrite the loaded verifying key,
which is read after transcript absorption (for example \yul{OMEGA\_MPTR} in the
Lagrange block). For this fixture the peak is $L_{\max} = \texttt{0x1ce0}$, a
margin of \texttt{0x19a0} ($6560$ bytes), so the condition is comfortably
satisfied; the bound is a whole-program property established by the generator's
memory arena (lines 106--118), not by a runtime guard. Recorded so that anyone
re-targeting the generator to a larger circuit -- more instances, more advice
columns, or many more evaluations, all of which land in segments 1 and 7 --
re-derives the margin rather than assuming it.
\end{finding}

\subsubsection{Open points}

\begin{enumerate}
  \item The generated comment at lines 766--770 attributes the
        "every absorbed commitment reaches a precompile" invariant to
        \yul{ProtocolPlan::validate}. The \emph{consequence} of that claim was
        verified here for all 34 points against the deployed fixture
        (Finding~TSH-3); the \emph{generator-side enforcement} was not read, so
        it is the comment, not the code, that asserts the property will hold for
        future plans.
  \item The native \yul{Hashable<Keccak256>} instance for the VK digest word --
        specifically whether it is absorbed as an $\Fq$ element and therefore
        necessarily $< r$ -- was inferred from the fixture comment at lines
        1475--1481 together with the pinned value \texttt{0x56c0824f...} in
        \texttt{Halo2VerifyingKey.sol}:68, which is $< r$. Since the VK runtime
        is pinned by codehash, a non-canonical digest could not vary per call in
        any case, and the worst outcome of a mismatch is fail-closed
        non-verification rather than a soundness break.
  \item The per-segment item counts in Table~\ref{tab:tsh:schedule} were derived
        by resolving the loop bounds at lines 1553--1773 and cross-checked
        against the pinned proof length:
        $34 \cdot 128 + 107 \cdot 32 = 7776 = \texttt{0x1e60}$,
        exactly the value asserted at line 1419. Those call
        sites are outside this section's range and belong to the
        transcript-schedule section.
\end{enumerate}


%% ==== 06-ec-helpers.tex =================================================
\section{EIP-2537 precompile wrappers: batch inversion and pairing}
\label{sec:ec:top}

This section specifies \texttt{Halo2Verifier.sol} lines 811--991, the block the
generator labels \emph{EC primitives (EIP-2537 wrappers)}. It contains exactly
two Yul functions:

\begin{itemize}
  \item \yul{batch\_invert(success, mptr\_start, mptr\_end, scratch\_mptr, r)}
        (lines 836--951), Montgomery batch inversion of a contiguous run of
        $\Fr$ words, in place, using a single call to the EIP-198 modexp
        precompile at address \texttt{0x05};
  \item \yul{ec\_pairing(success, lhs\_mptr, rhs\_mptr)} (lines 955--990), the
        terminal two-pair BLS12-381 pairing check through the EIP-2537
        \opcode{PAIRING\_CHECK} precompile at address \texttt{0x0f}.
\end{itemize}

The banner comment at lines 811--816 records the two structural premises the
rest of the block relies on: these helpers operate on \emph{four-word}
EIP-2537-padded $\GO$ points rather than BN254's two-word points, and they
allocate their scratch from planned windows above Solidity's reserved memory
prefix, reusing addresses whose earlier occupant (the streaming transcript
buffer) is provably dead once all challenges have been squeezed. Neither
premise is re-checked on chain inside this range; both are generation-time
properties of the memory arena (\texttt{src/lowering/layout/memory.rs}).

\subsection{Shared constants and the staticcall discipline}
\label{sec:ec:constants}

Table~\ref{tab:ec:constants} collects every named constant this range reads.
All are \yul{internal constant} declarations that \texttt{solc} inlines as
\opcode{PUSH} immediates; none is loaded from memory or calldata.

\begin{table}[htbp]
\centering
\begin{tabular}{@{}llp{6.6cm}@{}}
\toprule
Constant & Value & Role \\
\midrule
\mptr{MODEXP\_GAS}        & $4080$      & Gas forwarded to \texttt{0x05} (line 293) \\
\mptr{PAIRING\_GAS\_2PAIR}& $102900$    & Gas forwarded to \texttt{0x0f} (line 292) \\
\mptr{G2\_BASE\_MPTR}     & \texttt{0x3860} & $[1]_2$, 8 words, in the pinned VK payload (line 141) \\
\mptr{NEG\_S\_G2\_BASE\_MPTR} & \texttt{0x3960} & $-[s]_2$, 8 words (line 142) \\
\mptr{ERR\_PRECOMPILE\_FAILED} & \texttt{0x84e81692} & \revert{PrecompileFailed} selector (line 324) \\
\mptr{ERR\_PROOF\_REJECTED}    & \texttt{0xc3b0d8cd} & \revert{ProofRejected} selector (line 325) \\
\bottomrule
\end{tabular}
\caption{Constants referenced by lines 811--991, with their declaration sites.}
\label{tab:ec:constants}
\end{table}

Two disciplines are applied identically at all three precompile call sites in
this range (lines 873, 924, 980).

\paragraph{Exact-cost gas bounds, never \opcode{GAS}.}
Every \opcode{STATICCALL} forwards a literal constant rather than
\yul{gas()}. The rationale is recorded at lines 268--289 (outside this range):
a modexp or EIP-2537 precompile that rejects its input consumes \emph{all} gas
supplied to the child frame, so forwarding \yul{gas()} would let one malformed
input burn $63/64$ of the transaction budget, whereas forwarding the scheduled
cost caps the loss at that cost. The two bounds used here are:
\begin{align}
  \mptr{PAIRING\_GAS\_2PAIR} &= 37700 + 2 \cdot 32600 = 102900,
  \label{eq:ec:pairgas}\\
  \mptr{MODEXP\_GAS} &= 4080 \;\ge\; 16 \cdot 254 = 4064 .
  \label{eq:ec:modexpgas}
\end{align}
Equation~\eqref{eq:ec:pairgas} is the EIP-2537 \opcode{PAIRING\_CHECK} schedule
for $k=2$ pairs, exactly. Equation~\eqref{eq:ec:modexpgas} is the EIP-7883
(Osaka) price of the $32/32/32$-byte modexp frame used here: multiplication
complexity $16$ for operands of at most 32 bytes, iteration count
$\mathrm{bitlen}(r-2)-1 = 254$, and no $/3$ divisor. The declaration comment at
lines 277--282 computes the bound with $\mathrm{bitlen} = 256$, i.e.\ with
$255$ iterations instead of $254$, giving $4080$ instead of $4064$; the surplus
is harmless because unused gas is refunded on success. The same off-by-one
shows in that comment's EIP-2565 figure: it states $1360 = \lfloor 16\cdot
255/3\rfloor$ where the true cost of this frame is
$\lfloor 16\cdot 254/3\rfloor = 1354$ (multiplication complexity is
$\lceil 32/8\rceil^2 = 16$ under EIP-2565 as well, so only the $/3$ divisor and
the $200$ floor differ). Either way $4080$ over-forwards on a pre-Osaka
chain and is exact-or-generous on a post-Osaka one. Forwarding the smaller of
the two would have bricked the verifier at the fork; forwarding the larger is
free. This is the right choice.

\paragraph{Success \emph{and} return size.}
Precompile success is never taken from the \opcode{STATICCALL} return flag
alone. Each site immediately folds in an equality on \yul{returndatasize()}:

\begin{equation}
  \mathit{ok} \;=\; \opcode{STATICCALL}(\cdot) \;\wedge\;
  \bigl(\yul{returndatasize()} = \texttt{0x20}\bigr).
  \label{eq:ec:okdef}
\end{equation}

This is load-bearing for two distinct reasons. First, \opcode{STATICCALL} to an
address with no code returns $1$ with empty return data; on a chain that has
not activated EIP-2537, the call to \texttt{0x0f} at line 980 would therefore
\emph{succeed} and write nothing, leaving the output window holding the stale
input bytes that were just \yul{mcopy}'d there. Second, the EVM copies only
$\min(\yul{returndatasize()}, \mathit{outsize})$ bytes into the output window,
left-aligned, so a short-returning precompile leaves the \emph{low}-order bytes
of the result word stale --- the return data lands in the high bytes and the
tail keeps whatever was there before, which at these three sites is the frame
the caller just wrote. The equality in \eqref{eq:ec:okdef} closes both. In every case the
guard is evaluated \emph{before} the result word is loaded --- see lines
874--876, 925--933 and 981--987 --- so no stale word ever reaches the
arithmetic.

\paragraph{What \yul{fail} does.} Every revert named in this
section goes through the helper at \texttt{Halo2Verifier.sol}:690--693,
\yul{function fail(sel) \{ mstore(0x00, shl(224, sel)) revert(0x00, 0x04) \}}:
a bare 4-byte selector written into Solidity's own scratch word at
\texttt{0x00}, which never overlaps the generated arena. No revert in this
range carries arguments.

\subsection{\yul{batch\_invert}: Montgomery batch inversion}
\label{sec:ec:batchinvert}

\subsubsection{Interface and contract}

\begin{lstlisting}[style=solidity,caption={\texttt{Halo2Verifier.sol} line 836: signature},label={lst:ec:sig}]
function batch_invert(success, mptr_start, mptr_end, scratch_mptr, r) -> ret, precompile_failed {
\end{lstlisting}

The function inverts every word in the half-open memory range
$[\yul{mptr\_start}, \yul{mptr\_end})$ modulo \yul{r}, writing each inverse
back over its input. Let
\begin{equation}
  m \;=\; \frac{\yul{mptr\_end} - \yul{mptr\_start}}{32},
  \qquad
  d_i \;=\; \yul{mload}(\yul{mptr\_start} + 32i), \quad 0 \le i < m ,
\end{equation}
and write $S_k = \prod_{i=0}^{k} d_i$ for the prefix products, so
$S_{m-1}$ is the total product.

It returns two booleans rather than reverting. \yul{ret} is the conjunction of
the incoming \yul{success} with "this inversion is valid"; \yul{precompile\_failed}
is set only when the modexp \opcode{STATICCALL} itself failed the test in
\eqref{eq:ec:okdef}. The comment at lines 824--835 states the intent: callers
combine \yul{ret} with other \yul{success} plumbing and convert to a revert at
a section boundary, and the second flag lets that boundary distinguish a chain
or gas-schedule fault (\revert{PrecompileFailed}) from a rejected input
(\revert{ProofRejected}). For the Lagrange batch, a rejected input means the
squeezed $x$ landed on a domain point, an event of probability $\approx n/r$.

\subsubsection{Guards and the single-element fast path}

\begin{lstlisting}[style=solidity,caption={\texttt{Halo2Verifier.sol} lines 837--878: entry guards and the $m=1$ path (rationale comments elided; line 847 kept)},label={lst:ec:fast}]
                ret := success
                if iszero(ret) { leave }
                if lt(mptr_end, mptr_start) {
                    ret := 0
                    leave
                }

                let count_bytes := sub(mptr_end, mptr_start)
                // Empty batch is valid and leaves memory untouched.
                if iszero(count_bytes) { leave }

                if eq(count_bytes, 0x20) {
                    let x := mload(mptr_start)
                    if iszero(lt(x, r)) {
                        ret := 0
                        leave
                    }
                    if iszero(x) {
                        ret := 0
                        leave
                    }

                    let single_scratch := scratch_mptr
                    mstore(add(single_scratch, 0x00), 0x20)
                    mstore(add(single_scratch, 0x20), 0x20)
                    mstore(add(single_scratch, 0x40), 0x20)
                    mstore(add(single_scratch, 0x60), x)
                    mstore(add(single_scratch, 0x80), sub(r, 2))
                    mstore(add(single_scratch, 0xa0), r)
                    ret := staticcall(MODEXP_GAS, 0x05, single_scratch, 0xc0, single_scratch, 0x20)
                    ret := and(ret, eq(returndatasize(), 0x20))
                    precompile_failed := iszero(ret)
                    if ret { mstore(mptr_start, mload(single_scratch)) }
                    leave
                }
\end{lstlisting}

Four guards fire before any arithmetic:
\begin{enumerate}
  \item line 838, short-circuit on an already-failed \yul{success}: returns
        $(\yul{ret}, \yul{precompile\_failed}) = (0, 0)$, so an earlier failure
        is never re-attributed to the modexp precompile;
  \item lines 841--844, \yul{lt(mptr\_end, mptr\_start)}: a reversed range
        returns $(0,0)$. Note the comparison is strict, so
        $\yul{mptr\_end} = \yul{mptr\_start}$ falls through to the next guard;
  \item line 848, an empty range returns \yul{ret := success} unchanged and
        touches no memory;
  \item line 852, $\yul{count\_bytes} = \texttt{0x20}$ takes the single-element
        path.
\end{enumerate}

The $m=1$ path rejects $d_0 \ge r$ (line 857) \emph{and} $d_0 = 0$ (line 861).
The first check is not redundant: modexp reduces its base modulo the modulus,
so for any $d_0 \equiv 0 \pmod r$ --- $d_0 = r$, say --- it would return $0$,
and the caller would multiply by that zero as though it were a valid inverse.
The comment at lines 854--856 says exactly this.

Lines 867--872 lay out the EIP-198 frame
$[\,\mathit{base\_len}, \mathit{exp\_len}, \mathit{mod\_len},
\mathit{base}, \mathit{exp}, \mathit{mod}\,]$
$= [\texttt{0x20}, \texttt{0x20}, \texttt{0x20}, d_0, r-2, r]$ at
\yul{scratch\_mptr}, six words, and calls \texttt{0x05} with
$\mathit{insize} = \texttt{0xc0}$, $\mathit{outsize} = \texttt{0x20}$, output
aliased onto the input base. Fermat's little theorem gives
$d_0^{\,r-2} \equiv d_0^{-1} \pmod r$ for $d_0 \ne 0$, which is why the zero
check must precede the call. Aliasing input and output is safe: the child
frame's calldata is snapshotted at call time and the return data is copied back
only after the call returns.

Line 876 writes the result back \emph{only} when \yul{ret} holds. On failure
the input word is left intact, so a caller that ignores the return value reads
back its own denominator rather than a stale frame header --- an anti-footgun
that recurs in the batch path.

\subsubsection{Forward pass}

\begin{lstlisting}[style=solidity,caption={\texttt{Halo2Verifier.sol} lines 886--915: prefix products and the total-product zero check},label={lst:ec:fwd}]
                let gp_mptr := scratch_mptr
                let gp := mload(mptr_start)
                if iszero(lt(gp, r)) {
                    ret := 0
                    leave
                }
                let mptr := add(mptr_start, 0x20)
                for {} lt(mptr, sub(mptr_end, 0x20)) {} {
                    let x := mload(mptr)
                    if iszero(lt(x, r)) {
                        ret := 0
                        leave
                    }
                    gp := mulmod(gp, x, r)
                    mstore(gp_mptr, gp)
                    mptr := add(mptr, 0x20)
                    gp_mptr := add(gp_mptr, 0x20)
                }
                let x_last := mload(mptr)
                if iszero(lt(x_last, r)) {
                    ret := 0
                    leave
                }
                gp := mulmod(gp, x_last, r)
                if iszero(gp) {
                    ret := 0
                    leave
                }
\end{lstlisting}

The loop bound is \yul{lt(mptr, sub(mptr\_end, 0x20))}, i.e.\ it runs for
$i = 1, \dots, m-2$ and stops with \yul{mptr} pointing at $d_{m-1}$. On
iteration $i$ it stores $S_i$ at $\yul{scratch\_mptr} + 32(i-1)$, so scratch
ends up holding
\begin{equation}
  S_1, S_2, \dots, S_{m-2}
  \quad\text{at}\quad
  \yul{scratch\_mptr} + 32\cdot 0, \dots, \yul{scratch\_mptr} + 32(m-3),
  \label{eq:ec:prefix}
\end{equation}
which is $m-2$ words. $S_0 = d_0$ and $S_{m-1}$ are deliberately \emph{not}
stored: the first is recoverable from the batch itself and the last is carried
in \yul{gp}.

\paragraph{Canonicality coverage is complete.} The check
$\yul{x} < r$ is applied to $d_0$ at line 888, to $d_1 \dots d_{m-2}$ at line
895, and to $d_{m-1}$ at line 905 --- every element of the run, with no gap.
The comment at lines 883--885 gives the reason: without it \yul{mulmod} would
silently reduce a non-canonical word, making accept/reject semantics depend on
whether the batch happened to be of length one or more. As written, the
single-element and multi-element paths agree.

\paragraph{Zero detection.} Individual zeros are not tested. Instead line 912
rejects when the \emph{total} product is zero. Since $\Fr$ is a field and hence
has no zero divisors, $S_{m-1} = 0 \iff \exists i:\, d_i = 0$, so the single
test is exactly equivalent to testing each element, at $1$ comparison instead
of $m$. This is sound. The cost is that the rejection is not localized: the
caller learns that \emph{some} denominator vanished, not which.

\subsubsection{The single modexp and the backward pass}

\begin{lstlisting}[style=solidity,caption={\texttt{Halo2Verifier.sol} lines 918--950: one modexp, then the backward pass (comments elided)},label={lst:ec:bwd}]
                mstore(add(gp_mptr, 0x00), 0x20)
                mstore(add(gp_mptr, 0x20), 0x20)
                mstore(add(gp_mptr, 0x40), 0x20)
                mstore(add(gp_mptr, 0x60), gp)
                mstore(add(gp_mptr, 0x80), sub(r, 2))
                mstore(add(gp_mptr, 0xa0), r)
                ret := staticcall(MODEXP_GAS, 0x05, gp_mptr, 0xc0, gp_mptr, 0x20)
                ret := and(ret, eq(returndatasize(), 0x20))
                precompile_failed := iszero(ret)
                if iszero(ret) { leave }
                let all_inv := mload(gp_mptr)

                let first_mptr := mptr_start
                let second_mptr := add(first_mptr, 0x20)
                gp_mptr := sub(gp_mptr, 0x20)
                for {} lt(second_mptr, mptr) {} {
                    let inv := mulmod(all_inv, mload(gp_mptr), r)
                    all_inv := mulmod(all_inv, mload(mptr), r)
                    mstore(mptr, inv)
                    mptr := sub(mptr, 0x20)
                    gp_mptr := sub(gp_mptr, 0x20)
                }
                let inv_first := mulmod(all_inv, mload(second_mptr), r)
                let inv_second := mulmod(all_inv, mload(first_mptr), r)
                mstore(first_mptr, inv_first)
                mstore(second_mptr, inv_second)
\end{lstlisting}

The modexp frame is written at the \emph{current} \yul{gp\_mptr}, i.e.\ at
$\yul{scratch\_mptr} + 32(m-2)$, immediately past the last prefix product. It
therefore overlays six words of scratch that the forward pass has not used, and
its one-word output overwrites the frame's own \yul{base\_len} slot. The total
scratch footprint of the multi-element path is exactly
\begin{equation}
  \mathit{scratch\_bytes}(m) \;=\; 32(m-2) + \texttt{0xc0}.
  \label{eq:ec:scratchsize}
\end{equation}

Line 932, \yul{if iszero(ret) \{ leave \}}, is placed \emph{before}
\yul{mload(gp\_mptr)}. The comment at lines 927--931 explains why this early
exit is not cosmetic: a failed \opcode{STATICCALL} writes no output, so
\yul{mload(gp\_mptr)} would read back the stale \texttt{0x20} frame header, and
the backward pass would then overwrite every word of
$[\yul{mptr\_start}, \yul{mptr\_end})$ with garbage \emph{before} returning
$\yul{ret} = 0$. Any caller that inspected memory before checking the flag
would consume that garbage. The guard is correct and correctly ordered.

\paragraph{Correctness of the backward pass.} On entry, \yul{mptr} points at
$d_{m-1}$, \yul{gp\_mptr} has been decremented once to
$\yul{scratch\_mptr} + 32(m-3)$ (holding $S_{m-2}$), and
$\yul{all\_inv} = S_{m-1}^{-1}$. The loop invariant, at the top of the
iteration that processes index $i$ (running $i = m-1, m-2, \dots, 2$), is
\begin{equation}
  \yul{all\_inv} = S_i^{-1},
  \qquad
  \yul{gp\_mptr} \text{ holds } S_{i-1},
  \qquad
  \yul{mptr} = \yul{mptr\_start} + 32i .
  \label{eq:ec:inv}
\end{equation}
The body then computes
\begin{equation}
  \mathit{inv} = S_i^{-1} \cdot S_{i-1} = d_i^{-1},
  \qquad
  \yul{all\_inv} \leftarrow S_i^{-1}\cdot d_i = S_{i-1}^{-1},
\end{equation}
stores $\mathit{inv}$ over $d_i$ \emph{after} reading $d_i$, and decrements
both pointers, re-establishing \eqref{eq:ec:inv} for $i-1$. The loop condition
\yul{lt(second\_mptr, mptr)} with $\yul{second\_mptr} = \yul{mptr\_start}+32$
stops when \yul{mptr} reaches index $1$, so the last index handled inside the
loop is $i = 2$ and the loop exits with $\yul{all\_inv} = S_1^{-1} =
(d_0 d_1)^{-1}$. Lines 947--950 then finish the pair:
\begin{equation}
  d_0^{-1} = (d_0d_1)^{-1}\cdot d_1,
  \qquad
  d_1^{-1} = (d_0d_1)^{-1}\cdot d_0 ,
\end{equation}
computing both products before either store, so the read-after-write hazard on
the two adjacent slots is avoided. For $m = 2$ the loop body never executes and
these two lines are the whole backward pass, which is correct because
$\yul{all\_inv} = S_1^{-1}$ already.

Note that \yul{gp\_mptr} is decremented once more after the final iteration,
leaving it at $\yul{scratch\_mptr} - \texttt{0x20}$ (and, for $m=2$, it is
already $\yul{scratch\_mptr} - \texttt{0x20}$ before the loop). That value is
never dereferenced --- \yul{gp\_mptr} is read only inside the loop body --- so
the underflow below the scratch base is dead. This is worth stating because it
is the one pointer in the routine that leaves its own window.

\subsubsection{Formal statement of the routine}

\begin{algorithm}[htbp]
\caption{\yul{batch\_invert} (\texttt{Halo2Verifier.sol} 836--951)}
\label{alg:ec:batchinvert}
\begin{algorithmic}[1]
\Require $\mathit{success}$, base $A$, end $E$, scratch $S$, modulus $r$
\Ensure $(\mathit{ret}, \mathit{pf})$; on $\mathit{ret}=1$, $[A,E)$ holds the inverses
\State $\mathit{ret} \gets \mathit{success}$; $\mathit{pf} \gets 0$
\If{$\mathit{ret} = 0$} \Return $(0,0)$ \EndIf
\If{$E < A$} \Return $(0,0)$ \Comment{codegen fault: reversed range}\EndIf
\If{$E = A$} \Return $(\mathit{success},0)$ \Comment{empty batch, memory untouched}\EndIf
\If{$E - A = 32$}
  \State $x \gets \mathrm{mem}[A]$
  \If{$x \ge r$ \textbf{or} $x = 0$} \Return $(0,0)$ \EndIf
  \State write $[\,32,32,32,x,r-2,r\,]$ at $S$; $\mathit{ok} \gets \opcode{STATICCALL}(4080,\texttt{0x05},S,192,S,32)$
  \State $\mathit{ok} \gets \mathit{ok} \wedge (\yul{returndatasize()} = 32)$; $\mathit{pf} \gets \neg\,\mathit{ok}$
  \If{$\mathit{ok}$} $\mathrm{mem}[A] \gets \mathrm{mem}[S]$ \EndIf
  \State \Return $(\mathit{ok}, \mathit{pf})$
\EndIf
\State $g \gets \mathrm{mem}[A]$; \textbf{if} $g \ge r$ \textbf{then return} $(0,0)$
\State $q \gets S$; $\; a \gets A + 32$
\While{$a < E - 32$}
  \State $x \gets \mathrm{mem}[a]$; \textbf{if} $x \ge r$ \textbf{then return} $(0,0)$
  \State $g \gets g\cdot x \bmod r$; $\;\mathrm{mem}[q] \gets g$; $\;a \gets a+32$; $\;q \gets q+32$
\EndWhile
\State $x \gets \mathrm{mem}[a]$; \textbf{if} $x \ge r$ \textbf{then return} $(0,0)$
\State $g \gets g\cdot x \bmod r$; \textbf{if} $g = 0$ \textbf{then return} $(0,0)$
\State write $[\,32,32,32,g,r-2,r\,]$ at $q$; $\mathit{ok} \gets \opcode{STATICCALL}(4080,\texttt{0x05},q,192,q,32)$
\State $\mathit{ok} \gets \mathit{ok} \wedge (\yul{returndatasize()} = 32)$; $\mathit{pf} \gets \neg\,\mathit{ok}$
\If{$\neg\,\mathit{ok}$} \Return $(0,1)$ \Comment{batch left untouched}\EndIf
\State $u \gets \mathrm{mem}[q]$; $\; q \gets q - 32$
\While{$A + 32 < a$}
  \State $v \gets u\cdot \mathrm{mem}[q] \bmod r$; $\; u \gets u \cdot \mathrm{mem}[a] \bmod r$
  \State $\mathrm{mem}[a] \gets v$; $\; a \gets a-32$; $\; q \gets q-32$
\EndWhile
\State $v_0 \gets u\cdot \mathrm{mem}[A+32] \bmod r$; $\; v_1 \gets u\cdot \mathrm{mem}[A] \bmod r$
\State $\mathrm{mem}[A] \gets v_0$; $\;\mathrm{mem}[A+32] \gets v_1$
\State \Return $(1, 0)$
\end{algorithmic}
\end{algorithm}

\subsubsection{Exit table}

\begin{table}[htbp]
\centering
\begin{tabularx}{\textwidth}{@{}llccX@{}}
\toprule
Line & Condition & \yul{ret} & \yul{precompile\_failed} & Memory effect \\
\midrule
838 & \yul{success} $=0$ on entry & 0 & 0 & none \\
841 & $\yul{mptr\_end} < \yul{mptr\_start}$ & 0 & 0 & none \\
848 & $\yul{mptr\_end} = \yul{mptr\_start}$ & 1 & 0 & none \\
857 & $m=1$, $d_0 \ge r$ & 0 & 0 & none \\
861 & $m=1$, $d_0 = 0$ & 0 & 0 & none \\
874 & $m=1$, modexp failed or short return & 0 & 1 & scratch frame written \\
876 & $m=1$, success & 1 & 0 & $d_0 \gets d_0^{-1}$ \\
888/895/905 & some $d_i \ge r$ & 0 & 0 & nothing at 888; $i-1$ prefix products in scratch at 895/905 \\
912 & $S_{m-1} = 0$ & 0 & 0 & prefix products in scratch \\
925 & modexp failed or short return & 0 & 1 & scratch only; batch untouched \\
950 & fall-through & 1 & 0 & $d_i \gets d_i^{-1}$ for all $i$ \\
\bottomrule
\end{tabularx}
\caption{Every exit of \yul{batch\_invert}, \texttt{Halo2Verifier.sol} 836--951.}
\label{tab:ec:exits}
\end{table}

\subsubsection{The concrete memory window of the only call site}

\yul{batch\_invert} is called exactly once in the runtime, at line 1826 inside
the Lagrange block (that block is specified elsewhere):

\begin{lstlisting}[style=solidity,caption={\texttt{Halo2Verifier.sol} line 1826: the single call site},label={lst:ec:callsite}]
                success, lagrange_precompile_failed := batch_invert(success, LAGRANGE_DENOMS_MPTR, add(mptr_end, 0x20), BATCH_INV_SCRATCH_MPTR, r)
\end{lstlisting}

with $\mptr{LAGRANGE\_DENOMS\_MPTR} = \texttt{0xb580}$ (line 240),
$\mptr{BATCH\_INV\_SCRATCH\_MPTR} = \texttt{0xb140}$ (line 236) and
$\yul{mptr\_end} = \texttt{0xb580} + \texttt{0x3a0}$ (line 1819). The run is
therefore $[\texttt{0xb580}, \texttt{0xb940})$, that is
$\texttt{0x3c0}/32 = 30$ words: $29$ Lagrange denominators
($19$ public instances plus $|\mathrm{rotation}_{\mathrm{last}}| = 10$
negative rows) followed by $x^n - 1$. Substituting $m = 30$ into
\eqref{eq:ec:scratchsize}:

\begin{table}[htbp]
\centering
\begin{tabular}{@{}lll@{}}
\toprule
Window & Range & Contents \\
\midrule
prefix products & $[\texttt{0xb140}, \texttt{0xb4c0})$ & $S_1 \dots S_{28}$, 28 words \\
modexp frame    & $[\texttt{0xb4c0}, \texttt{0xb580})$ & $[32,32,32,S_{29},r-2,r]$, 6 words \\
\midrule
\emph{scratch total} & $[\texttt{0xb140}, \texttt{0xb580})$ & $\texttt{0x440}$ bytes \\
batch (in/out)  & $[\texttt{0xb580}, \texttt{0xb940})$ & 30 denominators, then 30 inverses \\
\bottomrule
\end{tabular}
\caption{Exact memory window of the Lagrange batch inversion, $m=30$.}
\label{tab:ec:window}
\end{table}

The scratch window ends at \texttt{0xb580}, which is precisely
\mptr{LAGRANGE\_DENOMS\_MPTR}. The modexp frame's last byte and the batch's
first byte are adjacent with \emph{zero} bytes of slack.

\begin{finding}[EC-1]{Batch-inversion scratch abuts its own input with zero slack, and the overrun would be silent}
\textbf{Severity: Low} (correct as generated; no runtime guard).

\texttt{Halo2Verifier.sol}:1826 passes $\yul{scratch\_mptr} = \texttt{0xb140}$
for a 30-word run based at \texttt{0xb580}. By
\eqref{eq:ec:scratchsize} the routine consumes $32\cdot 28 + \texttt{0xc0} =
\texttt{0x440}$ bytes, and $\texttt{0xb580} - \texttt{0xb140} = \texttt{0x440}$
exactly. If the run were one word longer, the modexp frame would begin at
\texttt{0xb4e0} and its sixth word --- the \emph{modulus} $r$ --- would land on
\texttt{0xb580}, i.e.\ on $d_0$. Only that one word is clobbered
($\texttt{0xb4e0}+\texttt{0xc0} = \texttt{0xb5a0}$, and $d_1$ sits at
\texttt{0xb5a0}). The failure is not detectable at runtime: the frame itself is
intact, so the modexp still succeeds, \eqref{eq:ec:okdef} still holds, and
\yul{ret} stays $1$. The forward pass had already folded the true $d_0$ into
\yul{gp}, so $\yul{all\_inv} = S_1^{-1}$ is still correct and line 947 still
yields $d_0^{-1}$; but line 948 reads the clobbered slot and returns
$(d_0d_1)^{-1}\cdot r$ in place of $d_1^{-1}$. One Lagrange denominator is
inverted wrongly, with no revert.

The generator is aware of this. \texttt{src/lowering/layout/memory.rs}:707--730
recomputes the template's run length independently and asserts both that the
two formulas agree and that
$\mathit{scratch\_bytes} \ge 32(n-2) + \texttt{0xc0}$, with a comment that
names the silent-corruption mode verbatim. So the property holds for this
fixture and is enforced at generation time. What is absent is any on-chain
expression of it: \yul{batch\_invert} receives \yul{scratch\_mptr} as an opaque
pointer and never bounds it against \yul{mptr\_start}. A one-line guard
--- reject when $\yul{scratch\_mptr} + 32(m-2) + \texttt{0xc0} >
\yul{mptr\_start}$ for the (sole, upward-growing) layout used here --- would
convert a silent miscomputation into a revert for a handful of gas. Recorded as
Low rather than higher because the deployed fixture is provably correct and the
only path to the bug is a regenerated artifact whose planner assertions were
also removed.
\end{finding}

\begin{finding}[EC-2]{No runtime word-alignment check on the batch range}
\textbf{Severity: Informational.}

\yul{batch\_invert} guards the \emph{direction} of the range
(\texttt{Halo2Verifier.sol}:841) but not its alignment. If
$\yul{mptr\_end}-\yul{mptr\_start}$ were not a multiple of $32$ --- say
\texttt{0x30} --- control would reach the multi-element path, the forward loop
bound \yul{lt(mptr, sub(mptr\_end, 0x20))} would be false immediately, and line
904 would \yul{mload} at $\yul{mptr\_start}+\texttt{0x20}$, reading
$\texttt{0x10}$ bytes past \yul{mptr\_end}. The comment at lines 839--840
asserts alignment is guaranteed "by construction" and, for the single call site
at line 1826, it is: the length is the constant $\texttt{0x3a0}+\texttt{0x20}$.
No exploitable path exists in this artifact. The observation is that the
routine's precondition is stronger than the guard it actually implements, so it
is not safe to reuse with a computed range.
\end{finding}

\begin{finding}[EC-3]{A codegen fault and a rejected proof exit through the same channel}
\textbf{Severity: Informational.}

The reversed-range exit (\texttt{Halo2Verifier.sol}:841--844) returns
$(\yul{ret}, \yul{precompile\_failed}) = (0,0)$, which is bit-for-bit the same
signal as "a denominator was zero or non-canonical". At the section boundary
(lines 1877--1883) that maps to \revert{ProofRejected}. The code's own comment
calls a reversed range "always a codegen error" --- the category the
\revert{MemoryLayoutViolated} selector (\texttt{0xc9888d23}, line 327) exists
for. Symmetrically, the empty-range exit at line 848 returns
$\yul{ret} := \yul{success} = 1$, so a mis-planned zero-length batch is
reported as a \emph{successful} inversion of nothing. Neither is reachable in
this fixture, where the range is a compile-time constant; both slightly blunt
the incident-response taxonomy that MF-4 (lines 828--835) was introduced to
sharpen.
\end{finding}

\paragraph{Deferred conversion to a revert.} At the call site, a
$\yul{ret} = 0$ does not stop the Lagrange block: lines 1828--1875 continue to
compute $L_i(x)$, $L_{\mathrm{blind}}$, $L_0$, $L_{\mathrm{last}}$ and
$\mathrm{instance\_eval}$ from words that are denominators rather than
inverses, and write them (lines 1869--1874) into \mptr{X\_N\_MPTR},
\mptr{X\_N\_MINUS\_1\_INV\_MPTR}, \mptr{L\_LAST\_MPTR},
\mptr{L\_BLIND\_MPTR}, \mptr{L\_0\_MPTR}, \mptr{INSTANCE\_EVAL\_MPTR}. Only at
line 1877 does \yul{if iszero(success)} fire, dispatching (lines 1881--1882) to
\revert{PrecompileFailed} when \yul{lagrange\_precompile\_failed} is set and
\revert{ProofRejected} otherwise. This is sound --- the revert happens before
the quotient section reads any of those slots --- but it does mean a rejected
batch pays for roughly $80$ wasted \yul{mulmod}s and $19$ \yul{calldataload}s
before reverting. Given that the only way to reach it is $x$ landing on a
domain point (probability $\approx 2^{20}/r \approx 2^{-235}$), the wasted gas
is irrelevant and the simpler control flow is the better trade.

\paragraph{What the batching buys.} The alternative --- one modexp per
denominator --- would forward $30 \times 4080 = 122{,}400$ gas of precompile
budget. The batched form forwards $4080$ once and pays roughly $3m-3 = 87$
extra \yul{mulmod}s at 8 gas each plus $28$ \yul{mstore}s, i.e.\ well under
$1{,}000$ gas. The saving is on the order of $10^5$ gas per proof, at the price
of a routine whose backward pass is genuinely subtle and whose scratch sizing
is a cross-file invariant (finding EC-1). For a verifier called once per proof
with a fixed, generator-checked shape, that is a good trade.

\paragraph{Address reuse.} \mptr{BATCH\_INV\_SCRATCH\_MPTR} is numerically
equal to \mptr{SELECTOR\_ACC\_MPTR} $=\texttt{0xb140}$ (lines 234--236) and its
window $[\texttt{0xb140},\texttt{0xb580})$ also contains
$\mptr{PCS\_FINAL\_MSM\_MPTR} = \texttt{0xb280}$ (line 201). This is documented
reuse, not a collision: the batch inversion completes at line 1826, before the
quotient VM writes the first selector accumulator and long before the PCS
final-MSM staging near line 4000. The disjointness is a lifetime argument
checked by the generator's memory arena, not by the deployed bytecode. The
contract's only on-chain memory-layout guards are two free-memory-pointer
checks --- \yul{if gt(mload(0x40), TRANSCRIPT\_MPTR)} at lines 674--677 in the
runtime and \yul{if gt(mload(0x40), 0x1000)} at lines 363--366 in the
constructor, both reverting with \revert{MemoryLayoutViolated} --- and both
protect the generated region from \emph{Solidity}, not the generated regions
from each other.

\subsection{\yul{ec\_pairing}: the terminal two-pair KZG check}
\label{sec:ec:pairing}

\begin{lstlisting}[style=solidity,caption={\texttt{Halo2Verifier.sol} lines 955--990: \yul{ec\_pairing} (comments elided)},label={lst:ec:pairing}]
            function ec_pairing(success, lhs_mptr, rhs_mptr) -> ret {
                ret := success
                if iszero(ret) { fail(ERR_PROOF_REJECTED) }
                let scratch := 0x1240
                mcopy(scratch,              lhs_mptr,                 0x80)
                mcopy(add(scratch, 0x80),   G2_BASE_MPTR,             0x100)
                mcopy(add(scratch, 0x180),  rhs_mptr,                 0x80)
                mcopy(add(scratch, 0x200),  NEG_S_G2_BASE_MPTR,       0x100)
                ret := staticcall(PAIRING_GAS_2PAIR, 0x0f, scratch, 0x0300, scratch, 0x20)
                ret := and(ret, eq(returndatasize(), 0x20))
                if iszero(ret) { fail(ERR_PRECOMPILE_FAILED) }
                ret := eq(mload(scratch), 1)
                if iszero(ret) { fail(ERR_PROOF_REJECTED) }
                ret := 1
            }
\end{lstlisting}

\subsubsection{Input staging}

The frame base is the literal \texttt{0x1240} --- a compile-time constant
inside the function body, not a parameter. In the generator it is
\yul{FINAL\_PAIRING\_SCRATCH\_START}, defined in
\texttt{src/lowering/layout/mod.rs}:121--122 (with the rationale at 111--120) as
\yul{PAIRING\_BATCH\_PTR} $+$ \yul{PAIRING\_BATCH\_HASH\_BYTES} $=$
\texttt{0x1000} $+$ \texttt{0x240}; the runtime pointer field
\yul{final\_pairing\_scratch\_mptr} (\texttt{memory.rs}:587, allocated at 769)
carries the resulting address. That base is deliberate: the accumulator block
immediately above hashes a \texttt{0x240}-byte frame based at \texttt{0x1000}
(\texttt{Halo2Verifier.sol}:4052--4066, the \yul{keccak256(batch\_ptr, 0x0240)}
at line 4066), and its result must survive into this call, so these two windows
are the one pair the phase-based lifetime model cannot separate. The generator
enforces address-level disjointness for them explicitly
(\texttt{src/lowering/layout/memory.rs}:1283--1302 and 1786--1803). Again, the
enforcement is generation-time only.

Because the frame base is hardcoded while \yul{lhs\_mptr} and \yul{rhs\_mptr}
are parameters, \yul{ec\_pairing} carries an unstated and unchecked
precondition: both argument points must lie outside
$[\texttt{0x1240}, \texttt{0x1540})$, or the second, third and fourth
\yul{mcopy} would overwrite a source the first three still need. It holds here
--- $\mptr{PAIRING\_LHS\_MPTR} = \texttt{0x7e80}$ and
$\mptr{PAIRING\_RHS\_MPTR} = \texttt{0x7f00}$ (lines 186--187), and the two
$\GT$ bases are at \texttt{0x3860}/\texttt{0x3960} --- but nothing in the
function tests it.

Four \yul{mcopy} instructions (Cancun) build the EIP-2537
\opcode{PAIRING\_CHECK} input. Each $\GO$ point is $\texttt{0x80}$ bytes
(two 64-byte coordinates, each a 16-byte zero pad followed by a 48-byte
big-endian $\Fp$ element); each $\GT$ point is $\texttt{0x100}$ bytes
($x_{c_0}, x_{c_1}, y_{c_0}, y_{c_1}$ in that order, same padding).

\begin{table}[htbp]
\centering
\begin{tabular}{@{}lllc@{}}
\toprule
Destination & Source & Bytes & Meaning \\
\midrule
$[\texttt{0x1240}, \texttt{0x12c0})$ & \yul{lhs\_mptr} & \texttt{0x80} & $\GO$ argument of pair 1 \\
$[\texttt{0x12c0}, \texttt{0x13c0})$ & \mptr{G2\_BASE\_MPTR} \texttt{0x3860} & \texttt{0x100} & $[1]_2$ \\
$[\texttt{0x13c0}, \texttt{0x1440})$ & \yul{rhs\_mptr} & \texttt{0x80} & $\GO$ argument of pair 2 \\
$[\texttt{0x1440}, \texttt{0x1540})$ & \mptr{NEG\_S\_G2\_BASE\_MPTR} \texttt{0x3960} & \texttt{0x100} & $-[s]_2$ \\
\midrule
\multicolumn{2}{@{}l}{call input} & \texttt{0x300} & $2 \times 384$ bytes \\
\multicolumn{2}{@{}l}{call output} & \texttt{0x20} & aliased onto \texttt{0x1240} \\
\bottomrule
\end{tabular}
\caption{\yul{ec\_pairing} staging window, \texttt{Halo2Verifier.sol} 969--980.}
\label{tab:ec:pairframe}
\end{table}

The window is contiguous and exactly \texttt{0x300} bytes, matching the
\texttt{0x0300} \textit{insize} argument at line 980; the two-pair size is also
what the constructor smoke probe exercises (line 565), so a chain whose
\texttt{0x0f} rejects this input length fails at deployment rather than at
proof time. The \yul{mcopy} opcode itself is probed at lines 374--376, so a
non-Cancun fork also fails at deployment. Output is aliased back onto the
frame base, which is safe for the reason given in \S\ref{sec:ec:batchinvert}.
The comment at lines 965--968 claims the \yul{mcopy} form saves roughly 500
gas per call against the previous \yul{mstore} chains. That is the code's own
figure, reported here unverified; it is at least consistent with the
$3 + 3\lceil \mathit{len}/32\rceil$ \opcode{MCOPY} schedule, which prices the
four copies at $15 + 27 + 15 + 27 = 84$ gas, against $2\cdot 4 + 2\cdot 8 = 24$
\yul{mstore}/\yul{mload} pairs ($144$ gas) plus the loop control and address
arithmetic the two eight-iteration $\GT$ loops used to carry. The saving is in
the eliminated loads and loop overhead, not in the stores themselves.

The two $\GT$ bases are read from the verifying-key payload, which is
\yul{extcodecopy}'d into $[\mptr{VK\_MPTR}, \dots)$ at line 1393 only after
\yul{extcodesize} and \yul{extcodehash} are checked against the pinned
\yul{EXPECTED\_VK\_LENGTH} and \yul{EXPECTED\_VK\_CODEHASH\_WORD} at lines
1387--1388, on pain of \revert{VkMismatch} at line 1389 (covered in the VK
section). \yul{ec\_pairing} performs no
validation of its own on those 16 words --- not on-curve, not subgroup, not
non-identity. That is the correct division of labour given the codehash pin,
and the precompile would in any case reject an off-curve or off-subgroup $\GT$
point.

\subsubsection{The equation checked}

EIP-2537 \opcode{PAIRING\_CHECK} on pairs $(P_1,Q_1),(P_2,Q_2)$ returns the
32-byte word $1$ iff $\prod_i e(P_i, Q_i) = 1$ in the target group $\Fp^{12}$, and $0$
otherwise; it fails the call outright if any point is off-curve or outside the
prime-order subgroup. So line 980 tests

\begin{equation}
  \pair{\yul{lhs\_mptr}}{[1]_2} \cdot \pair{\yul{rhs\_mptr}}{-[s]_2} \;=\; 1
  \iff
  \pair{\yul{lhs\_mptr}}{[1]_2} \;=\; \pair{\yul{rhs\_mptr}}{[s]_2}.
  \label{eq:ec:pairing}
\end{equation}

The sole call site is line 4114,
\yul{ec\_pairing(success, PAIRING\_RHS\_MPTR, PAIRING\_LHS\_MPTR)}, so
$\yul{lhs\_mptr} = \mptr{PAIRING\_RHS\_MPTR}$ and
$\yul{rhs\_mptr} = \mptr{PAIRING\_LHS\_MPTR}$. With the contents documented at
lines 4099--4111 --- $\mptr{PAIRING\_LHS\_MPTR} = \pi$ and
$\mptr{PAIRING\_RHS\_MPTR} = \mathsf{F} - v\,G + x_3\,\pi$, each already folded
with $\alpha$ times the corresponding IVC accumulator point --- \eqref{eq:ec:pairing}
instantiates to the KZG batch-opening identity

\begin{equation}
  \pair{\mathsf{F} - v\,G + x_3\,\pi}{[1]_2}
  \;=\;
  \pair{\pi}{[s]_2}.
  \label{eq:ec:kzg}
\end{equation}

\begin{finding}[EC-4]{\yul{lhs\_mptr}/\yul{rhs\_mptr} name the pairing slots, \mptr{PAIRING\_LHS\_MPTR}/\mptr{PAIRING\_RHS\_MPTR} name the dual-MSM sides}
\textbf{Severity: Informational.}

The two naming conventions are opposites, and the only call site
(\texttt{Halo2Verifier.sol}:4114) consequently passes its arguments crossed.
Swapping them would compute
$\pair{\pi}{[1]_2} = \pair{\mathsf{F}-vG+x_3\pi}{[s]_2}$, which is a different
and generally false equation --- so the crossing is required, not incidental.
The call site defends itself with a nine-line comment (lines 4099--4111)
deriving exactly this, which is the right mitigation for a hazard that cannot
be removed without renaming memory constants used across many sections. Flagged
only so that a reader of \yul{ec\_pairing} in isolation does not conclude the
call site is buggy.
\end{finding}

\subsubsection{Guards, exits and result strictness}

\begin{table}[htbp]
\centering
\begin{tabularx}{\textwidth}{@{}llX@{}}
\toprule
Line & Condition & Outcome \\
\midrule
962 & \yul{success} $=0$ on entry & \revert{ProofRejected} \\
981--982 & \opcode{STATICCALL} returned 0, \emph{or} \yul{returndatasize()} $\ne \texttt{0x20}$ & \revert{PrecompileFailed} \\
987--988 & result word $\ne 1$ & \revert{ProofRejected} \\
989 & otherwise & \yul{ret := 1} \\
\bottomrule
\end{tabularx}
\caption{\yul{ec\_pairing} exits. Every path except the last reverts.}
\label{tab:ec:pairexits}
\end{table}

Three properties of this table deserve to be called out as sound rather than
merely present.

\paragraph{No silent-failure path.} Unlike \yul{batch\_invert}, this helper
never returns a falsy value: the only normal return sets \yul{ret := 1}
unconditionally (line 989). The comment at lines 957--961 gives the reason ---
the terminal \yul{return(RETURN\_MPTR, 0x20)} in the epilogue (line 4128)
returns \texttt{true} without consulting \yul{success}, so a helper that
handed back $0$ would depend on the caller to notice. Converting the
$\yul{success}=0$ entry case into a revert at line 962 rather than an early
\yul{leave} closes that. The caller's own
\yul{if iszero(success) \{ fail(ERR\_PROOF\_REJECTED) \}} at line 4125 is
therefore dead code, which the comment at lines 4118--4124 acknowledges and
justifies as keeping acceptance a local property of the file. This is
defence-in-depth done correctly: the redundant check costs a few gas and
removes a whole class of refactoring hazard.

\paragraph{Strict comparison against 1.} Line 987 is
\yul{ret := eq(mload(scratch), 1)}, not \yul{and(ret, mload(scratch))}. The
comment at lines 983--986 states the difference: the bitwise form would accept
\emph{any} odd result word. EIP-2537 specifies the output as $0$ or $1$, so on
a conformant chain the two forms agree; on a non-conformant or short-returning
one they do not, and the strict form is the safe one. Combined with the
\yul{returndatasize()} equality at line 981, the accepted set is exactly the
single 32-byte word $\mathtt{0x00\ldots01}$.

\paragraph{Ordering.} \yul{mload(scratch)} at line 987 executes only after the
guard at line 982 has already reverted on a failed or short call, so the result
word is never read out of stale memory. Note that on a chain without EIP-2537
the stale memory in question would be the first word of the copied
$\yul{lhs\_mptr}$ point, i.e.\ the padded high half of $x$ --- which for a
degenerate identity point is $0$ and for a real point is essentially never $1$.
The verifier does not rely on that: line 982 rejects first.

\begin{finding}[EC-5]{\revert{PrecompileFailed} is reachable by gas starvation, so it is not conclusive evidence of a chain fault}
\textbf{Severity: Low.}

The MF-4 comments at \texttt{Halo2Verifier.sol}:828--835 and 974--979 define
\revert{PrecompileFailed} as "the chain could not run the precompile" and
direct incident responders at the node or fork. But \opcode{STATICCALL}
forwards $\min(\mathit{requested},\ \mathit{remaining} - \lfloor
\mathit{remaining}/64\rfloor)$: a caller who supplies a transaction gas limit
leaving fewer than $\tfrac{64}{63}\cdot 102900 \approx 104.5$k gas at line 980
causes the pairing precompile to run out of gas, \opcode{STATICCALL} to return
$0$, and line 982 to revert with \revert{PrecompileFailed} on a perfectly
healthy chain. The retained $\lfloor \mathit{remaining}/64 \rfloor \approx 1633$
gas is far more than the \yul{fail} helper's \yul{mstore}${}+{}$\yul{revert}
needs, so the spurious selector really does reach the caller.

The two modexp sites at lines 873 and 924 do \emph{not} share this property, and
the point is worth making precisely because it is the opposite of what one would
expect. Starving them requires
$\mathit{remaining} < \tfrac{64}{63}\cdot 4080 \approx 4145$ at the call, which
leaves the caller under $65$ gas --- nowhere near enough to run the $\approx 80$
remaining \yul{mulmod}s of lines 1828--1875 and reach the dispatch at
lines 1877--1882. Gas starvation there aborts the whole transaction with an
out-of-gas, not with \revert{PrecompileFailed}; only the pairing site can
manufacture the misleading selector. The distinction the selectors draw is
therefore
"precompile did not produce a well-formed answer" rather than
"the chain is broken"; a responder who treats a \revert{PrecompileFailed}
report as a fork incident without first checking the submitted gas limit will
chase the wrong cause --- the exact failure mode MF-4 set out to prevent. The
underlying design (exact-cost bounds, fail closed) is right; only the
interpretation attached to the selector is over-strong. Note also that the
severity is bounded: a starved caller cannot make the verifier \emph{accept}
anything, only mislabel its own refusal.
\end{finding}

\begin{finding}[EC-6]{\yul{ec\_pairing} accepts an all-identity input}
\textbf{Severity: Observation.}

$e(\id, Q) = 1$ for every $Q$, so \eqref{eq:ec:pairing} is satisfied whenever
both $\GO$ arguments are the point at infinity. The contract demonstrates this
itself: the constructor smoke probe at
\texttt{Halo2Verifier.sol}:561--567 calls \texttt{0x0f} on a zeroed
\texttt{0x300}-byte frame --- two (identity, identity) pairs --- and asserts
the result is $1$. \yul{ec\_pairing} adds no non-degeneracy test on
\yul{lhs\_mptr} or \yul{rhs\_mptr}. This is not by itself a defect: the
degenerate case corresponds to $\pi = \id$ and
$\mathsf{F} - vG + x_3\pi = \id$, and whether a prover can force that is a
property of how those two points are built (lines 4030--4097 and the PCS
sections, specified elsewhere), not of this wrapper. Recorded so that the
absence of the check is an explicit, located fact rather than an omission a
reader has to notice. The Fiat-Shamir binding of the accumulator folding
scalar $\alpha$ to \mptr{VK\_DIGEST\_MPTR} and to the fully materialized
pairing inputs (the domain tag, digest and four \yul{mcopy}'d $\GO$ points at
lines 4057--4062, hashed at line 4066), together with the
\yul{if iszero(acc\_pair\_alpha) \{ acc\_pair\_alpha := 1 \}} floor at line
4067, is the mechanism that is supposed to make a coordinated cancellation
infeasible.
\end{finding}

\paragraph{Curve and subgroup validation.} \yul{ec\_pairing} runs no explicit
on-curve or subgroup check on its $\GO$ inputs; it relies on
\opcode{PAIRING\_CHECK} to do so, which EIP-2537 guarantees. In this artifact
the two arguments are the outputs of the EIP-2537 \texttt{G1ADD}
(\texttt{0x0b}) calls at lines 4080 and 4094, each fed by a \texttt{G1MSM}
(\texttt{0x0c}) at lines 4075 and 4089. That does \emph{not} by itself
establish subgroup membership: EIP-2537 \texttt{G1ADD} validates only that its
inputs are on the curve and explicitly performs \emph{no} subgroup check --- it
is \texttt{G1MSM} and \opcode{PAIRING\_CHECK} that do. The property holds by
closure instead: the subgroup is closed under addition, so a \texttt{G1ADD} of
two subgroup elements is a subgroup element, and the argument bottoms out in
whatever validates the \emph{other} \texttt{G1ADD} operands
(\mptr{PAIRING\_RHS\_MPTR} at line 4078 and \mptr{PAIRING\_LHS\_MPTR} at
line 4092), which the PCS and accumulator sections own. Granting that, the
pairing precompile's validation cannot fire, and the only way for the call at
line 980 to fail is a chain fault or the gas shortfall of finding EC-5. Were
the closure argument to break anywhere, an invalid proof-derived point would
surface as \revert{PrecompileFailed} --- a proof fault wearing a chain fault's
selector, and a second instance of the EC-5 pattern.

\begin{finding}[EC-7]{The sole call site pre-empts \yul{ec\_pairing}'s entry guard with a \emph{different} selector}
\textbf{Severity: Informational} (cross-boundary; line 4113 is specified elsewhere).

\yul{ec\_pairing}'s own entry guard (line 962) converts $\yul{success}=0$ into
\revert{ProofRejected}. But the only call site is preceded, one line earlier, by
\yul{if iszero(success) \{ fail(ERR\_PRECOMPILE\_FAILED) \}}
(\texttt{Halo2Verifier.sol}:4113), so line 962 is unreachable from it and every
incoming $\yul{success}=0$ is reported as \revert{PrecompileFailed} instead.
The two guards therefore disagree about which incident an already-failed
\yul{success} represents. Which reading is right depends on what can clear
\yul{success} before line 4113 --- the accumulator pairing-batch
\texttt{G1MSM}/\texttt{G1ADD} calls at lines 4075--4095, whose operands are
partly derived from the public IVC accumulator inputs; that block is specified
elsewhere. Recorded here because it means the exit table above describes
\yul{ec\_pairing} as a routine, not the behaviour observable from this
contract: line 962's \revert{ProofRejected} never fires in the deployed
artifact.
\end{finding}

\subsection{Summary of the region}
\label{sec:ec:summary}

Lines 811--991 are, on the whole, careful code. The properties that matter most
for soundness are all present and correctly ordered: canonicality is checked on
every element of the batch and not just some; zero denominators are rejected
through an argument (no zero divisors in $\Fr$) that is exactly equivalent to
the per-element test; the modexp result is never read when the call failed, and
the batch is left untouched in that case; every precompile result is gated on
\yul{returndatasize()} equality before it is used; the pairing result is
compared strictly against $1$; and \yul{ec\_pairing} has no return path that
could let an unverified proof through. Gas is bounded by exact scheduled
constants at every call site, which caps the damage of a failing precompile at
one frame's cost and is the correct trade against the liveness risk of a future
repricing --- a risk the constructor probes convert from a silent brick into a
deployment failure.

The two things a reviewer should carry forward are structural rather than
local: the batch-inversion scratch window is sized with zero slack against its
own input array and the guarantee lives entirely in the generator
(finding EC-1), and the \revert{PrecompileFailed}/\revert{ProofRejected} split
is looser than its comments claim --- gas starvation at the pairing site,
a codegen fault in the batch range, and the call site's own pre-emption of
\yul{ec\_pairing}'s entry guard all land in the wrong bucket (findings EC-5,
EC-3, EC-7). None of the three is exploitable in the deployed fixture, and none
can cause an invalid proof to be accepted: every one of them mislabels a
refusal rather than weakening it.


%% ==== 07-accumulator.tex ================================================
\section{IVC public-accumulator decoding and validation}
\label{sec:acc:main}

This section specifies \texttt{Halo2Verifier.sol} lines 992--1351: the six Yul
helpers that decode the recursive KZG accumulator carried in the public inputs,
and the driver \yul{validate\_public\_accumulator} that turns them into two
validated \GO{} points at \mptr{ACC\_LHS\_MPTR} and \mptr{ACC\_RHS\_MPTR}. The
call site (lines 1446--1453) is quoted here because it fixes the failure
taxonomy of the whole region; everything after it -- transcript absorption of
the same instance words (line 1518) and the randomized pairing batch (lines
4050--4097) -- is covered elsewhere and is only cited. The remainder of the
assigned range, lines 1352--1360, contains no accumulator logic: it closes the
helper block, creates the two variables this section's call site consumes
(\yul{let r := FR\_MODULUS} and \yul{let success := true}, lines 1355--1356),
and opens the VK-loading phase at line 1360, which is covered elsewhere.

\subsection{Position in the verifier and failure model}
\label{sec:acc:position}

The accumulator phase is marked \yul{// \_\_phase:accumulator\_msm} and runs
\emph{before} the transcript. It is the first block that touches attacker
controlled calldata after the VK/ABI shape checks of lines 1400--1430.

\begin{lstlisting}[style=solidity,caption={\texttt{Halo2Verifier.sol} lines 1446--1453: the call site and the two-way failure taxonomy},label={lst:acc:callsite}]
            let acc_precompile_failed := 0
            success, acc_precompile_failed := validate_public_accumulator(success, r)
            if iszero(success) {
                // MF-4: a G1MSM that could not run at all is a chain fault,
                // not a malformed accumulator point.
                if acc_precompile_failed { fail(ERR_PRECOMPILE_FAILED) }
                fail(ERR_BAD_POINT_ENCODING)
            }
\end{lstlisting}

Two properties follow immediately and are relied on throughout this section.

\begin{enumerate}
  \item \textbf{Fail closed, fail early.} \yul{validate\_public\_accumulator}
        never reverts internally; it returns a boolean \yul{out} that the
        caller converts into \revert{BadPointEncoding} (selector
        \texttt{0xf27905ec}, line 323) or \revert{PrecompileFailed} (selector
        \texttt{0x84e81692}, line 324). No path returns \yul{out}$=1$ for a
        point that failed a check. \yul{out} is assigned exactly five times:
        \yul{out := success} (line 1247, and \yul{success} is pinned to $1$
        by the \revert{BadCalldataShape} guard at line 1430), the two
        monotone folds \yul{out := and(out, lhs\_ok)} and
        \yul{out := and(out, rhs\_ok)} (lines 1266, 1302), and the two raw
        \opcode{STATICCALL} results (lines 1288, 1337). Only the last two can
        raise \yul{out}, and each sits inside an \yul{if out} block, so its
        prior value was already $1$; lines 1289 and 1345 immediately fold the
        return-size check back on top with \yul{and}.
  \item \textbf{Bounded rejection cost.} A malformed accumulator is rejected
        after at most two \GO{}MSM staticcalls of \texttt{12000} gas each and
        no transcript, quotient, PCS or pairing work. This ordering is a
        deliberate DoS-cost choice and it also removes any question of an
        invalid accumulator influencing Fiat--Shamir challenge derivation.
\end{enumerate}

\subsection{Codec parameters and calldata layout}
\label{sec:acc:params}

All codec parameters are literals rendered into
\yul{validate\_public\_accumulator} (lines 1248--1257) and are additionally
cross-checked against the loaded verifying key at lines 1402--1405, so a linked
\texttt{Halo2VerifyingKey} whose accumulator shape differs is rejected with
\revert{VkMismatch} before any decoding happens.

\begin{table}[htbp]
\centering
\begin{tabular}{@{}llll@{}}
\toprule
Symbol & Yul name & Value & VK cross-check \\
\midrule
$b$        & \yul{bits}            & $56$        & \mptr{NUM\_ACC\_LIMB\_BITS\_MPTR} $=56$ (line 1405) \\
$n$        & \yul{n}               & $7$         & \mptr{NUM\_ACC\_LIMBS\_MPTR} $=7$ (line 1404) \\
$B$        & \yul{limb\_base}      & $2^{56}$    & \yul{shl(bits, 1)} (line 1252) \\
$\lambda$  & \yul{limbs\_per\_word}& $4$         & literal (line 1253) \\
$W$        & \yul{coord\_words}    & $\lceil n/\lambda\rceil = 2$ & computed (line 1254) \\
--         & \yul{acc\_instance\_ptr} & \texttt{0x2044} & \mptr{ACC\_OFFSET\_MPTR} $=11$ (line 1403) \\
--         & \mptr{HAS\_ACCUMULATOR\_MPTR} & $1$ & line 1402 \\
\bottomrule
\end{tabular}
\caption{Accumulator codec parameters as rendered in this fixture.}
\label{tab:acc:params}
\end{table}

With \mptr{INSTANCE\_CPTR} $=$ \texttt{0x1ee4} (line 103) and the rendered
offset \texttt{0x0160} $=11\cdot\texttt{0x20}$, the accumulator occupies
instance words $11..18$, i.e. calldata \texttt{0x2044}--\texttt{0x2143}. The
generated ABI-shape check at lines 1419--1430 pins
\yul{calldataload(NUM\_INSTANCE\_CPTR)} $=19$ and
\yul{calldatasize()} $=$ \mptr{INSTANCE\_CPTR} $+$ \texttt{0x0260} $=$
\texttt{0x2144}, so the eight accumulator words are exactly the tail of the
instance array and every \yul{calldataload} performed by this section is
in-bounds. This matters: \opcode{CALLDATALOAD} past the end returns zero rather
than reverting, and a short instance array would otherwise decode as the
encoded value $E=0$ for every coordinate.

\begin{table}[htbp]
\centering
\begin{tabular}{@{}lllll@{}}
\toprule
Instance idx & Calldata offset & Point & Coordinate & Packed word \\
\midrule
11 & \texttt{0x2044} & LHS & $x$ & $w_0$ (limbs $0..3$) \\
12 & \texttt{0x2064} & LHS & $x$ & $w_1$ (limbs $4..6$) \\
13 & \texttt{0x2084} & LHS & $y$ & $w_0$ \\
14 & \texttt{0x20a4} & LHS & $y$ & $w_1$ \\
15 & \texttt{0x20c4} & RHS & $x$ & $w_0$ \\
16 & \texttt{0x20e4} & RHS & $x$ & $w_1$ \\
17 & \texttt{0x2104} & RHS & $y$ & $w_0$ \\
18 & \texttt{0x2124} & RHS & $y$ & $w_1$ \\
\bottomrule
\end{tabular}
\caption{Accumulator public-input layout for this render
(\yul{point\_pair}: two carried \GO{} points, implicit unit scalars, no
fixed-base tail).}
\label{tab:acc:layout}
\end{table}

The base-field sentinels used below are the contract constants of lines
336--342. Writing $p$ for the BLS12-381 base modulus,
\begin{align}
  p - 1 &= \mathtt{BLS\_P\_HI}\cdot 2^{256} + \mathtt{BLS\_P\_MINUS\_ONE\_LO},\\
  \mathtt{BLS\_P\_MINUS\_ONE\_PACKED\_0} &= (p-1) \bmod 2^{224}, \label{eq:acc:packed0}\\
  \mathtt{BLS\_P\_MINUS\_ONE\_PACKED\_1} &= \lfloor (p-1) / 2^{224} \rfloor, \label{eq:acc:packed1}\\
  \mathtt{BLS\_P\_MINUS\_ONE\_PACKED\_0\_WITH\_ID\_FLAG} &= \mathtt{BLS\_P\_MINUS\_ONE\_PACKED\_0} + 2^{56}. \label{eq:acc:flag}
\end{align}
All four identities were verified numerically against the declared constants.
Equation~\eqref{eq:acc:flag} is exact -- adding one raw radix base to limb $1$
of $p-1$ does not carry, because limb $1$ of $p-1$ is
\texttt{0xfffeb153ffffb9} $\neq 2^{56}-1$; the only hex digits that change are
\texttt{b9} $\to$ \texttt{ba}.

\subsection{The shifted limb codec: \yul{load\_acc\_coord\_shifted}}
\label{sec:acc:shifted}

\subsubsection{Exact packing formula}

For one coordinate at calldata offset $s$ let
$w_0 = \mathrm{calldataload}(s)$ and $w_1 = \mathrm{calldataload}(s+\texttt{0x20})$.
The decoder extracts $n=7$ limbs
\begin{equation}
  \ell_j \;=\; \Big\lfloor \frac{w_{\lfloor j/4 \rfloor} - \delta\,[\,\lfloor j/4\rfloor = 0\,]}{2^{56\,(j \bmod 4)}} \Big\rfloor \bmod 2^{56},
  \qquad j = 0,\dots,6,
  \label{eq:acc:limb}
\end{equation}
where $\delta \in \{0, 2^{56}\}$ is the \yul{first\_adjust} argument, and
reassembles
\begin{equation}
  E \;=\; \sum_{j=0}^{6} \ell_j\, 2^{56 j}, \qquad
  \mathit{lo} = E \bmod 2^{256}, \quad \mathit{hi} = \lfloor E / 2^{256}\rfloor .
  \label{eq:acc:E}
\end{equation}
Concretely, for $n=7$, $b=56$, $\lambda=4$ the limb-to-word map is
\[
  \underbrace{\ell_0,\ell_1,\ell_2,\ell_3}_{\text{word }w_0
  \text{ at bit offsets }0,56,112,168}
  \qquad
  \underbrace{\ell_4,\ell_5,\ell_6}_{\text{word }w_1
  \text{ at bit offsets }0,56,112}
\]
and, \emph{provided the over-packing guard of \S\ref{sec:acc:packing} holds},
equations~\eqref{eq:acc:limb}--\eqref{eq:acc:E} collapse to the much simpler
\begin{equation}
  E \;=\; (w_0 - \delta) \;+\; 2^{224}\, w_1 .
  \label{eq:acc:Esimple}
\end{equation}
Limb $4$ is the only one that straddles the 256-bit split: its shift is
$224 < 256$, but $224 + 56 = 280 > 256$, so its low $32$ bits land in
$\mathit{lo}$ at bit $224$ and its top $24$ bits land in $\mathit{hi}$ at bit
$0$. Limbs $5$ and $6$ land wholly in $\mathit{hi}$ at bits $24$ and $80$, so
$\mathit{hi} < 2^{136}$ before any range check.

\begin{lstlisting}[style=solidity,caption={\texttt{Halo2Verifier.sol} lines 1010--1049: \yul{load\_acc\_coord\_shifted}},label={lst:acc:shifted}]
            function load_acc_coord_shifted(src, bits, n, base, limbs_per_word, first_adjust) -> hi, lo {
                // Mask for one radix limb, e.g. 2^56 - 1 for the current
                // BLS12-381 self-emulation parameters.
                let mask := sub(base, 1)
                for { let i := 0 } lt(i, n) { i := add(i, 1) } {
                    // Limb words are little-endian packed inside each Fr
                    // public input. `first_adjust` removes the identity flag
                    // base from the first x word when present.
                    let packed := calldataload(add(src, mul(div(i, limbs_per_word), 0x20)))
                    // `and` here is bitwise, so it must not be fed the raw
                    // `first_adjust` (a radix base, i.e. a high power of two):
                    // `iszero(...)` is 0 or 1 and shares no bit with it, which
                    // would make the guard false for every call. Subtracting is
                    // already a no-op when `first_adjust` is zero, so gate on
                    // the word index alone.
                    if iszero(div(i, limbs_per_word)) {
                        packed := sub(packed, first_adjust)
                    }
                    // Select limb i from its packed field word. The mod/div
                    // pair maps a limb index to an intra-word limb slot and
                    // the calldata word containing it.
                    let limb := and(shr(mul(mod(i, limbs_per_word), bits), packed), mask)

                    let shift := mul(i, bits)
                    // Split the reconstructed 384-bit coordinate into the
                    // EIP-2537 high/low words expected by the precompiles.
                    if lt(shift, 256) {
                        lo := add(lo, shl(shift, limb))
                        if gt(add(shift, bits), 256) {
                            // A limb can straddle the 256-bit low/high split.
                            // Move the overflow bits into hi.
                            hi := add(hi, shr(sub(256, shift), limb))
                        }
                    }
                    if iszero(lt(shift, 256)) {
                        // Once shift >= 256 the whole limb belongs to hi.
                        hi := add(hi, shl(sub(shift, 256), limb))
                    }
                }
            }
\end{lstlisting}

\subsubsection{The \yul{first\_adjust} identity-flag bit}

\yul{first\_adjust} is subtracted from \emph{word $0$ only}, and is re-applied
on every one of the four iterations that read word $0$ (the loop re-loads
\yul{packed} from calldata each iteration, so the subtraction is not
cumulative). It is $0$ or exactly one radix base $B = 2^{56}$; i.e. the
``identity flag'' is $+1$ added to limb $\ell_1$ of the $x$ coordinate.

Two implementation details are load-bearing:

\begin{itemize}
  \item \textbf{The guard is on the word index, not on \yul{first\_adjust}.}
        Lines 1019--1027 spell out why: \yul{and} in Yul is bitwise, so
        \yul{and(iszero(div(i, limbs\_per\_word)), first\_adjust)} would be
        \emph{always zero} -- \yul{iszero} yields $0$ or $1$, which shares no
        set bit with the high power of two $2^{56}$. Gating on
        \yul{iszero(div(i, limbs\_per\_word))} alone is correct because
        subtracting $0$ is a no-op. The generator lints for the buggy form
        (\texttt{src/lowering/tests.rs:5000--5012}).
  \item \textbf{The subtraction cannot wrap.} \yul{sub} is modular over
        $2^{256}$, so $\delta = 2^{56} > w_0$ would produce garbage limbs.
        Every call site guarantees $w_0 \ge B$ whenever $\delta = B$: the probe
        at line 1108 is guarded by \yul{iszero(lt(calldataload(src), base))},
        and the second decode at line 1121 passes \yul{mul(is\_id, base)},
        which is nonzero only if that same probe fired.
\end{itemize}

\subsection{The $p-1$ sentinel: \yul{is\_bls\_p\_minus\_one} and \yul{is\_acc\_encoded\_identity}}
\label{sec:acc:sentinels}

The circuit-side codec stores $(\mathit{coord} - 1) \bmod p$, so a coordinate
of $0$ is encoded as $p-1$. Both sentinel predicates below are pure comparisons
against compile-time constants.

\begin{lstlisting}[style=solidity,caption={\texttt{Halo2Verifier.sol} lines 1051--1076: the two sentinel predicates},label={lst:acc:sentinels}]
            // The shifted coordinate codec represents zero as p-1 before the
            // final +1 below, so keep this sentinel explicit.
            function is_bls_p_minus_one(hi, lo) -> yes {
                yes := and(eq(hi, BLS_P_HI), eq(lo, BLS_P_MINUS_ONE_LO))
            }

            // Canonical encoded accumulator identity:
            //   x = p-1 plus the identity flag in the first packed word,
            //   y = p-1 with no identity flag.
            // It decodes to the EIP-2537 point-at-infinity slot (all zeros).
            //
            // This fast path is deliberately stricter than "decodes to zero":
            // the point at infinity has exactly one accepted public-input
            // encoding. Non-canonical zero-like encodings are rejected later.
            function is_acc_encoded_identity(src) -> yes {
                yes := and(
                    and(
                        eq(calldataload(src), BLS_P_MINUS_ONE_PACKED_0_WITH_ID_FLAG),
                        eq(calldataload(add(src, 0x20)), BLS_P_MINUS_ONE_PACKED_1)
                    ),
                    and(
                        eq(calldataload(add(src, 0x40)), BLS_P_MINUS_ONE_PACKED_0),
                        eq(calldataload(add(src, 0x60)), BLS_P_MINUS_ONE_PACKED_1)
                    )
                )
            }
\end{lstlisting}

Writing $S = (\mathtt{PACKED\_0}, \mathtt{PACKED\_1})$ for the unflagged $p-1$
word pair and $S^{\ast} = (\mathtt{PACKED\_0\_WITH\_ID\_FLAG},
\mathtt{PACKED\_1})$ for the flagged one, \yul{is\_acc\_encoded\_identity}
returns $1$ iff the four calldata words of the point are exactly
$(S^{\ast}, S)$. Note that the nesting of \yul{and} is safe here because every
operand is a $0/1$ \yul{eq} result -- this is the benign form of the bitwise
\yul{and} trap the file warns about at lines 1019--1024.

\subsection{Over-packing rejection: \yul{check\_acc\_coord\_packing}}
\label{sec:acc:packing}

\begin{lstlisting}[style=solidity,caption={\texttt{Halo2Verifier.sol} lines 1081--1102: \yul{check\_acc\_coord\_packing}},label={lst:acc:packing}]
            function check_acc_coord_packing(src, bits, n, limbs_per_word) -> ok {
                ok := 1
                // Number of packed native-field public-input words occupied by
                // one coordinate.
                let coord_words := div(add(n, sub(limbs_per_word, 1)), limbs_per_word)
                for { let word_idx := 0 } lt(word_idx, coord_words) { word_idx := add(word_idx, 1) } {
                    // The final word may contain fewer than limbs_per_word
                    // limbs. Any unused high bits must be zero, otherwise the
                    // same coordinate would have multiple calldata encodings.
                    let remaining := sub(n, mul(word_idx, limbs_per_word))
                    let limbs_in_word := limbs_per_word
                    if lt(remaining, limbs_per_word) {
                        limbs_in_word := remaining
                    }
                    let used_bits := mul(limbs_in_word, bits)
                    if lt(used_bits, 256) {
                        // shl(used_bits, 1) == 2^used_bits. The packed word
                        // must be strictly less than that bound.
                        ok := and(ok, lt(calldataload(add(src, mul(word_idx, 0x20))), shl(used_bits, 1)))
                    }
                }
            }
\end{lstlisting}

Instantiated at $n=7$, $b=56$, $\lambda=4$ the loop runs $W=2$ times and checks
\emph{exactly} the following, per coordinate:

\begin{table}[htbp]
\centering
\begin{tabularx}{\textwidth}{@{}lllX@{}}
\toprule
\yul{word\_idx} & \yul{limbs\_in\_word} & \yul{used\_bits} & Guard \\
\midrule
$0$ & $\min(7,4)=4$ & $224$ & $w_0 < 2^{224}$, i.e. bits $224..255$ of $w_0$ are all zero (top $32$ bits) \\
$1$ & $7-4=3$ & $168$ & $w_1 < 2^{168}$, i.e. bits $168..255$ of $w_1$ are all zero (top $88$ bits) \\
\bottomrule
\end{tabularx}
\caption{Bits constrained by \yul{check\_acc\_coord\_packing} for this render.
Bits $0..223$ of $w_0$ and $0..167$ of $w_1$ are consumed by limbs
$\ell_0..\ell_3$ and $\ell_4..\ell_6$ respectively; the guard forces every
remaining bit of both words to zero.}
\label{tab:acc:packing}
\end{table}

The \yul{if lt(used\_bits, 256)} guard exists because a word that is exactly
full ($\lambda \cdot b = 256$) has no unused bits and \yul{shl(256, 1)} would
be $0$, turning the guard into an unsatisfiable \yul{lt(w, 0)}. For this
render neither word is full, so both guards are live.

\begin{finding}[ACC-CLEAN-1]{Every bit of every accumulator instance word is constrained}
\textbf{Severity: Informational (clean).} The extraction in
equation~\eqref{eq:acc:limb} silently \emph{discards} bits $\ge 224$ of $w_0$
and bits $\ge 168$ of $w_1$: only four limbs are read from word $0$ and three
from word $1$. Without \yul{check\_acc\_coord\_packing} the map
$(w_0,w_1)\mapsto E$ would be massively non-injective, and an attacker could
mint $2^{32+88}$ distinct calldata encodings of any accumulator point --
each producing a different transcript absorption at line 1521 while decoding to
the same \GO{} point. The guard closes exactly the discarded bit ranges
(Table~\ref{tab:acc:packing}), so under it equation~\eqref{eq:acc:Esimple}
holds and $(w_0,w_1)\mapsto E$ is injective. Every limb of every word is
therefore constrained, and no limb slot is left unchecked.

A useful side effect: $w_0 < 2^{224}$ and $w_1 < 2^{168}$ imply $w_i < r$, so
the accumulator instance words are automatically canonical \Fr{} elements even
though the generic instance canonicality check (\yul{lt(inst\_be, r)}, line
1518) runs \emph{after} this phase. There is no ordering hazard.
\end{finding}

\subsection{Decoding one coordinate: \yul{load\_acc\_coord}}
\label{sec:acc:coord}

\begin{lstlisting}[style=solidity,caption={\texttt{Halo2Verifier.sol} lines 1106--1147: \yul{load\_acc\_coord}},label={lst:acc:coord}]
            function load_acc_coord(src, allow_id, bits, n, base, limbs_per_word) -> ok, hi, lo, is_id {
                ok := check_acc_coord_packing(src, bits, n, limbs_per_word)
                if and(allow_id, iszero(lt(calldataload(src), base))) {
                    // Probe the x identity flag by removing one radix base and
                    // checking whether the adjusted coordinate is p-1.
                    //
                    // `calldataload(src) >= base` is a cheap prefilter: only x
                    // can carry this flag, and adding one radix base must make
                    // the first packed word at least base.
                    let adj_hi, adj_lo := load_acc_coord_shifted(src, bits, n, base, limbs_per_word, base)
                    is_id := is_bls_p_minus_one(adj_hi, adj_lo)
                }

                // Decode again with the identity adjustment applied only when
                // the canonical identity flag was actually detected.
                hi, lo := load_acc_coord_shifted(src, bits, n, base, limbs_per_word, mul(is_id, base))
                ok := and(
                    ok,
                    // Coordinate must be in the BLS12-381 base field, i.e.
                    // <= p - 1 in split hi/lo form.
                    or(lt(hi, BLS_P_HI), and(eq(hi, BLS_P_HI), iszero(gt(lo, BLS_P_MINUS_ONE_LO))))
                )

                let was_p_minus_one := is_bls_p_minus_one(hi, lo)
                if was_p_minus_one {
                    // Shifted encoding maps p-1 back to zero.
                    hi := 0
                    lo := 0
                }
                if iszero(was_p_minus_one) {
                    // All other coordinates are encoded as coord - 1, so add
                    // one back with carry into the high word.
                    let next_lo := add(lo, 1)
                    hi := add(hi, lt(next_lo, lo))
                    lo := next_lo
                }

                // EIP-2537 pads each 48-byte Fp coordinate to 64 bytes,
                // so the high word must fit in its low 128 bits.
                // This also catches impossible reconstructions above 384 bits.
                ok := and(ok, lt(hi, shl(128, 1)))
            }
\end{lstlisting}

The routine performs, in order:

\begin{enumerate}
  \item \textbf{Packing check} (line 1107): $\mathit{ok} \gets$
        \yul{check\_acc\_coord\_packing}.
  \item \textbf{Identity probe} (lines 1108--1117), only if
        \yul{allow\_id} $=1$ \emph{and} $w_0 \ge B$. Decodes with
        $\delta = B$ and sets \yul{is\_id} $\gets [\,E - B = p-1\,]$.
        This costs a second full seven-iteration decode for \emph{every} $x$
        coordinate whose $w_0 \ge 2^{56}$, i.e. essentially always.
  \item \textbf{Authoritative decode} (line 1121) with
        $\delta = \mathtt{is\_id}\cdot B$.
  \item \textbf{Base-field range check} (lines 1122--1127):
        $\mathit{ok} \mathrel{\&}= (\mathit{hi} < \mathtt{BLS\_P\_HI})
        \lor (\mathit{hi} = \mathtt{BLS\_P\_HI} \land
        \mathit{lo} \le \mathtt{BLS\_P\_MINUS\_ONE\_LO})$, i.e.
        $E \le p-1$. Note this is applied to the \emph{encoded} value, before
        the shift is undone. The \yul{or} and the inner \yul{and} at line
        1126 are bitwise, but every operand is a $0/1$ comparison result, so
        the expression is the intended logical one -- the benign form of the
        trap the file warns about at lines 1019--1024.
  \item \textbf{Un-shift} (lines 1129--1141):
        \begin{equation}
          \mathit{coord} =
          \begin{cases}
            0, & E = p-1,\\
            E + 1, & 0 \le E \le p-2,
          \end{cases}
          \qquad\text{i.e.}\qquad \mathit{coord} = (E+1) \bmod p .
          \label{eq:acc:unshift}
        \end{equation}
        The $E+1$ arm is a 512-bit increment: \yul{next\_lo := add(lo, 1)},
        \yul{hi := add(hi, lt(next\_lo, lo))} detects the wrap correctly
        because \yul{lt(next\_lo, lo)} is $1$ exactly when
        $\mathit{lo} = 2^{256}-1$.
  \item \textbf{EIP-2537 padding check} (line 1146):
        $\mathit{ok} \mathrel{\&}= (\mathit{hi} < 2^{128})$.
\end{enumerate}

Because equation~\eqref{eq:acc:unshift} is a bijection of $[0,p-1]$ onto
itself, and step 4 confines $E$ to $[0,p-1]$, the returned coordinate is always
a valid \Fp{} element on the accepting path -- so the precompile never sees an
out-of-range coordinate originating here.

\begin{finding}[ACC-1]{The final \yul{lt(hi, 2\^{}128)} guard is unreachable-as-a-rejection}
\textbf{Severity: Observation.} Line 1146 is implied by the base-field check at
lines 1122--1127: if $E \le p-1$ then after the increment
$\mathit{hi} \le \mathtt{BLS\_P\_HI} < 2^{125}$, and if $E > p-1$ then
$\mathit{ok}$ is already $0$ and stays $0$ under \yul{and}. The guard therefore
never changes the verdict. It costs a handful of gas per coordinate (four
coordinates per proof) and is worth keeping as defence in depth against a
future codec change that widens $n\cdot b$ past $384$ bits -- the file's own
comment says as much. No action required; recorded so a reader does not mistake
it for the load-bearing range check, which is line 1126.
\end{finding}

\begin{finding}[ACC-2]{Double decode of every \texorpdfstring{$x$}{x} coordinate}
\textbf{Severity: Observation (gas).} Lines 1108--1121 run
\yul{load\_acc\_coord\_shifted} twice for each $x$ coordinate whenever
$w_0 \ge 2^{56}$, which holds for all but a negligible set of honest inputs.
The probe result is used only as a one-bit predicate. It is \emph{not} the
case that line 1157 has already settled the question: \yul{is\_acc\_encoded\_identity}
returning $0$ still leaves \yul{x\_is\_id} $=1$ reachable (Finding~ACC-4).
What is true is that the probe cannot change the \emph{verdict}. Under the
packing guard the probe fires iff $x$ words $= S^{\ast}$ exactly, and on such
an input the canonical fast path has already been declined, so $y$ words
$\neq S$. With the probe, \yul{is\_id} $=1$ and line 1201 forces \yul{ok}
$\gets 0$; without it, $\delta = 0$ gives $E_x = (p-1) + 2^{56} > p-1$ and
line 1126 forces \yul{ok} $\gets 0$. The accepted set is identical either
way, so the second decode buys nothing but costs roughly $2 \times 7$ extra
\opcode{CALLDATALOAD}/shift sequences per proof. Keeping the probe makes \yul{load\_acc\_coord} correct as
a standalone unit, which is the stated design intent (lines 1104--1105). No
soundness impact.
\end{finding}

\subsection{Decoding one point: \yul{load\_acc\_point}}
\label{sec:acc:point}

\begin{lstlisting}[style=solidity,caption={\texttt{Halo2Verifier.sol} lines 1152--1187: \yul{load\_acc\_point}, canonical-identity fast path and coordinate decode},label={lst:acc:point-a}]
            function load_acc_point(dst, src, bits, n, base) -> ok, is_id {
                // Prefer the canonical all-coordinate identity encoding before
                // attempting coordinate-level shifted decoding. This accepts
                // the point at infinity only in the exact form generated by the
                // circuit's public-input codec.
                is_id := is_acc_encoded_identity(src)
                if is_id {
                    ok := 1
                    // EIP-2537 encodes G1 identity as four zero words:
                    // x_hi = x_lo = y_hi = y_lo = 0.
                    mstore(dst, 0)
                    mstore(add(dst, 0x20), 0)
                    mstore(add(dst, 0x40), 0)
                    mstore(add(dst, 0x60), 0)
                }
                if iszero(is_id) {
                    // x occupies coord_words packed public-input words; y
                    // starts immediately after x.
                    let limbs_per_word := 4
                    let coord_words := div(add(n, sub(limbs_per_word, 1)), limbs_per_word)
                    // Only x may carry the identity flag. y must decode as a
                    // normal shifted coordinate.
                    let x_ok, x_hi, x_lo, x_is_id := load_acc_coord(src, 1, bits, n, base, limbs_per_word)
                    let y_ok, y_hi, y_lo, y_id := load_acc_coord(
                        add(src, mul(coord_words, 0x20)),
                        0,
                        bits,
                        n,
                        base,
                        limbs_per_word
                    )
                    // y_id is always zero because allow_id was false, but the
                    // tuple shape is shared with x decoding.
                    pop(y_id)
                    ok := and(x_ok, y_ok)
                    is_id := x_is_id
\end{lstlisting}

\begin{lstlisting}[style=solidity,caption={\texttt{Halo2Verifier.sol} lines 1189--1223: \yul{load\_acc\_point}, the two zero-point guards (the audit comment at lines 1193--1200 is elided)},label={lst:acc:point-b}]
                    if is_id {
                        // If x carried the identity flag, both decoded
                        // coordinates must be zero after shifting. Any other y
                        // value would be a malformed infinity encoding.
                        // ...
                        ok := and(ok, iszero(or(or(x_hi, x_lo), or(y_hi, y_lo))))
                        mstore(dst, 0)
                        mstore(add(dst, 0x20), 0)
                        mstore(add(dst, 0x40), 0)
                        mstore(add(dst, 0x60), 0)
                    }
                    if iszero(is_id) {
                        // The coordinate codec maps encoded p-1 to decoded
                        // zero. EIP-2537 reserves affine (0,0) for the point
                        // at infinity, so a decoded infinity is only valid
                        // when the canonical accumulator identity encoding
                        // was used above.
                        let decoded_zero := iszero(or(or(x_hi, x_lo), or(y_hi, y_lo)))
                        ok := and(ok, iszero(decoded_zero))
                        // Store the affine point in the exact precompile input
                        // layout: x_hi, x_lo, y_hi, y_lo.
                        mstore(dst, x_hi)
                        mstore(add(dst, 0x20), x_lo)
                        mstore(add(dst, 0x40), y_hi)
                        mstore(add(dst, 0x60), y_lo)
                    }
                }
            }
\end{lstlisting}

Note the re-dispatch: \yul{is\_id} is an output variable that is
\emph{reassigned} at line 1187 from \yul{x\_is\_id} \emph{inside} the outer
\yul{if iszero(is\_id)} block, and the inner \yul{if is\_id} / \yul{if
iszero(is\_id)} at lines 1189/1207 then branch on the new value. The outer
condition was already evaluated, so this is well-defined Yul. In every branch
\yul{dst[0..0x80)} is written, so the destination slot never retains stale
data -- although on a rejecting path the caller reverts before it is read.

\subsubsection{Algorithm}

\begin{algorithm}[htbp]
\caption{Full accumulator point decode, \yul{load\_acc\_point}
(\texttt{Halo2Verifier.sol} lines 1152--1223), instantiated at
$b=56$, $n=7$, $\lambda=4$, $W=2$, $B=2^{56}$}
\label{alg:acc:decode}
\begin{algorithmic}[1]
\Function{LoadAccPoint}{$\mathit{dst}, s$}
  \State $(x_{w_0}, x_{w_1}, y_{w_0}, y_{w_1}) \gets$ calldata words at
         $s, s{+}\texttt{0x20}, s{+}\texttt{0x40}, s{+}\texttt{0x60}$
  \If{$(x_{w_0},x_{w_1}) = S^{\ast}$ \textbf{and} $(y_{w_0},y_{w_1}) = S$}
     \Comment{canonical identity, line 1157}
     \State $\mathit{dst}[0..4] \gets (0,0,0,0)$
     \State \Return $(\mathit{ok}=1,\ \mathit{is\_id}=1)$
  \EndIf
  \State $\mathit{ok} \gets [\,x_{w_0} < 2^{224}\,] \land [\,x_{w_1} < 2^{168}\,]
          \land [\,y_{w_0} < 2^{224}\,] \land [\,y_{w_1} < 2^{168}\,]$
     \Comment{lines 1107, 1081--1102}
  \State $\mathit{is\_id} \gets [\,x_{w_0} \ge B\,] \land
          [\,(x_{w_0} - B) + 2^{224} x_{w_1} = p-1\,]$
     \Comment{lines 1108--1117}
  \State $E_x \gets (x_{w_0} - \mathit{is\_id}\cdot B) + 2^{224} x_{w_1}$,\quad
         $E_y \gets y_{w_0} + 2^{224} y_{w_1}$
     \Comment{line 1121; equation~\eqref{eq:acc:Esimple}}
  \State $\mathit{ok} \mathrel{\&}= [\,E_x \le p-1\,] \land [\,E_y \le p-1\,]$
     \Comment{line 1126}
  \State $x \gets (E_x + 1) \bmod p$,\quad $y \gets (E_y + 1) \bmod p$
     \Comment{lines 1129--1141; equation~\eqref{eq:acc:unshift}}
  \State $\mathit{ok} \mathrel{\&}= [\,\lfloor x/2^{256}\rfloor < 2^{128}\,]
          \land [\,\lfloor y/2^{256}\rfloor < 2^{128}\,]$
     \Comment{line 1146}
  \If{$\mathit{is\_id}$} \Comment{line 1189}
     \State $\mathit{ok} \mathrel{\&}= [\,x = 0 \land y = 0\,]$
     \State $\mathit{dst}[0..4] \gets (0,0,0,0)$
  \Else \Comment{line 1207}
     \State $\mathit{ok} \mathrel{\&}= \lnot[\,x = 0 \land y = 0\,]$
     \State $\mathit{dst}[0..4] \gets (\lfloor x/2^{256}\rfloor,\ x \bmod 2^{256},\
             \lfloor y/2^{256}\rfloor,\ y \bmod 2^{256})$
  \EndIf
  \State \Return $(\mathit{ok}, \mathit{is\_id})$
\EndFunction
\end{algorithmic}
\end{algorithm}

The $\mathit{is\_id}$, $E_x/E_y$ and un-shift steps of
Algorithm~\ref{alg:acc:decode} are stated in the collapsed form of
equation~\eqref{eq:acc:Esimple}, which presupposes the packing step
immediately above them. When that guard fails, the true extraction is
equation~\eqref{eq:acc:limb} and $E_x$, $E_y$, \yul{is\_id} and the branch
taken at line 1189/1207 may all differ from what the algorithm states. That
divergence is unobservable: $\mathit{ok}$ is already $0$ and can only stay $0$,
and both branches write all four words of $\mathit{dst}$, so the caller sees
the same rejection either way.

\subsubsection{The four infinity encodings}
\label{sec:acc:infinity}

Four calldata shapes could plausibly be intended as ``the point at infinity''.
The contract's treatment of each is exhaustively determined by the guards
above; \emph{exactly one} is accepted.

\begin{table}[htbp]
\centering
\begin{tabularx}{\textwidth}{@{}llXl@{}}
\toprule
$x$ words & $y$ words & Contract behaviour & Line \\
\midrule
$S^{\ast}$ & $S$ &
  \textbf{Accepted} as \id{}. Stored as the EIP-2537 four-zero-word slot
  $(0,0,0,0)$ without ever entering the limb decoder. This is the unique
  accepted infinity encoding. & 1157--1166 \\
$S^{\ast}$ & $\neq S$ &
  \textbf{Rejected}. \yul{x\_is\_id} $=1$, so the guard at line 1201 demands
  $x = y = 0$; $y=0$ requires $E_y = p-1$, i.e. $y$ words $= S$, contradiction.
  \yul{ok} $\gets 0$. & 1189--1206 \\
$S$ & $S$ &
  \textbf{Rejected}. No flag, so \yul{is\_id} $=0$; both coordinates un-shift
  to $0$, \yul{decoded\_zero} $=1$ and line 1214 forces \yul{ok} $\gets 0$.
  This is the ``EIP-2537 infinity without the circuit's identity flag'' alias
  and it is closed. & 1207--1214 \\
$S$ & $\neq S$ &
  \textbf{Not an infinity encoding at all}: it decodes to the affine point
  $(0, y)$ with $y \neq 0$, which is stored and handed to \GO{}MSM. See
  Finding~ACC-3. & 1207--1221 \\
\bottomrule
\end{tabularx}
\caption{The four candidate infinity encodings. $S$ is the unflagged $p-1$ word
pair, $S^{\ast}$ the flagged one (\S\ref{sec:acc:sentinels}).}
\label{tab:acc:infinity}
\end{table}

\begin{finding}[ACC-3]{The \texorpdfstring{$p-1$}{p-1} sentinel does collide with the on-curve points \texorpdfstring{$(0,\pm 2)$}{(0, +-2)}, and the collision is closed by the EIP-2537 subgroup check -- not by this contract}
\textbf{Severity: Observation.} Row 4 of Table~\ref{tab:acc:infinity} is a real
decode: $x$ words $=S$ un-shift to $x = 0$, and line 1214 only rejects the
\emph{joint} zero $(0,0)$, so $x=0$ with $y \neq 0$ passes every in-contract
check and is written to the precompile slot as $(0,0,y_{hi},y_{lo})$.

On BLS12-381 ($y^2 = x^3 + 4$) the only points with $x=0$ are $(0, 2)$ and
$(0, p-2)$. Both are on the curve. Doubling gives
$\lambda = 3x^2/(2y) = 0$, hence $[2](0,y) = (0,-y) = -(0,y)$ and
$[3](0,y) = \id$: these are order-3 points, and since $r \equiv 1 \pmod 3$
they are \emph{not} in the $r$-order subgroup. This was recomputed rather
than inferred: $r\cdot(0,2) = (0,2) \neq \id$.
EIP-2537 mandates a subgroup check on \GO{}MSM inputs, so the staticcall at
line 1288/1337 reverts and the proof is rejected.

This is sound but it is a \emph{delegated} soundness property: the contract
does not itself exclude $x=0$, and it must not (a legitimate accumulator point
could in principle have any $x$; it is only $x=0$ that happens to be
subgroup-excluded). The dependency is explicitly acknowledged by the deployment
smoke test, which asserts that \GO{}MSM rejects an on-curve non-subgroup point:
\begin{lstlisting}[style=solidity,caption={\texttt{Halo2Verifier.sol} lines 495--500: the constructor probe that pins the property Finding~ACC-3 relies on},label={lst:acc:subgroup-probe}]
            mstore(add(scratch, 0x00), 0x0000000000000000000000000000000000000000000000000000000000000000)
            mstore(add(scratch, 0x20), 0x0000000000000000000000000000000000000000000000000000000000000004)
            mstore(add(scratch, 0x40), 0x000000000000000000000000000000000a989badd40d6212b33cffc3f3763e9b)
            mstore(add(scratch, 0x60), 0xc760f988c9926b26da9dd85e928483446346b8ed00e1de5d5ea93e354abe706c)
            mstore(add(scratch, 0x80), 1)
            if staticcall(200000, 0x0c, scratch, 0xa0, scratch, 0x80) { revert(0, 0) }
\end{lstlisting}
The probe's own precondition was verified independently of the comment at
lines 486--490: the point $P = (4, y)$ with $y$ as written at lines 497--498
satisfies $y^2 \equiv 4^3 + 4 \pmod p$, and $r\cdot P \neq \id$, so the
staticcall at line 500 really does exercise the subgroup check rather than a
mere off-curve rejection. Note also that the probe forwards a hard-coded
\texttt{200000} gas rather than \yul{G1MSM\_GAS\_1PAIR}, because a
rejecting precompile consumes everything forwarded to it (comment at lines
492--494). A deployment onto a chain whose \GO{}MSM omits the subgroup check
fails at construction time, so the delegation is checked rather than assumed. No change
recommended; recorded because it is the single strongest external assumption
this section makes.
\end{finding}

\begin{finding}[ACC-4]{The ``unreachable'' comment at lines 1193--1200 overstates the case}
\textbf{Severity: Informational.} The comment claims the \yul{if is\_id} branch
at line 1189 is ``[u]nreachable by construction'' because ``the whole-point
sentinel check above already accepted every encoding in which x carries the
identity flag''. That is not accurate: \yul{is\_acc\_encoded\_identity}
requires \emph{both} halves to match, so an input with $x$ words $=S^{\ast}$
and $y$ words $\neq S$ (row 2 of Table~\ref{tab:acc:infinity}) falls through to
line 1189 and does execute this branch. What is genuinely unreachable is the
branch \emph{succeeding}: on that input \yul{ok} is always forced to $0$ by
line 1201, because $y = 0$ would imply $y$ words $= S$ by the injectivity of
the codec. The behaviour is correct -- the input is rejected -- and the guard
is load-bearing rather than dead. Only the comment is wrong, and it is wrong in
the direction that invites a future maintainer to delete a live check.
Suggested fix: restate as ``reachable only with a mismatched $y$, in which case
it always rejects''.
\end{finding}

\subsection{Driver: \yul{validate\_public\_accumulator}}
\label{sec:acc:validate}

\begin{lstlisting}[style=solidity,caption={\texttt{Halo2Verifier.sol} lines 1246--1292: parameter setup, LHS decode, LHS \GO{}MSM validation (comment blocks elided)},label={lst:acc:validate-a}]
            function validate_public_accumulator(success, r) -> out, precompile_failed {
                out := success
                let bits := 56
                let n := 7
                // ...
                let limb_base := shl(bits, 1)
                let limbs_per_word := 4
                let coord_words := div(add(n, sub(limbs_per_word, 1)), limbs_per_word)
                // ...
                let acc_instance_ptr := add(INSTANCE_CPTR, 0x0160)
                // ...
                let lhs_scalar_ptr := add(acc_instance_ptr, mul(mul(2, coord_words), 0x20))
                let lhs_ok, lhs_is_id := load_acc_point(ACC_LHS_MPTR, acc_instance_ptr, bits, n, limb_base)
                out := and(out, lhs_ok)
                // Shared scratch for one-pair LHS validation and the later
                // variable-length RHS MSM.
                let acc_scratch := 0xb140
                {
                    // Already-collapsed point-pair layout: carried scalars are
                    // implicit one.
                    let lhs_scalar := 1
                    // Identity status is useful for decoding checks above, but
                    // validation still goes through G1MSM for all points.
                    pop(lhs_is_id)
                    // Always route the decoded carried point through G1MSM,
                    // even for identity points and zero/one scalars. The
                    // precompile is the on-curve/subgroup validator for this
                    // public-input point; skipping it would let a malformed
                    // non-identity point hide behind scalar 0.
                    mcopy(acc_scratch, ACC_LHS_MPTR, 0x80)
                    mstore(add(acc_scratch, 0x80), lhs_scalar)
                    if out {
                        // Single-pair MSM output overwrites ACC_LHS_MPTR with
                        // lhs_scalar * decoded_lhs. If lhs_scalar is one, this
                        // is also a curve/subgroup validation round-trip.
                        out := staticcall(G1MSM_GAS_1PAIR, 0x0c, acc_scratch, 0xa0, ACC_LHS_MPTR, 0x80)
                        out := and(out, eq(returndatasize(), 0x80))
                        precompile_failed := iszero(out)
                    }
                }
\end{lstlisting}

\begin{lstlisting}[style=solidity,caption={\texttt{Halo2Verifier.sol} lines 1293--1351: RHS decode, one-pair MSM input assembly, RHS fold (comment blocks elided)},label={lst:acc:validate-b}]
                // RHS layout for this generated verifier is an already
                // collapsed point pair: lhs point, rhs point. Both carried
                // scalars are implicit one, and there is no fixed-base scalar
                // tail.
                let rhs_instance_ptr := lhs_scalar_ptr
                // ...
                let rhs_scalar_ptr := add(rhs_instance_ptr, mul(mul(2, coord_words), 0x20))
                let rhs_ok, rhs_is_id := load_acc_point(ACC_RHS_MPTR, rhs_instance_ptr, bits, n, limb_base)
                out := and(out, rhs_ok)
                // acc_pair_ptr appends (G1, scalar) pairs into acc_scratch for
                // one final RHS MSM.
                let acc_pair_ptr := acc_scratch
                {
                    // Implicit unit scalar for already-collapsed point pairs.
                    let rhs_scalar := 1
                    pop(rhs_is_id)
                    // ...
                    mcopy(acc_pair_ptr, ACC_RHS_MPTR, 0x80)
                    mstore(add(acc_pair_ptr, 0x80), rhs_scalar)
                    // Move to the next (G1, scalar) pair slot.
                    acc_pair_ptr := add(acc_pair_ptr, 0xa0)
                }
                // ...
                let acc_msm_len := sub(acc_pair_ptr, acc_scratch)
                if acc_msm_len {
                    // ...
                    if out {
                        // ...
                        out := staticcall(
                            ACC_RHS_MSM_GAS,
                            0x0c,
                            acc_scratch,
                            acc_msm_len,
                            ACC_RHS_MPTR,
                            0x80
                        )
                        out := and(out, eq(returndatasize(), 0x80))
                        precompile_failed := iszero(out)
                    }
                }
                // ...
            }
\end{lstlisting}

\subsubsection{Pointer arithmetic}

With $W=2$ the pointer chain evaluates to constants:
\begin{align*}
  \texttt{acc\_instance\_ptr} &= \texttt{0x1ee4} + \texttt{0x0160} = \texttt{0x2044},\\
  \texttt{lhs\_scalar\_ptr} = \texttt{rhs\_instance\_ptr} &= \texttt{0x2044} + 2W\cdot\texttt{0x20} = \texttt{0x20c4},\\
  \texttt{rhs\_scalar\_ptr} &= \texttt{0x20c4} + 2W\cdot\texttt{0x20} = \texttt{0x2144} = \texttt{calldatasize()},\\
  \texttt{acc\_msm\_len} &= \texttt{0xa0}\ \text{(one $(\GO{},\text{scalar})$ pair)}.
\end{align*}
This render selects the \emph{already-collapsed point-pair} arm of the
generator template: both carried scalars are the implicit literal $1$, there is
no explicit scalar word, and there is no fixed-base scalar tail. Consequently
\yul{rhs\_scalar\_ptr} is never dereferenced (which is fortunate: it equals
\yul{calldatasize()}), the \yul{r} parameter is unused, and the
\yul{if acc\_msm\_len} guard at line 1323 is always taken. All three are
acknowledged in the generated comments (lines 1236--1239, 1262--1263,
1319--1321).

\subsubsection{Algorithm}

\begin{algorithm}[htbp]
\caption{\yul{validate\_public\_accumulator} (\texttt{Halo2Verifier.sol} lines 1246--1351) as rendered in this fixture}
\label{alg:acc:validate}
\begin{algorithmic}[1]
\Function{ValidatePublicAccumulator}{$\mathit{success}, r$}
  \State $\mathit{out} \gets \mathit{success}$; \quad
         $\mathit{precompile\_failed} \gets 0$
  \State $(\mathit{lhs\_ok}, \cdot) \gets
         \Call{LoadAccPoint}{\mptr{ACC\_LHS\_MPTR}, \texttt{0x2044}}$
  \State $\mathit{out} \gets \mathit{out} \land \mathit{lhs\_ok}$
  \State $\mathit{mem}[\texttt{0xb140}..\texttt{0xb1c0}) \gets
         \mathit{mem}[\mptr{ACC\_LHS\_MPTR}..{+}\texttt{0x80})$;\quad
         $\mathit{mem}[\texttt{0xb1c0}] \gets 1$
  \If{$\mathit{out}$}
     \State $\mathit{out} \gets
            \opcode{STATICCALL}(12000, \texttt{0x0c}, \texttt{0xb140},
            \texttt{0xa0}, \mptr{ACC\_LHS\_MPTR}, \texttt{0x80})$
     \State $\mathit{out} \gets \mathit{out} \land
            [\,\opcode{RETURNDATASIZE} = \texttt{0x80}\,]$
     \State $\mathit{precompile\_failed} \gets \lnot \mathit{out}$
  \EndIf
  \State $(\mathit{rhs\_ok}, \cdot) \gets
         \Call{LoadAccPoint}{\mptr{ACC\_RHS\_MPTR}, \texttt{0x20c4}}$
  \State $\mathit{out} \gets \mathit{out} \land \mathit{rhs\_ok}$
  \State $\mathit{mem}[\texttt{0xb140}..\texttt{0xb1c0}) \gets
         \mathit{mem}[\mptr{ACC\_RHS\_MPTR}..{+}\texttt{0x80})$;\quad
         $\mathit{mem}[\texttt{0xb1c0}] \gets 1$
  \If{$\mathit{out}$}
     \State $\mathit{out} \gets
            \opcode{STATICCALL}(12000, \texttt{0x0c}, \texttt{0xb140},
            \texttt{0xa0}, \mptr{ACC\_RHS\_MPTR}, \texttt{0x80})$
     \State $\mathit{out} \gets \mathit{out} \land
            [\,\opcode{RETURNDATASIZE} = \texttt{0x80}\,]$
     \State $\mathit{precompile\_failed} \gets \lnot \mathit{out}$
  \EndIf
  \State \Return $(\mathit{out}, \mathit{precompile\_failed})$
\EndFunction
\end{algorithmic}
\end{algorithm}

\subsection{Curve and subgroup validation via \texorpdfstring{\GO{}}{G1}MSM}
\label{sec:acc:msm}

Neither \yul{load\_acc\_point} nor any caller evaluates $y^2 = x^3 + 4$ or a
cofactor clearing. Curve membership and $r$-order subgroup membership are
delegated in full to the EIP-2537 \GO{}MSM precompile at address
\texttt{0x0c}, by multiplying each decoded point by the scalar $1$:

\begin{lstlisting}[style=solidity,caption={\texttt{Halo2Verifier.sol} lines 1282--1291: the scalar-1 \GO{}MSM round trip that validates the LHS point},label={lst:acc:msm-quote}]
                    mcopy(acc_scratch, ACC_LHS_MPTR, 0x80)
                    mstore(add(acc_scratch, 0x80), lhs_scalar)
                    if out {
                        // Single-pair MSM output overwrites ACC_LHS_MPTR with
                        // lhs_scalar * decoded_lhs. If lhs_scalar is one, this
                        // is also a curve/subgroup validation round-trip.
                        out := staticcall(G1MSM_GAS_1PAIR, 0x0c, acc_scratch, 0xa0, ACC_LHS_MPTR, 0x80)
                        out := and(out, eq(returndatasize(), 0x80))
                        precompile_failed := iszero(out)
                    }
\end{lstlisting}

Properties of this construction, in the order they matter:

\begin{itemize}
  \item \textbf{Unconditional routing.} The point is copied into the MSM input
        \emph{before} the \yul{if out} guard and regardless of
        \yul{lhs\_is\_id}/\yul{rhs\_is\_id} (both are \yul{pop}ped at lines
        1276 and 1309). The generated comment states the reason at lines
        1277--1281: with a conditional skip, ``a malformed non-identity point
        [could] hide behind scalar 0''. In this render the scalars are the
        literal $1$, so the concern is structural rather than live, but the
        invariant ``every decoded accumulator point is fed to \GO{}MSM'' holds
        unconditionally.
  \item \textbf{Output is the canonicalized point.} The MSM writes its
        \texttt{0x80}-byte result back over \mptr{ACC\_LHS\_MPTR} /
        \mptr{ACC\_RHS\_MPTR}. Since the scalar is $1$, the value is unchanged
        mathematically, but downstream consumers see the precompile's own
        canonical EIP-2537 encoding rather than the decoder's. Combined with
        the \yul{eq(returndatasize(), 0x80)} check this closes the
        short-return class of precompile misbehaviour.
  \item \textbf{Identity is accepted.} EIP-2537 encodes \id{} as four zero
        words and \GO{}MSM accepts it, returning \id{}. An accumulator whose
        LHS or RHS is the canonical encoded identity therefore validates and
        contributes $\id$ to the pairing batch.
  \item \textbf{Gas.} \yul{G1MSM\_GAS\_1PAIR} $=$
        \yul{ACC\_RHS\_MSM\_GAS} $= 12000$ (lines 291, 299), which is exactly
        the EIP-2537 cost of a one-pair \GO{}MSM ($12000 \cdot k \cdot
        \mathrm{discount}(1)/1000$). There is no headroom by design; a
        repricing fork makes the staticcall fail, which fails closed as
        \revert{PrecompileFailed}, and the constructor probes forward the same
        bounds so such a chain cannot host a fresh deployment (lines 284--288).
  \item \textbf{Scratch aliasing.} \texttt{0xb140} is a three-way alias:
        \mptr{SELECTOR\_ACC\_MPTR} (line 229),
        \mptr{BATCH\_INV\_SCRATCH\_MPTR} (line 235) and this
        \yul{acc\_scratch} (line 1269). This is not a collision, because the
        three lifetimes are ordered in time: the accumulator phase completes
        and reverts-or-proceeds at line 1448, the Lagrange batch inversion
        first uses the address at line 1826, and the quotient evaluator makes
        its first selector-accumulator write at line 2027. Only
        \texttt{0xb140}--\texttt{0xb1e0} is touched here, which is also clear
        of \mptr{PCS\_FINAL\_MSM\_MPTR} $=$ \texttt{0xb280}. Neither
        staticcall aliases its own input and output windows: input
        \texttt{0xb140}$+$\texttt{0xa0}, output \mptr{ACC\_LHS\_MPTR} /
        \mptr{ACC\_RHS\_MPTR}. \mptr{ACC\_LHS\_MPTR} $=$ \texttt{0x7b40} and
        \mptr{ACC\_RHS\_MPTR} $=$ \texttt{0x7bc0} are adjacent
        \texttt{0x80}-byte slots and do not overlap each other.
\end{itemize}

\begin{finding}[ACC-5]{An off-curve or out-of-subgroup accumulator point is reported as \revert{PrecompileFailed}, not \revert{BadPointEncoding}}
\textbf{Severity: Low (error taxonomy only; both paths fail closed).} The
generated comment at lines 1240--1245 states that \yul{precompile\_failed}
separates ``a G1MSM staticcall that could not run (chain/gas fault)'' from ``a
public-input point this verifier decoded and rejected (bad packing,
out-of-field coordinate, non-canonical identity encoding, \emph{or a point the
precompile found off-curve/out-of-subgroup})'', and that ``only the second is a
BadPointEncoding''. The code does not implement that split for the last
category: lines 1290 and 1346 assign
\yul{precompile\_failed := iszero(out)} where \yul{out} is the staticcall
result itself. EIP-2537 signals an off-curve or out-of-subgroup input by
reverting, so \yul{staticcall} returns $0$, \yul{precompile\_failed} becomes
$1$, and the caller at lines 1451--1452 emits \revert{PrecompileFailed}.

Impact is confined to incident triage -- exactly the concern the MF-4 comment
was written to address. An operator seeing \revert{PrecompileFailed} would look
at the node/fork rather than at the submitted accumulator. The decode-side
rejections (packing, out-of-field, non-canonical identity) \emph{are} correctly
classified, because they zero \yul{out} before the \yul{if out} guard so the
staticcall is skipped and \yul{precompile\_failed} stays $0$. Distinguishing
the two would require reading \opcode{RETURNDATASIZE}/return data or a probe
call, which is not free; the cheapest honest fix is to correct the comment.
\end{finding}

\begin{finding}[ACC-6]{Dead bindings in the rendered arm}
\textbf{Severity: Informational.} In this render \yul{rhs\_scalar\_ptr} (line
1300) is assigned and never read; the \yul{r} parameter (line 1246) is never
used; and \yul{if acc\_msm\_len} (line 1323) is a compile-time-true guard. All
three exist because the template supports carried-scalar and fixed-base-tail
layouts that this artifact does not render
(\texttt{templates/partials/verifier/AccumulatorHelpers.yul}, the
\texttt{expected\_acc\_has\_carried\_scalars} and \texttt{acc\_fixed\_bases}
branches). They cost bytecode but no security. Worth noting only because
\yul{rhs\_scalar\_ptr} evaluates to \yul{calldatasize()}: were a future render
to start dereferencing it without also extending the calldata-shape check at
lines 1424--1427, it would silently read zero.
\end{finding}

\begin{finding}[ACC-7]{Accepting the canonical identity encoding makes the accumulator pairing equation vacuous; nothing in this range constrains the accumulator's value}
\textbf{Severity: Observation (delegated soundness, out of range).} The fast
path at lines 1157--1166 accepts $(S^{\ast}, S)$ and stores the EIP-2537
all-zero slot. EIP-2537 \GO{}MSM accepts that slot and returns \id{}, so
\yul{out} stays $1$ and both \mptr{ACC\_LHS\_MPTR} and
\mptr{ACC\_RHS\_MPTR} can legitimately hold \id{}. The batching block then
adds $\alpha\cdot\id = \id$ into \mptr{PAIRING\_RHS\_MPTR}, and likewise
into \mptr{PAIRING\_LHS\_MPTR} (lines 4072--4096, covered elsewhere), which
leaves both untouched. An accumulator submitted as $(\id, \id)$ therefore
contributes nothing and the deferred KZG accumulator equation is satisfied
trivially.

This is not a defect in lines 992--1351: the accumulator instance words are
attacker-chosen public inputs, and this range deliberately validates only
their \emph{encoding} (canonical packing, in-field coordinates, curve and
subgroup membership) and never their \emph{value}. It is recorded because the
``Identity is accepted'' behaviour above must not be read as an independent
safety property. Soundness of the recursive accumulator rests entirely on the
proving circuit constraining instance words $11..18$ to the accumulator it
actually computed, plus the transcript binding at lines 1508--1522 that stops
those words from being swapped after the fact. Neither can be established from
this fixture.
\end{finding}

\subsection{Soundness assessment}
\label{sec:acc:soundness}

\subsubsection{Can a malicious accumulator alias a valid one?}

No. The accepted-encoding map is injective, and the argument is fully local to
this section:

\begin{enumerate}
  \item \emph{Word-level injectivity.} Under
        \yul{check\_acc\_coord\_packing}, $w_0 < 2^{224}$ and $w_1 < 2^{168}$,
        so equation~\eqref{eq:acc:Esimple} gives
        $E = (w_0 - \delta) + 2^{224} w_1$, hence
        $w_0 = (E \bmod 2^{224}) + \delta$ and
        $w_1 = \lfloor E/2^{224}\rfloor$ (and $\delta \in \{0, B\}$ is
        itself a function of $w_0$): the pair is recoverable from $E$ (\S\ref{sec:acc:packing},
        Finding~ACC-CLEAN-1).
  \item \emph{Coordinate-level injectivity.} The range check
        $E \le p-1$ plus equation~\eqref{eq:acc:unshift} make
        $E \mapsto \mathit{coord}$ a bijection of $[0,p-1]$ onto itself.
  \item \emph{Identity flag determinism.} $\mathtt{is\_id}=1$ requires
        $E_x - 2^{56} = p-1$, which by (1) forces
        $(x_{w_0}, x_{w_1}) = S^{\ast}$ exactly -- the addition in
        equation~\eqref{eq:acc:flag} does not carry out of limb $1$. Hence the
        flag detected by \yul{load\_acc\_coord} and the $x$ half of
        \yul{is\_acc\_encoded\_identity} agree on every input; there is no
        third state.
  \item \emph{One infinity.} \id{} is accepted only as $(S^{\ast}, S)$
        (Table~\ref{tab:acc:infinity}); the affine all-zero point is rejected
        at line 1214, so no non-identity encoding can produce the EIP-2537
        four-zero-word slot.
  \item \emph{Canonical storage.} The \GO{}MSM round trip rewrites both slots
        with the precompile's own encoding, so even the byte-level memory image
        handed to the pairing block is a function of the point alone.
\end{enumerate}

Therefore distinct accepted calldata always denote distinct \GO{} points, and
each \GO{} point has exactly one accepted calldata form. Two further bindings
(covered elsewhere) mean aliasing would not help an attacker even if it
existed: the same eight instance words are absorbed into the Fiat--Shamir
transcript at lines 1508--1522, and the pairing-batch randomizer $\alpha$ at
lines 4057--4067 is a Keccak digest over the materialized \mptr{ACC\_LHS\_MPTR}
/ \mptr{ACC\_RHS\_MPTR} words together with the VK digest.

\subsubsection{Is every limb of every word constrained?}

Yes, and more strongly every \emph{bit} of both packed words is either consumed
by a limb or forced to zero, per Table~\ref{tab:acc:packing}. There is no
unchecked limb slot, no unchecked high remainder, and no coordinate word that
escapes the range check -- with one intentional exception: when
\yul{is\_acc\_encoded\_identity} fires, the packing and range checks are
skipped entirely (line 1157 short-circuits before line 1107 is ever reached),
because the four words were compared against exact constants, which is strictly
stronger.

\subsubsection{Verdict}

The region is sound \emph{as an encoding validator}, which is all it claims to
be. The four decode-side guards (over-packing, base-field range, EIP-2537
high-word padding, joint-zero rejection) together with the mandatory \GO{}MSM
round trip give a complete, injective, canonical decode of the IVC accumulator,
and every failure mode reverts before any transcript state exists. The findings
raised above are one Low-severity error-classification defect (ACC-5), one
incorrect comment over a live check (ACC-4), and five
Observation/Informational notes; none of them admits an accumulator point that
is not a well-formed, canonically encoded element of \GO{}. Two soundness
properties are \emph{delegated} and cannot be discharged here: the EIP-2537
subgroup check (ACC-3, pinned at deployment by the probe at lines 495--500) and
the circuit's constraint on the accumulator's value (ACC-7, not pinned by
anything in this artifact).


%% ==== 08-vkload.tex =====================================================
\section{Verifying-key loading and ABI envelope validation}
\label{sec:vkload:top}

This section specifies \texttt{Halo2Verifier.sol} lines 1360--1456: the block that
materialises the pinned verifying key at \mptr{VK\_MPTR}, cross-checks its header
against verifier-side literals, pins the exact shape of the incoming calldata, and
then dispatches accumulator public-input validation. It is the first phase of
\yul{verifyProof} that performs observable work; everything before it is either a
cheap entry guard or a Yul \yul{function} definition that has not yet been called.
Lines 1454--1456 carry no code: line 1454 is blank, line 1455 is the
\yul{// \_\_phase:transcript} marker, and line 1456 opens the comment banner of the
next section.

\subsection{Entry state and scope}
\label{sec:vkload:entry}

On entry to line 1360 the following hold, all established outside this section:

\begin{itemize}
  \item \yul{vk} is a stack copy of the immutable \yul{AUTHORIZED\_VK}
        (line 646). It is an \emph{immutable}, so the value is a bytecode
        literal, not storage; no caller can influence it.
  \item The Solidity-level ABI head guard has already run (lines 636--641,
        covered elsewhere): it requires
        $\mathrm{calldataload}(\texttt{0x04}) = \texttt{0x40}$ and
        $\mathrm{calldataload}(\texttt{0x24}) =
        \mptr{NUM\_INSTANCE\_CPTR} - \texttt{0x04} = \texttt{0x1ec0}$ (the
        source writes the second bound as \yul{sub(NUM\_INSTANCE\_CPTR, 0x04)},
        not as a literal), i.e. the
        \texttt{proof} head points at \texttt{0x44} and the \texttt{instances}
        head points at \texttt{0x1ec4}. Failure there reverts
        \revert{BadCalldataShape} with an inline selector store.
  \item The MF-2 memory-layout guard has run (lines 674--677). The test is
        \yul{if gt(mload(0x40), TRANSCRIPT\_MPTR)}, i.e. it reverts
        \revert{MemoryLayoutViolated} only when solc's free-memory pointer
        \emph{strictly exceeds} \mptr{TRANSCRIPT\_MPTR} $= \texttt{0x1000}$;
        equality is accepted, correctly, because a free-memory pointer of
        exactly \texttt{0x1000} means the spill reservation ends just below the
        generated layout. Note that the NatSpec on \revert{MemoryLayoutViolated}
        (lines 78--82) describes the guard as firing when the pointer starts
        ``at or above'' \mptr{TRANSCRIPT\_MPTR}, which the code does not do; the
        code is the tighter and the correct of the two. That guard is outside
        this range and is covered elsewhere.
  \item \yul{let r := FR\_MODULUS} (line 1355) and \yul{let success := true}
        (line 1356). In the Yul dialect embedded in Solidity, \yul{true} is the
        literal $1$.
  \item Memory has not been touched above Solidity's reserved prefix. In
        particular nothing has yet been written into
        $[\texttt{0x1000}, \texttt{0x7900})$.
\end{itemize}

This fixture is the \emph{linked-VK} render. The generator template
\texttt{templates/partials/verifier/VkLoading.yul} also offers an embedded-VK
branch (\yul{embedded\_vk = Some(..)}) that \yul{mstore}s the payload words
inline; that branch is not taken here and is not specified. The fixture text
matches the template's \yul{None} arm verbatim -- there is no fixture/template
divergence to report in this range.

\subsection{Constants consumed by this block}
\label{sec:vkload:constants}

\begin{table}[h]
\centering
\begin{tabular}{@{}llll@{}}
\toprule
Constant & Value & Decimal & Declared \\
\midrule
\mptr{EXPECTED\_VK\_PAYLOAD\_LENGTH} & \texttt{0x4280} & $17024$ & line 92 \\
\mptr{EXPECTED\_VK\_LENGTH}          & \texttt{0x4281} & $17025$ & line 93 \\
\mptr{EXPECTED\_VK\_CODEHASH\_WORD}  & \multicolumn{2}{l}{\texttt{0xe68d8936...1c94cf52}} & line 94 \\
\mptr{VK\_MPTR}                      & \texttt{0x3680} & $13952$ & line 128 \\
\mptr{NUM\_INSTANCES\_MPTR}          & \texttt{0x36a0} & --      & line 130 \\
\mptr{K\_MPTR}                       & \texttt{0x36c0} & --      & line 131 \\
\mptr{HAS\_ACCUMULATOR\_MPTR}        & \texttt{0x3760} & --      & line 136 \\
\mptr{ACC\_OFFSET\_MPTR}             & \texttt{0x3780} & --      & line 137 \\
\mptr{NUM\_ACC\_LIMBS\_MPTR}         & \texttt{0x37a0} & --      & line 138 \\
\mptr{NUM\_ACC\_LIMB\_BITS\_MPTR}    & \texttt{0x37c0} & --      & line 139 \\
\mptr{CHALLENGE\_MPTR}               & \texttt{0x7900} & $30976$ & line 144 \\
\mptr{PROOF\_LEN\_CPTR}              & \texttt{0x44}   & $68$    & line 100 \\
\mptr{PROOF\_CPTR}                   & \texttt{0x64}   & $100$   & line 101 \\
\mptr{NUM\_INSTANCE\_CPTR}           & \texttt{0x1ec4} & $7876$  & line 102 \\
\mptr{INSTANCE\_CPTR}                & \texttt{0x1ee4} & $7908$  & line 103 \\
\mptr{ERR\_VK\_MISMATCH}             & \texttt{0xa447d73e} & -- & line 321 \\
\mptr{ERR\_BAD\_CALLDATA\_SHAPE}     & \texttt{0x1b99e37c} & -- & line 320 \\
\mptr{ERR\_BAD\_POINT\_ENCODING}     & \texttt{0xf27905ec} & -- & line 323 \\
\mptr{ERR\_PRECOMPILE\_FAILED}       & \texttt{0x84e81692} & -- & line 324 \\
\bottomrule
\end{tabular}
\caption{Constants read by the VK-loading block. Error constants are
\texttt{bytes4} selectors written by \yul{fail(sel)} (lines 690--693) as
\yul{mstore(0x00, shl(224, sel))} followed by \yul{revert(0x00, 0x04)}.}
\label{tab:vkload:constants}
\end{table}

\subsection{Per-proof re-pinning of the verifying-key runtime}
\label{sec:vkload:pin}

\begin{lstlisting}[style=solidity,caption={\texttt{Halo2Verifier.sol} lines 1376--1393: per-proof VK pin and payload copy. Dedented by 12 spaces; no other edits. The opening brace on line 1376 is the block scope, closed at line 1431 after the checks in Listings~\ref{sec:vkload:header} and~\ref{sec:vkload:abi}},label={lst:vkload:pin}]
{
    // Re-check the pinned VK dependency on every proof. The
    // constructor check catches normal deployment mistakes, while
    // this fresh check hardens forks or same-transaction edge
    // cases where code at the authorized address could differ
    // from the runtime originally pinned by this verifier.
    //
    // EXPECTED_VK_LENGTH includes the leading INVALID byte in the
    // Halo2VerifyingKey runtime. EXPECTED_VK_CODEHASH_WORD is the
    // full runtime hash, not only the payload hash.
    if iszero(and(
        eq(extcodesize(vk), EXPECTED_VK_LENGTH),
        eq(extcodehash(vk), EXPECTED_VK_CODEHASH_WORD)
    )) { fail(ERR_VK_MISMATCH) }
    // Runtime byte 0 is INVALID so direct calls cannot execute the
    // payload. Copy from byte 1 into VK_MPTR to reconstruct the
    // exact payload layout used by the embedded branch.
    extcodecopy(vk, VK_MPTR, 0x01, EXPECTED_VK_PAYLOAD_LENGTH)
\end{lstlisting}

Let $\mathrm{code}(a)$ denote the deployed runtime byte string of account $a$ at
the moment of the call. The guard at lines 1386--1389 is exactly

\[
  \bigl|\mathrm{code}(\yul{vk})\bigr| = 17025
  \;\wedge\;
  \opcode{EXTCODEHASH}(\yul{vk}) =
  \texttt{0xe68d8936}\ldots\texttt{1c94cf52},
\]

and its negation reverts \revert{VkMismatch}. Both operands of \yul{and} are
results of \yul{eq}, hence in $\{0,1\}$, so the bitwise \yul{and} is a correct
logical conjunction; this is the well-known Yul \yul{and} trap, and it is not
triggered here.

The three degenerate account states are all rejected by construction:

\begin{table}[h]
\centering
\begin{tabular}{@{}lll@{}}
\toprule
State of \yul{vk} & \opcode{EXTCODESIZE} & \opcode{EXTCODEHASH} \\
\midrule
non-existent account & $0$ & $0$ \\
exists, empty code   & $0$ & $\mathrm{keccak256}(\varepsilon) = \texttt{0xc5d2}\ldots\texttt{a470}$ \\
exists, wrong code   & arbitrary & $\ne$ pinned word (keccak preimage resistance) \\
\bottomrule
\end{tabular}
\caption{Account states at \yul{vk} and why each fails the pin.}
\label{tab:vkload:codestates}
\end{table}

\paragraph{What the pin buys.} \opcode{EXTCODEHASH} equality against a
generation-time literal means the whole $17025$-byte runtime -- and hence the
$17024$ payload bytes about to be copied out of it -- is, up to keccak256
preimage resistance, byte-identical to the \texttt{Halo2VerifyingKey}
runtime the generator produced alongside this verifier. That covers
\emph{everything} in the payload: \yul{vk\_digest}, the domain scalars
$n^{-1}, \omega, \omega^{-1}, \omega^{-\ell}$, the accumulator metadata, the
$\GO$ and $\GT$ base points used by the final pairing, the compact quotient-VM
constant pool and packed bytecode, and all fixed and permutation commitments.
No later phase re-validates any of those; the codehash pin is their sole
integrity anchor.

\paragraph{What the pin does not buy.} It is an \emph{identity} check, not a
\emph{semantic} one. It asserts ``these bytes are the bytes my generator
emitted''; it cannot assert ``these bytes are the correct verifying key for the
circuit the application means''. In particular \yul{vk\_digest} at
\mptr{VK\_DIGEST\_MPTR} is absorbed into the transcript (covered elsewhere) but
is never compared against an independent statement of the constraint system, and
the six header cross-checks in \S\ref{sec:vkload:header} do not reach it. Binding
the verifier instance to an application-meaningful circuit remains the
integrator's obligation, as the \yul{verifyProof} NatSpec states at lines
588--597.

\begin{finding}[VKLOAD-1]{\opcode{EXTCODESIZE} check is dominated by the \opcode{EXTCODEHASH} check}
\textbf{Severity: Observation.} At \texttt{Halo2Verifier.sol:1387}, the
\yul{eq(extcodesize(vk), EXPECTED\_VK\_LENGTH)} conjunct is implied by the
\yul{extcodehash} conjunct on line 1388: a matching keccak256 digest already
determines the byte string, hence its length. The size check therefore adds no
security beyond a keccak collision. Its cost is one extra account-access opcode
($100$ gas warm; under EIP-2929 one of the two accesses on lines 1387--1388 pays
the $2600$ cold charge and the other the warm $100$, in either evaluation order,
unless the caller already warmed the address via an EIP-2930 access list or an
earlier touch in the same transaction, in which case both are warm). It is retained
as documentation of the $17025 = 17024 + 1$ relationship and as a guard that
would fire independently if \mptr{EXPECTED\_VK\_CODEHASH\_WORD} were ever
hand-edited without regenerating. No action required; noted so a reader does not
mistake it for a second independent trust anchor.
\end{finding}

\begin{finding}[VKLOAD-2]{Self-destruct / redeploy at the pinned address: no TOCTOU window, but a permanent brick}
\textbf{Severity: Informational.} The re-check exists because
\yul{AUTHORIZED\_VK} is an immutable address, and an address is not a commitment
to code. Two properties matter.

\emph{Soundness.} Between the guard at lines 1386--1389 and the
\yul{extcodecopy} at line 1393 there is no \opcode{CALL}, \opcode{CREATE},
\opcode{DELEGATECALL} or any other operation that can transfer control; the two
are adjacent in the same call frame. Therefore no time-of-check/time-of-use skew
is possible and the bytes copied are provably the bytes hashed. This is correct
and complete. (Note also that both pre- and post-EIP-6780 semantics defer account
deletion to the end of the transaction, so even a same-transaction
\opcode{SELFDESTRUCT} cannot make \opcode{EXTCODECOPY} disagree with a preceding
\opcode{EXTCODEHASH}.)

\emph{Availability.} On a chain still running pre-Cancun \opcode{SELFDESTRUCT}
semantics -- an older L2 fork, a permissioned sidechain, or a replayed history --
the account at \yul{vk} can be destroyed and a different runtime redeployed there
via \opcode{CREATE2}. The per-proof check turns that from a soundness failure
(silently verifying against attacker-chosen VK bytes) into an availability
failure: every subsequent \yul{verifyProof} reverts \revert{VkMismatch}, forever,
because \yul{AUTHORIZED\_VK} cannot be repointed. Recovery requires deploying a
fresh \texttt{Halo2VerifyingKey} and a fresh \texttt{Halo2Verifier}, and the
application wrapper swapping the verifier address. This is precisely the
wrapper obligation (``replaceable verifier address'') that lines 594--597
declare a REQUIREMENT and defer to
\texttt{docs/reference/DEPLOYMENT\_AND\_INCIDENT\_RESPONSE.md}. An integrator
that hard-codes the verifier address has no recovery path. The trade is the
right one -- fail closed on availability rather than open on soundness -- but the
residual obligation is real and lives entirely outside this contract.
\end{finding}

\subsection{Payload transfer: \yul{extcodecopy} from byte 1}
\label{sec:vkload:copy}

Line 1393 executes
\[
  \mathrm{mem}\bigl[\texttt{0x3680},\,\texttt{0x3680}+\texttt{0x4280}\bigr)
  \;\leftarrow\;
  \mathrm{code}(\yul{vk})\bigl[\texttt{0x01},\,\texttt{0x4281}\bigr).
\]

The offset $1$ skips the \texttt{INVALID} (\texttt{0xfe}) prefix byte that
\texttt{Halo2VerifyingKey}'s constructor writes at
\texttt{Halo2VerifyingKey.sol:64} before returning
\yul{return(runtime, 0x4281)} at line 693. That prefix exists so that a direct
\opcode{CALL} to the VK address cannot execute payload words as bytecode; it is
not part of the data model, so the verifier drops it here and the resulting
memory image is identical to what the embedded-VK render would \yul{mstore}
directly.

\opcode{EXTCODECOPY} zero-pads if the requested window runs past the end of the
code. Here it cannot: the guard on line 1387 established
$|\mathrm{code}(\yul{vk})| = \mptr{EXPECTED\_VK\_LENGTH} = 17025$, and the two
literals satisfy $\mptr{EXPECTED\_VK\_PAYLOAD\_LENGTH} + 1 =
\mptr{EXPECTED\_VK\_LENGTH}$, so the window $[1, 17025)$ is exactly the payload
and no padding occurs. Worth stating explicitly, since it is the one place a
reader might expect a check and find none: that relation between the two
constants is a generator invariant documented in a comment (lines 89--91) and is
never asserted on chain. Were they emitted inconsistently, line 1387 would still
pass and line 1393 would silently zero-pad or truncate the VK image; both are
compile-time constants declared on adjacent lines 92--93, so the invariant is
checkable by inspection rather than by the contract. The two constants are
consistent with the VK contract, whose last payload word is written at
\texttt{payload + 0x4260} (\texttt{Halo2VerifyingKey.sol:689}), giving
$\texttt{0x4260} + \texttt{0x20} = \texttt{0x4280} = 17024$ bytes $= 532$ words.

\paragraph{Destination extent.} The copy lands in
$[\texttt{0x3680}, \texttt{0x7900})$ and
$\texttt{0x3680} + \texttt{0x4280} = \texttt{0x7900} = \mptr{CHALLENGE\_MPTR}$
(line 144). The VK image therefore abuts the challenge region exactly, with no
gap and no overlap. Nothing below \texttt{0x3680} is disturbed, which matters
because \yul{scalar\_inv} (lines 700--722) deliberately reuses the dead transcript
buffer immediately below \mptr{VK\_MPTR} as modexp scratch: its EIP-198 frame base
is \yul{p := 0x3580} (line 710) and the frame spans
$[\texttt{0x3580}, \texttt{0x3640})$, ending $\texttt{0x40}$ bytes short of
\mptr{VK\_MPTR}. Nothing live is
overwritten either: at this point in execution memory above Solidity's reserved
prefix is untouched.

\begin{table}[h]
\centering
\begin{tabular}{@{}rlll@{}}
\toprule
Payload word & Memory address & Named constant & Content \\
\midrule
0  & \texttt{0x3680} & \mptr{VK\_DIGEST\_MPTR}         & \yul{vk\_digest} \\
1  & \texttt{0x36a0} & \mptr{NUM\_INSTANCES\_MPTR}     & \yul{num\_instances} \\
2  & \texttt{0x36c0} & \mptr{K\_MPTR}                  & $k$ \\
3  & \texttt{0x36e0} & \mptr{N\_INV\_MPTR}             & $n^{-1}$ \\
4  & \texttt{0x3700} & \mptr{OMEGA\_MPTR}              & $\omega$ \\
5  & \texttt{0x3720} & \mptr{OMEGA\_INV\_MPTR}         & $\omega^{-1}$ \\
6  & \texttt{0x3740} & \mptr{OMEGA\_INV\_TO\_L\_MPTR}  & $\omega^{-\ell}$ \\
7  & \texttt{0x3760} & \mptr{HAS\_ACCUMULATOR\_MPTR}   & \yul{has\_accumulator} \\
8  & \texttt{0x3780} & \mptr{ACC\_OFFSET\_MPTR}        & \yul{acc\_offset} \\
9  & \texttt{0x37a0} & \mptr{NUM\_ACC\_LIMBS\_MPTR}    & \yul{num\_acc\_limbs} \\
10 & \texttt{0x37c0} & \mptr{NUM\_ACC\_LIMB\_BITS\_MPTR} & \yul{num\_acc\_limb\_bits} \\
11--14 & \texttt{0x37e0} & \mptr{G1\_BASE\_MPTR}       & $\GO$ generator, EIP-2537 padded \\
15--22 & \texttt{0x3860} & \mptr{G2\_BASE\_MPTR}       & $\GT$ generator \\
23--30 & \texttt{0x3960} & \mptr{NEG\_S\_G2\_BASE\_MPTR} & $-s\,\GT$ generator \\
31--$\ast$ & \texttt{0x3a60} & --                      & quotient-VM constants + packed program \\
$\ast$ & -- & -- & fixed commitments, then permutation commitments (4 words each) \\
\bottomrule
\end{tabular}
\caption{VK payload layout as landed by line 1393. Word indices are from
\texttt{Halo2VerifyingKey.sol} lines 22--39; addresses are
$\mptr{VK\_MPTR} + 32 i$ and match the named constants at lines 128--142.}
\label{tab:vkload:layout}
\end{table}

\paragraph{Cost.} The block pays, per proof: one cold account access
($2600$) plus one warm ($100$) for lines 1387--1388; \opcode{EXTCODECOPY} at
$100 + 3 \cdot 532 = 1696$; and the first memory expansion of the frame, from the
three or four words solc has touched (exactly how many depends on codegen) to
$968$ words $= \texttt{0x7900}$ bytes, at
$3 \cdot 968 + \lfloor 968^2/512 \rfloor - 12 = 2904 + 1830 - 12 = 4722$
(the subtracted term is the few gas already charged for those first words, so
this figure is good to $\pm 3$).
That is roughly $9.1\mathrm{k}$ gas before any header or ABI check. The
embedded-VK alternative trades this for a much larger verifier runtime and a
correspondingly larger deployment cost; for a 17\,KB payload the linked-contract
choice is clearly right, and the extra per-proof cost is small next to the
EIP-2537 precompile budget that follows.

\subsection{The six VK header cross-checks}
\label{sec:vkload:header}

\begin{lstlisting}[style=solidity,caption={\texttt{Halo2Verifier.sol} lines 1400--1406: VK header cross-checks (dedented by 12 spaces)},label={lst:vkload:header}]
    success := and(success, eq(mload(NUM_INSTANCES_MPTR), 19))
    success := and(success, eq(mload(K_MPTR), 20))
    success := and(success, eq(mload(HAS_ACCUMULATOR_MPTR), 1))
    success := and(success, eq(mload(ACC_OFFSET_MPTR), 11))
    success := and(success, eq(mload(NUM_ACC_LIMBS_MPTR), 7))
    success := and(success, eq(mload(NUM_ACC_LIMB_BITS_MPTR), 56))
    if iszero(success) { fail(ERR_VK_MISMATCH) }
\end{lstlisting}

Formally, with $H_j$ the $j$-th payload word now resident in memory, the block
enforces
\[
  H_1 = 19,\quad H_2 = 20,\quad H_7 = 1,\quad
  H_8 = 11,\quad H_9 = 7,\quad H_{10} = 56,
\]
and reverts \revert{VkMismatch} otherwise. All six comparisons are full 32-byte
\yul{eq}s, so no truncation or aliasing is possible. \yul{success} enters as $1$
and each \yul{eq} yields $0$ or $1$, so the fold stays in $\{0,1\}$ and
\yul{iszero(success)} is an exact ``some check failed'' test.

The essential point is that each literal on the right-hand side is the value the
\emph{code generator inlined elsewhere in this same file}. The check is what
binds the VK's declaration to the verifier's hard-coded assumption:

\begin{table}[h]
\centering
\begin{tabularx}{\textwidth}{@{}llX@{}}
\toprule
Header word & Literal & Verifier-side site the check binds it to \\
\midrule
\yul{num\_instances} & $19$ &
  \yul{eq(19, calldataload(NUM\_INSTANCE\_CPTR))} at line 1420; the
  \texttt{0x0260} instance-array span at lines 1426, 1511 and 1855;
  \mptr{NUM\_INSTANCE\_CPTR}/\mptr{INSTANCE\_CPTR} at lines 102--103. \\
$k$ & $20$ &
  \yul{let k := 20} at line 1800, the repeated-squaring bound for
  $x^n$, $n = 2^{20}$, and again at line 3381 for the joint
  $x^n$/$x^{n-1}$ walk. \mptr{K\_MPTR} is read \emph{nowhere else} in the
  contract than line 1401. \\
\yul{has\_accumulator} & $1$ &
  The existence of the \yul{\_\_phase:accumulator\_msm} block at lines
  1432--1453 at all; the template renders it under
  \yul{expected\_has\_accumulator}. \\
\yul{acc\_offset} & $11$ &
  \yul{add(INSTANCE\_CPTR, 0x0160)} at line 1257, since
  $\texttt{0x0160} = 352 = 11 \cdot 32$. \\
\yul{num\_acc\_limbs} & $7$ &
  \yul{let n := 7} at line 1249, feeding the limb-recomposition loop bound
  and $\mathrm{coord\_words} = \lceil 7/4 \rceil = 2$. \\
\yul{num\_acc\_limb\_bits} & $56$ &
  \yul{let bits := 56} at line 1248 and
  $\mathrm{limb\_base} = 2^{56}$ at line 1252. \\
\bottomrule
\end{tabularx}
\caption{Each header cross-check and the inlined verifier literal it pins.}
\label{tab:vkload:crosscheck}
\end{table}

Against the fixture VK these all hold:
\texttt{Halo2VerifyingKey.sol} lines 69, 70, 75, 76, 77 and 78 store
\texttt{0x13} $= 19$, \texttt{0x14} $= 20$, $1$, \texttt{0x0b} $= 11$, $7$ and
\texttt{0x38} $= 56$ respectively. The pair is self-consistent.

\begin{finding}[VKLOAD-3]{Header cross-checks are a build-consistency canary, not an independent trust anchor, and they cover 6 of 11 header words}
\textbf{Severity: Informational.} Because line 1388 already pinned the full
runtime by codehash, the words loaded at lines 1400--1405 are a deterministic
function of a generation-time constant. Consequently lines 1400--1406 can only
fail if \mptr{EXPECTED\_VK\_CODEHASH\_WORD} and the six inline literals were
produced from mutually inconsistent generator inputs -- i.e. it is a
\emph{generator-drift} detector evaluated on chain, once per proof, rather than a
defence against any runtime adversary. It is unreachable for a self-consistent
build.

Two consequences a reviewer should hold explicitly:
(i) the checks cost roughly $90$ gas per proof and are worth keeping precisely
because they turn a catastrophic silent mismatch into a typed revert; but
(ii) their coverage is asymmetric and must not be read as a general VK
validation. Header words $0$ (\yul{vk\_digest}), $3$ ($n^{-1}$), $4$ ($\omega$),
$5$ ($\omega^{-1}$) and $6$ ($\omega^{-\ell}$), plus all base points, the entire
quotient-VM program, and every fixed and permutation commitment, have no
verifier-side literal to be compared against and are covered by the codehash pin
alone. A generator bug that emitted, say, an $\omega$ of the wrong order together
with a matching codehash would pass every check in this block. Nothing in the
deployed contract can catch that class of fault; it is a generator-testing
obligation (\texttt{src/lowering/tests.rs}), not an on-chain one.
\end{finding}

\begin{finding}[VKLOAD-4]{\revert{VkMismatch} conflates ``wrong code at the pinned address'' with ``generator drift''}
\textbf{Severity: Informational.} Lines 1389 and 1406 revert with the same
selector \texttt{0xa447d73e}. An incident responder decoding a revert therefore
cannot distinguish a VK contract that was destroyed, redeployed, or never
existed (line 1389 -- an operational/chain event, remediated by redeploying the
pair) from a verifier and VK built from inconsistent generator inputs (line 1406
-- a build fault, remediated by regenerating and re-auditing). The contract's own
error taxonomy comment (lines 48--57) argues for exactly this kind of
distinction, and the NatSpec on \yul{VkMismatch} (lines 61--63) explicitly
lumps both cases together. Given VKLOAD-3, the second case is unreachable for a
correct build, so the practical impact is nil; splitting the selectors would cost
one additional \yul{error} declaration. Recorded for completeness only.
\end{finding}

\subsection{ABI envelope validation}
\label{sec:vkload:abi}

\begin{lstlisting}[style=solidity,caption={\texttt{Halo2Verifier.sol} lines 1419--1430: dynamic ABI envelope checks (dedented by 12 spaces)},label={lst:vkload:abi}]
    success := and(success, eq(0x1e60, calldataload(PROOF_LEN_CPTR)))
    success := and(success, eq(19, calldataload(NUM_INSTANCE_CPTR)))
    // Calldata must contain exactly the ABI selector, proof bytes,
    // instance-array length, and generated number of instance
    // words. Any trailing bytes fail closed.
    success := and(
        success,
        eq(calldatasize(), add(INSTANCE_CPTR, 0x0260))
    )
    // Stop before any transcript absorption if the ABI/proof shape
    // is not exactly the generated one.
    if iszero(success) { fail(ERR_BAD_CALLDATA_SHAPE) }
\end{lstlisting}

The three conditions are
\begin{align}
  \mathrm{calldataload}(\texttt{0x44})   &= \texttt{0x1e60} = 7776, \label{eq:vkload:plen}\\
  \mathrm{calldataload}(\texttt{0x1ec4}) &= 19,                     \label{eq:vkload:ninst}\\
  \mathrm{calldatasize}()                &= \texttt{0x1ee4} + \texttt{0x0260} = \texttt{0x2144} = 8516. \label{eq:vkload:size}
\end{align}
Note that \eqref{eq:vkload:size} is evaluated as an \opcode{ADD} of two
compile-time constants, which the optimiser folds; the source keeps the
\yul{add} form because the generator emits the instance span
$19 \cdot 32 = \texttt{0x0260}$ as a separate substitution.

\begin{table}[h]
\centering
\begin{tabular}{@{}llrl@{}}
\toprule
Region & Offset & Bytes & Pinned by \\
\midrule
selector \yul{verifyProof(bytes,uint256[])} & \texttt{0x00} & $4$ & solc dispatcher \\
\texttt{proof} head ($=\texttt{0x40}$)      & \texttt{0x04} & $32$ & line 636 \\
\texttt{instances} head ($=\texttt{0x1ec0}$)& \texttt{0x24} & $32$ & line 636 \\
\texttt{proof} length ($=\texttt{0x1e60}$)  & \texttt{0x44} & $32$ & line 1419 \\
\texttt{proof} bytes                        & \texttt{0x64} & $7776$ & length above \\
\texttt{instances} length ($=19$)           & \texttt{0x1ec4} & $32$ & line 1420 \\
\texttt{instances} words                    & \texttt{0x1ee4} & $608$ & line 1424 \\
\midrule
total                                       & & $8516$ & \\
\bottomrule
\end{tabular}
\caption{The one accepted calldata shape. The arithmetic closes exactly:
$\mptr{PROOF\_CPTR} + \texttt{0x1e60} = \texttt{0x64} + \texttt{0x1e60}
= \texttt{0x1ec4} = \mptr{NUM\_INSTANCE\_CPTR}$, and
$\mptr{NUM\_INSTANCE\_CPTR} + \texttt{0x20} = \mptr{INSTANCE\_CPTR}$, so the
proof bytes abut the instance-array header with no slack.}
\label{tab:vkload:calldata}
\end{table}

\paragraph{Short calldata is safe.} \opcode{CALLDATALOAD} zero-extends past
\opcode{CALLDATASIZE} rather than reverting, so on truncated input
\eqref{eq:vkload:plen} and \eqref{eq:vkload:ninst} read $0$ and fail, and
\eqref{eq:vkload:size} fails independently. There is no path on which a short
calldata is accepted and later read out of bounds.

\paragraph{Why the envelope must be pinned before line 1447.} The accumulator
validator invoked at line 1447 reads its public-input limbs straight from
calldata, starting at \yul{add(INSTANCE\_CPTR, 0x0160)} (line 1257). With
$\mathrm{coord\_words} = \lceil 7/4 \rceil = 2$ (line 1254) each coordinate
occupies two instance words and each G1 point four, and the render decodes two
points (lines 1265 and 1301). It therefore consumes instance words $11$ through
$18$, that is calldata
$[\texttt{0x1ee4}+\texttt{0x0160},\ \texttt{0x1ee4}+\texttt{0x0260})
= [\texttt{0x2044}, \texttt{0x2144})$ -- exactly the tail of the instance array,
ending precisely at the \opcode{CALLDATASIZE} pinned by \eqref{eq:vkload:size}.
Because \yul{success} is enforced at line 1430 \emph{before} the call at line
1447, every one of those \opcode{CALLDATALOAD}s is guaranteed to read a word the
caller actually supplied rather than \opcode{CALLDATALOAD}'s implicit zero
extension past \opcode{CALLDATASIZE}. Two things are worth being explicit about.
First, the containment $\yul{acc\_offset} + 8 \le \yul{num\_instances}$, i.e.
$11 + 8 \le 19$, is never asserted on chain: lines 1400 and 1403 pin the two
values independently against literals, and their compatibility is a
generator-side invariant. Second, \yul{rhs\_scalar\_ptr} (line 1300) evaluates to
instance word $19$, i.e. calldata offset \texttt{0x2144} $=$
\opcode{CALLDATASIZE}, one word past the end of the array; in this render it is
computed and never dereferenced, so the boundary is touched but not crossed. Both
facts are properties of the callee and are specified elsewhere; they are recorded
here because \eqref{eq:vkload:ninst} and \eqref{eq:vkload:size} are what make
them safe.

\paragraph{Redundancy and the load-bearing check.} Solidity's own calldata
decoder for \yul{bytes calldata proof, uint256[] calldata instances} already
validates that each head is in range and that
$\mathrm{head} + 32 + 32\cdot\mathrm{len} \le \mathrm{calldatasize}$. So
\eqref{eq:vkload:plen} and \eqref{eq:vkload:ninst} are largely -- though not
entirely, since solc only enforces $\le$, not the generated \emph{value} --
duplicated work. The genuinely load-bearing check is \eqref{eq:vkload:size}: the
ABI permits arbitrary trailing bytes, solc ignores them, and only line 1424
rejects them.

What the exact-size check buys is not verification soundness in the narrow
sense -- every read in this contract is at a fixed offset, so trailing bytes
could never alter the accept/reject decision -- but \emph{envelope
non-malleability}: for a given (proof, instances) pair there is exactly one
calldata string this contract accepts. An application that deduplicates
submissions on \yul{keccak256(msg.data)}, or that reasons about transaction
identity, gets that property for free instead of having to re-derive it. Note
the scope: this pins the \emph{envelope} only. Malleability inside the proof
bytes (EIP-2537 padding bytes, out-of-field coordinates, non-canonical scalars)
is rejected by \yul{common\_uncompressed\_g1} at lines 780--789 and by the
proof-scalar canonicality guards at lines 1523, 1706, 1759 and 1789, all covered
in other sections.

\begin{finding}[VKLOAD-5]{Exact \opcode{CALLDATASIZE} equality is incompatible with calldata-appending relayers}
\textbf{Severity: Observation.} Line 1424 makes \yul{verifyProof} uncallable by
any caller that appends bytes to calldata: ERC-2771 trusted forwarders (which
append a 20-byte sender suffix), and any multicall or paymaster wrapper that
does the same. Such a call reverts \revert{BadCalldataShape} regardless of proof
validity. This is a deliberate, documented trade -- the NatSpec states it
explicitly at lines 606--611 and directs integrators to route relayed traffic
through a wrapper that reassembles exact calldata. It is recorded here because
it is an integration hazard that is invisible at the ABI level and produces a
failure indistinguishable from a malformed proof unless the revert data is
decoded. No change recommended; the alternative (accepting a suffix) would
forfeit the non-malleability property above.
\end{finding}

\begin{finding}[VKLOAD-6]{The VK-load and ABI-envelope region is otherwise clean}
\textbf{Severity: Observation.} Stated positively, because the STYLE contract
asks for it. Lines 1376--1431 contain no unbounded loop, no attacker-influenced
memory offset, no attacker-influenced length, and no external call other than
the two account-metadata opcodes and \opcode{EXTCODECOPY} against an immutable
address. Every destination and length in the block is a compile-time constant;
the only attacker-controlled inputs read are the two \opcode{CALLDATALOAD} words
at lines 1419 and 1420 and \opcode{CALLDATASIZE} at line 1426, each compared for
exact $32$-byte equality against a constant.
The \yul{success} fold is confined to $\{0,1\}$ and is enforced twice (lines
1406 and 1430) before any value derived from calldata is used. The block is
sound.
\end{finding}

\subsection{Accumulator validation call site and the MF-4 error split}
\label{sec:vkload:accsite}

\begin{lstlisting}[style=solidity,caption={\texttt{Halo2Verifier.sol} lines 1446--1453: accumulator validation call site (dedented by 12 spaces)},label={lst:vkload:accsite}]
let acc_precompile_failed := 0
success, acc_precompile_failed := validate_public_accumulator(success, r)
if iszero(success) {
    // MF-4: a G1MSM that could not run at all is a chain fault,
    // not a malformed accumulator point.
    if acc_precompile_failed { fail(ERR_PRECOMPILE_FAILED) }
    fail(ERR_BAD_POINT_ENCODING)
}
\end{lstlisting}

\yul{validate\_public\_accumulator} is defined at lines 1246--1351; its interior
(limb recomposition from radix-$2^{56}$ instance words, canonicality checks, and
the two \opcode{STATICCALL}s to the EIP-2537 G1MSM precompile at address
\texttt{0x0c}) is specified elsewhere. This section specifies only the call-site
contract.

\begin{algorithm}[h]
\caption{Call-site error classification at lines 1446--1453}
\label{alg:vkload:mf4}
\begin{algorithmic}[1]
\State $\yul{acc\_precompile\_failed} \gets 0$
\State $(\yul{success}, \yul{acc\_precompile\_failed}) \gets
        \yul{validate\_public\_accumulator}(\yul{success}, r)$
\If{$\yul{success} = 0$}
  \If{$\yul{acc\_precompile\_failed} \ne 0$}
    \State \Return revert \revert{PrecompileFailed}
  \EndIf
  \State \Return revert \revert{BadPointEncoding}
\EndIf
\State \textbf{continue} to the transcript phase (line 1455 onward)
\end{algorithmic}
\end{algorithm}

The classification is exact because of how the callee sets its second return
value. \yul{precompile\_failed} is written only inside the two \yul{if out}
guards (lines 1284--1291 and 1327--1347), in each case as the third of three
consecutive statements:
\yul{out := staticcall(...)}, then
\yul{out := and(out, eq(returndatasize(), 0x80))}, then
\yul{precompile\_failed := iszero(out)} -- lines 1288--1290 for the LHS
single-pair MSM and lines 1337--1346 for the RHS MSM. Hence:

\begin{table}[h]
\centering
\begin{tabularx}{\textwidth}{@{}Xll@{}}
\toprule
Failure mode & \yul{acc\_precompile\_failed} & Revert at call site \\
\midrule
Limb decode / canonicality rejected a public-input point
  & $0$ (guards skipped, \yul{if out} false)
  & \revert{BadPointEncoding} \\
G1MSM \opcode{STATICCALL} returned failure (out of gas at the precompile,
  precompile absent on this chain, or point off-curve / out of subgroup)
  & $1$
  & \revert{PrecompileFailed} \\
G1MSM succeeded but $\mathrm{returndatasize} \ne \texttt{0x80}$
  & $1$
  & \revert{PrecompileFailed} \\
LHS MSM failed, so the RHS block's \yul{if out} never runs
  & stays $1$ from the LHS assignment
  & \revert{PrecompileFailed} \\
\bottomrule
\end{tabularx}
\caption{MF-4 error split as realised by the call site. Note that
\yul{success} is provably $1$ on entry to line 1447, having been enforced at
line 1430.}
\label{tab:vkload:mf4}
\end{table}

\begin{finding}[VKLOAD-7]{An off-curve accumulator point is reported as \revert{PrecompileFailed}, contradicting the MF-4 comment that governs the split}
\textbf{Severity: Low} (diagnostics only; the accept/reject decision is
unaffected and fail-closed). The EIP-2537 G1MSM precompile signals ``input point
is not on the curve or not in the prime-order subgroup'' by reverting, which at
the Yul level is indistinguishable from ``precompile unavailable'' or ``out of
gas'': all three yield \opcode{STATICCALL} $= 0$, hence \yul{out} $= 0$, hence
\yul{precompile\_failed} $= 1$ at line 1290 or 1346, hence
\revert{PrecompileFailed} at line 1451.

The MF-4 comment at lines 1240--1245 states the opposite. It enumerates the
\revert{BadPointEncoding} bucket as ``a public-input point this verifier decoded
and rejected (bad packing, out-of-field coordinate, non-canonical identity
encoding, \emph{or a point the precompile found off-curve/out-of-subgroup})'',
and closes ``only the second is a BadPointEncoding''. The fourth item in that
list is precisely the case the code routes to \revert{PrecompileFailed}. This is
not an acknowledged limitation in the comment; it is a contradiction between the
comment and the code, and the code is what ships. Only the first three items --
the ones \yul{load\_acc\_point} and its helpers reject by returning
\yul{ok} $= 0$ without ever entering an \yul{if out} guard -- actually produce
\revert{BadPointEncoding}.

Consequences. The gas budgets are fixed compile-time constants
(\mptr{G1MSM\_GAS\_1PAIR} $= 12000$ at line 291, \mptr{ACC\_RHS\_MSM\_GAS}
$= 12000$ at line 299), so the ambiguity cannot be induced by gas-griefing, and
\revert{PrecompileFailed}'s own NatSpec (lines 70--71, ``A precompile call failed
or returned an unexpected size'') is literally accurate. But the failure is
attacker-selectable in the wrong direction: any caller can submit a well-formed,
canonically packed, in-field accumulator point that simply is not on the curve
and thereby make the verifier emit the selector reserved for chain faults. An
integrator whose monitoring treats \revert{PrecompileFailed} as ``EIP-2537 is
missing or repriced on this chain'' -- which is exactly what the constants
comment at lines 285--288 trains a reader to believe -- can be made to raise a
false chain-capability alarm by an ordinary bad submission. Recommended fix:
either correct the MF-4 comment to match the code, or, if the documented split is
the intended one, have the call site distinguish the cases (a probe call, or a
cheaper separate on-curve check before the MSM). Behaviour is sound either way;
only the reported cause is wrong.
\end{finding}

\paragraph{Why the accumulator is validated here, before the transcript.}
The comment at lines 1433--1445 gives the rationale: an invalid accumulator
public input is rejected before any transcript absorption, quotient evaluation,
PCS work, or final pairing. Three things this ordering buys:

\begin{enumerate}
  \item \emph{Cost of rejection.} A malformed accumulator is rejected for at
        most the two budgeted G1MSM calls ($2 \times 12000$ gas) plus the load
        of this block, instead of the full verification. Note the asymmetry
        between the two rejection classes: a decode or canonicality failure
        clears \yul{out} at line 1266 or 1302 and therefore skips the
        \yul{if out} guards entirely, costing \emph{no} precompile call at all,
        while a curve/subgroup failure necessarily pays for the MSM that
        detects it.
  \item \emph{Challenge hygiene.} No value derived from an unvalidated
        accumulator point reaches the Keccak transcript. Since the accumulator
        limbs are public instances that \emph{are} absorbed later, this does not
        change the transcript contents on the accepting path; it only guarantees
        that the rejecting path never squeezes a challenge from unvalidated
        data.
  \item \emph{Output staging.} The two \opcode{STATICCALL}s write their
        $\texttt{0x80}$-byte results to \mptr{ACC\_LHS\_MPTR} $= \texttt{0x7b40}$
        and \mptr{ACC\_RHS\_MPTR} $= \texttt{0x7bc0}$ (lines 166--167), two
        dedicated four-word slots above the VK image
        $[\texttt{0x3680}, \texttt{0x7900})$, above the transcript buffer that
        starts at \texttt{0x1000}, and disjoint from every challenge slot
        (\texttt{0x7900}--\texttt{0x7a20}, lines 150--159) and from
        \mptr{F\_COM\_MPTR}/\mptr{PI\_MPTR} (lines 162--163). That the
        intervening transcript, quotient and PCS phases do not overwrite them is
        a property of code outside this range; the generator's arena rejects any
        two live regions sharing an address (lines 104--118), so this section
        \emph{asserts} rather than proves that the validated points reach the
        final pairing unchanged.
\end{enumerate}

What it costs: on a proof that would have failed for an unrelated reason
earlier in the transcript, roughly $24\mathrm{k}$ gas of precompile budget is
spent first. That is the correct trade for a verifier whose accept path is
dominated by pairing costs anyway.

\subsection{Guard summary}
\label{sec:vkload:summary}

\begin{table}[h]
\centering
\begin{tabularx}{\textwidth}{@{}llX@{}}
\toprule
Line & Revert & Condition that must hold to continue \\
\midrule
1386--1389 & \revert{VkMismatch} &
  $\opcode{EXTCODESIZE}(\yul{vk}) = 17025 \wedge
   \opcode{EXTCODEHASH}(\yul{vk}) = \texttt{0xe68d8936}\ldots$ \\
1400 & \revert{VkMismatch} & $\mathrm{mem}[\texttt{0x36a0}] = 19$ \\
1401 & \revert{VkMismatch} & $\mathrm{mem}[\texttt{0x36c0}] = 20$ \\
1402 & \revert{VkMismatch} & $\mathrm{mem}[\texttt{0x3760}] = 1$ \\
1403 & \revert{VkMismatch} & $\mathrm{mem}[\texttt{0x3780}] = 11$ \\
1404 & \revert{VkMismatch} & $\mathrm{mem}[\texttt{0x37a0}] = 7$ \\
1405 & \revert{VkMismatch} & $\mathrm{mem}[\texttt{0x37c0}] = 56$ \\
1419 & \revert{BadCalldataShape} & $\mathrm{calldataload}(\texttt{0x44}) = \texttt{0x1e60}$ \\
1420 & \revert{BadCalldataShape} & $\mathrm{calldataload}(\texttt{0x1ec4}) = 19$ \\
1424--1427 & \revert{BadCalldataShape} & $\mathrm{calldatasize}() = \texttt{0x2144}$ \\
1448--1452 & \revert{PrecompileFailed} / \revert{BadPointEncoding} &
  \yul{validate\_public\_accumulator} returned \yul{out} $\ne 0$ \\
\bottomrule
\end{tabularx}
\caption{Every revert reachable from lines 1360--1456, with its guard. The rows
for lines 1400--1405 and 1419--1427 name the line that computes each condition,
not a revert site: those conditions fold into \yul{success} and are enforced
jointly at line 1406 and line 1430 respectively, so a failure of any one of the
six header checks, or of any one of the three envelope checks, is
indistinguishable from a failure of its siblings in the revert data. No other
line in the range can terminate execution: in particular
\yul{validate\_public\_accumulator} (line 1246) and every helper it calls
(\yul{load\_acc\_coord\_shifted}, \yul{is\_bls\_p\_minus\_one},
\yul{is\_acc\_encoded\_identity}, \yul{check\_acc\_coord\_packing},
\yul{load\_acc\_coord}, \yul{load\_acc\_point}, lines 1010--1351) contain no
\yul{fail}, \yul{revert}, \yul{return}, \yul{stop} or \yul{invalid}, so the call
at line 1447 always returns and the only reverts it can cause are the two at
lines 1451--1452.}
\label{tab:vkload:guards}
\end{table}

The code \emph{in} this range contains no loop, no dynamic memory offset and no
dynamic length: every offset and every length on lines 1376--1453 is a
compile-time constant. This is a statement about the range's own instructions,
not about the whole dynamic extent of a proof: the call at line 1447 enters
\yul{validate\_public\_accumulator}, which does contain loops (limb
recomposition) and does pass a computed length \yul{acc\_msm\_len} to a
\opcode{STATICCALL} at line 1341; that callee is specified elsewhere, and as
noted in Table~\ref{tab:vkload:guards} it cannot terminate execution itself.
Control leaves the range only by falling through to the transcript phase, whose
first instruction follows the \yul{// \_\_phase:transcript} marker on line 1455,
or by one of the reverts in Table~\ref{tab:vkload:guards}.


%% ==== 09-proof-parse.tex ================================================
\section{Proof parsing and the Fiat--Shamir challenge schedule}
\label{sec:tp:main}

This section specifies \texttt{Halo2Verifier.sol} lines 1456--1792, the block the
generator tags \yul{// \_\_phase:transcript}. It is the soundness hinge of the
whole contract: it is the only place where proof calldata is turned into field
elements and group elements, the only place where the Fiat--Shamir transcript is
built, and the only place where the ten verifier challenges
$\theta,\beta,\gamma,\tau,y,x,x_1,x_2,x_3,x_4$ come into existence. Everything
after it (Lagrange evaluation, quotient reconstruction, PCS multi-open, final
pairing) consumes values this block wrote and never re-reads the raw proof except
through the single saved cursor \mptr{Q\_EVAL\_CPTR\_MPTR}.

The block does three jobs simultaneously, and the specification below keeps them
separate even though the code interleaves them:

\begin{enumerate}
  \item \textbf{Absorb} public data and proof bytes into a streaming Keccak-256
        transcript in exactly the order the native midnight-proofs verifier does,
        and squeeze ten challenges out of it.
  \item \textbf{Range-check} every scalar against $r$ and every $\GO$ coordinate
        against $p$, plus the EIP-2537 zero-padding discipline.
  \item \textbf{Distribute} the parsed commitments and evaluations into the
        planned memory windows consumed by the later phases.
\end{enumerate}

\subsection{Entry state}
\label{sec:tp:entry}

Three facts hold on entry at line 1456 and the parser depends on all three. They
are established earlier in the file and are covered elsewhere; they are restated
here because the parser has no runtime bound checks of its own.

\begin{enumerate}
  \item \textbf{The ABI envelope is exact} (lines 636--642 and 1419--1430). The
        two dynamic-argument heads must be \texttt{0x40} and
        $\texttt{NUM\_INSTANCE\_CPTR}-4 = \texttt{0x1ec0}$; the \yul{proof} byte length at
        \mptr{PROOF\_LEN\_CPTR} must equal \texttt{0x1e60}; the instance-array
        length at \mptr{NUM\_INSTANCE\_CPTR} must equal $19$; and
        \opcode{CALLDATASIZE} must equal $\texttt{INSTANCE\_CPTR}+\texttt{0x0260}
        = \texttt{0x2144}$ exactly, so trailing bytes fail closed with
        \revert{BadCalldataShape}. Consequently every \opcode{CALLDATALOAD} the
        parser performs is inside real calldata; none can silently read the
        zero-extension past \opcode{CALLDATASIZE}.
  \item \textbf{The VK is pinned} (lines 1382--1406). The linked
        \texttt{Halo2VerifyingKey} runtime is checked by \opcode{EXTCODESIZE} and
        \opcode{EXTCODEHASH} on every call before its payload is copied to
        \mptr{VK\_MPTR}, and six header words are cross-checked against generated
        constants ($\text{num\_instances}=19$, $k=20$, $\text{has\_accumulator}=1$,
        $\text{acc\_offset}=11$, $\text{num\_acc\_limbs}=7$,
        $\text{num\_acc\_limb\_bits}=56$). The digest absorbed first by the
        transcript is therefore not attacker-controlled.
  \item \textbf{\yul{success} is $1$ and \yul{r} is $r$}. Line 1355 binds
        \yul{let r := FR\_MODULUS} and line 1356 binds \yul{let success := true};
        the guards at lines 1406, 1430 and 1448--1453 each revert on
        \yul{iszero(success)}, so the flag is exactly $1$ when the parser starts
        folding into it.
\end{enumerate}

\begin{table}[t]
\centering
\caption{Calldata cursors and moduli used by the parser.}
\label{tab:tp:constants}
\begin{tabular}{@{}llll@{}}
\toprule
Constant & Value & Line & Meaning \\
\midrule
\texttt{PROOF\_LEN\_CPTR}     & \texttt{0x0044} & 100 & word holding the proof byte length \\
\texttt{PROOF\_CPTR}          & \texttt{0x0064} & 101 & first proof byte \\
\texttt{NUM\_INSTANCE\_CPTR}  & \texttt{0x1ec4} & 102 & instance-array length word \\
\texttt{INSTANCE\_CPTR}       & \texttt{0x1ee4} & 103 & first instance word \\
\texttt{TRANSCRIPT\_MPTR}     & \texttt{0x1000} & 119 & base of the Keccak buffer \\
\texttt{FR\_MODULUS} $=r$     & \texttt{0x73eda753\ldots00000001} & 331 & BLS12-381 scalar modulus \\
\texttt{BLS\_P\_HI}           & \texttt{0x1a0111ea397fe69a4b1ba7b6434bacd7} & 336 & top 16 bytes of $p$ (zero-padded to a word) \\
\texttt{BLS\_P\_MINUS\_ONE\_LO} & \texttt{0x64774b84\ldots ffffaaaa} & 337 & low 32 bytes of $p-1$ \\
\bottomrule
\end{tabular}
\end{table}

Note that $\texttt{PROOF\_CPTR} + \texttt{0x1e60} = \texttt{0x1ec4} =
\texttt{NUM\_INSTANCE\_CPTR}$ and
$\texttt{INSTANCE\_CPTR} + 19\cdot\texttt{0x20} = \texttt{0x2144}$: the ABI
envelope is fully determined by the generated constants, with no slack.

\subsection{The streaming Keccak transcript}
\label{sec:tp:helpers}

The transcript is not a sponge with a persistent state; it is a
\emph{hash-and-reseed} chain over a memory buffer that begins at
\mptr{TRANSCRIPT\_MPTR}. Four helpers implement it. The cursor \yul{buf\_len} is
an absolute memory address, not a length, despite its name.

\begin{lstlisting}[style=solidity,caption={\texttt{Halo2Verifier.sol} lines 733--743: \yul{transcript\_init} and \yul{common\_word}},label={lst:tp:common-word}]
            function transcript_init() -> buf_len {
                // Empty transcript buffer starts exactly at TRANSCRIPT_MPTR.
                buf_len := TRANSCRIPT_MPTR
            }

            // Append one 32-byte big-endian field/transcript word at the
            // current end of the transcript buffer.
            function common_word(buf_len, word) -> ret {
                mstore(buf_len, word)
                ret := add(buf_len, 32)
            }
\end{lstlisting}

\begin{lstlisting}[style=solidity,caption={\texttt{Halo2Verifier.sol} lines 775--796: \yul{common\_uncompressed\_g1} (the leading 30-line comment block, lines 745--774, is elided)},label={lst:tp:common-g1}]
            function common_uncompressed_g1(buf_len, cptr) -> ret {
                let x_hi_word := calldataload(cptr)
                let x_lo := calldataload(add(cptr, 0x20))
                let y_hi_word := calldataload(add(cptr, 0x40))
                let y_lo := calldataload(add(cptr, 0x60))
                if shr(128, x_hi_word) { fail(ERR_BAD_POINT_ENCODING) }
                if shr(128, y_hi_word) { fail(ERR_BAD_POINT_ENCODING) }

                let x_hi := and(x_hi_word, 0xffffffffffffffffffffffffffffffff)
                let y_hi := and(y_hi_word, 0xffffffffffffffffffffffffffffffff)
                if iszero(or(lt(x_hi, BLS_P_HI), and(eq(x_hi, BLS_P_HI), iszero(gt(x_lo, BLS_P_MINUS_ONE_LO))))) {
                    fail(ERR_BAD_POINT_ENCODING)
                }
                if iszero(or(lt(y_hi, BLS_P_HI), and(eq(y_hi, BLS_P_HI), iszero(gt(y_lo, BLS_P_MINUS_ONE_LO))))) {
                    fail(ERR_BAD_POINT_ENCODING)
                }

                // Memcpy the 4 calldata words (128 bytes) verbatim
                // into the keccak buffer.
                calldatacopy(buf_len, cptr, 0x80)
                ret := add(buf_len, 0x80)
            }
\end{lstlisting}

\begin{lstlisting}[style=solidity,caption={\texttt{Halo2Verifier.sol} lines 801--809: \yul{squeeze\_to}},label={lst:tp:squeeze}]
            function squeeze_to(buf_len, mptr) -> ret {
                let h0 := keccak256(TRANSCRIPT_MPTR, sub(buf_len, TRANSCRIPT_MPTR))
                // Reseed: write the 32-byte digest at start of buffer.
                mstore(TRANSCRIPT_MPTR, h0)
                let r := FR_MODULUS
                // Sample Fq as uint256(keccak_digest_be) mod r.
                mstore(mptr, mod(h0, r))
                ret := add(TRANSCRIPT_MPTR, 32)
            }
\end{lstlisting}

\paragraph{Semantics of \yul{common\_uncompressed\_g1}.}
A $\GO$ point occupies four calldata words in EIP-2537 padded uncompressed form,
$x_{\mathrm{hi}} \Vert x_{\mathrm{lo}} \Vert y_{\mathrm{hi}} \Vert
y_{\mathrm{lo}}$, where each coordinate is $16$ zero pad bytes followed by $48$
big-endian field bytes. The helper enforces two properties.

\begin{itemize}
  \item \emph{Padding discipline.} \yul{shr(128, x\_hi\_word)} must be zero, i.e.
        the top $16$ bytes of each \yul{\_hi} word are all zero. Without this,
        two distinct calldata encodings would map to the same point but different
        Keccak inputs, so the transcript would not be a function of the proof.
  \item \emph{Field range.} With
        $x = x_{\mathrm{hi}}\cdot 2^{256} + x_{\mathrm{lo}}$ where
        $x_{\mathrm{hi}}$ is the low $16$ bytes of the first word, the guard is
        \[
          x_{\mathrm{hi}} < P_{\mathrm{hi}}
          \;\lor\;
          \bigl(x_{\mathrm{hi}} = P_{\mathrm{hi}} \;\land\;
                x_{\mathrm{lo}} \le P^{-}_{\mathrm{lo}}\bigr),
        \]
        with $P_{\mathrm{hi}} = \lfloor p/2^{256}\rfloor$ and
        $P^{-}_{\mathrm{lo}} = (p-1) \bmod 2^{256}$. Since
        $p \bmod 2^{256}$ ends in \texttt{\ldots aaab}, the subtraction defining
        $P^{-}_{\mathrm{lo}}$ does not borrow, and the guard is exactly
        $x < p$. The same test is applied to $y$. Failure is
        \revert{BadPointEncoding}.
\end{itemize}

The absorbed bytes are the \emph{verbatim} $128$ calldata bytes
(\opcode{CALLDATACOPY}), not a re-serialisation, so the transcript input is byte-identical
to the calldata whatever the coordinate values are. There is deliberately no
on-curve or subgroup check here; see Finding~\ref{find:tp:curve}.

\paragraph{Semantics of \yul{squeeze\_to}.}
Let $B_i$ be the buffer contents at the $i$-th squeeze. Then
\[
  h_i \;=\; \mathrm{keccak256}(B_i), \qquad
  c_i \;=\; h_i \bmod r, \qquad
  B_{i+1} \;=\; h_i \Vert (\text{bytes absorbed after squeeze } i).
\]
The buffer is truncated to the $32$-byte digest and the cursor is reset to
$\texttt{TRANSCRIPT\_MPTR}+32$. Two consequences matter. First, the challenge
value stored at \yul{mptr} is the \emph{reduced} $c_i$, but the reseed uses the
\emph{unreduced} $h_i$, so two transcripts that happen to collide modulo $r$ still
diverge immediately afterwards. Second, the buffer between squeezes is short --
the largest is $2720$ bytes -- which bounds the \opcode{KECCAK256} word cost.

\paragraph{Equality with the native transcript primitives.}
All four helper semantics were checked directly against the native
implementation in \texttt{midfall/proofs/src/transcript/implementors.rs}, not
merely against the contract's own comments:

\begin{itemize}
  \item \yul{TranscriptHash for Keccak256}'s \yul{squeeze} (lines 198--211) finalises a
        clone of the state, then replaces the state with a \emph{fresh}
        \texttt{Keccak256} seeded with exactly the $32$-byte digest. That is
        precisely $B_{i+1} = h_i \Vert \cdots$, i.e.\ \yul{squeeze\_to}'s
        \yul{mstore(TRANSCRIPT\_MPTR, h0)} plus cursor reset. There is no
        retained sponge state on either side.
  \item The \yul{Sampleable} impl for \yul{Fq} over \yul{Keccak256} (lines 337--341)
        delegates its \yul{sample} to
        \yul{sample\_keccak\_digest\_be\_mod\_r} (lines 19--40), which computes
        \yul{BigUint::from\_bytes\_be(digest) \% modulus}. This is
        \yul{mod(h0, r)} exactly -- \emph{not} a $64$-byte
        \yul{from\_uniform\_bytes} wide reduction (which is what the
        \texttt{Blake2b} instances at lines 119--124 and 171--177 use). The
        contract's rule at line 807 therefore matches, and
        Finding~\ref{find:tp:sampling} applies to both sides equally.
  \item The \yul{Hashable} impl for \yul{G1Projective} over \yul{Keccak256} has a
        \yul{to\_input} (lines 283--298) that emits $128$ bytes: $x$ big-endian into \texttt{out[16..64]}, $y$ into
        \texttt{out[80..128]}, and all-zero for the identity. That is the same
        EIP-2537 padded layout \yul{common\_uncompressed\_g1} accepts, and it
        confirms the $0^{128}$ literal used for \yul{committed\_pi}.
  \item The \yul{Hashable} impl for \yul{Fq} over \yul{Keccak256} has a \yul{to\_input}
        (lines 316--320) that is
        \yul{to\_repr()} (little-endian) reversed, i.e.\ the canonical $32$-byte
        big-endian word -- what \yul{common\_word} stores.
\end{itemize}

One divergence between that file's documentation and the contract is worth
naming, because it is the sort of thing a reader will take on trust: the native
doc comment (line 278) says the EVM verifier masks ``the 16-byte zero pads to
defeat transcript malleability''. The contract does not mask, it \emph{rejects}
(lines 780--781). Rejecting is the strictly stronger behaviour -- masking would
map two calldata encodings onto one transcript input, which is exactly the
aliasing the comment claims to prevent -- so the code is right and the comment
is stale. The same comment (line 275) also names the helper
\yul{common\_g1\_uncompressed}, whereas the rendered function is
\yul{common\_uncompressed\_g1}.

\paragraph{No precompiles, no external calls.}
For completeness: lines 1456--1792 contain no \opcode{STATICCALL},
\opcode{CALL}, \opcode{EXTCODECOPY} or any other external interaction. The
entire block is \opcode{CALLDATALOAD}, \opcode{CALLDATACOPY}, \opcode{MSTORE},
\opcode{MLOAD}, \opcode{KECCAK256} and arithmetic. Every question about gas
bounds and return-size checks for precompiles therefore belongs to the
accumulator block before it (line 1288, covered elsewhere) and the PCS/pairing
blocks after it; the parser itself cannot fail on a precompile.

\paragraph{Why the concatenation is unambiguous.}
The transcript applies no domain separation and no length prefixing: absorbed
items are simply concatenated. This is safe here only because, for a fixed VK
digest, the entire schedule -- the number of items, their order, and each item's
byte width ($32$ or $128$) -- is a compile-time constant of the render, so the
byte string admits exactly one parse. The VK digest is the first thing absorbed
and pins which schedule is in force, so two different circuits cannot produce
the same transcript prefix by re-partitioning the same bytes. A future emitter
that made any section length depend on proof content would break this property
and would need explicit length framing.

\begin{table}[t]
\centering
\caption{Challenge slots written by the ten \yul{squeeze\_to} call sites. These
are the only memory words outside the transcript buffer that the squeeze path
touches; each is written exactly once.}
\label{tab:tp:challenges}
\begin{tabular}{@{}lllll@{}}
\toprule
$i$ & Challenge & Constant & Address & Squeeze line \\
\midrule
1  & $\theta$ & \texttt{THETA\_MPTR}           & \texttt{0x7900} & 1568 \\
2  & $\beta$  & \texttt{BETA\_MPTR}            & \texttt{0x7920} & 1586 \\
3  & $\gamma$ & \texttt{GAMMA\_MPTR}           & \texttt{0x7940} & 1587 \\
4  & $\tau$   & \texttt{TRASH\_CHALLENGE\_MPTR} & \texttt{0x7960} & 1646 \\
5  & $y$      & \texttt{Y\_MPTR}               & \texttt{0x7980} & 1663 \\
6  & $x$      & \texttt{X\_MPTR}               & \texttt{0x79a0} & 1686 \\
7  & $x_1$    & \texttt{X1\_MPTR}              & \texttt{0x79c0} & 1720 \\
8  & $x_2$    & \texttt{X2\_MPTR}              & \texttt{0x79e0} & 1721 \\
9  & $x_3$    & \texttt{X3\_MPTR}              & \texttt{0x7a00} & 1733 (masked at 1743) \\
10 & $x_4$    & \texttt{X4\_MPTR}              & \texttt{0x7a20} & 1767 \\
\bottomrule
\end{tabular}
\end{table}

The ten slots are contiguous on a \texttt{0x20} stride from \texttt{0x7900} to
\texttt{0x7a20} inclusive (constants at lines 150--159), and the next constant
after them is \mptr{F\_COM\_MPTR} $=\texttt{0x7a40}$, so no squeeze can reach
the $\GO$ slots written later in the same block.

\subsection{Proof calldata layout}
\label{sec:tp:layout}

Table~\ref{tab:tp:layout} is the complete proof layout, reconstructed from the
loop bounds in the source. Offsets are relative to \mptr{PROOF\_CPTR} and also
given as absolute calldata addresses. Every element is either a $\GO$
($\texttt{0x80}$ bytes) or an $\Fr$ scalar ($\texttt{0x20}$ bytes).

\begin{table}[t]
\centering\small
\caption{Proof layout. \texttt{rel} is the offset from \mptr{PROOF\_CPTR}
(\texttt{0x64}); \texttt{abs} is the calldata address. Every section length in
the source is a compile-time constant and equals (count $\times$ element size).}
\label{tab:tp:layout}
\begin{tabular}{@{}llrrlll@{}}
\toprule
\# & Section & Cnt & Type & rel & abs range & Memory base \\
\midrule
1  & advice ($\mathsf{A}_{0..14}$)            & 15  & $\GO$ & \texttt{0x0000} & \texttt{0x0064}--\texttt{0x07e4} & \texttt{0xa140} \\
2  & lookup multiplicities ($\mathsf{M}_{0,1}$) & 2  & $\GO$ & \texttt{0x0780} & \texttt{0x07e4}--\texttt{0x08e4} & \texttt{0xa8c0} \\
3  & permutation $Z$ ($\mathsf{Z}^{\mathrm{p}}_{0..5}$) & 6 & $\GO$ & \texttt{0x0880} & \texttt{0x08e4}--\texttt{0x0be4} & \texttt{0xa9c0} \\
4  & lookup 0 helper                           & 1   & $\GO$ & \texttt{0x0b80} & \texttt{0x0be4}--\texttt{0x0c64} & \texttt{0xacc0} \\
5  & lookup 0 accumulator                      & 1   & $\GO$ & \texttt{0x0c00} & \texttt{0x0c64}--\texttt{0x0ce4} & \texttt{0xadc0} \\
6  & lookup 1 helper                           & 1   & $\GO$ & \texttt{0x0c80} & \texttt{0x0ce4}--\texttt{0x0d64} & \texttt{0xad40} \\
7  & lookup 1 accumulator                      & 1   & $\GO$ & \texttt{0x0d00} & \texttt{0x0d64}--\texttt{0x0de4} & \texttt{0xae40} \\
8  & trashcan                                  & 1   & $\GO$ & \texttt{0x0d80} & \texttt{0x0de4}--\texttt{0x0e64} & \texttt{0xaec0} \\
9  & quotient limbs ($\mathsf{Q}_{0..3}$)      & 4   & $\GO$ & \texttt{0x0e00} & \texttt{0x0e64}--\texttt{0x1064} & \texttt{0xaf40} \\
10 & evaluations $\mathbf{e}$                  & 102 & $\Fr$ & \texttt{0x1000} & \texttt{0x1064}--\texttt{0x1d24} & \texttt{0x9480} \\
11 & $\mathsf{F}$ (\yul{f\_com})               & 1   & $\GO$ & \texttt{0x1cc0} & \texttt{0x1d24}--\texttt{0x1da4} & \texttt{0x7a40} \\
12 & \yul{q\_evals}                            & 5   & $\Fr$ & \texttt{0x1d40} & \texttt{0x1da4}--\texttt{0x1e44} & (cursor only) \\
13 & $\pi$ (\yul{pi})                          & 1   & $\GO$ & \texttt{0x1de0} & \texttt{0x1e44}--\texttt{0x1ec4} & \texttt{0x7ac0} \\
\midrule
   & \textbf{total}                            & 34 $\GO$ + 107 $\Fr$ & & & \texttt{0x1e60} bytes & \\
\bottomrule
\end{tabular}
\end{table}

\paragraph{Counts, verified against the source.} The assignment brief listed the
schedule; every count in it is confirmed by the fixture. Specifically:
advice $=\texttt{0x0780}/\texttt{0x80} = 15$ (line 1553);
multiplicities $=\texttt{0x0100}/\texttt{0x80} = 2$ (line 1573);
permutation $Z$ $=\texttt{0x0300}/\texttt{0x80} = 6$ (line 1592);
lookup helpers $+$ accumulators $= 4$ total, emitted as two unrolled
\emph{helper loop $+$ straight-line accumulator} pairs (lines 1610--1639);
trashcans $=\texttt{0x80}/\texttt{0x80} = 1$ (line 1651);
quotient limbs $=\texttt{0x0200}/\texttt{0x80} = 4$ (line 1674);
evaluations $=\texttt{0x0cc0}/\texttt{0x20} = 102$ (line 1699);
\yul{q\_evals} $=\texttt{0xa0}/\texttt{0x20} = 5$ (line 1754);
$\mathsf{F}$ and $\pi$ one each (lines 1727, 1773).
The totals close exactly:
$34\cdot 128 + 107\cdot 32 = 4352 + 3424 = 7776 = \texttt{0x1e60}$.
There is one presentational correction to the brief: the four lookup commitments
are interleaved on the wire as
\emph{helper$_0$, acc$_0$, helper$_1$, acc$_1$}, not as a block of four
``helpers $+$ $z$''; the wire order and the memory order differ (rows 4--7 of
Table~\ref{tab:tp:layout} show the memory bases out of address order for exactly
this reason).

\paragraph{Memory windows.} The per-category $\GO$ bases tile without gaps, and
the eval spill window abuts the first of them:
\[
\begin{aligned}
\texttt{0x9480} + 102\cdot\texttt{0x20} &= \texttt{0xa140} = \texttt{ADVICE\_COMMS\_MPTR\_BASE}, \\
\texttt{0xa140} + 15\cdot\texttt{0x80}  &= \texttt{0xa8c0} = \texttt{LOOKUP\_M\_COMMS\_MPTR\_BASE}, \\
\texttt{0xa8c0} + 2\cdot\texttt{0x80}   &= \texttt{0xa9c0} = \texttt{PERM\_Z\_COMMS\_MPTR\_BASE}, \\
\texttt{0xa9c0} + 6\cdot\texttt{0x80}   &= \texttt{0xacc0} = \texttt{LOOKUP\_HELPER\_COMMS\_MPTR\_BASE}, \\
\texttt{0xacc0} + 2\cdot\texttt{0x80}   &= \texttt{0xadc0} = \texttt{LOOKUP\_Z\_COMMS\_MPTR\_BASE}, \\
\texttt{0xadc0} + 2\cdot\texttt{0x80}   &= \texttt{0xaec0} = \texttt{TRASHCAN\_COMMS\_MPTR\_BASE}, \\
\texttt{0xaec0} + 1\cdot\texttt{0x80}   &= \texttt{0xaf40} = \texttt{QUOTIENT\_LIMB\_COMMS\_MPTR\_BASE}, \\
\texttt{0xaf40} + 4\cdot\texttt{0x80}   &= \texttt{0xb140} = \texttt{SELECTOR\_ACC\_MPTR}.
\end{aligned}
\]
Each walk pointer is advanced by \texttt{0x80} exactly once per commitment, in
lockstep with \yul{proof\_cptr}, so each window is written exactly to its
capacity and no walk overruns into the next category. \mptr{F\_COM\_MPTR}
(\texttt{0x7a40}), \mptr{PI\_MPTR} (\texttt{0x7ac0}) and \mptr{ACC\_LHS\_MPTR}
(\texttt{0x7b40}) are likewise contiguous and disjoint, so writing $\mathsf{F}$
and $\pi$ cannot disturb the accumulator points that the pre-transcript
accumulator validator (function at line 1246, called at line 1447) already
validated.

\subsection{The absorb/squeeze schedule}
\label{sec:tp:schedule}

Algorithm~\ref{alg:tp:schedule} is the exact ordered schedule executed by lines
1474--1789. $\mathrm{cw}(w)$ denotes \yul{common\_word}, $\mathrm{cg1}(c)$
denotes \yul{common\_uncompressed\_g1} at calldata cursor $c$, and
$\mathrm{sq}()$ denotes \yul{squeeze\_to}.

\begin{algorithm}[t]
\caption{Transcript schedule, \texttt{Halo2Verifier.sol} lines 1474--1789}
\label{alg:tp:schedule}
\begin{algorithmic}[1]
\State $B \gets \varepsilon$ \Comment{\yul{buf\_len := transcript\_init()}, line 1474}
\State $\mathrm{cw}(\mathrm{mload}(\texttt{VK\_DIGEST\_MPTR}))$ \Comment{line 1481}
\State $B \gets B \Vert 0^{128}$ \Comment{\yul{committed\_pi} $=\id$, lines 1496--1500}
\State $\mathrm{cw}(19)$ \Comment{instance count, line 1508}
\For{$i = 0$ \textbf{to} $18$} \Comment{lines 1511--1522}
  \State $t \gets \mathrm{calldataload}(\texttt{INSTANCE\_CPTR} + 32i)$
  \State $\mathit{success} \gets \mathit{success} \wedge (t < r)$
  \State $\mathrm{cw}(t)$
\EndFor
\State \textbf{if} $\neg\mathit{success}$ \textbf{then} \revert{NonCanonicalScalar} \Comment{line 1523}
\State $c \gets \texttt{PROOF\_CPTR}$
\For{$j = 0$ \textbf{to} $14$} \Comment{advice, lines 1553--1561}
  \State $\mathrm{cg1}(c)$; \; copy $128$ bytes to $\texttt{0xa140}+128j$; \; $c \mathrel{+}= 128$
\EndFor
\State $\theta \gets \mathrm{sq}()$ \Comment{line 1568}
\For{$j = 0$ \textbf{to} $1$} \Comment{multiplicities, lines 1573--1580}
  \State $\mathrm{cg1}(c)$; \; copy to $\texttt{0xa8c0}+128j$; \; $c \mathrel{+}= 128$
\EndFor
\State $\beta \gets \mathrm{sq}()$; \; $\gamma \gets \mathrm{sq}()$ \Comment{lines 1586--1587}
\For{$j = 0$ \textbf{to} $5$} \Comment{permutation $Z$, lines 1592--1599}
  \State $\mathrm{cg1}(c)$; \; copy to $\texttt{0xa9c0}+128j$; \; $c \mathrel{+}= 128$
\EndFor
\For{$\ell = 0$ \textbf{to} $1$} \Comment{lookups, unrolled at lines 1608--1639}
  \State $\mathrm{cg1}(c)$; \; copy to $\texttt{0xacc0}+128\ell$; \; $c \mathrel{+}= 128$ \Comment{helper}
  \State $\mathrm{cg1}(c)$; \; copy to $\texttt{0xadc0}+128\ell$; \; $c \mathrel{+}= 128$ \Comment{accumulator}
\EndFor
\State $\tau \gets \mathrm{sq}()$ \Comment{\yul{trash\_challenge}, line 1646}
\State $\mathrm{cg1}(c)$; \; copy to $\texttt{0xaec0}$; \; $c \mathrel{+}= 128$ \Comment{trashcan, a $1$-iteration \yul{for} loop with bound \texttt{0x80}, lines 1651--1658}
\State $y \gets \mathrm{sq}()$ \Comment{line 1663}
\For{$j = 0$ \textbf{to} $3$} \Comment{quotient limbs, lines 1674--1681}
  \State $\mathrm{cg1}(c)$; \; copy to $\texttt{0xaf40}+128j$; \; $c \mathrel{+}= 128$
\EndFor
\State $x \gets \mathrm{sq}()$ \Comment{line 1686}
\For{$k = 0$ \textbf{to} $101$} \Comment{evaluations, lines 1699--1714}
  \State $e_k \gets \mathrm{calldataload}(c)$
  \State \textbf{if} $e_k \ge r$ \textbf{then} \revert{NonCanonicalScalar}
  \State $\mathrm{mstore}(\texttt{0x9480}+32k,\, e_k)$; \; $\mathrm{cw}(e_k)$; \; $c \mathrel{+}= 32$
\EndFor
\State $x_1 \gets \mathrm{sq}()$; \; $x_2 \gets \mathrm{sq}()$ \Comment{lines 1720--1721}
\State $\mathrm{cg1}(c)$; \; copy to \mptr{F\_COM\_MPTR}; \; $c \mathrel{+}= 128$ \Comment{lines 1727--1729}
\State $x_3 \gets \mathrm{sq}()$; \; $x_3 \gets x_3 \wedge (2^{128}-1)$ \Comment{lines 1733, 1743}
\State $\mathrm{mstore}(\texttt{Q\_EVAL\_CPTR\_MPTR},\, c)$ \Comment{line 1753}
\For{$k = 0$ \textbf{to} $4$} \Comment{\yul{q\_evals}, lines 1754--1762}
  \State $v \gets \mathrm{calldataload}(c)$
  \State \textbf{if} $v \ge r$ \textbf{then} \revert{NonCanonicalScalar}
  \State $\mathrm{cw}(v)$; \; $c \mathrel{+}= 32$
\EndFor
\State $x_4 \gets \mathrm{sq}()$ \Comment{line 1767}
\State $\mathrm{cg1}(c)$; \; copy to \mptr{PI\_MPTR}; \; $c \mathrel{+}= 128$ \Comment{lines 1773--1775}
\State \textbf{if} $c \ne \texttt{NUM\_INSTANCE\_CPTR}$ \textbf{then} \revert{BadCalldataShape} \Comment{line 1785}
\State \textbf{if} $\neg\mathit{success}$ \textbf{then} \revert{NonCanonicalScalar} \Comment{line 1789}
\end{algorithmic}
\end{algorithm}

Exactly ten \yul{squeeze\_to} call sites appear in the range (lines 1568, 1586,
1587, 1646, 1663, 1686, 1720, 1721, 1733, 1767), matching the ten fixed
challenges the generator's native anchor enumerates
(\texttt{src/lowering/transcript\_anchor.rs:50}, \yul{FIXED\_CHALLENGES = 10},
whose doc comment lists them in exactly this order). This
render squeezes no user-phase challenges: the template emits
\yul{squeeze\_to(buf\_len, add(CHALLENGE\_MPTR, ...))} once per user challenge and
this circuit declares none, so the schedule begins at $\theta$ immediately after
the phase-1 advice loop.

Table~\ref{tab:tp:buffer} gives the exact number of bytes hashed at each squeeze,
which is what \opcode{KECCAK256} is charged for and what bounds the buffer's high-water
mark.

\begin{table}[t]
\centering
\caption{Bytes hashed per squeeze. ``Reseed'' is the $32$-byte digest carried
forward from the previous squeeze.}
\label{tab:tp:buffer}
\begin{tabular}{@{}llrr@{}}
\toprule
$i$ & Challenge & Absorbed since previous squeeze & $|B_i|$ (bytes) \\
\midrule
1  & $\theta$ & $32 + 128 + 32 + 608 + 1920$ (no reseed) & $2720$ \\
2  & $\beta$  & reseed $+\ 2\times128$    & $288$ \\
3  & $\gamma$ & reseed only               & $32$ \\
4  & $\tau$   & reseed $+\ 6\times128 + 4\times128$ & $1312$ \\
5  & $y$      & reseed $+\ 1\times128$    & $160$ \\
6  & $x$      & reseed $+\ 4\times128$    & $544$ \\
7  & $x_1$    & reseed $+\ 102\times32$   & $3296$ \\
8  & $x_2$    & reseed only               & $32$ \\
9  & $x_3$    & reseed $+\ 1\times128$    & $160$ \\
10 & $x_4$    & reseed $+\ 5\times32$     & $192$ \\
\midrule
-- & (none)   & reseed $+\ 1\times128$ ($\pi$, never hashed) & $160$ \\
\bottomrule
\end{tabular}
\end{table}

The buffer high-water mark is
$\texttt{TRANSCRIPT\_MPTR} + 3296 = \texttt{0x1ce0}$ at the $x_1$ squeeze,
comfortably below \mptr{VK\_MPTR} $=\texttt{0x3680}$ and below the
\yul{scalar\_inv} modexp frame at \texttt{0x3580}.

\subsection{Section-by-section walkthrough}
\label{sec:tp:walk}

\subsubsection{Header: digest, committed instances, instance vector}

\begin{lstlisting}[style=solidity,caption={\texttt{Halo2Verifier.sol} lines 1474--1524: transcript header (interleaved comments elided)},label={lst:tp:header}]
            let buf_len := transcript_init()
            buf_len := common_word(buf_len, mload(VK_DIGEST_MPTR))
            {
                // 128 zero bytes: zero out 4 consecutive 32-byte words
                mstore(buf_len, 0)
                mstore(add(buf_len, 0x20), 0)
                mstore(add(buf_len, 0x40), 0)
                mstore(add(buf_len, 0x60), 0)
                buf_len := add(buf_len, 0x80)
            }
            {
                buf_len := common_word(buf_len, 19)
                let instance_cptr := INSTANCE_CPTR
                for { let instance_cptr_end := add(instance_cptr, 0x0260) }
                    lt(instance_cptr, instance_cptr_end)
                    { instance_cptr := add(instance_cptr, 0x20) } {
                    let inst_be := calldataload(instance_cptr)
                    success := and(success, lt(inst_be, r))
                    buf_len := common_word(buf_len, inst_be)
                }
                if iszero(success) { fail(ERR_NON_CANONICAL_SCALAR) }
            }
\end{lstlisting}

Three points of exactness.

\emph{The digest.} \mptr{VK\_DIGEST\_MPTR} $=$ \mptr{VK\_MPTR} $=\texttt{0x3680}$,
the first word of the payload copied from the codehash-pinned VK contract. It is
absorbed as a raw big-endian word with no $<r$ check. That is correct here: the
value is fixed by \texttt{EXPECTED\_VK\_CODEHASH\_WORD} (line 94) and a prover
cannot vary it, so the absence of a canonicality check has no attacker-reachable
consequence. Note precisely \emph{what} makes it non-varying: the contract never
recomputes the digest from the VK payload -- \opcode{KECCAK256} appears in this
file only at line 802 (the transcript), line 4066 (the accumulator batching
scalar), and in comments -- so the only thing tying the absorbed word to the
commitments the later phases actually use is the \opcode{EXTCODEHASH} equality at
lines 1386--1389. See Finding~\ref{find:tp:digest}.

Also worth stating negatively, because a reader may expect otherwise: the VK
payload itself (fixed and permutation commitments, $\GT$ bases, the quotient-VM
program) is \emph{not} absorbed. Only the one digest word is. The transcript is
therefore a function of $(\text{digest}, \text{instances}, \text{proof})$ and
nothing else.

\emph{The committed-instances identity.} The native verifier, with the
\yul{committed-instances} feature, absorbs
$\texttt{committed\_pi} = \id \in \GO$ \emph{before} the instance count. Under
the patched \yul{Hashable}/\yul{Keccak256} \yul{to\_input} the identity hashes as $128$
zero bytes (the EIP-2537 $(0,0)$ convention), not as the $48$-byte ZCash
compressed form \texttt{0xc0}$\Vert 0^{47}$. The contract writes four zero words
directly rather than routing them through \yul{common\_uncompressed\_g1}, which
is equivalent (a $(0,0)$ input passes both padding and range guards) and $\approx
40$ gas cheaper. Note this is a hard-coded constant: the contract does not read a
committed-instance commitment from the proof, so there is nothing to check.

\emph{The instance loop.} The bound $\texttt{0x0260} = 19\cdot32$ matches the
generated instance count already pinned twice (VK header at line 1400, ABI length
at line 1420). Canonicality is folded into \yul{success} with a bitwise
\yul{and} rather than short-circuiting. This is safe only because both operands
are Booleans: \yul{success} is $1$ on entry (Section~\ref{sec:tp:entry}) and
\yul{lt} yields $0$ or $1$, so the bitwise \yul{and} coincides with logical
conjunction. The deferral costs nothing in soundness because line 1523 enforces
the flag \emph{before} the first \yul{squeeze\_to}: a proof with a non-canonical
public input can never observe a challenge. It buys a branch-free loop body
(one \yul{and} instead of a conditional jump per instance).

\subsubsection{Phase-1 advice and $\theta$}

\begin{lstlisting}[style=solidity,caption={\texttt{Halo2Verifier.sol} lines 1545--1580: advice, $\theta$, lookup multiplicities},label={lst:tp:advice}]
            let proof_cptr := PROOF_CPTR
            let advice_walk := ADVICE_COMMS_MPTR_BASE
            // ---- User phase 1 ----
            for { let end := add(proof_cptr, 0x0780) }
                lt(proof_cptr, end)
                {} {
                buf_len := common_uncompressed_g1(buf_len, proof_cptr)
                // Store the commitment at its phase-ordered advice slot.
                calldatacopy(advice_walk, proof_cptr, 0x80)
                advice_walk := add(advice_walk, 0x80)
                proof_cptr := add(proof_cptr, 0x80)
            }

            // ---- theta ----
            buf_len := squeeze_to(buf_len, THETA_MPTR)
            // ---- multiplicities (one G1 per lookup) ----
            let lookup_m_walk := LOOKUP_M_COMMS_MPTR_BASE
            for { let end := add(proof_cptr, 0x0100) }
                lt(proof_cptr, end)
                {} {
                buf_len := common_uncompressed_g1(buf_len, proof_cptr)
                calldatacopy(lookup_m_walk, proof_cptr, 0x80)
                lookup_m_walk := add(lookup_m_walk, 0x80)
                proof_cptr := add(proof_cptr, 0x80)
            }
\end{lstlisting}

Every $\GO$ loop in the block has this identical shape: the loop variable
\yul{end} is captured \emph{before} the body runs, the condition is
\yul{lt(proof\_cptr, end)}, the post-block is empty, and the body advances
\yul{proof\_cptr} by \texttt{0x80} exactly once. Because the section length is a
multiple of the stride, the walk is a half-open interval
$[\,c_0,\, c_0 + L\,)$ with no partial final iteration and no possibility of
overshoot. The same argument holds for the two $\Fr$ loops with stride
\texttt{0x20} and lengths \texttt{0x0cc0} and \texttt{0xa0}.

The double read of the calldata (once by \yul{common\_uncompressed\_g1} for the
transcript, once by \yul{calldatacopy} for the memory window) costs one extra
$128$-byte copy per commitment ($34$ copies total, roughly $34\times 15 = 510$
gas of copy cost) but keeps the transcript byte stream literally identical to
calldata, which is what makes the ``absorbed exactly once, verbatim'' argument in
Section~\ref{sec:tp:assess} mechanical rather than semantic.

\subsubsection{Permutation, lookups, $\tau$, trashcans, quotient}

\begin{lstlisting}[style=solidity,caption={\texttt{Halo2Verifier.sol} lines 1606--1639: the two unrolled lookup groups (comments elided)},label={lst:tp:lookups}]
            let lookup_helper_walk := LOOKUP_HELPER_COMMS_MPTR_BASE
            let lookup_z_walk := LOOKUP_Z_COMMS_MPTR_BASE
            // lookup 0: 1 helper(s) + 1 acc
            for { let end := add(proof_cptr, 0x80) }
                lt(proof_cptr, end)
                {} {
                buf_len := common_uncompressed_g1(buf_len, proof_cptr)
                calldatacopy(lookup_helper_walk, proof_cptr, 0x80)
                lookup_helper_walk := add(lookup_helper_walk, 0x80)
                proof_cptr := add(proof_cptr, 0x80)
            }
            buf_len := common_uncompressed_g1(buf_len, proof_cptr)
            calldatacopy(lookup_z_walk, proof_cptr, 0x80)
            lookup_z_walk := add(lookup_z_walk, 0x80)
            proof_cptr := add(proof_cptr, 0x80)
            // lookup 1: 1 helper(s) + 1 acc
            for { let end := add(proof_cptr, 0x80) }
                lt(proof_cptr, end)
                {} {
                buf_len := common_uncompressed_g1(buf_len, proof_cptr)
                calldatacopy(lookup_helper_walk, proof_cptr, 0x80)
                lookup_helper_walk := add(lookup_helper_walk, 0x80)
                proof_cptr := add(proof_cptr, 0x80)
            }
            buf_len := common_uncompressed_g1(buf_len, proof_cptr)
            calldatacopy(lookup_z_walk, proof_cptr, 0x80)
            lookup_z_walk := add(lookup_z_walk, 0x80)
            proof_cptr := add(proof_cptr, 0x80)
\end{lstlisting}

The per-lookup group is unrolled because the helper count is per-lookup data in
the plan; the accumulator is always exactly one $\GO$ and is emitted
straight-line. The interleaving on the wire is
$\text{helper}_0,\ \text{acc}_0,\ \text{helper}_1,\ \text{acc}_1$, but the two
walk pointers write into two disjoint contiguous windows
(\texttt{0xacc0}--\texttt{0xadc0} for helpers, \texttt{0xadc0}--\texttt{0xaec0}
for accumulators) because the quotient and PCS schedules address the two
categories independently. This costs a memory permutation relative to wire order
and buys constant-stride addressing in the two consumers.

$\tau$ (\yul{TRASH\_CHALLENGE\_MPTR}, line 1646) is squeezed
\emph{unconditionally}, even for circuits with no trash arguments, because the
native verifier does. This render does have one trashcan commitment, absorbed at
lines 1651--1658 \emph{after} $\tau$ and \emph{before} $y$, so $y$ binds it. The
squeezed $\tau$ is genuinely consumed: line 3115 reads
\mptr{TRASH\_CHALLENGE\_MPTR} into the quotient evaluator.

$y$ is squeezed at line 1663 strictly before the four quotient limb commitments
are read (lines 1674--1681), matching the Rust flow: the prover must commit to
$h(X)$ split into limbs \emph{after} learning the identity-batching challenge.

\subsubsection{$x$, the evaluations, and the \mptr{REVERSED\_EVALS\_MPTR} spill}

\begin{lstlisting}[style=solidity,caption={\texttt{Halo2Verifier.sol} lines 1686--1715: the evaluation loop},label={lst:tp:evals}]
            buf_len := squeeze_to(buf_len, X_MPTR)

            // ---- evaluations ----
            {
                let eval_buf := REVERSED_EVALS_MPTR
                for { let end := add(proof_cptr, 0x0cc0) }
                    lt(proof_cptr, end)
                    {} {
                    let eval := calldataload(proof_cptr)
                    if iszero(lt(eval, r)) { fail(ERR_NON_CANONICAL_SCALAR) }
                    // Spill for quotient numerator and PCS codegen.
                    mstore(eval_buf, eval)
                    eval_buf := add(eval_buf, 0x20)
                    buf_len := common_word(buf_len, eval)
                    proof_cptr := add(proof_cptr, 0x20)
                }
            }
\end{lstlisting}

This loop is the ``optimisation H3'' spill. Each of the $102$ evaluations is
loaded once, range-checked, written to memory, and absorbed, in a single pass.
The mapping is the identity on index:
\[
  \mathrm{mem}[\texttt{0x9480} + 32k] \;=\; e_k
  \;=\; \mathrm{calldataload}(\texttt{0x1064} + 32k),
  \qquad 0 \le k \le 101 .
\]
The spill costs $102$ \opcode{MSTORE}s ($\approx 306$ gas) plus the memory
expansion for a $\texttt{0xcc0}$-byte window; it buys $3$-gas \opcode{MLOAD} in
place of a calldata reference at every downstream use. All $102$ slots are in
fact referenced later by literal address in the quotient and PCS blocks (verified
by scanning every address in $[\texttt{0x9480}, \texttt{0xa140})$ that appears
after line 1800: $102$ of $102$ are read), so no proof evaluation is parsed and
then dropped.

Note that the failure mode here differs from the instance loop: an out-of-range
evaluation reverts \emph{immediately} with \revert{NonCanonicalScalar} rather
than being folded into \yul{success}. Both are fail-closed; the immediate form is
used wherever the value is also written to memory, so that no unchecked value can
be observed by a later phase even transiently.

\subsubsection{$x_1$, $x_2$, $\mathsf{F}$, $x_3$ and its truncation}

\begin{lstlisting}[style=solidity,caption={\texttt{Halo2Verifier.sol} lines 1720--1743: PCS batching challenges, \yul{f\_com}, and the $x_3$ mask (comments elided)},label={lst:tp:x3}]
            buf_len := squeeze_to(buf_len, X1_MPTR)
            buf_len := squeeze_to(buf_len, X2_MPTR)

            // ---- f_com (1 uncompressed G1) ----
            buf_len := common_uncompressed_g1(buf_len, proof_cptr)
            calldatacopy(F_COM_MPTR, proof_cptr, 0x80)
            proof_cptr := add(proof_cptr, 0x80)

            // ---- x3 ----
            buf_len := squeeze_to(buf_len, X3_MPTR)
            mstore(X3_MPTR, and(mload(X3_MPTR), 0xffffffffffffffffffffffffffffffff))
\end{lstlisting}

The mask keeps the low $128$ bits of $x_3$, i.e.
$x_3 \gets (h_9 \bmod r) \wedge (2^{128}-1)$. Crucially the mask is applied
\emph{after} \yul{squeeze\_to} has already reseeded the buffer with the
\emph{full} digest $h_9$, so the transcript state is unaffected and only the
challenge value the verifier arithmetic uses is truncated. This mirrors the
midnight-proofs \yul{truncated-challenges} rule for the KZG multi-open
(\texttt{proofs/src/poly/kzg/mod.rs}). Per the in-file comment at lines
1734--1742, $x_1$ and $x_4$ remain full $\Fr$ words here, but the PCS block later
stores $\mathrm{trunc}(x_1^i)$ and $\mathrm{trunc}(x_4^i)$ while keeping the power
accumulators at full precision -- visible in the fixture at lines 3819--3829,
where each \yul{x4\_pow\_i} is masked with the same $2^{128}-1$ constant while
\yul{x4\_pow\_full} is not. So this single \yul{mstore} is one instance of a
general truncation rule, not the whole rule.

\subsubsection{\yul{q\_evals}, $x_4$, $\pi$, and the terminal checks}

\begin{lstlisting}[style=solidity,caption={\texttt{Halo2Verifier.sol} lines 1753--1789: \yul{q\_evals}, $x_4$, $\pi$, and the parser's closing invariants (comments elided)},label={lst:tp:tail}]
            mstore(Q_EVAL_CPTR_MPTR, proof_cptr)
            for { let end := add(proof_cptr, 0xa0) }
                lt(proof_cptr, end)
                {} {
                let eval := calldataload(proof_cptr)
                if iszero(lt(eval, r)) { fail(ERR_NON_CANONICAL_SCALAR) }
                buf_len := common_word(buf_len, eval)
                proof_cptr := add(proof_cptr, 0x20)
            }

            // ---- x4 ----
            buf_len := squeeze_to(buf_len, X4_MPTR)

            // ---- pi (1 uncompressed G1) ----
            buf_len := common_uncompressed_g1(buf_len, proof_cptr)
            calldatacopy(PI_MPTR, proof_cptr, 0x80)
            proof_cptr := add(proof_cptr, 0x80)

            if iszero(eq(proof_cptr, NUM_INSTANCE_CPTR)) { fail(ERR_BAD_CALLDATA_SHAPE) }

            if iszero(success) { fail(ERR_NON_CANONICAL_SCALAR) }
\end{lstlisting}

The five \yul{q\_evals} are the only proof scalars \emph{not} spilled to memory.
Instead \yul{proof\_cptr} is saved to \mptr{Q\_EVAL\_CPTR\_MPTR} (\texttt{0x9220})
before the loop, holding the constant-derived value \texttt{0x1da4}, and the PCS
block re-reads them from calldata at lines 3648 and 3818 as
$\mathrm{calldataload}(\texttt{Q\_EVAL\_CPTR} + 32k)$ for $k \in \{0,\dots,4\}$.
Calldata is immutable within a call frame, so the words re-read are bit-for-bit
the words that were range-checked and absorbed here; the saved-cursor idiom
introduces no aliasing. It buys the $5$ \opcode{MSTORE}s and $160$ bytes of memory
that the spill would have cost, at the price of a slightly less local
data-flow argument.

The two closing guards are the parser's completeness statement.
\yul{eq(proof\_cptr, NUM\_INSTANCE\_CPTR)} asserts that the cursor -- advanced
once per element, by that element's exact width -- landed on the word immediately
after the dynamic \yul{bytes} payload. Combined with the entry check
$\texttt{calldataload(PROOF\_LEN\_CPTR)} = \texttt{0x1e60}$, this pins the walk to
the interval $[\texttt{0x64}, \texttt{0x1ec4})$ with no slack at either end. The
in-file comment correctly calls this redundant with the generated proof length
\emph{today}; its value is that a future emitter that changes a section count
without changing the length constant fails closed rather than silently
misinterpreting the tail. The final \yul{iszero(success)} re-check is likewise
redundant with line 1523 -- \yul{success} is not modified between the two -- and
keeps the parser's fail-closed property a local property of the block.

\subsection{Assessment}
\label{sec:tp:assess}

\subsubsection{Is every proof byte absorbed exactly once, in the native order?}

\textbf{Yes.} The argument is mechanical.

\emph{Coverage and disjointness.} There are $13$ sections, each a half-open
calldata interval walked with a constant stride that divides its length
(Section~\ref{sec:tp:walk}). The sections are laid out back to back starting at
\mptr{PROOF\_CPTR}: each section's start is the previous section's \yul{end},
because every loop and every straight-line read advances the single shared cursor
\yul{proof\_cptr} and nothing else ever assigns to it. Summing
Table~\ref{tab:tp:layout} gives \texttt{0x1e60}, and line 1785 asserts the cursor
ends at \texttt{0x64}$+$\texttt{0x1e60}. Hence the union of the $13$ intervals is
exactly $[\texttt{0x64}, \texttt{0x1ec4})$ and the intervals are pairwise
disjoint: every proof byte is read exactly once.

\emph{Absorbed verbatim.} $\GO$ elements are appended with
\yul{calldatacopy(buf\_len, cptr, 0x80)} (line 794), a byte-for-byte copy.
$\Fr$ elements are appended with \yul{mstore(buf\_len, calldataload(c))}, which
for a $32$-byte-aligned read and a $32$-byte store is also byte-for-byte. There
is no normalisation, no reduction, and no re-encoding anywhere in the path, so
the transcript input over the proof region is the concatenation of the proof
bytes in cursor order. The only non-proof material in the transcript is the
header (digest, $0^{128}$, the count $19$, the $19$ instance words), all of it
prepended before the first proof byte.

\emph{Order.} The interleaving of the $13$ sections with the $10$ squeezes is
Algorithm~\ref{alg:tp:schedule}. That this is the native order is not something
the fixture can prove about itself; it is pinned generator-side by
\texttt{src/lowering/transcript\_anchor.rs} (\yul{FIXED\_CHALLENGES = 10} at line
50; the read-sequence and squeeze-count assertions at lines 258--273), which
implements midnight-proofs'
\yul{Transcript} trait over the real \yul{CircuitTranscript}, drives the real
\yul{prepare} over a proof synthesised from the generator's own plan, and records
every \yul{read}/\yul{common}/\yul{squeeze} the native verifier performs. A plan
that disagreed about counts or order would desynchronise there. This
specification reports the fixture's schedule; the equality with the native
schedule rests on that anchor plus the trace differential.

\subsubsection{Is anything read but never absorbed, or absorbed but never checked?}

\textbf{Neither.} Table~\ref{tab:tp:checks} enumerates every value that reaches
the transcript.

\begin{table}[t]
\centering\small
\caption{Canonicality coverage of transcript inputs.}
\label{tab:tp:checks}
\begin{tabular}{@{}llll@{}}
\toprule
Input & Source & Check & Line \\
\midrule
\yul{vk\_digest}       & memory (pinned VK) & none needed -- fixed by codehash & 1481 \\
\yul{committed\_pi}    & literal $0^{128}$  & none needed -- constant          & 1496--1500 \\
instance count $19$    & literal            & cross-checked vs.\ ABI + VK      & 1400, 1420, 1508 \\
$19$ instances         & calldata           & $<r$, deferred then enforced     & 1518, 1523 \\
$34 \times \GO$        & calldata           & padding $=0$ and $x,y<p$         & 780--790 \\
$102$ evaluations      & calldata           & $<r$, immediate revert           & 1706 \\
$5$ \yul{q\_evals}     & calldata           & $<r$, immediate revert           & 1759 \\
\bottomrule
\end{tabular}
\end{table}

Conversely, nothing is read without being absorbed: every
\opcode{CALLDATALOAD}/\opcode{CALLDATACOPY} of proof bytes in the range is
immediately paired with a \yul{common\_*} call. The $19$ instance words are also
read \emph{before} the transcript by the accumulator validator (line 1257 reads
from $\texttt{INSTANCE\_CPTR}+\texttt{0x0160}$, i.e.\ instance index $11$, which
matches the pinned $\texttt{acc\_offset}=11$ and $\texttt{num\_acc\_limbs}=7$),
but those same words are all absorbed here as part of the $19$, so the
accumulator path adds no unabsorbed input.

\begin{finding}[TP-1]{Absorption is complete, exactly-once, and fully range-checked}
\textbf{Severity: Informational (clean).}
\texttt{fixtures/moonlight-wrap/Halo2Verifier.sol:1474--1789}. The parser walks
$[\texttt{0x64},\texttt{0x1ec4})$ with a single monotone cursor whose stride
divides every section length and whose terminal value is asserted equal to
\mptr{NUM\_INSTANCE\_CPTR}; every byte in that interval is absorbed verbatim
exactly once; every calldata-sourced scalar is checked $<r$ and every $\GO$
coordinate is checked $<p$ with zero EIP-2537 padding before absorption; and both
deferred-failure paths are enforced before any challenge is squeezed. No aliasing
of transcript inputs is reachable. This region is sound as written.
\end{finding}

\begin{finding}[TP-2]{Curve and subgroup membership are not checked at parse time}
\label{find:tp:curve}
\textbf{Severity: Informational.}
\texttt{Halo2Verifier.sol:775--796} deliberately omits an on-curve and
$r$-torsion check; the header comment (lines 766--770) states the invariant it
relies on: \yul{ProtocolPlan::validate} rejects generated plans in which an
absorbed proof commitment is not later consumed by an EIP-2537 \opcode{G1MSM} or
pairing, and those precompiles perform the validation.

The invariant was checked twice: mechanically against this fixture, and against
the generator that is supposed to enforce it.

\emph{In the fixture.} All $34$ absorbed $\GO$ elements reach a validating
precompile. The final MSM staging buffer at \mptr{PCS\_FINAL\_MSM\_MPTR}
($=\texttt{0xb280}$) is $78$ terms of
$[\,\text{point } \texttt{0x80} \Vert \text{scalar } \texttt{0x20}\,]$, stride
\texttt{0xa0}, total $78\cdot\texttt{0xa0} = \texttt{0x30c0}$, filled by lines
3836--3996 and consumed by the single
\yul{staticcall(525096, 0x0c, PCS\_FINAL\_MSM\_MPTR, 0x30c0, FINAL\_COM\_MPTR, 0x80)}
at line 3999, whose \opcode{RETURNDATASIZE} is checked against \texttt{0x80} at
line 4000. Enumerating every \yul{mcopy} source in that block gives exactly the
$15$ advice bases $\texttt{0xa140}+\texttt{0x80}j$, the $2$ multiplicity bases,
the $6$ permutation-$Z$ bases, the $2$ lookup-helper bases, the $2$
lookup-accumulator bases and the trashcan base -- $28$ of the $34$. The $4$
quotient limbs are staged as terms $41$--$44$, but by \emph{absolute} address:
\yul{mcopy(0xcc20, add(QUOTIENT\_LIMB\_COMMS\_MPTR\_BASE, 0x0), 0x80)} and its
three siblings at \texttt{0xccc0}, \texttt{0xcd60}, \texttt{0xce00} (lines 3920,
3923, 3926, 3929), which are $\texttt{0xb280}+\texttt{0x19a0}$,
$+\texttt{0x1a40}$, $+\texttt{0x1ae0}$, $+\texttt{0x1b80}$ -- the point slots
whose scalars are written at $+\texttt{0x1a20}$, $+\texttt{0x1ac0}$,
$+\texttt{0x1b60}$, $+\texttt{0x1c00}$ on the next line each time. A reader
grepping for \mptr{PCS\_FINAL\_MSM\_MPTR} will miss these four; that is worth
knowing, and it is the one place in this argument where the addressing is not
self-describing. $\mathsf{F}$ is term $78$, staged at
$\texttt{PCS\_FINAL\_MSM\_MPTR}+\texttt{0x3020}$ (line 3995). $\pi$ is copied to
\mptr{PAIRING\_LHS\_MPTR} (line 4011) and separately fed to a one-pair
\opcode{G1MSM} at lines 4023--4026 whose return size is also checked. The MSM
input length \texttt{0x30c0} is a compile-time constant with no conditional term
omission, so no commitment can be skipped at runtime by choosing a zero scalar.

\emph{In the generator.} \yul{ProtocolPlan::validate}
(\texttt{src/lowering/protocol/mod.rs:697--918}) does contain the coverage
checks the contract comment advertises, and they are per-category rather than
generic: advice columns absorbed but never opened by PCS are rejected by name
(lines 782--795), and the counts of \yul{PcsQuerySource::LookupMultiplicity},
\yul{LookupHelper}, \yul{LookupAccumulator}, \yul{PermutationZ} and \yul{Trash}
queries are each required to equal the corresponding absorbed commitment count
(lines 798--864). The quotient limbs are covered by the requirement that the PCS
query schedule end with a \yul{Linearization} query (lines 778--780). Two caveats
a reviewer should carry: (i) what \yul{validate} enforces is that every absorbed
commitment is \emph{queried by the PCS plan}, not that the rendered Yul actually
stages every queried commitment into an MSM -- the step from plan to emitted
\yul{mcopy} is trusted, which is why the fixture-level enumeration above is not
redundant; and (ii) the argument assumes EIP-2537 \opcode{G1MSM} validates every
input point regardless of its paired scalar, including a zero scalar. The
constructor smoke test does exercise the sharper half of (ii): lines 486--500
feed \opcode{G1MSM} the point $(4, y)$, which satisfies $y^2 = x^3 + 4$ over
$\Fp$ but lies outside the $r$-order subgroup, with gas bounded at $200000$, and
requires the call to \emph{fail}. It uses scalar $1$ (line 499), so the
zero-scalar case is untested and (ii) rests there on the EIP text rather than on
a test in this repository.

What this buys: no ladder, no $y^2 = x^3 + 4$ evaluation, and no cofactor
clearing in Yul -- a large gas and bytecode saving over $34$ in-contract curve
checks. What it costs: a non-local invariant. A future emitter change that
absorbs a commitment without routing it to a precompile would accept off-curve or
small-order points into the transcript. This is a standing property to re-verify
on every regeneration, not a defect in the present artifact.

A related observation: the on-chain calldata form is the $128$-byte uncompressed
encoding produced by an off-chain shim from the native $48$-byte compressed
stream. Because compression is a bijection on valid $\GO$ points (including
$\id$, compressed \texttt{0xc0}$\Vert 0^{47}$ vs.\ padded $(0,0)$), the accepted
sets coincide once the precompile curve check is applied -- the only inputs the
contract accepts that have no compressed preimage are exactly the ones the MSM
rejects.
\end{finding}

\begin{finding}[TP-3]{$x_3$ is truncated to 128 bits}
\textbf{Severity: Informational.}
\texttt{Halo2Verifier.sol:1743} masks the PCS evaluation point to
$x_3 \in [0, 2^{128})$. This is intentional and mirrors midnight-proofs'
\yul{truncated-challenges}; matching it is a hard \emph{requirement}, since a
verifier that did not truncate would reject every honest proof. On the EVM the
mask is not an optimisation -- \opcode{MULMOD} costs the same for a $128$-bit
operand -- so the cost is one-sided: the Schwartz--Zippel error for the batched
opening identity rises from $\deg/r \approx 2^{-235}$ to $\deg/2^{128}$, i.e.\
about $2^{-108}$ for a degree-$2^{20}$ domain. That remains far above any
practical threshold, but it is the tightest soundness margin the transcript
introduces and should be recorded as such. Note also that $x_3 = 0$ is reachable
with probability $2^{-128}$; the code does not special-case it here, and the
consequence is handled by the batch-inversion rejection path in the PCS block
(covered elsewhere).
\end{finding}

\begin{finding}[TP-4]{Challenges are sampled by modular reduction of a 256-bit digest}
\label{find:tp:sampling}
\textbf{Severity: Informational.}
\texttt{Halo2Verifier.sol:807} computes $c = h \bmod r$ from a full $256$-bit
Keccak digest. Write $2^{256} = 2r + s$; then $\lfloor 2^{256}/r \rfloor = 2$ and
$s = 2^{256} \bmod r \approx 0.0943 \cdot 2^{256}$. The $s$ smallest residues have
three digest preimages and the remaining $r - s$ have two, so the biased slice is
$s/r \approx 20.8\%$ of $\Fr$ (equivalently $9.43\%$ of digest space). The
resulting statistical distance from uniform is
\[
  \tfrac{1}{2}\Bigl(s\bigl|\tfrac{3}{2^{256}}-\tfrac{1}{r}\bigr|
  + (r-s)\bigl|\tfrac{2}{2^{256}}-\tfrac{1}{r}\bigr|\Bigr)
  \;\approx\; 0.0747 ,
\]
and the extreme likelihood ratios are $3r/2^{256} \approx 1.359$ (most likely) and
$2r/2^{256} \approx 0.906$ (least likely) -- i.e.\ at most $0.442$ bits of
advantage on any fixed target, not the $3/2$ that the three-versus-two preimage
count suggests at a glance. For Fiat--Shamir soundness this multiplies an
adversary's success probability against any fixed target by at most $1.36$, which
is negligible against the margins above. It is recorded because $0.0747$ is a
large statistical distance in absolute terms and a reader who sees only that
number could mistake it for a problem.

The rule must in any case match the native implementation exactly, and it does:
\texttt{proofs/src/transcript/implementors.rs:19--40} computes
\yul{BigUint::from\_bytes\_be(digest) \% modulus} and asserts the result is a
canonical repr, with the \yul{Sampleable} impl for \yul{Keccak256}/\yul{Fq} (lines 337--341)
delegating to it. Notably the \texttt{Blake2b} instances in the same file (lines
119--124, 171--177) use \yul{from\_uniform\_bytes} over $64$ bytes instead, whose
bias is $\approx 2^{-256}$; the Keccak path deliberately does not, and the
contract must mirror the Keccak path.
\end{finding}

\begin{finding}[TP-5]{$\pi$'s transcript absorption is a no-op; only its encoding check is load-bearing}
\textbf{Severity: Observation.}
\texttt{Halo2Verifier.sol:1773} calls \yul{common\_uncompressed\_g1} on $\pi$,
but no \yul{squeeze\_to} follows (the last squeeze is $x_4$ at line 1767) and
\yul{buf\_len} is dead from line 1775 onward (confirmed: \yul{buf\_len} does not
appear again anywhere after line 1775). The $128$ bytes are copied into a buffer
nothing ever hashes -- $15$ gas of \opcode{CALLDATACOPY}, with no memory
expansion, since the buffer was already grown to \texttt{0x1ce0} at the $x_1$
squeeze.

This is harmless and matches the native schedule, but it creates a trap: the call
is \emph{not} dead code, because it is the only place $\pi$'s EIP-2537 padding
and coordinate range are validated before $\pi$ is used as a pairing argument at
line 4011. An optimiser -- human or generated -- that removed the ``useless''
absorb would silently remove that validation. Worth an explicit comment at the
call site.
\end{finding}

\begin{finding}[TP-6]{\texttt{REVERSED\_EVALS\_MPTR} performs no reversal at runtime}
\textbf{Severity: Observation.}
\texttt{Halo2Verifier.sol:1698--1714} writes $e_k$ to
$\texttt{REVERSED\_EVALS\_MPTR} + 32k$ in ascending proof order; the mapping is
the identity on index. The name refers to the \emph{generator-side} choice of
which query each proof slot carries (the lowering plan assigns evaluation indices
so that the quotient VM's reverse fold reads them in ascending address order), as
the in-file comment at lines 1694--1696 explains. A reader auditing the runtime
will look for a reversal that does not exist. Purely a naming/documentation
observation; no behavioural impact.
\end{finding}

\begin{finding}[TP-7]{\mptr{CHALLENGE\_MPTR} and \mptr{THETA\_MPTR} coincide in this render}
\textbf{Severity: Observation.}
\texttt{Halo2Verifier.sol:144} and \texttt{150} both resolve to \texttt{0x7900}.
That is correct here: this circuit declares zero user-phase challenges, so the
user-challenge window is empty and no \yul{squeeze\_to(buf\_len, add(CHALLENGE\_MPTR, ...))}
site is rendered in the assigned range. The aliasing would matter only for a
circuit \emph{with} user challenges and an undersized window, in which case
user challenge $0$ would be overwritten by $\theta$ while the transcript still
matched the prover. The generator already treats this as a tracked hazard (P8/L-5)
and pins it: \texttt{src/lowering/layout/memory.rs:814--842} carries the P8/L-5
comment and two assertions -- \yul{challenge\_indices.len()} must equal the phase
sum of \yul{num\_user\_challenges}, and
$\texttt{challenge\_start} + 32\cdot\texttt{user\_challenges\_total} \le
\texttt{theta\_start}$ -- with the regression test
\yul{undersized\_user\_challenge\_window\_is\_rejected} at
\texttt{memory.rs:1631}. In this render \texttt{CHALLENGE\_MPTR} is declared at
line 144 and then referenced nowhere else in the entire contract, which is the
observable signature of the empty window. Recorded so a reader of the constants
block does not read the shared address as a bug.
\end{finding}

\begin{finding}[TP-8]{The transcript buffer's upper bound is generator-enforced, not runtime-checked}
\textbf{Severity: Observation.}
The buffer grows from \mptr{TRANSCRIPT\_MPTR} $=\texttt{0x1000}$ toward
\mptr{VK\_MPTR} $=\texttt{0x3680}$ with no runtime bound test. For this fixture
the high-water mark is $\texttt{0x1000}+3296=\texttt{0x1ce0}$ at the $x_1$
squeeze (Table~\ref{tab:tp:buffer}), leaving
$\texttt{0x3680}-\texttt{0x1ce0} = 6560$ bytes of headroom (and $6304$ bytes to
the \yul{scalar\_inv} modexp frame at \texttt{0x3580}, which is the nearer
neighbour), and every
contributing length is a compile-time constant, so the bound holds by
construction. The guard at lines 674--677 protects the \emph{lower} boundary (solc's
spill region reaching \mptr{TRANSCRIPT\_MPTR}) and reverts with
\revert{MemoryLayoutViolated}; there is no symmetric upper-boundary guard.
Overrunning into \mptr{VK\_MPTR} would corrupt the VK digest and commitments
after they were already absorbed -- a fail-open-shaped failure rather than a
fail-closed one -- so the invariant is worth keeping explicit even though it is
not reachable in this render. Nothing to change in the fixture.
\end{finding}

\begin{finding}[TP-9]{The absorbed VK digest is never checked against the VK it labels}
\label{find:tp:digest}
\textbf{Severity: Observation.}
Line 1481 absorbs \yul{mload(VK\_DIGEST\_MPTR)} as the transcript's first input,
under the comment ``this digest commits to the verifier key / constraint system''
(lines 1479--1480). It does not: it is simply the first word of the payload the
VK contract stores, and no code in this file recomputes it. The only
\opcode{KECCAK256} sites in \texttt{Halo2Verifier.sol} are line 802
(\yul{squeeze\_to}) and line 4066 (the accumulator pairing batching scalar);
neither hashes \mptr{VK\_MPTR}. A \texttt{Halo2VerifyingKey} whose stored digest
word disagreed with its own commitments would be accepted without complaint.

What actually closes the gap is the \opcode{EXTCODESIZE}/\opcode{EXTCODEHASH}
pair at lines 1386--1389, re-run on every call, which fixes the \emph{whole}
runtime -- digest word and commitments together -- to one generated artifact. So
the property the comment claims does hold in this render, but for a different
reason than the comment gives, and by a mechanism that lives $95$ lines away and
outside the transcript block. The distinction matters for two foreseeable
changes: an embedded-VK render (where the payload is bytecode-resident and the
codehash check has no analogue at this granularity) and any future upgradeable-VK
variant. In both, absorbing an unverified digest would become the load-bearing
step rather than a cheap label. Nothing to change in this artifact; the comment
at lines 1479--1480 overclaims and would be better phrased as ``pinned by
\texttt{EXPECTED\_VK\_CODEHASH}''.
\end{finding}

\subsubsection{Summary of the region}

The parser is tight. The three things a reviewer most wants -- exactly-once byte
coverage, canonical decoding of everything absorbed, and enforcement of both
deferred-failure flags before the first challenge is observable -- all hold, and
they hold for structural reasons (a single monotone cursor, verbatim copies,
compile-time-constant strides) rather than by coincidence of the constants. The
two genuine non-localities are the deferral of curve and subgroup validation to
the EIP-2537 precompiles (Finding~\ref{find:tp:curve}), which is verified to hold
for all $34$ commitments in this artifact and is backed by per-category coverage
assertions in \yul{ProtocolPlan::validate}, and the reliance on generator-side
anchoring for the claim that this schedule \emph{is} the native schedule. The
transcript \emph{primitives} no longer need that trust: they were read directly
out of \texttt{proofs/src/transcript/implementors.rs} and match the contract
line for line (Section~\ref{sec:tp:helpers}).


%% ==== 10-lagrange.tex ===================================================
\section{Domain evaluation: $x^n$, Lagrange values, instance evaluation, and the batch inversion}
\label{sec:lag:top}

This section specifies \texttt{Halo2Verifier.sol} lines 1792--1917: the block that turns
the squeezed evaluation challenge $x$ into every domain-dependent scalar the rest of the
verifier consumes -- $x^n$, $(x^n-1)^{-1}$, $L_0(x)$, $L_{\mathrm{last}}(x)$,
$L_{\mathrm{blind}}(x)$, and the public-instance evaluation $I(x)$ -- using a single
\texttt{modexp} call. It also specifies the two quotient-VM helper functions
\yul{q\_program\_fail} and \yul{q\_pow5}, which the generator renders at the tail of this
region (lines 1885--1915) even though both are called only from the quotient interpreter
covered elsewhere.

The whole block is pure $\Fr$ arithmetic. It reads no calldata except the 19 already
validated public-instance words, performs exactly one precompile call, and writes six
named memory slots.

\subsection{Entry preconditions}
\label{sec:lag:pre}

Control reaches line 1792 only after two unconditional guards immediately above it:

\begin{lstlisting}[style=solidity,caption={\texttt{Halo2Verifier.sol} lines 1785--1789: the guards immediately preceding this block},label={lst:lag:pre}]
            if iszero(eq(proof_cptr, NUM_INSTANCE_CPTR)) { fail(ERR_BAD_CALLDATA_SHAPE) }

            // `success` carries deferred canonicality failures from public
            // instance reads. G1/proof scalar helpers revert immediately.
            if iszero(success) { fail(ERR_NON_CANONICAL_SCALAR) }
\end{lstlisting}

Consequently, on entry:

\begin{enumerate}
  \item $\yul{success} = 1$. This matters because \yul{batch\_invert} (lines 836--838)
        opens with \yul{ret := success} followed by \yul{if iszero(ret) \{ leave \}}; that
        early-out is provably not taken here, because nothing between line 1789 and line
        1826 assigns to \yul{success}, so the helper always performs real work.
  \item $x$, read at line 1801 via \yul{mload(X\_MPTR)}, was written exactly once -- at
        line 1686, the sole \yul{squeeze\_to(buf\_len, X\_MPTR)} call site and the only
        write to \mptr{X\_MPTR} anywhere in the contract. Line 807 of \yul{squeeze\_to}
        stores \yul{mod(h0, r)}, so $0 \le x < r$ by construction. This block performs no
        further canonicality check on $x$, and none is needed.
  \item All 19 instance words at calldata offsets $\texttt{0x1ee4}$ through
        $\texttt{0x2124}$ were range-checked against $r$ at line 1518
        (\yul{success := and(success, lt(inst\_be, r))}) and enforced at line 1523, and \yul{calldatasize()} was pinned to
        $\mptr{INSTANCE\_CPTR} + \texttt{0x0260} = \texttt{0x2144}$ at line 1426. The
        instance loop at line 1860 therefore re-reads exactly the same, already
        validated, words.
  \item The verifying key was loaded by \yul{extcodecopy} at line 1393 after an
        \yul{extcodehash} pin (line 1388), and its header words were cross-checked at
        lines 1400--1405, including \yul{eq(mload(K\_MPTR), 20)}.
\end{enumerate}

\subsection{Constants and memory used}
\label{sec:lag:consts}

\begin{table}[h]
\centering
\begin{tabular}{@{}llll@{}}
\toprule
Constant & Address & Words & Role in this block \\
\midrule
\mptr{X\_MPTR}                & \texttt{0x79a0} & 1  & challenge $x$ (read) \\
\mptr{OMEGA\_MPTR}            & \texttt{0x3700} & 1  & $\omega$ (VK, read) \\
\mptr{OMEGA\_INV\_TO\_L\_MPTR}& \texttt{0x3740} & 1  & $\omega^{-10}$ (VK, read) \\
\mptr{N\_INV\_MPTR}           & \texttt{0x36e0} & 1  & $n^{-1}$ (VK, read) \\
\mptr{LAGRANGE\_DENOMS\_MPTR} & \texttt{0xb580} & 30 & denominators $\to$ inverses $\to$ $L_i(x)$ \\
\mptr{BATCH\_INV\_SCRATCH\_MPTR}& \texttt{0xb140} & 34 & prefix products + one modexp frame \\
\mptr{INSTANCE\_CPTR}         & \texttt{0x1ee4} & 19 & calldata, public inputs \\
\midrule
\mptr{X\_N\_MPTR}             & \texttt{0x7c40} & 1  & written: $x^n$ (dead store) \\
\mptr{X\_N\_MINUS\_1\_INV\_MPTR}& \texttt{0x7c60} & 1 & written: $(x^n-1)^{-1}$ (dead store) \\
\mptr{L\_LAST\_MPTR}          & \texttt{0x7c80} & 1  & written: $L_{\mathrm{last}}(x)$ \\
\mptr{L\_BLIND\_MPTR}         & \texttt{0x7ca0} & 1  & written: $L_{\mathrm{blind}}(x)$ \\
\mptr{L\_0\_MPTR}             & \texttt{0x7cc0} & 1  & written: $L_0(x)$ \\
\mptr{INSTANCE\_EVAL\_MPTR}   & \texttt{0x7ce0} & 1  & written: $I(x)$ \\
\bottomrule
\end{tabular}
\caption{Memory constants referenced by the Lagrange block. Definitions at
\texttt{Halo2Verifier.sol} lines 103, 132--135, 155, 170--175, 235, 239.}
\label{tab:lag:mem}
\end{table}

The scalar modulus is bound once for the whole \yul{verifyProof} assembly block as
\yul{let r := FR\_MODULUS} at line 1355, with
\[
r = \texttt{0x73eda753299d7d483339d80809a1d80553bda402fffe5bfeffffffff00000001}.
\]

The four domain parameters come from the pinned VK payload
(\texttt{Halo2VerifyingKey.sol} lines 70--74). Their concrete values, and the algebraic
identities this specification verified against them, are:

\begin{table}[h]
\centering
\begin{tabular}{@{}lll@{}}
\toprule
VK word & Value & Verified identity \\
\midrule
$k$              & \texttt{0x14} $= 20$ & matches the literal \yul{k := 20} at line 1800 \\
$n^{-1}$         & \texttt{0x73eda014...1a400fff00001001} & $n^{-1}\cdot 2^{20} \equiv 1$ \\
$\omega$         & \texttt{0x03e1c54b...fece7e0e39898d4b} & $\omega^{2^{20}}=1$, $\omega^{2^{19}}\neq 1$ \\
$\omega^{-1}$    & \texttt{0x6c39442e...ef97c3573a28fc8c} & $\omega\cdot\omega^{-1}\equiv 1$ \\
$\omega^{-L}$    & \texttt{0x2a0ccbaa...9e5604b8d6b68d45} & equals $(\omega^{-1})^{10}$, so $L=10$ \\
\bottomrule
\end{tabular}
\caption{Domain constants from the pinned VK, with the identities checked numerically
against $r$ while writing this section.}
\label{tab:lag:vk}
\end{table}

Hence $n = 2^{20} = 1{,}048{,}576$, the last-row rotation magnitude is $L = 10$, and the
number of blinding rows is $L-1 = 9$.

\subsection{Step 1: $x^n$ by repeated squaring}
\label{sec:lag:xn}

\begin{lstlisting}[style=solidity,caption={\texttt{Halo2Verifier.sol} lines 1798--1826: $x^n$, the denominator pass, and the batch-inversion call},label={lst:lag:pass1}]
            let lagrange_precompile_failed := 0
            {
                let k := 20
                let x := mload(X_MPTR)
                // Compute x^n by repeated squaring, with n = 2^k.
                let x_n := x
                for { let idx := 0 } lt(idx, k) { idx := add(idx, 1) } {
                    x_n := mulmod(x_n, x_n, r)
                }

                let omega := mload(OMEGA_MPTR)

                // First pass writes denominators (x - omega_i) for every
                // Lagrange value needed below, then appends x^n - 1. The
                // batch inversion pass turns all of them into inverses in one
                // modexp call. The run lives in the dedicated planner-registered
                // LAGRANGE_DENOMS_MPTR scratch region; only the distilled
                // results below are persisted into the named theta slots.
                let mptr := LAGRANGE_DENOMS_MPTR
                let mptr_end := add(mptr, 0x03a0)
                for { let pow_of_omega := mload(OMEGA_INV_TO_L_MPTR) }
                    lt(mptr, mptr_end)
                    { mptr := add(mptr, 0x20) } {
                    mstore(mptr, addmod(x, sub(r, pow_of_omega), r))
                    pow_of_omega := mulmod(pow_of_omega, omega, r)
                }
                let x_n_minus_1 := addmod(x_n, sub(r, 1), r)
                mstore(mptr_end, x_n_minus_1)
                success, lagrange_precompile_failed := batch_invert(success, LAGRANGE_DENOMS_MPTR, add(mptr_end, 0x20), BATCH_INV_SCRATCH_MPTR, r)
\end{lstlisting}

The loop at lines 1804--1806 initialises $\yul{x\_n} \leftarrow x$ and squares it
\emph{exactly} $k = 20$ times, so
\[
\yul{x\_n} \;=\; x^{2^{k}} \;=\; x^{2^{20}} \;=\; x^{n}, \qquad n = 2^{20}.
\]
This is correct: the loop body is $\yul{x\_n} \leftarrow \yul{x\_n}^2$, applied $k$ times
to the seed $x$, giving $x^{2^k}$, not $x^{2^{k}-1}$ or $x^{2^{k+1}}$. Twenty
\opcode{MULMOD}s replace a 20-bit exponentiation ladder; there is no branch and no data
dependence on $x$, so the cost is constant.

\yul{k} is rendered as the decimal literal \yul{20} rather than loaded from
\mptr{K\_MPTR}. That is a codegen specialisation, not an unchecked assumption: line 1401
asserts \yul{eq(mload(K\_MPTR), 20)} against the copied VK payload and fails with
\revert{VkMismatch}, and \mptr{EXPECTED\_VK\_CODEHASH} (line 95, enforced at line 1388)
pins the entire VK runtime. The literal buys one \opcode{MLOAD} and lets \texttt{solc}
keep the bound in a stack slot; it costs a second place that must change when $k$ changes,
which the generator and the line-1401 check together cover.

\subsection{Step 2: the denominator run and its index-to-rotation map}
\label{sec:lag:denoms}

The first pass (lines 1816--1823) fills a contiguous run of 29 words starting at
$\mptr{LAGRANGE\_DENOMS\_MPTR} = \texttt{0xb580}$, with
$\yul{mptr\_end} = \texttt{0xb580} + \texttt{0x03a0} = \texttt{0xb920}$ and
$\texttt{0x03a0} = 928 = 29 \times 32$. The running variable \yul{pow\_of\_omega} is
seeded from $\mptr{OMEGA\_INV\_TO\_L\_MPTR} = \omega^{-10}$ and multiplied by $\omega$
after each store, so word $j$ (for $j = 0,\dots,28$) receives
\[
d_j \;=\; x - \omega^{\,j-10} \pmod r ,
\]
computed as \yul{addmod(x, sub(r, pow\_of\_omega), r)}. Note that \yul{sub} here is the
plain 256-bit subtraction: since $\yul{pow\_of\_omega} < r$ always (it starts at a VK word
$< r$ and is thereafter produced by \opcode{MULMOD}), $r - \yul{pow\_of\_omega}$ is the
correct additive inverse and the \opcode{ADDMOD} reduces the sum. No underflow is
possible.

Line 1824 then forms $\yul{x\_n\_minus\_1} = x^n - 1 \bmod r$ and line 1825 stores it at
\yul{mptr\_end}, i.e.\ as word 29 -- \emph{one past} the denominator run. The batch
inversion at line 1826 is therefore invoked over $[\texttt{0xb580}, \texttt{0xb940})$, a
window of $\texttt{0x03c0} = 960$ bytes $= 30$ words.

\begin{table}[h]
\centering
\begin{tabularx}{\textwidth}{@{}llllX@{}}
\toprule
Index $j$ & Address & Denominator & Rotation & Consumer of the resulting $L$ value \\
\midrule
$0$        & \texttt{0xb580} & $x-\omega^{-10}$ & $-10$      & $L_{\mathrm{last}}$, read at line 1866 \\
$1\dots 9$ & \texttt{0xb5a0}--\texttt{0xb6a0} & $x-\omega^{-10+j}$ & $-9\dots-1$ & the nine $L_{\mathrm{blind}}$ summands, lines 1841--1847 \\
$10$       & \texttt{0xb6c0} & $x-1$            & $0$        & $L_0$ (line 1867) \emph{and} the coefficient of instance $0$ \\
$11\dots28$& \texttt{0xb6e0}--\texttt{0xb900} & $x-\omega^{\,j-10}$ & $1\dots 18$ & coefficients of instances $1\dots18$ \\
$29$       & \texttt{0xb920} & $x^n-1$          & ---        & $(x^n-1)^{-1}$, read at line 1865 \\
\bottomrule
\end{tabularx}
\caption{Layout of the 30-word batch-inversion window at
\mptr{LAGRANGE\_DENOMS\_MPTR}. Index $j$ corresponds to domain row $j-10 \bmod n$.}
\label{tab:lag:window}
\end{table}

The run length is exactly what is consumed: $1$ (last row) $+\;9$ (blinding rows)
$+\;19$ (instance rows, of which row $0$ doubles as $L_0$) $= 29$, plus the vanishing
term. Word 10 is deliberately shared between $L_0$ and the first instance coefficient
rather than duplicated -- the correct choice, since duplicating it would add a word to the
inversion window and a prefix product to the scratch region.

\subsection{Step 3: the single batch inversion}
\label{sec:lag:invert}

\yul{batch\_invert} (lines 836--951) is specified in full elsewhere; what matters here is
its memory geometry and its two-valued failure report. For an $N$-word window it stores
$N-2$ prefix products at \yul{scratch\_mptr} and then overlays one $\texttt{0xc0}$-byte
\texttt{modexp} frame at the current prefix pointer:

\begin{lstlisting}[style=solidity,caption={\texttt{Halo2Verifier.sol} lines 917--933: the single modexp inside \yul{batch\_invert} (forward and backward passes elided)},label={lst:lag:modexp}]
                // Invert the total product once.
                mstore(add(gp_mptr, 0x00), 0x20)
                mstore(add(gp_mptr, 0x20), 0x20)
                mstore(add(gp_mptr, 0x40), 0x20)
                mstore(add(gp_mptr, 0x60), gp)
                mstore(add(gp_mptr, 0x80), sub(r, 2))
                mstore(add(gp_mptr, 0xa0), r)
                ret := staticcall(MODEXP_GAS, 0x05, gp_mptr, 0xc0, gp_mptr, 0x20)
                ret := and(ret, eq(returndatasize(), 0x20))
                precompile_failed := iszero(ret)
                // Leave before the backward pass on a failed modexp. A failed
                // staticcall writes no output, so `mload(gp_mptr)` would read
                // back the stale frame header and the pass below would
                // overwrite every denominator in [mptr_start, mptr_end) with
                // garbage products before returning ret = 0.
                if iszero(ret) { leave }
                let all_inv := mload(gp_mptr)
\end{lstlisting}

With $N = 30$, the forward loop (line 893) runs while
$\yul{mptr} < \yul{mptr\_end} - \texttt{0x20}$, i.e.\ $28$ iterations, advancing
\yul{gp\_mptr} by one word each. The resulting geometry is:

\begin{table}[h]
\centering
\begin{tabular}{@{}lll@{}}
\toprule
Sub-region & Range & Size \\
\midrule
prefix products $\prod_{i\le j+1} d_i$, $j=0\dots27$ & \texttt{0xb140}--\texttt{0xb4c0} & $28 \times \texttt{0x20} = \texttt{0x380}$ \\
\texttt{modexp} frame (3 lengths, base, exp, mod)    & \texttt{0xb4c0}--\texttt{0xb580} & $\texttt{0xc0}$ \\
\midrule
next allocated region (\mptr{LAGRANGE\_DENOMS\_MPTR}) & \texttt{0xb580} & slack $=$ \textbf{0 bytes} \\
\bottomrule
\end{tabular}
\caption{Scratch geometry for the $N=30$ Lagrange inversion. The frame ends exactly where
the denominator run begins.}
\label{tab:lag:scratch}
\end{table}

The exponent is $r-2$, so the precompile returns $g^{r-2} = g^{-1}$ by Fermat, where
$g = \prod_{j=0}^{29} d_j$. The backward pass then recovers each $d_j^{-1}$ in place, so
after line 1826 word $j$ of the window holds $d_j^{-1}$ rather than $d_j$.

The scratch base \texttt{0xb140} is \emph{not} virgin memory when this call runs: the
accumulator block already used it as \yul{acc\_scratch} at line 1269 (G1MSM staging), and
\mptr{SELECTOR\_ACC\_MPTR} (line 229) and \mptr{PCS\_FINAL\_MSM\_MPTR} $=\texttt{0xb280}$
(line 204) are later tenants of the same arena. No stale word can be consumed here: the
backward pass reads exactly the $28$ prefix words the forward pass wrote, descending from
$\texttt{0xb4a0}$ to $\texttt{0xb140}$ (line 939 pre-decrements \yul{gp\_mptr} from
$\texttt{0xb4c0}$, and the loop at lines 940--946 runs $28$ times), and the modexp frame
is fully rewritten by lines 918--923 before the \opcode{STATICCALL}. Every read in the
scratch region is therefore preceded by a write in the same call.

\paragraph{Rejection semantics.} \yul{batch\_invert} has five $\yul{ret}=0$ exits, and
\emph{none} of them writes into $[\yul{mptr\_start}, \yul{mptr\_end})$: the forward pass
writes only to \yul{scratch\_mptr}, and the modexp-failure exit at line 932 deliberately
\yul{leave}s before the backward pass. Enumerated:
\begin{itemize}
  \item reversed range, $\yul{mptr\_end} < \yul{mptr\_start}$ (lines 841--844) --
        unreachable here: $\texttt{0xb940} > \texttt{0xb580}$ are both rendered constants;
  \item empty window, $\yul{count\_bytes} = 0$ (line 848) -- this exit leaves
        $\yul{ret}=1$, and is unreachable here anyway;
  \item a non-canonical word $d_j \ge r$ (lines 888, 895, 905) -- unreachable in this call,
        because every $d_j$ is an \opcode{ADDMOD} output and hence $< r$;
  \item a zero total product (line 912) -- reachable, see below;
  \item a failed or short-returning \opcode{STATICCALL} (lines 924--932) -- reachable.
\end{itemize}
Only the last sets \yul{precompile\_failed} (line 926); the single-element fast path at
lines 852--878 also sets it at line 875, but $N=30$ never takes that path. Note that the
line-848 exit returns $\yul{ret}=1$ on an empty window, so an accidentally empty run would
be silently accepted -- immaterial here, but it means the ``$\yul{ret}=0$'' report is not a
proof that any work was done.

Because $d_j = x - \omega^{j-10}$ vanishes iff $x$ is the corresponding $n$-th root of
unity, and $d_{29} = x^n-1$ vanishes iff $x^n=1$, the zero-product condition is exactly
\[
\prod_{j} d_j = 0 \iff x^n = 1 \iff x \in \mu_n \subset \Fr .
\]
Since $x$ is a Keccak squeeze reduced mod $r$, this happens with probability
$\approx n/r = 2^{20}/2^{255} \approx 2^{-235}$. The inline comment at lines 1878--1880
states this bound and it is correct.

\begin{finding}[lag:F1]{Zero-slack adjacency between the inversion scratch and the denominator run}
\textbf{Severity: Observation.}
$\mptr{BATCH\_INV\_SCRATCH\_MPTR} = \texttt{0xb140}$ and
$\mptr{LAGRANGE\_DENOMS\_MPTR} = \texttt{0xb580}$
(\texttt{Halo2Verifier.sol}:235, 239). For the $N=30$ window this block builds,
$\texttt{0xb140} + 28\cdot\texttt{0x20} + \texttt{0xc0} = \texttt{0xb580}$ exactly --
the \texttt{modexp} frame's last byte is the byte before denominator $d_0$. One
additional word in the run would place the frame's final \yul{mstore} on top of $d_0$,
and the failure would be \emph{silent}: the modexp would still succeed, \yul{ret} would
stay $1$, and the backward pass would return internally consistent but wrong inverses,
yielding wrong Lagrange values and a verifier that accepts or rejects on garbage rather
than reverting. There is no runtime bound check on this adjacency anywhere in the
fixture. The mitigation is entirely generation-time: \texttt{src/lowering/layout/memory.rs}
lines 707--730 re-derives the template's run length independently and asserts both
\yul{n == template\_run\_words} and
$\yul{batch\_invert\_len} \ge (N-2)\cdot 32 + \texttt{0xc0}$, with a comment naming this
exact silent-corruption mode. That is the right place for the check and it is present;
the observation is that the deployed artifact carries no defence in depth, so the whole
argument rests on two Rust formulas staying in agreement. The fixture as generated is
correct.
\end{finding}

\subsection{Step 4: converting inverses into Lagrange evaluations}
\label{sec:lag:convert}

\begin{lstlisting}[style=solidity,caption={\texttt{Halo2Verifier.sol} lines 1828--1861: Lagrange conversion, $L_{\mathrm{blind}}$, and the instance combination},label={lst:lag:pass2}]
                // Convert inverted denominators into Lagrange evaluations:
                // L_i(x) = (x^n - 1) * n^-1 * omega_i / (x - omega_i).
                mptr := LAGRANGE_DENOMS_MPTR
                let l_i_common := mulmod(x_n_minus_1, mload(N_INV_MPTR), r)
                for { let pow_of_omega := mload(OMEGA_INV_TO_L_MPTR) }
                    lt(mptr, mptr_end)
                    { mptr := add(mptr, 0x20) } {
                    mstore(mptr, mulmod(l_i_common, mulmod(mload(mptr), pow_of_omega, r), r))
                    pow_of_omega := mulmod(pow_of_omega, omega, r)
                }

                // l_blind is the sum of the negative-rotation Lagrange terms
                // used by the midnight-proofs blinding identity.
                let l_blind := mload(add(LAGRANGE_DENOMS_MPTR, 0x20))
                let l_i_cptr := add(LAGRANGE_DENOMS_MPTR, 0x40)
                for { let l_i_cptr_end := add(LAGRANGE_DENOMS_MPTR, 0x0140) }
                    lt(l_i_cptr, l_i_cptr_end)
                    { l_i_cptr := add(l_i_cptr, 0x20) } {
                    l_blind := addmod(l_blind, mload(l_i_cptr), r)
                }

                // Public instance polynomial evaluation at x. Instance words
                // have already been range-checked and absorbed in transcript
                // order; this loop only forms the linear combination.
                let instance_eval := 0
                for {
                        let instance_cptr := INSTANCE_CPTR
                        let instance_cptr_end := add(instance_cptr, 0x0260)
                    }
                    lt(instance_cptr, instance_cptr_end)
                    { instance_cptr := add(instance_cptr, 0x20)
                      l_i_cptr := add(l_i_cptr, 0x20) } {
                    instance_eval := addmod(instance_eval, mulmod(mload(l_i_cptr), calldataload(instance_cptr), r), r)
                }
\end{lstlisting}

The second pass (lines 1830--1837) walks the same 29 words $j = 0,\dots,28$ -- note it
stops at \yul{mptr\_end}, so word 29 is \emph{not} touched and keeps $(x^n-1)^{-1}$ -- and
replaces each $d_j^{-1}$ with
\[
\boxed{\;\Lambda_j \;=\; \underbrace{(x^n-1)\,n^{-1}}_{\yul{l\_i\_common}} \cdot\; \omega^{\,j-10} \cdot\; \bigl(x-\omega^{\,j-10}\bigr)^{-1}\;}
\]
where \yul{pow\_of\_omega} is re-seeded from $\omega^{-10}$ so that its value at word $j$
is again $\omega^{j-10}$, matching the denominator written there in the first pass.

Writing $i = j - 10 \bmod n$, this is precisely the standard Lagrange basis evaluation on
a multiplicative domain,
\[
L_i(X) \;=\; \frac{\omega^{i}}{n}\cdot\frac{X^{n}-1}{X-\omega^{i}},
\]
which is Halo2's \texttt{EvaluationDomain::l\_i\_range}: it inverts $x - \omega^{i}$ for
each rotation, multiplies by the common factor $(x^n-1)\cdot$ \texttt{barycentric\_weight}
$= (x^n-1)n^{-1}$, and then applies \texttt{rotate\_omega} (the $\omega^i$ factor). So
$\Lambda_j = L_{j-10}(x)$ with negative indices read modulo $n$, and the generated code
matches the reference definition term for term. Crucially, \yul{x\_n\_minus\_1} on line
1831 is the \emph{stack} copy captured at line 1824, so it is still $x^n-1$ and not the
inverse that \yul{batch\_invert} wrote over word 29 -- reusing the memory word here would
have been a subtle bug, and the code does not.

\paragraph{$L_{\mathrm{last}}$ and $L_0$.} Lines 1866--1867 read
$\Lambda_0 = L_{-10}(x) = L_{n-10}(x)$ as \yul{l\_last}, and
$\Lambda_{10} = L_{0}(x)$ as \yul{l\_0} (offset $\texttt{0x0140} = 10$ words). Halo2
defines \texttt{l\_last} as \texttt{l\_evals[0]} over the rotation range
$-(\textit{blinding\_factors}+1)\dots 0$ and \texttt{l\_0} as
\texttt{l\_evals[1 + blinding\_factors]}; with $\textit{blinding\_factors} = 9$ these are
rotations $-10$ and $0$ respectively. Both match.

\paragraph{$L_{\mathrm{blind}}$.} Lines 1841--1847 compute
\[
L_{\mathrm{blind}}(x) \;=\; \sum_{j=1}^{9} \Lambda_j \;=\; \sum_{i=-9}^{-1} L_{i}(x)
\;=\; \sum_{t=n-9}^{n-1} L_{t}(x),
\]
by seeding the accumulator from word 1 (\yul{add(LAGRANGE\_DENOMS\_MPTR, 0x20)}) and then
adding words $2$ through $9$ -- the loop starts at \yul{+0x40} and runs while
$\yul{l\_i\_cptr} < \yul{LAGRANGE\_DENOMS\_MPTR} + \texttt{0x0140}$, so its last iteration
reads offset $\texttt{0x0120}$, i.e.\ word $9$. Nine terms in total. This is exactly
Halo2's \texttt{l\_evals[1..1+blinding\_factors].sum()}: the sum over the blinding rows,
excluding both the last row and row $0$. The formula and the bound are correct.

\paragraph{Instance evaluation.} The loop exit condition leaves \yul{l\_i\_cptr} at
precisely $\mptr{LAGRANGE\_DENOMS\_MPTR} + \texttt{0x0140}$, i.e.\ word $10$, and the
instance loop (lines 1853--1861) advances it in lockstep with the calldata cursor. With
$\texttt{0x0260} = 608 = 19\times 32$, it therefore computes
\[
I(x) \;=\; \sum_{t=0}^{18} L_{t}(x)\cdot \mathrm{inst}_t, \qquad
\mathrm{inst}_t = \yul{calldataload}(\texttt{0x1ee4} + 32t),
\]
consuming words $10$ through $28$ of the window -- which is the last word, so the run
length $29$ is exactly saturated. This is the standard Halo2 instance-polynomial
evaluation at rotation $\mathrm{cur}$: the public inputs occupy rows $0\dots18$, so
$I(X) = \sum_t \mathrm{inst}_t L_t(X)$ and $I(x)$ is its evaluation. The verifier stores a
single instance-evaluation slot (\mptr{INSTANCE\_EVAL\_MPTR}), consistent with the
instance column being queried only at \texttt{Rotation::cur} -- its sole consumer is line
2489, which feeds it into the permutation identity.

Reading the instances straight from calldata rather than from a memory mirror is safe
because the calldata region is immutable within the call and was fully validated at lines
1419--1432 and 1510--1523. It saves 19 \opcode{MLOAD}s worth of staging with no loss of
soundness.

\begin{finding}[lag:F2]{Fall-through of \yul{l\_i\_cptr} presumes at least two negative-rotation words}
\textbf{Severity: Observation.}
\texttt{Halo2Verifier.sol}:1841--1859. The instance loop does not re-initialise
\yul{l\_i\_cptr}; it inherits whatever the $L_{\mathrm{blind}}$ loop left. That is correct
only if the blinding loop terminates with $\yul{l\_i\_cptr} = \yul{LAGRANGE\_DENOMS\_MPTR}
+ \texttt{0x0140}$, which requires the unrolled first summand (word 1) and the loop start
(word 2) to be inside the negative-rotation prefix -- i.e.\ at least two such words. The
fixture has ten, so the code is correct as generated; the numbers were verified against
$\omega^{-L}$ with $L=10$. The observation is that the coupling is implicit: the
generated shape encodes a precondition on \texttt{num\_neg\_lagranges} that no comment,
guard, or literal in the deployed contract states. It is recorded here so that a reader
re-deriving the code does not mistake the fall-through for an accident.
\end{finding}

\subsection{Step 5: distillation into the named slots}
\label{sec:lag:distill}

\begin{lstlisting}[style=solidity,caption={\texttt{Halo2Verifier.sol} lines 1863--1883: distillation and the failure classification},label={lst:lag:distill}]
                // Persist the derived values into named memory slots consumed
                // by quotient reconstruction and PCS preparation.
                let x_n_minus_1_inv := mload(mptr_end)
                let l_last := mload(LAGRANGE_DENOMS_MPTR)
                let l_0 := mload(add(LAGRANGE_DENOMS_MPTR, 0x0140))

                mstore(X_N_MPTR, x_n)
                mstore(X_N_MINUS_1_INV_MPTR, x_n_minus_1_inv)
                mstore(L_LAST_MPTR, l_last)
                mstore(L_BLIND_MPTR, l_blind)
                mstore(L_0_MPTR, l_0)
                mstore(INSTANCE_EVAL_MPTR, instance_eval)
            }

            if iszero(success) {
                // A zero or non-canonical denominator is a rejected input,
                // not a broken chain: the only way to reach it is a squeezed
                // x that coincides with a domain point (probability ~n/r).
                if lagrange_precompile_failed { fail(ERR_PRECOMPILE_FAILED) }
                fail(ERR_PROOF_REJECTED)
            }
\end{lstlisting}

Six values leave the block, all in the contiguous window
$\texttt{0x7c40}$--$\texttt{0x7ce0}$:
\[
\begin{aligned}
\mptr{X\_N\_MPTR} &\leftarrow x^{n} &
\mptr{L\_BLIND\_MPTR} &\leftarrow \textstyle\sum_{j=1}^{9}\Lambda_j \\
\mptr{X\_N\_MINUS\_1\_INV\_MPTR} &\leftarrow (x^{n}-1)^{-1} &
\mptr{L\_0\_MPTR} &\leftarrow \Lambda_{10} = L_0(x) \\
\mptr{L\_LAST\_MPTR} &\leftarrow \Lambda_{0} = L_{n-10}(x) &
\mptr{INSTANCE\_EVAL\_MPTR} &\leftarrow I(x)
\end{aligned}
\]

After line 1874 the entire 30-word run at \texttt{0xb580} is dead. This is what licenses
the arena's reuse of that address range by the quotient VM's operand stack, whose base
literal \texttt{0xb8e0} appears at line 2172 (\yul{let q\_sp := 0xb8e0}) and in the
underflow guard at line 2229, and whose ceiling \texttt{0xbfa0} appears in the overflow
guard at line 2218. Since $\texttt{0xb580} \le \texttt{0xb8e0} < \texttt{0xb940}$, the
stack's lower words sit on top of window words $27$--$29$. The reuse is a lifetime-disjoint
phase reuse, not a collision, and the ordering that makes it sound is visible in the
source: distillation completes at line 1874, the block's braces close at line 1875, and the
quotient VM does not begin until after line 1916 (\yul{// \_\_phase:quotient\_vm}).

\begin{finding}[lag:F3]{The named slots are written before \yul{success} is examined}
\textbf{Severity: Observation.}
\texttt{Halo2Verifier.sol}:1826 vs.\ 1877. On every $\yul{ret}=0$ exit the window still
holds the \emph{raw denominators} $d_j$ (established above: no $\yul{ret}=0$ path in
\yul{batch\_invert} writes into the window), yet the code unconditionally runs the
conversion pass, the two accumulation loops, and all six \yul{mstore}s before testing
\yul{success}. The six named slots therefore transiently contain meaningless field
elements -- concretely, $\Lambda_j$ computed from $d_j$ instead of $d_j^{-1}$. This is
sound as written -- \yul{fail(sel)} (line 690) is \yul{mstore}$+$\yul{revert},
unconditionally terminal, and the first read of any of the six slots anywhere in the
contract is line 2489 (\mptr{INSTANCE\_EVAL\_MPTR}), followed by lines 2527, 2531, 2557
and 2587--2589 (\mptr{L\_0\_MPTR}, \mptr{L\_LAST\_MPTR}, \mptr{L\_BLIND\_MPTR}); the other
two are never read at all -- so no garbage value can influence a decision. The
design buys straight-line code with no branch in the hot path, at the cost of a
correctness argument that depends on \yul{fail} never returning and on the checkpoint
staying at line 1877. Flagged for visibility, not as a defect.
\end{finding}

\begin{finding}[lag:F6]{Two of the six named slots are dead stores in this fixture}
\textbf{Severity: Observation.}
\mptr{X\_N\_MPTR} $=\texttt{0x7c40}$ and \mptr{X\_N\_MINUS\_1\_INV\_MPTR}
$=\texttt{0x7c60}$ are written at \texttt{Halo2Verifier.sol}:1869--1870 and are
\emph{never read} -- neither symbolically nor through the literal addresses
\texttt{0x7c40}/\texttt{0x7c60}, anywhere in \texttt{Halo2Verifier.sol} or
\texttt{Halo2VerifyingKey.sol}. The quotient VM cannot reach them either: its
absolute-pointer opcode bounds \yul{q\_ptr} to $[\texttt{0x3680}, \texttt{0xa120}]$ via
\yul{if gt(sub(q\_ptr, 0x3680), 0x6aa0) \{ q\_program\_fail() \}} at line 2216, and while
\texttt{0x7c40} lies inside that range, no rendered opcode operand references it. The four
remaining slots are live (\mptr{INSTANCE\_EVAL\_MPTR} at line 2489; \mptr{L\_0\_MPTR} at
2527, 2537, 2587; \mptr{L\_LAST\_MPTR} at 2531, 2557, 2588; \mptr{L\_BLIND\_MPTR} at 2557,
2589). The cost is two \opcode{MSTORE}s plus the memory-expansion the two words already
imply; the correctness impact is nil. It is recorded because the surrounding comment at
lines 1863--1864 asserts these values are ``consumed by quotient reconstruction and PCS
preparation'', which for these two words is not true of this artifact -- a reader auditing
the PCS section should not go looking for a consumer that does not exist, and a future
change that starts reading \mptr{X\_N\_MPTR} would be relying on a slot that no test
currently pins.
\end{finding}

\subsection{Failure classification at the section boundary}
\label{sec:lag:fail}

Lines 1877--1883 are the only exit from this region other than falling through. The
two-valued report is the point of \yul{lagrange\_precompile\_failed}, which is hoisted to
line 1798 -- outside the inner braces -- so it survives the block scope:

\begin{center}
\begin{tabular}{@{}llll@{}}
\toprule
\yul{success} & \yul{lagrange\_precompile\_failed} & Outcome & Selector \\
\midrule
$1$ & $0$ & continue to the quotient VM & --- \\
$0$ & $1$ & \revert{PrecompileFailed}  & \texttt{0x84e81692} \\
$0$ & $0$ & \revert{ProofRejected}     & \texttt{0xc3b0d8cd} \\
\bottomrule
\end{tabular}
\end{center}

Both selectors were recomputed as $\mathrm{keccak256}(\text{``PrecompileFailed()''})[0..4]
= \texttt{0x84e81692}$ and $\mathrm{keccak256}(\text{``ProofRejected()''})[0..4] =
\texttt{0xc3b0d8cd}$, matching the constants at lines 324--325.

This region is clean. The classification is sound because the two conditions are
genuinely disjoint at the source: \yul{precompile\_failed} is set only at lines 875 and
926, both immediately after a \opcode{STATICCALL} to address \texttt{0x05} whose result or
\yul{returndatasize} was wrong; every other $\yul{ret}=0$ path leaves it at its
zero-initialised value. Distinguishing them is operationally meaningful -- a
\revert{PrecompileFailed} here means the node or the gas schedule is broken (the
\texttt{modexp} stipend \yul{MODEXP\_GAS = 4080} at line 293 is a fixed constant, and the
constructor smoke test at line 351 forward-probes it at deployment), whereas
\revert{ProofRejected} means a legitimate but astronomically unlikely transcript outcome.
Both fail closed, so the distinction costs one stack slot and buys correct incident
triage.

\begin{finding}[lag:F4]{Domain constants are trusted from the VK without runtime algebraic cross-checks}
\textbf{Severity: Informational.}
\texttt{Halo2Verifier.sol}:1808, 1818, 1831 read $\omega$, $\omega^{-10}$, and $n^{-1}$
from \mptr{OMEGA\_MPTR}, \mptr{OMEGA\_INV\_TO\_L\_MPTR}, and \mptr{N\_INV\_MPTR} and use
them directly. Unlike $k$, \texttt{num\_instances}, and the accumulator parameters -- each
cross-checked against a rendered literal at lines 1400--1405 -- these three are never
validated at runtime (no $\omega^{n}\stackrel{?}{=}1$, no $n^{-1}\cdot n
\stackrel{?}{=}1$). A wrong $n^{-1}$ or $\omega$ would silently produce a wrong Lagrange
basis and hence a verifier that accepts proofs for a different domain. The mitigation is
\mptr{EXPECTED\_VK\_CODEHASH} (line 95) enforced by \yul{extcodehash} at line 1388, which
pins every byte of the VK runtime including these words, so for the deployed pair the
values cannot differ from the audited ones. This specification verified them numerically
against $r$: $\omega$ is a primitive $2^{20}$-th root of unity, $n^{-1}\cdot 2^{20}\equiv
1$, $\omega\cdot\omega^{-1}\equiv 1$, and $\omega^{-L}=(\omega^{-1})^{10}$. Recorded so
that a reviewer knows the codehash pin, not an on-chain algebraic check, is what carries
this weight.
\end{finding}

\subsection{Consolidated algorithm}
\label{sec:lag:alg}

\begin{algorithm}[h]
\caption{Lagrange and instance-evaluation block (\texttt{Halo2Verifier.sol} 1792--1883)}
\label{alg:lag:main}
\begin{algorithmic}[1]
\Require $\yul{success}=1$; $x < r$ at \mptr{X\_MPTR}; VK loaded and header-checked
\State $\yul{lagrange\_precompile\_failed} \gets 0$
\State $x_n \gets x$
\For{$\mathit{idx} = 0$ \textbf{to} $19$} \Comment{$k = 20$ literal}
  \State $x_n \gets x_n^2 \bmod r$
\EndFor
\State $\pi \gets \omega^{-10}$ \Comment{\mptr{OMEGA\_INV\_TO\_L\_MPTR}}
\For{$j = 0$ \textbf{to} $28$}
  \State $W[j] \gets (x - \pi) \bmod r$; \quad $\pi \gets \pi\,\omega \bmod r$
\EndFor
\State $\delta \gets (x_n - 1) \bmod r$; \quad $W[29] \gets \delta$
\State $(\yul{success}, \yul{lagrange\_precompile\_failed}) \gets \textsc{BatchInvert}(W[0..29])$
\Statex \Comment{one \texttt{modexp} with exponent $r-2$; on success $W[j] \gets W[j]^{-1}$}
\State $c \gets \delta \cdot n^{-1} \bmod r$; \quad $\pi \gets \omega^{-10}$
\For{$j = 0$ \textbf{to} $28$}
  \State $W[j] \gets c \cdot W[j] \cdot \pi \bmod r$; \quad $\pi \gets \pi\,\omega \bmod r$
\EndFor
\State $\ell_{\mathrm{blind}} \gets W[1]$
\For{$j = 2$ \textbf{to} $9$} \State $\ell_{\mathrm{blind}} \gets \ell_{\mathrm{blind}} + W[j] \bmod r$ \EndFor
\State $I \gets 0$
\For{$t = 0$ \textbf{to} $18$}
  \State $I \gets I + W[10+t]\cdot \yul{calldataload}(\texttt{0x1ee4}+32t) \bmod r$
\EndFor
\State store $x_n$, $W[29]$, $W[0]$, $\ell_{\mathrm{blind}}$, $W[10]$, $I$ into the six named slots
\If{$\yul{success} = 0$}
  \If{$\yul{lagrange\_precompile\_failed} \ne 0$} \State \revert{PrecompileFailed} \EndIf
  \State \revert{ProofRejected}
\EndIf
\end{algorithmic}
\end{algorithm}

\paragraph{Cost.} The block performs 20 squarings, 29 subtractions and 29 root-of-unity
updates in the first pass, roughly 87 \opcode{MULMOD}s inside \yul{batch\_invert}, 87 more
in the conversion pass, 8 additions for $L_{\mathrm{blind}}$, and 38 operations for the
instance combination -- about 300 modular operations plus exactly one \texttt{modexp}.
The alternative, one \texttt{modexp} per denominator, would cost
$30 \times 4080 = 122{,}400$ gas in stipends alone; Montgomery's trick reduces that to a
single $4{,}080$-gas call plus $\approx 700$ gas of \opcode{MULMOD}s. This is the single
largest arithmetic saving in the verifier and it is implemented correctly.

\subsection{\yul{q\_program\_fail}}
\label{sec:lag:qfail}

\begin{lstlisting}[style=solidity,caption={\texttt{Halo2Verifier.sol} lines 1890--1893: \yul{q\_program\_fail}},label={lst:lag:qfail}]
            function q_program_fail() {
                mstore(0x00, shl(224, 0x3cc81b89))
                revert(0x00, 0x04)
            }
\end{lstlisting}

An unconditional, argument-free terminator. It writes the 4-byte selector
$\texttt{0x3cc81b89}$ left-aligned in the first word of Solidity's scratch space
(\yul{shl(224, .)} moves a \texttt{bytes4} into the high-order position) and reverts with
those four bytes. This specification recomputed
$\mathrm{keccak256}(\text{``QuotientProgramInvalid()''})[0..4] = \texttt{0x3cc81b89}$,
confirming the literal and the inline comment at line 1886.

Writing at \texttt{0x00} is safe: Solidity reserves \texttt{0x00}--\texttt{0x3f} as
scratch and \texttt{0x40} as the free-memory pointer, while the generated layout starts at
$\mptr{TRANSCRIPT\_MPTR} = \texttt{0x1000}$ (line 119), a separation enforced on every
call by the guard at lines 674--677 -- \yul{if gt(mload(0x40), TRANSCRIPT\_MPTR)} inside
\yul{verifyProof}'s assembly block (opened at line 647) -- which reverts with
\revert{MemoryLayoutViolated} (\texttt{0xc9888d23}, recomputed here) if \yul{mload(0x40)}
ever exceeds \texttt{0x1000}. It is a runtime guard, not a deployment-time one. This is the
same discipline as the general \yul{fail(sel)} helper at line 690.

The function is defined here rather than reusing \yul{fail(ERR\_QUOTIENT\_PROGRAM\_INVALID)}
because, per the comment at lines 1885--1889, the quotient VM is rendered into
\emph{both} this verifier and the standalone \texttt{Halo2QuotientEvaluator}, and
\yul{fail} is not in scope in the latter. It has 26 call sites in this fixture, all inside
the quotient interpreter (lines 2216--3105), which is covered elsewhere.

\begin{finding}[lag:F5]{The \revert{QuotientProgramInvalid} selector is duplicated as a literal}
\textbf{Severity: Observation.}
\texttt{Halo2Verifier.sol}:1891 hard-codes \texttt{0x3cc81b89} even though line 326
already defines \yul{ERR\_QUOTIENT\_PROGRAM\_INVALID} with the identical value and line
690 provides \yul{fail(sel)}. The duplication is deliberate and documented (the shared
quotient-VM body must also compile inside a contract that has neither the constant nor
the helper), and both copies are pinned by the generator test
\texttt{p4\_error\_selectors\_match\_declared\_errors} in \texttt{src/lowering/tests.rs}.
Both values were independently recomputed here and agree. No action required; noted only
because a second literal encoding of an error selector is the kind of thing that drifts
if the pinning test is ever weakened.
\end{finding}

\subsection{\yul{q\_pow5}}
\label{sec:lag:qpow5}

\begin{lstlisting}[style=solidity,caption={\texttt{Halo2Verifier.sol} lines 1911--1915: \yul{q\_pow5}},label={lst:lag:qpow5}]
            function q_pow5(x) -> z {
                let q_r := FR_MODULUS
                let x2 := mulmod(x, x, q_r)
                z := mulmod(x, mulmod(x2, x2, q_r), q_r)
            }
\end{lstlisting}

Three \opcode{MULMOD}s compute
\[
z \;=\; x \cdot \bigl(x^{2}\bigr)^{2} \;=\; x \cdot x^{4} \;=\; x^{5} \pmod r,
\]
which is the minimum-length addition chain for the fifth power. The routine is total: it
has no guards, no reverts, and no memory access, and it is correct for every input word,
including $x \ge r$ (which \opcode{MULMOD} reduces) and $x = 0$. It rebinds
\yul{q\_r := FR\_MODULUS} rather than closing over the enclosing \yul{r} because Yul
functions have no access to enclosing local scopes; the value is the same constant, so the
duplication is a language requirement, not a divergence.

The comment block at lines 1900--1910 records the provenance: the generator emits this
helper after recognising five equal multiplicative factors in an \texttt{Expression} tree
drawn from \texttt{vk.cs.gates}, corresponding to the Poseidon S-box
(\texttt{poseidon\_chip.rs::sbox}). This is important for a reviewer to internalise and
the comment states it explicitly: \yul{q\_pow5} is a peephole optimisation over an already
committed constraint system, \emph{not} a verifier rule. The constraint being checked is
whatever the VK's pinned gate expression says; recognising $x\cdot x\cdot x\cdot x\cdot x$
and emitting $x\cdot(x^2)^2$ changes only the number of \opcode{MULMOD}s, never the
polynomial identity. There are six call sites, all at lines 3132--3157 in the quotient
numerator, covered elsewhere.

The optimisation buys exactly one \opcode{MULMOD} per invocation over the naive
four-multiply form $x\cdot x\cdot x\cdot x\cdot x$ -- three multiplications instead of four,
so $8$ gas $\times$ 6 sites $= 48$ gas -- plus smaller code than inlining the chain six
times. The saving is negligible; the real reason to prefer the helper is code size. It
costs one more shape the generator's recogniser must get right; the mitigation is that a
mis-recognition would change the numerator and be caught by the native-anchor tests
referenced in \texttt{src/lowering/quotient\_numerator/native\_anchor.rs}.

\subsection{Assessment summary}
\label{sec:lag:summary}

Every formula in this region matches the standard Halo2 definitions. Specifically:

\begin{itemize}
  \item $x^n$ is $x^{2^{20}}$, from exactly $k=20$ squarings, with $k$ cross-checked
        against the VK at line 1401.
  \item $L_i(x) = \omega^{i}(x^n-1)\,n^{-1}(x-\omega^{i})^{-1}$ reproduces
        \texttt{EvaluationDomain::l\_i\_range} exactly, including the re-seeding of
        $\omega^{i}$ so that numerator and denominator powers agree per word.
  \item $L_{\mathrm{last}} = L_{n-10}(x)$, $L_0 = L_0(x)$, and
        $L_{\mathrm{blind}} = \sum_{i=-9}^{-1}L_i(x)$ match
        \texttt{l\_evals[0]}, \texttt{l\_evals[1+bf]}, and
        \texttt{l\_evals[1..1+bf]} with $\textit{bf}=9$, which is what
        $\omega^{-L}=(\omega^{-1})^{10}$ in the VK implies.
  \item $I(x)=\sum_{t=0}^{18}L_t(x)\,\mathrm{inst}_t$ is the instance polynomial at
        rotation $\mathrm{cur}$, over the 19 words the calldata-shape check pinned.
  \item The batch inversion window is $[d_0,\dots,d_{28},x^n-1]$ in that order --
        denominators first in ascending rotation from $-10$, vanishing term last -- and
        the vanishing term is the only word the conversion pass leaves untouched, which
        is why $(x^n-1)^{-1}$ survives to line 1865.
  \item The failure classification at lines 1877--1883 is complete and fails closed on
        both branches.
\end{itemize}

\textbf{No deviation from the reference definitions was found.} All six findings in this
section are Observation or Informational: they concern zero-slack memory adjacency whose
only guard is generation-time (\ref{sec:lag:invert}), an implicit precondition in the
loop-cursor fall-through, the ordering of the fail-closed checkpoint relative to the
stores, reliance on the VK codehash pin rather than on-chain algebraic checks for
$\omega$ and $n^{-1}$, a documented duplicate error-selector literal, and two write-only
named slots whose accompanying comment overstates their use. None of them is exploitable
in the fixture as generated, and the numeric verification of the VK domain constants
reported in Table~\ref{tab:lag:vk} was performed directly against $r$ (including
$\omega^{2^{20}}=1$, $\omega^{2^{19}}\neq 1$, $n^{-1}\cdot 2^{20}\equiv 1$,
$\omega\,\omega^{-1}\equiv 1$ and $\omega^{-L}=(\omega^{-1})^{10}$), as were all four
error selectors appearing in this section (\texttt{0x84e81692}, \texttt{0xc3b0d8cd},
\texttt{0x3cc81b89}, \texttt{0xc9888d23}).

Two things this section deliberately does \emph{not} claim. First, \yul{MODEXP\_GAS} $=
4080$ (line 293) is cited as a fixed stipend but its derivation against a particular
modexp repricing is out of range and is not verified here; the block's correctness does not
depend on the figure being optimal, only on the call succeeding, which the two-way
classification at lines 1877--1883 already handles. Second, the arena's phase model that
licenses the \texttt{0xb140} and \texttt{0xb580} overlaps is specified elsewhere; what is
verified here is the local ordering and the absence of any read-before-write inside this
block.


%% ==== 11-qvm-setup.tex ==================================================
\section{The quotient identity VM: rationale, state, and the inline prefix}
\label{sec:qvms}

This section specifies \texttt{Halo2Verifier.sol} lines 1917--2210: the design
rationale for the quotient identity interpreter, the memory frame and register
set it runs on, the two VK-resident pointers that supply its constants and its
program, the $y$-power precomputation table, the selector-accumulator zeroing
loop, and the four identities that are emitted as straight-line Yul \emph{before}
the interpreter loop starts. The interpreter's opcode semantics (lines
2210--3092) and the post-VM trash suffix and finalization (lines 3106--3359) are
covered elsewhere; this section cites them only where a bound established here is
consumed there.

Throughout, $r$ is the BLS12-381 scalar modulus
\texttt{0x73eda753299d7d483339d80809a1d80553bda402fffe5bfeffffffff00000001}
(\yul{FR\_MODULUS}, \texttt{Halo2Verifier.sol}:331), bound to the Yul local
\yul{r} at line 1355 and in scope for the whole of this block.

\subsection{What this phase computes, and what it deliberately does not}
\label{sec:qvms:target}

The block's own header (lines 1917--1932) states the contract precisely, and it
is worth restating in the notation of this document because it is the single
most important structural decision in the verifier.

Let $I_0, I_1, \dots, I_{N-1}$ be the sequence of \emph{partially evaluated
identities} produced by the Rust reference verifier's
\yul{partially\_evaluate\_identities}, in the fixed order
gate, permutation, lookup, trash. Each $I_j \in \Fr$ is a function of the alleged
polynomial evaluations $\mathbf{e} = (e_0,\dots,e_{101})$ read from the proof
after the transcript sampled $x$, of the Lagrange values $L_0(x)$,
$L_{\mathrm{last}}(x)$, $L_{\mathrm{blind}}(x)$, of the instance evaluation, and
of the challenges $\theta,\beta,\gamma,\tau$. The $y$-batched constraint
numerator is the Horner fold

\begin{equation}
  \nu_y(x) \;=\; \sum_{j=0}^{N-1} y^{\,N-1-j}\, I_j
  \;=\; \Bigl(\cdots\bigl((0 \cdot y + I_0)\cdot y + I_1\bigr)\cdots\Bigr)\cdot y + I_{N-1}.
  \label{eq:qvms:nuy}
\end{equation}

The verifier never receives, parses, or trusts a scalar $h(x)$. Instead this
block reconstructs $\nu_y(x)$ from the evaluations and stores its
\emph{negation} into \mptr{LINEARIZATION\_EVAL\_MPTR} (\texttt{0x7d00}) at
lines 3358--3359:

\begin{equation}
  \mathrm{mem}[\texttt{0x7d00}] \;\gets\; -\,\nu_y(x) \bmod r .
  \label{eq:qvms:negate}
\end{equation}

The matching commitment side is assembled in the following block from the
quotient limb commitments as
$(1 - x^n)\sum_i x_{\mathrm{split}}^{\,i}\,\mathsf{Q}_i$ plus the simple-selector
commitments, and the KZG multi-open check later binds that linearized commitment
to the scalar of \eqref{eq:qvms:negate} at the point $x$. Because $h(x)$ is
never a proof-supplied scalar, there is no opportunity for a prover to assert a
quotient evaluation that the verifier must take on faith. The inputs to
\eqref{eq:qvms:nuy} are exactly the ones the block header enumerates at lines
1948--1954: the evaluation vector $\mathbf{e}$, the Lagrange values and $x^n$
derived from the squeezed $x$, \mptr{INSTANCE\_EVAL\_MPTR} (a linear
combination of the caller's public instance words), and the squeezed challenges
in \mptr{CHALLENGE\_MPTR}$\dots$\mptr{Y\_MPTR}. Of these only $\mathbf{e}$ is
freely chosen by the prover, and every one of its 102 words is range-checked
canonical ($e_i < r$) before it is spilled to memory (line 1706,
\revert{NonCanonicalScalar}; the spill is at line 1708). The remainder are
either public input or Keccak outputs the prover cannot steer independently of
the commitments it has already sent.

\begin{finding}[QVMS-1]{Reconstructing $\nu_y(x)$ instead of trusting $h(x)$ removes a whole class of forgery}
\textbf{Severity: Observation (clean).} \\
\texttt{Halo2Verifier.sol}:1920--1932, 3358--3359. The design contains no code
path in which a proof-carried $h(x)$ scalar is read, and the expected opening
scalar for the linearized commitment is a pure function of already-committed,
already-range-checked evaluations and of challenges squeezed from the
transcript. Soundness of this phase therefore reduces entirely to (i) the
faithfulness of the identity program to the circuit's constraint system, which
is pinned by \yul{EXPECTED\_VK\_CODEHASH}, and (ii) the correctness of the fold
bookkeeping specified in \S\ref{sec:qvms:prefix}. This is the right structure and
it is implemented as documented.
\end{finding}

\subsection{Why a bytecode VM rather than unrolled Yul}
\label{sec:qvms:why}

Lines 1994--2006 give the rationale in the source. The forces are:

\begin{enumerate}
  \item \textbf{EIP-170.} The verifier runtime must fit in 24{,}576 bytes. The
        generator's own pinned-pragma header (\texttt{Halo2Verifier.sol}:8--11,
        identical in \texttt{templates/contracts/Halo2Verifier.sol})
        records that measured builds already emit 29{,}567 bytes at
        \yul{solc 0.8.24} with \texttt{runs=1} and 29{,}836 bytes at
        \yul{solc 0.8.30} with \texttt{runs=100000} -- both undeployable -- and
        that only the pinned \texttt{(0.8.30, runs)} pair produces a deployable
        contract. The size budget is not comfortable; it is the binding
        constraint.
  \item \textbf{Field constants are expensive as code.} Every distinct $\Fr$
        constant in an unrolled identity costs a \texttt{PUSH32} plus its 32
        immediate bytes, i.e.\ 33 runtime bytes each. This artifact's identity
        program references 178 such constants
        (\S\ref{sec:qvms:pointers}), which alone would be 5{,}874 bytes of
        verifier runtime.
  \item \textbf{The VK is already a separate, codehash-pinned data contract.}
        Bytes moved into the VK payload are still integrity-protected -- the
        verifier re-checks \yul{extcodesize} and \yul{extcodehash} of the linked
        \yul{Halo2VerifyingKey} against generated constants on \emph{every}
        proof (lines 1386--1389, \revert{VkMismatch}) before
        \yul{extcodecopy}-ing the payload to \mptr{VK\_MPTR} (line 1393) -- but
        they no longer count against the verifier's own 24 kB budget. The VK
        data contract has its own EIP-170 headroom
        (\texttt{src/lowering/render/models.rs}:271--274); this artifact's
        payload is \yul{EXPECTED\_VK\_PAYLOAD\_LENGTH} $=17{,}024$ bytes plus the
        one-byte \texttt{INVALID} prefix, comfortably under the cap.
\end{enumerate}

Relocating both the constant table (5{,}696 bytes) and the packed program
(4{,}559 bytes) into the VK moves 10{,}255 bytes -- about 60\% of the entire VK
payload -- out of the verifier's own runtime, and replaces the unrolled
arithmetic they would have generated with a single \texttt{switch} loop of
eleven arms.

\paragraph{What it costs.} Three things.
\emph{Gas:} each interpreted operation pays dispatch overhead (a word
\yul{mload}, a \yul{byte} extraction, a \yul{switch}, and a \yul{q\_pc}
increment) on top of the one \yul{addmod}/\yul{mulmod} it performs, and each
operand is fetched from memory rather than being a stack immediate. The
generator mitigates this by leaving the heaviest identities as
\emph{native callbacks} (\opcode{NATIVE\_PERMUTATION} \texttt{0x19},
\opcode{NATIVE\_LOOKUP} \texttt{0x1f}, \opcode{NATIVE\_IDENTITY} \texttt{0x1b})
and by rendering a fixed-size direct prefix of gate identities as straight-line
Yul (\S\ref{sec:qvms:prefix}), so the VM handles the long tail rather than the
hot path. The two selections are made by different rules: the native callbacks
are chosen by a gas/byte knapsack over the remaining gates
(\texttt{src/lowering/config.rs}:14--19), whereas the prefix is simply the
\emph{first} \yul{DEFAULT\_HYBRID\_QUOTIENT\_INLINE\_IDENTITIES} $= 4$ gate
identities in Rust order (\texttt{src/lowering/config.rs}:12,
\texttt{src/lowering/quotient.rs}:338--339) -- it is a constant, not a
cost ranking.
\emph{Auditability:} the constraint system is no longer readable in the verifier
source. A reviewer must decode 4{,}559 bytes of \texttt{quotient\_program} in
\texttt{Halo2VerifyingKey.sol} against the opcode table to see what is actually
being checked.
\emph{A new failure mode:} malformed generator output can now produce a program
that mis-tracks the operand stack or runs off the end of the stream. The
interpreter answers this with explicit structural guards
(\S\ref{sec:qvms:frame}) that fail closed via \yul{q\_program\_fail}.

\paragraph{The revert path.} All VM structural violations funnel through one
helper at lines 1890--1893, defined outside the block (rather than in
\texttt{AssemblyHelpers.yul}) because the quotient VM renders into both the main
verifier and the optional standalone evaluator:

\begin{lstlisting}[style=solidity,caption={\texttt{Halo2Verifier.sol} lines 1890--1893: the VM's fail-closed exit},label={lst:qvms:fail}]
            function q_program_fail() {
                mstore(0x00, shl(224, 0x3cc81b89))
                revert(0x00, 0x04)
            }
\end{lstlisting}

\noindent
\texttt{0x3cc81b89} is $\mathrm{bytes4}(\mathrm{keccak256}(\texttt{"QuotientProgramInvalid()"}))$,
pinned by the generator test \yul{p4\_error\_selectors\_match\_declared\_errors}.
Note that this is a \emph{different} error from the ordinary
\revert{ProofRejected} path: a \revert{QuotientProgramInvalid} revert indicates a
broken artifact, not a bad proof, and cannot be triggered by calldata
(\S\ref{sec:qvms:controlflow}).

\paragraph{This artifact is the monolithic render.} Lines 1944--1946 note that
the same template serves an external-evaluator path in which
\yul{Halo2QuotientEvaluator} first copies the verifier memory frame into the
same generated addresses. No such contract is present in this fixture directory,
and \yul{Halo2QuotientEvaluator} appears in the source only inside that comment,
so every address in this section is a direct memory address in the
\yul{verifyProof} frame.

\subsection{VK-resident constant table and program stream}
\label{sec:qvms:pointers}

Two pointers are established, at lines 2013 and 2016 (the surrounding lines are
comments). Both are Yul literals; neither is derived from calldata, from a
challenge, or from any runtime quantity.

\begin{lstlisting}[style=solidity,caption={\texttt{Halo2Verifier.sol} lines 2007--2028: challenge load, VK pointers, accumulator init, selector zeroing},label={lst:qvms:setup}]
                // Load the quotient batching challenge used by every fold.
                let y := mload(Y_MPTR)

                // q_const_mptr points to Fr constants used by the VM.
                // q_program_mptr points to the bytecode stream.
                // Constants are stored as consecutive 32-byte Fr words.
                let q_const_mptr := 0x3a60
                // Program bytes are also stored in the VK payload, packed into
                // 32-byte words by PackedProgramCodec.
                let q_program_mptr := 0x50a0
                // Running Horner accumulator for fully evaluated identities.
                // After all identities, this is nu_y(x) for the `None`
                // identity group.
                // Initialize A = 0 before scanning the identity stream.
                mstore(0xb280, 0)
                // Simple selectors are grouped into separate linearization
                // buckets. They start at zero for every proof.
                // q_sel_zero_off walks selector bucket byte offsets.
                for { let q_sel_zero_off := 0 } lt(q_sel_zero_off, 0x0140) { q_sel_zero_off := add(q_sel_zero_off, 0x20) } {
                    // B_s = 0 for each simple selector bucket.
                    mstore(add(SELECTOR_ACC_MPTR, q_sel_zero_off), 0)
                }
\end{lstlisting}

The VK payload occupies $[\texttt{0x3680}, \texttt{0x7900})$: \mptr{VK\_MPTR}
$=\texttt{0x3680}$ and \yul{EXPECTED\_VK\_PAYLOAD\_LENGTH} $= 17{,}024 =
\texttt{0x4280}$, so the payload ends exactly at \mptr{CHALLENGE\_MPTR}
$=\texttt{0x7900}$. Table~\ref{tab:qvms:vkregions} locates the two VM sections
inside it, cross-checked against the emitting \texttt{mstore} lines in
\texttt{Halo2VerifyingKey.sol}.

\begin{table}[htbp]
\centering
\caption{Quotient VM sections of the VK payload.}
\label{tab:qvms:vkregions}
\begin{tabular}{@{}llrl@{}}
\toprule
Section & Memory extent & Size & Emitting source \\
\midrule
VK header + base points & $[\texttt{0x3680},\texttt{0x3a60})$ & \texttt{0x03e0} (31 words) & \texttt{Halo2VerifyingKey.sol}:68--98 \\
Quotient constant table  & $[\texttt{0x3a60},\texttt{0x50a0})$ & \texttt{0x1640} (178 words) & \texttt{Halo2VerifyingKey.sol}:99--276 \\
Packed program bytes     & $[\texttt{0x50a0},\texttt{0x626f})$ & \texttt{0x11cf} (4{,}559 B) & \texttt{Halo2VerifyingKey.sol}:277--419 \\
Zero pad to word boundary& $[\texttt{0x626f},\texttt{0x6280})$ & 17 bytes & (tail of the last word) \\
Fixed / permutation commitments & $[\texttt{0x6280},\texttt{0x7900})$ & --- & (covered elsewhere) \\
\bottomrule
\end{tabular}
\end{table}

The program is stored as 143 consecutive 32-byte words (payload offsets
\texttt{0x1a20}--\texttt{0x2be0}) holding $143 \times 32 = 4{,}576$ bytes, of
which the last 17 are zero padding beyond the declared program length
\texttt{0x11cf}. The interpreter's termination test is a strict
$\yul{q\_pc} < \yul{q\_end}$ against the \emph{declared} length, so the padding is
never dispatched; had it been, byte \texttt{0x00} is not a case in the
\texttt{switch} and would land in the \yul{default} arm at line 3090 and revert.

Constants are addressed as $\mathrm{mem}[\texttt{0x3a60} + 32k]$ for a
program-supplied index $k$ (e.g.\ line 2245, \opcode{MUL\_CONST\_U8}). Slot $0$
holds the field element $1$ (\texttt{Halo2VerifyingKey.sol}:99), which is why
the generator can express ``multiply by the identity scale'' as a single
one-byte operand.

\begin{finding}[QVMS-2]{The constant-table index carries no \emph{runtime} bound, unlike every other decoded operand class}
\textbf{Severity: Informational.} \\
\texttt{Halo2Verifier.sol}:2245, 2317, 2347, 2369, 2396, 2414, 2432. All seven
constant-table accesses in the rendered interpreter have the shape
\yul{mload(add(q\_const\_mptr, shl(5, qconst)))} with \yul{qconst} decoded from
the program stream, and no arm compares \yul{qconst} against \texttt{178}.
Every other decoded operand class \emph{is} clamped at runtime: memory pointers
by \yul{gt(sub(q\_ptr, 0x3680), 0x6aa0)} (ten sites, \S\ref{sec:qvms:controlflow}),
the \opcode{FOLD\_SELECTOR} bucket index and gap at lines 3066--3067, and the
\opcode{NATIVE\_IDENTITY} callback index by the sub-switch
\yul{default \{ q\_program\_fail() \}} at line 3051.

\emph{Why this is only Informational.} Three independent reasons, and the
first two were checked in the generator rather than inferred from the comment
next to the code.
(i) The emitted indexes \emph{are} validated, exhaustively, at build time:
\yul{validate\_quotient\_const\_slots}
(\texttt{src/lowering/quotient\_numerator/vm/mod.rs}:2987--3008) folds over
\yul{visit\_quotient\_operands} and rejects any \yul{ConstSlotU8} or
\yul{ConstSlotU16} operand with \yul{slot >= consts.len()}; it is called
unconditionally on the finalized byte stream (\texttt{mod.rs}:1211) and panics
codegen on failure, so a program with an out-of-range index cannot be pinned
into a VK.
(ii) The generator does have a runtime clamp macro for this operand
(\yul{const\_guard}, \texttt{vm/yul\_arms.rs}:231--237, rendering
\yul{if iszero(lt(qconst, num\_consts)) \{ q\_program\_fail() \}}); it is emitted
only for the \emph{u16} constant opcodes \opcode{PUSH\_CONST} \texttt{0x01},
\opcode{ADD\_CONST} \texttt{0x0e} and \opcode{MUL\_CONST} \texttt{0x0f}
(\texttt{yul\_arms.rs}:412, 630, 648), none of which this artifact renders. The
omission for the byte-indexed forms is therefore a deliberate size trade, not an
oversight: a one-byte index cannot escape the table by more than 78 slots.
(iii) The reachable read window for a one-byte index is
$[\texttt{0x3a60},\texttt{0x3a60}+32\cdot255] = [\texttt{0x3a60},\texttt{0x59e0}]$,
entirely inside the codehash-pinned VK payload -- so even a hypothetical
out-of-range index reads a VK word, never proof data and never an unmapped
address, and the result is a wrong but \emph{prover-independent} constant that
breaks every proof equally rather than weakening any.
Adding \yul{if iszero(lt(qconst, 178)) \{ q\_program\_fail() \}} would still make
the runtime operand validation uniform, at roughly 20 gas per
\opcode{MUL\_CONST\_U8}.
\end{finding}

\subsection{The quotient VM memory frame}
\label{sec:qvms:frame}

The interpreter's working frame runs from \mptr{SELECTOR\_ACC\_MPTR}
$=\texttt{0xb140}$ to \texttt{0xbfa0}. Table~\ref{tab:qvms:frame} gives the exact
layout; the second column is the generator's name for the region
(\texttt{src/lowering/layout/memory.rs}:980--994,
\texttt{src/lowering/quotient.rs}:31--54).

\begin{table}[htbp]
\centering
\caption{Quotient VM memory frame, with the phase that previously owned each address.}
\label{tab:qvms:frame}
\begin{tabularx}{\textwidth}{@{}llrX@{}}
\toprule
Extent & Role & Words & Prior occupant (same addresses, earlier phase) \\
\midrule
$[\texttt{0xb140},\texttt{0xb280})$ & selector accumulators $B_0..B_9$ & 10 &
  \mptr{BATCH\_INV\_SCRATCH\_MPTR}: \yul{batch\_invert} prefix products \\
$\texttt{0xb280}$ & Horner accumulator $A$ (\yul{eval\_numer\_mptr}) & 1 &
  \yul{batch\_invert} prefix product \#10 \\
$\texttt{0xb2a0}$ & \yul{trace\_id\_mptr} (trace renders only; unused here) & 1 &
  \yul{batch\_invert} prefix product \#11 \\
$[\texttt{0xb2c0},\texttt{0xb8e0})$ & $y$-power table $y^0 \dots y^{48}$ & 49 &
  \yul{batch\_invert} prefix products \#12--27 (to \texttt{0xb4c0}) and its
  six-word modexp frame $[\texttt{0xb4c0},\texttt{0xb580})$; then
  \mptr{LAGRANGE\_DENOMS\_MPTR} at \texttt{0xb580} \\
$[\texttt{0xb8e0},\texttt{0xbfa0})$ & VM operand stack / native callback scratch & 54 &
  tail of the Lagrange denominator run, up to \texttt{0xb940} \\
\bottomrule
\end{tabularx}
\end{table}

Two consequences follow, and both matter.

\paragraph{The zeroing at lines 2021 and 2025--2028 is load-bearing, not defensive.}
EVM memory is zero-initialised, so a reader might take
\yul{mstore(0xb280, 0)} and the ten-iteration selector loop for belt-and-braces.
They are not. The Lagrange phase (lines 1808--1882) calls
\yul{batch\_invert(success, LAGRANGE\_DENOMS\_MPTR, add(mptr\_end, 0x20), BATCH\_INV\_SCRATCH\_MPTR, r)}
at line 1826 over a run of $30$ words starting at \texttt{0xb580}. In the general
path (lines 886--933) that call writes $28$ prefix products to
$[\texttt{0xb140},\texttt{0xb4c0})$ and a six-word EIP-198 modexp frame to
$[\texttt{0xb4c0},\texttt{0xb580})$ ($28 = 30 - 2$ stores: the
forward loop whose condition is \yul{lt(mptr, sub(mptr\_end, 0x20))} (line 893) skips
the first element, whose value seeds \yul{gp} without a store, and the last,
which is folded in as \yul{x\_last} at line 909). Every address in
$[\texttt{0xb140},\texttt{0xb4c0})$ therefore holds a nonzero, $x$-dependent
prefix product when the quotient block begins, and
$[\texttt{0xb4c0},\texttt{0xb580})$ holds the six-word EIP-198 modexp frame
$(\texttt{0x20}, \texttt{0x20}, \texttt{0x20}, gp, r-2, r)$ with its first word
overwritten by the precompile's 32-byte output (lines 918--924). No word in
$[\texttt{0xb140},\texttt{0xb580})$ is zero: the batch inverse exists only when
\yul{gp} is nonzero (line 912), which forces every prefix product to be nonzero
too. Without line 2021 the Horner
accumulator would start at a stale prefix product instead of $0$; without the
loop, each $B_s$ would start at one. Both would corrupt $\nu_y(x)$ and the
selector scalars for every proof.

\paragraph{The rest of the residue is provably dead.} The $y$-power loop
overwrites all of $[\texttt{0xb2c0},\texttt{0xb8e0})$ unconditionally
(\S\ref{sec:qvms:ypow}), so no Lagrange residue survives there. Address
\texttt{0xb2a0} is \emph{not} cleared and retains prefix product \#11, but this
is a non-trace render, so nothing in this phase reads it: the identifier
\yul{trace\_id\_mptr} is referenced only under the generator's trace flag
(\texttt{src/lowering/quotient.rs}:38), and the literal \texttt{0xb2a0} appears
in the contract only as an \yul{mstore} destination, at lines 3464, 3552 and
3608. The PCS phase does later read that word, but only through
\yul{add(0xb280, 0x20)} (e.g.\ line 3508) and only after line 3464 has written
it, so the stale value never reaches an arithmetic result. The first three stack words $[\texttt{0xb8e0},\texttt{0xb940})$ also
retain Lagrange values, but the stack discipline (below) guarantees no slot is
read before it is written.

\paragraph{Stack window and its guards.} The window is
$[\texttt{0xb8e0}, \texttt{0xbfa0})$, i.e.\ $\texttt{0x6c0} = 1728$ bytes
$= 54$ words. Its upper bound is the planner's \yul{quotient\_stack\_hi}
(\texttt{src/lowering/layout/memory.rs}:569--577), sized as the maximum of the
interpreted operand-stack depth and the structured scratch used by the native
callbacks that reuse the same base -- notably the permutation callback, which
writes a scratch table at \yul{program.stack\_mptr} (lines 2163--2166, 2470).
In this artifact the 54 words are not slack: the permutation callback partitions
the whole window exactly, as $18 + 18 + 6 + 6 + 5 + 1 = 54$ words at
\texttt{0xb8e0} (\yul{q\_perm\_vals}), \texttt{0xbb20} (\yul{q\_perm\_sigmas}),
\texttt{0xbd60} (\yul{q\_perm\_z\_cur}), \texttt{0xbe20}
(\yul{q\_perm\_z\_next}), \texttt{0xbee0} (\yul{q\_perm\_z\_last}) and
\texttt{0xbf80} (\yul{q\_perm\_delta\_base\_ptr}), the last of which ends at
\texttt{0xbfa0} (lines 2470--2475; the callback body is covered elsewhere).
Note also that the spill guard is tested \emph{before} the store, so the highest
legal spill is at $\yul{q\_sp} = \texttt{0xbf80}$, writing the window's last
word and leaving $\yul{q\_sp} = \texttt{0xbfa0}$; that is exactly in bounds and
the next spill reverts.
Ten structural guard classes enforce the VM's invariants; all are established
here and consumed in the interpreter body (the ten memory-operand clamps are
separate again, and are treated in \S\ref{sec:qvms:controlflow}):

\begin{table}[htbp]
\centering
\caption{Structural guards of the quotient VM. All revert via \yul{q\_program\_fail} (\revert{QuotientProgramInvalid}). The ten memory-operand guards are tabulated separately in \S\ref{sec:qvms:controlflow}.}
\label{tab:qvms:guards}
\begin{tabularx}{\textwidth}{@{}llX@{}}
\toprule
Line(s) & Condition & Meaning \\
\midrule
2218, 2332 & \yul{if iszero(lt(q\_sp, 0xbfa0))} & spill would overflow the window \\
2229       & \yul{if eq(q\_sp, 0xb8e0)}          & binary \opcode{ADD} would pop below the base \\
2464, 2580, 2772 & \yul{if iszero(eq(q\_sp, 0xb8e0))} & native callbacks are only legal at an empty spill stack \\
3051       & \yul{default \{ q\_program\_fail() \}} & \opcode{NATIVE\_IDENTITY}'s u16 callback index names no generated sub-case \\
3066       & \yul{if iszero(lt(q\_sel\_idx, 10))} & \opcode{FOLD\_SELECTOR} bucket index outside the 10 accumulators \\
3067       & \yul{if gt(q\_sel\_gap, 0x30)}      & \opcode{FOLD\_SELECTOR} gap outside the 49-entry $y$-power table \\
3068       & \yul{if iszero(q\_has\_top)}        & \opcode{FOLD\_SELECTOR} with no value to fold \\
3098       & \yul{if iszero(eq(q\_pc, q\_end))}  & the last opcode over- or under-read its operands \\
3099       & \yul{if q\_has\_top}                & an identity was left unconsumed \\
3105       & \yul{if iszero(eq(q\_sp, 0xb8e0))}  & spilled operands silently abandoned by a fold \\
\bottomrule
\end{tabularx}
\end{table}

The comment at lines 3100--3104 is worth preserving verbatim in a specification
because it explains why the third epilogue check is not redundant: a
\opcode{FOLD} executed with more than one operand live consumes only the cached
top, leaving abandoned words below \yul{q\_sp} while \yul{q\_has\_top} is clear
-- so the \yul{q\_pc} and \yul{q\_has\_top} checks both pass while an operand of
the identity has been silently dropped from $\nu_y(x)$. Only the \yul{q\_sp}
balance check catches that.

\begin{finding}[QVMS-3]{The quotient frame overlays three earlier phases, and the disjointness argument is temporal, not spatial}
\textbf{Severity: Informational.} \\
\texttt{Halo2Verifier.sol}:204, 229, 235, 239; \texttt{src/lowering/layout/memory.rs}:960--994.
\mptr{BATCH\_INV\_SCRATCH\_MPTR} $=$ \mptr{SELECTOR\_ACC\_MPTR} $=\texttt{0xb140}$
and \mptr{PCS\_FINAL\_MSM\_MPTR} $=\texttt{0xb280}$ share addresses with the
Horner accumulator; \mptr{LAGRANGE\_DENOMS\_MPTR} $=\texttt{0xb580}$ sits inside
the $y$-power table and runs 3 words past it into the stack window. None of these
is a bug -- the regions are disjoint in \emph{time} and the generator's arena
rejects any two regions sharing an address while both are live -- but note that
only the \texttt{0xb140} overlap is documented at the declaration (lines
231--234, pinned by
\yul{selector\_accumulators\_are\_live\_from\_quotient\_to\_final\_msm}). The
\texttt{0xb580} overlap with the $y$-power table and the \texttt{0xb280} overlap
with the final-MSM staging buffer are not called out in the contract source, so
a reader modifying either the Lagrange run length or the $y$-power table size by
hand has no local warning. The end-to-end liveness argument checks out: $A$ is
read for the last time at line 3358 and \texttt{0xb280} is first rewritten as PCS
scratch at line 3463; the Lagrange denominators are distilled into
\mptr{L\_LAST\_MPTR} / \mptr{L\_BLIND\_MPTR} / \mptr{L\_0\_MPTR} at lines
1866--1867 and never read again.
\end{finding}

\subsection{VM registers}
\label{sec:qvms:regs}

The five interpreter registers are declared at lines 2167--2176, after the
inline prefix has run.

\begin{lstlisting}[style=solidity,caption={\texttt{Halo2Verifier.sol} lines 2167--2176: VM register initialization},label={lst:qvms:regs}]
                // q_pc starts at the first encoded instruction.
                let q_pc := q_program_mptr
                // q_end is an exclusive byte pointer for the VM loop.
                let q_end := add(q_program_mptr, 0x11cf)
                // q_sp starts at the first free stack word.
                let q_sp := 0xb8e0
                // q_top is meaningless until q_has_top is set.
                let q_top := 0
                // q_has_top = 0 means the VM stack is empty.
                let q_has_top := 0
\end{lstlisting}

\begin{table}[htbp]
\centering
\caption{Quotient VM register file.}
\label{tab:qvms:registers}
\begin{tabularx}{\textwidth}{@{}lllX@{}}
\toprule
Register & Init & Domain & Role \\
\midrule
\yul{q\_pc}      & \texttt{0x50a0} & $[\texttt{0x50a0},\texttt{0x626f})$ at each loop head & byte-addressed instruction cursor into the VK program stream; a \texttt{case} arm may advance it past \yul{q\_end}, which the epilogue check at line 3098 then rejects \\
\yul{q\_end}     & \texttt{0x626f} & constant & exclusive end of the program; $\texttt{0x50a0} + \texttt{0x11cf}$ \\
\yul{q\_sp}      & \texttt{0xb8e0} & $[\texttt{0xb8e0},\texttt{0xbfa0})$ & word-aligned pointer to the first \emph{free} spill slot \\
\yul{q\_top}     & $0$ & $\Fr$ & cached top of the operand stack \\
\yul{q\_has\_top}& $0$ & $\{0,1\}$ & whether \yul{q\_top} holds a live value \\
\bottomrule
\end{tabularx}
\end{table}

The operand stack is thus a one-element register cache over a memory stack. The
invariant is: the logical stack has depth
$(\yul{q\_sp} - \texttt{0xb8e0})/32 + \yul{q\_has\_top}$; a push with
\yul{q\_has\_top} already set first spills \yul{q\_top} to
$\mathrm{mem}[\yul{q\_sp}]$ and advances \yul{q\_sp} by \texttt{0x20}; a binary
\opcode{ADD} retreats \yul{q\_sp} by \texttt{0x20} and combines the
spilled word with \yul{q\_top}. (The block header at lines 1978--1980 describes
this for ``\yul{ADD}/\yul{MUL}''; the ISA's binary \opcode{MUL} is
\texttt{0x07}, and this render emits no \texttt{case} for it, so \opcode{ADD}
\texttt{0x06} is the only pop site in the artifact.) The cache buys one avoided \yul{mstore} plus
\yul{mload} pair (6 gas plus the store's dirty-slot accounting) on the very
common ``push, then immediately combine'' pattern that dominates a linear
constraint expression.

\subsection{Opcode inventory}
\label{sec:qvms:ops}

Lines 2178--2196 render an opcode summary. The comment states that it is
generated from the same \yul{program.op\_usage} predicates that gate the
interpreter's \texttt{case} arms, so the artifact documents exactly the opcodes
its program can contain. That claim is verifiable against the fixture: the
\texttt{switch} at lines 2208--3092 has exactly eleven \texttt{case} arms
(\texttt{0x05}, \texttt{0x06}, \texttt{0x08}, \texttt{0x0d}, \texttt{0x10},
\texttt{0x11}, \texttt{0x21}, \texttt{0x19}, \texttt{0x1f}, \texttt{0x1b},
\texttt{0x0b}), matching the eleven listed opcodes one-for-one, plus a
\yul{default} that reverts.

\begin{table}[htbp]
\centering
\caption{Opcodes rendered for this artifact (\texttt{Halo2Verifier.sol} lines 2182--2192), with the line of the corresponding \texttt{case} arm. Semantics are specified elsewhere.}
\label{tab:qvms:opcodes}
\begin{tabular}{@{}llr@{}}
\toprule
Byte & Mnemonic & \texttt{case} at line \\
\midrule
\texttt{0x05} & \opcode{PUSH\_MEM\_U16}      & 2210 \\
\texttt{0x06} & \opcode{ADD}                 & 2226 \\
\texttt{0x08} & \opcode{NEG}                 & 2234 \\
\texttt{0x0b} & \opcode{FOLD\_SELECTOR}      & 3054 \\
\texttt{0x0d} & \opcode{MUL\_CONST\_U8}      & 2240 \\
\texttt{0x10} & \opcode{ADD\_MEM\_U16}       & 2248 \\
\texttt{0x11} & \opcode{MUL\_MEM\_U16}       & 2257 \\
\texttt{0x19} & \opcode{NATIVE\_PERMUTATION} & 2453 \\
\texttt{0x1b} & \opcode{NATIVE\_IDENTITY}    & 2762 \\
\texttt{0x1f} & \opcode{NATIVE\_LOOKUP}      & 2570 \\
\texttt{0x21} & \opcode{MODARITH7}           & 2281 \\
\bottomrule
\end{tabular}
\end{table}

\begin{finding}[QVMS-4]{The block header describes a \opcode{FOLD\_MAIN} opcode that this artifact does not contain}
\textbf{Severity: Informational (documentation).} \\
\texttt{Halo2Verifier.sol}:1982--1986 states that ``\yul{FOLD\_MAIN} or
\yul{FOLD\_SELECTOR} consumes it and advances the global $y$-batch position''.
The rendered opcode inventory (lines 2182--2192) and the \texttt{switch} contain
\opcode{FOLD\_SELECTOR} (\texttt{0x0b}) but no \opcode{FOLD\_MAIN}. In this
artifact every \emph{fully evaluated} (non-selector) identity is folded either by
a native callback -- lines 2528--2559 (permutation), 2597--2751 (lookup) and the
\opcode{NATIVE\_IDENTITY} body, whose four sub-cases at lines 2915, 2966, 3004
and 3043 in fact all fold into selector buckets -- or by the post-VM trash suffix
at lines 3303--3304, all of which perform the $A \gets A\cdot y$ then
$A \gets A + I$ pair inline rather than through a VM opcode.
\opcode{FOLD\_MAIN} is a genuine ISA opcode, \texttt{0x0a}
(\texttt{src/lowering/quotient\_numerator/vm/mod.rs}:551); the header is
template prose written for the ISA, while the opcode summary and the
\texttt{switch} are rendered from this program's \yul{op\_usage}, so the two
legitimately disagree. No functional impact; noted so a reader does not go
looking for a dispatch arm that is absent by design, and so that the header's
inventory is not mistaken for a bound on what the pinned program may contain --
that bound comes from the \yul{default} arm at line 3090, not from the comment.
\end{finding}

\subsection{The $y$-power precomputation table}
\label{sec:qvms:ypow}

Simple-selector identities are sparse in the global identity stream: a given
selector bucket contributes at positions $j_1 < j_2 < \dots$ with arbitrary
gaps. Advancing such a bucket's local accumulator across a gap of $g$ positions
requires a multiplication by $y^{g}$. The generator knows every gap at codegen
time, so rather than maintain a running per-selector scale (or compute $y^{-1}$
with a runtime \texttt{modexp}), the block precomputes a dense table of $y^i$.

\begin{lstlisting}[style=solidity,caption={\texttt{Halo2Verifier.sol} lines 2033--2050: the $y$-power table},label={lst:qvms:ypow}]
                {
                    // q_y_power holds y^i at the current loop index.
                    let q_y_power := 1
                    // Slot 0 holds y^0 = 1. Codegen never emits a read of it
                    // (FOLD_SELECTOR guards on a nonzero gap, and
                    // selector_tail_updates drops zero tails), but the tail
                    // block multiplies by mload(selector_power_mptr + offset)
                    // unconditionally -- so initialize the slot rather than
                    // leaving correctness to two filters in another file.
                    mstore(0xb2c0, 1)
                    // Start at i=1 because y^0 = 1 is written above.
                    for { let q_y_power_i := 1 } lt(q_y_power_i, 49) { q_y_power_i := add(q_y_power_i, 1) } {
                        // Advance from y^(i-1) to y^i modulo Fr.
                        q_y_power := mulmod(q_y_power, y, r)
                        // Store y^i at selector_power_mptr + 32*i.
                        mstore(add(0xb2c0, shl(5, q_y_power_i)), q_y_power)
                    }
                }
\end{lstlisting}

Formally, after this block:
\begin{equation}
  \mathrm{mem}[\texttt{0xb2c0} + 32i] \;=\; y^{i} \bmod r,
  \qquad i = 0,1,\dots,48 .
  \label{eq:qvms:ypow}
\end{equation}
The loop bound \texttt{49} is $\yul{selector\_fold.max\_power} + 1$
(\texttt{src/lowering/quotient.rs}:229--234) with
$\yul{max\_power} = 48$; the table therefore occupies exactly
$49 \times 32 = \texttt{0x620}$ bytes, ending at
$\texttt{0xb2c0} + \texttt{0x620} = \texttt{0xb8e0}$, which is precisely the
stack base. There is no slack word between the table and the stack.

\paragraph{The table's bounds are matched by the consumers' guards.}
This is the part a reviewer should check, and it holds exactly:

\begin{itemize}
  \item \opcode{FOLD\_SELECTOR} reads
        \yul{mload(add(0xb2c0, shl(5, q\_sel\_gap)))} at line 3085, guarded by
        \yul{if gt(q\_sel\_gap, 0x30) \{ q\_program\_fail() \}} at line 3067.
        $\texttt{0x30} = 48$, so the maximum reachable address is
        $\texttt{0xb2c0} + 32\cdot48 = \texttt{0xb8c0}$, the last table slot.
        The guard and the table size are exactly aligned -- no over-read, no
        wasted slot.
  \item The post-VM selector tail updates (lines 3315--3352) read fixed offsets
        \texttt{0x0600}, \texttt{0x05e0}, \texttt{0x0580}, \texttt{0x0520},
        \texttt{0x04c0}, \texttt{0x0460}, \texttt{0x0400}, \texttt{0x03a0},
        \texttt{0x0340}, \texttt{0x0280}, i.e.\ exponents
        $48,47,44,41,38,35,32,29,26,20$, one per bucket in \texttt{0x00} order.
        Each is $N - 1 - \mathrm{last}(s)$, where $N$ is the total number of
        $y$-batch positions and $\mathrm{last}(s)$ is bucket $s$'s final
        contribution; that is what re-aligns a sparse bucket with the global
        reverse fold. The largest, $48$, is again the last table slot, and it
        equals \yul{selector\_fold.max\_power}, which is what sized the table in
        the first place -- so the tail set and the table bound are the same
        number by construction, not by coincidence. $N$ itself cannot be read off
        from this range (it needs a decode of the program stream, covered
        elsewhere); the values are consistent with $N = 49$ and with buckets 0
        and 1 contributing only at their inline-prefix positions $j = 0$ and
        $j = 1$.
  \item The inline prefix (\S\ref{sec:qvms:prefix}) reads offset \texttt{0x20},
        exponent $1$.
\end{itemize}

\paragraph{On slot 0.} The comment at lines 2036--2041 is candid about why
$y^0 = 1$ is stored even though nothing is supposed to read it: the tail block
multiplies by \yul{mload(selector\_power\_mptr + offset)} unconditionally, and
the author preferred initializing the slot over relying on two independent
filters in the generator (\opcode{FOLD\_SELECTOR}'s
\yul{if q\_sel\_gap} guard at line 3081, and \yul{selector\_tail\_updates}
dropping zero tails). This is the right call: it makes the rendered artifact
correct on its own terms rather than correct-because-of-an-invariant-elsewhere.
Note also that without it, slot 0 would hold a \yul{batch\_invert} prefix-product
residue rather than zero (\S\ref{sec:qvms:frame}), so an accidental read would be
silently wrong rather than obviously so.

\paragraph{Cost.} 48 \yul{mulmod}s and 49 \yul{mstore}s, unconditionally, on
every proof -- roughly 700--900 gas -- against the alternative of one runtime
\texttt{modexp} for $y^{-1}$ (\yul{MODEXP\_GAS} $= 4080$ forwarded, per line 293)
plus per-identity scale maintenance. The precomputation is clearly cheaper, and
it also removes a \texttt{modexp} call site, which removes a
precompile-repricing liveness dependency.

\subsection{The selector-accumulator zeroing loop}
\label{sec:qvms:selzero}

Lines 2025--2028 (Listing~\ref{sec:qvms:frame}) run
$\lceil \texttt{0x0140}/\texttt{0x20} \rceil = 10$ iterations with
\yul{q\_sel\_zero\_off} taking the values
$\texttt{0x00}, \texttt{0x20}, \dots, \texttt{0x0120}$, writing $0$ to
$\mathrm{mem}[\texttt{0xb140} + \mathit{off}]$. The bound \texttt{0x0140} is
$32 \times \yul{num\_simple\_selectors}$ with
$\yul{num\_simple\_selectors} = 10$, and it is matched exactly by the
\opcode{FOLD\_SELECTOR} bucket guard
\yul{if iszero(lt(q\_sel\_idx, 10)) \{ q\_program\_fail() \}} at line 3066 and by
the ten tail updates at lines 3315--3352. The region cleared is
$[\texttt{0xb140}, \texttt{0xb280})$, ending exactly where the Horner
accumulator begins.

Note that the bound is written as a byte count (\texttt{0x0140}) rather than as a
count of selectors, while the interpreter's guard is written as a selector count
(\texttt{10}). The two are consistent here, but they are two independently
rendered numbers describing one quantity; a divergence would leave the highest
bucket(s) holding \yul{batch\_invert} residue while \opcode{FOLD\_SELECTOR} still
accepted writes to them. This is generator-internal and there is no runtime
cross-check, but nothing calldata-controlled can cause a divergence.

\subsection{The direct inline prefix}
\label{sec:qvms:prefix}

Lines 2052--2153 emit four identities as straight-line Yul before the VM loop
starts (the comment occupies 2052--2054; the code is 2055--2153). They are not
selected by cost: the generator takes the \emph{first}
$\min(\lvert\mathrm{gates}\rvert, 4)$ gate identities in the Rust
\yul{partially\_evaluate\_identities} order, where $4$ is the fixed constant
\yul{DEFAULT\_HYBRID\_QUOTIENT\_INLINE\_IDENTITIES}
(\texttt{src/lowering/config.rs}:12, \texttt{src/lowering/quotient.rs}:338--339).
That they happen to be small expression trees over evaluation slots is a property
of this circuit, not of the selection rule. Critically, they use the
\emph{same fold snippets} as VM and native identities, so they occupy the same
global $y$-batch positions ($j = 0,1,2,3$) and no special-casing is needed
downstream.

\subsubsection{The fold snippets}

Two fold shapes appear. Both begin by advancing the global Horner position:
\begin{equation}
  A \;\gets\; A \cdot y \bmod r ,
  \qquad A \equiv \mathrm{mem}[\texttt{0xb280}] .
  \label{eq:qvms:hornerstep}
\end{equation}
This happens for \emph{every} identity, selector or not, which is what keeps
later main identities landing on the same powers of $y$ as the Rust reference
verifier's reverse fold. A fully evaluated identity then additionally does
$A \gets A + I_j$; a simple-selector identity instead routes its value to bucket
$s$:
\begin{equation}
  B_s \;\gets\; B_s \cdot y^{g} + I_j \bmod r ,
  \label{eq:qvms:selfold}
\end{equation}
where $g$ is the number of $y$-batch positions since bucket $s$ last
contributed, and the $y^g$ factor is omitted entirely from the rendered code when
$g = 0$. All four prefix identities are simple-selector identities, so
\eqref{eq:qvms:hornerstep} runs four times with no addition into $A$: after the
prefix, $A = 0$ still, but the batch position has advanced to $4$.

This is exactly the semantics of \opcode{FOLD\_SELECTOR} (lines 3054--3088), with
one mechanical difference: the interpreter takes $I_j$ from the \yul{q\_top}
register, whereas the inline prefix writes $I_j$ to the scratch word
$\mathrm{mem}[\texttt{0xb8e0}]$ and reads it back. That word is the stack base,
and \yul{q\_sp} is not initialized until line 2172 -- after the prefix -- so
there is no conflict; the generator floors the stack region at one word precisely
so this eval-scratch write is always in bounds
(\texttt{src/lowering/quotient.rs}:201--219).

\subsubsection{Identity 0 (bucket 0, gap 0)}

\begin{lstlisting}[style=solidity,caption={\texttt{Halo2Verifier.sol} lines 2055--2098: prefix identity 0 and its selector fold},label={lst:qvms:prefix0}]
                {
                let var0 := 0x1
                let f_3 := mload(0x9ac0)
                let f_4 := mload(0x99c0)
                let a_0 := mload(0x94a0)
                let var1 := mulmod(f_4, a_0, r)
                let var2 := addmod(f_3, var1, r)
                let f_5 := mload(0x99e0)
                let a_1 := mload(0x94c0)
                let var3 := mulmod(f_5, a_1, r)
                let var4 := addmod(var2, var3, r)
                let f_6 := mload(0x9a00)
                let a_2 := mload(0x94e0)
                let var5 := mulmod(f_6, a_2, r)
                let var6 := addmod(var4, var5, r)
                let f_7 := mload(0x9a20)
                let a_3 := mload(0x9500)
                let var7 := mulmod(f_7, a_3, r)
                let var8 := addmod(var6, var7, r)
                let f_8 := mload(0x9a40)
                let a_4 := mload(0x9520)
                let var9 := mulmod(f_8, a_4, r)
                let var10 := addmod(var8, var9, r)
                let f_0 := mload(0x9a60)
                let a_0_next_1 := mload(0x9540)
                let var11 := mulmod(f_0, a_0_next_1, r)
                let var12 := addmod(var10, var11, r)
                let f_1 := mload(0x9a80)
                let var13 := mulmod(f_1, a_0, r)
                let var14 := mulmod(var13, a_1, r)
                let var15 := addmod(var12, var14, r)
                let f_2 := mload(0x9aa0)
                let var16 := mulmod(f_2, a_0, r)
                let var17 := mulmod(var16, a_2, r)
                let var18 := addmod(var15, var17, r)
                let var19 := mulmod(var0, var18, r)
                mstore(0xb8e0, var19)
                }
                mstore(0xb280, mulmod(mload(0xb280), y, r))
                {
                let q_selector_ptr := add(SELECTOR_ACC_MPTR, 0x0)
                let q_selector_acc := mload(q_selector_ptr)
                mstore(q_selector_ptr, addmod(q_selector_acc, mload(0xb8e0), r))
                }
\end{lstlisting}

Every \yul{mload} address lies in the decoded evaluation buffer
$[\mptr{REVERSED\_EVALS\_MPTR}, \mptr{ADVICE\_COMMS\_MPTR\_BASE}) =
[\texttt{0x9480}, \texttt{0xa140})$, which holds the $102$ proof evaluations
$e_0,\dots,e_{101}$ spilled at line 1708 after the canonicality check at line
1706. Table~\ref{tab:qvms:slots} decodes the generated variable names.

\begin{table}[htbp]
\centering
\caption{Evaluation slots referenced by the inline prefix. Index $i$ is $(\mathrm{addr} - \texttt{0x9480})/32$.}
\label{tab:qvms:slots}
\begin{tabular}{@{}llll@{}}
\toprule
Yul name & Address & $e_i$ & Yul name / Address / $e_i$ \\
\midrule
\yul{a\_0}        & \texttt{0x94a0} & $e_{1}$ & \yul{f\_4} / \texttt{0x99c0} / $e_{42}$ \\
\yul{a\_1}        & \texttt{0x94c0} & $e_{2}$ & \yul{f\_5} / \texttt{0x99e0} / $e_{43}$ \\
\yul{a\_2}        & \texttt{0x94e0} & $e_{3}$ & \yul{f\_6} / \texttt{0x9a00} / $e_{44}$ \\
\yul{a\_3}        & \texttt{0x9500} & $e_{4}$ & \yul{f\_7} / \texttt{0x9a20} / $e_{45}$ \\
\yul{a\_4}        & \texttt{0x9520} & $e_{5}$ & \yul{f\_8} / \texttt{0x9a40} / $e_{46}$ \\
\yul{a\_0\_next\_1} & \texttt{0x9540} & $e_{6}$ & \yul{f\_0} / \texttt{0x9a60} / $e_{47}$ \\
\yul{a\_1\_next\_1} & \texttt{0x9560} & $e_{7}$ & \yul{f\_1} / \texttt{0x9a80} / $e_{48}$ \\
                  &                 &         & \yul{f\_2} / \texttt{0x9aa0} / $e_{49}$ \\
                  &                 &         & \yul{f\_3} / \texttt{0x9ac0} / $e_{50}$ \\
\bottomrule
\end{tabular}
\end{table}

Writing $a_k$ for the advice evaluation in slot $k$, $a_k'$ for its
$\omega$-rotated companion, and $f_k$ for the fixed-column evaluation, identity 0
is the standard PLONK-style custom gate

\begin{equation}
  I_0 \;=\; 1 \cdot \bigl(
    f_3
    + f_4 a_0 + f_5 a_1 + f_6 a_2 + f_7 a_3 + f_8 a_4
    + f_0 a_0'
    + f_1 a_0 a_1
    + f_2 a_0 a_2
  \bigr) \bmod r .
  \label{eq:qvms:i0}
\end{equation}
The evaluation order in the code is strict left-to-right Horner-free
accumulation (\yul{var2}, \yul{var4}, \dots, \yul{var18}), each step reduced by
\yul{addmod}/\yul{mulmod} modulo $r$, so no intermediate can exceed $r$.

\subsubsection{Identities 1--3}

\begin{align}
  I_1 &= 1 \cdot \bigl( a_1 + a_2 - a_3 - a_4 \bigr) \bmod r
        &&\text{(lines 2099--2112, bucket 1, gap 0)} \label{eq:qvms:i1}\\
  I_2 &= 1 \cdot \bigl( a_0 + f_4 - a_0' \bigr) \bmod r
        &&\text{(lines 2119--2129, bucket 2, gap 0)} \label{eq:qvms:i2}\\
  I_3 &= 1 \cdot \bigl( a_1 + f_5 - a_1' \bigr) \bmod r
        &&\text{(lines 2136--2146, bucket 2, gap 1)} \label{eq:qvms:i3}
\end{align}

Negation is rendered as \yul{addmod(0, sub(r, v), r)}. This is correct for the
whole input range including $v = 0$: $\yul{sub}(r, 0) = r$ and
$\yul{addmod}(0, r, r) = 0$. It is also correct for a hypothetical non-canonical
$v \geq r$, though that cannot occur here because line 1706 rejects any
$e_i \geq r$ before it reaches memory.

Identity 3 is the only prefix fold with a nonzero gap, and it is the one place in
this range where the $y$-power table is consumed:

\begin{lstlisting}[style=solidity,caption={\texttt{Halo2Verifier.sol} lines 2136--2153: prefix identity 3 and the only prefix fold with a $y^g$ gap factor},label={lst:qvms:prefix3}]
                {
                let var0 := 0x1
                let a_1 := mload(0x94c0)
                let f_5 := mload(0x99e0)
                let var1 := addmod(a_1, f_5, r)
                let a_1_next_1 := mload(0x9560)
                let var2 := addmod(0, sub(r, a_1_next_1), r)
                let var3 := addmod(var1, var2, r)
                let var4 := mulmod(var0, var3, r)
                mstore(0xb8e0, var4)
                }
                mstore(0xb280, mulmod(mload(0xb280), y, r))
                {
                let q_selector_ptr := add(SELECTOR_ACC_MPTR, 0x40)
                let q_selector_acc := mload(q_selector_ptr)
                q_selector_acc := mulmod(q_selector_acc, mload(add(0xb2c0, 0x20)), r)
                mstore(q_selector_ptr, addmod(q_selector_acc, mload(0xb8e0), r))
                }
\end{lstlisting}

\begin{table}[htbp]
\centering
\caption{The inline prefix's four identities and their fold targets.}
\label{tab:qvms:prefixfolds}
\begin{tabular}{@{}rlllll@{}}
\toprule
$j$ & Identity lines & Fold lines & Bucket $s$ & \yul{SELECTOR\_ACC} offset & Gap $g$ \\
\midrule
0 & 2055--2092 & 2093--2098 & 0 & \texttt{0x00} & 0 \\
1 & 2099--2112 & 2113--2118 & 1 & \texttt{0x20} & 0 \\
2 & 2119--2129 & 2130--2135 & 2 & \texttt{0x40} & 0 \\
3 & 2136--2146 & 2147--2153 & 2 & \texttt{0x40} & 1 \\
\bottomrule
\end{tabular}
\end{table}

The gap of $1$ on identity 3 is exactly $j_2 - j_1 = 3 - 2$ for bucket 2, which
is the invariant \eqref{eq:qvms:selfold} requires. Since $B_2 = I_2$ after
identity 2, the fold yields $B_2 = I_2 \cdot y + I_3$, i.e.\ bucket 2's own local
Horner accumulation, independent of the three intervening global positions. The
final tail multiplications at lines 3315--3352 then lift every bucket into the
same global reverse-fold scale before the buckets become MSM scalars.

\begin{finding}[QVMS-5]{The inline prefix is guard-free by construction, and that is sound}
\textbf{Severity: Observation (clean).} \\
\texttt{Halo2Verifier.sol}:2055--2153. Unlike the interpreter, the prefix
contains no bounds checks whatsoever. It does not need any: every memory address
it touches is a compile-time literal, every one of them resolves inside the
102-word evaluation buffer $[\texttt{0x9480},\texttt{0xa140})$ or the two
generated accumulator regions, every value it loads was range-checked
canonical at parse time (line 1706), and every arithmetic step is an
\yul{addmod}/\yul{mulmod} modulo $r$. There is no index, no cursor, and no
prover-influenced quantity anywhere in the block. The only way to break it is to
regenerate the contract from a different constraint system, which changes
\yul{EXPECTED\_VK\_CODEHASH} and the deployed bytecode. The remaining risk --
that the rendered Yul does not compute the identity the circuit declares -- is
addressed generator-side rather than at runtime, by
\yul{certify\_inline\_prefix} (\texttt{src/lowering/plan.rs}:355), which
value-checks this prefix against an independent typed lowering of the same gates
and fails codegen on a mismatch. A reviewer verifying this section by hand
should therefore re-derive
\eqref{eq:qvms:i0}--\eqref{eq:qvms:i3} from the circuit, not from the generated
variable names: \yul{f\_k} and \yul{a\_k} are emitter labels with no runtime
meaning, and Table~\ref{tab:qvms:slots} is the only thing that ties them to
concrete evaluation slots.
\end{finding}

\begin{finding}[QVMS-6]{Each inline identity pays a semantically dead \yul{mulmod} by the literal 1}
\textbf{Severity: Observation.} \\
\texttt{Halo2Verifier.sol}:2056/2090, 2100/2110, 2120/2127, 2137/2144. Each
prefix identity opens with \yul{let var0 := 0x1} and closes with
\yul{mulmod(var0, <expr>, r)}. Since \yul{<expr>} is always the output of an
\yul{addmod(..., r)} and therefore already less than $r$, this multiplication is
the identity map. It is a codegen artifact: the emitter renders the identity's
leading constant scale uniformly, and here that scale happens to be $1$. The Yul
optimizer cannot fold it, because proving $\yul{mulmod}(1,v,r) = v$ requires
knowing $v < r$, which is not visible to it. Cost is 4 identities $\times$
(3 gas \yul{PUSH} + 8 gas \yul{MULMOD}) $\approx$ 44 gas per proof -- negligible,
but it is dead code in the most scrutinized arithmetic in the contract, and a
reviewer should know it is inert rather than wonder what scale is being applied.
\end{finding}

\subsection{The interpreter loop head}
\label{sec:qvms:loophead}

The final ten lines of this range open the interpreter.

\begin{lstlisting}[style=solidity,caption={\texttt{Halo2Verifier.sol} lines 2200--2210: loop condition, opcode fetch, and dispatch},label={lst:qvms:loophead}]
                for { } lt(q_pc, q_end) { } {
                    // The bytecode table is byte-addressed, but EVM memory
                    // loads whole words. `byte(0, mload(q_pc))` extracts the
                    // opcode at the current byte cursor; each case advances
                    // q_pc by exactly its operand width.
                    let q_op := byte(0, mload(q_pc))
                    q_pc := add(q_pc, 1)

                    switch q_op
                    // VM 0x05 PUSH_MEM_U16 (bytes): next two bytes are a short memory pointer.
                    case 0x05 {
\end{lstlisting}

Three properties of the fetch are worth stating explicitly:

\begin{enumerate}
  \item The loop has an empty post-body clause; \yul{q\_pc} is advanced solely by
        the opcode fetch (one byte) and by each \texttt{case} arm's own operand
        width. There is no branch, jump, or call opcode in the ISA -- no
        \texttt{case} assigns \yul{q\_pc} anything other than
        \yul{add(q\_pc, <literal>)}. The program is therefore a strictly
        forward, straight-line byte stream, executed exactly once end to end.
  \item \yul{mload(q\_pc)} reads a full 32-byte word starting at an unaligned
        byte cursor. Near the end of the program this reads past
        \yul{q\_end} into the 17 zero pad bytes and, for the last few
        instructions, past the padding into the fixed-commitment words at
        \texttt{0x6280}. Those bytes are discarded by the \yul{byte}/\yul{shr}
        extraction as long as each \texttt{case} consumes only the operand width
        it declares -- which is exactly what the epilogue check
        \yul{if iszero(eq(q\_pc, q\_end))} at line 3098 verifies. The read itself
        is always in bounds: \texttt{0x6280} is inside the VK payload, and memory
        up to \texttt{0xbfa0} is already touched by this frame, so no memory
        expansion is triggered.
  \item An unknown opcode byte -- including \texttt{0x00} and, as line 3089 notes
        deliberately, \texttt{0x1a} (\yul{Q\_OP\_RESERVED},
        \texttt{src/lowering/quotient\_numerator/vm/yul\_arms.rs}:62) -- falls
        into the \yul{default} arm at line 3090 and reverts. The dispatch fails
        closed. Note that this is what actually bounds the opcode set the pinned
        program may use; the eleven-line summary comment at lines 2182--2192 is
        documentation and enforces nothing.
\end{enumerate}

\subsection{Assessment: can calldata influence control flow?}
\label{sec:qvms:controlflow}

This is the central security question for the region, and the answer is
\textbf{no}. The argument has four legs, each of which is checkable in the
fixture.

\begin{enumerate}
  \item \textbf{The program cursor and its bound are literals.}
        $\yul{q\_pc} \gets \texttt{0x50a0}$ and
        $\yul{q\_end} \gets \texttt{0x50a0} + \texttt{0x11cf}$ (lines
        2168--2170) via the literal \yul{q\_program\_mptr := 0x50a0} at line 2016
        and the immediate \texttt{0x11cf} at line 2170. No calldata value, no
        challenge, and no memory read participates in either.
  \item \textbf{The ISA has no control-flow instruction.} Every \texttt{case}
        advances \yul{q\_pc} by a fixed literal width; the only data-dependent
        branches inside the interpreter body are (a) the \texttt{switch} on the
        program byte itself and (b) \yul{if q\_sel\_gap} at line 3081, whose
        condition is also a program byte. Neither reads proof data.
  \item \textbf{The program bytes cannot be modified after they are loaded.}
        They arrive by \yul{extcodecopy} at line 1393 from an address whose
        \yul{extcodesize} and \yul{extcodehash} are re-verified against the
        generated constants \yul{EXPECTED\_VK\_LENGTH} and
        \yul{EXPECTED\_VK\_CODEHASH\_WORD} on \emph{every} call (lines
        1386--1389) -- not merely at construction (line 582). Between that copy
        and the VM, nothing writes into $[\texttt{0x3680},\texttt{0x7900})$: a
        grep of the contract finds no \yul{mstore} to a literal in that range and
        no \yul{mstore}, \yul{mstore8} or \yul{mcopy} whose destination is a
        literal in that range (the \yul{mcopy}s in the
        final-MSM staging block, e.g.\ lines 3848--3862, read \emph{from} the
        VK region into \mptr{PCS\_FINAL\_MSM\_MPTR}, never the reverse), and all twelve
        \yul{calldatacopy} sites (lines 794, 1558, 1577, 1596, 1614, 1621, 1630,
        1637, 1655, 1678, 1728, 1774) fall into three destination families: the transcript buffer at
        \mptr{TRANSCRIPT\_MPTR} $=\texttt{0x1000}$, the padded G1 commitment
        windows at \texttt{0xa140} and above, and \mptr{F\_COM\_MPTR} /
        \mptr{PI\_MPTR} at \texttt{0x7a40} / \texttt{0x7ac0}. All are outside the
        VK region.
  \item \textbf{The transcript buffer cannot grow into the VK region.} The proof
        is fixed-length: line 1419 requires
        \yul{eq(0x1e60, calldataload(PROOF\_LEN\_CPTR))} and line 1420 requires
        exactly 19 instances, line 636 pins the ABI head words, and line 1785
        requires the proof cursor to land exactly on \mptr{NUM\_INSTANCE\_CPTR}
        (\revert{BadCalldataShape}). The transcript absorption length is
        therefore a compile-time constant, and the buffer stays below
        \mptr{VK\_MPTR}.
\end{enumerate}

What calldata \emph{can} influence is the \emph{data} the program operates on:
the values at the memory pointers its \opcode{*\_MEM\_U16} operands name. Those
pointers are themselves program bytes, and each is guarded by
\yul{if gt(sub(q\_ptr, 0x3680), 0x6aa0) \{ q\_program\_fail() \}}. There are
\emph{ten} such sites, not the seven a reader might infer from the three
single-pointer opcodes: lines 2216 (\opcode{PUSH\_MEM\_U16}),
2253 (\opcode{ADD\_MEM\_U16}), 2261 (\opcode{MUL\_MEM\_U16}), and inside
\opcode{MODARITH7} lines 2308 (condition pointer), 2344 (LIN7 term), 2357 and
2364 (BILIN7\_ROW left-hand and right-hand pointers), 2411 (extra memory term),
and 2426 and 2427 (extra product term, both operands). The check is a single
\yul{gt} on an unsigned difference, so a pointer below \texttt{0x3680} wraps to
a value far above the bound and is rejected; the admitted set is exactly
$\yul{q\_ptr} \in [\texttt{0x3680}, \texttt{0x3680}+\texttt{0x6aa0}] =
[\texttt{0x3680}, \texttt{0xa120}]$ -- the VK payload, the challenge and
Lagrange scratch, and the evaluation buffer, but not the commitment windows at
\texttt{0xa140} and above and not the VM's own frame at \texttt{0xb140} and
above. (The operands are \yul{u16} in any case, so no pointer can exceed
\texttt{0xffff}.) Every one of these is a \yul{mload}: no \opcode{*\_MEM\_U16}
or \opcode{MODARITH7} operand is ever a store destination. The VM does store
through one decoded operand, but not through a pointer: \opcode{FOLD\_SELECTOR}
writes \yul{mstore(add(SELECTOR\_ACC\_MPTR, shl(5, q\_sel\_idx)), ...)} at lines
3079 and 3087, with \yul{q\_sel\_idx} a program byte clamped to
$[0,10)$ at line 3066 -- so the write set is exactly the ten accumulator words.
\opcode{MODARITH7} additionally guards its two \emph{base} pointers with
the tighter bound \yul{if gt(sub(q\_lhs\_base, 0x3680), 0x69e0)} at lines
2384--2385: because that arm indexes $\mathrm{base} + 32i$ for $i \in [0,7)$, the
bound is reduced by the $\texttt{0xc0}$-byte span of the run so its \emph{last}
word, not its first, is what must stay under \texttt{0xa120}
($\texttt{0x6aa0} - \texttt{0xc0} = \texttt{0x69e0}$). That is the
correct adjustment and it is applied to both operands.

\begin{finding}[QVMS-7]{Quotient VM control flow is derived exclusively from codehash-pinned data}
\textbf{Severity: Observation (clean).} \\
\texttt{Halo2Verifier.sol}:1386--1393, 2013--2016, 2168--2176, 2200--2206,
3090--3105. The interpreter's program pointer, end bound, dispatch byte, operand
widths, stack window, selector-bucket count, and $y$-power table size are all
either Yul literals in the verifier runtime or bytes read out of a VK payload
whose \yul{extcodehash} is checked on every proof. A malicious prover can choose
the field elements the program consumes, but cannot alter which instructions
execute, how many execute, or in what order. Combined with the fail-closed
\yul{default} arm and the three epilogue balance checks at lines 3098--3105, the
worst a prover can achieve inside this VM is a wrong $\nu_y(x)$ derived from its
own evaluations -- which the subsequent PCS opening check rejects. The region is
sound as written.
\end{finding}

\begin{finding}[QVMS-8]{The \texttt{0x6aa0} memory-operand window is far wider than the addresses a correct program names}
\textbf{Severity: Informational.} \\
\texttt{Halo2Verifier.sol}:2216, 2253, 2261, 2308, 2344, 2357, 2364, 2411, 2426,
2427. The guard admits any \yul{q\_ptr} in $[\texttt{0x3680}, \texttt{0xa120}]$
-- a window of \texttt{0x6aa0} bytes spanning the entire VK payload, the
challenge slots at \texttt{0x7900}, the PCS \mptr{ROT\_POINTS\_MPTR} /
\mptr{X1\_POWERS\_MPTR} / \mptr{Q\_EVAL\_SET\_MPTR} scratch, and the evaluation
buffer. It also imposes no word alignment. That is much looser than the set of
addresses a correct program actually names (VK constants, challenges, Lagrange
slots, and evaluation slots), and it is looser than the comment beside it
claims: line 2250 asserts that ``the pointed word is an already range-checked Fr
scalar in verifier memory'', but the admitted window also covers the packed
program bytes themselves $[\texttt{0x50a0},\texttt{0x6280})$, the VK base-point
limbs, and the EIP-2537 padded G1 coordinate words at \mptr{F\_COM\_MPTR}
\texttt{0x7a40} and \mptr{PI\_MPTR} \texttt{0x7ac0} -- none of which is a
canonical $\Fr$ element. A mis-emitted operand therefore lands on an unrelated
but still readable word and produces a well-formed-looking wrong result rather
than reverting. (The arithmetic itself stays well-defined: every use is inside an
\yul{addmod}/\yul{mulmod} modulo $r$, so a non-canonical word is silently
reduced, not misinterpreted.)

\emph{Why this is Informational rather than Low.} The runtime clamp is
explicitly the \emph{second} line of defence, and the first one is exact. At
build time \yul{validate\_quotient\_mem\_ptrs}
(\texttt{src/lowering/quotient\_numerator/vm/mod.rs}, invoked from
\texttt{src/lowering/plan.rs}:347) checks every pointer the finalized bytecode
loads from for membership in a specific named
\yul{QuotientReadWindow} built from the converged memory layout -- window by
window, not against a coarse union -- and
\yul{certify\_quotient\_program} (\texttt{plan.rs}:351) then re-proves that the
bytecode still evaluates the identity tree it was lowered from. Both fail
codegen, so the class of faults this finding describes cannot reach a deployed
artifact through the supported generator path. The runtime bound is deliberately
coarse because it is rendered from a single
$[\yul{operand\_lo}, \yul{operand\_hi}]$ pair
(\texttt{src/lowering/plan.rs}:451--471,
\texttt{src/lowering/render/models.rs}:804--812) and each additional
comparison costs verifier runtime bytes against the EIP-170 budget of
\S\ref{sec:qvms:why}. Note also that the policy is applied uniformly: all ten
sites use the same base \texttt{0x3680} and the same span, and
\opcode{MODARITH7} correctly tightens it to \texttt{0x69e0} for its two strided
base pointers, so this is a single deliberate window rather than an oversight at
any one site. A per-token bound (the plan already carries
\yul{used\_mem\_tokens}) or an alignment check would convert a further class of
silent codegen faults into \revert{QuotientProgramInvalid} reverts at runtime as
well as at build time.
\end{finding}

\subsection{Summary of the phase's pre-VM state}
\label{sec:qvms:summary}

Algorithm~\ref{alg:qvms:setup} states the whole of lines 2007--2176 as
pseudocode. Its postcondition is the precondition for the interpreter section.

\begin{algorithm}[htbp]
\caption{Quotient VM setup and inline prefix (lines 2007--2176)}
\label{alg:qvms:setup}
\begin{algorithmic}[1]
\State $y \gets \mathrm{mem}[\mptr{Y\_MPTR}]$ \Comment{\texttt{0x7980}}
\State $A \gets 0$ \Comment{$A \equiv \mathrm{mem}[\texttt{0xb280}]$}
\For{$\mathit{off} = \texttt{0x00}, \texttt{0x20}, \dots, \texttt{0x0120}$}
  \State $\mathrm{mem}[\texttt{0xb140} + \mathit{off}] \gets 0$ \Comment{10 selector buckets}
\EndFor
\State $\mathrm{mem}[\texttt{0xb2c0}] \gets 1$
\State $t \gets 1$
\For{$i = 1$ \textbf{to} $48$}
  \State $t \gets t \cdot y \bmod r$; \quad $\mathrm{mem}[\texttt{0xb2c0} + 32i] \gets t$
\EndFor
\For{$j = 0$ \textbf{to} $3$} \Comment{inline prefix, unrolled in the source}
  \State $\mathrm{mem}[\texttt{0xb8e0}] \gets I_j$ \Comment{Eqs.~\eqref{eq:qvms:i0}--\eqref{eq:qvms:i3}}
  \State $A \gets A \cdot y \bmod r$
  \State $s \gets \mathit{bucket}(j)$; \; $g \gets \mathit{gap}(j)$ \Comment{Table~\ref{tab:qvms:prefixfolds}}
  \If{$g \neq 0$} \State $B_s \gets B_s \cdot \mathrm{mem}[\texttt{0xb2c0} + 32g] \bmod r$ \EndIf
  \State $B_s \gets B_s + \mathrm{mem}[\texttt{0xb8e0}] \bmod r$
\EndFor
\State $\yul{q\_pc} \gets \texttt{0x50a0}$; \; $\yul{q\_end} \gets \texttt{0x626f}$
\State $\yul{q\_sp} \gets \texttt{0xb8e0}$; \; $\yul{q\_top} \gets 0$; \; $\yul{q\_has\_top} \gets 0$
\end{algorithmic}
\end{algorithm}

\noindent
On entry to the interpreter loop at line 2200 the following all hold:
$A = 0$ with the global batch position at $4$;
$B_0 = I_0$, $B_1 = I_1$, $B_2 = I_2 y + I_3$, $B_3 = \dots = B_9 = 0$;
$\mathrm{mem}[\texttt{0xb2c0} + 32i] = y^i$ for $0 \le i \le 48$;
the operand stack is empty and balanced at its base;
and \yul{q\_pc} addresses the first byte of a 4{,}559-byte program whose contents
are fixed by \yul{EXPECTED\_VK\_CODEHASH}.


%% ==== 12-qvm-arith.tex ==================================================
\section{Quotient VM arithmetic opcodes}
\label{sec:qvma:main}

This section specifies the arithmetic instruction set of the compact quotient
virtual machine (VM) interpreted inside \yul{verifyProof}, covering
\texttt{Halo2Verifier.sol} lines 2200--2452: the dispatch loop and the seven
opcodes \opcode{PUSH\_MEM\_U16} (\texttt{0x05}), \opcode{ADD} (\texttt{0x06}),
\opcode{NEG} (\texttt{0x08}), \opcode{MUL\_CONST\_U8} (\texttt{0x0d}),
\opcode{ADD\_MEM\_U16} (\texttt{0x10}), \opcode{MUL\_MEM\_U16} (\texttt{0x11})
and \opcode{MODARITH7} (\texttt{0x21}). The remaining opcodes the program may
contain --- \opcode{FOLD\_SELECTOR} (\texttt{0x0b}, lines 3053--3088),
\opcode{NATIVE\_PERMUTATION} (\texttt{0x19}, line 2453),
\opcode{NATIVE\_IDENTITY} (\texttt{0x1b}, line 2762) and
\opcode{NATIVE\_LOOKUP} (\texttt{0x1f}, line 2570) --- are covered elsewhere;
they are referenced here only where the arithmetic opcodes' invariants depend
on them.

All excerpts below are copied verbatim from the fixture and dedented by a
uniform per-excerpt amount; no other change is made. In particular the visibly
inconsistent indentation of the \yul{if gt(sub(...)) \{ q\_program\_fail() \}}
guard lines is reproduced as it appears in the artifact --- the generator emits
that guard from a macro with a fixed 24-space indent regardless of nesting
depth (\texttt{src/lowering/quotient\_numerator/vm/yul\_arms.rs} lines
186--197, \yul{ptr\_guard}), so the misalignment is cosmetic, not a scoping
error. The two \texttt{BILIN7\_PAIRWISE} guards come from a different macro
(\yul{vec7\_guard}, lines 198--207) with a fixed 28-space indent, which is why
those two --- and only those two --- appear correctly aligned with their
surrounding block.

\subsection{Machine state and environment constants}
\label{sec:qvma:state}

The VM is a single-accumulator stack machine over $\Fr$ with a cached
top-of-stack. Its five registers are declared at lines 2168--2176, immediately
above the assigned range:

\begin{itemize}
  \item \yul{q\_pc} --- byte cursor into the packed bytecode, initialised to
        \yul{q\_program\_mptr};
  \item \yul{q\_end} --- exclusive end cursor,
        $\yul{q\_program\_mptr} + \texttt{0x11cf}$;
  \item \yul{q\_sp} --- memory stack pointer for spilled operands, initialised
        to \texttt{0xb8e0};
  \item \yul{q\_top} --- the cached top-of-stack value, initialised to $0$;
  \item \yul{q\_has\_top} --- a Boolean flag, initialised to $0$, recording
        whether \yul{q\_top} currently holds a live stack value.
\end{itemize}

Write $T$ for \yul{q\_top}, $h$ for \yul{q\_has\_top}, $\sigma$ for \yul{q\_sp},
$M[a]$ for the 32-byte EVM memory word at address $a$, and
$K[c] = M[\yul{q\_const\_mptr} + 32c]$ for entry $c$ of the quotient constant
table. All field arithmetic is performed with \yul{addmod}/\yul{mulmod} modulo
$r = \texttt{FR\_MODULUS}$ (line 331), which is bound to the local \yul{r} at
line 1355.

\begin{table}[htbp]
\centering
\caption{Environment constants read by the arithmetic opcodes, with the
fixture line that fixes each one.}
\label{tab:qvma:env}
\begin{tabular}{@{}llll@{}}
\toprule
Symbol & Value & Meaning & Line \\
\midrule
\yul{q\_const\_mptr}   & \texttt{0x3a60} & base of the 178-word $\Fr$ constant table & 2013 \\
\yul{q\_program\_mptr} & \texttt{0x50a0} & base of the packed bytecode & 2016 \\
program length       & \texttt{0x11cf} & 4559 bytes & 2170 \\
\yul{q\_end}          & \texttt{0x626f} & exclusive end cursor & 2170 \\
stack base           & \texttt{0xb8e0} & initial \yul{q\_sp}; also the underflow sentinel & 2172, 2229 \\
stack ceiling        & \texttt{0xbfa0} & exclusive; overflow sentinel & 2218, 2332 \\
$\mathrm{lo}$        & \texttt{0x3680} & operand-window low bound ($=\mptr{VK\_MPTR}$) & 2216 et al. \\
$\mathrm{hi}-\mathrm{lo}$ & \texttt{0x6aa0} & operand-window span (scalar form) & 2216 et al. \\
$\mathrm{hi}$        & \texttt{0xa120} & inclusive top of the window (derived) & --- \\
$\mathrm{hi}-\mathrm{lo}-\texttt{0xc0}$ & \texttt{0x69e0} & window span for 7-word limb vectors & 2384--2385 \\
\bottomrule
\end{tabular}
\end{table}

The constant table and the bytecode both live inside the VK payload:
\mptr{VK\_MPTR} $=\texttt{0x3680}$ and \texttt{EXPECTED\_VK\_PAYLOAD\_LENGTH}
$=17024=\texttt{0x4280}$ (lines 92, 128), so the payload occupies
$[\texttt{0x3680}, \texttt{0x7900})$. Within it, \texttt{Halo2VerifyingKey.sol}
places 178 \texttt{quotient\_const} words at payload offsets
$\texttt{0x03e0}\dots\texttt{0x1a00}$ (first word addresses
$\texttt{0x3a60}\dots\texttt{0x5080}$ in memory, so the table occupies
$[\texttt{0x3a60}, \texttt{0x50a0})$) and 143 \texttt{quotient\_program}
words at payload offsets $\texttt{0x1a20}\dots\texttt{0x2be0}$ (first word
addresses $\texttt{0x50a0}\dots\texttt{0x6260}$, so the program region is
$[\texttt{0x50a0}, \texttt{0x6280})$, of which the first \texttt{0x11cf}
bytes are program and the trailing 17 bytes are zero padding). The two
regions are contiguous: the constant table ends exactly where the bytecode
begins, which is what makes the containment argument of
Finding~\ref{find:qvma:constidx} work. Both figures were re-derived from
\texttt{Halo2VerifyingKey.sol} by counting the \texttt{quotient\_const} /
\texttt{quotient\_program} \yul{mstore} lines (178 and 143 respectively,
offsets contiguous with no gap).

The whole payload is copied into memory by \yul{extcodecopy} at line 1393,
after a per-proof re-check of both \yul{extcodesize} and \yul{extcodehash}
against \texttt{EXPECTED\_VK\_LENGTH} / \texttt{EXPECTED\_VK\_CODEHASH\_WORD}
(lines 1386--1389). This is the pin that every ``the program is trusted''
argument below rests on: the bytecode, the constant table, and every operand
byte the VM decodes are covered by a hash checked on every single call, not
only at construction (line 582).

\subsection{The dispatch loop}
\label{sec:qvma:dispatch}

\begin{lstlisting}[style=solidity,caption={\texttt{Halo2Verifier.sol} lines 2200--2208: the VM dispatch loop header},label={lst:qvma:dispatch}]
for { } lt(q_pc, q_end) { } {
    // The bytecode table is byte-addressed, but EVM memory
    // loads whole words. `byte(0, mload(q_pc))` extracts the
    // opcode at the current byte cursor; each case advances
    // q_pc by exactly its operand width.
    let q_op := byte(0, mload(q_pc))
    q_pc := add(q_pc, 1)

    switch q_op
\end{lstlisting}

The loop condition is \yul{lt(q\_pc, q\_end)}: strictly less-than, so an
instruction is entered only while at least one byte remains. The opcode is
extracted with \yul{byte(0, mload(q\_pc))} --- a full 32-byte load of which only
the leading byte is used --- and \yul{q\_pc} is advanced by one before the
\yul{switch}. Each case then advances \yul{q\_pc} by exactly its own operand
width, so the encoding is self-delimiting: there is no length prefix and no
per-operand tag.

Two structural properties follow, and both matter for the guard analysis:

\begin{enumerate}
  \item The loop bound is checked \emph{only} at instruction boundaries.
        Nothing inside a case checks that
        $\yul{q\_pc} + \text{width} \le \yul{q\_end}$, so a final instruction
        whose operands over-run the program reads bytes beyond \yul{q\_end}.
  \item Every arithmetic case that terminates normally leaves \yul{q\_pc}
        pointing at the next opcode; the epilogue at lines 3098--3105
        reconciles the three terminal invariants
        $\yul{q\_pc} = \yul{q\_end}$, $h = 0$, and
        $\sigma = \texttt{0xb8e0}$, calling \yul{q\_program\_fail()} otherwise.
\end{enumerate}

Every guard in this section funnels into the same helper, lines 1890--1893:

\begin{lstlisting}[style=solidity,caption={\texttt{Halo2Verifier.sol} lines 1890--1893: the single failure exit of the quotient VM},label={lst:qvma:fail}]
function q_program_fail() {
    mstore(0x00, shl(224, 0x3cc81b89))
    revert(0x00, 0x04)
}
\end{lstlisting}

\noindent
This reverts with the 4-byte selector \texttt{0x3cc81b89}. That value was
recomputed independently for this specification: the first four bytes of
$\mathrm{keccak256}$ of the signature string
\texttt{QuotientProgramInvalid()} are
\texttt{3c c8 1b 89}, matching both the literal at line 1891 and the error
declared as \revert{QuotientProgramInvalid} at line 77. The store is
\yul{mstore(0x00, shl(224, 0x3cc81b89))} followed by
\yul{revert(0x00, 0x04)}: the selector is left-aligned into the first word of
the EVM scratch area and exactly four bytes are returned, so the revert data
is a bare selector with no ABI arguments. There is no recovery path: the VM
is success-or-revert, and no arithmetic opcode has a ``skip'' or ``clamp''
branch.

\subsection{Instruction encoding summary}
\label{sec:qvma:encoding}

\begin{table}[htbp]
\centering
\caption{Encoding and stack effect of the seven arithmetic opcodes. ``Operand
bytes'' excludes the opcode byte; ``Width'' includes it. $S$ denotes the
spilled stack, $T$ the cached top.}
\label{tab:qvma:opcodes}
\begin{tabular}{@{}llllll@{}}
\toprule
Byte & Mnemonic & Operand bytes & Width & Stack effect & Semantics \\
\midrule
\texttt{0x05} & \opcode{PUSH\_MEM\_U16}  & u16 pointer                  & 3        & push $1$      & $T \gets M[p]$ \\
\texttt{0x06} & \opcode{ADD}             & ---                          & 1        & pop $1$ spill & $T \gets M[\sigma-32] + T$ \\
\texttt{0x08} & \opcode{NEG}             & ---                          & 1        & none          & $T \gets -T$ \\
\texttt{0x0d} & \opcode{MUL\_CONST\_U8}  & u8 constant index            & 2        & none          & $T \gets T \cdot K[c]$ \\
\texttt{0x10} & \opcode{ADD\_MEM\_U16}   & u16 pointer                  & 3        & none          & $T \gets T + M[p]$ \\
\texttt{0x11} & \opcode{MUL\_MEM\_U16}   & u16 pointer                  & 3        & none          & $T \gets T \cdot M[p]$ \\
\texttt{0x21} & \opcode{MODARITH7}       & variable, see \S\ref{sec:qvma:modarith7} & $\ge 7$ & push $1$ & $T \gets A$ (Eq.~\ref{eq:qvma:modarith7}) \\
\bottomrule
\end{tabular}
\end{table}

Every arithmetic result is reduced modulo $r$ by the \yul{addmod}/\yul{mulmod}
that produces it. The operand fields are decoded with the following exact
expressions; each row lists the Yul the fixture uses, the bytes it consumes,
and the byte positions (relative to \yul{q\_pc}) it extracts.

\begin{table}[htbp]
\centering
\caption{Bit-exact operand extraction. \yul{w} abbreviates \yul{mload(q\_pc)}.}
\label{tab:qvma:extract}
\begin{tabular}{@{}llll@{}}
\toprule
Field & Yul expression & Bytes consumed & Bytes extracted \\
\midrule
u16 pointer            & \yul{shr(240, w)}                  & 2 & $0..1$ \\
u8 constant index      & \yul{byte(0, w)}                   & 1 & $0$ \\
(coeff, ptr) triple    & \yul{byte(0, w)}; \yul{and(shr(232, w), 0xffff)} & 3 & $0$; $1..2$ \\
pairwise base pair     & \yul{shr(240, w)}; \yul{and(shr(224, w), 0xffff)} & 4 & $0..1$; $2..3$ \\
(coeff, lhs, rhs) quint& \yul{byte(0, w)}; \yul{and(shr(232, w), 0xffff)}; \yul{and(shr(216, w), 0xffff)} & 5 & $0$; $1..2$; $3..4$ \\
pairwise coefficient   & \yul{byte(0, mload(add(q\_coeff\_pc, add(q\_i, q\_j))))} & 13 per block & offset $i+j$ \\
counts word            & \yul{byte(k, mload(q\_pc))}, $k = 0..4$ & 5 & $0..4$ \\
\bottomrule
\end{tabular}
\end{table}

All of these read a full 32-byte word and discard the tail. That is the whole
point of the byte-oriented encoding: EVM memory has no sub-word load, so the
generator packs operands adjacently and pays one \yul{mload} plus a shift or
\yul{byte} per field instead of one word per field. The cost is that a load
near the end of the program reads bytes belonging to the next instruction or,
at the very end, past \yul{q\_end}; since only the extracted bits are used, this
is harmless except for the over-read case discussed in
Finding~\ref{find:qvma:overread}.

\subsection{\texorpdfstring{\texttt{0x05} \opcode{PUSH\_MEM\_U16}}{0x05 PUSH\_MEM\_U16}}
\label{sec:qvma:push}

\begin{lstlisting}[style=solidity,caption={\texttt{Halo2Verifier.sol} lines 2210--2224: \opcode{PUSH\_MEM\_U16}},label={lst:qvma:push}]
case 0x05 {
    // Operand layout: u16 absolute memory pointer. The
    // memory planner keeps the hot quotient frame below
    // 64 KiB when this compact form is emitted.
    let q_ptr := shr(240, mload(q_pc))
    q_pc := add(q_pc, 2)
    if gt(sub(q_ptr, 0x3680), 0x6aa0) { q_program_fail() }
    if q_has_top {
        if iszero(lt(q_sp, 0xbfa0)) { q_program_fail() }
        mstore(q_sp, q_top)
        q_sp := add(q_sp, 0x20)
    }
    q_top := mload(q_ptr)
    q_has_top := 1
}
\end{lstlisting}

Let $p = \yul{shr(240, mload(q\_pc))}$ be the 16-bit operand. The opcode
performs, in this order:

\begin{enumerate}
  \item advance $\yul{q\_pc} \gets \yul{q\_pc} + 2$;
  \item \textbf{operand-window guard}: if
        $(p - \texttt{0x3680}) \bmod 2^{256} > \texttt{0x6aa0}$ then
        \yul{q\_program\_fail()}. Because \yul{sub} wraps, this single unsigned
        comparison enforces
        $\texttt{0x3680} \le p \le \texttt{0xa120}$;
  \item if $h \ne 0$: \textbf{stack-overflow guard} --- if
        $\sigma \ge \texttt{0xbfa0}$ then \yul{q\_program\_fail()}; otherwise
        $M[\sigma] \gets T$ and $\sigma \gets \sigma + 32$;
  \item $T \gets M[p]$ and $h \gets 1$.
\end{enumerate}

Note that the spill is conditional on $h$: pushing onto an empty cached top
costs no memory traffic, which is the reason for carrying $h$ at all. The
loaded word is \emph{not} reduced modulo $r$ (see
Finding~\ref{find:qvma:unreduced}).

\subsection{\texorpdfstring{\texttt{0x06} \opcode{ADD}, \texttt{0x08} \opcode{NEG}, \texttt{0x0d} \opcode{MUL\_CONST\_U8}}{0x06 ADD, 0x08 NEG, 0x0d MUL\_CONST\_U8}}
\label{sec:qvma:unary}

\begin{lstlisting}[style=solidity,caption={\texttt{Halo2Verifier.sol} lines 2226--2246: \opcode{ADD}, \opcode{NEG} and \opcode{MUL\_CONST\_U8}},label={lst:qvma:unary}]
case 0x06 {
    // The safety validator guarantees a spilled operand
    // exists before ADD. q_top is the right operand.
    if eq(q_sp, 0xb8e0) { q_program_fail() }
    q_sp := sub(q_sp, 0x20)
    q_top := addmod(mload(q_sp), q_top, r)
}
// VM 0x08 NEG: replace q_top with its Fr negation.
case 0x08 {
    // addmod(0, r - x, r) maps zero back to zero and every
    // nonzero scalar to its canonical additive inverse.
    q_top := addmod(0, sub(r, q_top), r)
}
// VM 0x0d MUL_CONST_U8: multiply q_top by a small constant-table slot.
case 0x0d {
    // One-byte constant-index multiply, used by short
    // affine chains after an initial PUSH.
    let qconst := byte(0, mload(q_pc))
    q_pc := add(q_pc, 1)
    q_top := mulmod(q_top, mload(add(q_const_mptr, shl(5, qconst))), r)
}
\end{lstlisting}

\paragraph{\opcode{ADD} (\texttt{0x06}).}
\textbf{Stack-underflow guard}: if $\sigma = \texttt{0xb8e0}$ ---
i.e.\ the spilled stack is empty --- then \yul{q\_program\_fail()}. Otherwise
\[
  \sigma \gets \sigma - 32, \qquad
  T \gets \bigl(M[\sigma] + T\bigr) \bmod r .
\]
The spilled word is the left operand and the cached top is the right operand;
addition is commutative so the order is immaterial here, but it fixes the
convention that \yul{q\_top} is always the most recently produced value. The
arm does not read $h$ and does not set $h$
(Finding~\ref{find:qvma:addtop}).

\paragraph{\opcode{NEG} (\texttt{0x08}).}
\[
  T \gets \bigl(0 + (r - T)\bigr) \bmod r ,
\]
computed as \yul{addmod(0, sub(r, q\_top), r)}. For $T \in [0, r]$ this is the
canonical additive inverse and maps $0 \mapsto 0$ (since
$\yul{sub}(r, 0) = r$ and $r \bmod r = 0$). Correctness depends on
$T \le r$: for $T > r$ the \yul{sub} wraps to $2^{256} + r - T$ and the result
is $(2^{256} - T) \bmod r$, which is not $-T$. See
Finding~\ref{find:qvma:unreduced}.

\paragraph{\opcode{MUL\_CONST\_U8} (\texttt{0x0d}).}
Let $c = \yul{byte(0, mload(q\_pc))}$; then
$\yul{q\_pc} \gets \yul{q\_pc} + 1$ and
\[
  T \gets \bigl(T \cdot M[\texttt{0x3a60} + 32c]\bigr) \bmod r .
\]
The index is a raw byte with no range check; see
Finding~\ref{find:qvma:constidx}.

\subsection{\texorpdfstring{\texttt{0x10} \opcode{ADD\_MEM\_U16} and \texttt{0x11} \opcode{MUL\_MEM\_U16}}{0x10 ADD\_MEM\_U16 and 0x11 MUL\_MEM\_U16}}
\label{sec:qvma:memops}

\begin{lstlisting}[style=solidity,caption={\texttt{Halo2Verifier.sol} lines 2247--2263: the fused memory-operand arms},label={lst:qvma:memops}]
// VM 0x10 ADD_MEM_U16: add a short memory load into q_top.
case 0x10 {
    // Operand layout: u16 pointer. The pointed word is an
    // already range-checked Fr scalar in verifier memory.
    let q_ptr := shr(240, mload(q_pc))
    q_pc := add(q_pc, 2)
    if gt(sub(q_ptr, 0x3680), 0x6aa0) { q_program_fail() }
    q_top := addmod(q_top, mload(q_ptr), r)
}
// VM 0x11 MUL_MEM_U16: multiply q_top by a short memory load.
case 0x11 {
    // In-place multiply by a planned memory word.
    let q_ptr := shr(240, mload(q_pc))
    q_pc := add(q_pc, 2)
    if gt(sub(q_ptr, 0x3680), 0x6aa0) { q_program_fail() }
    q_top := mulmod(q_top, mload(q_ptr), r)
}
\end{lstlisting}

Both decode $p = \yul{shr(240, mload(q\_pc))}$, advance \yul{q\_pc} by 2, apply
the identical operand-window guard
$(p - \texttt{0x3680}) \bmod 2^{256} \le \texttt{0x6aa0}$, and then compute
\[
  \texttt{0x10}: \; T \gets (T + M[p]) \bmod r,
  \qquad
  \texttt{0x11}: \; T \gets (T \cdot M[p]) \bmod r .
\]
Neither touches $\sigma$, $h$, or the spilled stack: they are strictly in-place
updates of the cached top. This is what makes the cached-top design pay ---
the dominant expression shape (a Horner-style chain of alternating
multiply-by-memory and add-memory) executes with zero memory stores.

\subsection{\texorpdfstring{\texttt{0x21} \opcode{MODARITH7}}{0x21 MODARITH7}}
\label{sec:qvma:modarith7}

\opcode{MODARITH7} is a single fused instruction that evaluates an entire
affine/bilinear foreign-field identity. It is by far the most complex opcode
and, in this artifact, the dominant one: it accounts for 4331 of the 4559
program bytes (95.0\%).

\subsubsection{Header: flags, optional condition, optional seed}

\begin{lstlisting}[style=solidity,caption={\texttt{Halo2Verifier.sol} lines 2300--2335: \opcode{MODARITH7} header decode, counts word, and the push spill},label={lst:qvma:ma-header}]
let q_flags := byte(0, mload(q_pc))
q_pc := add(q_pc, 1)
let q_cond_ptr := 0
if and(q_flags, 0x01) {
    // Optional condition pointer. When present, the
    // whole identity is gated by mload(q_cond_ptr).
    q_cond_ptr := shr(240, mload(q_pc))
    q_pc := add(q_pc, 2)
if gt(sub(q_cond_ptr, 0x3680), 0x6aa0) { q_program_fail() }
}

let q_acc := 0
if and(q_flags, 0x02) {
    // Optional constant seed for affine identities
    // with a standalone constant term.
    let qconst := byte(0, mload(q_pc))
    q_pc := add(q_pc, 1)
    q_acc := mload(add(q_const_mptr, shl(5, qconst)))
}

// Five one-byte counters describe the blocks that
// follow. Each block has a fixed-width internal layout,
// so q_pc can advance without per-term tags.
let q_counts_word := mload(q_pc)
let q_lin_count := byte(0, q_counts_word)
let q_row_count := byte(1, q_counts_word)
let q_pairwise_count := byte(2, q_counts_word)
let q_mem_count := byte(3, q_counts_word)
let q_product_count := byte(4, q_counts_word)
q_pc := add(q_pc, 5)

if q_has_top {
    if iszero(lt(q_sp, 0xbfa0)) { q_program_fail() }
    mstore(q_sp, q_top)
    q_sp := add(q_sp, 0x20)
}
\end{lstlisting}

The header is decoded in this exact order.

\begin{enumerate}
  \item $f \gets \yul{byte(0, mload(q\_pc))}$; $\yul{q\_pc} \mathrel{+}= 1$.
        Only two bits are defined:
        \begin{itemize}
          \item bit 0 (\texttt{and(q\_flags, 0x01)}): a condition pointer
                follows, and the final sum is multiplied by the word it
                addresses;
          \item bit 1 (\texttt{and(q\_flags, 0x02)}): a one-byte constant-table
                index follows, seeding the accumulator.
        \end{itemize}
        Bits 2--7 are never examined (Finding~\ref{find:qvma:flagbits}).
  \item If bit 0 is set:
        $p_{\mathrm{cond}} \gets \yul{shr(240, mload(q\_pc))}$;
        $\yul{q\_pc} \mathrel{+}= 2$; then the operand-window guard at line
        2308. If bit 0 is clear, $p_{\mathrm{cond}}$ stays $0$ and is never
        dereferenced (line 2439 re-tests the same bit).
  \item $A \gets 0$. If bit 1 is set:
        $s \gets \yul{byte(0, mload(q\_pc))}$; $\yul{q\_pc} \mathrel{+}= 1$;
        $A \gets K[s]$.
  \item The five block counters are read from one word:
        \[
          n_L = \mathrm{byte}_0(w_c), \quad
          n_R = \mathrm{byte}_1(w_c), \quad
          n_P = \mathrm{byte}_2(w_c), \quad
          n_M = \mathrm{byte}_3(w_c), \quad
          n_X = \mathrm{byte}_4(w_c),
        \]
        where $w_c$ is the word \yul{q\_counts\_word} loaded at line 2323 and
        $\mathrm{byte}_k$ is the Yul \yul{byte(k, ...)} selector; then
        $\yul{q\_pc} \mathrel{+}= 5$.
        Each counter is an unguarded \texttt{u8} in $[0, 255]$.
  \item \textbf{Stack-overflow guard and spill}: if $h \ne 0$ then, exactly as
        in \opcode{PUSH\_MEM\_U16}, fail unless $\sigma < \texttt{0xbfa0}$,
        then $M[\sigma] \gets T$, $\sigma \gets \sigma + 32$.
\end{enumerate}

Placing the spill after the counts word --- rather than at the very start of
the arm --- has no semantic effect (nothing between the two points reads $T$ or
$\sigma$); it only changes which guard fires first on a malformed program.

\begin{table}[htbp]
\centering
\caption{\opcode{MODARITH7} counts word: byte position, extraction, and the
fixed record width of the block it counts.}
\label{tab:qvma:counts}
\begin{tabular}{@{}lllll@{}}
\toprule
Byte & Extraction & Counter & Block & Bytes per block \\
\midrule
0 & \yul{byte(0, q\_counts\_word)} & \yul{q\_lin\_count}      & \texttt{LIN7}            & $7 \times 3 = 21$ \\
1 & \yul{byte(1, q\_counts\_word)} & \yul{q\_row\_count}      & \texttt{BILIN7\_ROW}     & $2 + 7 \times 3 = 23$ \\
2 & \yul{byte(2, q\_counts\_word)} & \yul{q\_pairwise\_count} & \texttt{BILIN7\_PAIRWISE}& $4 + 13 = 17$ \\
3 & \yul{byte(3, q\_counts\_word)} & \yul{q\_mem\_count}      & extra linear term        & $3$ \\
4 & \yul{byte(4, q\_counts\_word)} & \yul{q\_product\_count}  & extra product term       & $5$ \\
\bottomrule
\end{tabular}
\end{table}

The total encoded width of one \opcode{MODARITH7} instruction is therefore
\begin{equation}
  \mathrm{width} = 1 + 1 + 2\,[f_0] + 1\,[f_1] + 5
  + 21 n_L + 23 n_R + 17 n_P + 3 n_M + 5 n_X ,
  \label{eq:qvma:width}
\end{equation}
where $[f_0], [f_1] \in \{0,1\}$ are the flag bits. The minimum is 7 bytes
(both flags clear, all counters zero); the maximum a \texttt{u8}-counted
instruction can reach is
$10 + 255\,(21+23+17+3+5) = 17{,}605$ bytes.

\subsubsection{\texttt{LIN7} and \texttt{BILIN7\_ROW} blocks}

\begin{lstlisting}[style=solidity,caption={\texttt{Halo2Verifier.sol} lines 2337--2375: the \texttt{LIN7} and \texttt{BILIN7\_ROW} loops},label={lst:qvma:ma-lin}]
// LIN7 blocks: q_acc += sum_i c_i * limb_i.
for { let q_lin_block := 0 } lt(q_lin_block, q_lin_count) { q_lin_block := add(q_lin_block, 1) } {
    for { let q_i := 0 } lt(q_i, 7) { q_i := add(q_i, 1) } {
        let q_word := mload(q_pc)
        let qconst := byte(0, q_word)
        let q_ptr := and(shr(232, q_word), 0xffff)
        q_pc := add(q_pc, 3)
if gt(sub(q_ptr, 0x3680), 0x6aa0) { q_program_fail() }
        q_acc := addmod(
            q_acc,
            mulmod(mload(add(q_const_mptr, shl(5, qconst))), mload(q_ptr), r),
            r
        )
    }
}

// BILIN7_ROW blocks: q_acc += lhs * sum_i c_i * rhs_i.
for { let q_row_block := 0 } lt(q_row_block, q_row_count) { q_row_block := add(q_row_block, 1) } {
    let q_lhs := shr(240, mload(q_pc))
    q_pc := add(q_pc, 2)
if gt(sub(q_lhs, 0x3680), 0x6aa0) { q_program_fail() }
    let q_lhs_value := mload(q_lhs)
    for { let q_i := 0 } lt(q_i, 7) { q_i := add(q_i, 1) } {
        let q_word := mload(q_pc)
        let qconst := byte(0, q_word)
        let q_rhs := and(shr(232, q_word), 0xffff)
        q_pc := add(q_pc, 3)
if gt(sub(q_rhs, 0x3680), 0x6aa0) { q_program_fail() }
        q_acc := addmod(
            q_acc,
            mulmod(
                mulmod(q_lhs_value, mload(q_rhs), r),
                mload(add(q_const_mptr, shl(5, qconst))),
                r
            ),
            r
        )
    }
}
\end{lstlisting}

A \texttt{LIN7} block is exactly seven $(\text{coeff}, \text{ptr})$ triples ---
the loop bound \yul{lt(q\_i, 7)} is a literal, not an operand, which is what
makes the block width fixed at 21 bytes. Each triple contributes
$K[c] \cdot M[p]$ to $A$, with the operand-window guard at line 2344 applied to
$p$ before the \yul{mload}.

A \texttt{BILIN7\_ROW} block is a 2-byte left-hand pointer $\ell$ (guarded at
line 2357, then loaded \emph{once} into \yul{q\_lhs\_value}) followed by seven
$(\text{coeff}, \text{ptr})$ triples (each guarded at line 2364). Each triple
contributes $M[\ell] \cdot M[q] \cdot K[c]$. Hoisting \yul{mload(q\_lhs)} out of
the inner loop saves six \yul{mload}s per block; the guard on $\ell$ is
likewise evaluated once.

\subsubsection{\texttt{BILIN7\_PAIRWISE} blocks}

\begin{lstlisting}[style=solidity,caption={\texttt{Halo2Verifier.sol} lines 2377--2403: the 7-by-7 weighted convolution},label={lst:qvma:ma-pair}]
// BILIN7_PAIRWISE blocks: q_acc += weighted 7-by-7
// product convolution.
for { let q_pair_block := 0 } lt(q_pair_block, q_pairwise_count) { q_pair_block := add(q_pair_block, 1) } {
    let q_pair_word := mload(q_pc)
    let q_lhs_base := shr(240, q_pair_word)
    let q_rhs_base := and(shr(224, q_pair_word), 0xffff)
    q_pc := add(q_pc, 0x04)
    if gt(sub(q_lhs_base, 0x3680), 0x69e0) { q_program_fail() }
    if gt(sub(q_rhs_base, 0x3680), 0x69e0) { q_program_fail() }
    let q_coeff_pc := q_pc
    q_pc := add(q_pc, 13)
    for { let q_i := 0 } lt(q_i, 7) { q_i := add(q_i, 1) } {
        let q_lhs_value := mload(add(q_lhs_base, shl(5, q_i)))
        for { let q_j := 0 } lt(q_j, 7) { q_j := add(q_j, 1) } {
            let qconst := byte(0, mload(add(q_coeff_pc, add(q_i, q_j))))
            q_acc := addmod(
                q_acc,
                mulmod(
                    mulmod(q_lhs_value, mload(add(q_rhs_base, shl(5, q_j))), r),
                    mload(add(q_const_mptr, shl(5, qconst))),
                    r
                ),
                r
            )
        }
    }
}
\end{lstlisting}

This is the only opcode form that addresses memory as a \emph{vector}. The
4-byte header supplies two base pointers,
$\alpha = \yul{shr(240, w)}$ and $\beta = \yul{and(shr(224, w), 0xffff)}$,
each of which names a run of seven consecutive words
$M[\alpha], M[\alpha+32], \dots, M[\alpha + 192]$. Accordingly the guard uses
the tightened span \texttt{0x69e0}
$= \texttt{0x6aa0} - \texttt{0xc0}$, so that
$\texttt{0x3680} \le \alpha \le \texttt{0xa060}$ and the highest word actually
loaded, $M[\alpha + \texttt{0xc0}]$, still ends at \texttt{0xa140} --- exactly
the top of the scalar window. The same holds for $\beta$. This is a precisely
tight bound, not a conservative one. Note that the guard is a \emph{range}
check only: it does not require $\alpha$ or $\beta$ to be 32-byte aligned
(Finding~\ref{find:qvma:align}).

Unlike \texttt{BILIN7\_ROW}, this arm hoists only the left vector: line 2389
loads $M[\alpha + 32i]$ once per outer iteration, but line 2395 reloads
$M[\beta + 32j]$ inside the inner loop, so the seven right-hand limbs are
loaded 49 times rather than 7. The coefficient byte is likewise re-loaded per
term at line 2391. That is 91 \yul{mload}s per block where 63 would do; it
costs roughly 100 gas per block and is invisible in this artifact because the
block never executes (Finding~\ref{find:qvma:deadpair}).

The 13 bytes following the bases are constant-table indexes
$d_0, \dots, d_{12}$. \yul{q\_coeff\_pc} is snapshotted \emph{before}
\yul{q\_pc} is advanced by 13, and the coefficient for the limb pair $(i, j)$ is
read as \yul{byte(0, mload(add(q\_coeff\_pc, add(q\_i, q\_j))))}: byte $i+j$ of the
13-byte table, with $i + j \in [0, 12]$. That indexing is the signature of a
base-$2^{\mathrm{LOG2\_BASE}}$ limb convolution --- the coefficient of
$\text{lhs}_i \cdot \text{rhs}_j$ depends only on $i+j$, matching
$\mathrm{base}^{i+j} \bmod m$ (the fixture cites
\texttt{circuits/src/field/foreign/params.rs::double\_base\_powers} at line
2273). So 17 encoded bytes drive 49 weighted product terms, i.e.\ 98
\yul{mulmod}s and 49 \yul{addmod}s:
\[
  A \mathrel{+}= \sum_{i=0}^{6} \sum_{j=0}^{6}
      K[d_{i+j}] \cdot M[\alpha + 32i] \cdot M[\beta + 32j] .
\]

\subsubsection{Extra linear and product terms, and the condition}

\begin{lstlisting}[style=solidity,caption={\texttt{Halo2Verifier.sol} lines 2405--2447: extra terms, the optional gate condition, and the push},label={lst:qvma:ma-tail}]
    // Extra linear memory terms outside the 7-limb shapes.
    for { let q_mem_block := 0 } lt(q_mem_block, q_mem_count) { q_mem_block := add(q_mem_block, 1) } {
        let q_word := mload(q_pc)
        let qconst := byte(0, q_word)
        let q_ptr := and(shr(232, q_word), 0xffff)
        q_pc := add(q_pc, 3)
    if gt(sub(q_ptr, 0x3680), 0x6aa0) { q_program_fail() }
        q_acc := addmod(
            q_acc,
            mulmod(mload(add(q_const_mptr, shl(5, qconst))), mload(q_ptr), r),
            r
        )
    }

    // Extra binary product terms outside the 7-limb shapes.
    for { let q_product_block := 0 } lt(q_product_block, q_product_count) { q_product_block := add(q_product_block, 1) } {
        let q_word := mload(q_pc)
        let qconst := byte(0, q_word)
        let q_lhs := and(shr(232, q_word), 0xffff)
        let q_rhs := and(shr(216, q_word), 0xffff)
        q_pc := add(q_pc, 5)
    if gt(sub(q_lhs, 0x3680), 0x6aa0) { q_program_fail() }
    if gt(sub(q_rhs, 0x3680), 0x6aa0) { q_program_fail() }
        q_acc := addmod(
            q_acc,
            mulmod(
                mulmod(mload(q_lhs), mload(q_rhs), r),
                mload(add(q_const_mptr, shl(5, qconst))),
                r
            ),
            r
        )
    }

    if and(q_flags, 0x01) {
        // Apply the optional gate condition last so every
        // subterm shares the same selector/condition.
        q_acc := mulmod(mload(q_cond_ptr), q_acc, r)
    }
    // MODARITH7 pushes its fused identity value.
    q_top := q_acc
    q_has_top := 1
}
\end{lstlisting}

An ``extra linear'' record is a 3-byte $(\text{coeff}, \text{ptr})$ triple
contributing $K[c] \cdot M[p]$ --- identical in shape to one \texttt{LIN7}
term, but individually counted rather than grouped in sevens. An ``extra
product'' record is a 5-byte $(\text{coeff}, \text{lhs}, \text{rhs})$ quintuple
contributing $K[c] \cdot M[u] \cdot M[v]$; both pointers are guarded (lines
2426, 2427).

Finally, if flag bit 0 is set, $A \gets M[p_{\mathrm{cond}}] \cdot A \bmod r$,
and the instruction pushes: $T \gets A$, $h \gets 1$. Applying the condition
last, once, rather than distributing it over every subterm, is both cheaper
(one \yul{mulmod} instead of one per term) and structurally safer --- it is
impossible for a subterm to escape the gate.

\subsubsection{Full semantics}

Collecting all five block families, one \opcode{MODARITH7} instruction computes
\begin{equation}
\begin{aligned}
A \;=\; \Bigl(M[p_{\mathrm{cond}}]\Bigr)^{[f_0]} \cdot \Biggl(
  & \;[f_1]\,K[s]
  \;+\; \sum_{b=0}^{n_L-1} \sum_{i=0}^{6} K\bigl[c^{L}_{b,i}\bigr]\, M\bigl[p^{L}_{b,i}\bigr] \\
  + & \; \sum_{b=0}^{n_R-1} \sum_{i=0}^{6} M\bigl[\ell_b\bigr]\, M\bigl[q^{R}_{b,i}\bigr]\, K\bigl[c^{R}_{b,i}\bigr] \\
  + & \; \sum_{b=0}^{n_P-1} \sum_{i=0}^{6} \sum_{j=0}^{6} K\bigl[d_{b,i+j}\bigr]\, M\bigl[\alpha_b + 32i\bigr]\, M\bigl[\beta_b + 32j\bigr] \\
  + & \; \sum_{k=0}^{n_M-1} K\bigl[c^{M}_{k}\bigr]\, M\bigl[p^{M}_{k}\bigr]
  \;+\; \sum_{k=0}^{n_X-1} K\bigl[c^{X}_{k}\bigr]\, M\bigl[u_k\bigr]\, M\bigl[v_k\bigr]
\Biggr) \bmod r ,
\end{aligned}
\label{eq:qvma:modarith7}
\end{equation}
where $(\cdot)^{[f_0]}$ means the factor is present only when flag bit 0 is
set, $[f_1] \in \{0,1\}$, and every $+$ and $\cdot$ is \yul{addmod}/
\yul{mulmod} modulo $r$. The accumulation order in the fixture is exactly the
order of Equation~\ref{eq:qvma:modarith7} read left to right, top to bottom;
since $(\Fr, +, \cdot)$ is a commutative ring the order does not affect the
result, but it does fix which guard fires first on malformed input.

\begin{algorithm}[htbp]
\caption{\opcode{MODARITH7} (\texttt{0x21}), \texttt{Halo2Verifier.sol} lines 2281--2447}
\label{alg:qvma:modarith7}
\begin{algorithmic}[1]
\State $f \gets \mathrm{byte}_0(M[\mathit{pc}])$; $\mathit{pc} \gets \mathit{pc} + 1$
\State $p_{\mathrm{cond}} \gets 0$
\If{$f \wedge 1$}
  \State $p_{\mathrm{cond}} \gets \mathrm{u16}(M[\mathit{pc}])$; $\mathit{pc} \gets \mathit{pc} + 2$
  \State \textbf{require} $p_{\mathrm{cond}} - \texttt{0x3680} \le \texttt{0x6aa0}$ \Comment{else \yul{q\_program\_fail}}
\EndIf
\State $A \gets 0$
\If{$f \wedge 2$}
  \State $s \gets \mathrm{byte}_0(M[\mathit{pc}])$; $\mathit{pc} \gets \mathit{pc} + 1$; $A \gets K[s]$
\EndIf
\State $(n_L, n_R, n_P, n_M, n_X) \gets (\mathrm{byte}_0, \dots, \mathrm{byte}_4)(M[\mathit{pc}])$; $\mathit{pc} \gets \mathit{pc} + 5$
\If{$h \ne 0$}
  \State \textbf{require} $\sigma < \texttt{0xbfa0}$; $M[\sigma] \gets T$; $\sigma \gets \sigma + 32$
\EndIf
\For{$b = 0$ \textbf{to} $n_L - 1$}
  \For{$i = 0$ \textbf{to} $6$}
    \State decode $(c, p)$ from $M[\mathit{pc}]$; $\mathit{pc} \gets \mathit{pc} + 3$; \textbf{require} $p - \texttt{0x3680} \le \texttt{0x6aa0}$
    \State $A \gets A + K[c] \cdot M[p]$
  \EndFor
\EndFor
\For{$b = 0$ \textbf{to} $n_R - 1$}
  \State $\ell \gets \mathrm{u16}(M[\mathit{pc}])$; $\mathit{pc} \gets \mathit{pc} + 2$; \textbf{require} $\ell - \texttt{0x3680} \le \texttt{0x6aa0}$; $L \gets M[\ell]$
  \For{$i = 0$ \textbf{to} $6$}
    \State decode $(c, q)$ from $M[\mathit{pc}]$; $\mathit{pc} \gets \mathit{pc} + 3$; \textbf{require} $q - \texttt{0x3680} \le \texttt{0x6aa0}$
    \State $A \gets A + L \cdot M[q] \cdot K[c]$
  \EndFor
\EndFor
\For{$b = 0$ \textbf{to} $n_P - 1$}
  \State $(\alpha, \beta) \gets \mathrm{u16}, \mathrm{u16}$ of $M[\mathit{pc}]$; $\mathit{pc} \gets \mathit{pc} + 4$
  \State \textbf{require} $\alpha - \texttt{0x3680} \le \texttt{0x69e0}$ \textbf{and} $\beta - \texttt{0x3680} \le \texttt{0x69e0}$
  \State $\mathit{cpc} \gets \mathit{pc}$; $\mathit{pc} \gets \mathit{pc} + 13$
  \For{$i = 0$ \textbf{to} $6$; $j = 0$ \textbf{to} $6$}
    \State $A \gets A + K[\mathrm{byte}_0(M[\mathit{cpc} + i + j])] \cdot M[\alpha + 32i] \cdot M[\beta + 32j]$
  \EndFor
\EndFor
\For{$k = 0$ \textbf{to} $n_M - 1$}
  \State decode $(c, p)$; $\mathit{pc} \gets \mathit{pc} + 3$; \textbf{require} $p - \texttt{0x3680} \le \texttt{0x6aa0}$; $A \gets A + K[c] \cdot M[p]$
\EndFor
\For{$k = 0$ \textbf{to} $n_X - 1$}
  \State decode $(c, u, v)$; $\mathit{pc} \gets \mathit{pc} + 5$
  \State \textbf{require} $u - \texttt{0x3680} \le \texttt{0x6aa0}$ \textbf{and} $v - \texttt{0x3680} \le \texttt{0x6aa0}$
  \State $A \gets A + K[c] \cdot M[u] \cdot M[v]$
\EndFor
\If{$f \wedge 1$} \State $A \gets M[p_{\mathrm{cond}}] \cdot A$ \EndIf
\State $T \gets A$; $h \gets 1$
\end{algorithmic}
\end{algorithm}

\subsection{Complete guard inventory for lines 2200--2452}
\label{sec:qvma:guards}

Table~\ref{tab:qvma:guards} lists every \yul{q\_program\_fail()} call site in the
assigned range. There are exactly fifteen (counted by grepping
\yul{q\_program\_fail()} over lines 2200--2452), in three classes:
operand-window checks (twelve), stack-overflow checks (two), and a
stack-underflow check (one). The table adds one row for line 2464, which lies
just \emph{outside} the assigned range --- it is the \opcode{NATIVE\_PERMUTATION}
arm's empty-stack guard --- because the stack-discipline argument below refers
to it; it is covered in the section on that opcode.

\begin{table}[htbp]
\centering
\caption{Every \yul{q\_program\_fail()} guard in \texttt{Halo2Verifier.sol} lines
2200--2452 (the first fifteen rows), plus the adjacent line-2464 guard
referenced by the stack-discipline argument.}
\label{tab:qvma:guards}
\begin{tabular}{@{}lllll@{}}
\toprule
Line & Opcode & Guarded quantity & Condition that fails & Class \\
\midrule
2216 & \texttt{0x05} & \yul{q\_ptr}       & \yul{gt(sub(q\_ptr, 0x3680), 0x6aa0)}       & window \\
2218 & \texttt{0x05} & \yul{q\_sp}        & \yul{iszero(lt(q\_sp, 0xbfa0))}             & overflow \\
2229 & \texttt{0x06} & \yul{q\_sp}        & \yul{eq(q\_sp, 0xb8e0)}                     & underflow \\
2253 & \texttt{0x10} & \yul{q\_ptr}       & \yul{gt(sub(q\_ptr, 0x3680), 0x6aa0)}       & window \\
2261 & \texttt{0x11} & \yul{q\_ptr}       & \yul{gt(sub(q\_ptr, 0x3680), 0x6aa0)}       & window \\
2308 & \texttt{0x21} & \yul{q\_cond\_ptr}  & \yul{gt(sub(q\_cond\_ptr, 0x3680), 0x6aa0)}  & window \\
2332 & \texttt{0x21} & \yul{q\_sp}        & \yul{iszero(lt(q\_sp, 0xbfa0))}             & overflow \\
2344 & \texttt{0x21} & \texttt{LIN7} ptr & \yul{gt(sub(q\_ptr, 0x3680), 0x6aa0)}       & window \\
2357 & \texttt{0x21} & row \yul{q\_lhs}   & \yul{gt(sub(q\_lhs, 0x3680), 0x6aa0)}       & window \\
2364 & \texttt{0x21} & row \yul{q\_rhs}   & \yul{gt(sub(q\_rhs, 0x3680), 0x6aa0)}       & window \\
2384 & \texttt{0x21} & \yul{q\_lhs\_base}  & \yul{gt(sub(q\_lhs\_base, 0x3680), 0x69e0)}  & window (vec7) \\
2385 & \texttt{0x21} & \yul{q\_rhs\_base}  & \yul{gt(sub(q\_rhs\_base, 0x3680), 0x69e0)}  & window (vec7) \\
2411 & \texttt{0x21} & extra mem ptr     & \yul{gt(sub(q\_ptr, 0x3680), 0x6aa0)}       & window \\
2426 & \texttt{0x21} & product \yul{q\_lhs} & \yul{gt(sub(q\_lhs, 0x3680), 0x6aa0)}     & window \\
2427 & \texttt{0x21} & product \yul{q\_rhs} & \yul{gt(sub(q\_rhs, 0x3680), 0x6aa0)}     & window \\
2464 & \texttt{0x19} & \yul{q\_sp}        & \yul{iszero(eq(q\_sp, 0xb8e0))}             & stack-empty \\
\bottomrule
\end{tabular}
\end{table}

\paragraph{Stack bounds are exactly tight.}
The stack region is $[\texttt{0xb8e0}, \texttt{0xbfa0})$, i.e.\ $\texttt{0x6c0}$
bytes $=54$ words. The overflow guard \yul{iszero(lt(q\_sp, 0xbfa0))} permits a
store at $\sigma = \texttt{0xbf80}$ (writing the words
$[\texttt{0xbf80}, \texttt{0xbfa0})$) and rejects the next one, so the VM can
never write outside its reserved region and can never write one word short of
it either. Maximum VM depth is 55 values (54 spilled plus the cached top). The
underflow guard \yul{eq(q\_sp, 0xb8e0)} is the exact complement. Every push site
in the range (lines 2217--2221 and 2331--2335) carries the overflow guard, and
the only pop site (line 2229--2230) carries the underflow guard, so the stack
discipline in this range is complete.

Both sentinels are planner constants, not interpreter constants. The base
$\texttt{0xb8e0}$ is \yul{program.stack\_mptr}; the ceiling $\texttt{0xbfa0}$ is
\yul{program.stack\_hi}, which
\texttt{src/lowering/layout/memory.rs} line 1159 defines as
$\yul{quotient\_stack\_mptr} + \max(\yul{quotient\_stack\_len},
\texttt{MODEXP\_FRAME\_BYTES})$, and
\texttt{src/lowering/quotient.rs} lines 202--220 set
$\yul{quotient\_stack\_words} = \max(\text{max interpreted stack},
\text{native-callback scratch words}, 1)$. The 54 words are therefore sized by
the \emph{native callback scratch}, not by the interpreted stack: the
permutation callback's scratch table saturates the region exactly (lines
2470--2475 place \yul{q\_perm\_vals} at \texttt{0xb8e0} and
\yul{q\_perm\_delta\_base\_ptr} at \texttt{0xbf80}, whose single word ends at
\texttt{0xbfa0}). The interpreted VM's own maximum depth in this artifact is
one word.

\paragraph{The window is a range check, not an alignment check.}
None of the fifteen guards constrains the low five bits of a pointer. A u16
operand may name any byte address in $[\texttt{0x3680}, \texttt{0xa120}]$, and
an unaligned one makes \yul{mload} return a word straddling two planned slots
--- a value that is arbitrary, is generally $\ge r$, and (through
\opcode{NEG}, Finding~\ref{find:qvma:unreduced}) need not even behave like a
field element. This is contained rather than prevented: the straddling read
still lies wholly inside $[\texttt{0x3680}, \texttt{0xa140})$, so an unaligned
pointer buys an attacker nothing the aligned window does not already grant.
Alignment is a build-time property, established by
\yul{validate\_quotient\_mem\_ptrs} (cited in
\texttt{src/lowering/quotient\_numerator/vm/yul\_arms.rs} line 18), not a
runtime one. Static decoding confirms all 41 distinct pointer operands in the
pinned program are 32-byte aligned.

\paragraph{The operand window.}
$[\texttt{0x3680}, \texttt{0xa140})$ is the union of the planned quotient read
windows, rendered as a single unsigned comparison. Its low bound is
\mptr{VK\_MPTR} and its high bound is exactly
$\mptr{REVERSED\_EVALS\_MPTR} + 32 \cdot 102 = \texttt{0x9480} + \texttt{0xcc0}
= \texttt{0xa140} = \mptr{ADVICE\_COMMS\_MPTR\_BASE}$ (lines 226, 257). The
window therefore admits the VK payload (header, constant table, and the VM's
own bytecode), the challenge block at \texttt{0x7900}, the Lagrange /
linearization scratch, the PCS staging areas, and the decoded evaluation
vector $\mathbf{e}$ --- and excludes:
\begin{itemize}
  \item all memory below \mptr{VK\_MPTR}, including the EVM scratch words at
        \texttt{0x00}, the free-memory pointer, and the transcript buffer at
        \mptr{TRANSCRIPT\_MPTR} $=\texttt{0x1000}$;
  \item the G1 commitment tables at \texttt{0xa140} and above;
  \item \emph{every accumulator the VM itself writes}: the simple-selector
        buckets at \mptr{SELECTOR\_ACC\_MPTR} $=\texttt{0xb140}$ (line 229),
        the fully-evaluated numerator accumulator at \texttt{0xb280} ---
        which is \mptr{PCS\_FINAL\_MSM\_MPTR} (line 204), reused by the VM
        after being zeroed at line 2021 and folded by $y$ at sites such as 2093,
        3043 and 3078 ---
        the $y$-power table at \texttt{0xb2c0} (lines 2042, 2048), and the VM
        stack at \texttt{0xb8e0}.
\end{itemize}
That last exclusion is the load-bearing one: no operand pointer, however
malformed, can make the VM read back a value it is in the middle of
accumulating, so the interpreter has no feedback path from its own outputs into
its own inputs.

\subsection{Bounds-check assessment: are all operand pointers checked?}
\label{sec:qvma:bounds}

Enumerating every address-forming expression in lines 2200--2452 that feeds an
\yul{mload} or \yul{mstore}:

\begin{enumerate}
  \item \textbf{Data pointers} --- \yul{q\_ptr} (2214, 2251, 2259, 2342, 2409),
        \yul{q\_cond\_ptr} (2306), \yul{q\_lhs}/\yul{q\_rhs} (2355, 2362, 2423,
        2424), \yul{q\_lhs\_base}/\yul{q\_rhs\_base} (2381, 2382). \textbf{All are
        window-checked before the corresponding \yul{mload}}, with the vector
        forms using the correctly tightened \texttt{0x69e0} span. No exceptions.
  \item \textbf{Stack addresses} --- \yul{mstore(q\_sp, ...)} and
        \yul{mload(q\_sp)}. Both bounded, see \S\ref{sec:qvma:guards}.
  \item \textbf{Constant-table addresses} ---
        \yul{add(q\_const\_mptr, shl(5, qconst))} at 2245, 2317, 2347, 2369,
        2396, 2414, 2432. \textbf{Not bounds-checked}
        (Finding~\ref{find:qvma:constidx}).
  \item \textbf{The program cursor} --- \yul{mload(q\_pc)} and
        \yul{mload(add(q\_coeff\_pc, add(q\_i, q\_j)))}. \textbf{Not bounds-checked
        within an instruction} (Finding~\ref{find:qvma:overread}).
\end{enumerate}

So the answer is: every pointer that names a \emph{field element operand} is
range-checked before it is dereferenced --- range-checked, not alignment-checked
(Finding~\ref{find:qvma:align}). The two unchecked address
computations are the constant-table index and the instruction cursor, and
neither can produce a memory \emph{write} --- the only \yul{mstore} in the
range is \yul{mstore(q\_sp, q\_top)}, whose address is separately bounded. The
findings below analyse each in turn.

\begin{finding}[QVMA-1]{Constant-table indexes are decoded from program bytes without a range check}
\label{find:qvma:constidx}
\textbf{Severity: Informational.}
\textbf{Where:} \texttt{Halo2Verifier.sol}:2245, 2317, 2347, 2369, 2396, 2414,
2432.

Every constant-table access computes
\yul{mload(add(q\_const\_mptr, shl(5, qconst)))} where \yul{qconst} is a raw
\texttt{u8} taken from the program with no comparison against the rendered
table length. The table holds 178 words at
$[\texttt{0x3a60}, \texttt{0x50a0})$, so indexes 178--255 read words the
generator did not intend as constants.

\emph{Why this is nevertheless sound in this artifact.} The \texttt{u8} width
caps the reachable offset at $255 \cdot 32 = \texttt{0x1fe0}$, so the load
address lies in $[\texttt{0x3a60}, \texttt{0x5a40}]$ and the loaded word ends
at \texttt{0x5a60} --- entirely inside the VK payload
$[\texttt{0x3680}, \texttt{0x7900})$. Because the constant table ends at
\texttt{0x50a0} and the bytecode begins there, \emph{every} out-of-range index
$c \in [178, 255]$ lands in $[\texttt{0x50a0}, \texttt{0x5a40}]$, i.e.\ inside
the VM's own program bytes: an out-of-range index reads the VM's own bytecode
as a field element, but can never leave the pinned payload, never touch proof
calldata, never touch an accumulator, and never write. Decoding the fixture's bytecode confirms the used index set is exactly
$\{0, \dots, 177\}$: all 178 table entries are referenced and none is
exceeded. Because the program is re-pinned by \yul{extcodehash} on every proof
(lines 1386--1389), an attacker cannot alter an index; only a generator bug
could, and the generator's build-time
\yul{validate\_quotient\_const\_slots} check covers exactly this. The generator
documents the omission as deliberate
(\texttt{src/lowering/quotient\_numerator/vm/yul\_arms.rs} lines 25--29) and
notes that the wider \texttt{u16} constant forms --- which are \emph{not}
emitted in this artifact --- \emph{are} clamped, because for those the same
containment argument fails. That asymmetry is correctly reasoned.

\emph{Cost/benefit.} Omitting the check saves roughly 6 gas per constant
access, and static decoding puts the exact count at 1187 constant accesses per
executed program --- so about 7\,000 gas per proof. (For scale, the 1552
pointer operands that \emph{are} guarded cost about 9\,300 gas per proof at the
same 6 gas each, which is what makes the asymmetry a deliberate trade rather
than an oversight.) The residual risk is a generator bug producing a silently
wrong numerator rather than a revert.
\end{finding}

\begin{finding}[QVMA-2]{Operand over-read past \texttt{q\_end} is detected only after the instruction has executed}
\label{find:qvma:overread}
\textbf{Severity: Informational.}
\textbf{Where:} \texttt{Halo2Verifier.sol}:2200 (loop condition) and every
\yul{mload(q\_pc)} in the range; reconciled at 3098.

The loop tests \yul{lt(q\_pc, q\_end)} only at instruction boundaries. A
malformed trailing instruction --- in the worst case a \opcode{MODARITH7} whose
five counters are all \texttt{0xff}, consuming 17{,}605 bytes per
Equation~\ref{eq:qvma:width} --- would decode operands from memory beyond
\yul{q\_end} $=\texttt{0x626f}$, i.e.\ from the VK payload's fixed and
permutation commitments and then from zero-initialised memory.

\emph{Why this fails closed.} Every operand pointer so decoded is still
window-checked, and a zero pointer (the value read from untouched memory) fails
the window check immediately. If the over-read somehow produced only in-window
pointers, the arithmetic would still be discarded, because the epilogue at line
3098 requires $\yul{q\_pc} = \yul{q\_end}$ \emph{exactly} --- a strict equality,
not $\ge$ --- and calls \yul{q\_program\_fail()} otherwise. The transaction
therefore reverts in all cases; the only observable effect is wasted gas in a
call that was going to revert anyway. Adding a per-operand bound would cost gas
on every instruction to convert one revert into a slightly earlier revert.
Decoding the fixture's program confirms it terminates at exactly
$\yul{q\_pc} - \yul{q\_program\_mptr} = \texttt{0x11cf}$ with $h = 0$ and
$\sigma = \texttt{0xb8e0}$, so all three epilogue invariants hold with no
slack.

\emph{Gas, not availability.} The five counters are unbounded \texttt{u8}s, so
a single malformed \opcode{MODARITH7} could drive $255 \times 49 = 12\,495$
weighted pairwise terms ($\approx 25\,000$ \yul{mulmod}s, on the order of
$10^6$ gas). That is a cost paid by the caller of a transaction that reverts
regardless, not a liveness problem for other users: there is no unbounded
\emph{loop over attacker input}, because the counters live in the
codehash-pinned VK and every pointer the loop touches is still window-checked
before it is read.
\end{finding}

\begin{finding}[QVMA-3]{\opcode{ADD} neither reads nor sets \yul{q\_has\_top}}
\label{find:qvma:addtop}
\textbf{Severity: Low.}
\textbf{Where:} \texttt{Halo2Verifier.sol}:2226--2232.

The \opcode{ADD} arm guards only the spilled side ($\sigma \ne \texttt{0xb8e0}$)
and then computes $T \gets M[\sigma] + T$ without testing $h$, and without
setting $h \gets 1$ afterwards. The sibling fold arms do carry the
corresponding guard --- \opcode{FOLD\_SELECTOR} rejects $h = 0$ at line 3068.
The generator's ``top guard'' is documented as a fail-closed containment
measure (\texttt{src/lowering/quotient\_numerator/vm/yul\_arms.rs} lines
31--41), but note that its stated scope is \emph{fold} opcodes
(``FOLD\_MAIN/FOLD\_SELECTOR consume the cached top''), not \opcode{ADD}; the
generator does not claim an $h$ guard on \opcode{ADD} and does not emit one
(\yul{pop\_guard}, \texttt{yul\_arms.rs} lines 209--215, emits only the
$\sigma$ test).

\emph{Reachable state.} $h = 0$ with a non-empty spill stack is constructible:
\opcode{FOLD\_SELECTOR} clears $h$ (line 3070) without requiring
$\sigma = \texttt{0xb8e0}$, unlike the native-callback opcodes which do require
it (lines 2464, 2580, 2772). A program of the form
\texttt{PUSH; PUSH; FOLD\_SELECTOR; ADD} would therefore pop the abandoned word
and add it to a \emph{stale} $T$ (\opcode{FOLD\_SELECTOR} does not clear
\yul{q\_top}, only $h$), leaving $h = 0$ while $T$ holds a value, and $\sigma$
back at the base. Both the $h = 0$ and $\sigma = \texttt{0xb8e0}$ epilogue
checks (lines 3099, 3105) then pass.

Two points sharpen this. First, the corrupted $T$ itself cannot reach
$\nu_y(x)$ by this route: with $h = 0$ the only opcode that consumes $T$
(\opcode{FOLD\_SELECTOR}) reverts at 3068, and the only opcodes that set
$h \gets 1$ (\opcode{PUSH\_MEM\_U16}, \opcode{MODARITH7}) overwrite $T$ first.
The real damage is that \opcode{ADD} \emph{launders} the stack imbalance:
lines 3100--3104 state in so many words that check 3105 exists to catch ``a
FOLD executed with more than one operand live \dots\ while an operand of the
identity has been silently dropped''. An unguarded \opcode{ADD} drains the
abandoned word and defeats exactly that detector, so the dropped operand
--- a genuinely wrong $\nu_y(x)$ --- passes unnoticed.

Second, and worse, an $h$ guard on \opcode{ADD} would not close the general
case. Consider \texttt{PUSH a; PUSH b; FOLD\_SELECTOR; PUSH c; ADD;
FOLD\_SELECTOR}. Here the second \opcode{PUSH} sees $h = 0$ and does not spill,
so \opcode{ADD} runs with $h = 1$ (legal under any $h$ test) and folds
$a + c$: the abandoned operand $a$ of the \emph{first} identity is injected
into a \emph{later} one, and all three epilogue checks pass. The root gap is
that \opcode{FOLD\_SELECTOR} clears $h$ without requiring
$\sigma = \texttt{0xb8e0}$ (lines 3054--3088, covered elsewhere), unlike the
native-callback opcodes at 2464, 2580 and 2772 which do require it.

\emph{Exploitability: none on-chain.} The opcode stream is not attacker
controlled --- it is part of the VK runtime whose \yul{extcodehash} is
re-verified on every call (line 1388) --- so reaching either state requires a
generator bug, not a crafted proof. Decoding the fixture's 4559-byte program
confirms it never occurs: the maximum spill depth is one word, every
\opcode{FOLD\_SELECTOR} executes with $\sigma = \texttt{0xb8e0}$, and every
\opcode{ADD} executes with $h = 1$. This is therefore a gap in
defence-in-depth, not a live vulnerability. A one-line
\yul{if iszero(q\_has\_top) \{ q\_program\_fail() \}} in the \opcode{ADD} arm
closes the first variant for roughly 3 gas per \opcode{ADD}; only a
$\sigma = \texttt{0xb8e0}$ requirement on \opcode{FOLD\_SELECTOR} closes both.

The same absence of an $h$ test applies to \opcode{NEG} (2234--2238) and
\opcode{MUL\_CONST\_U8} (2240--2246); those are strictly less dangerous, since
neither consumes a spilled word --- with $h = 0$ and a stale $T$ they merely
transform a value that a subsequent \opcode{PUSH} would overwrite.
\end{finding}

\begin{finding}[QVMA-4]{Loaded words are not reduced modulo $r$, and \opcode{NEG} is only correct for inputs $\le r$}
\label{find:qvma:unreduced}
\textbf{Severity: Informational.}
\textbf{Where:} \texttt{Halo2Verifier.sol}:2222 (\yul{q\_top := mload(q\_ptr)})
and 2237 (\yul{addmod(0, sub(r, q\_top), r)}).

\opcode{PUSH\_MEM\_U16} stores the raw memory word into $T$ with no modular
reduction. Every other arithmetic opcode is indifferent to this ---
\yul{addmod} and \yul{mulmod} reduce their operands --- but \opcode{NEG}
computes $\yul{sub}(r, T)$ in $\mathbb{Z}_{2^{256}}$ before reducing. For
$T \le r$ this yields the canonical inverse; for $T > r$ it yields
$(2^{256} - T) \bmod r$, which is not $-T \bmod r$.

The operand window admits words that are not canonical scalars: for example
the padded G1 base-point low word at \texttt{0x3800}
(\mptr{G1\_BASE\_MPTR} $+ \texttt{0x20}$, holding the low 32 bytes of a
48-byte $\Fp$ coordinate) and the quotient-commitment coordinates at
\mptr{QUOTIENT\_MPTR} $=\texttt{0x7d20}$ all lie inside
$[\texttt{0x3680}, \texttt{0xa120}]$ and can exceed $r$. So the window alone
does not establish the precondition.

\emph{Why it holds here.} Decoding the fixture's program shows all 41 distinct
pointer operands lie in $[\texttt{0x94a0}, \texttt{0x9aa0}]$ and are 32-byte
aligned, entirely inside \mptr{REVERSED\_EVALS\_MPTR} $=\texttt{0x9480}$. That
buffer is written only by the transcript's evaluation loop, lines
1697--1715, which does
\yul{if iszero(lt(eval, r)) \{ fail(ERR\_NON\_CANONICAL\_SCALAR) \}} at line
1706 \emph{before} the \yul{mstore(eval\_buf, eval)} at line 1708; the loop
runs over \texttt{0x0cc0} calldata bytes $=102$ words, so the range-checked
buffer is exactly $[\texttt{0x9480}, \texttt{0xa140})$ --- which is also
precisely where the operand window ends. Every VM operand in
this artifact is therefore a canonical $\Fr$ element in $[0, r)$ and
\opcode{NEG} is exact. This is a precondition that the runtime does not
re-establish, only the generator's window plan does; it should be recorded as
an invariant of the artifact rather than assumed from the guard.
\end{finding}

\begin{finding}[QVMA-5]{\texttt{BILIN7\_PAIRWISE} is unreachable in the deployed program}
\label{find:qvma:deadpair}
\textbf{Severity: Observation.}
\textbf{Where:} \texttt{Halo2Verifier.sol}:2377--2403.

Decoding the pinned bytecode shows 16 \opcode{MODARITH7} instructions, and
\yul{q\_pairwise\_count} is $0$ in every one of them (the observed count vectors
are given in Table~\ref{tab:qvma:ma-census}). The entire pairwise loop body ---
including its distinct \texttt{0x69e0} vector guard, the 13-byte convolution
coefficient table, and the 49-term double loop --- is therefore dead code in
this artifact: it occupies runtime bytecode in a contract whose own header
records that the two nearby \texttt{(solc version, optimize-runs)} settings it
measured emit 29\,567 and 29\,836 bytes, both \emph{over} the 24\,576-byte
EIP-170 limit, so that only the pinned pair is deployable at all (lines
8--11). It is never executed by any proof this VK can verify.

This is not a defect --- the interpreter is rendered from the opcode table with
usage gating (\texttt{yul\_arms.rs}), and the emission of a
\opcode{MODARITH7} arm necessarily brings all five of its sub-blocks --- but a
reviewer should know that the pairwise path receives no on-chain coverage from
this deployment and its correctness rests entirely on the generator's own
tests. Note in particular that its guard is the only one in the range that uses
a different span constant, and that span is exactly right
($\texttt{0x6aa0}-\texttt{0xc0}$, matching \yul{vec7\_guard} in
\texttt{yul\_arms.rs}); the arithmetic of the tightening was verified by hand
above and does not depend on the path being exercised.
\end{finding}

\begin{finding}[QVMA-6]{Undefined \opcode{MODARITH7} flag bits are ignored rather than rejected}
\label{find:qvma:flagbits}
\textbf{Severity: Informational.}
\textbf{Where:} \texttt{Halo2Verifier.sol}:2300--2303, 2312, 2439.

The flags byte is tested only with \yul{and(q\_flags, 0x01)} and
\yul{and(q\_flags, 0x02)}; bits 2--7 are neither used nor checked, so
\texttt{0x03} and \texttt{0xff} decode identically. By contrast the opcode
dispatch itself fails closed --- the \yul{default} arm at line 3090 calls
\yul{q\_program\_fail()}, and line 3089 records that the reserved opcode
\texttt{0x1a} lands there deliberately. Applying the same fail-closed
discipline to the flags byte (\yul{if gt(q\_flags, 0x03) \{ q\_program\_fail() \}},
about 6 gas per \opcode{MODARITH7}, so under 100 gas per proof here) would make
the encoding self-validating in the same way the opcode space already is. The
observed flag values in the pinned program are only \texttt{0x02} (4
instructions) and \texttt{0x03} (12 instructions).
\end{finding}

\begin{finding}[QVMA-7]{The operand window is a range check only; pointer alignment is never enforced at runtime}
\label{find:qvma:align}
\textbf{Severity: Observation.}
\textbf{Where:} \texttt{Halo2Verifier.sol}:2216, 2253, 2261, 2308, 2344, 2357,
2364, 2384, 2385, 2411, 2426, 2427 (all twelve window guards).

Each guard is a single \yul{gt(sub(ptr, 0x3680), span)} test. It constrains no
low-order bit, so a u16 operand may name any \emph{byte} address in
$[\texttt{0x3680}, \texttt{0xa120}]$, and the subsequent \yul{mload} may return
a word straddling two planned 32-byte slots. Such a word is arbitrary, is
usually $\ge r$, and therefore also triggers the \opcode{NEG} precondition of
Finding~\ref{find:qvma:unreduced}. The same is true of the two
\texttt{BILIN7\_PAIRWISE} base pointers, where a misaligned $\alpha$ makes all
seven limb loads straddle.

This is contained, not prevented, and deliberately so: the straddling read
still lies wholly inside $[\texttt{0x3680}, \texttt{0xa140})$, so alignment
buys nothing an aligned in-window pointer does not already grant, and adding
an \yul{and(ptr, 0x1f)} test would cost about 6 gas on each of the 1552
pointer operands this program decodes ($\approx 9\,300$ gas per proof) to
convert one class of nonsense into a revert. Alignment is instead a build-time
property proved by \yul{validate\_quotient\_mem\_ptrs}
(\texttt{src/lowering/quotient\_numerator/vm/yul\_arms.rs} line 18). Static
decoding confirms all 41 distinct pointer operands in the pinned program are
32-byte aligned. Recorded here only because the guard's single-comparison form
invites the reader to assume it validates a \emph{slot} index when it
validates a byte address.
\end{finding}

\subsection{What the pinned program actually contains}
\label{sec:qvma:census}

Because the bytecode is a fixed byte string in the VK runtime, the instruction
stream can be decoded statically. Doing so is the strongest available check
that the specification above is faithful: an independent decoder driven by the
widths in Table~\ref{tab:qvma:opcodes} and Equation~\ref{eq:qvma:width}
consumes the 4559-byte program and lands on exactly \yul{q\_end}, with the
terminal state $h = 0$, $\sigma = \texttt{0xb8e0}$ required by lines
3098--3105. Any error in the stated operand widths would have desynchronised
the decode long before the end.

\begin{table}[htbp]
\centering
\caption{Opcode census of the pinned 4559-byte quotient program.}
\label{tab:qvma:census}
\begin{tabular}{@{}llr@{}}
\toprule
Byte & Mnemonic & Occurrences \\
\midrule
\texttt{0x05} & \opcode{PUSH\_MEM\_U16}      & 12 \\
\texttt{0x06} & \opcode{ADD}                 & 7 \\
\texttt{0x08} & \opcode{NEG}                 & 5 \\
\texttt{0x0b} & \opcode{FOLD\_SELECTOR}      & 21 \\
\texttt{0x0d} & \opcode{MUL\_CONST\_U8}      & 20 \\
\texttt{0x10} & \opcode{ADD\_MEM\_U16}       & 2 \\
\texttt{0x11} & \opcode{MUL\_MEM\_U16}       & 12 \\
\texttt{0x19} & \opcode{NATIVE\_PERMUTATION} & 1 \\
\texttt{0x1b} & \opcode{NATIVE\_IDENTITY}    & 4 \\
\texttt{0x1f} & \opcode{NATIVE\_LOOKUP}      & 1 \\
\texttt{0x21} & \opcode{MODARITH7}           & 16 \\
\bottomrule
\end{tabular}
\end{table}

\begin{table}[htbp]
\centering
\caption{The 16 \opcode{MODARITH7} instructions of the pinned program: flags,
seed index, block counts, and encoded width (Eq.~\ref{eq:qvma:width}). Offsets
are relative to \yul{q\_program\_mptr}.}
\label{tab:qvma:ma-census}
\begin{tabular}{@{}lllrrrrrr@{}}
\toprule
Offset & Flags & Seed $s$ & $n_L$ & $n_R$ & $n_P$ & $n_M$ & $n_X$ & Width \\
\midrule
\texttt{0x0014} & \texttt{0x02} & 1   & 3 & 7  & 0 & 2  & 0  & 238 \\
\texttt{0x0106} & \texttt{0x02} & 23  & 3 & 7  & 0 & 1  & 0  & 235 \\
\texttt{0x01f8} & \texttt{0x02} & 39  & 1 & 0  & 0 & 9  & 0  & 56 \\
\texttt{0x0234} & \texttt{0x02} & 40  & 1 & 0  & 0 & 8  & 0  & 53 \\
\texttt{0x026d} & \texttt{0x03} & 41  & 0 & 0  & 0 & 11 & 43 & 258 \\
\texttt{0x0375} & \texttt{0x03} & 61  & 3 & 8  & 0 & 2  & 21 & 368 \\
\texttt{0x04eb} & \texttt{0x03} & 77  & 3 & 8  & 0 & 1  & 21 & 365 \\
\texttt{0x065e} & \texttt{0x03} & 93  & 0 & 0  & 0 & 12 & 57 & 331 \\
\texttt{0x07af} & \texttt{0x03} & 97  & 2 & 15 & 0 & 10 & 0  & 427 \\
\texttt{0x0960} & \texttt{0x03} & 98  & 2 & 15 & 0 & 9  & 0  & 424 \\
\texttt{0x0b0e} & \texttt{0x03} & 99  & 0 & 0  & 0 & 11 & 43 & 258 \\
\texttt{0x0c16} & \texttt{0x03} & 122 & 3 & 8  & 0 & 2  & 21 & 368 \\
\texttt{0x0d8c} & \texttt{0x03} & 147 & 3 & 8  & 0 & 1  & 21 & 365 \\
\texttt{0x0eff} & \texttt{0x03} & 172 & 0 & 0  & 0 & 14 & 16 & 132 \\
\texttt{0x0f89} & \texttt{0x03} & 173 & 4 & 1  & 0 & 2  & 21 & 228 \\
\texttt{0x1073} & \texttt{0x03} & 174 & 4 & 1  & 0 & 1  & 21 & 225 \\
\bottomrule
\end{tabular}
\end{table}

Salient facts, each of which bears on the assessment above:

\begin{itemize}
  \item \textbf{The stack is barely used.} Maximum spill depth is one word,
        against a 54-word reservation. The \opcode{PUSH\_MEM\_U16} spill path
        (lines 2217--2221) is taken 7 times per proof; the \opcode{MODARITH7}
        spill path (lines 2331--2335) is never taken. The overflow guards are
        thus pure containment, and the stack window is over-provisioned by
        roughly $54\times$ against the interpreted stack --- but the region is not
        sized for the interpreter at all. Native callbacks share the same base
        and saturate it exactly: the permutation callback places
        \yul{q\_perm\_vals} at \texttt{0xb8e0} and
        \yul{q\_perm\_delta\_base\_ptr} at \texttt{0xbf80} (lines 2470--2475),
        whose word ends at the ceiling \texttt{0xbfa0}; the lookup callback
        places \yul{q\_lookup\_prefix} at \texttt{0xb960} and
        \yul{q\_lookup\_suffix} at \texttt{0xb9e0} (lines 2584--2586). That
        reuse is safe precisely because every native callback
        first requires $\sigma = \texttt{0xb8e0}$ (lines 2464, 2580, 2772), so
        no live VM operand can be resident in the scratch range.
  \item \textbf{Every VM memory operand is a proof evaluation.} All 41 distinct
        pointer values fall in $[\texttt{0x94a0}, \texttt{0x9aa0}]$, i.e.\
        within evaluation slots $e_1 \dots e_{49}$ of the decoded eval vector
        at \mptr{REVERSED\_EVALS\_MPTR} (the run is $e_1 \dots e_{39}$ plus
        $e_{44}$ and $e_{49}$; the remaining slots of $\mathbf{e}$ are consumed
        by the PCS, not by the VM). The VM never reads a challenge, a
        Lagrange value, or a VK word through a pointer operand in this artifact
        --- those enter through the constant table or through the natively
        emitted prefix/suffix blocks.
  \item \textbf{The constant table is exactly saturated.} Indexes $0$ through
        $177$ all appear and none is exceeded, matching the 178 rendered
        \texttt{quotient\_const} words.
  \item \textbf{\opcode{MODARITH7} dominates.} 16 instructions occupy 4331 of
        4559 bytes. Its densest form is the extra-product block: instruction
        \texttt{0x065e} encodes 57 products in 285 bytes ($5$ bytes per
        product, each driving three \yul{mload}s --- left operand, right
        operand, coefficient --- and two \yul{mulmod}s plus one
        \yul{addmod}). Across the whole program that is 1187 constant-table
        loads and 1552 guarded pointer loads per proof.
\end{itemize}

\subsection{Summary assessment}
\label{sec:qvma:summary}

The arithmetic core of the quotient VM is clean in the sense that matters most:
\textbf{every operand that becomes a field element is bounds-checked against a
window that excludes the interpreter's own accumulators, and every stack
transition is bounds-checked against a region the VM cannot escape.} The
window is enforced with one wrapping unsigned comparison per operand, the
vector form is tightened by exactly the $\texttt{0xc0}$ bytes its seven-word
stride requires, and all fifteen guard sites in the range funnel into a single
typed revert. No opcode in this range performs an external call: there is no
\yul{staticcall}, \yul{call}, or precompile invocation anywhere in lines
2200--2452, so no gas bound or return-size check arises here.

The design's principal trade is that structural well-formedness of the
\emph{program} (correct stack discipline, in-range constant indexes, operands
that stay inside \yul{q\_end}) is enforced at generation time and at deployment
by the codehash pin, not at runtime. That is a defensible split --- the program
is not attacker-controlled, and re-validating it on every proof would cost gas
in a contract already at the EIP-170 limit --- but it means the runtime guards
are containment for generator bugs, and Finding~\ref{find:qvma:addtop}
identifies one containment gap where a generator bug could yield a silently
wrong numerator rather than a revert. No issue in this range is reachable by
any choice of proof or public input.


%% ==== 13-qvm-native.tex =================================================
\section{Quotient VM native-callback opcodes: permutation, lookup, heavy gates,
and selector folds}
\label{sec:qvmn}

This section specifies \texttt{Halo2Verifier.sol} lines 2452--3095: the four
remaining arms of the quotient-VM opcode dispatch
(\opcode{native\_permutation} \texttt{0x19}, \opcode{native\_lookup}
\texttt{0x1f}, \opcode{native\_identity} \texttt{0x1b},
\opcode{fold\_selector} \texttt{0x0b}) together with the \yul{default} arm that
closes the \yul{switch}. The stack/arithmetic opcodes that precede them
(\texttt{0x05}, \texttt{0x06}, \texttt{0x08}, \texttt{0x0d}, \texttt{0x10},
\texttt{0x11}, \texttt{0x21}) are covered elsewhere; where their behaviour is
load-bearing for the invariants below it is cited by line number.

\subsection{The runtime frame these arms operate on}
\label{sec:qvmn:frame}

Every arm in this range works over the same six memory regions. Every address
below is either a generated literal or a named pointer constant; none is
derived from proof calldata. Nothing in \texttt{Halo2Verifier.sol}
2452--3095 performs a \yul{call}, \yul{staticcall}, \yul{calldataload},
\yul{sload}/\yul{sstore}, \yul{keccak256} or \yul{mcopy}: the whole range is
pure \Fq{} arithmetic over memory, and its only exit is
\yul{q\_program\_fail}.

\begin{table}[h]
\centering
\small
\begin{tabularx}{\textwidth}{llX}
\toprule
Address & Size & Role \\
\midrule
\texttt{0xb140} (\mptr{SELECTOR\_ACC\_MPTR}) & $10 \times 32$\,B &
  Ten simple-selector linearization accumulators $B_0,\dots,B_9$. Zeroed by the
  VM prologue (line 2025--2028, loop bound \texttt{0x0140}); consumed as MSM
  scalars at lines 3932--3950. \\
\texttt{0xb280} & $32$\,B &
  The running Horner accumulator $A$ for fully-evaluated identities. Set to $0$
  at line 2021. Negated into \mptr{LINEARIZATION\_EVAL\_MPTR} at line 3358. The
  same address is \mptr{PCS\_FINAL\_MSM\_MPTR}, reused in a strictly later
  phase. \\
\texttt{0xb2c0} & $49 \times 32$\,B &
  Precomputed y-power table $Y[i] = y^i$ for $i = 0,\dots,48$, written by lines
  2042--2049. Word start addresses \texttt{0xb2c0} through \texttt{0xb8c0}. \\
\texttt{0xb8e0} & $54 \times 32$\,B &
  The registered VM stack window, base value of \yul{q\_sp}. The
  \texttt{0x05} and \texttt{0x21} spill guards (lines 2218 and 2332) refuse to
  push unless \yul{q\_sp} $<$ \texttt{0xbfa0}, so the window is
  exactly $[\mathtt{0xb8e0},\mathtt{0xbfa0})$. The native callbacks re-purpose
  this same window as a structured scratch table. \\
\texttt{0x9480} (\mptr{REVERSED\_EVALS\_MPTR}) & $102 \times 32$\,B &
  The decoded evaluation vector $\mathbf{e} = (e_0,\dots,e_{101})$, spilled and
  range-checked (\yul{iszero(lt(eval, r))} $\Rightarrow$
  \revert{NonCanonicalScalar}) at lines 1698--1714. \\
\texttt{0x7900}--\texttt{0x7ce0} & $8$ words &
  Named scalar words the arms load and never write:
  \mptr{THETA\_MPTR} (\texttt{0x7900}), \mptr{BETA\_MPTR} (\texttt{0x7920}),
  \mptr{GAMMA\_MPTR} (\texttt{0x7940}), \mptr{X\_MPTR} (\texttt{0x79a0}),
  \mptr{L\_LAST\_MPTR} (\texttt{0x7c80}), \mptr{L\_BLIND\_MPTR}
  (\texttt{0x7ca0}), \mptr{L\_0\_MPTR} (\texttt{0x7cc0}),
  \mptr{INSTANCE\_EVAL\_MPTR} (\texttt{0x7ce0}). \\
\bottomrule
\end{tabularx}
\caption{Memory regions read and written by the native-callback arms.}
\label{tab:qvmn:frame}
\end{table}

\paragraph{The fold snippet.}
Every identity in the whole quotient block --- inline prefix, interpreted,
native callback, and structured trash tail alike --- advances the global y-batch
position by exactly one, using one of two textually identical snippets.
A \emph{main} fold (used by \texttt{0x19} and \texttt{0x1f}) is
\begin{center}
\yul{mstore(0xb280, mulmod(mload(0xb280), y, r))} \quad followed by \quad
\yul{mstore(0xb280, addmod(mload(0xb280), <the identity value>, r))},
\end{center}
i.e.\ $A \leftarrow A\cdot y + \varepsilon$. A \emph{selector} fold (used by
\texttt{0x0b} and by the four \texttt{0x1b} sub-cases) performs the same
$A \leftarrow A\cdot y$ but then adds $\varepsilon$ into a selector bucket
instead of into $A$. Writing $N$ for the total identity count and indexing
identities $k = 0,\dots,N-1$ in generated order, this yields
\begin{equation}
A_{\text{final}} \;=\; \sum_{k \in \mathcal{M}} y^{\,N-1-k}\,\varepsilon_k,
\qquad
B_s \;=\; \sum_{k \in \mathcal{S}_s} y^{\,N-1-k}\,\varepsilon_k,
\label{eq:qvmn:reverse-fold}
\end{equation}
where $\mathcal{M}$ is the set of fully-evaluated (\texttt{None}-tagged)
identity positions and $\mathcal{S}_s$ the positions routed to simple selector
$s$. Equation~\eqref{eq:qvmn:reverse-fold} is the on-chain transcription of
midnight-proofs' reverse fold by powers of $y$ in
\texttt{linearization/verifier.rs}. The second half of
\eqref{eq:qvmn:reverse-fold} is realised in two pieces: the arms in this range
maintain $B_s$ only up to a trailing $y$ power, and the post-VM tail block
(lines 3314--3353, covered elsewhere) applies the missing factor
$y^{\,N-1-k_{\mathrm{last}}(s)}$.

\paragraph{Semantics of a selector fold.}
The value $\varepsilon_k$ folded into $B_s$ is the gate polynomial \emph{with
its simple selector removed}: every such expression is emitted as
\yul{let var0 := 0x1} $\ldots$ \yul{mulmod(var0, <body>, r)} in the inline and
native forms, and as a terminal \yul{MUL\_CONST\_U8} with constant slot $0$ in
the interpreted form. Constant slot $0$ of the VK quotient constant table
(memory \texttt{0x3a60}, VK payload offset \texttt{0x03e0}) is the literal
$1$, so both forms multiply by one. The selector itself is reinstated on the
commitment side: $B_s$ becomes the MSM scalar of the selector's fixed-column
commitment at lines 3932--3950. This is why a selector fold must \emph{not}
add into $A$ --- doing so would double-count the identity, once as a scalar and
once as a commitment.

\subsection{The identity budget of this artifact}
\label{sec:qvmn:budget}

The pinned bytecode is \texttt{0x11cf} bytes long (\yul{q\_end := add(q\_program\_mptr,
0x11cf)}, line 2170) and lives at memory \texttt{0x50a0}, immediately after the
178-entry constant table at \texttt{0x3a60}. Decoding it with the exact operand
widths implemented by the eleven rendered arms yields a program that consumes
every byte and terminates precisely at \yul{q\_end}:

\begin{table}[h]
\centering
\small
\begin{tabular}{llrr}
\toprule
Opcode & Mnemonic & Operand bytes & Occurrences \\
\midrule
\texttt{0x05} & \opcode{push\_mem\_u16}   & 2 & 12 \\
\texttt{0x06} & \opcode{add}              & 0 & 7 \\
\texttt{0x08} & \opcode{neg}              & 0 & 5 \\
\texttt{0x0b} & \opcode{fold\_selector}   & 3 & 21 \\
\texttt{0x0d} & \opcode{mul\_const\_u8}   & 1 & 20 \\
\texttt{0x10} & \opcode{add\_mem\_u16}    & 2 & 2 \\
\texttt{0x11} & \opcode{mul\_mem\_u16}    & 2 & 12 \\
\texttt{0x19} & \opcode{native\_permutation} & 0 & 1 \\
\texttt{0x1b} & \opcode{native\_identity} & 2 & 4 \\
\texttt{0x1f} & \opcode{native\_lookup}   & 0 & 1 \\
\texttt{0x21} & \opcode{modarith7}        & variable & 16 \\
\bottomrule
\end{tabular}
\caption{Static opcode census of the VK-pinned quotient program
(\texttt{Halo2VerifyingKey.sol} \texttt{quotient\_program} bytes, memory
$[\mathtt{0x50a0},\mathtt{0x626f})$, last byte \texttt{0x626e}). The census
matches the opcode summary comment at lines 2182--2192 exactly, and was
reproduced here by decoding the VK payload bytes directly.}
\label{tab:qvmn:census}
\end{table}

Two structural facts follow immediately and are used throughout the rest of this
section.

\begin{enumerate}
\item \textbf{This artifact renders no \opcode{fold\_main} arm.} Opcode
\texttt{0x0a} is \opcode{fold\_main} in the VM ABI, but the interpreter's
\yul{switch} has no \texttt{0x0a} case, so a \texttt{0x0a} byte would hit the
\yul{default} arm at line 3090 and revert. Consequently \emph{every}
fully-evaluated identity in this verifier is produced by non-interpreted code:
the inline prefix (lines 2055--2153), the two whole-family native callbacks, and
the structured trash tail (lines 3114--3305). The interpreted portion of the
program exists solely to evaluate simple-selector gate bodies and hand them to
\opcode{fold\_selector}.
\item \textbf{$N = 49$.} Counting dynamic folds: 4 from the inline prefix,
21 from \opcode{fold\_selector}, 4 from \opcode{native\_identity}, 13 from
\opcode{native\_permutation}, 6 from \opcode{native\_lookup}, and 1 from the
trash tail. The y-power table is sized to exactly $N$ entries (loop bound
\texttt{49} at line 2044), and the \opcode{fold\_selector} gap clamp is
\texttt{0x30} $= 48 = N-1$ (line 3067). Both bounds are tight, not
conservative.
\end{enumerate}

Table~\ref{tab:qvmn:stream} gives the complete identity stream. It is the
reference used below to check that native fold positions agree with interpreted
ones.

\begin{table}[h]
\centering
\small
\begin{tabularx}{\textwidth}{llX}
\toprule
Positions $k$ & Producer & Target \\
\midrule
$0,1,2,3$ & inline prefix, lines 2055--2153 & $B_0$; $B_1$; $B_2$; $B_2$ with $y^1$ \\
$4$   & \opcode{fold\_selector} @ \texttt{0x000d} & $B_2$, gap 1 \\
$5$   & \opcode{native\_identity} 0 @ \texttt{0x0011} & $B_3$, no gap \\
$6,7$ & \opcode{fold\_selector} @ \texttt{0x0102}, \texttt{0x01f1} & $B_3$, gap 1 \\
$8$   & \opcode{native\_identity} 1 @ \texttt{0x01f5} & $B_4$, no gap \\
$9,10$ & \opcode{fold\_selector} & $B_4$, gap 1 \\
$11,12,13$ & \opcode{fold\_selector} & $B_5$ (gaps $0,1,1$) \\
$14,15,16$ & \opcode{fold\_selector} & $B_6$ (gaps $0,1,1$) \\
$17,18,19$ & \opcode{fold\_selector} & $B_7$ (gaps $0,1,1$) \\
$20,21,22$ & \opcode{fold\_selector} & $B_8$ (gaps $0,1,1$) \\
$23,24,25$ & \opcode{fold\_selector} & $B_9$ (gaps $0,1,1$) \\
$26$  & \opcode{native\_identity} 2 @ \texttt{0x1196} & $B_9$, $y^1$ \\
$27$  & \opcode{native\_identity} 3 @ \texttt{0x1199} & $B_9$, $y^1$ \\
$28$  & \opcode{fold\_selector} @ \texttt{0x11c9} & $B_9$, gap 1 \\
$29$--$41$ & \opcode{native\_permutation} @ \texttt{0x11cd} & $A$ (13 folds) \\
$42$--$47$ & \opcode{native\_lookup} @ \texttt{0x11ce} & $A$ (6 folds) \\
$48$  & structured trash tail, lines 3114--3305 & $A$ (1 fold) \\
\bottomrule
\end{tabularx}
\caption{The complete 49-identity stream. Byte offsets are relative to
\yul{q\_program\_mptr} $=$ \texttt{0x50a0}.}
\label{tab:qvmn:stream}
\end{table}

The ordering in Table~\ref{tab:qvmn:stream} is exactly the order produced by
\texttt{proofs/src/plonk/mod.rs::partially\_evaluate\_identities}: all gate
polynomials first (each tagged with the index of its first simple selector),
then the permutation family, then the lookup families, then trash. Gate
identities carry a selector tag and are routed to buckets; permutation, lookup
and trash identities carry \texttt{None} and are routed to $A$. This is
precisely what Table~\ref{tab:qvmn:stream} shows.

\paragraph{Cross-check of the y-batch positions.} The tail block at lines
3314--3353 multiplies bucket $s$ by $y^{p_s}$ with
$(p_0,\dots,p_9) = (48,47,44,41,38,35,32,29,26,20)$, read from the y-power table
offsets \texttt{0x0600}, \texttt{0x05e0}, \texttt{0x0580}, \texttt{0x0520},
\texttt{0x04c0}, \texttt{0x0460}, \texttt{0x0400}, \texttt{0x03a0},
\texttt{0x0340}, \texttt{0x0280}. Equation~\eqref{eq:qvmn:reverse-fold} requires
$p_s = N-1-k_{\mathrm{last}}(s) = 48 - k_{\mathrm{last}}(s)$, giving
$k_{\mathrm{last}} = (0,1,4,7,10,13,16,19,22,28)$. Every one of these matches
Table~\ref{tab:qvmn:stream}, including the three buckets whose last or
intermediate contribution comes from a native \texttt{0x1b} callback rather than
from the interpreter ($B_3$ at $k=5$, $B_4$ at $k=8$, $B_9$ at $k=26,27$). The
native callbacks are therefore \emph{position-consistent} with the VM: they
consume exactly one y-batch slot per identity, they apply the same gap
convention, and the tail correction derived from the interpreted stream is
correct for them too.

\subsection{Shared preconditions of the three native arms}
\label{sec:qvmn:preconditions}

All three native arms open with the same three statements before running any
generated Yul (lines 2457--2464, 2574--2580, 2770--2772):

\begin{lstlisting}[style=solidity,caption={\texttt{Halo2Verifier.sol} lines 2452--2468: the \texttt{0x19} entry preamble. The \texttt{0x1f} (2570--2583) and \texttt{0x1b} (2762--2774) preambles are identical modulo comments and the \texttt{0x1b} operand read.},label={lst:qvmn:preamble}]
                    // VM 0x19 NATIVE_PERMUTATION: marker for the generated permutation callback.
                    case 0x19 {
                        // Native callbacks are identity-boundary opcodes. They
                        // must not inherit any partially evaluated VM stack
                        // state from the previous expression.
                        q_top := 0
                        q_has_top := 0
                        // The generated loop below uses program.stack_mptr as
                        // its scratch-table base, not as a conventional VM
                        // stack. The Rust memory planner must reserve enough
                        // words for structured_permutation_scratch_words(meta)
                        // whenever this opcode can appear.
                        if iszero(eq(q_sp, 0xb8e0)) { q_program_fail() }
                        // The generated lines below call the same fold snippets
                        // used by interpreted expressions, so trace IDs and
                        // y-batch positions remain contiguous.
                        {
\end{lstlisting}

\paragraph{The stack-boundary requirement.} The single runtime guard is
\yul{if iszero(eq(q\_sp, 0xb8e0)) \{ q\_program\_fail() \}}. It asserts that the
\emph{spilled} operand stack is empty, i.e.\ that the callback is invoked at an
identity boundary. This is not decorative: the callbacks reuse
$[\mathtt{0xb8e0},\mathtt{0xbfa0})$ --- the very words the interpreter spills
into --- as a structured scratch table, so a non-empty spilled stack would be
silently overwritten. \yul{q\_program\_fail} (defined at lines 1890--1893)
stores the 4-byte selector \texttt{0x3cc81b89} and reverts with
\revert{QuotientProgramInvalid}.

\paragraph{Why the scratch aliasing is safe here.} The generator plans the stack
window as $\max(\text{VM depth},\,\text{native scratch words})$
(\texttt{src/lowering/quotient.rs::native\_callback\_scratch\_words}). For this
circuit the permutation table needs
$2\cdot 18 + 2\cdot 6 + 5 + 1 = 54$ words and the lookup table needs
$3 \cdot 4 = 12$ words, so $54$ words are reserved. The push guards in the
\texttt{0x05} and \texttt{0x21} arms use the literal \texttt{0xbfa0}
$=$ \texttt{0xb8e0} $+\ 54\cdot 32$, and the permutation table's last word
(\yul{q\_perm\_delta\_base\_ptr} $=$ \texttt{0xbf80}) is exactly the last legal
stack word.

The \yul{q\_sp} guard separates a callback's scratch from live \emph{interpreter}
operands; it does \emph{not} separate the callbacks from each other. The
permutation table (\texttt{0xb8e0}--\texttt{0xbfa0}) and the three lookup tables
(\texttt{0xb8e0}--\texttt{0xba60}) overlap outright, the \texttt{0x1b} sub-cases
use \texttt{0xb8e0} as a one-word result slot, and the inline prefix already
wrote \texttt{0xb8e0} at line 2145 before the VM started. The reason this is
sound is not disjointness but total initialisation: every scratch word an arm
reads is written by that same arm first. All $18+18+6+6+5$ permutation slots are
filled at lines 2480--2525 before family (A) runs; the lookup arm writes all
four \yul{q\_lookup\_f} slots in the loop at 2604--2609, seeds
\yul{q\_lookup\_prefix} slot $0$ at 2614 and \yul{q\_lookup\_suffix} slot $3$ at
2619, and only then runs the two directional loops. No stale word is ever read.
Outside the window nothing else is live during the quotient phase: the y-power
table's last word occupies \texttt{0xb8c0}--\texttt{0xb8df}, and
\mptr{PCS\_FINAL\_MSM\_MPTR} staging
(\texttt{0xb280} onward) belongs to a strictly later phase.

\begin{finding}[QVMN-1]{Native arms discard a live cached top instead of failing closed --- Informational}
All three native arms execute \yul{q\_top := 0} and \yul{q\_has\_top := 0}
unconditionally (\texttt{Halo2Verifier.sol}:2457--2458, 2574--2575,
2770--2771) \emph{before} the \yul{q\_sp} check, and there is no
\yul{if q\_has\_top \{ q\_program\_fail() \}}. The asymmetry is deliberate
(``must not inherit any partially evaluated VM stack state'') but it removes a
check the same block relies on elsewhere: the epilogue at line 3099 treats a
live \yul{q\_has\_top} as a program error, and the comment at lines 3100--3104
explains that the companion \yul{q\_sp} check exists precisely because a fold
that runs with more operands live ``silently drops an operand of the identity
from $\nu_y(x)$''. A generator bug that emitted a native callback immediately
after an unconsumed expression would produce exactly that failure --- one gate
identity vanishes from the y-batch, weakening soundness --- and no runtime
check would fire. \\[2pt]
\textbf{Reachability.} None from proof calldata. The bytecode is copied from an
address whose \yul{extcodesize} and \yul{extcodehash} are pinned and re-checked
on every proof (lines 1386--1393), so the opcode stream is immutable relative to
the deployed verifier; the exposure is entirely to generator regressions.
Contrast \opcode{fold\_selector} at line 3068, which \emph{does} fail closed on
the opposite condition (\yul{iszero(q\_has\_top)}). \\[2pt]
\textbf{Cost of fixing.} One \yul{if} per arm, three arms, roughly 15 gas per
executed callback. Given that \yul{q\_sp} is already compared against a literal
two lines later, the omission buys almost nothing.
\end{finding}

\subsection{\texttt{0x19} \opcode{native\_permutation}}
\label{sec:qvmn:perm}

\paragraph{Operand encoding.} None. The opcode is a single byte; \yul{q\_pc} is
advanced only by the 1 byte consumed by the dispatch fetch at line 2205. It
occurs once, at program byte \texttt{0x11cd}.

\paragraph{Scratch layout.} Lines 2469--2479 bind eleven generated literals.
\begin{lstlisting}[style=solidity,caption={\texttt{Halo2Verifier.sol} lines 2469--2479: permutation scratch pointers and shape constants.},label={lst:qvmn:perm-consts}]
                        let delta := 0x8634d0aa021aaf843cab354fabb0062f6502437c6a09c006c083479590189d7
                        let q_perm_vals := 0xb8e0
                        let q_perm_sigmas := 0xbb20
                        let q_perm_z_cur := 0xbd60
                        let q_perm_z_next := 0xbe20
                        let q_perm_z_last := 0xbee0
                        let q_perm_delta_base_ptr := 0xbf80
                        let q_perm_num_cols := 18
                        let q_perm_num_sets := 6
                        let q_perm_chunk_len := 3
                        let q_perm_delta_chunk := 0x4285088329c399ea457a8ca1d30f8957e74c7f529842a1579b4fee55b3982923
\end{lstlisting}

\begin{table}[h]
\centering
\small
\begin{tabular}{llll}
\toprule
Table & Base & Words & Contents \\
\midrule
\yul{q\_perm\_vals}           & \texttt{0xb8e0} & 18 & $v_0,\dots,v_{17}$, the permuted columns' evals at $x$ \\
\yul{q\_perm\_sigmas}         & \texttt{0xbb20} & 18 & $\sigma_0(x),\dots,\sigma_{17}(x)$ \\
\yul{q\_perm\_z\_cur}         & \texttt{0xbd60} & 6  & $z_0(x),\dots,z_5(x)$ \\
\yul{q\_perm\_z\_next}        & \texttt{0xbe20} & 6  & $z_0(\omega x),\dots,z_5(\omega x)$ \\
\yul{q\_perm\_z\_last}        & \texttt{0xbee0} & 5  & $z_0(\omega^{-(t+1)}x),\dots,z_4(\omega^{-(t+1)}x)$ \\
\yul{q\_perm\_delta\_base\_ptr} & \texttt{0xbf80} & 1 & running $\beta x \delta^{3i}$ \\
\midrule
\multicolumn{2}{l}{total} & 54 & $= 2\cdot18 + 2\cdot 6 + (6-1) + 1$ \\
\bottomrule
\end{tabular}
\caption{Native permutation scratch table. Bases are consecutive; the table
exactly fills the registered stack window $[\mathtt{0xb8e0},\mathtt{0xbfa0})$.}
\label{tab:qvmn:perm-scratch}
\end{table}

\paragraph{Table fill.} Lines 2480--2525 populate the six tables from the eval
buffer with a mixture of single \yul{mstore}s and counted copy loops. The
generated copy loops carry two independent stride variables so that a source run
of stride $\ne 32$ can be gathered into a stride-32 destination:

\begin{lstlisting}[style=solidity,caption={\texttt{Halo2Verifier.sol} lines 2480--2497: the unit-stride table fills for \yul{q\_perm\_vals}. Lines 2498--2504 repeat the same pattern for \yul{q\_perm\_sigmas}; lines 2505--2525 use the same loop skeleton with \yul{src\_off := mul(i, 0x60)} to gather the three permutation-$z$ tables out of their interleaved triples.},label={lst:qvmn:perm-fill}]
                        mstore(add(q_perm_vals, 0x0), mload(0x99a0))
                        {
                        for { let q_perm_val_load_i := 0 } lt(q_perm_val_load_i, 5) { q_perm_val_load_i := add(q_perm_val_load_i, 1) } {
                        let q_perm_val_load_dst_off := shl(5, q_perm_val_load_i)
                        let q_perm_val_load_src_off := q_perm_val_load_dst_off
                        mstore(add(add(q_perm_vals, 0x20), q_perm_val_load_dst_off), mload(add(0x94a0, q_perm_val_load_src_off)))
                        }
                        }
                        mstore(add(q_perm_vals, 0xc0), mload(0x9480))
                        mstore(add(q_perm_vals, 0xe0), mload(INSTANCE_EVAL_MPTR))
                        {
                        for { let q_perm_val_load_i := 0 } lt(q_perm_val_load_i, 9) { q_perm_val_load_i := add(q_perm_val_load_i, 1) } {
                        let q_perm_val_load_dst_off := shl(5, q_perm_val_load_i)
                        let q_perm_val_load_src_off := q_perm_val_load_dst_off
                        mstore(add(add(q_perm_vals, 0x100), q_perm_val_load_dst_off), mload(add(0x95a0, q_perm_val_load_src_off)))
                        }
                        }
                        mstore(add(q_perm_vals, 0x220), mload(0x9980))
\end{lstlisting}

The resulting source map is given in Table~\ref{tab:qvmn:perm-src}. Note that
column $7$ is the \emph{instance} column: its value is not read from the proof
eval buffer at all but from \mptr{INSTANCE\_EVAL\_MPTR}, the barycentric
evaluation of the public inputs computed by the verifier itself. The
advice annotations are the ones the artifact's own generated variable names
disclose (\yul{a\_i} for advice, \yul{f\_i} for fixed); \texttt{0x99a0} and
\texttt{0x9480} carry no such name anywhere in the file, so their column kind
cannot be established from the fixture and is not asserted here. It does not
affect the identity: \eqref{eq:qvmn:permD} treats advice, fixed and instance
columns identically, exactly as \texttt{permutation.rs::expressions} does after
its \texttt{column\_type()} match.

\begin{table}[h]
\centering
\small
\begin{tabular}{lll}
\toprule
Table slot & Source address & Eval index \\
\midrule
\yul{vals[0]}      & \texttt{0x99a0} & $e_{41}$ (kind not disclosed) \\
\yul{vals[1..5]}   & \texttt{0x94a0} $+\,32i$, $i<5$ & $e_1,\dots,e_5$ (advice) \\
\yul{vals[6]}      & \texttt{0x9480} & $e_0$ (kind not disclosed) \\
\yul{vals[7]}      & \mptr{INSTANCE\_EVAL\_MPTR} (\texttt{0x7ce0}) & --- (instance) \\
\yul{vals[8..16]}  & \texttt{0x95a0} $+\,32i$, $i<9$ & $e_9,\dots,e_{17}$ (advice) \\
\yul{vals[17]}     & \texttt{0x9980} & $e_{40}$ (advice) \\
\yul{sigmas[0..17]}& \texttt{0x9bc0} $+\,32i$, $i<18$ & $e_{58},\dots,e_{75}$ \\
\yul{z\_cur[0..5]} & \texttt{0x9e00} $+\,\texttt{0x60}\,i$, $i<6$ & $e_{76},e_{79},e_{82},e_{85},e_{88},e_{91}$ \\
\yul{z\_next[0..5]}& \texttt{0x9e20} $+\,\texttt{0x60}\,i$, $i<6$ & $e_{77},e_{80},e_{83},e_{86},e_{89},e_{92}$ \\
\yul{z\_last[0..4]}& \texttt{0x9e40} $+\,\texttt{0x60}\,i$, $i<5$ & $e_{78},e_{81},e_{84},e_{87},e_{90}$ \\
\bottomrule
\end{tabular}
\caption{Permutation scratch source map, derived from lines 2480--2525. Eval
indices are $(\text{address}-\texttt{0x9480})/32$.}
\label{tab:qvmn:perm-src}
\end{table}

The stride-\texttt{0x60} gathers encode the fact that the permutation $z$
evaluations are stored in the proof as consecutive
$(z_i(x), z_i(\omega x), z_i(\omega^{-(t+1)}x))$ triples. The final set has no
last-rotation query, so the \yul{z\_last} loop bound is $5$ rather than $6$;
had it been $6$ it would have read \texttt{0xa020}, which holds the first lookup
multiplicity eval. The bound is therefore load-bearing and correct.

\paragraph{Identity families.} Lines 2526--2561 evaluate four families in the
order fixed by \texttt{proofs/src/plonk/permutation.rs::expressions}.

\begin{lstlisting}[style=solidity,caption={\texttt{Halo2Verifier.sol} lines 2526--2541: families A, B and C, and the initialisation of the running $\delta$ base.},label={lst:qvmn:perm-abc}]
                        let q_perm_eval := 0
                        q_perm_eval := mulmod(mload(L_0_MPTR), addmod(1, sub(r, mload(q_perm_z_cur)), r), r)
                        mstore(0xb280, mulmod(mload(0xb280), y, r))
                        mstore(0xb280, addmod(mload(0xb280), q_perm_eval, r))
                        let q_perm_zn := mload(add(q_perm_z_cur, 0xa0))
                        q_perm_eval := mulmod(mload(L_LAST_MPTR), addmod(mulmod(q_perm_zn, q_perm_zn, r), sub(r, q_perm_zn), r), r)
                        mstore(0xb280, mulmod(mload(0xb280), y, r))
                        mstore(0xb280, addmod(mload(0xb280), q_perm_eval, r))
                        for { let q_perm_i := 1 } lt(q_perm_i, 6) { q_perm_i := add(q_perm_i, 1) } {
                        let q_perm_cur := mload(add(q_perm_z_cur, shl(5, q_perm_i)))
                        let q_perm_prev := mload(add(q_perm_z_last, shl(5, sub(q_perm_i, 1))))
                        q_perm_eval := mulmod(mload(L_0_MPTR), addmod(q_perm_cur, sub(r, q_perm_prev), r), r)
                        mstore(0xb280, mulmod(mload(0xb280), y, r))
                        mstore(0xb280, addmod(mload(0xb280), q_perm_eval, r))
                        }
                        mstore(q_perm_delta_base_ptr, mulmod(mload(BETA_MPTR), mload(X_MPTR), r))
\end{lstlisting}

\begin{align}
\text{(A), }k=29:\quad & \varepsilon \;=\; L_0(x)\,\bigl(1 - z_0(x)\bigr) \label{eq:qvmn:permA}\\
\text{(B), }k=30:\quad & \varepsilon \;=\; L_{\mathrm{last}}(x)\,\bigl(z_5(x)^2 - z_5(x)\bigr) \label{eq:qvmn:permB}\\
\text{(C), }k=30+i,\ i=1..5:\quad & \varepsilon \;=\; L_0(x)\,\bigl(z_i(x) - z_{i-1}(\omega^{-(t+1)}x)\bigr) \label{eq:qvmn:permC}
\end{align}

Family (B) uses \yul{mload(add(q\_perm\_z\_cur, 0xa0))}, i.e.\ slot
$5 = \texttt{num\_sets}-1$, so it is the \emph{last} set's product evaluated at
$x$, exactly \texttt{sets.last()} in the reference. Family (C) chains set $i$ at
$x$ to set $i-1$ at the last-rotation point, which is what makes six independent
grand products behave as one.

\begin{lstlisting}[style=solidity,caption={\texttt{Halo2Verifier.sol} lines 2542--2561: family D, the per-set active-row product equality.},label={lst:qvmn:perm-d}]
                        for { let q_perm_set := 0 } lt(q_perm_set, 6) { q_perm_set := add(q_perm_set, 1) } {
                        let q_perm_start := mul(q_perm_set, q_perm_chunk_len)
                        let q_perm_end := add(q_perm_start, q_perm_chunk_len)
                        if gt(q_perm_end, q_perm_num_cols) { q_perm_end := q_perm_num_cols }
                        let q_perm_left := mload(add(q_perm_z_next, shl(5, q_perm_set)))
                        let q_perm_right := mload(add(q_perm_z_cur, shl(5, q_perm_set)))
                        let q_perm_delta_pow := mload(q_perm_delta_base_ptr)
                        for { let q_perm_j := q_perm_start } lt(q_perm_j, q_perm_end) { q_perm_j := add(q_perm_j, 1) } {
                        let q_perm_off := shl(5, q_perm_j)
                        let q_perm_v := mload(add(q_perm_vals, q_perm_off))
                        let q_perm_s := mload(add(q_perm_sigmas, q_perm_off))
                        q_perm_left := mulmod(q_perm_left, addmod(addmod(q_perm_v, mulmod(mload(BETA_MPTR), q_perm_s, r), r), mload(GAMMA_MPTR), r), r)
                        q_perm_right := mulmod(q_perm_right, addmod(addmod(q_perm_v, q_perm_delta_pow, r), mload(GAMMA_MPTR), r), r)
                        q_perm_delta_pow := mulmod(q_perm_delta_pow, delta, r)
                        }
                        q_perm_eval := mulmod(addmod(1, sub(r, addmod(mload(L_LAST_MPTR), mload(L_BLIND_MPTR), r)), r), addmod(q_perm_left, sub(r, q_perm_right), r), r)
                        mstore(0xb280, mulmod(mload(0xb280), y, r))
                        mstore(0xb280, addmod(mload(0xb280), q_perm_eval, r))
                        mstore(q_perm_delta_base_ptr, mulmod(mload(q_perm_delta_base_ptr), q_perm_delta_chunk, r))
                        }
\end{lstlisting}

For each set $i = 0,\dots,5$ with column range $[3i, 3i+3)$, at y-batch
position $k = 36+i$:
\begin{equation}
\varepsilon_{36+i} \;=\;
\Bigl(1 - \bigl(L_{\mathrm{last}}(x)+L_{\mathrm{blind}}(x)\bigr)\Bigr)
\left(
 z_i(\omega x)\!\!\prod_{j=3i}^{3i+2}\!\!\bigl(v_j + \beta\,\sigma_j(x) + \gamma\bigr)
 \;-\;
 z_i(x)\!\!\prod_{j=3i}^{3i+2}\!\!\bigl(v_j + \beta\,x\,\delta^{\,j} + \gamma\bigr)
\right).
\label{eq:qvmn:permD}
\end{equation}

The $\delta$ bookkeeping deserves a line-by-line justification, because the code
computes $\delta^{\,j}$ incrementally where the reference recomputes it. Before
the outer loop, \yul{q\_perm\_delta\_base\_ptr} holds $\beta x$. Inside set $i$
the local \yul{q\_perm\_delta\_pow} is seeded from that word and multiplied by
$\delta$ once per column, so at column $j$ it equals
$\beta x \delta^{\,3i} \cdot \delta^{\,j-3i} = \beta x \delta^{\,j}$; at the end
of set $i$ the base word is multiplied by
\yul{q\_perm\_delta\_chunk}, advancing it to $\beta x \delta^{\,3(i+1)}$. This
matches the reference's
\texttt{current\_delta = (beta * x) * DELTA.pow(chunk\_index * chunk\_len)}
followed by \texttt{current\_delta *= DELTA} per column.

\paragraph{Constant verification.} Both generated field constants were checked
against the BLS12-381 scalar field:
\begin{itemize}
\item $\delta = \texttt{0x8634d0aa021aaf843cab354fabb0062f6502437c6a09c006c083479590189d7}$
satisfies $\delta < r$ and $\delta^{\,(r-1)/2^{32}} = 1$, and its multiplicative
order is exactly $(r-1)/2^{32}$ --- an \emph{odd} integer of about $223$ bits,
\emph{not} $2^{32}$. That is the defining shape of \texttt{PrimeField::DELTA},
which \texttt{ff} derives as $g^{2^{S}}$ from the field generator $g$ and the
two-adicity $S = 32$ of \Fq{}: it generates the odd-order part of the
multiplicative group. The column-separation argument needs exactly that, not a
large order. Writing $H$ for the multiplicative subgroup of order $n = 2^{20}$
used as the evaluation domain, $\langle\delta\rangle \cap H = \{1\}$ because
$|\langle\delta\rangle|$ is odd while $|H|$ is a power of two; hence
$\delta^{\,j}H = \delta^{\,j'}H$ forces $\delta^{\,j-j'} = 1$, which for
$0 \le j,j' < 18$ forces $j = j'$. The $18$ cosets $\delta^{\,j}H$ are therefore
pairwise disjoint, and $\delta^0,\dots,\delta^{17}$ are pairwise distinct
(checked numerically).
\item $\texttt{q\_perm\_delta\_chunk} = \delta^{3} \bmod r$ exactly.
\end{itemize}

\paragraph{Faithfulness.} Line-for-line, listings~\ref{lst:qvmn:perm-abc} and
\ref{lst:qvmn:perm-d} reproduce
\texttt{proofs/src/plonk/permutation.rs::expressions} --- same four families,
same order within families, same operand order, same use of $L_0$,
$L_{\mathrm{last}}$, $L_{\mathrm{blind}}$. The only structural difference is the
incremental $\delta$ base, justified above. The 13 folds land at positions
$29..41$, contiguously and in reference order.

\begin{finding}[QVMN-2]{The permutation chunk clamp is one-sided: column coverage is not checked at runtime --- Observation}
The outer loop at line 2542 runs $6$ times over chunks of
\yul{q\_perm\_chunk\_len} $= 3$, and line 2545 clamps a chunk that would run
past \yul{q\_perm\_num\_cols} $= 18$. The clamp handles a short \emph{final}
chunk, but nothing checks the opposite direction: if the generated shape
constants ever satisfied $\text{num\_sets} \cdot \text{chunk\_len} <
\text{num\_cols}$, the trailing columns would be dropped from \emph{both}
products in \eqref{eq:qvmn:permD}, and the identity would still be satisfiable
--- a silent soundness loss with no revert. \\[2pt]
Worth noting while reading lines 2469--2479: of the three shape constants only
\yul{q\_perm\_num\_cols} and \yul{q\_perm\_chunk\_len} are ever read
(lines 2543--2545). \yul{q\_perm\_num\_sets} (line 2477) is a dead binding ---
both loops that need the set count, family (C) at line 2534 and family (D) at
line 2542, use the bare literal \texttt{6}. The set count therefore appears
three times in the artifact and only the emitter keeps the three copies equal.
\\[2pt]
For this artifact the condition is vacuous: $6 \cdot 3 = 18$ columns
exactly, so no clamp ever fires and every column is covered once. All three
values are generator literals baked into VK-pinned runtime code, so the
exposure is to generator regressions only, and it is the kind of regression the
per-identity trace differential would catch. Recorded so a reader of the
generated code does not mistake the one-sided clamp for a completeness check.
\end{finding}

\subsection{\texttt{0x1f} \opcode{native\_lookup}}
\label{sec:qvmn:lookup}

\paragraph{Operand encoding.} None; a single byte, occurring once at program
byte \texttt{0x11ce}. Like \texttt{0x19} it is a \emph{whole-family} opcode: one
byte stands for all six LogUp identities of both lookup arguments.

\paragraph{Scratch layout and hoisted values.} Lines 2584--2593 bind three
scratch bases and seven hoisted scalars.

\begin{lstlisting}[style=solidity,caption={\texttt{Halo2Verifier.sol} lines 2584--2599: lookup scratch bases, hoisted Lagrange/challenge values, and the first boundary identity.},label={lst:qvmn:lookup-head}]
                        let q_lookup_f := 0xb8e0
                        let q_lookup_prefix := 0xb960
                        let q_lookup_suffix := 0xb9e0
                        let q_lookup_l0 := mload(L_0_MPTR)
                        let q_lookup_llast := mload(L_LAST_MPTR)
                        let q_lookup_lblind := mload(L_BLIND_MPTR)
                        let q_lookup_lsum := addmod(q_lookup_l0, q_lookup_llast, r)
                        let q_lookup_active := addmod(1, sub(r, addmod(q_lookup_llast, q_lookup_lblind, r)), r)
                        let q_lookup_beta := mload(BETA_MPTR)
                        let q_lookup_theta := mload(THETA_MPTR)
                        {
                        {
                        let q_lookup_eval := mulmod(q_lookup_lsum, mload(0xa060), r)
                        mstore(0xb280, mulmod(mload(0xb280), y, r))
                        mstore(0xb280, addmod(mload(0xb280), q_lookup_eval, r))
                        }
\end{lstlisting}

The three tables are $4$ words each (\yul{max\_parallel} $= 4$, the widest
helper chunk in the circuit): the $f{+}\beta$ table at \texttt{0xb8e0}, prefix products at
\texttt{0xb960}, suffix products at \texttt{0xb9e0}, ending at \texttt{0xba60}
--- $12$ of the $54$ reserved words.

\paragraph{Notation.} Write $\langle g_0,\dots,g_{c-1}\rangle_\theta$ for the
$\theta$-compression $\sum_{i} g_i\,\theta^{\,c-1-i}$, which the generated code
computes by Horner from a zero seed:
\yul{addmod(mulmod(acc, theta, r), g\_i, r)} with \yul{acc} initialised as
\yul{mulmod(0, theta, r)}. For lookup $\ell$ let $m_\ell$ be the multiplicity
eval, $h_\ell$ the helper eval, $\varphi_\ell,\varphi_\ell^{\omega}$ the
accumulator evals at $x$ and $\omega x$, $s_\ell$ the selector expression,
$t_\ell$ the compressed table, and $f_{\ell,0},\dots,f_{\ell,k-1}$ the
compressed inputs of the (single) helper chunk.

\begin{table}[h]
\centering
\small
\begin{tabular}{lllll}
\toprule
$\ell$ & $m_\ell$ & $h_\ell$ & $\varphi_\ell$ & $\varphi_\ell^{\omega}$ \\
\midrule
$0$ & \texttt{0xa020} ($e_{93}$) & \texttt{0xa040} ($e_{94}$) & \texttt{0xa060} ($e_{95}$) & \texttt{0xa080} ($e_{96}$) \\
$1$ & \texttt{0xa0a0} ($e_{97}$) & \texttt{0xa0c0} ($e_{98}$) & \texttt{0xa0e0} ($e_{99}$) & \texttt{0xa100} ($e_{100}$) \\
\bottomrule
\end{tabular}
\caption{Lookup eval addresses used by the \texttt{0x1f} arm.}
\label{tab:qvmn:lookup-evals}
\end{table}

\paragraph{The three identities, per lookup.} In generated order (matching
\texttt{proofs/src/plonk/logup.rs::expressions}, which chains
\texttt{once(boundary)} then the helper constraints then the accumulator
constraint):
\begin{align}
\text{boundary:}\quad & \varepsilon \;=\; \bigl(L_0(x)+L_{\mathrm{last}}(x)\bigr)\,\varphi_\ell(x) \label{eq:qvmn:lkbnd}\\
\text{helper:}\quad & \varepsilon \;=\; h_\ell(x)\prod_{j=0}^{k-1}\bigl(f_{\ell,j}(x)+\beta\bigr) \;-\; \sum_{j=0}^{k-1}\ \prod_{\substack{m=0\\ m\ne j}}^{k-1}\bigl(f_{\ell,m}(x)+\beta\bigr) \label{eq:qvmn:lkhlp}\\
\text{accumulator:}\quad & \varepsilon \;=\; \Bigl(1-\bigl(L_{\mathrm{last}}(x)+L_{\mathrm{blind}}(x)\bigr)\Bigr)\Bigl[\bigl(\varphi_\ell^{\omega}-\varphi_\ell-s_\ell h_\ell\bigr)\bigl(t_\ell+\beta\bigr)+m_\ell\Bigr] \label{eq:qvmn:lkacc}
\end{align}
Together, \eqref{eq:qvmn:lkhlp} pins $h_\ell = \sum_j (f_{\ell,j}+\beta)^{-1}$
(as a polynomial identity that stays total when some $f_{\ell,j}+\beta$
vanishes), and \eqref{eq:qvmn:lkacc} is the LogUp telescoping relation
$\varphi(\omega x) - \varphi(x) = s\,h(x) - m(x)/(t(x)+\beta)$ restricted to
active rows, with \eqref{eq:qvmn:lkbnd} closing the telescope at both ends.

\paragraph{Lookup 0: the shared-prefix chunk ($k=4$).}

\begin{lstlisting}[style=solidity,caption={\texttt{Halo2Verifier.sol} lines 2601--2631: shared-prefix compression, product, prefix/suffix tables, and the helper fold for lookup 0.},label={lst:qvmn:lookup-chunk}]
                        let f_10 := mload(0x9ae0)
                        let var0 := addmod(mulmod(0, q_lookup_theta, r), f_10, r)
                        let var1 := mulmod(var0, q_lookup_theta, r)
                        for { let q_lookup_shared_i := 0 } lt(q_lookup_shared_i, 4) { q_lookup_shared_i := add(q_lookup_shared_i, 1) } {
                        let q_lookup_shared_off := shl(5, q_lookup_shared_i)
                        let q_lookup_shared_tail := mload(add(0x94c0, q_lookup_shared_off))
                        let q_lookup_shared_compressed := addmod(var1, q_lookup_shared_tail, r)
                        mstore(add(q_lookup_f, q_lookup_shared_off), addmod(q_lookup_shared_compressed, q_lookup_beta, r))
                        }
                        let q_lookup_product := 1
                        for { let q_lookup_prod_i := 0 } lt(q_lookup_prod_i, 4) { q_lookup_prod_i := add(q_lookup_prod_i, 1) } {
                        q_lookup_product := mulmod(q_lookup_product, mload(add(q_lookup_f, shl(5, q_lookup_prod_i))), r)
                        }
                        mstore(q_lookup_prefix, 1)
                        for { let q_lookup_pref_i := 1 } lt(q_lookup_pref_i, 4) { q_lookup_pref_i := add(q_lookup_pref_i, 1) } {
                        let q_lookup_pref_prev := sub(q_lookup_pref_i, 1)
                        mstore(add(q_lookup_prefix, shl(5, q_lookup_pref_i)), mulmod(mload(add(q_lookup_prefix, shl(5, q_lookup_pref_prev))), mload(add(q_lookup_f, shl(5, q_lookup_pref_prev))), r))
                        }
                        mstore(add(q_lookup_suffix, 0x60), 1)
                        for { let q_lookup_suf_i := sub(4, 1) } gt(q_lookup_suf_i, 0) { q_lookup_suf_i := sub(q_lookup_suf_i, 1) } {
                        let q_lookup_suf_prev := sub(q_lookup_suf_i, 1)
                        mstore(add(q_lookup_suffix, shl(5, q_lookup_suf_prev)), mulmod(mload(add(q_lookup_suffix, shl(5, q_lookup_suf_i))), mload(add(q_lookup_f, shl(5, q_lookup_suf_i))), r))
                        }
                        let q_lookup_sum := 0
                        for { let q_lookup_sum_i := 0 } lt(q_lookup_sum_i, 4) { q_lookup_sum_i := add(q_lookup_sum_i, 1) } {
                        q_lookup_sum := addmod(q_lookup_sum, mulmod(mload(add(q_lookup_prefix, shl(5, q_lookup_sum_i))), mload(add(q_lookup_suffix, shl(5, q_lookup_sum_i))), r), r)
                        }
                        let q_lookup_eval := addmod(mulmod(mload(0xa040), q_lookup_product, r), sub(r, q_lookup_sum), r)
                        mstore(0xb280, mulmod(mload(0xb280), y, r))
                        mstore(0xb280, addmod(mload(0xb280), q_lookup_eval, r))
                        }
\end{lstlisting}

The four parallel inputs of lookup 0 are the 2-tuples $(f_{10}, a_{1+j})$ for
$j = 0,\dots,3$: they share their entire prefix and their differing tail evals
sit at consecutive addresses \texttt{0x94c0}, \texttt{0x94e0}, \texttt{0x9500},
\texttt{0x9520} ($e_2..e_5$). The generator's shared-prefix rule
(\texttt{yul\_emit.rs::lookup\_shared\_prefix\_base\_ptr}: chunk length
$\ge 3$, identical first $c-1$ elements, last elements at
$\text{base}+32j$) fires, so the prefix is compressed once
(\yul{var1} $= f_{10}\theta$) and the loop only adds the tail:
\begin{equation}
f_{0,j} \;+\; \beta \;=\; \bigl(f_{10}\,\theta + a_{1+j}\bigr) + \beta,
\qquad j = 0,\dots,3 .
\end{equation}
This is bit-for-bit the same value the general path would compute, since
$\langle f_{10}, a_{1+j}\rangle_\theta = (0\cdot\theta + f_{10})\theta + a_{1+j}$.

Writing $F_j$ for the word stored at \yul{q\_lookup\_f} slot $j$, namely
$f_{0,j}+\beta$, the partial products $\prod_{m\ne j} F_m$ are built as
$\mathrm{pre}[j]\cdot\mathrm{suf}[j]$ with
$\mathrm{pre}[0]=1$, $\mathrm{pre}[j]=\mathrm{pre}[j-1]\cdot F_{j-1}$ and
$\mathrm{suf}[3]=1$, $\mathrm{suf}[j-1]=\mathrm{suf}[j]\cdot F_{j}$, so
$\mathrm{pre}[j]\cdot\mathrm{suf}[j] = \prod_{m<j} F_m \cdot \prod_{m>j} F_m$
as required. The literal \texttt{0x60} at the \yul{q\_lookup\_suffix} seed is
$32\cdot(k-1)$ for $k=4$; the descending loop uses \yul{gt(..., 0)} so it
terminates at index $1$ and writes index $0$ last.

\paragraph{Lookup 0: the accumulator.}
\begin{lstlisting}[style=solidity,caption={\texttt{Halo2Verifier.sol} lines 2632--2646: the lookup-0 accumulator identity.},label={lst:qvmn:lookup-acc}]
                        {
                        let q_lookup_sum_h := mload(0xa040)
                        let f_17 := mload(0x9b60)
                        let f_11 := mload(0x9b00)
                        let var0 := addmod(mulmod(0, q_lookup_theta, r), f_11, r)
                        let f_12 := mload(0x9b20)
                        let var1 := addmod(mulmod(var0, q_lookup_theta, r), f_12, r)
                        let q_lookup_s_sum_h := mulmod(f_17, q_lookup_sum_h, r)
                        let q_lookup_diff := addmod(mload(0xa080), sub(r, addmod(mload(0xa060), q_lookup_s_sum_h, r)), r)
                        let q_lookup_t_beta := addmod(var1, q_lookup_beta, r)
                        let q_lookup_core := addmod(mulmod(q_lookup_diff, q_lookup_t_beta, r), mload(0xa020), r)
                        let q_lookup_eval := mulmod(q_lookup_active, q_lookup_core, r)
                        mstore(0xb280, mulmod(mload(0xb280), y, r))
                        mstore(0xb280, addmod(mload(0xb280), q_lookup_eval, r))
                        }
\end{lstlisting}
Here $s_0 = f_{17} = e_{55}$, $t_0 = \langle f_{11}, f_{12}\rangle_\theta$ with
$f_{11} = e_{52}$, $f_{12} = e_{53}$, and $\sum h = h_0$ (one chunk). The
selector multiplies only the input side, exactly as the reference documents
(``the selector gates only the input side; multiplicities are always subtracted
so the table-side balance is maintained'').

\paragraph{Lookup 1: the $k=1$ specialisation.} Lines 2648--2753. The boundary
identity (2650) is \eqref{eq:qvmn:lkbnd} with $\varphi_1 = e_{99}$. The helper
chunk has $k=1$, so the general prefix/suffix machinery collapses --- the
product is $f_{1,0}+\beta$ and the empty-product sum is $1$ --- and the emitter
renders \eqref{eq:qvmn:lkhlp} directly, with no scratch use at all:
\begin{lstlisting}[style=solidity,caption={\texttt{Halo2Verifier.sol} lines 2686--2690: the $k=1$ helper identity for lookup 1. Lines 2655--2686 compress the 16-element input tuple $\langle a_{14},a_0,\dots,a_{13},f_{13}\rangle_\theta$ into \yul{var15} by straight-line Horner; that compression is elided here.},label={lst:qvmn:lookup-k1}]
                        let var15 := addmod(mulmod(var14, q_lookup_theta, r), f_13, r)
                        let q_lookup_eval := addmod(mulmod(mload(0xa0c0), addmod(var15, q_lookup_beta, r), r), sub(r, 1), r)
                        mstore(0xb280, mulmod(mload(0xb280), y, r))
                        mstore(0xb280, addmod(mload(0xb280), q_lookup_eval, r))
                        }
\end{lstlisting}
The accumulator identity (2691--2752) has $s_1 = 1$ (\yul{var0 := 0x1}, line
2693) and a table expression that is the same 16-element tuple with every
element scaled by $1 - f_{24}$ (\yul{var2}, line 2696, with
$f_{24} = e_{56}$). Because $\theta$-compression is linear in its inputs, the
compressed table is
\begin{equation}
t_1 \;=\; \bigl(1 - f_{24}\bigr)\cdot\bigl\langle a_{14},a_0,\dots,a_{13},f_{13}\bigr\rangle_\theta ,
\end{equation}
but the code does \emph{not} exploit that: it emits 16 explicit
\yul{mulmod(var2, ...)} scalings and re-Horners them (lines 2697--2744), which
is the faithful transcription of the expression tree and costs
$16$ extra \yul{mulmod}s. That is a deliberate choice --- the emitter walks the
tree rather than factoring it --- and it removes any risk of an algebraic
rewrite diverging from the reference.

\paragraph{Faithfulness.} The order (boundary, helper chunks, accumulator; per
lookup, lookups in index order), every operand, and every sign match
\texttt{logup.rs::expressions}. The six folds occupy positions $42..47$.

\begin{finding}[QVMN-3]{The on-chain helper identity is total where the native reference is partial --- Observation}
The reference computes the partial products as
\texttt{product * f.invert().unwrap()}
(\texttt{proofs/src/plonk/logup.rs}), which panics when any compressed input
satisfies $f_{\ell,j}(x) + \beta = 0$. The Solidity arm computes the same
quantity as $\mathrm{pre}[j]\cdot\mathrm{suf}[j]$ (lines 2614--2627), which is
defined for every input. \\[2pt]
The two agree on every input where the reference is defined, and the polynomial
identity \eqref{eq:qvmn:lkhlp} is the same in both, so this is not a divergence
in what is enforced --- it is the on-chain implementation being total where the
native one is partial. Worth recording for two reasons: a reviewer diffing the
two implementations will see structurally different code computing the same
thing, and the $f_{\ell,j}+\beta = 0$ case (probability $\approx k/r$ over
$\beta$) is one where the on-chain verifier keeps evaluating rather than
aborting. It remains sound: with $f_{\ell,j}+\beta = 0$ the identity degenerates
to $-\prod_{k\ne j}\bigl(f_{\ell,k}+\beta\bigr) = 0$, which still constrains the
prover (a second compressed input must also hit $-\beta$).
\end{finding}

\subsection{\texttt{0x1b} \opcode{native\_identity}}
\label{sec:qvmn:natid}

\paragraph{Operand encoding.} Two bytes: a big-endian \texttt{u16} callback
index, read as \yul{shr(240, mload(q\_pc))} and followed by
\yul{q\_pc := add(q\_pc, 2)}. Unlike \texttt{0x19}/\texttt{0x1f} this is a
\emph{per-identity} opcode: each occurrence produces exactly one identity.

\begin{lstlisting}[style=solidity,caption={\texttt{Halo2Verifier.sol} lines 2762--2775: operand decode, boundary reset, and the sub-dispatch.},label={lst:qvmn:natid-head}]
                    case 0x1b {
                        // Operand layout: u16 native callback index. The
                        // manifest validates that callback indexes appear in
                        // generated order and target existing switch cases.
                        let q_native_idx := shr(240, mload(q_pc))
                        q_pc := add(q_pc, 2)
                        // Heavy identities are whole expressions, so clear the
                        // interpreter stack before dispatching.
                        q_top := 0
                        q_has_top := 0
                        if iszero(eq(q_sp, 0xb8e0)) { q_program_fail() }
                        // Native identity sub-cases are generated from selected heavy gate identities.
                        switch q_native_idx
                        case 0 {
\end{lstlisting}

The operand is read \emph{before} the boundary reset and \emph{before} the
\yul{q\_sp} check, but nothing in the reset depends on it, so the ordering is
immaterial. Note that the read is a full 32-byte \yul{mload} of which only the
top 2 bytes are used; near \yul{q\_end} this reads past the program into the
adjacent VK payload region, which is allocated verifier memory and therefore
harmless. In this artifact all four occurrences are interior
(\texttt{0x0011}, \texttt{0x01f5}, \texttt{0x1196}, \texttt{0x1199}) and the
program's last two bytes are the zero-operand \texttt{0x19}/\texttt{0x1f}
opcodes, so \yul{q\_pc} lands exactly on \yul{q\_end}.

\paragraph{Sub-case dispatch and routing.} The four sub-cases each evaluate one
gate polynomial into a straight-line chain of \yul{var}s, store the result to
the scratch word \texttt{0xb8e0}, advance $A$ by $y$, and then add the stored
value into a selector bucket. The bucket index and gap factor are generated
literals; they bypass the \texttt{0x0b} bounds checks entirely (there is no
untrusted index to check).

\begin{table}[h]
\centering
\small
\begin{tabular}{lllll}
\toprule
\yul{q\_native\_idx} & Lines & Bucket offset & Bucket & Gap factor \\
\midrule
0 & 2775--2921 & \texttt{0x60}  & $B_3$ & none (first contribution) \\
1 & 2922--2972 & \texttt{0x80}  & $B_4$ & none (first contribution) \\
2 & 2973--3011 & \texttt{0x120} & $B_9$ & $y^1$ (\yul{mload(add(0xb2c0, 0x20))}) \\
3 & 3012--3050 & \texttt{0x120} & $B_9$ & $y^1$ \\
\bottomrule
\end{tabular}
\caption{The four \texttt{0x1b} sub-cases and their selector routing.}
\label{tab:qvmn:natid}
\end{table}

\begin{lstlisting}[style=solidity,caption={\texttt{Halo2Verifier.sol} lines 2912--2921 (sub-case 0, no gap) and 3001--3011 (sub-case 2, with a $y^1$ gap). The two fold shapes are the only two the generator emits.},label={lst:qvmn:natid-fold}]
                            let var116 := mulmod(var0, var115, r)
                            mstore(0xb8e0, var116)
                            }
                            mstore(0xb280, mulmod(mload(0xb280), y, r))
                            {
                            let q_selector_ptr := add(SELECTOR_ACC_MPTR, 0x60)
                            let q_selector_acc := mload(q_selector_ptr)
                            mstore(q_selector_ptr, addmod(q_selector_acc, mload(0xb8e0), r))
                            }
                        }
// ... sub-case 1 in full and sub-case 2's expression body elided ...
                            let var18 := mulmod(var0, var17, r)
                            mstore(0xb8e0, var18)
                            }
                            mstore(0xb280, mulmod(mload(0xb280), y, r))
                            {
                            let q_selector_ptr := add(SELECTOR_ACC_MPTR, 0x120)
                            let q_selector_acc := mload(q_selector_ptr)
                            q_selector_acc := mulmod(q_selector_acc, mload(add(0xb2c0, 0x20)), r)
                            mstore(q_selector_ptr, addmod(q_selector_acc, mload(0xb8e0), r))
                            }
                        }
\end{lstlisting}

The gap-free form is emitted only where the bucket is provably still zero (first
contribution), which Table~\ref{tab:qvmn:stream} confirms for $B_3$ at $k=5$ and
$B_4$ at $k=8$; multiplying zero by $y^{\text{gap}}$ would be a no-op anyway, so
the omission is a pure gas saving with no semantic content. The $y^1$ factor in
sub-cases 2 and 3 is correct because their predecessors in $B_9$ are at $k=25$
and $k=26$ respectively.

\paragraph{What the four identities are.} Each body is a plain evaluation of a
gate polynomial with its simple selector replaced by the literal $1$
(\yul{var0 := 0x1}, then a final \yul{mulmod(var0, ..., r)}). Structurally:

\begin{itemize}
\item \textbf{Sub-cases 0 and 1} are emulated-field limb identities over a
base-$2^{56}$ representation. Sub-case 0 (147 lines, \yul{var0}--\yul{var116})
reads two seven-limb operands --- $a_0..a_6$ at \texttt{0x94a0}, \texttt{0x94c0},
\texttt{0x94e0}, \texttt{0x9500}, \texttt{0x9520}, \texttt{0x95a0},
\texttt{0x95c0}, and the next-rotation limbs $a'_0..a'_6$ at \texttt{0x9540},
\texttt{0x9560}, \texttt{0x9580}, \texttt{0x96c0}, \texttt{0x96e0},
\texttt{0x9700}, \texttt{0x9720} --- and forms
\begin{equation}
\sum_{(i,j)\in I} c_{\,i+j}\; a_i\,a'_j
\;+\; \bigl(a_0 + 2^{56}a_1 + 2^{112}a_2\bigr)
\;+\; \bigl(a'_0 + 2^{56}a'_1 + 2^{112}a'_2\bigr)
\;-\; \bigl(a_7 + 2^{56}a_8 + 2^{112}a_9\bigr)
\;-\; c_7 \;-\; \kappa\,a'_7 \;-\; 2^{134}\bigl(a'_8 + \lambda\bigr),
\label{eq:qvmn:natid0}
\end{equation}
where $I = \{(i,j) : i+j \le 2\} \cup \{(i,j) : i+j \ge 7\}$ over
$0\le i,j\le 6$, the coefficient depends only on $i+j$
(Table~\ref{tab:qvmn:limbcoef}),
$\kappa = \texttt{0x241eabfffeb153ffffb9feffffffffaaab}$ --- which is exactly
$p \bmod 2^{134}$ for the BLS12-381 base field --- and
$\lambda = \texttt{0x73eda753299d7d483339d80809a1d80553b9202d7ffe85d4800008bb20000001}$,
i.e.\ $-\texttt{0x483d57fffd62a7ffff743e0000000} \bmod r$ (both checked
numerically; $\texttt{0x4000000000000000000000000000000000} = 2^{134}$ exactly).
The index set was read off the emitted \yul{var} chain term by term: the six
pairs with $i+j \le 2$ and all $21$ pairs with $i+j \ge 7$ appear exactly once,
and the band $3 \le i+j \le 6$ is absent. What the missing band means at the
circuit level is not disclosed by the fixture and is not asserted here.
\\[2pt]
Sub-case 1 (51 lines) is the affine sibling: the same base-$2^{56}$
recompositions with a per-limb $2^{112}$ bias, the same subtracted
$(a_7 + 2^{56}a_8 + 2^{112}a_9)$, the same $\kappa\,a'_7$ and
$2^{134}(a'_8 + \lambda')$ tail terms, and two extra constant offsets
$-2^{112}$ and $-\texttt{0xd9d44a30b019261257667fde3844a8cd6}$, with no
bilinear block at all. \\[2pt]
\textbf{Caveat.} The six large coefficients in
Table~\ref{tab:qvmn:limbcoef} are the generator's ``double base powers'' from
\texttt{circuits/src/field/foreign/params.rs} (see the \texttt{0x21} comment at
lines 2264--2279). They are \emph{not} $2^{56k} \bmod p$, and no single modulus
$m$ satisfies $c_{k+1} \equiv c_k \cdot 2^{56} \pmod m$ for $k=7..11$, so the
emulated modulus cannot be recovered from the fixture alone. This
specification therefore records the constants and the index set exactly and
does not assert what field they reduce.
\item \textbf{Sub-cases 2 and 3} (39 lines each) have the shape
\begin{equation}
f_j - a'_j \;+\; c_{j,0}\,a_0^{2}a_3 \;+\; c_{j,1}\,a_1^{2}a_4 \;+\; c_{j,2}\,a_2^{2}a_5,
\qquad j \in \{0,1\},
\end{equation}
reading $f_0$ at \texttt{0x9a60} and $f_1$ at \texttt{0x9a80}. The six
coefficients are distinct elements of $[0,r)$ --- bit lengths $253,254,254$ in
sub-case 2 and $255,252,252$ in sub-case 3. Combined with the three interpreted
identities at $k = 23,24,25$ (program bytes \texttt{0x115a}, \texttt{0x116e},
\texttt{0x1182}) --- each emitted as \opcode{push\_mem}, \opcode{mul\_mem}
twice, \opcode{push\_mem}, \opcode{neg}, \opcode{add},
\opcode{mul\_const\_u8} slot $0$, i.e.\ $a_i^{3} - a_{i+3}$ --- the cube
witnesses $a_{i+3} = a_i^{3}$ turn $a_i^{2}a_{i+3}$ into $a_i^{5}$, so
sub-cases 2 and 3 are two rows of a $3\times 3$-mixed $x^{5}$ S-box. The third
row is the interpreted identity at $k = 28$ (program bytes
\texttt{0x119c}--\texttt{0x11c8}), which has exactly the same shape with
$j = 2$: $f_2$ at \texttt{0x9aa0}, $a'_2$ at \texttt{0x9580}, and constant
slots $175,176,177$ --- a third distinct coefficient triple, and the last three
slots of the $178$-entry constant table.
\end{itemize}

\begin{table}[h]
\centering
\small
\begin{tabular}{ll}
\toprule
$k = i+j$ & $c_k$ \\
\midrule
$0$  & \texttt{0x1} \\
$1$  & \texttt{0x100000000000000} $= 2^{56}$ \\
$2$  & \texttt{0x10000000000000000000000000000} $= 2^{112}$ \\
$3,4,5,6$ & (no term emitted) \\
$7$  & \texttt{0x3212e00cde6d2002b119d800000347fcb8} \\
$8$  & \texttt{0x297784894e27525bc342b7fde37dba9366} \\
$9$  & \texttt{0x340f2ebe380a0f5eff4360543988a61dc2} \\
$10$ & \texttt{0x13af65741744bd7bb2c6872df2b800320} \\
$11$ & \texttt{0x2cb9b546d20373eaf85e8f53db883cb548} \\
$12$ & \texttt{0xc8557e86f90d0d89eed6eb5349a0f8820} \\
\bottomrule
\end{tabular}
\caption{Limb-product coefficients of \texttt{0x1b} sub-case 0
(\texttt{Halo2Verifier.sol} lines 2777--2878). The coefficient of $a_i a'_j$
depends only on $i+j$; the band $3 \le i+j \le 6$ is absent.}
\label{tab:qvmn:limbcoef}
\end{table}

\paragraph{The invalid-sub-index arm.} Line 3051:
\begin{lstlisting}[style=solidity,caption={\texttt{Halo2Verifier.sol} lines 3050--3052: the \texttt{0x1b} sub-dispatch default.},label={lst:qvmn:natid-default}]
                        }
                        default { q_program_fail() }
                    }
\end{lstlisting}
Any \yul{q\_native\_idx} outside $\{0,1,2,3\}$ reverts with
\revert{QuotientProgramInvalid}. This is the correct closure: the operand is a
16-bit field with $65\,536$ encodable values and only four legal ones, and a
\yul{switch} without a \yul{default} would fall through silently, consuming a
y-batch slot's worth of nothing --- except that it would not even do that, since
the fold lives inside the sub-cases. A missing \yul{default} would therefore
have dropped an identity entirely. It is present.

\paragraph{Assessment.} Unlike \texttt{0x19} and \texttt{0x1f}, whose reference
implementations are single self-contained Rust functions that can be diffed
line-by-line (and were, above), the \texttt{0x1b} bodies are generic expression
walks over \texttt{vk.cs.gates}. Their faithfulness cannot be established from
the fixture alone; it rests on the emitter and on the per-identity trace
differential described in \S\ref{sec:qvmn:assurance}. What \emph{can} be
established from the fixture, and is, is the routing: each sub-case performs
exactly one $A\cdot y$ and one bucket add, at the y-batch position
Table~\ref{tab:qvmn:stream} predicts, into a bucket index within the ten that
the prologue zeroed.

\subsection{\texttt{0x0b} \opcode{fold\_selector}}
\label{sec:qvmn:foldsel}

\paragraph{Operand encoding.} Three bytes, read as a single
\yul{shr(232, mload(q\_pc))} and then split: the high byte is the selector bucket
index, the low \texttt{u16} is the y-power gap since this bucket's previous
contribution.

\begin{lstlisting}[style=solidity,caption={\texttt{Halo2Verifier.sol} lines 3054--3088: the complete \texttt{0x0b} arm.},label={lst:qvmn:foldsel}]
                    case 0x0b {
                        // Operand layout packed into three bytes:
                        //   high byte: selector bucket index;
                        //   low u16 : y-power gap since this selector's
                        //             previous contribution.
                        let q_selector_payload := shr(232, mload(q_pc))
                        q_pc := add(q_pc, 3)
                        let q_sel_idx := shr(16, q_selector_payload)
                        let q_sel_gap := and(q_selector_payload, 0xffff)
                        // P12: the bucket index addresses the SELECTOR_ACC
                        // region and the gap indexes the y-power table; both
                        // are codegen-known sizes, so clamp before the writes.
                        if iszero(lt(q_sel_idx, 10)) { q_program_fail() }
                        if gt(q_sel_gap, 0x30) { q_program_fail() }
                        if iszero(q_has_top) { q_program_fail() }
                        let q_eval := q_top
                        q_has_top := 0
                        // Simple-selector identity: keep the same y-batch
                        // position as main identities, then advance only this
                        // selector bucket by its codegen-known gap.
                        //
                        // The global fully-evaluated accumulator is still
                        // multiplied by y so later main identities land at the
                        // same y powers as Rust's reverse fold.
                        mstore(0xb280, mulmod(mload(0xb280), y, r))
                        let q_target_ptr := add(SELECTOR_ACC_MPTR, shl(5, q_sel_idx))
                        let q_sel_acc := mload(q_target_ptr)
                        if q_sel_gap {
                            // Selector buckets are sparse in the global
                            // identity stream. Precomputed y^gap advances only
                            // this selector's local accumulator.
                            q_sel_acc := mulmod(q_sel_acc, mload(add(0xb2c0, shl(5, q_sel_gap))), r)
                        }
                        mstore(q_target_ptr, addmod(q_sel_acc, q_eval, r))
                    }
\end{lstlisting}

\paragraph{Guards.} Three, all before any state change:
\begin{enumerate}
\item \yul{iszero(lt(q\_sel\_idx, 10))} $\Rightarrow$ fail. Ten is the number of
buckets the prologue zeroed (\yul{lt(q\_sel\_zero\_off, 0x0140)}, line 2025;
$\texttt{0x0140}/32 = 10$), and those ten words
\texttt{0xb140}--\texttt{0xb280} are exactly what the linearization MSM later
consumes (lines 3932--3950). The bound is exact, and the first word past it is
the worst one available: \yul{q\_sel\_idx} is the high byte of a three-byte
operand, so it ranges over $0..255$, and index $10$ writes
\mptr{SELECTOR\_ACC\_MPTR} $+\ \texttt{0x140} =$ \texttt{0xb280} --- the
numerator accumulator $A$ itself. Indices $12..60$ walk the y-power table
\texttt{0xb2c0}--\texttt{0xb8c0}, index $61$ hits the stack base
\texttt{0xb8e0}, and index $255$ reaches \texttt{0xd120}, inside the region the
final MSM stages into.
\item \yul{gt(q\_sel\_gap, 0x30)} $\Rightarrow$ fail. \texttt{0x30} $=48$ is the
top index of the y-power table, which holds $y^0..y^{48}$; it is also exactly
$N-1$, the largest gap any bucket can legitimately have. Without it, a
\texttt{u16} gap up to $65\,535$ would read
$\mathtt{0xb2c0} + 32\cdot 65535 = \mathtt{0x20b2a0}$: never-written memory. An
\yul{mload} there does not revert --- it returns $0$ after expanding memory to
roughly two megabytes, which costs on the order of $10^{7}$ gas, so in practice
the call would run out of gas long before it mis-verified --- and a returned $0$
would zero the bucket, silently dropping every prior contribution to that
selector. Note also that the guard literal \texttt{0x30} and the table size
literal \texttt{49} at line 2044 are two independently generated numbers that
must stay in step; the guard is safe only because $\texttt{0x30} = 49 - 1$.
\item \yul{iszero(q\_has\_top)} $\Rightarrow$ fail. The opcode consumes an
identity value; an empty stack means the preceding expression did not produce
one.
\end{enumerate}
Note what is \emph{not} checked: \yul{q\_sp} is not compared against
\texttt{0xb8e0}. This is deliberate and is the case the epilogue comment at
lines 3100--3104 describes --- a fold executed with spilled operands still live
consumes only the cached top, so the residue is caught by
\yul{if iszero(eq(q\_sp, 0xb8e0))} at line 3105 (or earlier, by a subsequent
native callback's own \yul{q\_sp} check).

\paragraph{Semantics.} Write $s$ for the decoded \yul{q\_sel\_idx}, $g$ for the
decoded \yul{q\_sel\_gap} and $\varepsilon$ for the consumed \yul{q\_top}. One
execution performs
\begin{equation}
A \;\leftarrow\; A\cdot y,
\qquad
B_s \;\leftarrow\; B_s \cdot y^{\,g} \;+\; \varepsilon ,
\label{eq:qvmn:foldsel}
\end{equation}
and clears \yul{q\_has\_top}.
The \yul{if q\_sel\_gap} guard skips the multiply when $g = 0$, which is both a
gas saving and the reason slot $Y[0] = 1$ is never read by this arm --- although
the prologue initialises it anyway (line 2042, with a comment explaining that
the tail block reads its slot unconditionally). Iterating
\eqref{eq:qvmn:foldsel} over a bucket's contributions at positions
$k_1 < \dots < k_p$ with $g_i = k_i - k_{i-1}$ gives
$B_s = \sum_i y^{\,k_p - k_i}\varepsilon_{k_i}$, which the tail factor
$y^{\,N-1-k_p}$ turns into the second half of
\eqref{eq:qvmn:reverse-fold}.

\paragraph{Assessment.} This arm is clean. Every index that reaches memory is
bounds-checked against a size the prologue actually allocated, both bounds are
tight rather than generous, the stack precondition is checked in the direction
that matters for this opcode, and the residual stack-imbalance case is closed by
the epilogue. The design decision worth naming is the \emph{gap encoding}
itself: carrying $y^{\text{gap}}$ per selector, from a table precomputed once,
avoids both a runtime modular inverse of $y$ and a per-identity multiply of ten
bucket accumulators. It costs $49\cdot 32 = 1568$ bytes of scratch and $48$
\yul{mulmod}s in the prologue, against roughly $10 \times 49$ avoided
\yul{mulmod}s --- a clear win, and it keeps the selector arithmetic in the same
$y$-graded ring as $A$, which is what makes the cross-check in
\S\ref{sec:qvmn:budget} possible at all.

\subsection{The \yul{default} arm}
\label{sec:qvmn:default}

\begin{lstlisting}[style=solidity,caption={\texttt{Halo2Verifier.sol} lines 3089--3093: the opcode-dispatch default and the end of the VM loop.},label={lst:qvmn:default}]
                    // Invalid generated bytecode should fail closed. 0x1a intentionally lands here.
                    default {
                        q_program_fail()
                    }
                }
\end{lstlisting}

Any opcode byte other than the eleven in Table~\ref{tab:qvmn:census} reverts
with \revert{QuotientProgramInvalid} (selector \texttt{0x3cc81b89}). Three
properties are worth stating explicitly.

\begin{enumerate}
\item \textbf{The accepted set is a strict subset of the VM ABI.} The ABI
defines opcodes \texttt{0x01}--\texttt{0x22}; this render emits case arms only
for the eleven its own program uses, driven by \yul{program.op\_usage}. Bytes
such as \texttt{0x07} (\opcode{mul}), \texttt{0x0a} (\opcode{fold\_main}) and
\texttt{0x20} (\opcode{pow5}) are perfectly legal ABI opcodes that this
verifier will reject. Narrowing the accepted alphabet per artifact costs
nothing and shrinks the reachable state space.
\item \textbf{\texttt{0x1a} is an ABI hole.} The comment calls it out because
\texttt{0x1a} sits between \opcode{native\_permutation} (\texttt{0x19}) and
\opcode{native\_identity} (\texttt{0x1b}) and is unassigned; a reader scanning
the case labels might otherwise assume the range is dense and that a
\texttt{0x1a} case was accidentally dropped. It is not the only hole:
\texttt{0x17} and \texttt{0x18} are unassigned as well, but they fall outside
the span of rendered case labels, so nothing draws the eye to them.
\item \textbf{Fail-closed is the only correct choice here.} There is no
``skip unknown opcode'' semantics available: an unknown byte has unknown operand
width, so \yul{q\_pc} could not be advanced correctly, and continuing would
desynchronise the entire remaining instruction stream and hence the y-batch
positions of every later identity.
\end{enumerate}

Immediately after the loop (lines 3098--3105, covered elsewhere) the epilogue
enforces \yul{eq(q\_pc, q\_end)}, \yul{iszero(q\_has\_top)}, and
\yul{eq(q\_sp, 0xb8e0)}. Together with the \yul{default} arm these make
``the program ran to completion, consumed every byte, and left no operand
behind'' a checked property rather than an assumption --- with the single
exception recorded in Finding QVMN-1.

\subsection{Assurance chain for the native arms}
\label{sec:qvmn:assurance}

The generator runs a render-time self-certification pass
(\texttt{src/lowering/quotient\_numerator/vm/certify.rs}) that executes the
finalized bytecode in an independent reference interpreter at
challenge-derived random assignments and compares each identity against direct
evaluation of the expression tree it was lowered from. That pass is what makes
the \texttt{0x21} superinstruction and the run-compaction rewrites trustworthy.
It is important to be precise about its scope for \emph{this} section: as
\texttt{vm/reference.rs} states, ``Identities executed as inline Yul, native
callbacks, or the structured tail are not lowered to bytecode at all, so this
module evaluates them directly from their expression trees''. In other words,
the certifier confirms that the native callbacks appear at the right positions
in the identity stream, but it does not read the Yul in lines 2452--3050.

\begin{finding}[QVMN-4]{The native-callback Yul is outside the render-time bytecode certifier's scope --- Informational}
Lines 2452--3050 are the largest block of hand-shaped generated Yul in the
quotient phase (roughly 600 lines, three whole identity families), and they are
the block the render-time certifier explicitly does not cover
(\texttt{vm/reference.rs} lines 20--27). Their correctness rests on two other
mechanisms: (i) the per-identity Rust/Solidity trace differential
(\texttt{quotient\_numerator/native\_anchor.rs}), which compares every identity
value, the y-batched numerator, and every selector fold against
\texttt{midnight\_proofs::plonk::verifier}'s own \texttt{solidity\_trace}
events on real proofs; and (ii) the shape-corpus coverage assertions
(\texttt{shape\_corpus.rs::lookup\_chunk\_arms},
\texttt{assert\_case\_coverage}), which pin that the \texttt{Single},
\texttt{SharedPrefix} and \texttt{Staged} lookup-chunk arms and the
permutation-set counts are all exercised. \\[2pt]
Both mechanisms run at $k \le 6$ on corpus circuits, not on this fixture's
$k = 20$ shape. That is appropriate --- the identities are shape-dependent, not
size-dependent --- but it means the assurance argument for this range is
``the emitter is anchored on shapes equivalent to this one'', not ``this exact
Yul was checked''. A reviewer should confirm the corpus contains a case with
$\texttt{num\_permutation\_zs} \ge 2$ and a shared-prefix chunk of width
$\ge 3$ before treating the anchor as covering
listings~\ref{lst:qvmn:perm-d} and \ref{lst:qvmn:lookup-chunk}. Independently of
that, this document re-derived both families against
\texttt{proofs/src/plonk/permutation.rs} and
\texttt{proofs/src/plonk/logup.rs} by hand and found them faithful.
\end{finding}

\subsection{Summary of this range}
\label{sec:qvmn:summary}

\begin{itemize}
\item \textbf{Operand encodings.} \texttt{0x19}: 0 bytes. \texttt{0x1f}: 0
bytes. \texttt{0x1b}: 2 bytes, big-endian \texttt{u16} sub-index, legal values
$\{0,1,2,3\}$, everything else reverts. \texttt{0x0b}: 3 bytes, high byte
bucket index ($<10$ enforced), low \texttt{u16} gap ($\le 48$ enforced).
\item \textbf{Stack boundary.} All three native arms require
$\yul{q\_sp} = \texttt{0xb8e0}$ and reset the cached top. \texttt{0x0b} requires
a live cached top and leaves \yul{q\_sp} untouched, deferring the spilled-stack
check to the epilogue.
\item \textbf{Identity faithfulness.} The permutation arm reproduces
\texttt{permutation.rs::expressions} family-for-family and operand-for-operand;
the lookup arm reproduces \texttt{logup.rs::expressions} identically, with a
total prefix/suffix formulation of the partial products in place of the
reference's field inversion. Both generated field constants verify against \Fq{}:
$\delta$ is \texttt{PrimeField::DELTA} (odd multiplicative order
$(r-1)/2^{32}$, \emph{not} order $2^{32}$) and
\yul{q\_perm\_delta\_chunk} $= \delta^{3} \bmod r$. The four \texttt{0x1b} bodies are generic
expression walks whose routing is verified here and whose algebra is anchored
elsewhere.
\item \textbf{External surface.} None. Lines 2452--3095 contain no
\yul{call}, \yul{staticcall}, \yul{delegatecall}, \yul{calldataload},
\yul{calldatacopy}, \yul{extcodecopy}, \yul{sload}, \yul{sstore},
\yul{keccak256}, \yul{log}$n$ or \yul{mcopy}, and no precompile invocation. The
only control-flow exit other than falling through to the epilogue is
\yul{q\_program\_fail}, which reverts with the 4-byte selector
\texttt{0x3cc81b89} (\revert{QuotientProgramInvalid}) and no return data beyond
it.
\item \textbf{Fold-position consistency.} Verified exhaustively by decoding the
VK-pinned bytecode: $N = 49$ identities, native callbacks at positions
$5, 8, 26, 27$ (selector) and $29$--$47$ (main), and the ten tail exponents
$(48,47,44,41,38,35,32,29,26,20)$ agree with
$48 - k_{\mathrm{last}}(s)$ for every bucket, including the three buckets whose
contributions come from native callbacks.
\item \textbf{Findings.} QVMN-1 (Informational): native arms discard rather
than reject a live cached top. QVMN-2 (Observation): one-sided chunk clamp in
the permutation loop. QVMN-3 (Observation): total-vs-partial divergence from the
reference in the helper identity, in the safe direction. QVMN-4
(Informational): the native Yul lies outside the render-time certifier's scope.
No finding in this range is reachable from proof calldata: every branch target,
every loop bound, every memory address and every constant here is either a
Solidity literal or a byte of VK payload whose \yul{extcodehash} is re-verified
on each call (lines 1386--1389).
\end{itemize}


%% ==== 14-qvm-epilogue.tex ===============================================
\section{Quotient VM epilogue: selector tail updates and the linearization scalar}
\label{sec:qeps:main}

This section specifies \texttt{Halo2Verifier.sol} lines \textbf{3094--3402}: everything
the quotient block does after the bytecode interpreter's dispatch loop terminates, up to
(but not including) the first generated PCS sub-block at line 3414. Lines 3403--3413 are
the PCS banner comment and are specified with the PCS section. Concretely this section covers

\begin{enumerate}
  \item the three post-loop VM invariant checks (lines 3094--3105);
  \item the structured post-VM suffix that folds the \emph{trash} identity into the
        numerator accumulator (lines 3107--3306);
  \item the ten selector-bucket tail updates (lines 3307--3353);
  \item the negation of the accumulator into \mptr{LINEARIZATION\_EVAL\_MPTR}
        (lines 3355--3361);
  \item the transition block that overwrites \mptr{QUOTIENT\_MPTR} with the two scalars
        $x^{n-1}$ and $1-x^{n}$ (lines 3363--3402).
\end{enumerate}

\subsection{Entry state}
\label{sec:qeps:entry}

The interpreter loop (\texttt{for \{ \} lt(q\_pc, q\_end) \{ \}}, line 2200) has just
exited. The following runtime state is live at line 3094; everything in this table is
established by earlier sections and is cited here only as an entry condition.

\begin{table}[h]
\centering
\begin{tabularx}{\textwidth}{llX}
\toprule
Name & Value / address & Meaning \\
\midrule
\yul{y}            & \mptr{Y\_MPTR} $=$ \texttt{0x7980} & quotient batching challenge, $y \in \Fr$ \\
\yul{r}            & \texttt{FR\_MODULUS} (line 331)     & BLS12-381 scalar modulus $r$ \\
\yul{q\_pc}        & register                            & bytecode cursor \\
\yul{q\_end}       & \texttt{0x50a0 + 0x11cf} $=$ \texttt{0x626f} (line 2170) & exclusive end of the VK-resident bytecode \\
\yul{q\_sp}        & register, base \texttt{0xb8e0}      & spill-stack pointer \\
\yul{q\_has\_top}  & register                            & $1$ iff \yul{q\_top} holds a live operand \\
$A$                & word at \texttt{0xb280} $=$ \mptr{PCS\_FINAL\_MSM\_MPTR} & Horner accumulator for fully evaluated identities \\
$B_0..B_9$         & \mptr{SELECTOR\_ACC\_MPTR} $=$ \texttt{0xb140}, $10$ words & one linearization accumulator per simple selector \\
$Y[0..48]$         & \texttt{0xb2c0} $+ 32i$             & precomputed table $Y[i]=y^{i}$, $0\le i\le 48$ (lines 2042--2049) \\
\bottomrule
\end{tabularx}
\caption{Live state at the entry of the quotient VM epilogue.}
\label{tab:qeps:entry}
\end{table}

Two address facts matter throughout and are used again in
Section~\ref{sec:qeps:assessment}: the $y$-power table occupies exactly
$[\texttt{0xb2c0},\,\texttt{0xb2c0}+49\cdot 32) = [\texttt{0xb2c0},\,\texttt{0xb8e0})$,
i.e.\ it abuts the VM spill stack at \texttt{0xb8e0} without overlapping it; and the ten
selector accumulators occupy $[\texttt{0xb140},\,\texttt{0xb280})$, i.e.\ they sit
\emph{below} \mptr{PCS\_FINAL\_MSM\_MPTR} and are therefore outside the final-MSM staging
window.

\subsection{The identity schedule}
\label{sec:qeps:schedule}

Every routine in this section is parameterised by the \emph{global identity index}. The
verifier folds a fixed sequence of $N$ identities; identity $i$ carries batching weight
$y^{N-1-i}$. Identities are of two kinds:

\begin{description}
  \item[Main (fully evaluated).] The identity is a pure scalar $v_i \in \Fr$. The fold is
        \[ A \;\leftarrow\; A\cdot y + v_i . \]
  \item[Selector-targeted.] The identity is $s_j(x)\cdot u_i$ where $s_j$ is a simple
        selector \emph{column} whose commitment $\mathsf{S}_j$ is pinned in the verifying
        key. The scalar $s_j(x)$ is not available, so the verifier keeps only $u_i$ and
        moves $s_j$ to the commitment side. The fold is
        \[ A \;\leftarrow\; A\cdot y, \qquad
           B_j \;\leftarrow\; B_j\cdot y^{g} + u_i, \]
        where $g$ is the codegen-known gap to this bucket's previous contribution
        (\opcode{FOLD\_SELECTOR}, lines 3053--3088). Note that $A$ is still advanced, so
        the global $y$ ladder stays aligned for later main identities.
\end{description}

The fixture's schedule is fully determined by the deployed pair (verifier runtime plus
codehash-pinned VK payload). Decoding the VK-resident bytecode at payload offset
\texttt{0x1a20}, length \texttt{0x11cf} --- with the exact operand widths implemented by
the dispatch \yul{switch} at lines 2208--3093 --- consumes the stream and lands the cursor
precisely on \yul{q\_end}, and yields the schedule of Table~\ref{tab:qeps:schedule}. The
decode was re-run independently for this audit: replaying the eleven case arms over the VK
payload bytes ends at exactly \texttt{0x11cf} and emits, in order, the
$(\text{opcode},\text{bucket},\text{gap})$ stream recorded in the table.

\begin{table}[h]
\centering
\begin{tabular}{llll}
\toprule
Global $i$ & Emitter & Source lines & Target \\
\midrule
$0$        & inline prefix, fold \#1        & 2093--2098 & selector $0$ \\
$1$        & inline prefix, fold \#2        & 2113--2118 & selector $1$ \\
$2$        & inline prefix, fold \#3        & 2130--2135 & selector $2$ \\
$3$        & inline prefix, fold \#4        & 2147--2153 & selector $2$ (gap $1$) \\
$4$        & VM \opcode{FOLD\_SELECTOR}     & 3053--3088 & selector $2$ (gap $1$) \\
$5$        & VM \opcode{NATIVE\_IDENTITY} $0$ & 2915--2920 & selector $3$ (gap $0$) \\
$6,7$      & VM \opcode{FOLD\_SELECTOR}     & 3053--3088 & selector $3$ (gap $1$ each) \\
$8$        & VM \opcode{NATIVE\_IDENTITY} $1$ & 2966--2971 & selector $4$ (gap $0$) \\
$9,10$     & VM \opcode{FOLD\_SELECTOR}     & 3053--3088 & selector $4$ (gap $1$ each) \\
$11,12,13$ & VM \opcode{FOLD\_SELECTOR}     & 3053--3088 & selector $5$ (gaps $0,1,1$) \\
$14,15,16$ & VM \opcode{FOLD\_SELECTOR}     & 3053--3088 & selector $6$ (gaps $0,1,1$) \\
$17,18,19$ & VM \opcode{FOLD\_SELECTOR}     & 3053--3088 & selector $7$ (gaps $0,1,1$) \\
$20,21,22$ & VM \opcode{FOLD\_SELECTOR}     & 3053--3088 & selector $8$ (gaps $0,1,1$) \\
$23,24,25$ & VM \opcode{FOLD\_SELECTOR}     & 3053--3088 & selector $9$ (gaps $0,1,1$) \\
$26$       & VM \opcode{NATIVE\_IDENTITY} $2$ & 3004--3010 & selector $9$ (gap $1$) \\
$27$       & VM \opcode{NATIVE\_IDENTITY} $3$ & 3043--3049 & selector $9$ (gap $1$) \\
$28$       & VM \opcode{FOLD\_SELECTOR}     & 3053--3088 & selector $9$ (gap $1$) \\
$29\ldots41$ & VM \opcode{NATIVE\_PERMUTATION} ($13$ folds) & 2528--2559 & main \\
$42\ldots47$ & VM \opcode{NATIVE\_LOOKUP} ($6$ folds)       & 2597--2750 & main \\
$48$       & post-VM trash suffix           & 3303--3304 & main \\
\bottomrule
\end{tabular}
\caption{The complete identity schedule of the fixture, $N = 49$. Derived by decoding the
codehash-pinned VK bytecode (\texttt{Halo2VerifyingKey.sol}, payload bytes
\texttt{0x1a20}--\texttt{0x2bee} inclusive) and statically counting the folds emitted by each
inline and native block. The stream contains $21$ \opcode{FOLD\_SELECTOR}s, $4$
\opcode{NATIVE\_IDENTITY}s, one \opcode{NATIVE\_PERMUTATION} and one
\opcode{NATIVE\_LOOKUP}; the last two are the final two instructions, at program-relative
byte offsets \texttt{0x11cd} and \texttt{0x11ce}. Every other opcode in the stream is a
pure expression opcode that folds nothing. The $13$ permutation folds are $1 + 1 + 5$ (loop at line 2534) $+\, 6$
(loop at line 2542); the $6$ lookup folds are all straight-line.}
\label{tab:qeps:schedule}
\end{table}

Writing $\mathrm{Main} = \{29,\ldots,48\}$ and $S_j \subseteq \{0,\ldots,28\}$ for the
positions targeting bucket $j$, the two accumulators satisfy, immediately after the last
fold:
\begin{align}
A &= \sum_{i \in \mathrm{Main}} v_i \, y^{48-i}
   = \sum_{j=0}^{19} v_{29+j}\, y^{19-j},
   \label{eq:qeps:A}\\
B_j &= \sum_{i \in S_j} u_i \, y^{\,p_j - i},
   \qquad p_j := \max S_j .
   \label{eq:qeps:Bpre}
\end{align}
Equation~\eqref{eq:qeps:Bpre} is the \emph{bucket-local} weighting produced by the gap
scheme: the newest contribution always lands at $y^{0}$. Aligning it with the global
ladder of \eqref{eq:qeps:A} is exactly the job of the tail updates in
Section~\ref{sec:qeps:tails}.

\subsection{Post-loop invariant checks (lines 3094--3105)}
\label{sec:qeps:invariants}

\begin{lstlisting}[style=solidity,caption={\texttt{Halo2Verifier.sol} lines 3094--3105: post-loop VM invariants},label={lst:qeps:invariants}]
                // The VK-pinned bytecode must end exactly at q_end and every
                // identity must have been consumed by a fold/native callback.
                // This catches malformed generator output whose final opcode
                // over-reads operands or leaves a partial expression live.
                if iszero(eq(q_pc, q_end)) { q_program_fail() }
                if q_has_top { q_program_fail() }
                // The spilled stack must also be balanced. A FOLD executed
                // with more than one operand live consumes only the cached
                // top, leaving abandoned words below q_sp with q_has_top
                // clear -- so both checks above pass while an operand of the
                // identity has been silently dropped from nu_y(x).
                if iszero(eq(q_sp, 0xb8e0)) { q_program_fail() }
\end{lstlisting}

All three checks call \yul{q\_program\_fail()} (lines 1890--1893), whose body is
\yul{mstore(0x00, shl(224, 0x3cc81b89))} followed by \yul{revert(0x00, 0x04)}: the selector
is left-aligned in the word at offset $0$ and exactly four bytes are returned. This audit
computed the Keccak-256 digest of the ASCII signature \texttt{QuotientProgramInvalid()}
and confirms its first four bytes are \texttt{0x3cc81b89}; the error is declared at line 77 and the constant
\texttt{ERR\_QUOTIENT\_PROGRAM\_INVALID} at line 326. There is no recovery path and no state
is written before the revert.

\begin{enumerate}
  \item \textbf{$\yul{q\_pc} = \yul{q\_end}$.} The loop guard is \yul{lt(q\_pc, q\_end)},
        so an instruction that begins at $\yul{q\_end}-1$ may legally advance the cursor
        \emph{past} \yul{q\_end}. The bytecode is byte-addressed but read with
        \yul{mload}, which always pulls $32$ bytes, so a trailing instruction reads past the
        program buffer; and \opcode{MODARITH7} (line 2281), whose five one-byte counters are
        themselves read from the stream, can advance \yul{q\_pc} by up to
        $255\cdot(7\cdot 3) + 255\cdot 23 + 255\cdot 17 + 255\cdot 3 + 255\cdot 5$ bytes in a
        single instruction. Three things contain this. The program ends at
        $\texttt{0x50a0} + \texttt{0x11cf} = \texttt{0x626f}$, inside the
        VK image $[\texttt{0x3680},\,\texttt{0x7900})$ that
        \yul{extcodecopy} populated (line 1393), so a modest over-read stays in already-paid
        memory and a wild one costs only memory-expansion gas (or runs out of gas, which is
        also a revert); any operand decoded from an over-read is still subjected to the
        per-opcode pointer bound $\yul{q\_ptr} - \texttt{0x3680} \le \texttt{0x6aa0}$
        (lines 2216, 2253, 2261, 2308, 2344, 2357, 2364, 2411, 2426--2427), tightened to
        $\le \texttt{0x69e0}$ for the two \opcode{MODARITH7} pairwise bases (lines
        2384--2385); and the resulting $\yul{q\_pc} \neq \yul{q\_end}$ makes this check
        revert. The construction is therefore fail-closed rather than merely fail-fast.
  \item \textbf{$\yul{q\_has\_top} = 0$.} No expression may be left un-folded. The two
        opcodes that materialise a fresh value set the flag ---
        \opcode{PUSH\_MEM\_U16} at line 2223 and \opcode{MODARITH7} at line 2446 --- and the
        five in-place opcodes (\opcode{ADD}, \opcode{NEG}, \opcode{MUL\_CONST\_U8},
        \opcode{ADD\_MEM\_U16}, \opcode{MUL\_MEM\_U16}) leave it as they found it. Only
        \opcode{FOLD\_SELECTOR} (line 3070, guarded by \yul{if iszero(q\_has\_top)
        \{ q\_program\_fail() \}} at line 3068) and the three native-callback opcodes
        (lines 2458, 2575, 2771) clear it. A program ending mid-expression therefore leaves
        the flag set and reverts here.
  \item \textbf{$\yul{q\_sp} = \texttt{0xb8e0}$.} The spill stack must be empty. The
        source comment states the precise attack this closes: a fold consumes only the
        cached top, so a program that pushed two operands and folded once would satisfy
        checks (1) and (2) while having silently dropped a summand. The same predicate is
        also asserted \emph{before} each native callback (lines 2464, 2580, 2772), because
        the native blocks reuse \texttt{0xb8e0} onward as their own scratch table.
\end{enumerate}

\begin{finding}[qeps-1]{Post-loop invariants guard the generator, not the prover -- Informational}
The three checks in Listing~\ref{lst:qeps:invariants} cannot be provoked by proof
calldata. The bytecode lives in the VK payload, whose \yul{extcodesize} and
\yul{extcodehash} are re-verified on every call (lines 1386--1389). The precise invariant
is narrower than the source comment at line 2005 suggests, and worth stating exactly:
\emph{opcodes and immediates} come only from the VK image, so the three checked registers
\yul{q\_pc}, \yul{q\_sp} and \yul{q\_has\_top} are functions of the bytecode alone. The
\emph{values} the VM manipulates are not calldata-free --- \opcode{PUSH\_MEM\_U16} and
friends dereference pointers that legitimately land in the decoded evaluation buffer
$[\texttt{0x9480},\texttt{0xa120}]$ --- but no dereferenced word can feed back into control
flow or into these three registers. They are integrity guards against a malformed generator or a mismatched
deployment, cost a few tens of gas in total, and fail closed. This is cheap and correct
defence in depth. Note, however, what they do \emph{not} check: a program that folds the
wrong \emph{number} of identities in a self-consistent way passes all three. See
finding~\ref{fnd:qeps-3}.
\label{fnd:qeps-1}
\end{finding}

\subsection{The post-VM structured suffix: the trash identity (lines 3107--3306)}
\label{sec:qeps:trash}

After the interpreter, the generator emits one structured Yul block that participates in
the same identity order and writes into the same accumulators. In this fixture that block
holds the single \emph{trash} identity, global index $48$ --- the last identity, hence
batching weight $y^{0}$.

\subsubsection{Operand map}

Every operand is an \yul{mload} of the decoded evaluation buffer
$\mptr{REVERSED\_EVALS\_MPTR} = \texttt{0x9480}$, which holds
$\mathbf{e} = (e_0,\dots,e_{101})$ --- $102$ words, canonicality-checked and transcript-
absorbed by the loop at lines 1697--1714 (\texttt{0x0cc0} calldata bytes $= 102\cdot 32$).
The slot at address $a$ is $e_k$ with $k = (a - \texttt{0x9480})/32$.

\begin{table}[h]
\centering
\begin{tabular}{lllllll}
\toprule
Addr & $e_k$ & Name & \phantom{xx} & Addr & $e_k$ & Name \\
\midrule
\texttt{0x94a0} & $e_1$    & \yul{a\_0}          && \texttt{0x99c0} & $e_{42}$ & \yul{f\_4} \\
\texttt{0x94c0} & $e_2$    & \yul{a\_1}          && \texttt{0x99e0} & $e_{43}$ & \yul{f\_5} \\
\texttt{0x94e0} & $e_3$    & \yul{a\_2}          && \texttt{0x9a00} & $e_{44}$ & \yul{f\_6} \\
\texttt{0x9500} & $e_4$    & \yul{a\_3}          && \texttt{0x9a20} & $e_{45}$ & \yul{f\_7} \\
\texttt{0x9520} & $e_5$    & \yul{a\_4}          && \texttt{0x9a60} & $e_{47}$ & \yul{f\_0} \\
\texttt{0x9540} & $e_6$    & \yul{a\_0\_next\_1} && \texttt{0x9a80} & $e_{48}$ & \yul{f\_1} \\
\texttt{0x9560} & $e_7$    & \yul{a\_1\_next\_1} && \texttt{0x9aa0} & $e_{49}$ & \yul{f\_2} \\
\texttt{0x9580} & $e_8$    & \yul{a\_2\_next\_1} && \texttt{0x9ac0} & $e_{50}$ & \yul{f\_3} \\
\texttt{0x95a0} & $e_9$    & \yul{a\_5}          && \texttt{0x9ba0} & $e_{57}$ & \yul{f\_26} \\
\texttt{0x95c0} & $e_{10}$ & \yul{a\_6}          && \texttt{0xa120} & $e_{101}$ & trash eval \\
\texttt{0x95e0} & $e_{11}$ & \yul{a\_7}          &&                 &          & \\
\bottomrule
\end{tabular}
\caption{Evaluation slots read by the trash suffix. \yul{f\_8} $=e_{46}$ at
\texttt{0x9a40} is used by the inline prefix but not here. The generated \yul{f\_j} names
are fixed-column indices; the memory order is query order, hence the non-monotone pairing.}
\label{tab:qeps:evals}
\end{table}

Three further facts fix the semantics of the block:
\begin{itemize}
  \item \yul{var0} $=$ \texttt{0x73eda753299d7d483339d80809a1d80553bda402fffe5bfeffffffff00000000}
        $= r - 1 \equiv -1 \pmod r$; every \yul{mulmod(v, var0, r)} is a negation.
  \item \yul{q\_pow5(x)} (lines 1911--1915) computes $x^{5}$ as
        $x\cdot(x^{2})^{2}$ with three \yul{mulmod}s.
  \item \yul{q\_trash\_tau} $=$ \yul{mload(TRASH\_CHALLENGE\_MPTR)} $=$ the challenge
        $\tau$ squeezed at line 1646, between $\gamma$ and $y$.
\end{itemize}

\subsubsection{Structure}

\begin{lstlisting}[style=solidity,caption={\texttt{Halo2Verifier.sol} lines 3114--3135 (head of the trash suffix; lines 3136--3298, holding the rest of $t_0$ and rows $t_1..t_7$ with their Horner steps, are elided)},label={lst:qeps:trashhead}]
                {
                let q_trash_tau := mload(TRASH_CHALLENGE_MPTR)
                {
                let f_0 := mload(0x9a60)
                let a_0_next_1 := mload(0x9540)
                let var0 := 0x73eda753299d7d483339d80809a1d80553bda402fffe5bfeffffffff00000000
                let var1 := mulmod(a_0_next_1, var0, r)
                let var2 := addmod(f_0, var1, r)
                let var3 := 0x590ba402032e82eb1f660ef09796c5686345a5054ed96dae8e2d233633788771
                let a_0 := mload(0x94a0)
                let var4 := mulmod(var3, a_0, r)
                let var5 := addmod(var2, var4, r)
                let var6 := 0x52f789e4afc3801f7411102ee2f47cc5954a744e71cac98e75ea962a55a0a76f
                let a_1 := mload(0x94c0)
                let var7 := mulmod(var6, a_1, r)
                let var8 := addmod(var5, var7, r)
                let var9 := 0x3509dd2fe3aac0080783557fec090fb1cb4b2b0901253c55282024331d1fe1a8
                let a_2 := mload(0x94e0)
                let var10 := q_pow5(a_2)
                let var11 := mulmod(var9, var10, r)
                let var12 := addmod(var8, var11, r)
                let var13 := 0x333f8046ece5579cbd6872449c57f2703dfc8864cfadc06d587ff104a0d0c1f2
\end{lstlisting}

The block computes eight row polynomials $t_0,\dots,t_7$ (Yul temporaries
\yul{var32}, \yul{var59}, \yul{var71}, \yul{var85}, \yul{var101}, \yul{var119},
\yul{var139}, \yul{var161}) and batches them by $\tau$ with an explicit Horner chain
(lines 3160, 3189, 3202, 3217, 3234, 3253, 3274, 3298):
\begin{align}
H_0 &= 0\cdot\tau + t_0, \nonumber\\
H_{m} &= H_{m-1}\cdot\tau + t_{m}, \qquad 1 \le m \le 7,
  \label{eq:qeps:horner}
\end{align}
so that \yul{var162} $= H_7 = \sum_{m=0}^{7} t_m\,\tau^{\,7-m}$. Note the seed
\yul{addmod(mulmod(0, q\_trash\_tau, r), var32, r)} at line 3160: the generator emits the
degenerate first Horner step verbatim rather than special-casing it.

Each row has the shape
\begin{equation}
t_m \;=\; f_{\sigma(m)} \;-\; z_m \;+\; \alpha_m\, a_0 \;+\; \beta_m\, a_1
        \;+\; \sum_{c} \gamma_{m,c}\, a_c^{5},
\label{eq:qeps:row}
\end{equation}
with $\sigma$, $z_m$ and the coefficient sets given in Table~\ref{tab:qeps:trashrows}.
Rows $0$ and $1$ carry six independent $a_c^{5}$ coefficients each ($c=2,\dots,7$). The
$a_c^{5}$ coefficients of rows $2\ldots7$ are instead \emph{Toeplitz} and lower-triangular:
the generated code reuses six constants $D_0,\dots,D_5$ so that the coefficient of
$a_c^{5}$ in row $m$ is
\[ \gamma_{m,c} = D_{\,m-c} \quad\text{for } 2 \le c \le m, \qquad
   \gamma_{m,c} = 0 \quad\text{for } c > m \;(m \ge 2), \]
so row $m$ contributes $m-1$ such terms --- one in row $2$, six in row $7$. The generator
emits no multiplication for a zero coefficient, which is why the rows have different
lengths. The shape is consistent with a Poseidon partial-round / round-skip gate, although
this specification does not cross-check the constants against the circuit's Poseidon
parameters.

\begin{table}[h]
\centering
\begin{tabularx}{\textwidth}{llllX}
\toprule
$m$ & $f_{\sigma(m)}$ & $z_m$ & $\alpha_m,\ \beta_m$ & $a_c^{5}$ coefficients \\
\midrule
$0$ & \yul{f\_0}\,$=e_{47}$ & $e_6$    & \texttt{0x590ba402..788771}, \texttt{0x52f789e4..a0a76f}
    & $a_2^5$:\,\texttt{0x3509dd2f..1fe1a8}; $a_3^5$:\,\texttt{0x333f8046..d0c1f2};
      $a_4^5$:\,\texttt{0x412c9823..a60ca5}; $a_5^5$:\,\texttt{0x53fded36..c9790e};
      $a_6^5$:\,\texttt{0x6ccb1c7d..dd6d34}; $a_7^5$:\,\texttt{0x3f05c4df..e1edc1} \\
\addlinespace
$1$ & \yul{f\_1}\,$=e_{48}$ & $e_7$    & \texttt{0x5b1fc262..0f1a5d}, \texttt{0x4d0ea7f9..358028}
    & $a_2^5$:\,\texttt{0x26cc223e..18eb10}; $a_3^5$:\,\texttt{0x31e823a4..838d25};
      $a_4^5$:\,\texttt{0x275a2036..69ba85}; $a_5^5$:\,\texttt{0x5f3a15ba..33f701};
      $a_6^5$:\,\texttt{0x301cf56f..fb2ecd}; $a_7^5$:\,\texttt{0x0fdf664d..7f1a1f} \\
\addlinespace
$2$ & \yul{f\_2}\,$=e_{49}$ & $e_4$    & \texttt{0x5e1d3dbe..c49140}, \texttt{0x6bd72f9c..5aa78a} & $a_2^5$:\,$D_0$ \\
$3$ & \yul{f\_3}\,$=e_{50}$ & $e_5$    & \texttt{0x222e83e7..7934c1}, \texttt{0x26c2cc87..d07191} & $a_2^5$:\,$D_1$; $a_3^5$:\,$D_0$ \\
$4$ & \yul{f\_4}\,$=e_{42}$ & $e_9$    & \texttt{0x726df150..202e4b}, \texttt{0x24822e1a..f05520} & $a_2^5$:\,$D_2$; $a_3^5$:\,$D_1$; $a_4^5$:\,$D_0$ \\
$5$ & \yul{f\_5}\,$=e_{43}$ & $e_{10}$ & \texttt{0x2f5908b1..7751a7}, \texttt{0x23a6684b..c6e827} & $a_2^5$:\,$D_3$; $a_3^5$:\,$D_2$; $a_4^5$:\,$D_1$; $a_5^5$:\,$D_0$ \\
$6$ & \yul{f\_6}\,$=e_{44}$ & $e_{11}$ & \texttt{0x6d05a419..c7274a}, \texttt{0x27e71192..340849} & $a_2^5$:\,$D_4$; $a_3^5$:\,$D_3$; $a_4^5$:\,$D_2$; $a_5^5$:\,$D_1$; $a_6^5$:\,$D_0$ \\
$7$ & \yul{f\_7}\,$=e_{45}$ & $e_8$    & \texttt{0x70d8f2a7..386633}, \texttt{0x40fa389f..560261} & $a_2^5$:\,$D_5$; $a_3^5$:\,$D_4$; $a_4^5$:\,$D_3$; $a_5^5$:\,$D_2$; $a_6^5$:\,$D_1$; $a_7^5$:\,$D_0$ \\
\bottomrule
\end{tabularx}
\caption{The eight trash rows of \eqref{eq:qeps:row}, lines 3117--3298. The shared
Toeplitz constants are
$D_0=\texttt{0x4997c5aa..9de0db}$ (\yul{var69}),
$D_1=\texttt{0x4382d093..780272}$ (\yul{var81}),
$D_2=\texttt{0x4e528010..b1ea29}$ (\yul{var95}),
$D_3=\texttt{0x1981b4b3..f98d4e}$ (\yul{var111}),
$D_4=\texttt{0x00d94c46..2f0533}$ (\yul{var129}),
$D_5=\texttt{0x1f61345b..c50d3e}$ (\yul{var149}). Values here are shown zero-padded to
$32$ bytes; Yul literals drop leading zero nibbles, so \yul{var129} (as well as
\yul{var25} and \yul{var57} in rows $0$ and $1$) appears shorter in the source. Hex values
are abbreviated; the full literals appear verbatim in the source at the cited lines.}
\label{tab:qeps:trashrows}
\end{table}

\subsubsection{Trash gating and the fold (lines 3299--3304)}

\begin{lstlisting}[style=solidity,caption={\texttt{Halo2Verifier.sol} lines 3299--3306: trash gating and fold},label={lst:qeps:trashfold}]
                let f_26 := mload(0x9ba0)
                let q_trash_one_minus_selector := addmod(1, sub(r, f_26), r)
                let q_trash_scaled := mulmod(q_trash_one_minus_selector, mload(0xa120), r)
                let q_trash_eval := addmod(var162, sub(r, q_trash_scaled), r)
                mstore(0xb280, mulmod(mload(0xb280), y, r))
                mstore(0xb280, addmod(mload(0xb280), q_trash_eval, r))
                }
                }
\end{lstlisting}

Writing $\mathrm{sel}_{26} = e_{57}$ for the trash-selector evaluation and
$T = e_{101}$ for the trash-polynomial evaluation, the identity folded at global index
$48$ is
\begin{equation}
v_{48} \;=\; \Big(\textstyle\sum_{m=0}^{7} t_m\,\tau^{\,7-m}\Big)
             \;-\; \big(1 - \mathrm{sel}_{26}\big)\cdot T .
\label{eq:qeps:trash}
\end{equation}
Both subtractions use the canonical pattern \yul{addmod(u, sub(r, v), r)}, which is exact
for all $u,v \in [0,r)$ including $v=0$ (then $\yul{sub(r,0)}=r$ and
$\yul{addmod(u,r,r)}=u$). The inputs $e_{57}$ and $e_{101}$ were range-checked
$< r$ at line 1706, so no operand of \yul{sub} can underflow.

\begin{finding}[qeps-2]{The prover-supplied trash evaluation is bound to a commitment -- Informational}
Equation~\eqref{eq:qeps:trash} subtracts $T = e_{101}$, a scalar the prover chooses. On its
own that would let a prover absorb any residual. It is sound here because $e_{101}$ is not
free: it is transcript-absorbed before $x_1,\dots,x_4$ are squeezed (line 1712), and
\texttt{0xa120} is registered as an evaluation pointer in the PCS point-set table at line
3469, so the batched Horner at line 3511 forces $e_{101}$ to be the opening of the
corresponding committed polynomial at $x$. The gating factor $1 - \mathrm{sel}_{26}$ is
itself an evaluation of a VK-committed fixed column ($e_{57}$), likewise bound. This region
is clean; the binding argument lives outside this section (lines 3469, 3486, 3511) and is
specified there.
\label{fnd:qeps-2}
\end{finding}

\subsection{Selector tail updates (lines 3307--3353)}
\label{sec:qeps:tails}

\begin{lstlisting}[style=solidity,caption={\texttt{Halo2Verifier.sol} lines 3307--3326: the first three of ten tail updates (buckets $3..9$ at lines 3326--3353 follow the identical shape; the trailing brace is the opening of block $3$)},label={lst:qeps:tails}]
                // Finish selector buckets by applying the codegen-known tail
                // from each selector's last identity to the end of the global
                // y-batch.
                //
                // After this step, every selector bucket is aligned with the
                // final global y position and can be multiplied by its fixed
                // selector commitment in the linearized MSM.
                {
                    let q_sel_ptr := add(SELECTOR_ACC_MPTR, 0x00)
                    mstore(q_sel_ptr, mulmod(mload(q_sel_ptr), mload(add(0xb2c0, 0x0600)), r))
                }
                {
                    let q_sel_ptr := add(SELECTOR_ACC_MPTR, 0x20)
                    mstore(q_sel_ptr, mulmod(mload(q_sel_ptr), mload(add(0xb2c0, 0x05e0)), r))
                }
                {
                    let q_sel_ptr := add(SELECTOR_ACC_MPTR, 0x40)
                    mstore(q_sel_ptr, mulmod(mload(q_sel_ptr), mload(add(0xb2c0, 0x0580)), r))
                }
                {
\end{lstlisting}

Ten unconditional, fully unrolled blocks. Block $j$ performs
\begin{equation}
B_j \;\leftarrow\; B_j \cdot Y[T_j], \qquad Y[i] = y^{i} \text{ at } \texttt{0xb2c0}+32i,
\label{eq:qeps:tailupdate}
\end{equation}
with the pair (bucket byte offset, $y$-power byte offset) hard-coded as an immediate.
There is no loop, no bound check, and no data dependence on the proof: each block is two
\yul{mload}s, one \yul{mulmod} and one \yul{mstore} on compile-time constant addresses.
\yul{SELECTOR\_ACC\_MPTR} and the base \texttt{0xb2c0} are constants, so \yul{solc} folds
both \yul{add}s away.

\begin{table}[h]
\centering
\begin{tabular}{llllll}
\toprule
$j$ & \mptr{SELECTOR\_ACC\_MPTR}$+$ & $y$-power offset & $T_j$ & $p_j$ (last position) & $p_j+T_j$ \\
\midrule
$0$ & \texttt{0x00}   & \texttt{0x0600} & $48$ & $0$  & $48$ \\
$1$ & \texttt{0x20}   & \texttt{0x05e0} & $47$ & $1$  & $48$ \\
$2$ & \texttt{0x40}   & \texttt{0x0580} & $44$ & $4$  & $48$ \\
$3$ & \texttt{0x60}   & \texttt{0x0520} & $41$ & $7$  & $48$ \\
$4$ & \texttt{0x80}   & \texttt{0x04c0} & $38$ & $10$ & $48$ \\
$5$ & \texttt{0xa0}   & \texttt{0x0460} & $35$ & $13$ & $48$ \\
$6$ & \texttt{0xc0}   & \texttt{0x0400} & $32$ & $16$ & $48$ \\
$7$ & \texttt{0xe0}   & \texttt{0x03a0} & $29$ & $19$ & $48$ \\
$8$ & \texttt{0x0100} & \texttt{0x0340} & $26$ & $22$ & $48$ \\
$9$ & \texttt{0x0120} & \texttt{0x0280} & $20$ & $28$ & $48$ \\
\bottomrule
\end{tabular}
\caption{The ten tail updates of lines 3314--3353. $T_j$ is the power offset divided by
$32$. $p_j$ is taken from Table~\ref{tab:qeps:schedule}. The last column is constant,
which is precisely the codegen contract $T_j = N - 1 - p_j$ with $N = 49$.}
\label{tab:qeps:tails}
\end{table}

Substituting \eqref{eq:qeps:tailupdate} into \eqref{eq:qeps:Bpre} and using
$T_j = 48 - p_j$ gives the \textbf{final bucket scalar}
\begin{equation}
\boxed{\;
B_j \;=\; \sum_{i \in S_j} u_i \, y^{\,48 - i}\;}
\label{eq:qeps:Bfinal}
\end{equation}
which is exactly the global weighting of \eqref{eq:qeps:A}. Explicitly, from
Table~\ref{tab:qeps:schedule}:
\[
\begin{aligned}
B_0 &= u_0\,y^{48}, &
B_1 &= u_1\,y^{47}, \\
B_2 &= u_2 y^{46} + u_3 y^{45} + u_4 y^{44}, &
B_3 &= u_5 y^{43} + u_6 y^{42} + u_7 y^{41}, \\
B_4 &= u_8 y^{40} + u_9 y^{39} + u_{10} y^{38}, &
B_5 &= u_{11} y^{37} + u_{12} y^{36} + u_{13} y^{35}, \\
B_6 &= u_{14} y^{34} + u_{15} y^{33} + u_{16} y^{32}, &
B_7 &= u_{17} y^{31} + u_{18} y^{30} + u_{19} y^{29}, \\
B_8 &= u_{20} y^{28} + u_{21} y^{27} + u_{22} y^{26}, &
B_9 &= \textstyle\sum_{i=23}^{28} u_i\,y^{\,48-i}.
\end{aligned}
\]

All ten tails are non-zero, so none is dropped: the generator
(\texttt{src/lowering/quotient.rs:238--250}) filters out zero tails, and a zero tail can
only arise for a bucket whose last contribution is the very last identity --- impossible
here because identity $48$ is the main-targeted trash identity.

The largest tail is $T_0 = 48$, and the largest $y$-power table index is $48$
(loop bound \yul{lt(q\_y\_power\_i, 49)}, line 2044). The tail read for bucket $0$ is
therefore \yul{mload(0xb8c0)}, the last word of the table, one word below the spill-stack
base \texttt{0xb8e0}. Every tail read is in bounds; the same bound $48 =$ \texttt{0x30} is
the runtime clamp \yul{if gt(q\_sel\_gap, 0x30)} applied to \opcode{FOLD\_SELECTOR} gaps
at line 3067.

\begin{finding}[qeps-3]{Tail exponents are unvalidated compile-time immediates -- Informational}
The ten $y$-power offsets at lines 3316--3352 are literal immediates in the verifier
runtime, whereas the identity schedule that determines them is split across the verifier
runtime (inline prefix, native callbacks) and the \emph{VK} payload (the interpreted
bytecode). Unlike the \opcode{FOLD\_SELECTOR} gap, which is clamped to
$\le \texttt{0x30}$ at line 3067 and whose bucket index is clamped to $< 10$ at line 3066,
the tails receive no runtime bound check and no runtime cross-check against the number of
folds actually performed. A tail of, say, $60$ would read \texttt{0xba40}, inside the
native-callback scratch area, and silently corrupt a bucket scalar.

Two things contain this. First, the verifier re-verifies \yul{extcodesize} and
\yul{extcodehash} of the VK on every call (lines 1386--1389), so verifier and VK cannot be
mismatched at runtime; the pair is frozen at deployment. Second, the constants are
statically auditable, and this specification audits them:
Table~\ref{tab:qeps:tails} was produced by decoding the pinned VK bytecode and statically
counting folds --- and that decode was reproduced independently during the fidelity audit
of this section --- and all ten satisfy $p_j + T_j = 48 = N-1$ exactly. The region is
therefore correct as deployed. The finding is recorded because the correctness argument is
entirely external to the contract --- nothing in the runtime would reject a generator that
got these ten numbers wrong.
\label{fnd:qeps-3}
\end{finding}

\begin{finding}[qeps-4]{Bucket accumulators alias the batch-inversion scratch and neighbour the MSM staging window -- Observation}
\mptr{SELECTOR\_ACC\_MPTR} and \mptr{BATCH\_INV\_SCRATCH\_MPTR} are both \texttt{0xb140}
(lines 229 and 235); the accumulator $A$ sits at \texttt{0xb280} $=$
\mptr{PCS\_FINAL\_MSM\_MPTR}; and the $y$-power table at \texttt{0xb2c0}--\texttt{0xb8e0}
lies \emph{inside} the region the PCS block later reuses. Three live ranges therefore
overlap in address space. Each separation checks out:
(i) the Lagrange batch inversion completes at line 1826, before $A$ is zeroed at line 2021,
before the selector buckets are zeroed at line 2027 and before the $y$-power table is
written at lines 2042--2049 --- which matters because its prefix-product scratch is $28$
words wide ($\texttt{0xb140}$--$\texttt{0xb4c0}$, forward pass at lines 886--903) and so
transiently covers $A$'s slot and the low half of the $y$-power table, not only the ten
buckets (the source names the pinning test
\yul{selector\_accumulators\_are\_live\_from\_quotient\_to\_final\_msm} at line 234;
it lives at \texttt{src/lowering/layout/memory.rs:1857});
(ii) $A$ is consumed at line 3358, and \texttt{0xb280} is first overwritten at line 3463 ---
\emph{not} at the final-MSM \yul{mcopy} of line 3836, because the PCS eval-source table of
lines 3463--3505 already occupies $[\texttt{0xb280},\texttt{0xb7e0})$;
(iii) the $y$-power table is dead after line 3353, and the earliest write into it is line
3465 (\yul{mstore(0xb2c0, 0xa020)}), again part of that eval-source table.
Every \yul{mstore} and \yul{mcopy} between lines 3362 and 4200 targets either
$\ge \texttt{0xb280}$ or the disjoint return window at \texttt{0x1000}; none lands in
$[\texttt{0xb140},\texttt{0xb280})$, so the ten bucket scalars survive intact until they
are read at lines 3932--3950. The layout is sound but its soundness is purely a scheduling
argument; it is worth keeping the pinning test.
\label{fnd:qeps-4}
\end{finding}

\subsection{The linearization scalar (lines 3355--3361)}
\label{sec:qeps:lineval}

\begin{lstlisting}[style=solidity,caption={\texttt{Halo2Verifier.sol} lines 3355--3361: negation into \mptr{LINEARIZATION\_EVAL\_MPTR}},label={lst:qeps:lineval}]
                // Fully evaluated identities are the constant-polynomial side
                // of the linearization query. Rust subtracts that grouped
                // scalar into expected_eval, so Solidity stores -nu_y(x).
                let linearization_expected_eval := addmod(0, sub(r, mload(0xb280)), r)
                mstore(LINEARIZATION_EVAL_MPTR, linearization_expected_eval)
                pop(y)
            }
\end{lstlisting}

\begin{equation}
\mptr{LINEARIZATION\_EVAL\_MPTR} \;\leftarrow\; (-A) \bmod r,
\qquad A = \sum_{i=29}^{48} v_i\,y^{48-i}.
\label{eq:qeps:lin}
\end{equation}

The negation is \yul{addmod(0, sub(r, A), r)}. For $A = 0$ this evaluates
$\yul{sub}(r,0) = r$ and $\yul{addmod}(0,r,r) = 0$, so the canonical representative is
produced without a special case; for $A \in [1, r)$ it yields $r - A$. $A$ is a
\yul{mulmod}/\yul{addmod} result and hence always $< r$, so \yul{sub} cannot underflow.

\yul{pop(y)} at line 3360 discards the challenge. It is a no-op emitted unconditionally by
the template (\texttt{templates/partials/quotient\_numerator/QuotientNumeratorBlock.yul:420})
so that artifacts whose fold set is empty do not raise an unused-variable warning; here
\yul{y} is used many times. Yul's optimiser removes \yul{pop} of a plain variable, so the
runtime cost is zero.

The enclosing brace at line 3361 closes the whole quotient scope opened at line 1993. Only
two values escape it: \mptr{LINEARIZATION\_EVAL\_MPTR} and the ten words at
\mptr{SELECTOR\_ACC\_MPTR}.

\subsubsection{What $A$ is, precisely}

The comment at lines 2017--2019 calls $A$ ``\yul{nu\_y(x)} for the \yul{None} identity
group'', and that qualification is load-bearing. $A$ is \emph{not} the full batched
numerator: the $29$ selector-targeted identities contribute $0$ to it. Writing
$\nu_y(x)$ for the complete batched numerator,
\begin{equation}
\nu_y(x) \;=\; \underbrace{\sum_{i \in \mathrm{Main}} v_i\,y^{48-i}}_{A}
             \;+\; \sum_{j=0}^{9} \Big(\sum_{i \in S_j} u_i\,y^{48-i}\Big)\, s_j(x)
      \;=\; A \;+\; \sum_{j=0}^{9} B_j\, s_j(x),
\label{eq:qeps:nusplit}
\end{equation}
using \eqref{eq:qeps:Bfinal}. The second sum is never formed as a scalar --- the verifier
has no $s_j(x)$ --- and instead reappears on the commitment side as
$\sum_j B_j\,\mathsf{S}_j$.

\subsection{Transition block: \texorpdfstring{$x^{n-1}$}{x\^{}(n-1)} and \texorpdfstring{$1-x^{n}$}{1-x\^{}n} (lines 3363--3402)}
\label{sec:qeps:transition}

\begin{lstlisting}[style=solidity,caption={\texttt{Halo2Verifier.sol} lines 3379--3402: linearization scalar metadata},label={lst:qeps:transition}]
            {
                let x := mload(X_MPTR)
                let k := 20
                // Compute both x^n and x^(n-1) with the same squaring walk:
                // x_pow_2i tracks x^(2^i), while x_pow_2i_minus1 tracks
                // x^(2^i - 1).
                let x_pow_2i := x
                let x_pow_2i_minus1 := 1
                for { let idx := 0 } lt(idx, k) { idx := add(idx, 1) } {
                    x_pow_2i_minus1 := mulmod(
                        mulmod(x_pow_2i_minus1, x_pow_2i_minus1, r),
                        x,
                        r
                    )
                    x_pow_2i := mulmod(x_pow_2i, x_pow_2i, r)
                }
                let x_split := x_pow_2i_minus1
                let one_minus_x_n := addmod(1, sub(r, x_pow_2i), r)

                // PCS block 5 interprets this 2-word payload as scalar
                // metadata, not as a materialized G1 point.
                mstore(QUOTIENT_MPTR, x_split)
                mstore(add(QUOTIENT_MPTR, 0x20), one_minus_x_n)
            }
\end{lstlisting}

\begin{algorithm}[h]
\caption{Combined $x^{n}$ / $x^{n-1}$ squaring walk, $n = 2^{k}$, $k=20$}
\label{alg:qeps:powwalk}
\begin{algorithmic}[1]
\State $P \gets x$; $Q \gets 1$
\For{$\mathrm{idx} = 0$ to $k-1$}
  \State $Q \gets Q^{2}\cdot x \bmod r$
  \State $P \gets P^{2} \bmod r$
\EndFor
\State \Return $x_{\mathrm{split}} \gets Q$, $\;1 - x^{n} \gets (1 - P) \bmod r$
\end{algorithmic}
\end{algorithm}

\paragraph{Correctness.} The loop invariant at the top of iteration $i$ is
$P = x^{2^{i}}$ and $Q = x^{2^{i}-1}$. It holds at $i=0$ ($P=x^{1}$, $Q=x^{0}=1$) and is
preserved: $Q^{2}x = x^{2(2^{i}-1)+1} = x^{2^{i+1}-1}$ and $P^{2} = x^{2^{i+1}}$. After
$k=20$ iterations,
\begin{equation}
x_{\mathrm{split}} = x^{\,n-1}, \qquad \yul{one\_minus\_x\_n} = 1 - x^{n} \bmod r,
\qquad n = 2^{20} = 1\,048\,576 .
\label{eq:qeps:xpow}
\end{equation}
Cost: $3k = 60$ \yul{mulmod}s. \yul{sub(r, x\_pow\_2i)} is safe because $x < r$
(\yul{squeeze\_to} reduces the digest modulo $r$, line 807) and \yul{mulmod} outputs are
reduced; the $x^{n}=0$ edge case degenerates correctly to $1$ as in
Section~\ref{sec:qeps:lineval}.

\paragraph{Payload reinterpretation.} \mptr{QUOTIENT\_MPTR} $=$ \texttt{0x7d20} is declared
as a four-word EIP-2537 padded $\GO$ slot (line 182). On this code path it is never a point:
the block overwrites word $0$ with $x^{n-1}$ and word $1$ with $1-x^{n}$ and leaves words
$2$ and $3$ untouched. No read of \texttt{0x7d60} or \texttt{0x7d80} exists anywhere in the
contract, so the stale/zero tail is inert. The consumer is the fused final MSM at lines
3816--3817, which loads exactly the two words back.

\begin{finding}[qeps-5]{$x^{n}$ is recomputed although \mptr{X\_N\_MPTR} already holds it -- Observation}
The Lagrange block computes $x^{n}$ with the identical $k=20$ squaring walk and persists it
at \mptr{X\_N\_MPTR} (lines 1800--1806, persisted at 1869). Lines 3385--3394 recompute it. The
redundancy is $20$ \yul{mulmod}s, roughly $160$ gas, and it is not free of coupling risk:
two independent constants $k=20$ must stay in step. In mitigation, both are cross-checked
against the VK header, which is asserted to carry $k = 20$ at line 1401, and the shared
walk genuinely does buy $x^{n-1}$ for the price of one extra multiply per iteration, which
is much cheaper than deriving $x^{n-1} = x^{n}\cdot x^{-1}$ (a modexp inversion). Reading
$P$'s final value from \mptr{X\_N\_MPTR} instead of recomputing it would save the $20$
squarings while keeping the $Q$ recurrence, at the cost of one more cross-block dependency.
No correctness impact.
\label{fnd:qeps-5}
\end{finding}

\subsection{The resulting linearization claim}
\label{sec:qeps:claim}

The two outputs of this section are consumed by the fused final MSM at lines 3816--3950
(specified elsewhere). Write
\[ w \;:=\; \yul{mload(add(X1\_POWERS\_MPTR, 0x540))} \]
for the weight the PCS assigns to the linearization query. This is entry $42$ of the
$43$-entry $x_1$-power table built at lines 3444--3455; note that entries $1\ldots42$ are
\emph{truncated} to their low $128$ bits by \yul{and(acc, 0xffffffffffffffffffffffffffffffff)}
at line 3453, so $w$ is the low half of $x_1^{42}$, not $x_1^{42}$ itself. What matters here
is only that the identical word is used on both sides: the commitment side reads it into
\yul{lin\_query\_scalar\_41} at line 3918, and the evaluation side pairs
\mptr{LINEARIZATION\_EVAL\_MPTR} --- stored at \texttt{0xb7c0} $=$ \texttt{0xb280} $+$
\texttt{0x540} at line 3505 --- with the same table offset in the Horner accumulation of
lines 3506--3514: the sum is seeded from term $0$ at line 3506 and the loop at line 3509
(bound \yul{lt(i, 0x2b)}, so $i = 1,\dots,42$) advances \yul{eval\_p} from
$\texttt{0xb280}+\texttt{0x20}$ and \yul{pow\_p} from
$\mptr{X1\_POWERS\_MPTR}+\texttt{0x20}$ in lockstep, pairing \texttt{0xb7c0} with
$\mptr{X1\_POWERS\_MPTR}+\texttt{0x540}$ at $i=42$. Reading the two outputs back:

\begin{align}
\text{quotient terms:}\quad
&\text{scalar}_i = w\,(1-x^{n})\,(x^{\,n-1})^{i},
 && \text{point } \mathsf{Q}_i,\quad i=0,\dots,3 && \text{(lines 3918--3930)}
 \label{eq:qeps:msmq}\\
\text{selector terms:}\quad
&\text{scalar}_j = w\, B_j,
 && \text{point } \mathsf{S}_j,\quad j=0,\dots,9 && \text{(lines 3931--3950)}
 \label{eq:qeps:msms}\\
\text{eval term:}\quad
&\text{weight } w\cdot(-A)
 && && \text{(lines 3505, 3511)}
 \label{eq:qeps:msme}
\end{align}

That is, the linearized commitment
\begin{equation}
\mathsf{L} \;=\; (1-x^{n})\sum_{i=0}^{3} \big(x^{\,n-1}\big)^{i}\,\mathsf{Q}_i
            \;+\; \sum_{j=0}^{9} B_j\,\mathsf{S}_j
\label{eq:qeps:L}
\end{equation}
is entered into the KZG multi-open batch as the last member of point set $0$ (the
$43$rd of its $43$ evaluation terms, the $42$nd of its $42$ commitment terms, per the
generated comment at line 3460) with claimed opening value $-A$ at $x$. Because $w$ divides
out of both sides of the resulting equation, it does not appear in
\eqref{eq:qeps:assert}. Writing $h(X) = \sum_{i=0}^{3} X^{(n-1)i}\,h_i(X)$ for the
polynomial the four quotient limb commitments $\mathsf{Q}_i$ claim to commit to, and
$s_j$ for the VK-committed simple selector columns, the pairing check ultimately asserts
\begin{equation}
\sum_{j=0}^{9} B_j\,s_j(x) \;+\; (1-x^{n})\,h(x) \;=\; -A .
\label{eq:qeps:assert}
\end{equation}
Rearranging with \eqref{eq:qeps:nusplit},
\begin{equation}
\boxed{\;\nu_y(x) \;=\; \big(x^{n}-1\big)\cdot h(x)\;}
\label{eq:qeps:final}
\end{equation}
which is exactly the PLONK/Halo2 quotient relation $\nu_y(X) = Z_H(X)\,h(X)$ with
vanishing polynomial $Z_H(X) = X^{n}-1$, checked at the random point $x$.

\subsection{Assessment}
\label{sec:qeps:assessment}

\begin{finding}[qeps-6]{The $-\nu_y(x)$ convention is sound and self-consistent -- Informational (clean)}
The convention is: put the constant part of the linearization on the \emph{evaluation}
side with a minus sign (line 3358), and put the vanishing factor on the \emph{commitment}
side as $1-x^{n}$ rather than $x^{n}-1$ (line 3396). Derivation~\eqref{eq:qeps:assert}
$\Rightarrow$ \eqref{eq:qeps:final} shows the two negations cancel and reproduce the
intended relation exactly. Four properties make this sound:
\begin{enumerate}
  \item \textbf{No $h(x)$ scalar is trusted.} The proof carries no evaluation of the
        quotient polynomial. The comment at lines 176--180 asserts this, and the code bears
        it out: the opening value of the linearization query is the verifier-computed word
        \mptr{LINEARIZATION\_EVAL\_MPTR}, staged into the point-set table at line 3505, and
        no \yul{calldataload} anywhere contributes to it. $A$ is recomputed by the verifier
        from evaluations that are themselves bound to commitments through the same PCS
        batch, so the prover cannot choose the right-hand side of \eqref{eq:qeps:assert}
        independently of the left.
  \item \textbf{$y$ is sampled before the quotient commitments are read.} $y$ is the
        fifth transcript squeeze (after $\theta$ at line 1568, $\beta$ and $\gamma$ at
        1586--1587 and $\tau$ at 1646), performed at line 1663; the quotient limb
        commitments $\mathsf{Q}_i$ are absorbed only afterwards (lines 1665 onwards), and
        $x$ at line 1686 later still. So the $49$-term $y$-batch is a genuine random linear
        combination: a prover who
        satisfies \eqref{eq:qeps:final} for a random $y$ satisfies all $49$ identities
        except with probability $\le 48/r$.
  \item \textbf{The selector split does not weaken the batch.} Moving $s_j(x)$ from the
        scalar side to the commitment side is sound precisely because $\mathsf{S}_j$ is a
        VK-pinned fixed commitment, not a prover-supplied point: the selector commitments
        staged at lines 3931--3949 are read from
        $[\texttt{0x6980},\texttt{0x6f80})$, inside the VK image
        $[\texttt{0x3680},\texttt{0x7900})$ populated by \yul{extcodecopy} at line 1393.
        With \eqref{eq:qeps:Bfinal} every identity, selector-targeted or not, ends up
        weighted by a \emph{distinct} power of $y$ in $\{y^{0},\dots,y^{48}\}$, so no two
        identities can cancel.
  \item \textbf{The negation itself is exact.} \yul{addmod(0, sub(r, A), r)} is the
        canonical additive inverse on all of $[0,r)$, with no zero special case.
\end{enumerate}
The cost of the convention is one \yul{addmod}. The benefit is structural: the
linearization becomes an ordinary point-set query with a verifier-computed evaluation, so
it flows through the same $x_1/x_2/x_4$ batching machinery as every prover-supplied
evaluation instead of needing a bespoke pairing term. A sign slip in a future emitter (for
instance switching \eqref{eq:qeps:L} to $x^{n}-1$ without flipping line 3358) would assert
$\nu_y(x) = -(x^{n}-1)h(x)$; because $\mathsf{Q}_i$ are prover-chosen this would remain a
\emph{sound} statement about $-h$, and would show up as a completeness failure against
honest provers rather than as an acceptance of false proofs. That is the benign failure
direction, which is worth noting in the convention's favour.
\label{fnd:qeps-6}
\end{finding}

\begin{finding}[qeps-7]{Trash gating uses $1-\mathrm{sel}_{26}$ and is not range-restricted -- Observation}
Line 3300 forms $1 - e_{57}$ where $e_{57}$ is the evaluation of the fixed column
\yul{f\_26}. Nothing in the verifier constrains $e_{57} \in \{0,1\}$; it is an arbitrary
element of $\Fr$ bound only to its VK commitment. This is the standard Halo2 situation --
the selector column is boolean \emph{as a polynomial over the domain}, and the identity is
checked at a random $x$ where it takes an arbitrary value -- so no check is missing here.
It is recorded only because a reader may expect a booleanity guard at this line and will
not find one: the booleanity lives in the circuit and in the commitment, not in the
verifier.
\label{fnd:qeps-7}
\end{finding}


\begin{finding}[qeps-8]{\mptr{QUOTIENT\_MPTR} is repurposed from a $\GO$ slot to a scalar pair -- Observation}
Line 182 declares \mptr{QUOTIENT\_MPTR} as ``$4$ words'' -- an EIP-2537 padded G1 slot --
and the comment block at lines 3374--3377 documents the reinterpretation. Words $2$ and $3$
are never written by this path and never read by any path. The overload saves $64$ bytes of
memory plan and keeps the pointer model stable, at the cost of a name that no longer
describes its contents; a reader who follows \mptr{QUOTIENT\_MPTR} expecting a curve point
would mis-read lines 3816--3817. The generated comments mitigate this well. No correctness
or security impact.
\label{fnd:qeps-8}
\end{finding}

\begin{finding}[qeps-9]{The transcript order around $\tau$ is what makes the eight-row trash batch sound -- Informational (clean)}
The trash suffix compresses eight independent row identities into one folded scalar with
the $\tau$-Horner chain of \eqref{eq:qeps:horner}, and then subtracts a
\emph{prover-committed} polynomial evaluation $T=e_{101}$ gated by $1-\mathrm{sel}_{26}$.
Two opposite requirements meet at $\tau$, and the generated transcript satisfies both by
ordering alone:
\begin{itemize}
  \item $\tau$ must be squeezed \emph{after} everything that determines $t_0,\dots,t_7$.
        The rows read only advice columns ($a_0..a_7$ and three rotations) and VK-fixed
        columns, whose commitments are absorbed at or before line 1639. $\tau$ is squeezed
        at line 1646. So collapsing eight rows into one $\tau$-combination costs at most
        $7/r$ soundness: on any row where the gate is active the prover would have to hit a
        root of a degree-$7$ polynomial in $\tau$ fixed before $\tau$ was known.
  \item $\tau$ must be squeezed \emph{before} the trashcan commitment, because the honest
        prover builds the trashcan polynomial \emph{from} $\tau$. The trashcan commitments
        are absorbed at lines 1650--1658, immediately after the squeeze, and the source
        comment at line 1649 states the requirement in the other direction ($y$ must bind
        them). Both hold: the absorption sits between the $\tau$ squeeze (1646) and the $y$
        squeeze (1663), so the trashcan is chosen knowing $\tau$ but not $y$ and not $x$
        (line 1686).
\end{itemize}
This is the correct ordering, and it is the reason \eqref{eq:qeps:trash} is not circular
even though $T$ is prover-chosen after $\tau$: the identity is still checked at an $x$
sampled after the trashcan is committed. Note that the squeeze at line 1646 is
unconditional (comment at lines 1642--1645), so an artifact with no trash columns still
consumes the same transcript position --- a compatibility property, not a soundness one.
\label{fnd:qeps-9}
\end{finding}

\paragraph{Clean regions.} Beyond the findings above, the following properties of lines
3094--3402 were checked and hold:
\begin{itemize}
  \item Every arithmetic operation in the range is \yul{addmod}/\yul{mulmod} modulo $r$ or
        the canonical \yul{sub(r,\,\cdot)} negation pattern; there is no unreduced
        arithmetic and no operand can exceed $r$ (all inputs are either \yul{mulmod}
        outputs, VK constants, or evaluation words range-checked at line 1706).
  \item The trash suffix and the tail updates are entirely straight-line: no loop, no
        branch, no calldata read, no external call. Gas is constant and independent of the
        proof.
  \item Exhaustively, over the whole range 3094--3402 there are \emph{no} precompile calls
        (no \yul{staticcall}, \yul{call}, \yul{delegatecall}), no \yul{calldataload} or
        \yul{calldatacopy}, no \yul{extcodecopy}, no \yul{keccak256}, no \yul{mcopy}, and
        exactly three revert sites --- the three \yul{q\_program\_fail()} calls at lines
        3098, 3099 and 3105, all raising \revert{QuotientProgramInvalid}. Consequently there
        is no return-size check to specify in this range, and the only loop is the
        $k=20$ squaring walk of Listing~\ref{lst:qeps:transition}, whose bound is a literal.
  \item Every memory address touched in the range is a literal or a constant offset from
        a declared pointer constant --- \yul{add(0xb2c0, imm)},
        \yul{add(SELECTOR\_ACC\_MPTR, imm)}, \yul{add(QUOTIENT\_MPTR, 0x20)}, with
        \texttt{imm} a hex literal ---
        all of which \yul{solc} folds. Nothing in the range computes an address from proof
        data. The largest $y$-power address reached is \texttt{0xb2c0} $+$ \texttt{0x0600}
        $=$ \texttt{0xb8c0}, in bounds for the $49$-entry table.
  \item The tail exponents are exhaustively consistent with the deployed VK bytecode
        (Table~\ref{tab:qeps:tails}, last column). This audit re-decoded the payload
        independently of the section author: the instruction stream consumes exactly
        \texttt{0x11cf} bytes and lands on \yul{q\_end}, it leaves the operand stack empty,
        and its $21$ \opcode{FOLD\_SELECTOR} operands plus the four
        \opcode{NATIVE\_IDENTITY} sub-cases plus the four inline-prefix folds place the last
        contribution of bucket $j$ at exactly $p_j = 48 - T_j$ for all ten $j$. So the check
        at line 3098 passes for the honest artifact, it is not a latent liveness hazard, and
        \eqref{eq:qeps:Bfinal} holds as deployed.
\end{itemize}


%% ==== 15-pcs.tex ========================================================
\section{KZG multi-open: rotation point sets, the \texttt{q\_eval} source table, and the batched opening scalar}
\label{sec:pcs}

This section specifies \texttt{Halo2Verifier.sol} lines 3413--3812: the ``PCS
computation'' region that the generator emits as ten numbered sub-blocks. Lines
3413--3806 (sub-blocks 1--8) are $\Fr$ arithmetic plus the five
\yul{scalar\_inv} \texttt{modexp} \texttt{staticcall}s they need
(\S\ref{sec:pcs:scalarinv}); they turn the $102$
alleged polynomial evaluations, the linearization scalar, and the five
prover-supplied point-set evaluations into a single scalar
$f_{\mathrm{eval}}$ at \mptr{F\_EVAL\_MPTR}. Lines 3808--3835 open sub-block 9
and build the batched opening scalar $v$ that is stored to \mptr{V\_MPTR} at
line 4002. The commitment side of sub-block 9 (the fused 78-term \texttt{G1MSM}
staged from \mptr{PCS\_FINAL\_MSM\_MPTR}, lines 3836--4001) and sub-block 10
(the pairing assembly, lines 4004--4048) are specified elsewhere; this section
cites them only where the scalar side must be shown to match them term for
term.

\subsection{Protocol being implemented}
\label{sec:pcs:protocol}

The generated code implements the SHPLONK-style multi-open reduction used by
\texttt{midnight-proofs}. Write $\mathcal{Q}$ for the list of verifier queries,
each a triple (commitment, rotation, claimed evaluation at the rotated point).
Queries are partitioned into $J = 5$ \emph{point sets}
$S_0,\dots,S_4$. Set $S_j$ carries a rotation set
$R_j = \{z_{j,0},\dots,z_{j,c_j-1}\} \subset \Fr$ of cardinality $c_j$ and a
commitment list $C_{j,0},\dots,C_{j,m_j-1}$; every commitment in $S_j$ is
queried at every point of $R_j$, so $S_j$ contributes exactly $m_j c_j$ query
triples. Let $e_{j,i,t}$ denote the alleged evaluation of the polynomial behind
$C_{j,i}$ at $z_{j,t}$.

The reduction proceeds in four challenge-separated stages.

\begin{align}
  \mathsf{Q}_j &= \sum_{i=0}^{m_j-1} \chi_1^{\,i}\, C_{j,i}
    &&\text{($x_1$: intra-set batching, commitment side)} \label{eq:pcs:qcom}\\
  Q_j(z_{j,t}) &= \sum_{i=0}^{m_j-1} \chi_1^{\,i}\, e_{j,i,t}
    &&\text{($x_1$: intra-set batching, evaluation side)} \label{eq:pcs:qset}\\
  Z_j(X) &= \prod_{t=0}^{c_j-1} (X - z_{j,t})
    &&\text{(per-set vanishing polynomial)} \label{eq:pcs:vanish}\\
  F(X) &= \sum_{j=0}^{4} x_2^{\,j}\, \frac{Q_j(X) - r_j(X)}{Z_j(X)}
    &&\text{($x_2$: inter-set batching)} \label{eq:pcs:F}
\end{align}

where $r_j$ is the degree-$(c_j-1)$ interpolant of $Q_j$ over $R_j$. The prover
commits to $F$ as $\mathsf{F}$ (\yul{f\_com}, absorbed at line 1727 and stored at
\mptr{F\_COM\_MPTR}), the verifier squeezes $x_3$, and the prover sends the five
scalars $q_j = Q_j(x_3)$ (\yul{q\_evals}, absorbed at lines 1754--1762). The
verifier reconstructs
\begin{equation}
  f_{\mathrm{eval}} \;=\; \sum_{j=0}^{4} x_2^{\,j}\,
    \frac{q_j - r_j(x_3)}{Z_j(x_3)} \;=\; F(x_3)
  \label{eq:pcs:feval}
\end{equation}
and finally batches the six claims $\{Q_j(x_3)=q_j\}_{j=0}^4$ and
$F(x_3)=f_{\mathrm{eval}}$ with $x_4$:
\begin{equation}
  \mathsf{final\_com} = \sum_{j=0}^{4}\chi_4^{\,j}\,\mathsf{Q}_j
      + \chi_4^{\,5}\,\mathsf{F},
  \qquad
  v = \sum_{j=0}^{4}\chi_4^{\,j}\, q_j + \chi_4^{\,5}\, f_{\mathrm{eval}}.
  \label{eq:pcs:v}
\end{equation}
Sub-block 10 then checks the single KZG opening of $\mathsf{final\_com}$ to $v$
at $x_3$ against $\pi$.

\paragraph{Truncated challenges.} The symbols $\chi_1$ and $\chi_4$ above are
\emph{not} literal powers of $x_1$ and $x_4$. Following
\texttt{midnight-proofs}' \texttt{proofs/src/poly/kzg/mod.rs}, the verifier
stores
\begin{equation}
  \chi_k^{\,i} \;:=\; \bigl(x_k^{\,i} \bmod r\bigr) \bmod 2^{128},
  \qquad k \in \{1,4\},
  \label{eq:pcs:trunc}
\end{equation}
i.e.\ the power accumulator runs at full $\Fr$ precision but every value handed
to a \yul{mulmod} is masked to its low $128$ bits. $x_3$ is likewise truncated
in place at line 1743 (transcript section). $x_2$ is \emph{not} truncated. The
consequences are analysed in
\S\ref{sec:pcs:assessment} (finding PCS-3).

\subsection{Memory map used by this region}
\label{sec:pcs:memmap}

\begin{table}[htbp]
\centering
\caption{Memory constants read or written by lines 3413--3812. Declaration line
numbers refer to \texttt{Halo2Verifier.sol}.}
\label{tab:pcs:memmap}
\begin{tabular}{@{}llll@{}}
\toprule
Constant & Address & Words used & Role \\
\midrule
\mptr{OMEGA\_MPTR}        & \texttt{0x3700} & 1  & $\omega$ (VK) \\
\mptr{OMEGA\_INV\_MPTR}   & \texttt{0x3720} & 1  & $\omega^{-1}$ (VK) \\
\mptr{X\_MPTR}            & \texttt{0x79a0} & 1  & $x$ \\
\mptr{X1\_MPTR}           & \texttt{0x79c0} & 1  & $x_1$ \\
\mptr{X2\_MPTR}           & \texttt{0x79e0} & 1  & $x_2$ \\
\mptr{X3\_MPTR}           & \texttt{0x7a00} & 1  & $x_3$ (already 128-bit) \\
\mptr{X4\_MPTR}           & \texttt{0x7a20} & 1  & $x_4$ \\
\mptr{F\_COM\_MPTR}       & \texttt{0x7a40} & 4  & $\mathsf{F}$, MSM term 77 \\
\mptr{LINEARIZATION\_EVAL\_MPTR} & \texttt{0x7d00} & 1 & set-0 eval index 42 \\
\mptr{QUOTIENT\_MPTR}     & \texttt{0x7d20} & 2 of 4 & $x^{n-1}$, $1-x^{n}$ (scalars, not a point) \\
\mptr{F\_EVAL\_MPTR}      & \texttt{0x7dc0} & 1  & $f_{\mathrm{eval}}$ (written 3806) \\
\mptr{V\_MPTR}            & \texttt{0x7de0} & 1  & $v$ (written 4002) \\
\mptr{ROT\_POINTS\_MPTR}  & \texttt{0x7f80} & 4  & rotated evaluation points \\
\mptr{X1\_POWERS\_MPTR}   & \texttt{0x8300} & 43 & $\chi_1^{0..42}$ \\
\mptr{Q\_EVAL\_SET\_MPTR} & \texttt{0x8b20} & 11 & $Q_j(z_{j,t})$, all $j,t$ \\
\mptr{Q\_EVAL\_CPTR\_MPTR}& \texttt{0x9220} & 1  & calldata cursor of the 5 $q_j$ \\
\mptr{REVERSED\_EVALS\_MPTR} & \texttt{0x9480} & 102 & $\mathbf{e}=(e_0,\dots,e_{101})$ \\
\mptr{SELECTOR\_ACC\_MPTR}& \texttt{0xb140} & 10 & simple-selector accumulators \\
\mptr{PCS\_FINAL\_MSM\_MPTR} & \texttt{0xb280} & $78\times 5$ & MSM staging + eval-source scratch \\
(no constant) & \texttt{0x3580} & 6 & \yul{scalar\_inv} \texttt{modexp} frame (line 710) \\
\bottomrule
\end{tabular}
\end{table}

Two aliasing facts matter for reading this region and are checked in
\S\ref{sec:pcs:assessment}: the eval-source scratch table written by sub-blocks
3, 6 and 7 lives at \mptr{PCS\_FINAL\_MSM\_MPTR} itself, and
\mptr{LAGRANGE\_DENOMS\_MPTR} (\texttt{0xb580}, line 239) falls inside the range
that sub-block 3 overwrites.

A third address is written by this region without appearing in any constant:
each \yul{scalar\_inv} call materialises an EIP-198 \texttt{modexp} frame at the
hard-coded \yul{let p := 0x3580} (line 710), i.e.\ \texttt{0x3580}--\texttt{0x3640}.
That window lies inside the transcript buffer
(\mptr{TRANSCRIPT\_MPTR} $=\texttt{0x1000}$, line 119), which is dead by the time
this region runs, and its top is exactly \texttt{0x40} below
\mptr{VK\_MPTR} $=\texttt{0x3680}$ (line 128), so it cannot corrupt the
codehash-pinned VK words this region reads.

Two boundary facts are worth recording because they make several of the counts
below self-checking: \mptr{REVERSED\_EVALS\_MPTR} $+\ 102\times\texttt{0x20} =
\texttt{0xa140} = \mptr{ADVICE\_COMMS\_MPTR\_BASE}$ (line 257), so the eval
buffer ends exactly where the advice commitments begin; and
\mptr{PCS\_FINAL\_MSM\_MPTR} $+\ \texttt{0x30c0} = \texttt{0xe340} =
\mptr{TRACE\_U256\_MPTR}$ (line 240), so the $78$-term staging buffer ends
exactly at the next allocated region.

\subsection{Sub-block 1: the rotation point table}
\label{sec:pcs:rotpoints}

\begin{lstlisting}[style=solidity,caption={\texttt{Halo2Verifier.sol} lines
3419--3439: the four distinct rotation points. Leading indentation is stripped
uniformly in all listings in this section; the text is otherwise verbatim.},
label={lst:pcs:rotpoints}]
// 4 distinct rotation(s)
let x := mload(X_MPTR)
let omega := mload(OMEGA_MPTR)
let omega_inv := mload(OMEGA_INV_MPTR)
let x_pow_of_omega := x
mstore(add(ROT_POINTS_MPTR, 0x40), x_pow_of_omega)
x_pow_of_omega := mulmod(x_pow_of_omega, omega, r)
mstore(add(ROT_POINTS_MPTR, 0x60), x_pow_of_omega)
x_pow_of_omega := x
x_pow_of_omega := mulmod(x_pow_of_omega, omega_inv, r)
mstore(add(ROT_POINTS_MPTR, 0x20), x_pow_of_omega)
x_pow_of_omega := mulmod(x_pow_of_omega, omega_inv, r)
x_pow_of_omega := mulmod(x_pow_of_omega, omega_inv, r)
x_pow_of_omega := mulmod(x_pow_of_omega, omega_inv, r)
x_pow_of_omega := mulmod(x_pow_of_omega, omega_inv, r)
x_pow_of_omega := mulmod(x_pow_of_omega, omega_inv, r)
x_pow_of_omega := mulmod(x_pow_of_omega, omega_inv, r)
x_pow_of_omega := mulmod(x_pow_of_omega, omega_inv, r)
x_pow_of_omega := mulmod(x_pow_of_omega, omega_inv, r)
x_pow_of_omega := mulmod(x_pow_of_omega, omega_inv, r)
mstore(add(ROT_POINTS_MPTR, 0x0), x_pow_of_omega)
\end{lstlisting}

The emitter walks the positive rotations upward from $x$ and then restarts at
$x$ and walks the negative rotations downward, writing each rotation's point to
a fixed slot. Counting the \yul{mulmod} instructions exactly: line 3425 applies
$\omega$ once; line 3428 applies $\omega^{-1}$ once before the store at
\texttt{+0x20}; lines 3430--3438 apply $\omega^{-1}$ nine further times before
the store at \texttt{+0x00}. Hence

\begin{table}[htbp]
\centering
\caption{The rotation point table. $n = 2^{20}$ (the exponent $k=20$ is a
literal at line 3381).}
\label{tab:pcs:rotpoints}
\begin{tabular}{@{}llll@{}}
\toprule
Slot & Address & Value & Halo2 rotation \\
\midrule
\yul{rot\_pt\_0} & \mptr{ROT\_POINTS\_MPTR}\texttt{+0x00} & $x\,\omega^{-10}$ & $\mathrm{Rotation}(-10)$ (``last'') \\
\yul{rot\_pt\_1} & \mptr{ROT\_POINTS\_MPTR}\texttt{+0x20} & $x\,\omega^{-1}$  & $\mathrm{Rotation}(-1)$ (``prev'') \\
\yul{rot\_pt\_2} & \mptr{ROT\_POINTS\_MPTR}\texttt{+0x40} & $x$              & $\mathrm{Rotation}(0)$ (``cur'') \\
\yul{rot\_pt\_3} & \mptr{ROT\_POINTS\_MPTR}\texttt{+0x60} & $x\,\omega^{+1}$  & $\mathrm{Rotation}(+1)$ (``next'') \\
\bottomrule
\end{tabular}
\end{table}

The slot ordering is by signed rotation, not by write order, which is why the
stores are emitted out of address order. $-10$ is the permutation argument's
``last'' rotation $-(\text{blinding factors}+1)$; the VK carries the matching
$\omega^{-10}$ as \mptr{OMEGA\_INV\_TO\_L\_MPTR} (word 6 of the VK header),
used by the Lagrange block rather than here.

\paragraph{Assessment.} Only four rotations exist in this circuit, so the table
is four words although \mptr{ROT\_POINTS\_MPTR} reserves the $28$-word gap up
to \mptr{X1\_POWERS\_MPTR} ($\texttt{0x8300}-\texttt{0x7f80} = \texttt{0x380}$).
Two separate choices are visible here and should not be conflated. Unrolling the
negative walk (rather than looping it) removes a loop counter and its comparison
and costs nothing in \yul{mulmod}s. Walking down from $x$ at all, however, spends
ten \yul{mulmod}s ($10 \times 8 = 80$ gas) to rebuild $x\,\omega^{-10}$, even
though the VK already carries $\omega^{-10}$ in
\mptr{OMEGA\_INV\_TO\_L\_MPTR} (\texttt{0x3740}, line 135) --- a single
\yul{mulmod} would do, at the cost of making the emitter depend on that VK word
matching the largest negative rotation. The emitter deliberately does not take
that dependency: it derives every rotation point from $\omega^{-1}$ alone, so a
wrong \mptr{OMEGA\_INV\_TO\_L\_MPTR} cannot desynchronise the point sets from
the Lagrange block. There is no range check on $x$, $\omega$ or $\omega^{-1}$
here: $x$ is a squeezed challenge reduced mod $r$ by the transcript, and
$\omega,\omega^{-1}$ come from the codehash-pinned VK, so every value entering
\yul{mulmod} is already canonical. This is sound but relies entirely on the
upstream invariants; the block itself performs no validation.

\subsection{Sub-block 2: the truncated $x_1$ power table}
\label{sec:pcs:x1powers}

\begin{lstlisting}[style=solidity,caption={\texttt{Halo2Verifier.sol} lines
3445--3454: the $\chi_1$ table},label={lst:pcs:x1powers}]
// pre-compute 43 x1 power(s)
let x1 := mload(X1_MPTR)
mstore(X1_POWERS_MPTR, 1)
let acc := 1
let p := X1_POWERS_MPTR
for { let i := 0 } lt(i, 0x2a) { i := add(i, 1) } {
    p := add(p, 0x20)
    acc := mulmod(acc, x1, r)
    mstore(p, and(acc, 0xffffffffffffffffffffffffffffffff))
}
\end{lstlisting}

The loop bound is \texttt{0x2a} $= 42$ iterations, writing slots $1$ through
$42$; slot $0$ is stored separately as the literal $1$. The table therefore
holds exactly $43$ words:
\begin{equation}
  \mathtt{X1\_POWERS\_MPTR}[i] \;=\;
  \begin{cases}
    1, & i = 0,\\[2pt]
    (x_1^{\,i} \bmod r) \wedge (2^{128}-1), & 1 \le i \le 42.
  \end{cases}
  \label{eq:pcs:x1table}
\end{equation}
The accumulator \yul{acc} is never masked, so each stored word is the
truncation of the exact power, not a power of a truncation. Slot $0$ is
consistent with \eqref{eq:pcs:x1table} because $\chi_1^0 = 1 < 2^{128}$.

The largest index consumed anywhere is $42$: set $0$'s accumulation loop reads
up to \mptr{X1\_POWERS\_MPTR}\texttt{+0x540} (the \yul{mload(pow\_p)} at line
3510 on its final iteration, $i=42$), and the
linearization term reads the same slot at line 3918
(\yul{lin\_query\_scalar\_41 := mload(add(X1\_POWERS\_MPTR, 0x540))}). The table is
therefore exactly the size required, with no slack and no out-of-range read;
$43 = \max_j m_j$ for this circuit. The table occupies
\texttt{0x8300}--\texttt{0x8860}, well clear of \mptr{Q\_EVAL\_SET\_MPTR} at
\texttt{0x8b20}.

\paragraph{Assessment.} Materialising the powers costs $43$ words of memory and
$42$ \yul{mulmod}s once, replacing what would otherwise be a per-term running
multiply in five separate Horner loops \emph{and} in the MSM staging block. It
is what makes the eval side and the commitment side provably share the same
coefficient words: sub-blocks 3--7 and lines 3836--4001 both \yul{mload} from
this table rather than each recomputing the powers. Any drift between the two
sides is thus impossible by construction rather than by review.

\subsection{The point-set decomposition}
\label{sec:pcs:sets}

Table~\ref{tab:pcs:sets} summarises the five sets. The rotation membership is
read off the \yul{dx\_t} definitions in sub-block 8; the commitment membership
and the $x_4$ exponent are read off the staging order in sub-block 9.

\begin{table}[htbp]
\centering
\caption{The five point sets. $q_j$ is read from
\yul{calldataload(add(Q\_EVAL\_CPTR, 0x20*j))}.}
\label{tab:pcs:sets}
\begin{tabularx}{\textwidth}{@{}llllllX@{}}
\toprule
$j$ & $c_j$ & $R_j$ & $m_j$ & \mptr{Q\_EVAL\_SET\_MPTR} slots & $\chi_4$ & Commitments \\
\midrule
0 & 1 & $\{x\}$ & 43 & \texttt{0x00} & $\chi_4^0=1$ &
  \texttt{advice[14]}, committed instance (identity), \texttt{lookup\_m[0..1]},
  \texttt{lookup\_helper[0..1]}, \texttt{trashcan[0]},
  17 \texttt{fixed\_comms}, 18 \texttt{permutation\_comms}, linearized commitment \\
1 & 2 & $\{x,\,x\omega^{-1}\}$ & 3 & \texttt{0x20},\texttt{0x40} & $\chi_4^1$ & \texttt{advice[11..13]} \\
2 & 2 & $\{x,\,x\omega\}$ & 3 & \texttt{0x60},\texttt{0x80} & $\chi_4^2$ & \texttt{perm\_z[5]}, \texttt{lookup\_z[0..1]} \\
3 & 3 & $\{x,\,x\omega,\,x\omega^{-1}\}$ & 11 & \texttt{0xa0},\texttt{0xc0},\texttt{0xe0} & $\chi_4^3$ & \texttt{advice[0..10]} \\
4 & 3 & $\{x,\,x\omega,\,x\omega^{-10}\}$ & 5 & \texttt{0x100},\texttt{0x120},\texttt{0x140} & $\chi_4^4$ & \texttt{perm\_z[0..4]} \\
\bottomrule
\end{tabularx}
\end{table}

Inside a set the slot order follows the \yul{dx\_t} order, which is the order
the rotation points are first referenced: (\yul{rot\_pt\_2}, \yul{rot\_pt\_3},
\yul{rot\_pt\_0}) for set 4, (\yul{rot\_pt\_2}, \yul{rot\_pt\_3},
\yul{rot\_pt\_1}) for set 3, (\yul{rot\_pt\_2}, \yul{rot\_pt\_3}) for set 2,
(\yul{rot\_pt\_2}, \yul{rot\_pt\_1}) for set 1, and (\yul{rot\_pt\_2}) for
set 0.

The decomposition is the standard Halo2 one. Set $3$ is the fixed/advice
``cur, next, prev'' family; set $1$ holds the three advice columns queried only
at cur and prev; set $4$ is the permutation $z$ chunks $0..4$, which need the
extra $x\omega^{-10}$ opening to link chunk $i$ to chunk $i{+}1$; set $2$ holds
the final permutation chunk (no link, hence no ``last'' rotation) together with
the two lookup $z$ polynomials; set $0$ collects everything queried only at $x$.

\paragraph{Number of $q$ evaluations.} The transcript loop at lines 1754--1762
absorbs exactly \texttt{0xa0} bytes $= 5$ words starting at the cursor saved in
\mptr{Q\_EVAL\_CPTR\_MPTR}, each range-checked at line 1759 with
\yul{iszero(lt(eval, r))} raising \revert{NonCanonicalScalar}. That fixed byte count is the only thing
tying the proof shape to $J = 5$; the PCS blocks then index
\texttt{+0x00..+0x80} with no further bound check, which is safe precisely
because both numbers are codegen constants derived from the same plan.

\subsection{Sub-blocks 3--7: the \texttt{q\_eval} source table and Horner accumulation}
\label{sec:pcs:qeval}

Each set computes $c_j$ accumulators according to \eqref{eq:pcs:qset}. Two
emission forms appear, selected by $m_j$.

\paragraph{Unrolled form ($m_j < 4$).} Sets 1 and 2 emit straight-line code
with immediate eval addresses.

\begin{lstlisting}[style=solidity,caption={\texttt{Halo2Verifier.sol} lines
3521--3529: sub-block 4, $q_{\mathrm{eval\_set}}[1]$},label={lst:pcs:set1}]
// q_eval_set[1]: 3 evaluation term(s), 3 commitment term(s)
let q_eval_set_0 := mload(0x9660)
let q_eval_set_1 := mload(0x9920)
q_eval_set_0 := addmod(q_eval_set_0, mulmod(mload(0x9680), mload(add(X1_POWERS_MPTR, 0x20)), r), r)
q_eval_set_1 := addmod(q_eval_set_1, mulmod(mload(0x9940), mload(add(X1_POWERS_MPTR, 0x20)), r), r)
q_eval_set_0 := addmod(q_eval_set_0, mulmod(mload(0x96a0), mload(add(X1_POWERS_MPTR, 0x40)), r), r)
q_eval_set_1 := addmod(q_eval_set_1, mulmod(mload(0x9960), mload(add(X1_POWERS_MPTR, 0x40)), r), r)
mstore(add(Q_EVAL_SET_MPTR, 0x20), q_eval_set_0)
mstore(add(Q_EVAL_SET_MPTR, 0x40), q_eval_set_1)
\end{lstlisting}

Sub-block 5 (lines 3535--3543) is structurally identical with the addresses
\texttt{0x9fe0}/\texttt{0xa000} seeding the accumulators and
\texttt{0xa060},\texttt{0xa080},\texttt{0xa0e0},\texttt{0xa100} as the tail
terms, storing to \mptr{Q\_EVAL\_SET\_MPTR}\texttt{+0x60} and \texttt{+0x80}.
Note that these are \emph{not} Horner forms: the $i$-th term multiplies the
pre-computed $\chi_1^i$ directly, so no intermediate depends on the previous
one and the ordering carries no rounding subtleties.

\paragraph{Rolled form ($m_j \ge 4$).} Sets 0, 3 and 4 first stage a table of
\emph{eval source addresses} into memory at \texttt{0xb280}, then run a loop
that dereferences twice: \yul{mload(mload(eval\_p))}.

\begin{lstlisting}[style=solidity,caption={\texttt{Halo2Verifier.sol} lines
3460--3466 and 3503--3515: sub-block 3, $q_{\mathrm{eval\_set}}[0]$. Lines
3467--3502 are elided; they are 36 further \yul{mstore} instructions filling
\texttt{0xb300}..\texttt{0xb760} with the addresses of $e_{97},e_{98},e_{101},
e_{41},\dots,e_{73}$ in that order (table indices 4--39).},label={lst:pcs:set0}]
// q_eval_set[0]: 43 evaluation term(s), 42 commitment term(s) (rolled, m>=4)
// __phase:pcs_q_eval_source_table
// stage per-(commit, rotation) eval source addresses
mstore(0xb280, 0x9980)
mstore(0xb2a0, 0x9480)
mstore(0xb2c0, 0xa020)
mstore(0xb2e0, 0xa040)
// ... lines 3467-3502 elided ...
mstore(0xb780, 0x9dc0)
mstore(0xb7a0, 0x9de0)
mstore(0xb7c0, LINEARIZATION_EVAL_MPTR)
let q_eval_set_0 := mload(0x9980)
let pow_p := add(X1_POWERS_MPTR, 0x20)
let eval_p := add(0xb280, 0x20)
for { let i := 1 } lt(i, 0x2b) { i := add(i, 1) } {
    let pow := mload(pow_p)
    q_eval_set_0 := addmod(q_eval_set_0, mulmod(mload(mload(eval_p)), pow, r), r)
    pow_p := add(pow_p, 0x20)
    eval_p := add(eval_p, 0x20)
}
mstore(add(Q_EVAL_SET_MPTR, 0x0), q_eval_set_0)
\end{lstlisting}

The loop bound \texttt{0x2b} $= 43$ with $i$ starting at $1$ gives $42$
iterations. \yul{eval\_p} starts at \texttt{0xb2a0} and advances by
\texttt{0x20}, so the last dereference is \texttt{0xb7c0}; \yul{pow\_p} starts at
$\chi_1^1$ and ends at $\chi_1^{42}$. Together with the un-looped seed term
(\texttt{0x9980} with implicit coefficient $\chi_1^0$) the block realises
\eqref{eq:pcs:qset} for $j=0$, $m_0 = 43$, exactly.

For $c_j > 1$ the table is stored \emph{row-major by commitment}: the $c_j$
source addresses for one commitment are adjacent, and the stride is
$c_j \times \texttt{0x20}$.

\begin{lstlisting}[style=solidity,caption={\texttt{Halo2Verifier.sol} lines
3584--3599: sub-block 6's accumulation loop for
$q_{\mathrm{eval\_set}}[3]$},label={lst:pcs:set3}]
let q_eval_set_0 := mload(0x94a0)
let q_eval_set_1 := mload(0x9540)
let q_eval_set_2 := mload(0x97c0)
let pow_p := add(X1_POWERS_MPTR, 0x20)
let eval_p := add(0xb280, 0x60)
for { let i := 1 } lt(i, 0xb) { i := add(i, 1) } {
    let pow := mload(pow_p)
    q_eval_set_0 := addmod(q_eval_set_0, mulmod(mload(mload(eval_p)), pow, r), r)
    q_eval_set_1 := addmod(q_eval_set_1, mulmod(mload(mload(add(eval_p, 0x20))), pow, r), r)
    q_eval_set_2 := addmod(q_eval_set_2, mulmod(mload(mload(add(eval_p, 0x40))), pow, r), r)
    pow_p := add(pow_p, 0x20)
    eval_p := add(eval_p, 0x60)
}
mstore(add(Q_EVAL_SET_MPTR, 0xa0), q_eval_set_0)
mstore(add(Q_EVAL_SET_MPTR, 0xc0), q_eval_set_1)
mstore(add(Q_EVAL_SET_MPTR, 0xe0), q_eval_set_2)
\end{lstlisting}

The staged table for set 3 (lines 3551--3583) holds $33 = 11 \times 3$ words at
\texttt{0xb280}..\texttt{0xb680}; the loop bound \texttt{0xb} $= 11$ gives $10$
iterations covering commitments $1..10$, with commitment $0$ seeding the three
accumulators from \texttt{0x94a0}, \texttt{0x9540}, \texttt{0x97c0}. Sub-block
7 (lines 3605--3637) is the same shape with $15 = 5 \times 3$ staged words at
\texttt{0xb280}..\texttt{0xb440}, loop bound \texttt{0x5}, seeds
\texttt{0x9e00}, \texttt{0x9e20}, \texttt{0x9e40}, and destinations
\mptr{Q\_EVAL\_SET\_MPTR}\texttt{+0x100..+0x140}. All three rolled blocks share
one \emph{critically important} property: each rewrites the table from
\texttt{0xb280} upward, so the shorter tables of sub-blocks 6 and 7 leave stale
entries above their own extent, and each block reads only the prefix it wrote.

\paragraph{Source of every evaluation.} The eval buffer holds
$e_i = \mathbf{e}[i]$ at \mptr{REVERSED\_EVALS\_MPTR}$+\texttt{0x20}\,i$ for
$0 \le i \le 101$; the transcript loop at lines 1699--1714 walks exactly
\texttt{0x0cc0} bytes $= 102$ words of calldata, range-checks each against $r$
at line 1706 (\revert{NonCanonicalScalar} on failure) and spills it there. Table~\ref{tab:pcs:map} gives the full
$(\text{set},\ \chi_1\ \text{index},\ \text{commitment},\ \text{evals})$ map
recovered from the staged addresses and the MSM staging order. The commitment
names follow the per-category bases declared at lines 257--263
(\mptr{ADVICE\_COMMS\_MPTR\_BASE} $=\texttt{0xa140}$,
\mptr{LOOKUP\_M\_COMMS\_MPTR\_BASE} $=\texttt{0xa8c0}$,
\mptr{PERM\_Z\_COMMS\_MPTR\_BASE} $=\texttt{0xa9c0}$,
\mptr{LOOKUP\_HELPER\_COMMS\_MPTR\_BASE} $=\texttt{0xacc0}$,
\mptr{LOOKUP\_Z\_COMMS\_MPTR\_BASE} $=\texttt{0xadc0}$,
\mptr{TRASHCAN\_COMMS\_MPTR\_BASE} $=\texttt{0xaec0}$,
\mptr{QUOTIENT\_LIMB\_COMMS\_MPTR\_BASE} $=\texttt{0xaf40}$, each entry
\texttt{0x80} bytes) together with the VK payload layout
(\texttt{fixed\_comms} at \texttt{0x6280}$+\texttt{0x80}\,i$ for $0\le i\le 26$,
\texttt{permutation\_comms} at \texttt{0x7000}$+\texttt{0x80}\,i$ for
$0\le i\le 17$; the VK region itself is specified elsewhere). The two VK bases
are self-checking against the adjacent constants: the consecutive category
bases give category sizes $15,2,6,2,2,1$, and
$\texttt{0x6280}+27\times\texttt{0x80} = \texttt{0x7000}$ while
$\texttt{0x7000}+18\times\texttt{0x80} = \texttt{0x7900} =
\mptr{CHALLENGE\_MPTR}$ (line 144), so the $27$ fixed and $18$ permutation
commitments tile the VK payload exactly up to the challenge region.

\begin{table}[htbp]
\centering
\caption{Point-set membership map. Column $i$ is the $\chi_1$ exponent inside
the set; the ``MSM'' column is the term index in the fused \texttt{G1MSM}
staged at lines 3836--3996.}
\label{tab:pcs:map}
\begin{tabular}{@{}lllll@{}}
\toprule
Set & $i$ & Commitment & Evaluations (in $R_j$ order) & MSM \\
\midrule
0 & 0 & \texttt{advice[14]} (\texttt{0xa840}) & $e_{40}$ & 0 \\
0 & 1 & committed instance $=\id$ & $e_{0}$ & --- (omitted) \\
0 & 2 & \texttt{lookup\_m[0]} (\texttt{0xa8c0}) & $e_{93}$ & 1 \\
0 & 3 & \texttt{lookup\_helper[0]} (\texttt{0xacc0}) & $e_{94}$ & 2 \\
0 & 4 & \texttt{lookup\_m[1]} (\texttt{0xa940}) & $e_{97}$ & 3 \\
0 & 5 & \texttt{lookup\_helper[1]} (\texttt{0xad40}) & $e_{98}$ & 4 \\
0 & 6 & \texttt{trashcan[0]} (\texttt{0xaec0}) & $e_{101}$ & 5 \\
0 & 7--12 & \texttt{fixed\_comms[9,4,5,6,7,8]} & $e_{41}$--$e_{46}$ & 6--11 \\
0 & 13--16 & \texttt{fixed\_comms[0,1,2,3]} & $e_{47}$--$e_{50}$ & 12--15 \\
0 & 17--20 & \texttt{fixed\_comms[10,11,12,13]} & $e_{51}$--$e_{54}$ & 16--19 \\
0 & 21--23 & \texttt{fixed\_comms[17,24,26]} & $e_{55}$--$e_{57}$ & 20--22 \\
0 & 24--41 & \texttt{permutation\_comms[0..17]} & $e_{58}$--$e_{75}$ & 23--40 \\
0 & 42 & linearized commitment & \mptr{LINEARIZATION\_EVAL\_MPTR} & 41--54 \\
1 & 0--2 & \texttt{advice[11,12,13]} & $(e_{15},e_{37})$, $(e_{16},e_{38})$, $(e_{17},e_{39})$ & 55--57 \\
2 & 0--2 & \texttt{perm\_z[5]}, \texttt{lookup\_z[0]}, \texttt{lookup\_z[1]} & $(e_{91},e_{92})$, $(e_{95},e_{96})$, $(e_{99},e_{100})$ & 58--60 \\
3 & 0--4 & \texttt{advice[0..4]} & $(e_{1},e_{6},e_{26})\dots(e_{5},e_{19},e_{30})$ & 61--65 \\
3 & 5--10 & \texttt{advice[5..10]} & $(e_{9},e_{20},e_{31})\dots(e_{14},e_{25},e_{36})$ & 66--71 \\
4 & 0--4 & \texttt{perm\_z[0..4]} & $(e_{76+3i},e_{77+3i},e_{78+3i})$ & 72--76 \\
\bottomrule
\end{tabular}
\end{table}

Set 3's evaluation triples are, in full,
$(e_1,e_6,e_{26})$, $(e_2,e_7,e_{27})$, $(e_3,e_8,e_{28})$,
$(e_4,e_{18},e_{29})$, $(e_5,e_{19},e_{30})$, $(e_9,e_{20},e_{31})$,
$(e_{10},e_{21},e_{32})$, $(e_{11},e_{22},e_{33})$, $(e_{12},e_{23},e_{34})$,
$(e_{13},e_{24},e_{35})$, $(e_{14},e_{25},e_{36})$, in the rotation order
$(x,\,x\omega,\,x\omega^{-1})$. The non-monotone eval indices reflect the
upstream query order, which the generator is required to reproduce exactly
(see \texttt{src/lowering/kzg/mod.rs} lines 136--146 and the function
\yul{assert\_permutation\_query\_order\_is\_upstream\_equivalent} it documents at
line 147, where the plan refuses to proceed if the permutation-$z$ query order
would change the intermediate set structure).

The linearized commitment at set-0 index 42 is not a stored point. Sub-block 9
expands it into $14$ MSM terms with coefficients derived from
$\chi_1^{42}$: four quotient limb commitments (MSM terms $41..44$, copied from
\mptr{QUOTIENT\_LIMB\_COMMS\_MPTR\_BASE}$+\texttt{0x80}\,\ell$) scaled by
$\chi_1^{42}\,(1-x^{n})\,(x^{n-1})^{\ell}$ for $\ell = 0..3$ (lines
3918--3930, using the two scalars staged at \mptr{QUOTIENT\_MPTR} by lines
3400--3401, where \yul{x\_split} $= x^{n-1}$ and
\yul{one\_minus\_x\_n} $= 1-x^{n}$), and ten simple-selector
\texttt{fixed\_comms} --- indices $14,15,16,18,19,20,21,22,23,25$, MSM terms
$45..54$ --- scaled by
$\chi_1^{42}\cdot \mathtt{SELECTOR\_ACC\_MPTR}[s]$ (lines 3931--3950). The
corresponding evaluation is the single scalar
\mptr{LINEARIZATION\_EVAL\_MPTR}, written at line 3359.

\paragraph{Assessment: the double dereference.} \yul{mload(mload(eval\_p))} is
an indirect load and would be an arbitrary-memory-read primitive if the table
could be influenced. It cannot: every word in the table is an immediate
constant written by the same sub-block a few instructions earlier, and
\yul{eval\_p} is a loop-local cursor with a constant bound. No proof-derived
value ever reaches the address position. The construct trades one extra
\opcode{MLOAD} ($3$ gas) per term for a single emitted loop body in place of
$m_j$ inlined \opcode{PUSH}/\opcode{MLOAD}/\opcode{MULMOD}/\opcode{ADDMOD}
sequences; the three rolled sets ($m_j = 43, 11, 5$) are exactly the ones where
that trade pays.

The construct is safe but not self-checking, which is the subject of finding
PCS-5 below.

\subsection{Sub-block 8: vanishing scalars, batch inversion, and $f_{\mathrm{eval}}$}
\label{sec:pcs:feval}

\begin{lstlisting}[style=solidity,caption={\texttt{Halo2Verifier.sol} lines
3643--3652: sub-block 8 preamble},label={lst:pcs:feval-head}]
// __phase:scalar_inv
// f_eval via Horner over 5 reversed set(s)
let x2 := mload(X2_MPTR)
let x3 := mload(X3_MPTR)
let f_eval := 0
let Q_EVAL_CPTR := mload(Q_EVAL_CPTR_MPTR)
let rot_pt_0 := mload(add(ROT_POINTS_MPTR, 0x0))
let rot_pt_1 := mload(add(ROT_POINTS_MPTR, 0x20))
let rot_pt_2 := mload(add(ROT_POINTS_MPTR, 0x40))
let rot_pt_3 := mload(add(ROT_POINTS_MPTR, 0x60))
\end{lstlisting}

The five sets are then processed in \emph{descending} index order, $4,3,2,1,0$,
each in its own Yul block so the \yul{let} names can be reused. The reversal is
what makes a Horner fold produce ascending powers of $x_2$:
\begin{equation}
  f_{\mathrm{eval}} \;\leftarrow\; f_{\mathrm{eval}}\cdot x_2 + \mathrm{eval}_j,
  \qquad j = 4,3,2,1,0
  \;\;\Longrightarrow\;\;
  f_{\mathrm{eval}} = \sum_{j=0}^{4} x_2^{\,j}\,\mathrm{eval}_j.
  \label{eq:pcs:horner}
\end{equation}
Because \yul{f\_eval} is initialised to $0$, the first fold's multiply is a
no-op; it is emitted anyway so all five blocks are textually identical in
shape.

\subsubsection{One set of cardinality $c \ge 2$}

\begin{lstlisting}[style=solidity,caption={\texttt{Halo2Verifier.sol} lines
3653--3696: set 4, the $c=3$ template},label={lst:pcs:set4-eval}]
// --- set 4 (cardinality 3) ---
{
let dx_0 := addmod(x3, sub(r, rot_pt_2), r)
let dx_1 := addmod(x3, sub(r, rot_pt_3), r)
let dx_2 := addmod(x3, sub(r, rot_pt_0), r)
let lbasis_0 := 1
lbasis_0 := mulmod(lbasis_0, addmod(rot_pt_2, sub(r, rot_pt_3), r), r)
lbasis_0 := mulmod(lbasis_0, addmod(rot_pt_2, sub(r, rot_pt_0), r), r)
let lbasis_1 := 1
lbasis_1 := mulmod(lbasis_1, addmod(rot_pt_3, sub(r, rot_pt_2), r), r)
lbasis_1 := mulmod(lbasis_1, addmod(rot_pt_3, sub(r, rot_pt_0), r), r)
let lbasis_2 := 1
lbasis_2 := mulmod(lbasis_2, addmod(rot_pt_0, sub(r, rot_pt_2), r), r)
lbasis_2 := mulmod(lbasis_2, addmod(rot_pt_0, sub(r, rot_pt_3), r), r)
let bp_0 := dx_0
let bp_1 := mulmod(bp_0, dx_1, r)
let bp_2 := mulmod(bp_1, dx_2, r)
let bp_3 := mulmod(bp_2, lbasis_0, r)
let bp_4 := mulmod(bp_3, lbasis_1, r)
let bp_5 := mulmod(bp_4, lbasis_2, r)
let bq := scalar_inv(bp_5)
let lbasis_inv_2 := mulmod(bq, bp_4, r)
bq := mulmod(bq, lbasis_2, r)
let lbasis_inv_1 := mulmod(bq, bp_3, r)
bq := mulmod(bq, lbasis_1, r)
let lbasis_inv_0 := mulmod(bq, bp_2, r)
bq := mulmod(bq, lbasis_0, r)
let dx_inv_2 := mulmod(bq, bp_1, r)
bq := mulmod(bq, dx_2, r)
let dx_inv_1 := mulmod(bq, bp_0, r)
bq := mulmod(bq, dx_1, r)
let dx_inv_0 := bq
let den_inv := dx_inv_0
den_inv := mulmod(den_inv, dx_inv_1, r)
den_inv := mulmod(den_inv, dx_inv_2, r)
let eval := mulmod(calldataload(add(Q_EVAL_CPTR, 0x80)), den_inv, r)
let term_0 := mulmod(mulmod(mload(add(Q_EVAL_SET_MPTR, 0x100)), dx_inv_0, r), lbasis_inv_0, r)
eval := addmod(eval, sub(r, term_0), r)
let term_1 := mulmod(mulmod(mload(add(Q_EVAL_SET_MPTR, 0x120)), dx_inv_1, r), lbasis_inv_1, r)
eval := addmod(eval, sub(r, term_1), r)
let term_2 := mulmod(mulmod(mload(add(Q_EVAL_SET_MPTR, 0x140)), dx_inv_2, r), lbasis_inv_2, r)
eval := addmod(eval, sub(r, term_2), r)
f_eval := addmod(mulmod(f_eval, x2, r), eval, r)
}
\end{lstlisting}

Writing $z_t$ for the $t$-th rotation point of the set, the quantities are
\begin{align}
  \mathtt{dx\_t} &= x_3 - z_t, &
  \mathtt{lbasis\_t} &= \prod_{s \ne t} (z_t - z_s), \label{eq:pcs:dxlb}\\
  \mathtt{den\_inv} &= \Bigl(\prod_{t} (x_3 - z_t)\Bigr)^{-1} = Z_j(x_3)^{-1}. &&
  \label{eq:pcs:deninv}
\end{align}
Note that \yul{lbasis\_t} is the denominator of the $t$-th Lagrange basis
polynomial \emph{on the set's own points}, not on the evaluation domain.

The identity the block computes is the standard collapse of a barycentric
interpolant against its own vanishing polynomial:
\begin{equation}
  \frac{r_j(x_3)}{Z_j(x_3)}
  = \sum_{t} Q_j(z_t)\,
    \frac{\prod_{s\ne t}(x_3 - z_s)}{\prod_{s\ne t}(z_t - z_s)\, Z_j(x_3)}
  = \sum_{t} \frac{Q_j(z_t)}{(x_3-z_t)\,\prod_{s\ne t}(z_t-z_s)},
  \label{eq:pcs:bary}
\end{equation}
because $\prod_{s\ne t}(x_3-z_s)/Z_j(x_3) = (x_3-z_t)^{-1}$. Hence
\yul{term\_t} $=$ \eqref{eq:pcs:bary}'s $t$-th summand and
\begin{equation}
  \mathrm{eval}_j \;=\; \frac{q_j}{Z_j(x_3)} - \sum_{t=0}^{c_j-1}
    \frac{Q_j(z_{j,t})}{(x_3 - z_{j,t})\,\prod_{s\ne t}(z_{j,t}-z_{j,s})}
  \;=\; \frac{q_j - r_j(x_3)}{Z_j(x_3)},
  \label{eq:pcs:evalj}
\end{equation}
which is exactly the $j$-th summand of \eqref{eq:pcs:feval}. Subtraction is
always spelled \yul{addmod(a, sub(r, b), r)}; since every operand is already
reduced mod $r$, \yul{sub(r, b)} cannot underflow.

\subsubsection{The batch inversion}

Lines 3667--3684 are a Montgomery batch inversion over the $2c$ values
$(\mathtt{dx}_0,\dots,\mathtt{dx}_{c-1},\mathtt{lbasis}_0,\dots,\mathtt{lbasis}_{c-1})$,
in that order, using a single \yul{scalar\_inv} call. The prefix products are
\begin{equation}
  \mathtt{bp}_k =
  \begin{cases}
    \prod_{t\le k}\mathtt{dx}_t, & 0 \le k \le c-1,\\
    \bigl(\prod_{t}\mathtt{dx}_t\bigr)\prod_{u \le k-c}\mathtt{lbasis}_u, & c \le k \le 2c-1,
  \end{cases}
  \label{eq:pcs:bp}
\end{equation}
and the unwinding loop maintains the invariant that after emitting the inverse
of element $k$, the running quotient \yul{bq} equals $\mathtt{bp}_{k-1}^{-1}$.

\begin{algorithm}[htbp]
\caption{Batch inversion emitted per point set (\texttt{Halo2Verifier.sol}
3667--3684 for $c=3$; 3749--3760 and 3778--3789 for $c=2$)}
\label{alg:pcs:batchinv}
\begin{algorithmic}[1]
\State $a \gets (\mathtt{dx}_0,\dots,\mathtt{dx}_{c-1},\mathtt{lbasis}_0,\dots,\mathtt{lbasis}_{c-1})$
\State $\mathtt{bp}_0 \gets a_0$
\For{$k = 1$ to $2c-1$} \State $\mathtt{bp}_k \gets \mathtt{bp}_{k-1} \cdot a_k \bmod r$ \EndFor
\State $\mathtt{bq} \gets \mathtt{scalar\_inv}(\mathtt{bp}_{2c-1})$ \Comment{reverts if $\mathtt{bp}_{2c-1}=0$}
\For{$k = 2c-1$ down to $1$}
  \State $a_k^{-1} \gets \mathtt{bq}\cdot \mathtt{bp}_{k-1} \bmod r$
  \State $\mathtt{bq} \gets \mathtt{bq}\cdot a_k \bmod r$
\EndFor
\State $a_0^{-1} \gets \mathtt{bq}$
\end{algorithmic}
\end{algorithm}

The cost is $2c-1$ multiplications for the prefix pass, $2(2c-1)$ for the
unwinding pass, and one \yul{scalar\_inv}, versus $2c$ \yul{scalar\_inv} calls
for a naive implementation. Since \yul{scalar\_inv} is a
$\texttt{MODEXP\_GAS} = 4080$ gas \texttt{staticcall} to precompile
\texttt{0x05} (lines 700--722), the saving across the whole region is
$(2\cdot3+2\cdot3+2\cdot2+2\cdot2+1) - 5 = 16$ avoided modexp calls, roughly
$65\,000$ gas.

\subsubsection{\texttt{scalar\_inv} usages and the implied guards}
\label{sec:pcs:scalarinv}

Lines 3413--3812 contain no \yul{fail}, no \yul{revert} and no reference to the
running \yul{success} flag at all: grepping the range returns nothing. Every
revert reachable from this region therefore happens inside \yul{scalar\_inv},
and every precompile call it makes is that function's \texttt{modexp}
\texttt{staticcall}.

The region makes exactly five \yul{scalar\_inv} calls: line 3673 (set 4), line
3717 (set 3), line 3753 (set 2), line 3782 (set 1), line 3802 (set 0). The
first four invert a batch product; the fifth inverts $x_3 - x$ directly.
\yul{scalar\_inv} (lines 700--722) computes $x^{r-2} \bmod r$ and carries four
guards, in this order:

\begin{table}[htbp]
\centering
\caption{\yul{scalar\_inv} guards and precompile parameters
(\texttt{Halo2Verifier.sol} lines 700--722). All four fire on the critical path
of this region, since every inversion here goes through this one function.}
\label{tab:pcs:scalarinv}
\begin{tabular}{@{}lll@{}}
\toprule
Line & Condition & Error \\
\midrule
708 & \yul{iszero(lt(x, FR\_MODULUS))} & \revert{NonCanonicalScalar} \\
709 & \yul{iszero(x)} & \revert{NonCanonicalScalar} \\
719 & \texttt{staticcall} returned $0$ & \revert{PrecompileFailed} \\
720 & \yul{iszero(eq(returndatasize(), 0x20))} & \revert{PrecompileFailed} \\
\bottomrule
\end{tabular}
\end{table}

The \texttt{staticcall} at line 719 targets precompile address \texttt{0x05}
with the fixed gas bound $\texttt{MODEXP\_GAS} = 4080$ (line 293 --- the
EIP-7883 price of this frame, chosen over the EIP-2565 price of $1360$ so the
call also succeeds on a repriced chain), input \texttt{0xc0} bytes at
$p = \texttt{0x3580}$ (line 710) laid out as the EIP-198 frame
$(\texttt{0x20}, \texttt{0x20}, \texttt{0x20}, x, r-2, r)$, and output
\texttt{0x20} bytes back to $p$. The return-size check at line 720 is the one
that matters for a would-be-silent failure mode: \texttt{modexp} on a
non-existent or repriced chain could return short data, and without line 720 the
stale word already at $p$ (the modulus $r$, written at line 718) would be read
back as the ``inverse''.

The first guard is unreachable from this region --- every argument at the five
call sites is a \yul{mulmod}/\yul{addmod} output, hence already in $[0,r)$; the
comment at lines 701--707 says as much and keeps the check as a guard against a
future emitter. The zero test at line 709 is the load-bearing one:

\begin{equation}
  \mathtt{bp}_{2c-1} = 0
  \iff
  \underbrace{\exists\,t:\; x_3 = z_{j,t}}_{\text{some }\mathtt{dx}_t = 0}
  \;\;\lor\;\;
  \underbrace{\exists\,t\ne s:\; z_{j,t} = z_{j,s}}_{\text{some }\mathtt{lbasis}_t = 0}.
  \label{eq:pcs:zeroguard}
\end{equation}

The first disjunct is the real check: if the squeezed $x_3$ happens to land on
one of the four rotated domain points $x\omega^{\{-10,-1,0,1\}}$, the multi-open
denominators are undefined and the block reverts rather than producing a
meaningless scalar. The second disjunct cannot occur for a well-formed VK
($\omega$ has order $n$, and $-10,-1,0,1$ are distinct mod $n$) but is caught
structurally anyway. Set 0's direct \yul{scalar\_inv(dx0)} at line 3802 gives the
same guarantee for the single-point set.

\begin{finding}[PCS-1]{\texttt{scalar\_inv} conflates ``$x_3$ hit a rotation
point'' with ``non-canonical scalar'' --- Informational}
All five inversions in this region reject a zero argument by raising
\revert{NonCanonicalScalar} (\texttt{Halo2Verifier.sol}:709, reached from
3673/3717/3753/3782/3802). Failing closed is correct --- a zero denominator
would otherwise let \texttt{modexp} return $0$ and silently annihilate the
whole term chain, which the comment at lines 701--707 explicitly calls out.
But the reverting condition here is $x_3 \in \{x\omega^{-10},x\omega^{-1},x,
x\omega\}$, which is a transcript accident, not a malformed scalar. Its
probability is smaller than a first reading suggests: $x_3$ is masked to $128$
bits at line 1743 while $x$ is not masked at all, so a collision requires the
full-width point $x\omega^{j}$ to fall below $2^{128}$ \emph{and} to equal the
$128$-bit $x_3$, i.e.\ $\Pr \approx 4/r \approx 2^{-253}$, not
$4\cdot 2^{-128}$. An operator debugging a rejected proof nonetheless gets a
selector that points at calldata decoding rather than at challenge sampling ---
and the same selector is raised by the genuinely calldata-driven checks at lines
1706 and 1759, so the two cases are indistinguishable from the revert data. No
soundness impact; purely a diagnostics concern.
\end{finding}

\subsubsection{Set 0 and the store}

\begin{lstlisting}[style=solidity,caption={\texttt{Halo2Verifier.sol} lines
3799--3806: the $c=1$ specialisation and the result store},label={lst:pcs:set0-eval}]
// --- set 0 (cardinality 1) ---
{
let dx0 := addmod(x3, sub(r, rot_pt_2), r)
let dx0_inv := scalar_inv(dx0)
let eval := mulmod(addmod(calldataload(add(Q_EVAL_CPTR, 0x0)), sub(r, mload(add(Q_EVAL_SET_MPTR, 0x0))), r), dx0_inv, r)
f_eval := addmod(mulmod(f_eval, x2, r), eval, r)
}
mstore(F_EVAL_MPTR, f_eval)
\end{lstlisting}

For $c_0 = 1$ the interpolant is the constant $r_0 \equiv Q_0(x)$, so
$Z_0(x_3) = x_3 - x$ and \eqref{eq:pcs:evalj} degenerates to
$\mathrm{eval}_0 = (q_0 - Q_0(x))/(x_3-x)$ with no Lagrange basis at all. The
emitter specialises the template accordingly, saving one \yul{lbasis} chain and
the whole batch-inversion scaffold.

\subsection{The batched opening scalar $v$}
\label{sec:pcs:v}

\begin{lstlisting}[style=solidity,caption={\texttt{Halo2Verifier.sol} lines
3812--3835: the $\chi_4$ table and the formation of $v$},label={lst:pcs:v}]
// __phase:pcs_final_msm
// build final_com and v (KZG single-opening proof, fused MSM)
// final MSM input length from circuit/VK shape: 78 term(s)
let x4 := mload(X4_MPTR)
let lin_x_split := mload(QUOTIENT_MPTR)
let lin_one_minus_x_n := mload(add(QUOTIENT_MPTR, 0x20))
let Q_EVAL_CPTR := mload(Q_EVAL_CPTR_MPTR)
let x4_pow_full := 1
x4_pow_full := mulmod(x4_pow_full, x4, r)
let x4_pow_1 := and(x4_pow_full, 0xffffffffffffffffffffffffffffffff)
x4_pow_full := mulmod(x4_pow_full, x4, r)
let x4_pow_2 := and(x4_pow_full, 0xffffffffffffffffffffffffffffffff)
x4_pow_full := mulmod(x4_pow_full, x4, r)
let x4_pow_3 := and(x4_pow_full, 0xffffffffffffffffffffffffffffffff)
x4_pow_full := mulmod(x4_pow_full, x4, r)
let x4_pow_4 := and(x4_pow_full, 0xffffffffffffffffffffffffffffffff)
x4_pow_full := mulmod(x4_pow_full, x4, r)
let x4_pow_5 := and(x4_pow_full, 0xffffffffffffffffffffffffffffffff)
let v := calldataload(Q_EVAL_CPTR)
v := addmod(v, mulmod(calldataload(add(Q_EVAL_CPTR, 0x20)), x4_pow_1, r), r)
v := addmod(v, mulmod(calldataload(add(Q_EVAL_CPTR, 0x40)), x4_pow_2, r), r)
v := addmod(v, mulmod(calldataload(add(Q_EVAL_CPTR, 0x60)), x4_pow_3, r), r)
v := addmod(v, mulmod(calldataload(add(Q_EVAL_CPTR, 0x80)), x4_pow_4, r), r)
v := addmod(v, mulmod(mload(F_EVAL_MPTR), x4_pow_5, r), r)
\end{lstlisting}

Exactly five $\chi_4$ powers are materialised ($\chi_4^1 \dots \chi_4^5$);
$\chi_4^0$ is the implicit coefficient of the first term. As with $x_1$, the
accumulator \yul{x4\_pow\_full} carries the exact power and only the masked copy
is used. This realises the right-hand equation of \eqref{eq:pcs:v}:
\begin{equation}
  v = q_0 + \chi_4^{1} q_1 + \chi_4^{2} q_2 + \chi_4^{3} q_3
      + \chi_4^{4} q_4 + \chi_4^{5}\, f_{\mathrm{eval}}.
  \label{eq:pcs:vexplicit}
\end{equation}
$v$ is stored to \mptr{V\_MPTR} at line 4002, after the \texttt{G1MSM}
\texttt{staticcall}, and consumed by sub-block 10 at line 4013 as
\yul{addmod(0, sub(r, mload(V\_MPTR)), r)}, i.e. as $-v$ scaling
\mptr{G1\_BASE\_MPTR}. Both $v$ and $f_{\mathrm{eval}}$ are outputs of
\yul{addmod}/\yul{mulmod} chains, hence canonical in $\Fr$ by construction; no
explicit range check is emitted for either, and none is needed. The five $q_j$
were range-checked at line 1759 during transcript absorption and are re-read
here straight from calldata without re-validation --- correct, because calldata
is immutable within a transaction and the cursor
\mptr{Q\_EVAL\_CPTR\_MPTR} is written once at line 1753.

\paragraph{Coefficient agreement with the commitment side.} The MSM staging
block (lines 3836--3996) writes, for set $j$ commitment $i$, the scalar
\yul{mulmod(mload(add(X1\_POWERS\_MPTR, 0x20*i)), x4\_pow\_j, r)} --- the same two
truncated words that \eqref{eq:pcs:qset} and \eqref{eq:pcs:vexplicit} use,
multiplied in $\Fr$. Set 0 omits the $x_4$ factor entirely (coefficient
$\chi_1^i$ alone), matching $\chi_4^0 = 1$, and $\mathsf{F}$ receives
\yul{x4\_pow\_5}, matching the last term of \eqref{eq:pcs:vexplicit}. There is no
place where one side uses $x_k^i$ and the other $\chi_k^i$.

\subsection{Assessment}
\label{sec:pcs:assessment}

\subsubsection{Does the decomposition cover every commitment exactly once?}

Yes, and it can be checked mechanically from the fixture. Two independent
counts were reconstructed from the source text.

\emph{Evaluation side.} Summing the staged source tables of sub-blocks 3, 6 and
7 with the immediates of sub-blocks 4 and 5 gives $103$ evaluation terms: the
$102$ words of \mptr{REVERSED\_EVALS\_MPTR} each appearing \emph{exactly once},
plus \mptr{LINEARIZATION\_EVAL\_MPTR}. Formally, with $\mathcal{E}_j$ the
multiset of eval addresses consumed by set $j$,
\begin{equation}
  \biguplus_{j=0}^{4}\mathcal{E}_j
  \;=\; \{\,e_0,\dots,e_{101}\,\}\ \uplus\ \{\,\mathtt{LINEARIZATION\_EVAL\_MPTR}\,\},
  \label{eq:pcs:partition}
\end{equation}
with no repetitions and no gaps. Per set the counts are
$43 + 2\cdot 3 + 2\cdot 3 + 3\cdot 11 + 3\cdot 5 = 43+6+6+33+15 = 103$.

\emph{Commitment side.} The $78$ \yul{mcopy} sources staged into
\mptr{PCS\_FINAL\_MSM\_MPTR} land on all $78$ term slots
$\texttt{0xa0}\,k$, $0 \le k \le 77$, with no slot written twice and none
skipped; the sources are pairwise distinct addresses, and they exhaust
the circuit's commitment inventory: $15$ advice, $2$ \texttt{lookup\_m}, $2$
\texttt{lookup\_helper}, $2$ \texttt{lookup\_z}, $6$ \texttt{perm\_z}, $1$
\texttt{trashcan}, $4$ quotient limbs, all $27$ \texttt{fixed\_comms}, all $18$
\texttt{permutation\_comms}, and $\mathsf{F}$ --- $78$ in total, matching the
generated comment at line 3814 and the \texttt{staticcall} input length
$78 \times \texttt{0xa0} = \texttt{0x30c0}$ at line 3999. Of the $27$
\texttt{fixed\_comms}, $17$ carry an evaluation of their own in set 0 and $10$
are the simple selectors that enter only through the linearization term (their
queries are deliberately not in $\mathcal{Q}$ --- see
\texttt{src/lowering/kzg/mod.rs}:89--95).

Per-set $\chi_1$ exponents are dense: $0..42$ for set 0 (with one gap, below),
$0..2$ for sets 1 and 2, $0..10$ for set 3, $0..4$ for set 4. No exponent is
reused inside a set and none exceeds the $43$-entry table.

\begin{finding}[PCS-2]{The committed-instance query has an evaluation but no
commitment term --- Informational}
\label{finding:pcs:trunc-instance}
Set 0 declares ``43 evaluation term(s), 42 commitment term(s)''
(\texttt{Halo2Verifier.sol}:3460). The eval at table index 1
(\texttt{mstore(0xb2a0, 0x9480)}, line 3464 --- that is $e_0$) is multiplied by
$\chi_1^{1}$ and folded into $Q_0$, but $\chi_1^{1}$ appears nowhere in the MSM
scalar sequence of lines 3836--3917, which runs
$1,\chi_1^2,\chi_1^3,\dots,\chi_1^{41}$ and then jumps to $\chi_1^{42}$ for the
linearization. (Both halves of that statement were checked mechanically against
the fixture: the $78$ \yul{mcopy} sources of lines 3836--3995 are pairwise
distinct and none of them is \mptr{G1\_IDENTITY\_MPTR}, and the
\mptr{X1\_POWERS\_MPTR} offsets they pair with are
$\texttt{0x0},\texttt{0x40},\texttt{0x60},\dots,\texttt{0x520}$ --- every
multiple of \texttt{0x20} up to \texttt{0x520} except \texttt{0x20}.) The
generator pins the reason at
\texttt{src/lowering/kzg/mod.rs}:104--131: this query is the committed-instance
column, whose commitment is the $\GO$ identity (absorbed as 128 zero bytes),
and the MSM emitter drops identity points because they contribute nothing.
The multi-open equation still balances --- $\id$ scaled by anything is $\id$ ---
and the omission is in fact what \emph{forces} $e_0 = 0$: the eval stays in the
batched claim with a coefficient the prover cannot predict, so any nonzero
$e_0$ breaks the opening with probability $1 - 2^{-128}$. Two properties a
reviewer should note. First, the enforcement is indirect: there is no
\yul{if iszero(eq(e\_0, 0))} check anywhere, and the error probability is that of
the truncated coefficient ($2^{-128}$), not of a hard equality. Second, the
safety depends on the identity pointer being reserved for exactly this query
source; \texttt{src/lowering/kzg/mod.rs}:123--131 asserts that at generation
time, but the deployed bytecode carries no trace of the assertion. A future
emitter change that let a real commitment resolve to
\mptr{G1\_IDENTITY\_MPTR} would silently drop a live MSM term and produce an
unsound verifier that still passes honest proofs.
\end{finding}

\begin{finding}[PCS-3]{128-bit truncation of the $x_1$ and $x_4$ batching
coefficients --- Observation}
\label{finding:pcs:trunc}
Line 3453 and lines 3821--3829 mask every stored batching coefficient to
$2^{128}-1$; line 1743 masks $x_3$ itself. Each batched claim therefore has
soundness error $\approx 2^{-128}$ per term rather than $\approx 2^{-255}$, and
the multi-open evaluation point $x_3$ ranges over $2^{128}$ values instead of
$r$. This is a deliberate mirror of \texttt{midnight-proofs}'
\texttt{poly/kzg/mod.rs} (the fixture documents the rule at lines 1734--1743),
and $2^{-128}$ remains comfortably beyond any practical attack budget. On chain
the truncation is a pure cost, not an optimisation: \opcode{MULMOD} is priced at
$8$ gas regardless of operand width, so each mask is $3$ gas of \opcode{AND}
spent for nothing local ($47$ masks in this region). The reason it is there is
\emph{compatibility}: the native implementation truncates these challenges, and
the fixture's own comment (lines 1734--1742) says only that it mirrors
\texttt{proofs/src/poly/kzg/mod.rs}. Whatever the upstream motivation, the
on-chain verifier has no choice --- it must reproduce the native transcript's
scalars bit for bit or the pairing check cannot close. What it costs is $127$ bits of statistical margin and a subtle
invariant: the truncation must be applied
identically on the eval side and the commitment side. The fixture achieves that
by materialising one table and reading it from both (\S\ref{sec:pcs:v}), which
is the right structural answer. Note also the asymmetry: $x_2$ is used at full
width (line 3645) because it never appears on the commitment side --- the
$x_2$-batched polynomial is the prover-supplied $\mathsf{F}$, so no
verifier-side scalar has to match it.

One consequence of truncation that the code does not guard is that a stored
coefficient can be \emph{zero}: $\chi_1^{i} = 0$ whenever the low $128$ bits of
$x_1^{i}$ vanish, and likewise for $\chi_4^{j}$. No \yul{iszero} check exists on
any entry of \mptr{X1\_POWERS\_MPTR} or on \yul{x4\_pow\_1..5}. If it happened,
the affected query would drop out of \emph{both} the eval side and the
commitment side simultaneously --- the pairing check would still pass, but that
one commitment would be entirely unconstrained by the proof. The probability is
$\approx 2^{-128}$ per index and $\approx 47\cdot 2^{-128} \approx 2^{-122.4}$
across the region even for a prover grinding the transcript, so this is far
below any threshold that would justify a runtime check; it is recorded because
it is the only way a term can silently leave the equation, and because it is
invisible in the generated code.
\end{finding}

\subsubsection{Does the batching use independent challenges?}

Yes. The four challenges play four distinct roles and are squeezed at four
distinct transcript points, each after the data it must bind:

\begin{table}[htbp]
\centering
\caption{Fiat--Shamir ordering of the PCS challenges (transcript block, lines
1686--1767; specified in detail elsewhere).}
\label{tab:pcs:fs}
\begin{tabular}{@{}llll@{}}
\toprule
Challenge & Squeezed at & Binds & Role in this section \\
\midrule
$x$   & line 1686 & all commitments & evaluation point, defines $R_j$ \\
$x_1$ & line 1720 & $x$, all $102$ evals & intra-set batching \eqref{eq:pcs:qcom},\eqref{eq:pcs:qset} \\
$x_2$ & line 1721 & same prefix as $x_1$ & inter-set batching \eqref{eq:pcs:F},\eqref{eq:pcs:horner} \\
$x_3$ & line 1733 & $\mathsf{F}$ & opening point for $Q_j$ and $F$ \\
$x_4$ & line 1767 & the five $q_j$ & final batching \eqref{eq:pcs:v} \\
\bottomrule
\end{tabular}
\end{table}

The ordering is the sound one. $x_1$ and $x_2$ are fixed before the prover
commits to $\mathsf{F}$, so $\mathsf{F}$ cannot be adapted to them. $x_3$ is
fixed after $\mathsf{F}$, so $F$ cannot be chosen knowing where it will be
opened. $x_4$ is fixed after the five $q_j$, so the prover cannot tune $q_j$ to
cancel in the batch. $x_1$ and $x_2$ are drawn back to back from the same
prefix, which is fine: \yul{squeeze\_to} reseeds the buffer with the previous
digest, so the two are distinct outputs of a chained hash, and no prover
message is interposed that could be adapted to $x_1$ alone.

The role separation is also genuine. $x_1$ combines commitments \emph{within} a
set --- where it must, because the $m_j$ polynomials share a denominator
$Z_j$; $x_2$ combines the $c_j$-dependent quotients \emph{across} sets, where a
single challenge suffices because each summand is already a well-formed
quotient; $x_4$ then reduces $J+1 = 6$ independent opening claims at the
\emph{same} point $x_3$ to one, which is the only place a plain linear
combination is sound without further structure. Reusing $x_1$ in place of $x_4$
(or vice versa) would break that separation; the fixture does not.

\begin{finding}[PCS-4]{Sub-blocks 3, 6 and 7 stage scratch inside the
final-MSM buffer --- Observation, sound}
\label{finding:pcs:alias}
The eval-source tables are written to absolute addresses starting at
\texttt{0xb280}, which is \mptr{PCS\_FINAL\_MSM\_MPTR} itself (line 204), and
extend to \texttt{0xb7c0} (sub-block 3), \texttt{0xb680} (sub-block 6) and
\texttt{0xb440} (sub-block 7). They also overlap
\mptr{LAGRANGE\_DENOMS\_MPTR} (\texttt{0xb580}, line 239). Both reuses are
sound here because the lifetimes are disjoint in time: the Lagrange block makes
its last read of \mptr{LAGRANGE\_DENOMS\_MPTR} at line 1867 and has distilled
everything it needs into the named \mptr{L\_LAST\_MPTR} /
\mptr{L\_BLIND\_MPTR} / \mptr{L\_0\_MPTR} / \mptr{INSTANCE\_EVAL\_MPTR} slots
by line 1874, long before line 3463; and every table is
fully consumed by the loop in its own sub-block before the next one overwrites
it, with the last consumer (sub-block 7, whose final store is line 3637) preceding the first MSM
store (line 3836). The load-bearing detail is the \emph{lower} bound: no
eval-source store in this region goes below \texttt{0xb280} (the only write this
region makes below that address is \yul{scalar\_inv}'s \texttt{modexp} frame at
\texttt{0x3580}, in the dead transcript window), so the ten simple-selector
accumulators at \mptr{SELECTOR\_ACC\_MPTR} $= \texttt{0xb140}..\texttt{0xb280}$
--- which stay live from the quotient VM until they are read at lines
3932--3950 --- are never touched. The three tables sit exactly one word above
the top of that window. The MSM buffer's upper end is equally tight:
$\texttt{0xb280} + \texttt{0x30c0} = \texttt{0xe340} =
\mptr{TRACE\_U256\_MPTR}$. The cost of this arena packing is that a reader
cannot tell from the fixture alone that the overlaps are safe; the argument
lives in the generator's allocator (lines 104--118 assert it is checked at
generation time) and, for the selector accumulators, in a named regression test
cited at line 234.
\end{finding}

\begin{finding}[PCS-5]{A loop bound that disagreed with its staged table would
read stale-but-valid addresses instead of reverting --- Observation}
\label{finding:pcs:staletable}
Sub-blocks 3, 6 and 7 all rewrite the eval-source table from \texttt{0xb280}
upward, and each writes a shorter table than the one before
(\texttt{0xb7c0}, \texttt{0xb680}, \texttt{0xb440}). The words above a shorter
table's extent are therefore not garbage: they are still the valid
\mptr{REVERSED\_EVALS\_MPTR} addresses that sub-block 3 left there. Correctness
rests entirely on the loop bound matching the table length --- \texttt{0x2b} with
$43$ staged words at line 3509, \texttt{0xb} with $33$ at line 3589,
\texttt{0x5} with $15$ at line 3627 --- and both numbers come from the same
codegen plan, so today they agree (verified term by term in
\S\ref{sec:pcs:qeval}).

The observation is about the failure mode if they ever stopped agreeing. A loop
that over-ran its table would call \yul{mload(mload(eval\_p))} on a stale word
that is itself a legal eval address, so it would fold a wrong evaluation into
$Q_j$ with a legal $\chi_1$ coefficient and produce a wrong-but-well-formed
$f_{\mathrm{eval}}$. Nothing reverts, nothing reads out of bounds, and no
memory-safety check fires; the only symptom is that honest proofs stop
verifying. That is the benign direction --- it fails closed against honest
provers --- but it means the double-dereference construct carries no
self-consistency signal of its own, unlike the \yul{scalar\_inv} guards, which
do fail loudly. A cheap structural defence (a sentinel word above each table, or
deriving \yul{eval\_p}'s end pointer from the last \yul{mstore} address rather
than from an independently emitted count) would cost a few gas and remove the
dependency on two emitter outputs agreeing.
\end{finding}

\subsubsection{Regions that are clean, and why}

\begin{itemize}
\item \textbf{Loop bounds.} Every loop in the region has a literal bound whose
value is forced by the staged table it walks: \texttt{0x2a} at line 3450
($42$ iterations, powers $1..42$, plus the literal slot $0$ $=$ $43$ words),
\texttt{0x2b} at line 3509 ($42$ iterations over the $42$ table entries after
the seed $=$ $43$ set-0 terms), \texttt{0xb} at line 3589 ($10$ iterations $+$
seed $=$ $11$ set-3 commitments), \texttt{0x5} at line 3627 ($4$ iterations $+$
seed $=$ $5$ set-4 commitments). Each was checked against the number of
\yul{mstore} instructions immediately above it and against the highest
\mptr{X1\_POWERS\_MPTR} offset consumed; all four agree. No loop bound is
proof-derived.

\item \textbf{No missing range checks.} Every scalar entering an arithmetic
chain in this region is either a transcript-squeezed challenge (reduced mod $r$
by \yul{squeeze\_to}), a VK word (codehash-pinned), a proof evaluation
(range-checked at line 1706), a $q_j$ (range-checked at line 1759), or the
output of an \yul{addmod}/\yul{mulmod}. \yul{scalar\_inv} re-checks its own
argument's canonicality defensively. There is no reachable path on which a
non-canonical word reaches \yul{mulmod} with modulus $r$, and \yul{mulmod}
would in any case reduce it.

\item \textbf{No unchecked memory arithmetic.} All memory addresses in the
region are either codegen immediates or a base constant plus a codegen
immediate. The one indirect load, \yul{mload(mload(eval\_p))}, dereferences only
words that the same sub-block wrote from immediates.

\item \textbf{Determinism of the set/eval assignment.} The mapping in
Table~\ref{tab:pcs:map} is baked into immediates. Nothing in the proof can move
an evaluation from one set to another or change which $\chi_1$ power multiplies
it, so an adversary's only freedom is the \emph{values} $e_i$ and $q_j$ --- and
those are exactly what the pairing check constrains.
\end{itemize}

Two structural caveats are worth restating. Finding PCS-2: because the
commitment side is emitted from the same plan that emits the eval side, a term
dropped on the commitment side (for whatever reason) leaves no trace in the
deployed bytecode; the $43$-versus-$42$ comment at line 3460 is currently the
only in-contract signal that this has happened, and it is a comment. Finding
PCS-5: the same is true, in the other direction, of the rolled sub-blocks ---
the loop bound and the staged table length are two independently emitted numbers
whose agreement nothing in the deployed code checks. Both caveats are properties
of the emitter, not defects of this fixture, in which every count was verified
term by term above.


%% ==== 16-final.tex ======================================================
\section{Fused Final MSM, KZG Pairing, Randomized Accumulator Batching, and Return}
\label{sec:fin:top}

This section specifies \texttt{Halo2Verifier.sol} lines 3812--4130: the last five
generated phases of \yul{verifyProof}, tagged in the source as
\yul{\_\_phase:pcs\_final\_msm} (3812--4003), \yul{\_\_phase:pcs\_pairing} (4008--4034),
\yul{\_\_phase:accumulator\_pairing\_batch} (4050--4097),
\yul{\_\_phase:final\_pairing} (4112--4114), and \yul{\_\_phase:verifier\_return}
(4126--4128). Together they turn the scalar state produced by the preceding PCS
blocks into two $\GO$ points, batch a second (accumulator) pairing equation into
them with a verifier-local randomizer, run one two-pair EIP-2537 pairing check,
and return \texttt{true}.

Listings in this section are copied verbatim from the fixture; only the common
leading indentation of each excerpt has been removed so the code fits the page.

\subsection{Entry state and the \texttt{success} invariant}
\label{sec:fin:entry}

On entry at line 3812 the following are live (all established by blocks
specified elsewhere, cited by line):

\begin{center}
\begin{tabular}{@{}llp{7.4cm}@{}}
\toprule
Symbol & Slot & Meaning \\
\midrule
$x_3$ & \mptr{X3\_MPTR} $=$ \texttt{0x7a00} & PCS evaluation point, masked to its
low 128 bits at line 1743 \\
$x_4$ & \mptr{X4\_MPTR} $=$ \texttt{0x7a20} & final batching challenge, full
$\Fr$ word (line 1767) \\
$x_{\mathrm{sp}}$ & \mptr{QUOTIENT\_MPTR}$+$\texttt{0x00} & $x^{n-1}$, written at
line 3400 \\
$1-x^n$ & \mptr{QUOTIENT\_MPTR}$+$\texttt{0x20} & written at line 3401 \\
$\tau(x_1^{j})$ & \mptr{X1\_POWERS\_MPTR}$+32j$ & 43 truncated powers of $x_1$,
$j=0,\dots,42$, filled at lines 3446--3454 \\
$\sigma_j$ & \mptr{SELECTOR\_ACC\_MPTR}$+32j$ & 10 simple-selector accumulators
from the quotient VM \\
$f_{\mathrm{eval}}$ & \mptr{F\_EVAL\_MPTR} & Horner fold over the five point
sets, written at line 3806 \\
\yul{Q\_EVAL\_CPTR} & \mptr{Q\_EVAL\_CPTR\_MPTR} & calldata cursor to the five
$q$-evals (set at line 1753) \\
$\mathsf{F}$ & \mptr{F\_COM\_MPTR} & \yul{f\_com}, copied from calldata at line
1728 \\
$\pi$ & \mptr{PI\_MPTR} & KZG opening proof, copied from calldata at line 1774 \\
$\mathsf{acc}_L,\mathsf{acc}_R$ & \mptr{ACC\_LHS\_MPTR}, \mptr{ACC\_RHS\_MPTR} &
public IVC accumulator points, each a G1MSM output (lines 1288, 1337) \\
\bottomrule
\end{tabular}
\end{center}

\paragraph{\texttt{success} is provably $1$ here.}
\yul{success} is declared \yul{let success := true} at line 1356. There are
exactly six sites that write it before this section --- lines 1400--1405,
1419--1427, 1518, 1447 (tuple assignment from
\yul{validate\_public\_accumulator}) and 1826 (tuple assignment from
\yul{batch\_invert}) --- and each is followed, within the same block, by an
\yul{if iszero(success) \{ fail(...) \}} guard: 1406 ($\Rightarrow$
\revert{VkMismatch}), 1430 ($\Rightarrow$ \revert{BadCalldataShape}),
1448--1453 ($\Rightarrow$ \revert{PrecompileFailed} if
\yul{acc\_precompile\_failed}, else \revert{BadPointEncoding}), 1523 and 1789
($\Rightarrow$ \revert{NonCanonicalScalar}), and 1877--1883 ($\Rightarrow$
\revert{PrecompileFailed} if \yul{lagrange\_precompile\_failed}, else
\revert{ProofRejected}). No assignment to \yul{success} occurs between line 1826
and line 3999.
Consequently the first \yul{if success \{ ... \}} guard in this section (line
3997) is always taken, and from line 3997 onward \yul{success} is cleared only by
an EIP-2537 \opcode{STATICCALL} that failed or returned the wrong number of
bytes. That narrows the error attribution at line 4113
(\revert{PrecompileFailed}) to ``one of the nine EIP-2537 calls in this section
did not complete''; it does \emph{not} narrow it to a chain fault, because a
caller that supplies too little gas produces the same signal (FIN-11).

\subsection{Phase \yul{pcs\_final\_msm}: batching scalars and the claimed value \texorpdfstring{$v$}{v}}
\label{sec:fin:v}

\begin{lstlisting}[style=solidity,caption={\texttt{Halo2Verifier.sol} lines 3812--3835: truncated $x_4$ powers and the batched opening value $v$},label={lst:fin:v}]
// __phase:pcs_final_msm
// build final_com and v (KZG single-opening proof, fused MSM)
// final MSM input length from circuit/VK shape: 78 term(s)
let x4 := mload(X4_MPTR)
let lin_x_split := mload(QUOTIENT_MPTR)
let lin_one_minus_x_n := mload(add(QUOTIENT_MPTR, 0x20))
let Q_EVAL_CPTR := mload(Q_EVAL_CPTR_MPTR)
let x4_pow_full := 1
x4_pow_full := mulmod(x4_pow_full, x4, r)
let x4_pow_1 := and(x4_pow_full, 0xffffffffffffffffffffffffffffffff)
x4_pow_full := mulmod(x4_pow_full, x4, r)
let x4_pow_2 := and(x4_pow_full, 0xffffffffffffffffffffffffffffffff)
x4_pow_full := mulmod(x4_pow_full, x4, r)
let x4_pow_3 := and(x4_pow_full, 0xffffffffffffffffffffffffffffffff)
x4_pow_full := mulmod(x4_pow_full, x4, r)
let x4_pow_4 := and(x4_pow_full, 0xffffffffffffffffffffffffffffffff)
x4_pow_full := mulmod(x4_pow_full, x4, r)
let x4_pow_5 := and(x4_pow_full, 0xffffffffffffffffffffffffffffffff)
let v := calldataload(Q_EVAL_CPTR)
v := addmod(v, mulmod(calldataload(add(Q_EVAL_CPTR, 0x20)), x4_pow_1, r), r)
v := addmod(v, mulmod(calldataload(add(Q_EVAL_CPTR, 0x40)), x4_pow_2, r), r)
v := addmod(v, mulmod(calldataload(add(Q_EVAL_CPTR, 0x60)), x4_pow_3, r), r)
v := addmod(v, mulmod(calldataload(add(Q_EVAL_CPTR, 0x80)), x4_pow_4, r), r)
v := addmod(v, mulmod(mload(F_EVAL_MPTR), x4_pow_5, r), r)
\end{lstlisting}

Write $\tau(z) := z \bmod 2^{128}$ for the mask
\texttt{0xffffffffffffffffffffffffffffffff} (32 hex digits $=$ 128 bits). Lines
3819--3829 implement the recurrence
\[
  a_0 = 1,\qquad a_s = a_{s-1}\cdot x_4 \bmod r,\qquad
  \yul{x4\_pow\_s} = \tau(a_s) = \tau(x_4^{\,s}),\quad s = 1,\dots,5 .
\]
The accumulator \yul{x4\_pow\_full} keeps full precision; only the five emitted
copies are truncated. There are exactly $n_{\mathrm{sets}} = 5$ point sets, hence
five powers.

The claimed batched opening value at $x_3$ is (lines 3830--3835), with
$q_s := \mathtt{calldataload}(\mathtt{Q\_EVAL\_CPTR} + 32s)$ the five $q$-evals
range-checked against $r$ by the transcript loop at lines 1754--1762:
\begin{equation}
  v \;=\; q_0 \;+\; \sum_{s=1}^{4} \tau(x_4^{\,s})\, q_s \;+\;
  \tau(x_4^{5})\, f_{\mathrm{eval}} \pmod r .
  \label{eq:fin:v}
\end{equation}
$v$ is written to \mptr{V\_MPTR} at line 4002, after the MSM call.

\begin{finding}[FIN-1]{Batching scalars are truncated to 128 bits (severity: Informational)}
Every scalar the PCS batching uses is masked to its low 128 bits: $x_3$ at line
1743, $\tau(x_1^{j})$ at line 3453 (the mask is applied to the
stored copy only; the accumulator keeps full precision), and $\tau(x_4^{s})$ at lines 3821--3829. The
generator documents this as deliberate parity with \texttt{midnight-proofs}
(\texttt{proofs/src/poly/kzg/mod.rs}); see
\texttt{src/lowering/kzg/mod.rs:991--994} and lines 1735--1742 of the fixture.
\emph{Consequence:} the multi-open batching argument has soundness error
$\approx 2^{-128}$ per batched term rather than $\approx 2^{-255}$. That is the
intended security level of the native protocol, so this is not a divergence --
but a reviewer sizing the overall soundness margin of the deployed verifier must
use $2^{-128}$, not $\log_2 r$.
\emph{What it buys:} bit-exact agreement with the native Rust verifier, which is
what makes the differential fixtures meaningful. \emph{What it costs:} half the
statistical margin, and a masking idiom (\yul{and} with a 128-bit literal) that
looks like a bug at a glance.
\emph{Why the implementation is nonetheless internally consistent:} the identical
truncated values are used on \emph{both} sides of the opening equation --
$\tau(x_4^{s})$ multiplies $q_s$ in~\eqref{eq:fin:v} and the set-$s$ commitments
in the MSM; $\tau(x_1^{j})$ multiplies the $j$-th eval in the $q$-eval fold
(lines 3506--3515) and the $j$-th commitment in the MSM. Completeness therefore
holds exactly.
\end{finding}

\subsection{The fused final-MSM staging buffer}
\label{sec:fin:msm}

\subsubsection{Layout}

\mptr{PCS\_FINAL\_MSM\_MPTR} $=$ \texttt{0xb280} is the base of a flat run of
\emph{78} EIP-2537 G1MSM \emph{(point, scalar)} pairs. Each pair occupies
\texttt{0xa0} bytes:

\begin{center}
\begin{tabular}{@{}llp{7.6cm}@{}}
\toprule
Offset in pair & Size & Content \\
\midrule
\texttt{0x00} & \texttt{0x80} & G1 point, EIP-2537 padded
$(x_{\mathrm{hi}}, x_{\mathrm{lo}}, y_{\mathrm{hi}}, y_{\mathrm{lo}})$, written by
one \opcode{MCOPY} \\
\texttt{0x80} & \texttt{0x20} & big-endian scalar, written by one
\opcode{MSTORE} \\
\bottomrule
\end{tabular}
\end{center}

Hence term $k$ has its point at \mptr{PCS\_FINAL\_MSM\_MPTR}$+$\texttt{0xa0}$\cdot k$
and its scalar at \mptr{PCS\_FINAL\_MSM\_MPTR}$+$\texttt{0xa0}$\cdot k +$\texttt{0x80}.
The last scalar is at offset \texttt{0x30a0} ($k=77$), and the total input length
is $78 \times \mathtt{0xa0} = \mathtt{0x30c0} = 12480$ bytes, exactly the length
passed to the precompile at line 3999. The buffer spans
\texttt{0xb280}--\texttt{0xe340}; \mptr{SELECTOR\_ACC\_MPTR} $=$ \texttt{0xb140}
ends exactly at \texttt{0xb280} and \mptr{TRACE\_U256\_MPTR} $=$ \texttt{0xe340}
begins exactly where the buffer ends, so the staging writes clobber nothing live.
Earlier phases reuse the same addresses but are dead by line 3836: the set-0
$q$-eval source table occupies \texttt{0xb280}--\texttt{0xb7e0} (43 one-word
entries, lines 3463--3505, consumed by the fold loop at 3506--3515), the
per-set tables for sets 1--4 reuse the same base (lines 3552, 3608, \dots), and
the Lagrange denominators sit at \mptr{LAGRANGE\_DENOMS\_MPTR} $=$
\texttt{0xb580} (written 1816--1835, last read at line 1867).

\subsubsection{Emission pattern}

The block is fully unrolled: lines 3836--3996 are 161 statements, namely 78
\opcode{MCOPY}/\opcode{MSTORE} pairs with no loop and no branch, interleaved
with five pure scalar-arithmetic statements inside the linearization expansion
(3918, 3919, 3922, 3925, 3928). Listing~\ref{lst:fin:stage} shows the first four
terms; terms 4--40 and 45--77 follow the identical two-instruction shape, terms
41--44 do not (Section~\ref{sec:fin:lin}), and all 78 are enumerated in
Tables~\ref{tab:fin:termsA}--\ref{tab:fin:termsC}.

Decoding every destination in the block shows that the 78 pair slots are each
written exactly once, with no gap and no double write: the \opcode{MCOPY}
destinations enumerate $\mathtt{0xa0}\cdot k$ for $k=0,\dots,77$ and the
\opcode{MSTORE} destinations enumerate $\mathtt{0xa0}\cdot k + \mathtt{0x80}$
for the same $k$, each value appearing once. That total coverage is load-bearing
rather than incidental --- see FIN-12.

\begin{lstlisting}[style=solidity,caption={\texttt{Halo2Verifier.sol} lines 3836--3843: staging terms 0--3 of 78 (terms 4--40 and 45--77, lines 3844--3917 and 3931--3996, follow the same two-instruction shape; terms 41--44 do not, see Listing~\ref{lst:fin:lin})},label={lst:fin:stage}]
mcopy(PCS_FINAL_MSM_MPTR, 0xa840, 0x80)
mstore(add(PCS_FINAL_MSM_MPTR, 0x80), 1)
mcopy(add(PCS_FINAL_MSM_MPTR, 0xa0), 0xa8c0, 0x80)
mstore(add(PCS_FINAL_MSM_MPTR, 0x120), mload(add(X1_POWERS_MPTR, 0x40)))
mcopy(add(PCS_FINAL_MSM_MPTR, 0x140), 0xacc0, 0x80)
mstore(add(PCS_FINAL_MSM_MPTR, 0x1c0), mload(add(X1_POWERS_MPTR, 0x60)))
mcopy(add(PCS_FINAL_MSM_MPTR, 0x1e0), 0xa940, 0x80)
mstore(add(PCS_FINAL_MSM_MPTR, 0x260), mload(add(X1_POWERS_MPTR, 0x80)))
\end{lstlisting}

\subsubsection{The mathematical object being built}

Let the five point sets be indexed $s = 0,\dots,4$, and within set $s$ let
$\mathsf{C}_{s,0},\mathsf{C}_{s,1},\dots$ be the queried commitments in the
generator's order. The MSM computes
\begin{equation}
  \mathsf{final\_com}
  \;=\;
  \sum_{s=0}^{4}\ \sum_{i}\ \tau(x_1^{\,i})\,\tau(x_4^{\,s})\cdot \mathsf{C}_{s,i}
  \;+\;
  \tau(x_4^{5})\cdot \mathsf{F},
  \label{eq:fin:finalcom}
\end{equation}
with the conventions $\tau(x_1^{0}) = \tau(x_4^{0}) = 1$ (the generator folds
these to the literal \yul{1} or omits the factor entirely, which is why term 0
carries the immediate \yul{1} and terms 55, 58, 61, 72 carry a bare
\yul{x4\_pow\_s}), and with two structural exceptions:

\begin{enumerate}
  \item \textbf{Identity commitments are omitted.} Set 0 has 43 queried
  commitments, $i = 0,\dots,42$, but only 41 of them appear as MSM terms plus the
  expansion described next. Index $i = 1$ produces \emph{no} MSM term: the
  emitted scalars jump from the literal \yul{1} ($i=0$) straight to
  \mptr{X1\_POWERS\_MPTR}$+$\texttt{0x40} ($i=2$) at line 3839.
  \item \textbf{The linearized commitment is expanded in place.} Index $i = 42$
  of set 0 is the linearization commitment. Instead of materializing it with its
  own G1MSM, the generator splices its 14 constituent terms into the fused MSM
  (terms 41--54, Section~\ref{sec:fin:lin}).
\end{enumerate}

\begin{finding}[FIN-2]{The committed-instance commitment is omitted from the MSM while its evaluation stays in the batch (severity: Informational)}
\emph{What the fixture alone shows:} the $x_1$ exponent sequence staged by the
MSM is $0, 2, 3, \dots, 42$ --- the jump at line 3839 from the literal \yul{1}
to \mptr{X1\_POWERS\_MPTR}$+$\texttt{0x40} skips exponent~$1$ --- while the
set-0 fold loop at 3509 runs $i = 1,\dots,42$ and so consumes 43 evaluations.
Nothing in \texttt{Halo2Verifier.sol} references \mptr{G1\_IDENTITY\_MPTR}
(\texttt{0x9320}) at all: it is declared at line 218 and never read, precisely
because the emitter deletes the term rather than staging four zero words.
\emph{What the generator adds:} the skipped query is a committed-instance
column, plan-pinned to the identity commitment. The MSM emitter drops such
terms, but the corresponding evaluation
\emph{is} folded into $q_{\mathrm{eval}}[0]$ with coefficient $\tau(x_1^{1})$ (the
source-table entry \yul{mstore(0xb2a0, 0x9480)} at line 3464 and the fold loop at
lines 3506--3515; that block is specified elsewhere). The generator's own comment
at \texttt{src/lowering/kzg/mod.rs:104--122} names the consequence: the block
header comment at line 3460 reads ``43 evaluation term(s), 42 commitment
term(s)''.
\emph{Assessment:} this is sound and deliberate -- omitting an $\id$ commitment
changes nothing on the commitment side, and leaving the eval in the batched claim
is precisely what forces the committed-instance evaluation to be zero. But the
enforcement is \emph{indirect}: nowhere does the verifier compare that eval
against $0$. A nonzero value is caught only because it perturbs the batched
opening by $\tau(x_1^{1})\cdot e$, which by FIN-1 gives error $2^{-128}$, not by
an equality check. A reviewer looking for ``where is the committed instance
constrained?'' will not find a guard, and must be pointed at this argument. The
generator hard-asserts that no other query source may acquire the identity
pointer (same file, the \yul{assert!} at lines 124--131); that assertion is what
keeps the omission from silently dropping a real MSM term.
\end{finding}

\subsubsection{The linearization expansion (terms 41--54)}
\label{sec:fin:lin}

\begin{lstlisting}[style=solidity,caption={\texttt{Halo2Verifier.sol} lines 3918--3932: expanding the linearized commitment into 4 quotient-limb terms and the first of 10 simple-selector terms},label={lst:fin:lin}]
let lin_query_scalar_41 := mload(add(X1_POWERS_MPTR, 0x540))
let lin_cur_scalar_41 := mulmod(lin_query_scalar_41, lin_one_minus_x_n, r)
mcopy(0xcc20, add(QUOTIENT_LIMB_COMMS_MPTR_BASE, 0x0), 0x80)
mstore(add(PCS_FINAL_MSM_MPTR, 0x1a20), lin_cur_scalar_41)
lin_cur_scalar_41 := mulmod(lin_cur_scalar_41, lin_x_split, r)
mcopy(0xccc0, add(QUOTIENT_LIMB_COMMS_MPTR_BASE, 0x80), 0x80)
mstore(add(PCS_FINAL_MSM_MPTR, 0x1ac0), lin_cur_scalar_41)
lin_cur_scalar_41 := mulmod(lin_cur_scalar_41, lin_x_split, r)
mcopy(0xcd60, add(QUOTIENT_LIMB_COMMS_MPTR_BASE, 0x100), 0x80)
mstore(add(PCS_FINAL_MSM_MPTR, 0x1b60), lin_cur_scalar_41)
lin_cur_scalar_41 := mulmod(lin_cur_scalar_41, lin_x_split, r)
mcopy(0xce00, add(QUOTIENT_LIMB_COMMS_MPTR_BASE, 0x180), 0x80)
mstore(add(PCS_FINAL_MSM_MPTR, 0x1c00), lin_cur_scalar_41)
mcopy(add(PCS_FINAL_MSM_MPTR, 0x1c20), 0x6980, 0x80)
mstore(add(PCS_FINAL_MSM_MPTR, 0x1ca0), mulmod(lin_query_scalar_41, mload(add(SELECTOR_ACC_MPTR, 0x0)), r))
\end{lstlisting}

\yul{lin\_query\_scalar\_41} $= \tau(x_1^{42})$ is the batching coefficient the
linearized commitment would have received as a single MSM term; the suffix
\texttt{41} is the \emph{MSM pair index} at the point of emission, not the query
index. The variable \yul{lin\_cur\_scalar\_41} is \emph{mutated between the stores}
(lines 3922, 3925, 3928), so although all four \opcode{MSTORE}s render the same
identifier they write four different values:
\[
  \text{term } 41+j:\quad
  \tau(x_1^{42})\,(1 - x^{n})\, x_{\mathrm{sp}}^{\,j},\qquad j = 0,1,2,3,
  \qquad x_{\mathrm{sp}} = x^{\,n-1}.
\]
The ten simple-selector terms 45--54 use the unmutated
\yul{lin\_query\_scalar\_41}: term $45+j$ carries scalar
$\tau(x_1^{42})\cdot\sigma_j$ against VK fixed commitment
$\mathsf{S}_{c_j}$ with $(c_0,\dots,c_9) = (14,15,16,18,19,20,21,22,23,25)$. The
resulting sub-sum is exactly $\tau(x_1^{42})$ times the linearized commitment
described at lines 3363--3378,
\[
  \mathsf{L} \;=\; (1 - x^{n}) \sum_{j=0}^{3} x_{\mathrm{sp}}^{\,j}\,\mathsf{Q}_j
  \;+\; \sum_{j=0}^{9} \sigma_j\, \mathsf{S}_{c_j}.
\]

\begin{finding}[FIN-3]{Four staging destinations are emitted as absolute literals, breaking the self-describing addressing convention (severity: Observation)}
Lines 3920, 3923, 3926, 3929 write the four quotient-limb \emph{points} with
\yul{mcopy(0xcc20, ...)}, \yul{mcopy(0xccc0, ...)}, \yul{mcopy(0xcd60, ...)},
\yul{mcopy(0xce00, ...)}, while the matching \emph{scalars} on lines 3921, 3924,
3927, 3930 use the symbolic \yul{add(PCS\_FINAL\_MSM\_MPTR, 0x1a20)} form, and every
one of the other 74 terms uses the symbolic form for both halves. The literals
are correct today: $\mathtt{0xb280} + \mathtt{0x19a0} = \mathtt{0xcc20}$,
$+\mathtt{0x1a40} = \mathtt{0xccc0}$, $+\mathtt{0x1ae0} = \mathtt{0xcd60}$,
$+\mathtt{0x1b80} = \mathtt{0xce00}$, each exactly \texttt{0x80} below its
scalar. The cause is generator-side: \texttt{src/lowering/kzg/mod.rs:1874--1882}
renders the \opcode{MCOPY} destination straight from the raw
\texttt{pair\_base} value, while the \opcode{MSTORE} destination goes through
the \texttt{final\_msm\_addr} helper. Both are
derived from the same \yul{memory.pcs\_final\_msm\_scratch\_mptr}, so they cannot
desynchronize within one generation run; the defect is auditability, not
correctness. It is nevertheless the one place in a 78-term buffer where a reader
cannot see which term a store belongs to without doing hex arithmetic --- exactly
what commit \texttt{41b78eb9} set out to eliminate.
\end{finding}

\subsubsection{Complete term table}

Tables~\ref{tab:fin:termsA}--\ref{tab:fin:termsC} enumerate all 78 terms.
Notation: $\mathsf{A}_i$ advice, $\mathsf{M}_i$ lookup multiplicity,
$\mathsf{H}_i$ lookup helper, $\mathsf{Z}^{\mathrm{p}}_i$ permutation grand
product, $\mathsf{Z}^{\ell}_i$ lookup grand product, $\mathsf{T}_0$ trashcan,
$\mathsf{Q}_j$ quotient limb, $\mathsf{S}_i$ VK fixed commitment, $\mathsf{P}_j$
VK permutation commitment, $\mathsf{F}$ \yul{f\_com}. ``Offset'' is the point
offset from \mptr{PCS\_FINAL\_MSM\_MPTR}; the scalar sits \texttt{0x80} higher.
$\dagger$ marks the four destinations emitted as absolute literals (FIN-3).

\begin{table}[p]
\centering
\small
\begin{tabular}{@{}rlllll@{}}
\toprule
$k$ & Offset & Source & Point & Origin & Scalar \\
\midrule
0 & \texttt{0x0000} & \texttt{0xa840} & $\mathsf{A}_{14}$ & advice 14 & $1$ \\
1 & \texttt{0x00a0} & \texttt{0xa8c0} & $\mathsf{M}_{0}$ & lookup $m$ 0 & $\tau(x_1^{2})$ \\
2 & \texttt{0x0140} & \texttt{0xacc0} & $\mathsf{H}_{0}$ & lookup helper 0 & $\tau(x_1^{3})$ \\
3 & \texttt{0x01e0} & \texttt{0xa940} & $\mathsf{M}_{1}$ & lookup $m$ 1 & $\tau(x_1^{4})$ \\
4 & \texttt{0x0280} & \texttt{0xad40} & $\mathsf{H}_{1}$ & lookup helper 1 & $\tau(x_1^{5})$ \\
5 & \texttt{0x0320} & \texttt{0xaec0} & $\mathsf{T}_{0}$ & trashcan 0 & $\tau(x_1^{6})$ \\
6 & \texttt{0x03c0} & \texttt{0x6700} & $\mathsf{S}_{9}$ & VK fixed 9 & $\tau(x_1^{7})$ \\
7 & \texttt{0x0460} & \texttt{0x6480} & $\mathsf{S}_{4}$ & VK fixed 4 & $\tau(x_1^{8})$ \\
8 & \texttt{0x0500} & \texttt{0x6500} & $\mathsf{S}_{5}$ & VK fixed 5 & $\tau(x_1^{9})$ \\
9 & \texttt{0x05a0} & \texttt{0x6580} & $\mathsf{S}_{6}$ & VK fixed 6 & $\tau(x_1^{10})$ \\
10 & \texttt{0x0640} & \texttt{0x6600} & $\mathsf{S}_{7}$ & VK fixed 7 & $\tau(x_1^{11})$ \\
11 & \texttt{0x06e0} & \texttt{0x6680} & $\mathsf{S}_{8}$ & VK fixed 8 & $\tau(x_1^{12})$ \\
12 & \texttt{0x0780} & \texttt{0x6280} & $\mathsf{S}_{0}$ & VK fixed 0 & $\tau(x_1^{13})$ \\
13 & \texttt{0x0820} & \texttt{0x6300} & $\mathsf{S}_{1}$ & VK fixed 1 & $\tau(x_1^{14})$ \\
14 & \texttt{0x08c0} & \texttt{0x6380} & $\mathsf{S}_{2}$ & VK fixed 2 & $\tau(x_1^{15})$ \\
15 & \texttt{0x0960} & \texttt{0x6400} & $\mathsf{S}_{3}$ & VK fixed 3 & $\tau(x_1^{16})$ \\
16 & \texttt{0x0a00} & \texttt{0x6780} & $\mathsf{S}_{10}$ & VK fixed 10 & $\tau(x_1^{17})$ \\
17 & \texttt{0x0aa0} & \texttt{0x6800} & $\mathsf{S}_{11}$ & VK fixed 11 & $\tau(x_1^{18})$ \\
18 & \texttt{0x0b40} & \texttt{0x6880} & $\mathsf{S}_{12}$ & VK fixed 12 & $\tau(x_1^{19})$ \\
19 & \texttt{0x0be0} & \texttt{0x6900} & $\mathsf{S}_{13}$ & VK fixed 13 & $\tau(x_1^{20})$ \\
20 & \texttt{0x0c80} & \texttt{0x6b00} & $\mathsf{S}_{17}$ & VK fixed 17 & $\tau(x_1^{21})$ \\
21 & \texttt{0x0d20} & \texttt{0x6e80} & $\mathsf{S}_{24}$ & VK fixed 24 & $\tau(x_1^{22})$ \\
22 & \texttt{0x0dc0} & \texttt{0x6f80} & $\mathsf{S}_{26}$ & VK fixed 26 & $\tau(x_1^{23})$ \\
23 & \texttt{0x0e60} & \texttt{0x7000} & $\mathsf{P}_{0}$ & VK perm.\ 0 & $\tau(x_1^{24})$ \\
24 & \texttt{0x0f00} & \texttt{0x7080} & $\mathsf{P}_{1}$ & VK perm.\ 1 & $\tau(x_1^{25})$ \\
25 & \texttt{0x0fa0} & \texttt{0x7100} & $\mathsf{P}_{2}$ & VK perm.\ 2 & $\tau(x_1^{26})$ \\
26 & \texttt{0x1040} & \texttt{0x7180} & $\mathsf{P}_{3}$ & VK perm.\ 3 & $\tau(x_1^{27})$ \\
27 & \texttt{0x10e0} & \texttt{0x7200} & $\mathsf{P}_{4}$ & VK perm.\ 4 & $\tau(x_1^{28})$ \\
28 & \texttt{0x1180} & \texttt{0x7280} & $\mathsf{P}_{5}$ & VK perm.\ 5 & $\tau(x_1^{29})$ \\
29 & \texttt{0x1220} & \texttt{0x7300} & $\mathsf{P}_{6}$ & VK perm.\ 6 & $\tau(x_1^{30})$ \\
30 & \texttt{0x12c0} & \texttt{0x7380} & $\mathsf{P}_{7}$ & VK perm.\ 7 & $\tau(x_1^{31})$ \\
31 & \texttt{0x1360} & \texttt{0x7400} & $\mathsf{P}_{8}$ & VK perm.\ 8 & $\tau(x_1^{32})$ \\
32 & \texttt{0x1400} & \texttt{0x7480} & $\mathsf{P}_{9}$ & VK perm.\ 9 & $\tau(x_1^{33})$ \\
33 & \texttt{0x14a0} & \texttt{0x7500} & $\mathsf{P}_{10}$ & VK perm.\ 10 & $\tau(x_1^{34})$ \\
34 & \texttt{0x1540} & \texttt{0x7580} & $\mathsf{P}_{11}$ & VK perm.\ 11 & $\tau(x_1^{35})$ \\
35 & \texttt{0x15e0} & \texttt{0x7600} & $\mathsf{P}_{12}$ & VK perm.\ 12 & $\tau(x_1^{36})$ \\
36 & \texttt{0x1680} & \texttt{0x7680} & $\mathsf{P}_{13}$ & VK perm.\ 13 & $\tau(x_1^{37})$ \\
37 & \texttt{0x1720} & \texttt{0x7700} & $\mathsf{P}_{14}$ & VK perm.\ 14 & $\tau(x_1^{38})$ \\
38 & \texttt{0x17c0} & \texttt{0x7780} & $\mathsf{P}_{15}$ & VK perm.\ 15 & $\tau(x_1^{39})$ \\
39 & \texttt{0x1860} & \texttt{0x7800} & $\mathsf{P}_{16}$ & VK perm.\ 16 & $\tau(x_1^{40})$ \\
40 & \texttt{0x1900} & \texttt{0x7880} & $\mathsf{P}_{17}$ & VK perm.\ 17 & $\tau(x_1^{41})$ \\
\bottomrule
\end{tabular}
\caption{Final-MSM terms 0--40 (\texttt{Halo2Verifier.sol} lines 3836--3917): point set 0, direct queries, common factor $\tau(x_4^{0}) = 1$. The $x_1$ exponent skips $1$: see FIN-2.}
\label{tab:fin:termsA}
\end{table}


\begin{table}[p]
\centering
\small
\begin{tabular}{@{}rlllll@{}}
\toprule
$k$ & Offset & Source & Point & Origin & Scalar \\
\midrule
41 & \texttt{0x19a0}$^{\dagger}$ & \texttt{0xaf40} & $\mathsf{Q}_{0}$ & quotient limb 0 & $\tau(x_1^{42})\,(1-x^n)\,x_{\mathrm{sp}}^{0}$ \\
42 & \texttt{0x1a40}$^{\dagger}$ & \texttt{0xafc0} & $\mathsf{Q}_{1}$ & quotient limb 1 & $\tau(x_1^{42})\,(1-x^n)\,x_{\mathrm{sp}}^{1}$ \\
43 & \texttt{0x1ae0}$^{\dagger}$ & \texttt{0xb040} & $\mathsf{Q}_{2}$ & quotient limb 2 & $\tau(x_1^{42})\,(1-x^n)\,x_{\mathrm{sp}}^{2}$ \\
44 & \texttt{0x1b80}$^{\dagger}$ & \texttt{0xb0c0} & $\mathsf{Q}_{3}$ & quotient limb 3 & $\tau(x_1^{42})\,(1-x^n)\,x_{\mathrm{sp}}^{3}$ \\
45 & \texttt{0x1c20} & \texttt{0x6980} & $\mathsf{S}_{14}$ & VK fixed 14 & $\tau(x_1^{42})\cdot\sigma_{0}$ \\
46 & \texttt{0x1cc0} & \texttt{0x6a00} & $\mathsf{S}_{15}$ & VK fixed 15 & $\tau(x_1^{42})\cdot\sigma_{1}$ \\
47 & \texttt{0x1d60} & \texttt{0x6a80} & $\mathsf{S}_{16}$ & VK fixed 16 & $\tau(x_1^{42})\cdot\sigma_{2}$ \\
48 & \texttt{0x1e00} & \texttt{0x6b80} & $\mathsf{S}_{18}$ & VK fixed 18 & $\tau(x_1^{42})\cdot\sigma_{3}$ \\
49 & \texttt{0x1ea0} & \texttt{0x6c00} & $\mathsf{S}_{19}$ & VK fixed 19 & $\tau(x_1^{42})\cdot\sigma_{4}$ \\
50 & \texttt{0x1f40} & \texttt{0x6c80} & $\mathsf{S}_{20}$ & VK fixed 20 & $\tau(x_1^{42})\cdot\sigma_{5}$ \\
51 & \texttt{0x1fe0} & \texttt{0x6d00} & $\mathsf{S}_{21}$ & VK fixed 21 & $\tau(x_1^{42})\cdot\sigma_{6}$ \\
52 & \texttt{0x2080} & \texttt{0x6d80} & $\mathsf{S}_{22}$ & VK fixed 22 & $\tau(x_1^{42})\cdot\sigma_{7}$ \\
53 & \texttt{0x2120} & \texttt{0x6e00} & $\mathsf{S}_{23}$ & VK fixed 23 & $\tau(x_1^{42})\cdot\sigma_{8}$ \\
54 & \texttt{0x21c0} & \texttt{0x6f00} & $\mathsf{S}_{25}$ & VK fixed 25 & $\tau(x_1^{42})\cdot\sigma_{9}$ \\
\bottomrule
\end{tabular}
\caption{Final-MSM terms 41--54 (lines 3918--3950): the in-place expansion of set 0's linearization query at $\tau(x_1^{42})$. The \texttt{Source} column gives the resolved address; terms 41--44 are the only ones whose source is written symbolically in the fixture, as \yul{add(QUOTIENT\_LIMB\_COMMS\_MPTR\_BASE, 0x0/0x80/0x100/0x180)} with \mptr{QUOTIENT\_LIMB\_COMMS\_MPTR\_BASE} $=$ \texttt{0xaf40} (line 263).}
\label{tab:fin:termsB}
\end{table}


\begin{table}[p]
\centering
\small
\begin{tabular}{@{}rlllll@{}}
\toprule
$k$ & Offset & Source & Point & Origin & Scalar \\
\midrule
55 & \texttt{0x2260} & \texttt{0xa6c0} & $\mathsf{A}_{11}$ & advice 11 & $\tau(x_4^{1})$ \\
56 & \texttt{0x2300} & \texttt{0xa740} & $\mathsf{A}_{12}$ & advice 12 & $\tau(x_1^{1})\,\tau(x_4^{1})$ \\
57 & \texttt{0x23a0} & \texttt{0xa7c0} & $\mathsf{A}_{13}$ & advice 13 & $\tau(x_1^{2})\,\tau(x_4^{1})$ \\
58 & \texttt{0x2440} & \texttt{0xac40} & $\mathsf{Z}^{\mathrm{p}}_{5}$ & perm.\ $z$ 5 & $\tau(x_4^{2})$ \\
59 & \texttt{0x24e0} & \texttt{0xadc0} & $\mathsf{Z}^{\ell}_{0}$ & lookup $z$ 0 & $\tau(x_1^{1})\,\tau(x_4^{2})$ \\
60 & \texttt{0x2580} & \texttt{0xae40} & $\mathsf{Z}^{\ell}_{1}$ & lookup $z$ 1 & $\tau(x_1^{2})\,\tau(x_4^{2})$ \\
61 & \texttt{0x2620} & \texttt{0xa140} & $\mathsf{A}_{0}$ & advice 0 & $\tau(x_4^{3})$ \\
62 & \texttt{0x26c0} & \texttt{0xa1c0} & $\mathsf{A}_{1}$ & advice 1 & $\tau(x_1^{1})\,\tau(x_4^{3})$ \\
63 & \texttt{0x2760} & \texttt{0xa240} & $\mathsf{A}_{2}$ & advice 2 & $\tau(x_1^{2})\,\tau(x_4^{3})$ \\
64 & \texttt{0x2800} & \texttt{0xa2c0} & $\mathsf{A}_{3}$ & advice 3 & $\tau(x_1^{3})\,\tau(x_4^{3})$ \\
65 & \texttt{0x28a0} & \texttt{0xa340} & $\mathsf{A}_{4}$ & advice 4 & $\tau(x_1^{4})\,\tau(x_4^{3})$ \\
66 & \texttt{0x2940} & \texttt{0xa3c0} & $\mathsf{A}_{5}$ & advice 5 & $\tau(x_1^{5})\,\tau(x_4^{3})$ \\
67 & \texttt{0x29e0} & \texttt{0xa440} & $\mathsf{A}_{6}$ & advice 6 & $\tau(x_1^{6})\,\tau(x_4^{3})$ \\
68 & \texttt{0x2a80} & \texttt{0xa4c0} & $\mathsf{A}_{7}$ & advice 7 & $\tau(x_1^{7})\,\tau(x_4^{3})$ \\
69 & \texttt{0x2b20} & \texttt{0xa540} & $\mathsf{A}_{8}$ & advice 8 & $\tau(x_1^{8})\,\tau(x_4^{3})$ \\
70 & \texttt{0x2bc0} & \texttt{0xa5c0} & $\mathsf{A}_{9}$ & advice 9 & $\tau(x_1^{9})\,\tau(x_4^{3})$ \\
71 & \texttt{0x2c60} & \texttt{0xa640} & $\mathsf{A}_{10}$ & advice 10 & $\tau(x_1^{10})\,\tau(x_4^{3})$ \\
72 & \texttt{0x2d00} & \texttt{0xa9c0} & $\mathsf{Z}^{\mathrm{p}}_{0}$ & perm.\ $z$ 0 & $\tau(x_4^{4})$ \\
73 & \texttt{0x2da0} & \texttt{0xaa40} & $\mathsf{Z}^{\mathrm{p}}_{1}$ & perm.\ $z$ 1 & $\tau(x_1^{1})\,\tau(x_4^{4})$ \\
74 & \texttt{0x2e40} & \texttt{0xaac0} & $\mathsf{Z}^{\mathrm{p}}_{2}$ & perm.\ $z$ 2 & $\tau(x_1^{2})\,\tau(x_4^{4})$ \\
75 & \texttt{0x2ee0} & \texttt{0xab40} & $\mathsf{Z}^{\mathrm{p}}_{3}$ & perm.\ $z$ 3 & $\tau(x_1^{3})\,\tau(x_4^{4})$ \\
76 & \texttt{0x2f80} & \texttt{0xabc0} & $\mathsf{Z}^{\mathrm{p}}_{4}$ & perm.\ $z$ 4 & $\tau(x_1^{4})\,\tau(x_4^{4})$ \\
77 & \texttt{0x3020} & \texttt{0x7a40} & $\mathsf{F}$ & \yul{f\_com} & $\tau(x_4^{5})$ \\
\bottomrule
\end{tabular}
\caption{Final-MSM terms 55--77 (lines 3951--3996): point sets 1--4 and the \yul{f\_com} tail. Term 77's source is written as \mptr{F\_COM\_MPTR} $=$ \texttt{0x7a40} (line 162).}
\label{tab:fin:termsC}
\end{table}


Set cardinalities read off the tables: set 0 contributes 55 MSM terms (41 direct
$+$ 14 from the linearization expansion) out of 43 queries; set 1 contributes 3
($\mathsf{A}_{11}$ through $\mathsf{A}_{13}$); set 2 contributes 3
($\mathsf{Z}^{\mathrm{p}}_{5}, \mathsf{Z}^{\ell}_{0}, \mathsf{Z}^{\ell}_{1}$);
set 3 contributes 11 ($\mathsf{A}_{0}$ through $\mathsf{A}_{10}$); set 4 contributes 5
($\mathsf{Z}^{\mathrm{p}}_{0}$ through $\mathsf{Z}^{\mathrm{p}}_{4}$); plus one \yul{f\_com} tail term:
$55+3+3+11+5+1 = 78$.

\subsubsection{The G1MSM call}

\begin{lstlisting}[style=solidity,caption={\texttt{Halo2Verifier.sol} lines 3993--4002: last staged term, the fused G1MSM, and the $v$ spill},label={lst:fin:call}]
mcopy(add(PCS_FINAL_MSM_MPTR, 0x2f80), 0xabc0, 0x80)
mstore(add(PCS_FINAL_MSM_MPTR, 0x3000), mulmod(mload(add(X1_POWERS_MPTR, 0x80)), x4_pow_4, r))
mcopy(add(PCS_FINAL_MSM_MPTR, 0x3020), F_COM_MPTR, 0x80)
mstore(add(PCS_FINAL_MSM_MPTR, 0x30a0), x4_pow_5)
if success {
    // exact EIP-2537 G1MSM cost for 78 pair(s)
    success := staticcall(525096, 0x0c, PCS_FINAL_MSM_MPTR, 0x30c0, FINAL_COM_MPTR, 0x80)
    success := and(success, eq(returndatasize(), 0x80))
}
mstore(V_MPTR, v)
\end{lstlisting}

One \opcode{STATICCALL} to precompile \texttt{0x0c} (EIP-2537 G1MSM) with
\texttt{0x30c0} input bytes, output \texttt{0x80} bytes into
\mptr{FINAL\_COM\_MPTR} $=$ \texttt{0x7e00}. Two guards: the call's own success
flag, and \yul{eq(returndatasize(), 0x80)} --- the latter rejects a
short-returning or stubbed precompile whose partial output would otherwise leave
stale words in \mptr{FINAL\_COM\_MPTR}.

The forwarded gas is the literal $525{,}096$. EIP-2537 prices a $k$-pair G1MSM at
$\lfloor k \cdot 12000 \cdot d(k) / 1000 \rfloor$ with $d(78) = 561$ from the
discount table; $78 \cdot 12000 \cdot 561 / 1000 = 525{,}096$, so the bound is
exact, not padded. Because the schedule is the specified worst case, a conformant
chain cannot run out of gas \emph{inside} the call once the call is entered with
its full bound; and because a failing EIP-2537 call consumes all forwarded gas,
forwarding the exact figure rather than \yul{gas()} caps the burn of a malformed
proof at $525$k instead of $63/64$ of the transaction budget. The converse case
--- the frame not \emph{having} $525{,}096$ gas to forward, so EIP-150 hands the
precompile less than its bound --- is not excluded by anything here and is
covered by FIN-11.

\begin{finding}[FIN-4]{The final MSM's gas bound is a bare literal rather than the named constant it equals (severity: Observation)}
Line 3999 forwards \yul{525096}. The contract already declares
\yul{G1MSM\_GAS\_SMOKE = 525096} at line 295, and the constructor smoke probe at
line 555 forwards that constant over the same \texttt{0x30c0}-byte frame
(line 550) --- the probe is sized from this very MSM, which is what makes
deployment onto a repriced chain fail at construction rather than at proof time.
The runtime site does not reference the constant, so the two are linked only by
the generator (\texttt{src/lowering/artifacts.rs:307--315} versus
\texttt{src/lowering/kzg/mod.rs:1942--1946}; both call
\texttt{layout::gas::g1msm\_gas}). No live defect: the values are equal and
both derive from the same term count.
\emph{Related, and clean here:} the discount table
(\texttt{src/lowering/layout/mod.rs:179--188}) caps at $k = 128$ and reuses
$d = 519$ beyond it, which would \emph{under}-fund a larger MSM. This artifact's
$k = 78 \le 128$ uses the exact table entry $d(78) = 561$, so that hazard is not
reachable in this fixture.
\end{finding}

\begin{finding}[FIN-5]{Every point entering the pairing is precompile-validated (severity: Informational, clean)}
EIP-2537 G1MSM validates each input point's curve and subgroup membership
regardless of the associated scalar, and the fused call at line 3999 feeds it all
27 VK fixed commitments used here, all 18 VK permutation commitments, and every
proof-carried commitment listed in
Tables~\ref{tab:fin:termsA}--\ref{tab:fin:termsC}. $\pi$ is not in that MSM, but
is validated by the one-pair G1MSM at line 4026; $\mathsf{acc}_L$ and
$\mathsf{acc}_R$ were validated by the G1MSMs at lines 1288 and 1337. Every point
subsequently handed to the pairing precompile is therefore itself a G1ADD or
G1MSM \emph{output}. Staging a point with scalar $0$ does not exempt it from
validation, which is the property the accumulator block at lines 1277--1281
explicitly relies on. This region is sound for that reason.
\end{finding}

\begin{finding}[FIN-12]{Total, non-overlapping coverage of the staging buffer is what keeps stale scratch out of the MSM (severity: Informational)}
Immediately before line 3836 the region \texttt{0xb280}--\texttt{0xb7e0} holds
the set-0 $q$-eval \emph{source-address table} (43 words such as
\yul{mstore(0xb280, 0x9980)} at line 3463), and \texttt{0xb580} onward held the
Lagrange denominators. The former are small integers in the range
\texttt{0x7d00}--\texttt{0xa120} (they are memory addresses), the latter are
arbitrary $\Fr$ elements; neither is a valid EIP-2537 coordinate pair. The staging block never
zeroes the buffer; it relies entirely on writing every one of the
$78 \times 5 = 390$ words. Decoding all 156 destinations in lines 3836--3996
confirms exactly that: each of the 78 \opcode{MCOPY} destinations
$\mathtt{0xa0}k$ and each of the 78 \opcode{MSTORE} destinations
$\mathtt{0xa0}k + \mathtt{0x80}$ occurs once and only once, for
$k = 0,\dots,77$, and the highest address touched is
$\mathtt{0xb280} + \mathtt{0x30a0} + \mathtt{0x20} = \mathtt{0xe340}$, exactly
the declared end of the buffer.
\emph{Why it matters and why it is nonetheless sound:} a missed pair slot would
feed a leftover source-address word to the precompile as a coordinate. That is a
liveness/completeness bug (G1MSM would reject the input, \yul{success} would
clear, and line 4113 would revert with \revert{PrecompileFailed}), not a
soundness bug --- there is no leftover word in this buffer that is a valid
curve point. The property is a consequence of the emitter's counter, not of a
runtime check: nothing in the contract verifies that all 78 slots were written.
The generator does assert the count is stable across emission
(\texttt{src/lowering/kzg/mod.rs:1933--1936}, ``final MSM input term count
changed during emission''), which is where this invariant is actually enforced.
\end{finding}

\subsection{Phase \yul{pcs\_pairing}: staging the two KZG pairing inputs}
\label{sec:fin:pairing}

\begin{lstlisting}[style=solidity,caption={\texttt{Halo2Verifier.sol} lines 4008--4033: building $\pi$ and $\mathsf{final\_com} - v G + x_3 \pi$},label={lst:fin:pcspairing}]
// __phase:pcs_pairing
// Scale z*pi - vG before the final pairing check
// pairing inputs (LHS = pi; RHS = final_com - v*G + x3*pi)
mcopy(PAIRING_LHS_MPTR, PI_MPTR, 0x80)
mcopy(0x1000, G1_BASE_MPTR, 0x80)
mstore(0x1080, addmod(0, sub(r, mload(V_MPTR)), r))
if success {
    success := staticcall(G1MSM_GAS_1PAIR, 0x0c, 0x1000, 0xa0, 0x1000, 0x80)
    success := and(success, eq(returndatasize(), 0x80))
}
mcopy(0x1080, FINAL_COM_MPTR, 0x80)
if success {
    success := staticcall(G1ADD_GAS, 0x0b, 0x1000, 0x100, 0x1000, 0x80)
    success := and(success, eq(returndatasize(), 0x80))
}
mcopy(0x1080, PI_MPTR, 0x80)
mstore(0x1100, mload(X3_MPTR))
if success {
    success := staticcall(G1MSM_GAS_1PAIR, 0x0c, 0x1080, 0xa0, 0x1080, 0x80)
    success := and(success, eq(returndatasize(), 0x80))
}
if success {
    success := staticcall(G1ADD_GAS, 0x0b, 0x1000, 0x100, 0x1000, 0x80)
    success := and(success, eq(returndatasize(), 0x80))
}
mcopy(PAIRING_RHS_MPTR, 0x1000, 0x80)
\end{lstlisting}

\begin{algorithm}
\caption{\yul{pcs\_pairing} (lines 4011--4033). Scratch base $S = $ \texttt{0x1000}.}
\label{alg:fin:pcspairing}
\begin{algorithmic}[1]
\State $[\mptr{PAIRING\_LHS\_MPTR}] \gets \pi$ \Comment{plain \opcode{MCOPY}, no precompile}
\State $[S,\ S{+}\mathtt{0x80}) \gets G$; $[S{+}\mathtt{0x80}] \gets (-v) \bmod r$
\State $[S] \gets \opcode{G1MSM}(S, \mathtt{0xa0})$ \Comment{$(-v)\,G$}
\State $[S{+}\mathtt{0x80}] \gets \mathsf{final\_com}$
\State $[S] \gets \opcode{G1ADD}(S, \mathtt{0x100})$ \Comment{$\mathsf{final\_com} - vG$}
\State $[S{+}\mathtt{0x80}] \gets \pi$; $[S{+}\mathtt{0x100}] \gets x_3$
\State $[S{+}\mathtt{0x80}] \gets \opcode{G1MSM}(S{+}\mathtt{0x80}, \mathtt{0xa0})$ \Comment{$x_3\,\pi$}
\State $[S] \gets \opcode{G1ADD}(S, \mathtt{0x100})$ \Comment{$\mathsf{final\_com} - vG + x_3\pi$}
\State $[\mptr{PAIRING\_RHS\_MPTR}] \gets [S]$
\end{algorithmic}
\end{algorithm}

Exactly four precompile calls, in order: G1MSM($\mathtt{0x0c}$),
G1ADD($\mathtt{0x0b}$), G1MSM, G1ADD, each guarded by \yul{if success} and each
followed by a \yul{returndatasize} equality check. Gas is the exact schedule:
\yul{G1MSM\_GAS\_1PAIR} $= 12000$ ($k=1$, discount $1000$) and \yul{G1ADD\_GAS}
$= 375$. On exit,
\[
  \mptr{PAIRING\_LHS\_MPTR} = \pi, \qquad
  \mptr{PAIRING\_RHS\_MPTR} = \mathsf{final\_com} - v\,G + x_3\,\pi .
\]

\paragraph{Detail notes.}
\begin{itemize}
\item The negation on line 4013 is \yul{addmod(0, sub(r, mload(V\_MPTR)), r)}, not
\yul{sub(r, v)}. This matters exactly at $v = 0$: \yul{sub(r,0)} $= r$, and the
outer \yul{addmod} reduces it to $0$, so the scalar stays canonical and the MSM
returns $\id$. Without the \yul{addmod} the precompile would receive the
non-reduced scalar $r$; EIP-2537 accepts it, but the artifact would then differ
from the native verifier's representation. The guard is cheap and correct.
\item Memory reuse is in-place: the G1MSM at line 4015 reads
\texttt{0x1000}--\texttt{0x10a0} and writes \texttt{0x1000}--\texttt{0x1080}; the
G1MSM at line 4026 reads \texttt{0x1080}--\texttt{0x1120} and writes
\texttt{0x1080}--\texttt{0x1100}. \opcode{STATICCALL} copies return data only
after the callee has read its input, so aliasing input and output is safe.
\item The scratch base \texttt{0x1000} is \mptr{TRANSCRIPT\_MPTR} $=$
\mptr{RETURN\_MPTR}, reused here in a later, disjoint memory phase; the generator
allocates these from one arena that rejects overlapping live regions (see the
constants commentary at lines 104--118).
\item The comment on line 4009 calls the scalar \texttt{z}; the code uses $x_3$.
Stale GWC-era naming, no behavioural consequence.
\item If \yul{success} were already $0$ the four \yul{if success} bodies are
skipped and \mptr{PAIRING\_RHS\_MPTR} retains its zero-initialised value ($\id$).
That state is unreachable (Section~\ref{sec:fin:entry}), and even if reached the
guard at line 4113 reverts before the pairing.
\end{itemize}

\subsection{Phase \yul{accumulator\_pairing\_batch}}
\label{sec:fin:batch}

The block (lines 4050--4097) folds a \emph{second} pairing equation --- the
public IVC accumulator's --- into the two KZG pairing inputs, using a
verifier-local randomizer $\alpha$.

\subsubsection{Domain tag, preimage, and \texorpdfstring{$\alpha$}{alpha}}

\begin{lstlisting}[style=solidity,caption={\texttt{Halo2Verifier.sol} lines 4051--4067: the $\alpha$ Fiat--Shamir frame},label={lst:fin:alpha}]
{
    let batch_ptr := 0x1000

    // Domain || vk_digest || KZG rhs/lhs || accumulator rhs/lhs.
    // vk_digest makes alpha's binding to the verifying key local
    // instead of transitive-through-the-points (audit I-7).
    mstore(batch_ptr, 0x70616972696e672d62617463682d6163632d6b7a670000000000000000)
    mstore(add(batch_ptr, 0x20), mload(VK_DIGEST_MPTR))
    mcopy(add(batch_ptr, 0x40),  PAIRING_RHS_MPTR, 0x80)
    mcopy(add(batch_ptr, 0xc0),  PAIRING_LHS_MPTR, 0x80)
    mcopy(add(batch_ptr, 0x0140), ACC_RHS_MPTR,     0x80)
    mcopy(add(batch_ptr, 0x01c0), ACC_LHS_MPTR,     0x80)
    // alpha is Fiat-Shamir over the fully materialized pairing
    // inputs. Replace the negligible zero draw with one so the
    // accumulator equation cannot be accidentally dropped.
    let acc_pair_alpha := mod(keccak256(batch_ptr, 0x0240), r)
    if iszero(acc_pair_alpha) { acc_pair_alpha := 1 }
\end{lstlisting}

\paragraph{The domain tag.}
The literal \texttt{0x70616972696e672d62617463682d6163632d6b7a670000000000000000} is
29 bytes (58 hex digits): the ASCII string \texttt{"pairing-batch-acc-kzg"}
(21 bytes) followed by eight zero bytes. \opcode{MSTORE} left-pads a numeric
literal to a full word, so the 32 bytes actually hashed are
\[
  \underbrace{\texttt{00 00 00}}_{\text{3 pad bytes}}
  \;\Vert\;
  \texttt{"pairing-batch-acc-kzg"}
  \;\Vert\;
  \underbrace{\texttt{00}\times 8}_{\text{literal tail}} .
\]
This is \emph{not} the right-padded \yul{bytes32("pairing-batch-acc-kzg")} form.
Domain separation is unaffected (any fixed unique word works), but any
re-implementation or differential fixture that derives $\alpha$ from the
right-padded form will disagree with the deployed verifier on accept/reject. The
generator pins the exact literal and documents the discrepancy at
\texttt{src/lowering/layout/mod.rs:311--328}.

\paragraph{Preimage layout (\texttt{0x0240} $= 576$ bytes).}
\begin{center}
\begin{tabular}{@{}lllp{5.6cm}@{}}
\toprule
Offset & Size & Source & Content \\
\midrule
\texttt{0x0000} & \texttt{0x20} & literal & domain tag word (above) \\
\texttt{0x0020} & \texttt{0x20} & \mptr{VK\_DIGEST\_MPTR} & $\mathsf{vk\_digest}$ \\
\texttt{0x0040} & \texttt{0x80} & \mptr{PAIRING\_RHS\_MPTR} & $\mathsf{kzg}_R = \mathsf{final\_com} - vG + x_3\pi$ \\
\texttt{0x00c0} & \texttt{0x80} & \mptr{PAIRING\_LHS\_MPTR} & $\mathsf{kzg}_L = \pi$ \\
\texttt{0x0140} & \texttt{0x80} & \mptr{ACC\_RHS\_MPTR} & $\mathsf{acc}_R$ \\
\texttt{0x01c0} & \texttt{0x80} & \mptr{ACC\_LHS\_MPTR} & $\mathsf{acc}_L$ \\
\midrule
\texttt{0x0240} & --- & \multicolumn{2}{@{}l}{total, matching \yul{keccak256(batch\_ptr, 0x0240)}} \\
\bottomrule
\end{tabular}
\end{center}
$2 \times \mathtt{0x20} + 4 \times \mathtt{0x80} = \mathtt{0x240}$; the frame
occupies \texttt{0x1000}--\texttt{0x1240} and ends exactly where
\yul{ec\_pairing}'s scratch begins (\texttt{0x1240}, line 969). Then
\[
  \alpha \;=\; \operatorname{keccak256}(\text{frame}) \bmod r,
  \qquad
  \alpha \gets 1 \ \text{if}\ \alpha = 0 .
\]

\subsubsection{The two randomized additions}

\begin{lstlisting}[style=solidity,caption={\texttt{Halo2Verifier.sol} lines 4069--4097: $\mathsf{kzg}_R \mathrel{+}= \alpha\,\mathsf{acc}_R$ and $\mathsf{kzg}_L \mathrel{+}= \alpha\,\mathsf{acc}_L$},label={lst:fin:add}]
    // PAIRING_RHS_MPTR += alpha * ACC_RHS_MPTR.
    // First compute alpha * ACC_RHS with a one-pair G1MSM, then
    // add it into the KZG RHS point.
    mcopy(batch_ptr, ACC_RHS_MPTR, 0x80)
    mstore(add(batch_ptr, 0x80), acc_pair_alpha)
    if success {
        success := staticcall(G1MSM_GAS_1PAIR, 0x0c, batch_ptr, 0xa0, batch_ptr, 0x80)
        success := and(success, eq(returndatasize(), 0x80))
    }
    mcopy(add(batch_ptr, 0x80), PAIRING_RHS_MPTR, 0x80)
    if success {
        success := staticcall(G1ADD_GAS, 0x0b, batch_ptr, 0x0100, PAIRING_RHS_MPTR, 0x80)
        success := and(success, eq(returndatasize(), 0x80))
    }

    // PAIRING_LHS_MPTR += alpha * ACC_LHS_MPTR.
    // Mirror the same randomized batching on the KZG LHS point.
    mcopy(batch_ptr, ACC_LHS_MPTR, 0x80)
    mstore(add(batch_ptr, 0x80), acc_pair_alpha)
    if success {
        success := staticcall(G1MSM_GAS_1PAIR, 0x0c, batch_ptr, 0xa0, batch_ptr, 0x80)
        success := and(success, eq(returndatasize(), 0x80))
    }
    mcopy(add(batch_ptr, 0x80), PAIRING_LHS_MPTR, 0x80)
    if success {
        success := staticcall(G1ADD_GAS, 0x0b, batch_ptr, 0x0100, PAIRING_LHS_MPTR, 0x80)
        success := and(success, eq(returndatasize(), 0x80))
    }
}
\end{lstlisting}

Four more precompile calls (G1MSM, G1ADD, G1MSM, G1ADD), same gas bounds and same
two-guard pattern. The $\alpha$ frame at \texttt{0x1000} is overwritten by the
first \opcode{MCOPY} on line 4072 --- after $\alpha$ has been computed, so the
digest is unaffected. Each G1ADD writes its \texttt{0x80}-byte result directly
into the destination slot, so the updates are in place:
\[
  \mptr{PAIRING\_RHS\_MPTR} \gets \mathsf{kzg}_R + \alpha\,\mathsf{acc}_R,
  \qquad
  \mptr{PAIRING\_LHS\_MPTR} \gets \mathsf{kzg}_L + \alpha\,\mathsf{acc}_L .
\]
If a G1ADD fails, no return data is copied, the destination keeps its previous
(pre-batch) value, and \yul{success} is cleared --- which line 4113 turns into a
revert before any pairing runs.

\subsubsection{Soundness of the randomized batching}
\label{sec:fin:soundness}

Fix a generator of $\GO$ and write $a, b, c, d$ for the discrete logs of
$\mathsf{kzg}_R, \mathsf{kzg}_L, \mathsf{acc}_R, \mathsf{acc}_L$, and $s$ for the
SRS secret behind $[s]_2$. Define the two \emph{residuals}
\[
  U := a - s\,b \quad(\text{KZG equation}), \qquad
  V := c - s\,d \quad(\text{accumulator equation}).
\]
Each equation holds iff its residual is $0$. The check actually performed
(Section~\ref{sec:fin:swap}) is
$\pair{\mathsf{kzg}_R + \alpha\,\mathsf{acc}_R}{[1]_2}
 = \pair{\mathsf{kzg}_L + \alpha\,\mathsf{acc}_L}{[s]_2}$, i.e.
\begin{equation}
  U + \alpha V \;=\; 0 \pmod r .
  \label{eq:fin:batch}
\end{equation}

\begin{itemize}
\item $U = V = 0$: \eqref{eq:fin:batch} holds for every $\alpha$ --- completeness.
\item $V \neq 0$: \eqref{eq:fin:batch} holds for exactly one $\alpha = -U/V$.
\item $V = 0,\ U \neq 0$: \eqref{eq:fin:batch} holds for \emph{no} $\alpha$.
\end{itemize}
So the source comment's claim (lines 4048--4049) --- ``if either original equation
is bad, this combined equation holds for at most one alpha in $\Fr$'' --- is
literally correct, in both the $V \neq 0$ and $V = 0$ cases.

\begin{finding}[FIN-6]{The ``two bad equations could cancel'' rationale is correct, and the mitigation is the right one (severity: Informational, clean)}
The rejected alternative --- multiplying the two pairing equations, i.e.
checking $U + V = 0$ --- is genuinely broken: an adversary who can produce a KZG
residual $U \neq 0$ and choose the public accumulator so that $V = -U$ passes a
combined check while both equations individually fail. The accumulator points
come from \emph{public inputs} (decoded at lines 1246--1351), so the adversary
does control them directly.
\emph{One precision the source comment does not make:} mounting that attack in
\emph{this} verifier is not free, because the accumulator public inputs are
absorbed into the transcript (lines 1510--1522) \emph{before} $x, x_1, x_3, x_4$
are squeezed, and $U$ is a function of those challenges. Choosing
$\mathsf{acc}$ to satisfy $V = -U$ is therefore a transcript fixed-point
problem, not a one-shot algebraic solve. The mitigation is still the right one:
it removes the degree of freedom outright instead of relying on that ordering
argument, at a cost of four precompile calls. Both design and implementation are
correct here.
\end{finding}

\begin{finding}[FIN-7]{$\alpha$ binds everything the check depends on (severity: Informational, clean)}
The combined equation~\eqref{eq:fin:batch} is a statement about four $\GO$ points
and two fixed $\GT$ bases and nothing else. All four points are in the preimage,
each in full \texttt{0x80}-byte EIP-2537 form and each in a fixed position, so the
frame is unambiguously parsed (no length-prefix or concatenation ambiguity is
possible with fixed-width fields). $\mathsf{vk\_digest}$ is also absorbed
(audit item I-7), making the binding to the verifying key local rather than
transitive through the points. The two $\GT$ bases are \emph{not} in the
preimage, and do not need to be: \yul{ec\_pairing} loads them from the VK payload,
whose length and \yul{extcodehash} are re-pinned on every proof at lines
1386--1389, so they are constants of the deployment. $\alpha$ is drawn strictly
\emph{after} all four points are materialized (lines 4033 and 1337 precede line
4066), so no point can be adapted to a known $\alpha$.
\emph{One precision note on the framing.} $\alpha$ is a deterministic, publicly
computable function of the points, not an independent uniform draw; the ``at most
one $\alpha$'' statement therefore has to be completed with a random-oracle
argument: modelling keccak as a random oracle, each candidate 4-tuple an
adversary tries yields a fresh $\alpha$, hitting the unique bad value with
probability $\approx 2^{-254.9}$ (the $\bmod\ r$ reduction of a 256-bit digest
is mildly non-uniform: $2^{256}/r \approx 2.21$, so a residue has either two or
three preimages and the per-value probability is off by at most a factor
$3/2.21 \approx 1.36$), so $Q$ hash queries succeed with probability
$\approx Q\,2^{-254.9}$.
The conclusion is unchanged; the comment is just terser than the argument.
\end{finding}

\begin{finding}[FIN-8]{The $\alpha = 0$ fixup is safe, and it is load-bearing (severity: Informational)}
Line 4067, \yul{if iszero(acc\_pair\_alpha) \{ acc\_pair\_alpha := 1 \}}, is not
cosmetic: with $\alpha = 0$ the accumulator terms vanish from both sides and
\eqref{eq:fin:batch} degenerates to the KZG equation alone --- the public
accumulator would be silently \emph{unchecked}. Forcing $\alpha \neq 0$ preserves
the ``at most one $\alpha$'' property in the $V \neq 0$ branch.
\emph{Is the deterministic fallback exploitable?} The fallback value $1$ is
public, and $\alpha = 1$ is exactly the coefficient of the cancellation attack
($U + V = 0$). But reaching the branch requires
$\operatorname{keccak256}(\text{frame}) \equiv 0 \pmod r$, an event of probability
$\approx 2^{-254}$ per attempt, and the adversary must satisfy it
\emph{simultaneously} with $U = -V$. Since the digest condition alone already
costs $\approx 2^{254}$ oracle queries, the fallback gives no shortcut over
guessing $\alpha$ directly. The construction is sound as written.
\emph{Recorded hardening option:} \texttt{docs/audit/AUDIT\_FINDINGS.md} tracks
this as S6 and suggests re-hashing with a fallback domain prefix instead of
substituting $1$; that would remove the ``known fallback value'' shape at the
cost of one keccak on a $2^{-254}$ branch. Not required for soundness.
\end{finding}

\subsection{Final pairing and the argument-order swap}
\label{sec:fin:swap}

\begin{lstlisting}[style=solidity,caption={\texttt{Halo2Verifier.sol} lines 4112--4128: the pairing dispatch, the terminal guard, and the return (lines 4115--4117 are blank or whitespace-only in the fixture)},label={lst:fin:final}]
// __phase:final_pairing
if iszero(success) { fail(ERR_PRECOMPILE_FAILED) }
success := ec_pairing(success, PAIRING_RHS_MPTR, PAIRING_LHS_MPTR)



// Success path is terminal. Invalid inputs have already reverted,
// so the Solidity ABI observes `true`.
//
// The guard is redundant today -- every failure path above reverts
// rather than clearing `success` -- but it keeps acceptance a local
// property of this file instead of an invariant split across
// FinalPairing.yul and ec_pairing.
if iszero(success) { fail(ERR_PROOF_REJECTED) }
// __phase:verifier_return
mstore(RETURN_MPTR, 1)
return(RETURN_MPTR, 0x20)
\end{lstlisting}

\yul{ec\_pairing} is defined at lines 955--990 and is specified elsewhere; its body
is reproduced here only to justify the argument order.

\begin{lstlisting}[style=solidity,caption={\texttt{Halo2Verifier.sol} lines 969--989 (reference; \yul{ec\_pairing} is specified elsewhere)},label={lst:fin:ecpairing}]
let scratch := 0x1240
mcopy(scratch,              lhs_mptr,                 0x80)
mcopy(add(scratch, 0x80),   G2_BASE_MPTR,             0x100)
mcopy(add(scratch, 0x180),  rhs_mptr,                 0x80)
mcopy(add(scratch, 0x200),  NEG_S_G2_BASE_MPTR,       0x100)
// MF-4: separate "the chain could not run the pairing" from
// "the pairing ran and rejected this proof". Both fail closed,
// but they are different incidents: the first points at the
// node/fork (a missing, repriced, or short-returning
// precompile), the second at the proof. Collapsing them into
// ProofRejected sent every responder looking at the wrong one.
ret := staticcall(PAIRING_GAS_2PAIR, 0x0f, scratch, 0x0300, scratch, 0x20)
ret := and(ret, eq(returndatasize(), 0x20))
if iszero(ret) { fail(ERR_PRECOMPILE_FAILED) }
// Compare against 1 rather than truncating to the low bit:
// `and(ret, word)` would accept any odd result word. EIP-2537
// only ever returns 0 or 1, so this matches the strict form
// the constructor smoke test already uses.
ret := eq(mload(scratch), 1)
if iszero(ret) { fail(ERR_PROOF_REJECTED) }
ret := 1
\end{lstlisting}

\yul{ec\_pairing(success, X, Y)} lays out
$[\,X \mid [1]_2 \mid Y \mid [-s]_2\,]$ and asserts
$\pair{X}{[1]_2}\cdot\pair{Y}{[-s]_2} = 1$, i.e.
\begin{equation}
  \pair{X}{[1]_2} \;=\; \pair{Y}{[s]_2}.
  \label{eq:fin:ecp}
\end{equation}

\paragraph{Why the call at line 4114 passes the slots swapped.}
The KZG single-opening identity for a commitment $C$ opened at $x_3$ to value $v$
with proof $\pi$ is $C - vG = (s - x_3)\,\pi$ in the exponent, hence
\begin{equation}
  \pair{C - vG + x_3\pi}{[1]_2} \;=\; \pair{\pi}{[s]_2}.
  \label{eq:fin:kzg}
\end{equation}
Matching~\eqref{eq:fin:kzg} against~\eqref{eq:fin:ecp} forces
$X = C - vG + x_3\pi$ and $Y = \pi$. Section~\ref{sec:fin:pairing} put
$C - vG + x_3\pi$ in \mptr{PAIRING\_RHS\_MPTR} and $\pi$ in
\mptr{PAIRING\_LHS\_MPTR}, so the call must be
\yul{ec\_pairing(success, PAIRING\_RHS\_MPTR, PAIRING\_LHS\_MPTR)} --- which is exactly
line 4114. Including the accumulator batching, the equation the contract
ultimately enforces is
\begin{equation}
  \pair{\;\mathsf{final\_com} - vG + x_3\pi + \alpha\,\mathsf{acc}_R\;}{[1]_2}
  \;=\;
  \pair{\;\pi + \alpha\,\mathsf{acc}_L\;}{[s]_2},
  \label{eq:fin:full}
\end{equation}
equivalently $\pair{\cdot}{[1]_2}\cdot\pair{\cdot}{[-s]_2} = 1$ as the comment at
lines 4045--4046 states, and equivalently $U + \alpha V = 0$ as in
Section~\ref{sec:fin:soundness}. The swap is therefore correct, and the
accumulator equation it batches in is
$\pair{\mathsf{acc}_R}{[1]_2} = \pair{\mathsf{acc}_L}{[s]_2}$.

\paragraph{The accumulator orientation agrees with the native verifier.}
Because $\mathsf{acc}_R$ is added into \mptr{PAIRING\_RHS\_MPTR} (line 4080,
the $[1]_2$ side) and $\mathsf{acc}_L$ into \mptr{PAIRING\_LHS\_MPTR} (line
4094, the $[s]_2$ side), the batched accumulator equation is
$\pair{\mathsf{acc}_L}{[s]_2} = \pair{\mathsf{acc}_R}{[1]_2}$. That is exactly
the native convention: \texttt{midnight-proofs}'
\texttt{DualMSM::check} (\texttt{proofs/src/poly/kzg/msm.rs:284--296})
evaluates \texttt{multi\_miller\_loop} over
$(\mathtt{left},\ \mathtt{s\_g2}),\ (\mathtt{right},\ \mathtt{n\_g2})$ and
requires the result to be one, i.e.
$\pair{\mathtt{left}}{[s]_2} = \pair{\mathtt{right}}{[1]_2}$, with
$\mathtt{left}$ the $\pi$/LHS accumulator and $\mathtt{right}$ the combined/RHS
one. So $\mathsf{acc}_L \leftrightarrow \mathtt{left}$ and
$\mathsf{acc}_R \leftrightarrow \mathtt{right}$, and the sides are \emph{not}
transposed by the batching.

\begin{finding}[FIN-9]{The \texttt{PAIRING\_LHS} / \texttt{PAIRING\_RHS} slots are named opposite to their pairing roles (severity: Observation)}
\mptr{PAIRING\_LHS\_MPTR} holds the point that is paired against $[s]_2$ (the
\emph{second} \yul{ec\_pairing} argument), and \mptr{PAIRING\_RHS\_MPTR} holds the
one paired against $[1]_2$ (the \emph{first}). The names follow the historical
dual-MSM accumulator convention (left $=\pi$, right $=$ combined), not the
pairing argument order (and it matches the native
\texttt{DualMSM::check}, Section~\ref{sec:fin:swap}). The code is correct and
the fixture spends thirteen comment lines (4099--4111) explaining the inversion,
plus six more in the generator
(\texttt{src/lowering/kzg/mod.rs:1978--1983}).
\emph{How bad is the hazard, precisely?} A future edit that ``fixes'' the
apparent swap at line 4114 would turn the check into
$\pair{\pi}{[1]_2} = \pair{C - vG + x_3\pi}{[s]_2}$. It is worth being exact
about the consequence, because the natural fear --- a silent soundness break ---
is \emph{not} what happens. On an honest proof the KZG relation gives
$C - vG + x_3\pi = s\,\pi$ in the exponent, so the swapped check would demand
$\pi = s^{2}\pi$, i.e.\ $(s^{2}-1)\log\pi = 0$, which fails for every
$\pi \neq \id$. The swap therefore breaks \emph{completeness} and is caught by
the first valid-proof test; it is not a change that ``passes the happy path and
opens a hole''. The finding stands as an auditability observation --- two
identifiers whose names say the opposite of what they hold, in the one place
where the whole check is assembled --- not as a latent soundness trap. Renaming
the slots to, e.g., \texttt{PAIRING\_G2\_BASE\_SIDE\_MPTR} /
\texttt{PAIRING\_S\_G2\_SIDE\_MPTR} would make the call site self-checking. No
defect in the current artifact.
\end{finding}

\subsection{Terminal check and return}
\label{sec:fin:return}

Three guards stand between the last precompile and \texttt{return}:

\begin{enumerate}
  \item Line 4113, \yul{if iszero(success) \{ fail(ERR\_PRECOMPILE\_FAILED) \}}.
  By Section~\ref{sec:fin:entry}, \yul{success} at this point can only have been
  cleared by one of the nine EIP-2537 \opcode{STATICCALL}s in this section
  (lines 3999, 4015, 4020, 4026, 4030, 4075, 4080, 4089, 4094) or by their
  \yul{returndatasize} checks, so \revert{PrecompileFailed} is the correct
  classification rather than a catch-all.
  \item Inside \yul{ec\_pairing}: \revert{ProofRejected} if the incoming flag is
  clear (line 962), \revert{PrecompileFailed} if the pairing \opcode{STATICCALL}
  fails or does not return exactly \texttt{0x20} bytes (line 982), and
  \revert{ProofRejected} if the returned word is not exactly $1$ (lines 987--988).
  The strict \yul{eq(mload(scratch), 1)} rather than a low-bit test is what makes
  a hypothetical precompile returning an odd non-one word fail closed.
  \item Line 4125, \yul{if iszero(success) \{ fail(ERR\_PROOF\_REJECTED) \}}.
  \yul{ec\_pairing} returns only the literal $1$ (line 989), so this guard is
  unreachable today; the source says so and keeps it deliberately, so that
  ``accept'' is a locally checkable property of this file rather than an
  invariant split across two partials.
\end{enumerate}

Then \yul{mstore(RETURN\_MPTR, 1)} and \yul{return(RETURN\_MPTR, 0x20)} --- one
word, ABI-encoding \texttt{true} for \yul{verifyProof(bytes,uint256[]) returns (bool)}.
\mptr{RETURN\_MPTR} $=$ \texttt{0x1000} is the same address the pairing scratch
and $\alpha$ frame used; it is written immediately before the \opcode{RETURN}, so
the reuse is a pure lifetime reuse.

\begin{finding}[FIN-10]{No path returns without a completed pairing check (severity: Informational, clean)}
\yul{return} appears exactly once in \texttt{Halo2Verifier.sol}, at line 4128, and
it is dominated by line 4114. \yul{ec\_pairing} has no exit that hands a
falsy value back to its caller: every non-accepting path inside it calls
\yul{fail(...)}, which reverts with a 4-byte selector, and the only fallthrough
sets \yul{ret := 1}. There is no early \yul{return}, no \yul{leave} that skips
the pairing, and no branch in lines 3812--4128 that reaches line 4127 without
having executed line 4114. The one structural hazard the generator guards against
elsewhere --- an empty point-set list, which would leave
\mptr{PAIRING\_LHS\_MPTR}/\mptr{PAIRING\_RHS\_MPTR} zero (both $\id$), making the
pairing trivially $1$ and the verifier accept-all --- is rejected at generation
time by an \yul{assert!} in \texttt{src/lowering/kzg/mod.rs:997--1010}; this
artifact has $n_{\mathrm{sets}} = 5$, so the block is emitted. This region is
sound.
\end{finding}

\begin{finding}[FIN-11]{\revert{PrecompileFailed} at line 4113 is not exclusively a chain-fault signal (severity: Observation)}
The constants commentary at lines 269--275 and the \yul{ec\_pairing} comment at
lines 974--979 both present the \revert{PrecompileFailed} / \revert{ProofRejected}
split (audit item MF-4) as separating ``the chain could not run the pairing''
from ``the pairing ran and rejected this proof''. Within the region assigned
here that separation is only \emph{almost} exact. Every EIP-2537 call site
forwards a fixed literal gas bound (lines 3999, 4015, 4020, 4026, 4030, 4075,
4080, 4089, 4094), but EIP-150 caps a \opcode{STATICCALL} at $63/64$ of the gas
remaining in the frame. If the caller supplies a transaction gas limit low
enough that fewer than $\frac{64}{63}\cdot 525{,}096$ gas remain at line 3999,
the precompile is entered with less than its scheduled cost, exhausts it,
returns $0$, and \yul{success} is cleared --- producing
\revert{PrecompileFailed} on a perfectly healthy chain with a perfectly valid
proof. The same applies to \yul{ec\_pairing}'s guard at line 982.
\emph{Consequence:} \revert{PrecompileFailed} means ``an EIP-2537 call did not
complete'', which covers (i) a repriced or absent precompile, (ii) a
short-returning implementation, and (iii) an under-funded caller. Only (i) and
(ii) are chain incidents. An operator triaging this selector must check the
transaction's gas limit before concluding the node or fork is at fault. No
soundness consequence --- all three cases fail closed --- and no code change is
implied; the entry in this specification and any runbook derived from the
comment should list the third cause.
\end{finding}

\subsection{Cost accounting}
\label{sec:fin:cost}

\begin{center}
\begin{tabular}{@{}llrl@{}}
\toprule
Line & Call & Forwarded gas & Purpose \\
\midrule
3999 & G1MSM, 78 pairs & $525{,}096$ & $\mathsf{final\_com}$ \\
4015 & G1MSM, 1 pair & $12{,}000$ & $(-v)G$ \\
4020 & G1ADD & $375$ & $+\,\mathsf{final\_com}$ \\
4026 & G1MSM, 1 pair & $12{,}000$ & $x_3\pi$ \\
4030 & G1ADD & $375$ & $+\,x_3\pi$ \\
4075 & G1MSM, 1 pair & $12{,}000$ & $\alpha\,\mathsf{acc}_R$ \\
4080 & G1ADD & $375$ & into \mptr{PAIRING\_RHS\_MPTR} \\
4089 & G1MSM, 1 pair & $12{,}000$ & $\alpha\,\mathsf{acc}_L$ \\
4094 & G1ADD & $375$ & into \mptr{PAIRING\_LHS\_MPTR} \\
4114 & PAIRING, 2 pairs & $102{,}900$ & $32600\cdot 2 + 37700$ \\
\midrule
\multicolumn{2}{@{}l}{total precompile} & $677{,}496$ & \\
\bottomrule
\end{tabular}
\end{center}

Non-precompile costs in the region are dominated by staging: $78$
\opcode{MCOPY}s of \texttt{0x80} bytes ($15$ gas each) plus $78$
\opcode{MSTORE}s, $\approx 1.4$k gas, and the memory expansion to
\texttt{0xe340} ($1818$ words, $\approx 11.9$k gas). The
\yul{keccak256} over \texttt{0x240} bytes costs $30 + 6\cdot 18 = 138$.

The generator's own note (\texttt{src/lowering/kzg/mod.rs:1967--1970}) records why
the three-point RHS is built as $2\times$G1MSM-1 $+\ 2\times$G1ADD rather than
one 3-pair G1MSM: the EIP-2537 discount only overtakes the add-based form at
larger $k$. The conclusion is right, but note that the two figures the comment
quotes are both wrong. It says ``3-pair G1MSM ($\approx 27.5$k gas) is more
expensive than $2 \times$ G1MSM-1 $+\ 2 \times$ G1ADD ($= 25$k gas)''. From the
schedule the generator itself implements, the true costs are
$3 \cdot 12000 \cdot d(3)/1000 = 3 \cdot 12000 \cdot 848/1000 = 30{,}528$ and
$2 \cdot 12000 + 2 \cdot 375 = 24{,}750$. The ordering --- and hence the design
choice --- is unaffected, and the margin is in fact larger than claimed
($5{,}778$ rather than $\approx 2{,}500$); only the numbers in the comment are
stale. This is the one place in the region where a comment's arithmetic does not
reproduce from the code beside it.

\subsection{Assessment summary}
\label{sec:fin:summary}

\begin{description}
\item[(a) Is $\alpha$ bound to everything it must be?] Yes. The preimage covers
all four $\GO$ inputs to the two batched equations in fixed-width form, plus
$\mathsf{vk\_digest}$ and a domain tag; the $\GT$ bases are deployment constants
re-pinned by codehash on every call. $\alpha$ is drawn only after all four points
are final. See FIN-7 for the one framing precision (deterministic $\alpha$
requires a random-oracle/grinding argument, which goes through).
\item[(b) Is the ``two bad equations could cancel'' argument correct?] Yes, and
the attack it describes is real, because the accumulator points are attacker-chosen
public inputs. The randomized form is the standard fix and is stated correctly:
if either residual is nonzero, at most one $\alpha$ satisfies the combined
equation (FIN-6).
\item[(c) Is the zero-$\alpha$ fixup safe?] Yes, and it is necessary: $\alpha = 0$
would silently drop the accumulator equation. The deterministic fallback value
$1$ is reachable only through a $2^{-254}$ digest event, so it offers no shortcut
over guessing $\alpha$ (FIN-8).
\item[(d) Is the LHS/RHS swap correct?] Yes; the derivation is in
Section~\ref{sec:fin:swap}. The naming is inverted relative to the pairing
argument order, which is a maintenance hazard rather than a defect (FIN-9).
\item[(e) Does any path return without a pairing check?] No. One \yul{return},
dominated by the \yul{ec\_pairing} call; every failure path reverts with a typed
selector (FIN-10).
\item[(f) Do the region's own comments hold up?] Mostly, with three
qualifications recorded above, all found by checking the code rather than the
prose next to it: the batching rationale at lines 4040--4049 is correct but
understates the transcript-ordering obstacle to the attack it rejects (FIN-6);
the \revert{PrecompileFailed} / \revert{ProofRejected} split advertised at lines
269--275 and 974--979 does not distinguish a chain fault from an under-funded
caller (FIN-11); and the gas comparison in
\texttt{src/lowering/kzg/mod.rs:1967--1970} quotes two figures
($\approx 27.5$k, $25$k) that neither matches the schedule the same file
implements ($30{,}528$, $24{,}750$), though its conclusion survives
(Section~\ref{sec:fin:cost}).
\end{description}

Findings raised in this section: FIN-1 (Informational, 128-bit truncated batching
scalars), FIN-2 (Informational, committed-instance eval constrained only
indirectly), FIN-3 (Observation, four absolute-literal staging destinations),
FIN-4 (Observation, bare gas literal), FIN-11 (Observation,
\revert{PrecompileFailed} also fires on caller gas starvation), FIN-5 / FIN-6 /
FIN-7 / FIN-8 / FIN-10 / FIN-12 (Informational, clean-region justifications),
FIN-9 (Observation, inverted slot naming). No Critical, High, Medium, or Low
finding was identified in lines 3812--4130. Two claims that a reader might
expect to see raised as findings are explicitly \emph{not} raised, because they
do not survive checking: the \mptr{PAIRING\_LHS}/\mptr{PAIRING\_RHS} naming
inversion is not a latent soundness trap (undoing it breaks completeness and any
valid-proof test catches it --- FIN-9), and the accumulator equation is not
folded in with its sides transposed (the orientation matches
\texttt{DualMSM::check} --- Section~\ref{sec:fin:swap}).


%% ==== 17-testing.tex ====================================================
\section{Adversarial testing of these artifacts}
\label{sec:testing}

\subsection{Method}

A specification says what the code should do; the results below say what it
actually did when handed inputs designed to break it. The suite lives at
\texttt{tests/moonlight\_wrap\_fixture\_adversarial.rs} (13 tests plus one
opt-in exhaustive sweep) and runs the fixture contracts under \texttt{revm} at
\texttt{SpecId::PRAGUE}, so the EIP-2537 precompiles are the real
implementations rather than stubs. The default set completes in about
40 seconds.

Every case asserts three properties, not one:

\begin{enumerate}
\item the input is \textbf{rejected};
\item it is rejected with the \textbf{exact typed error} the design assigns to
  that failure class --- a guard that fires but reports the wrong taxonomy is a
  real defect, because incident response is routed off these selectors;
\item it is rejected at the right \textbf{depth}, measured in gas.
\end{enumerate}

The third point is what makes the suite meaningful. Every mutation inside the
transcript region also breaks Fiat--Shamir, so a deleted encoding guard would
still cause a rejection --- just at the pairing, roughly $1.34\,$M gas later.
Anchoring on gas separates \emph{the guard under test fired} from \emph{the
algebra caught it anyway}. Two helpers encode this:
\texttt{reject\_early} requires the rejection to cost less than half of an
accepted run, and \texttt{reject\_at\_pairing} requires it to cost more than
three quarters.

The accept baseline is what gives every rejection meaning: the harness first
verifies the honest fixture calldata and records its gas
($1{,}336{,}972$). A suite that rejects everything would otherwise pass
vacuously.

\subsection{Coverage and results}

\begin{center}
\begin{tabularx}{\textwidth}{@{}X r l@{}}
\toprule
Attack class & Cases & Outcome \\
\midrule
ABI envelope: selector, both heads, both length words, \texttt{calldatasize}
  $\pm$, truncation & 26 & all rejected \\
Verifying-key binding: tampered header word, base point, VM constant, program
  word, commitment; empty account & 7 & construction reverts \\
Scalar canonicality: all $102$ evaluations, all $5$ $q$-evaluations, all $19$
  instances, at $r$, $r+1$, $2^{256}-1$ & 331 & all rejected early \\
G1 encoding: all $34$ proof points $\times$ both coordinates $\times$
  \{pad bytes $0,7,15$; coordinate $=p$; $=p+1$\} & 340 & all rejected early \\
Curve and subgroup: non-subgroup point, the $h$-torsion point $(0,2)$, an
  off-curve point, the generator, point negation, point at infinity, over
  $8$ proof categories & 48 & all rejected \\
Fiat--Shamir binding: every evaluation $+1$, every $q$-evaluation $+1$, every
  non-accumulator instance $+1$ and $:=0$, commitment swaps & 135 & all
  rejected at the pairing \\
Accumulator codec: over-packing of both word widths, shifted value $=p$ and
  $=2^{392}-1$, unflagged $(0,0)$, wrong-but-valid point, non-subgroup point
  & 26 & all rejected \\
Sentinel collision: the accumulator codec can express $(0,\pm 2)$, the only
  curve points with $x=0$, via the $p-1$ zero sentinel & 2 & rejected, but only
  by the subgroup check \\
Committed-instance evaluation: $e_0$ is structurally zero and enforced only by
  the PCS batching & 5 & all rejected at the pairing \\
Cross-artifact: the IVC fixture's calldata, an all-zero proof body, all-zero
  instances & 3 & all rejected \\
Deployability: runtime sizes, pinned VK length, pinned VK codehash & 4 & all
  hold \\
\midrule
\textbf{Exhaustive single-bit sweep}: every one of the $8{,}516$ calldata
  bytes, low bit and high bit & \textbf{17{,}032} & \textbf{none accepted} \\
\bottomrule
\end{tabularx}
\end{center}

No input in any class was accepted. Nor did any input cause
\texttt{verifyProof} to \emph{return} \texttt{false}: across all 17{,}032
mutations plus the targeted cases, every non-accepting outcome was a revert.
That is the empirical form of the NatSpec's success-or-revert claim, and it is
worth stating because an integrator writing
\yul{if (!verifier.verifyProof(...))} would be writing dead code. The single-bit sweep is the strongest
statement the harness can make about completeness: it assumes nothing about
where a mutation lands and reaches regions no named test targets --- interior
bytes of coordinate low words, the middle of the evaluation block, the padding
inside the ABI length words.

\subsection{Where the rejections come from}
\label{sec:testing:taxonomy}

The exhaustive sweep also measures which layer catches what. This is the part
of the result that is not a simple pass:

\begin{center}
\begin{tabular}{@{}r r l l@{}}
\toprule
Count & Share & Selector & Error \\
\midrule
7{,}433 & 43.6\,\% & \texttt{0xc3b0d8cd} & \revert{ProofRejected} (reached the pairing) \\
6{,}840 & 40.2\,\% & \texttt{0x84e81692} & \revert{PrecompileFailed} \\
2{,}376 & 14.0\,\% & \texttt{0xf27905ec} & \revert{BadPointEncoding} \\
119 & 0.7\,\% & \texttt{0x77530042} & \revert{NonCanonicalScalar} \\
4 & 0.0\,\% & \texttt{0x1b99e37c} & \revert{BadCalldataShape} \\
260 & 1.5\,\% & --- & bare revert from solc's ABI decoder \\
\bottomrule
\end{tabular}
\end{center}

Two things stand out.

\begin{finding}[T-1]{\revert{PrecompileFailed} is dominated by proof faults, not chain faults --- Low}
Forty percent of all rejections carry \revert{PrecompileFailed}. The contract
documents that selector as a chain-level fault: the MF-4 comment at
\texttt{Halo2Verifier.sol:1240--1245} says \yul{precompile\_failed} separates
``a G1MSM staticcall that could not run (chain/gas fault)'' from ``a
public-input point this verifier decoded and rejected \dots\ or a point the
precompile found off-curve/out-of-subgroup'', and states that ``only the second
is a \revert{BadPointEncoding}''.

The code cannot make that distinction. An off-curve point and a missing
precompile both surface as a reverted \yul{staticcall}, both set
\yul{precompile\_failed}, and the call site at
\texttt{Halo2Verifier.sol:1447--1451} reports \revert{PrecompileFailed} for
both. Since the verifier performs no on-curve check of its own, \emph{every}
malformed point --- the overwhelmingly common attacker input --- lands in the
selector reserved for node problems.

\textbf{Impact}: none on soundness; every such input is still rejected. The cost
is operational: a monitoring rule that pages on \revert{PrecompileFailed} will
fire on ordinary garbage traffic, and a real precompile repricing will be lost
in that noise.

\textbf{Recommendation}: either (a) perform an explicit on-curve test on
decoded accumulator coordinates before the G1MSM round-trip and report
\revert{BadPointEncoding} when it fails, leaving \revert{PrecompileFailed} for
the residual case; or (b) correct the comment and the incident-response
document to say that \revert{PrecompileFailed} means ``the precompile rejected
this input \emph{or} could not run'', and route monitoring accordingly. Option
(b) is free; option (a) costs a few hundred gas on the accumulator path and
makes the taxonomy mean what it says.
\end{finding}

\begin{finding}[T-2]{The typed \revert{BadCalldataShape} guard is largely shadowed by solc's ABI decoder --- Informational}
Of the eight dynamic-argument-head mutations exercised, seven are stopped by
solc's own calldata decoder with a bare \texttt{revert(0, 0)} and no selector,
before the generated guard at \texttt{Halo2Verifier.sol:636--642} runs; only
one reaches it. In the exhaustive sweep, \revert{BadCalldataShape} accounts for
4 rejections out of 17{,}032, and bare reverts for 260.

This is not a defect --- fail-closed is preserved either way, and the generated
guard is still the only thing that catches the cases solc cannot see (trailing
calldata bytes, a wrong instance count that is still in bounds). But an
integrator who expects malformed envelopes to be identifiable by selector will
find most of them indistinguishable from any other empty revert.

\textbf{Recommendation}: document in the \texttt{verifyProof} NatSpec that
envelope faults may surface either as \revert{BadCalldataShape} or as an
untyped revert from the ABI decoder, so callers do not build error handling on
the assumption that all rejections are typed.
\end{finding}

\subsection{How the suite avoids drifting from the artifact}
\label{sec:testing:anchoring}

A replay suite that carries its own copy of the proof layout stops testing the
contract the moment the contract is re-rendered: the offsets still tile, the
sweeps still run, and every assertion still passes --- against the wrong
addresses. \texttt{pinned\_layout\_matches\_fixture\_source} is what prevents
that. It re-derives, from the artifact source rather than from constants in the
test file:

\begin{itemize}
\item the ABI cursors, the proof length, the instance count and the exact
  \texttt{calldatasize} equality, by string match against the rendered guards;
\item the eight typed-error selector constants;
\item the six verifying-key header words the verifier cross-checks
  ($19$, $20$, $1$, $11$, $7$, $56$), read back out of
  \texttt{Halo2VerifyingKey.sol};
\item the two field moduli the test needs, matched against
  \mptr{FR\_MODULUS}, \mptr{BLS\_P\_HI} and \mptr{BLS\_P\_MINUS\_ONE\_LO};
\item and --- the strongest of them --- the section table itself, against the
  parser's own loop bounds. The verifier walks the proof with
  \yul{for \{ let end := add(proof\_cptr, N) \}} per contiguous run; the test
  extracts every such $N$ in source order and asserts the exact sequence
  $[15,2,6,1,1,1,4]\cdot\texttt{0x80}$ followed by
  $102\cdot\texttt{0x20}$ and $5\cdot\texttt{0x20}$.
\end{itemize}

The remaining hardcoded values are the two curve moduli (anchored as above) and
the section \emph{names}, which carry no semantics. Every adversarial vector is
constructed rather than pasted: the non-subgroup point is found by search and
self-checked with \yul{on\_curve}/\yul{in\_subgroup}, the identity-flagged
accumulator word is derived as the packed $p-1$ word plus one radix base and
then matched against the verifier's own
\mptr{BLS\_P\_MINUS\_ONE\_PACKED\_0\_WITH\_ID\_FLAG}, and the over-packing
vectors are built from the declared limb widths.

\subsection{What the suite does not cover}

Three limits are worth stating explicitly, because a green suite invites
over-reading.

\begin{itemize}
\item \textbf{No forged-proof test.} Every case here mutates an honest proof.
  Building a proof that is internally consistent but false requires a prover,
  and the properties that would catch it --- the quotient identity, the
  permutation and lookup arguments --- are exercised only in the sense that
  every mutation reaches them. Soundness against a malicious prover rests on
  the differential anchors in \texttt{src/lowering/} that pin the generated
  arithmetic value-by-value against the native Rust verifier, not on this file.
\item \textbf{EIP-170 is lifted in the harness.} \texttt{src/evm.rs} sets
  \texttt{limit\_contract\_code\_size = usize::MAX}, so every test would pass
  for an undeployable artifact. That is why the sizes are measured explicitly
  in \texttt{fixture\_artifacts\_fit\_in\_a\_deployable\_contract}.
\item \textbf{One modexp schedule.} The pinned \texttt{revm 19} prices
  \texttt{0x05} by EIP-2565 even under \texttt{SpecId::OSAKA}, so no test in
  this repository exercises a chain that prices modexp by EIP-7883. The
  contract's response is to forward the larger EIP-7883 bound everywhere and
  probe it at deployment; that reasoning is sound but untested here.
\end{itemize}


%% ==== 18a-assessment.tex ================================================
\section{Consolidated assessment}
\label{sec:assess:summary}

\subsection{Method}

Twelve independent reviews were run over the two artifacts --- eight on
security dimensions (Fiat--Shamir soundness, KZG multi-open and pairing,
quotient VM integrity, the public-input codec, the memory model, precompile
usage, the trust model, and control flow) and four on code quality
(readability, maintainability, correctness hygiene, documentation accuracy).
Every finding was then handed to a second reviewer whose instructions were to
\emph{refute} it: re-read the cited code, show the claim is wrong, the path
unreachable, the arithmetic sound, or the severity inflated, and default to
refuted when a claim could not be substantiated. Several of the load-bearing
results were re-derived from scratch during that pass rather than re-read ---
the quotient program was decoded independently out of the verifying-key
contract, the final-MSM term census was re-extracted, and the memory map was
re-checked pairwise.

\begin{center}
\begin{tabular}{@{}l r l r@{}}
\toprule
Verdict after refutation & Count & Final severity & Count \\
\midrule
Confirmed & 49 & Critical & 0 \\
Partial (facts hold, severity trimmed) & 36 & High & 0 \\
Added by the refuting pass & 26 & Medium & 4 \\
Refuted & 2 & Low & 23 \\
& & Informational & 50 \\
& & Observation & 29 \\
\midrule
Total & 113 & & \\
\bottomrule
\end{tabular}
\end{center}

\subsection{Verdict}
\label{sec:assess:verdict}

\textbf{No soundness defect was found, and several of the properties that
matter most were positively established rather than merely not-falsified.}
The results below survived a deliberate attempt to break them.

\begin{enumerate}
\item \textbf{Acceptance implies a successful pairing.} \yul{return} appears
  exactly once in the artifact (line 4128). Reaching it requires passing
  \yul{if iszero(success) \{ fail(ERR\_PROOF\_REJECTED) \}} at 4124, and
  \yul{success} is last written at 4118 by \yul{ec\_pairing}, whose only
  non-reverting exit is gated on the staticcall succeeding \emph{and}
  \yul{returndatasize() == 0x20} \emph{and} the strict \yul{eq(mload(scratch), 1)}.
  No path returns \texttt{true} without the \texttt{0x0f} precompile having
  returned the word $1$.
\item \textbf{No live memory overlap.} All 30 registered regions were checked
  pairwise with explicit arithmetic, including every bare literal used as a
  memory operand. Every shared base (\texttt{0x1000}, \texttt{0xb140},
  \texttt{0xb280}) is a phase-disjoint reuse with a first-write/last-read proof.
  No memory address anywhere in the file is derived from calldata, and no
  literal write destination falls below \texttt{0x1000} except the five
  error-selector stores in Solidity's own scratch.
\item \textbf{Every absorbed commitment is validated.} All 34 proof \GO{}
  elements were mechanically traced from their transcript absorption to an
  EIP-2537 precompile input: 34 of 34. There is no absorbed-but-unvalidated
  commitment, so the documented decision to defer curve and subgroup checks to
  the precompiles is discharged rather than merely asserted.
\item \textbf{The multi-open closes exactly.} The evaluation slots form a
  bijection (103 slots, 103 distinct, set-equal to
  \texttt{0x9480}--\texttt{0xa120} with no gap or duplicate, plus
  \mptr{LINEARIZATION\_EVAL\_MPTR}), and the fused MSM covers indices $0..77$ at
  stride \texttt{0xa0} with 78 distinct point sources filling all
  \texttt{0x30c0} bytes. The $x_1$, $x_2$ and $x_4$ batching axes are
  independent and consistent between the evaluation and commitment sides.
\item \textbf{The pinned quotient program is well-formed.} Decoded
  independently from \texttt{Halo2VerifyingKey.sol}: 532 \yul{mstore} literals
  reconstructing a 17{,}024-byte payload with no gaps; the program at payload
  offset \texttt{0x1a20}, length \texttt{0x11cf}, terminating with
  \yul{q\_pc == q\_end} exactly; 101 instructions; an empty final stack peaking
  at one spilled word; 41 distinct operand pointers, all word-aligned and all
  inside \texttt{0x94a0}--\texttt{0x9aa0}; and constant indices exactly
  $0..177$ against a 178-entry table.
\item \textbf{Every precompile bound is correct.} Independently re-derived:
  G1ADD $375$, one-pair G1MSM $12{,}000$, pairing $32{,}600\cdot2+37{,}700 =
  102{,}900$, the 78-pair MSM $78\cdot12{,}000\cdot561/1000 = 525{,}096$, and
  modexp $4{,}080 \ge$ both live schedules. Twenty-three of the 24 staticcalls
  assert \yul{returndatasize}; the sole exception is the negative probe at line
  500, which expects failure. No loop bound is attacker-controlled.
\end{enumerate}

The findings that follow are therefore about hygiene, taxonomy, documentation
and provenance --- not about whether an invalid proof can be accepted.

\subsection{Findings at Medium and Low severity}
\label{sec:assess:table}

Informational and Observation items are summarised by theme in
\S\ref{sec:assess:themes} rather than tabulated. ``Partial'' means the facts
were reproduced but the severity was reduced by the refuting pass.

\begingroup
\footnotesize
\setlength{\tabcolsep}{4pt}
\begin{longtable}{@{}l l l p{0.50\textwidth} l@{}}
\toprule
ID & Severity & Area & Finding & Verdict \\
\midrule
\endfirsthead
\toprule
ID & Severity & Area & Finding & Verdict \\
\midrule
\endhead
\texttt{ADD-1} & Medium & maintainability & The P14 staleness gate is currently RED: 41b78eb9 changed five templates without re-pinning the digest & Added \\
\addlinespace
\texttt{ADD-2} & Medium & maintainability & fixtures/ivc/Halo2Verifier.sol is provably rendered from a superseded template tree & Added \\
\addlinespace
\texttt{MD-1} & Medium & maintainability & Fixture provenance stamp names a commit that cannot produce this artifact & Confirmed \\
\addlinespace
\texttt{R-3} & Medium & readability & 73 of 74 point sources in the fused final MSM are bare addresses, 45 of them unlabeled VK-interior offsets, while the naming data exists in the VK… & Confirmed \\
\addlinespace
\texttt{KZG-A1} & Low & PCS & The reference spec's alpha preimage omits vk\_digest entirely — a reimplementation built from the spec derives a different alpha and a shorter frame & Added \\
\addlinespace
\texttt{ACC-1} & Low & accumulator & The guard documented as "unreachable by construction" is live and load-bearing; deleting it admits a family of non-canonical identity encodings & Partial \\
\addlinespace
\texttt{ADD-1} & Low & documentation & The constructor's VK check reverts Error(string) "invalid vk", falsifying the taxonomy comment's claim that the memory-layout guard is the one… & Added \\
\addlinespace
\texttt{DOC-1} & Low & documentation & Header's EIP-170 / optimiser paragraph is stale and misdirects on the one build parameter that decides deployability & Confirmed \\
\addlinespace
\texttt{DOC-3} & Low & documentation & verifyProof's failure-taxonomy NatSpec is incomplete in two independent ways: it omits MemoryLayoutViolated and it over-promises that every… & Confirmed \\
\addlinespace
\texttt{CF-1} & Low & flow & The nine-staticcall PCS/pairing region has no precompile\_failed discriminator at all, so line 4113 attributes every off-curve proof commitment to a… & Confirmed \\
\addlinespace
\texttt{ADD-01} & Low & hygiene & ERR\_QUOTIENT\_PROGRAM\_INVALID is declared and never used; q\_program\_fail hardcodes the raw selector and open-codes fail()'s body & Added \\
\addlinespace
\texttt{DEAD-01} & Low & hygiene & The `if is\_id` branch in load\_acc\_point is documented "unreachable by construction" but is reachable and is the sole guard for x-flagged /… & Confirmed \\
\addlinespace
\texttt{DEAD-02} & Low & hygiene & X\_N\_MPTR and X\_N\_MINUS\_1\_INV\_MPTR are dead stores; x\^{}n is then recomputed with a redundant 20-squaring walk & Confirmed \\
\addlinespace
\texttt{MD-2} & Low & maintainability & The staleness digest hashes templates/ only; most emitted Yul comes from src/lowering/ & Partial \\
\addlinespace
\texttt{MD-5} & Low & maintainability & The error-selector pin checks name and selector independently, so a swapped assignment would pass & Confirmed \\
\addlinespace
\texttt{MD-6} & Low & maintainability & --optimize-runs is load-bearing for EIP-170 but is not recorded in the artifact, and the header's cited measurements do not cover the pinned pair & Partial \\
\addlinespace
\texttt{ADD-1} & Low & operations & The deployment script emits a four-field record that omits every field needed to reproduce or re-identify the deployed bytecode & Added \\
\addlinespace
\texttt{TMO-1} & Low & operations & BUILD\_ID identifies the verifying key, not the verifier — verifier-side drift produces an identical build id & Partial \\
\addlinespace
\texttt{TMO-3} & Low & operations & The pragma-pin rationale states measurements that contradict the artifact and a deployability claim that is false & Confirmed \\
\addlinespace
\texttt{QVM-A1} & Low & quotient VM & Native callback entry silently discards a live cached top instead of rejecting it; the guard covers only the spill stack, and both the rendered… & Added \\
\addlinespace
\texttt{A-1} & Low & readability & QUOTIENT\_RETURN\_MPTR is a dead constant whose only mention in the artifact is a header comment describing it as a live phase in a lifetime argument & Added \\
\addlinespace
\texttt{A-2} & Low & readability & 0xb140 carries four distinct roles -- two named, two bare literals -- and is the densest alias in the artifact, denser than either address R-1 reports & Added \\
\addlinespace
\texttt{A-3} & Low & readability & The 128-bit truncation mask is an undeclared literal that means two unrelated things, and carries no comment at the six sites where it truncates… & Added \\
\addlinespace
\texttt{R-1} & Low & readability & Addresses 0xb280 and 0xb2c0 each serve several distinct roles, only one is named, and the header's alias-disclosure list omits both reuses & Partial \\
\addlinespace
\texttt{R-10} & Low & readability & The quotient VM's inline-prefix blocks are anonymous var0..var162 chains with no per-identity label & Confirmed \\
\addlinespace
\texttt{R-2} & Low & readability & The header's final-MSM addressing policy is violated at exactly the four quotient-limb terms, and with the symbolic/absolute roles inverted & Partial \\
\addlinespace
\texttt{R-6} & Low & readability & The quotient VM's operand-window and stack-bound guards are written with undeclared magic numbers repeated ten times & Confirmed \\\\
\bottomrule
\end{longtable}
\endgroup

\subsection{The four Medium findings}
\label{sec:assess:medium}

Three of the four are about \emph{provenance}: they say that the tree from
which these artifacts came can no longer be identified. That matters more for a
generated verifier than it would for hand-written code, because the only
practical way to check a 4{,}130-line generated file is to re-render it and
diff.

\begin{finding}[ADD-1]{The repository's own staleness gate is currently failing --- Medium}
\label{finding:assess:digest}
\texttt{tests/template\_digest.rs::template\_tree\_digest\_is\_pinned} fails on
the current \texttt{HEAD}. Reproduced directly:
\begin{lstlisting}[style=solidity,caption={\texttt{cargo test --test template\_digest}, run against \texttt{HEAD} (\texttt{41b78eb9})},label={lst:assess:digest}]
The templates/ tree changed (digest 0xb22bc20e0263f0f1bedd07ded7346703946a0bad2d1a160d46bc0012803b517a).
Every committed fixture and dump was rendered from the previous tree and is now stale. ...
  left: "0xb22bc20e0263f0f1bedd07ded7346703946a0bad2d1a160d46bc0012803b517a"
 right: "0xf0a67125fae3f11061cd1b6582088ca5d925ab33c6e0bac462aec8af099c59da"
\end{lstlisting}
This is the single automated drift gate in the package, and it fired on the most
recent codegen commit and was left red. Everything the gate is meant to protect
--- ``every committed fixture and dump was rendered from the previous tree and
is now stale'' --- is therefore currently unverified.

\textbf{Recommendation}: resolve it in one commit rather than by re-pinning the
digest alone (re-pinning silently discharges steps 1 and 3 of the message's own
checklist): regenerate what can be regenerated, fix the source-commit stamps,
re-run the replay tests, \emph{then} pin. Longer term, make the digest update
part of the same commit as any \texttt{templates/} edit so the gate cannot be
left red across a merge.
\end{finding}

\begin{finding}[MD-1]{The fixture's provenance stamp names a commit that cannot have produced it --- Medium}
\label{finding:assess:stamp}
\texttt{fixtures/moonlight-wrap/README.md}:17 records source commit
\texttt{6e83e8e4}. That commit's tree contains no \mptr{PCS\_FINAL\_MSM\_MPTR}
anywhere in \texttt{src/} or \texttt{templates/}, yet the artifact declares it
at line 204 and uses it throughout lines 3812--4007.
\texttt{docs/reference/REPRODUCIBLE\_BUILDS.md}:39 independently records a
different commit (\texttt{3fb6d84}) for the same directory. The stamp is
readable as ``the parent of the committing commit'', a convention that
re-derives only when the render commit is fixture-only --- which
\texttt{41b78eb9} was not.

The fixture README states that staleness is ``tracked by the commit stamp
rather than an assertion'', so the stamp is the whole mechanism, and it does not
currently work.

\textbf{Recommendation}: stamp the commit whose tree actually rendered the
artifact, make the two documents agree, and replace the human step with a CI
check --- \texttt{git log -1 --format=\%h -- fixtures/<name>/Halo2Verifier.sol}
needs neither the SRS nor a Moonlight checkout. Better still, have the generator
emit the source commit into the artifact header next to \texttt{BUILD\_ID}, so
the stamp travels with the file it describes.
\end{finding}

\begin{finding}[ADD-2]{The sibling IVC fixture is provably from a superseded template tree --- Medium}
\label{finding:assess:ivc}
\texttt{fixtures/ivc/Halo2Verifier.sol} declares constants
(\mptr{QUOTIENT\_EVAL\_MPTR}, \mptr{Q\_COM\_MPTR}) that no longer exist anywhere
in \texttt{templates/}. It cannot be produced by the current generator, and
nothing in the repository says so. The two committed fixtures therefore come
from different generator trees, and ``the fixtures'' is no longer a single
snapshot of codegen --- which is exactly the property the replay tests rely on
when they treat a fixture as representative.

\textbf{Recommendation}: regenerate, or restore an explicit staleness banner
naming the two constants that prove it. A cheap machine-checkable guard in the
meantime: assert that no committed fixture \texttt{.sol} contains an identifier
absent from \texttt{templates/}. It needs no SRS, no solc and no prover, and it
would have caught this.
\end{finding}

\begin{finding}[R-3]{The block that decides which commitments are checked is the least labelled block in the file --- Medium (auditability)}
\label{finding:assess:msm}
Of the 78 point sources in the fused final MSM (lines 3812--4007), 73 are bare
hex addresses and exactly one is symbolic (\mptr{F\_COM\_MPTR}, line 3995).
Forty-five of the bare ones fall inside the verifying-key window
$[\texttt{0x3680}, \texttt{0x7900})$ --- which matches the verifying-key
contract's own census of $27$ \texttt{fixed\_comms} $+$ $18$
\texttt{permutation\_comms} $= 45$ exactly. Six more collide with declared
constants (\mptr{ADVICE\_COMMS\_MPTR\_BASE} and its five siblings) and are still
emitted as literals. The block carries seven comment lines in 196 lines of code,
and the three that precede the term list give only the phase name, the purpose
and the term count.

So the one block that determines \emph{which} commitments enter the pairing ---
the block where a dropped or duplicated term would be a soundness break --- is
the block a reviewer can least check by reading. The label data exists: the
generator already emits \texttt{// fixed\_comms[4].x\_hi} into the verifying-key
contract for the very word at \texttt{0x6480} that line 3840 copies from.

\textbf{Recommendation}: two independent, zero-bytecode-cost changes. Emit
\yul{add(<BASE>\_MPTR\_BASE, <offset>)} wherever a declared constant already
matches, and append a trailing comment to each of the 78 term pairs naming the
term (\texttt{// term 22: VK fixed\_comms[4] * x1\^{}22}). Comments do not
appear in bytecode, so neither change touches the EIP-170 headroom.
\end{finding}

\subsection{Low-severity findings worth acting on first}
\label{sec:assess:low}

\begin{finding}[ACC-1 / DEAD-01]{A branch documented as unreachable is live and is the sole guard for a family of malformed identity encodings --- Low}
\label{finding:assess:acc1}
\yul{load\_acc\_point} computes \yul{is\_id} from the four-word sentinel at line
1157, then \emph{reassigns} it to \yul{x\_is\_id} at line 1187 ---
and \yul{x\_is\_id} is decided by a probe that reads only $x$'s two words. So an
$x$ equal to the identity-flagged sentinel with \emph{any} $y$ enters the
\yul{if is\_id} branch at line 1189, and line 1201,
\yul{ok := and(ok, iszero(or(or(x\_hi, x\_lo), or(y\_hi, y\_lo))))}, is the only
thing that rejects it. The comment at lines 1194--1200 says the branch is
``unreachable by construction \dots{} kept as defence in depth \dots{} rather
than as a live branch'', and invites a future author to delete it. Deleting it
would admit identity-flagged $x$ with arbitrary $y$ --- i.e.\ reopen accumulator
identity malleability.

The difference is arithmetically exact:
\mptr{BLS\_P\_MINUS\_ONE\_PACKED\_0\_WITH\_ID\_FLAG} minus
\mptr{BLS\_P\_MINUS\_ONE\_PACKED\_0} is exactly $2^{56}$, the radix base, so the
flag probe fires on $x$ alone while the sentinel check also constrains $y$.

\textbf{Status}: now pinned by an executing test. The vector (identity-flagged
$x$, $y = 2$) is rejected with \revert{BadPointEncoding} before the transcript;
see \texttt{fixture\_verifier\_validates\_the\_public\_accumulator\_codec} case
(c2). \textbf{Recommendation}: replace the comment with the true statement, and
correct audit note I-3 in
\texttt{docs/audit/HALO2\_VERIFIER\_REVIEW\_2026-08.md}:657, which is where the
false claim originates.
\end{finding}

\begin{finding}[CF-1]{The whole PCS and pairing region attributes every failure to a chain fault --- Low}
\label{finding:assess:cf1}
Nine \yul{staticcall}s between lines 3997 and 4096 write \yul{success}, and
line 4113 is their sole consumer:
\yul{if iszero(success) \{ fail(ERR\_PRECOMPILE\_FAILED) \}}, with no
discriminator variable anywhere in the region. Every earlier \yul{and}-fold is
consumed before that point (at 1406, 1430, 1448, 1523, 1877), so
\yul{success == 0} at 4113 provably means one of the nine EC precompile calls
returned $0$ --- which is either a chain fault \emph{or} an attacker-supplied
point that EIP-2537 rejected. The selector is unconditionally the former. Since
\yul{common\_uncompressed\_g1} checks only the 16-byte pad and the two field
bounds, an off-curve proof commitment reaches the MSM and lands here.

This is the same defect as the accumulator-side MF-4 case, one region larger,
and together they explain the measured 40\,\% share of
\revert{PrecompileFailed} (\S\ref{sec:testing:taxonomy}).

\textbf{Recommendation}: on the failure path only --- which reverts anyway, so
cost is irrelevant --- run a liveness probe before choosing the selector:
re-issue a known-good one-pair G1MSM over the verifying key's own
\mptr{G1\_BASE\_MPTR} with scalar $1$. If the probe succeeds, the precompile is
live and the earlier failure came from the input, so raise
\revert{BadPointEncoding}; if it also fails, raise \revert{PrecompileFailed}.
Applying the same probe at lines 1288, 1337 and 1451 fixes the confirmed MF-4
defect by construction rather than by comment.
\end{finding}

\begin{finding}[DEAD-02]{Two named memory slots are written and never read, and $x^n$ is computed twice --- Low}
\label{finding:assess:dead2}
A full-file grep gives \mptr{X\_N\_MPTR} at lines 170 (declaration) and 1869
(store) only, and \mptr{X\_N\_MINUS\_1\_INV\_MPTR} at 171 and 1870 only --- no
\yul{mload}, no \yul{mcopy} source, and no appearance in the eval-source address
tables. Both are dead stores. The linearization block then recomputes $x^n$ with
a second 20-squaring walk at lines 3383--3396.

The gas cost is negligible (about 20 \opcode{MULMOD}s, roughly $0.015\,\%$ of a
verification). The real cost is that two independent computations of the same
value now have to be kept in agreement by hand, and a named slot that nothing
reads invites a reader to assume it is consumed somewhere.

\textbf{Recommendation}: in the linearization block, drop \yul{x\_pow\_2i} and
use \yul{addmod(1, sub(r, mload(X\_N\_MPTR)), r)}. This is safe --- the Lagrange
block at 1869 runs unconditionally first, and no literal write targets
\texttt{0x7c40} in between. Then either drop
\mptr{X\_N\_MINUS\_1\_INV\_MPTR} or add a generator assertion that every
declared \mptr{*\_MPTR} is either read somewhere or explicitly registered as
reserved.
\end{finding}

\begin{finding}[QVM-A1]{Native callbacks discard a live cached top instead of rejecting it --- Low}
\label{finding:assess:qvm}
All three native-callback arms (lines 2453--2464, 2570--2580, 2762--2772) clear
\yul{q\_has\_top} unconditionally and then assert only that the spill pointer is
at its base. A fully evaluated identity sitting in the cached top with
\yul{q\_sp == 0xb8e0} --- the ordinary single-value state the VM spends most of
its time in --- would therefore be dropped from $\nu_y(x)$ with no revert, and
the epilogue triple at 3098--3105 would still pass. The rendered comment at
2455--2457 claims the arm ensures the callback does ``not inherit any partially
evaluated VM stack state'', which this check does not deliver.

No exploit follows: an independent decode of the pinned program shows every
native callback is preceded by a \opcode{FOLD} or another native, so
\yul{q\_has\_top} is provably $0$ at all six dynamic entry points. The defect is
that the guard does not enforce what it claims, in the one place where a
generator bug would silently weaken $\nu_y(x)$ rather than break it.

\textbf{Recommendation}: assert before clearing ---
\yul{if q\_has\_top \{ q\_program\_fail() \}} ahead of the existing
\yul{q\_sp} check. Cost is three executions of about nine gas.
\end{finding}

\begin{finding}[MD-5]{The test cited as the pin for the error selectors does not pin them --- Low}
\label{finding:assess:md5}
\texttt{Halo2Verifier.sol}:314--317 cites
\texttt{p4\_error\_selectors\_match\_declared\_errors} as the pin for the eight
typed-error selectors. That test
(\texttt{src/lowering/tests.rs}:2929--2938) asserts that the selector appears in
the template \emph{and}, as a wholly separate statement, that the constant name
appears --- neither binds the pair. Since the template puts each name and value
on its own line, swapping \texttt{0xa447d73e} and \texttt{0x77530042} between
\mptr{ERR\_VK\_MISMATCH} and \mptr{ERR\_NON\_CANONICAL\_SCALAR} satisfies both
assertions. The neighbouring test at lines 2948--2951 \emph{does} bind name to
value, which shows the stronger form was available.

All eight selectors were recomputed independently for this document and are
correct in the current artifact, so this is a defect in the guarantee, not in
the values.

\textbf{Recommendation}: assert on the rendered pair, or parse the constants
block into a name-to-value map and compare against recomputed
\texttt{keccak256}. Two lines, cannot break a correct render.
\end{finding}

\subsection{The Informational and Observation findings, by theme}
\label{sec:assess:themes}

The remaining 79 items cluster into six themes. None changes what the contract
accepts.

\begin{center}
\begin{tabularx}{\textwidth}{@{}l r X@{}}
\toprule
Theme & Items & Character \\
\midrule
Protocol parameters worth documenting & 12 &
The \texttt{mod r} sampling bias (max ratio $1.36$, statistical distance
$2^{-3.74}$), the 128-bit truncation of $x_3$ and the $x_1$/$x_4$ power tables
(PCS soundness $\approx 2^{-107.6}$ at $k=20$), and the fact that the mask width
is a property of the scalar field rather than a tunable parameter. All are
forced by native compatibility; none is a defect; none is in the header. \\
\addlinespace
Error-taxonomy precision & 14 &
Guards that raise a plausible-but-wrong selector: \yul{scalar\_inv} reporting
\revert{NonCanonicalScalar} when $x_3$ lands on a rotation point; the
constructor's \texttt{Error(string)} ``invalid vk'' outside the typed taxonomy;
\revert{VkMismatch} conflating wrong-code-at-address with generator drift;
\revert{PrecompileFailed} reachable by caller under-gassing. \\
\addlinespace
Dead or unreachable code & 11 &
\mptr{QUOTIENT\_RETURN\_MPTR}, \mptr{G1\_IDENTITY\_MPTR} and
\mptr{TRACE\_U256\_MPTR} with no reader in this render; the redundant
$\yul{lt(hi, 2^{128})}$ check; the final \yul{if iszero(success)}; the
\yul{batch\_invert} reversed-range guard; \yul{pi}'s transcript copy after the
last squeeze. \\
\addlinespace
Comment accuracy & 15 &
Claims that overstate the code: the MF-4 taxonomy promise, the header's stale
size figures, the native-callback stack claim, the modexp price comment
($1354/4064$ actual against $1360/4080$ stated), and the ``unreachable by
construction'' note. \\
\addlinespace
Naming and address hygiene & 18 &
Bare literals for structural addresses; four roles under two names at
\texttt{0xb140}; the \mptr{PAIRING\_LHS}/\mptr{RHS} inversion; the undeclared
128-bit mask literal used for two unrelated purposes. Detailed in
\S\ref{sec:quality}. \\
\addlinespace
Operational and build hygiene & 9 &
\texttt{BUILD\_ID} identifying the verifying key more strongly than the
verifier; \texttt{--optimize-runs} being load-bearing but unrecorded in the
artifact; the deployment record omitting the fields needed to recompute
\texttt{BUILD\_ID}; no post-deployment way to re-run the capability probes. \\
\bottomrule
\end{tabularx}
\end{center}

\subsection{Prioritised recommendations}
\label{sec:assess:recs}

Ordered by value per unit of effort, not by severity.

\begin{enumerate}
\item \textbf{Get the staleness gate green} (ADD-1, MD-1, ADD-2). Until
  \texttt{template\_tree\_digest\_is\_pinned} passes and the provenance stamps
  are correct, nobody can re-derive these artifacts, which is the only real
  check on a generated file. Add the two machine-checkable guards proposed
  there --- the stamp check and the ``no identifier absent from
  \texttt{templates/}'' check --- both of which run without the SRS.
\item \textbf{Fix the failure taxonomy where it is measurably wrong} (CF-1,
  MF-4). The liveness-probe pattern is a small, local, failure-path-only change
  that makes \revert{PrecompileFailed} mean what the runbook says it means and
  moves 40\,\% of rejections into a proof-class error.
\item \textbf{Correct the three comments that overstate the code} (ACC-1's
  ``unreachable by construction'', MF-4's taxonomy promise, the header's size
  figures). These cost nothing and one of them --- the ``unreachable'' note ---
  actively invites a change that would reopen a malleability hole.
\item \textbf{Apply the address-naming policy uniformly} (Q-1, R-3). Emit
  \yul{add(BASE, offset)} instead of bare literals in the quotient VM and PCS
  emitters, name the three anonymous structural bases, and label the 78 MSM
  terms. All zero-bytecode-cost, and it is the single largest improvement
  available to a reviewer.
\item \textbf{Turn the invariants the code relies on into generator lints}
  (Q-2, Q-7, MD-5, DEAD-02): no literal gas operand in a \yul{staticcall}; every
  \yul{success := staticcall} inside an \yul{if success}; error name bound to
  error selector; every declared \mptr{*\_MPTR} either read or explicitly
  reserved. Each is a few lines and each closes a class rather than an instance.
\item \textbf{Record the truncation and the optimiser setting in the header}
  (theme 1, MD-6). Both are protocol- or deployability-critical parameters that
  a reader of the artifact alone cannot currently discover.
\end{enumerate}


%% ==== 18b-quality.tex ===================================================
\section{Critique: readability, maintainability and auditability}
\label{sec:quality}

The preceding sections judge the contract as a protocol implementation. This
one judges it as an artifact that humans and auditors have to read. Both matter,
and they pull in different directions: several of the things that make this
verifier fast are exactly the things that make it hard to review.

Two framing points before the specifics. First, this file is \emph{generated},
so every criticism below is really a criticism of an emitter, and every fix is a
change to a template rather than to a contract --- which makes them cheap and
uniform. Second, the artifact already sets a high bar for itself: 35\,\% of its
lines are comments, it carries a stated address-naming policy, a typed error
taxonomy, and inline provenance for its Rust counterparts. The findings below
are mostly about \emph{consistency} with the standard the file already sets,
not about the absence of a standard.

\subsection{What is done well}

\begin{itemize}
\item \textbf{The failure model is explicit and total in intent.}
  Success-or-revert with seven named errors, documented in NatSpec, is the right
  shape for a verifier. Contrast this with the common pattern of returning
  \texttt{bool} and letting callers forget to check.
\item \textbf{Guards are ordered by cost.} Accumulator decoding runs before the
  transcript; the ABI envelope runs before that. A malformed input is rejected
  in a few hundred thousand gas rather than 1.34\,M, and --- more importantly ---
  invalid public inputs cannot influence Fiat--Shamir derivation.
\item \textbf{Every precompile call forwards an exact scheduled gas bound}
  rather than \yul{gas()}, so a failing call burns the scheduled cost instead of
  $63/64$ of the transaction budget. This is a real design decision and it is
  applied nearly everywhere.
\item \textbf{Deployment-time capability probes.} \yul{require\_eip2537\_precompiles}
  executes \opcode{MCOPY}, \texttt{0x05}, and the EIP-2537 precompiles with
  known answers at the same gas bounds the runtime uses, so a chain that cannot
  run this verifier fails at deployment rather than on the first proof.
\item \textbf{The comments explain \emph{why}, not \emph{what}.} The MF-2,
  MF-4 and MF-12 notes, the LHS/RHS naming apology at lines 4101--4111, and the
  ``two bad equations could cancel'' argument at 4050--4062 are the kind of
  commentary that actually helps a reviewer.
\item \textbf{The verifying-key blob is arithmetically self-checking.} Its
  region sizes sum exactly to the pinned payload length, and every derived
  runtime pointer falls out of that sum (\S\ref{sec:appendix:vkmath}).
\end{itemize}

\subsection{Bare addresses: the single biggest readability cost}
\label{sec:quality:addresses}

The constants block states the policy explicitly: ``Every generated address in
that block is written as an offset from this base so a reader can see which term
a store belongs to; solc folds the addition away, so the symbolic form costs no
runtime bytes.'' That reasoning is correct and it generalises --- yet it is
applied to one block out of three.

\begin{center}
\begin{tabular}{@{}l r r r r@{}}
\toprule
Region & Lines & Bare \texttt{0xNNNN} & Named \mptr{*\_MPTR} & Bare share \\
\midrule
Helpers (683--1359) & 677 & 5 & 13 & 28\,\% \\
VK load + parser (1360--1791) & 432 & 8 & 24 & 25\,\% \\
Lagrange (1792--1916) & 125 & 4 & 20 & 17\,\% \\
Quotient VM (1917--3362) & 1{,}446 & 308 & 39 & \textbf{89\,\%} \\
PCS (3363--3811) & 449 & 203 & 55 & \textbf{79\,\%} \\
Final MSM + pairing (3812--4130) & 319 & 199 & 253 & 44\,\% \\
\midrule
Total & & \textbf{727} & \textbf{404} & \textbf{64\,\%} \\
\bottomrule
\end{tabular}
\end{center}

\begin{finding}[Q-1]{Two thirds of memory references are bare literals, and the worst-affected regions are the least reviewable ones --- Medium (auditability)}
\label{finding:q:addresses}
The hand-shaped regions --- helpers, parser, Lagrange --- are 17--28\,\% bare.
The two machine-generated regions a reviewer most needs help with are 79\,\% and
89\,\% bare. The final-MSM block, where the policy \emph{is} applied, sits at
44\,\% and is visibly easier to follow, which is the proof that the fix works.

The concrete harm is not abstract. \texttt{0xb280} appears 33 times in the
quotient VM as the running numerator accumulator, 6 times in the PCS block as
the base of the \yul{q\_eval} source table, and is \emph{also} the value of
\mptr{PCS\_FINAL\_MSM\_MPTR}. A reviewer who greps for
\texttt{PCS\_FINAL\_MSM\_MPTR} sees a fraction of its uses; a reviewer who reads
\yul{mstore(0xb280, 0)} in the VM has no way to learn that this address is
documented elsewhere under a different name and a different role. The same holds
for \texttt{0xb140}, which carries \emph{four} roles under \emph{two} names:
\mptr{SELECTOR\_ACC\_MPTR} (line 229), \mptr{BATCH\_INV\_SCRATCH\_MPTR}
(line 235), and the bare literals \yul{let msm\_scratch := 0xb140} in the
constructor smoke probe (line 549) and \yul{let acc\_scratch := 0xb140} in the
accumulator validator (line 1269). Half its uses are invisible to a reader
grepping either name.

Two structural addresses have no name at all despite being load-bearing:
\texttt{0xb2c0}, the base of the 49-entry $y$-power table, and \texttt{0xb8e0},
the base of the VM operand stack (which is also its lower bound, checked by
\yul{if eq(q\_sp, 0xb8e0) \{ q\_program\_fail() \}}).

\textbf{Recommendation}, in priority order:
\begin{enumerate}
\item Render every phase base as a named constant --- \texttt{0xb2c0} as
  \mptr{Y\_POWERS\_MPTR}, \texttt{0xb8e0} as \mptr{QVM\_STACK\_MPTR},
  \texttt{0xbfa0} as \mptr{QVM\_STACK\_END\_MPTR}, and give the numerator
  accumulator its own name (\mptr{QVM\_NUMERATOR\_MPTR}) even though it aliases
  \mptr{PCS\_FINAL\_MSM\_MPTR}. Aliases are fine; \emph{anonymous} aliases are
  not.
\item Extend the ``offset from a named base'' convention from the final-MSM
  block to the VM and PCS emitters. \yul{mload(0x94a0)} becomes
  \yul{mload(add(REVERSED\_EVALS\_MPTR, 0x20))}, which tells the reader it is
  evaluation~1. solc folds the addition, so this is free at runtime.
\item In the constants block, replace the prose note about phase reuse with a
  table: address, phase, role, lifetime. The information exists in the
  generator's arena allocator; publishing it in the artifact is what lets a
  reader check the claim instead of trusting it.
\end{enumerate}
\end{finding}

\subsection{Duplicated constants}

\begin{finding}[Q-2]{Named constants exist but are bypassed at three call sites --- Low}
\label{finding:q:dupes}
\begin{itemize}
\item Line 3999: \yul{staticcall(525096, 0x0c, PCS\_FINAL\_MSM\_MPTR, 0x30c0, ...)}.
  \texttt{525096} is exactly \mptr{G1MSM\_GAS\_SMOKE}, declared at line 295 and
  used by name in the constructor probe at line 555. The runtime path --- the
  one that matters --- uses the literal. A reviewer comparing the two has to
  notice they are the same number.
\item Line 500: \yul{staticcall(200000, 0x0c, scratch, 0xa0, scratch, 0x80)} in
  the negative smoke probe, a bare bound that corresponds to no declared
  constant and to no scheduled price.
\item Lines 1890--1893: \yul{q\_program\_fail} writes \texttt{0x3cc81b89}
  inline rather than \mptr{ERR\_QUOTIENT\_PROGRAM\_INVALID}. The comment
  explains why (the helper is shared with the standalone evaluator render), but
  the consequence is that \mptr{ERR\_QUOTIENT\_PROGRAM\_INVALID} has no reader
  in this artifact --- the constant and its only use have drifted apart.
\end{itemize}

\textbf{Recommendation}: emit a shared constants partial that both the verifier
and the evaluator include, and forbid bare numeric gas bounds in the emitters.
A generator-side assertion (``no \yul{staticcall} literal gas operand'') is a
five-line lint that makes this class impossible.
\end{finding}

\subsection{The quotient VM as a review object}

\begin{finding}[Q-3]{The inline-prefix and native-callback blocks are unreviewable line by line --- Medium (auditability)}
\label{finding:q:vm}
The bytecode VM is a good code-size trade and its interpreter is genuinely
readable: eleven opcodes, an explicit register set, and structural guards on
every stack transition. The problem is what surrounds it. The direct inline
prefix (lines 2052--2153) and the native callback bodies are long chains of
\yul{var0 \dots var162} over \yul{f\_3}, \yul{a\_0\_next\_1}, \yul{mload(0x9ac0)}.
Each line is trivially correct; the \emph{sequence} encodes a constraint-system
identity that no reader can recover from the artifact. Nothing in the file says
which gate \yul{var19} belongs to.

This is the region where correctness actually lives, and it is the region a
reviewer can do the least with. The mitigation that exists --- the differential
anchors in \texttt{src/lowering/quotient\_numerator/native\_anchor.rs} --- is
the right one, but it lives outside the artifact.

\textbf{Recommendation}: emit, for each identity block, a one-line header naming
its source (\texttt{gate "poseidon full round", polynomial 3 of 7}) and its
position in the $y$-batch. The generator knows both. This costs nothing at
runtime and converts an opaque block into a checkable claim. Second, name the
evaluation loads: \yul{mload(add(REVERSED\_EVALS\_MPTR, 0x20))} with a comment
\texttt{// advice[0] at rotation 0} beats \yul{mload(0x94a0)}.
\end{finding}

\subsection{Comment accuracy}

A file that is 35\,\% comments is making a large implicit promise. Most of it is
kept; the exceptions matter more than they would in a sparsely commented file,
because a reviewer calibrates on the average.

\begin{finding}[Q-4]{Three comments claim more than the code delivers --- Low}
\label{finding:q:comments}
\begin{enumerate}
\item \textbf{MF-4, lines 1240--1245}: states that a point ``the precompile
  found off-curve/out-of-subgroup'' is reported as \revert{BadPointEncoding}.
  It cannot be, and measurably is not: 40\,\% of all rejections carry
  \revert{PrecompileFailed} instead (\S\ref{sec:testing:taxonomy}).
\item \textbf{Header, lines 8--11}: the quoted runtime sizes (29{,}567 and
  29{,}836 bytes) do not reproduce with the pinned toolchain, which yields
  21{,}880 and 31{,}176 (\S\ref{sec:appendix:size}). The conclusion survives;
  the numbers do not.
\item \textbf{Header, lines 30--45}: the summary of differences from the BN254
  template does not mention that $x_3$ and the $x_1$/$x_4$ power tables are
  truncated to 128 bits, even though that is a protocol parameter which must
  match the prover exactly and which the crate's own \texttt{Cargo.toml} warns
  about.
\end{enumerate}

\textbf{Recommendation}: treat header claims as testable. The sizes should be
computed at render time; the taxonomy claim should be either fixed in code or
corrected in prose; truncation belongs in the header's feature summary.
\end{finding}

\subsection{Naming}

\begin{finding}[Q-5]{\mptr{PAIRING\_LHS\_MPTR} and \mptr{PAIRING\_RHS\_MPTR} do not name pairing arguments --- Low}
\label{finding:q:naming}
The file admits this itself at lines 4101--4111: the names follow the dual-MSM
accumulator (left $=\pi$, right $=$ combined), not the pairing argument order,
so the call site passes them \emph{swapped}:
\yul{ec\_pairing(success, PAIRING\_RHS\_MPTR, PAIRING\_LHS\_MPTR)}. A ten-line
comment explaining why an argument order looks wrong is a strong signal that
the names are wrong. \yul{ec\_pairing}'s own parameters are called
\yul{lhs\_mptr}/\yul{rhs\_mptr} and mean the \emph{opposite} thing, so the two
naming systems collide inside one call.

\textbf{Recommendation}: rename to \mptr{KZG\_PROOF\_MPTR} ($\pi$) and
\mptr{KZG\_CLAIM\_MPTR} ($\mathsf{C}-vG+x_3\pi$), or name them for their pairing
role (\mptr{PAIRING\_G1\_A}, \mptr{PAIRING\_G1\_B}). Either removes the swap and
the comment that defends it. The same applies to \yul{trash}/\yul{trashcan},
which is upstream vocabulary but reads as a placeholder; a rendered glossary
line would cost nothing.
\end{finding}

\subsection{Defensive code that cannot fire}

Several guards are documented as unreachable. Keeping them is defensible ---
they are cheap and they protect against future emitter changes --- but the
artifact should distinguish ``this rejects malicious input'' from ``this
rejects a codegen bug'', because a reviewer budgets attention differently for
the two.

\begin{center}
\begin{tabularx}{\textwidth}{@{}l X l@{}}
\toprule
Location & Guard & Reachable by an attacker? \\
\midrule
1188--1210 & \yul{if is\_id} branch in \yul{load\_acc\_point} & No (comment says so) \\
1143--1146 & \yul{lt(hi, shl(128, 1))} after the field check & No (implied by the field check) \\
2218, 2229 & VM stack overflow / underflow & No (program is codehash-pinned) \\
3094--3105 & VM epilogue invariants & No (same) \\
4124 & final \yul{if iszero(success)} & No (every failure already reverted) \\
839--844 & \yul{batch\_invert} reversed-range check & No (ranges are generated) \\
\bottomrule
\end{tabularx}
\end{center}

\begin{finding}[Q-6]{Attacker-reachable and codegen-only guards are indistinguishable --- Low}
\label{finding:q:guards}
All six guards above raise errors from the same taxonomy as guards that do
reject attacker input, so an incident responder cannot tell from a selector
whether they are looking at a hostile transaction or a broken build.
\revert{QuotientProgramInvalid} is the one error that is purely a build fault,
and it is documented as such --- which shows the distinction is already
understood.

\textbf{Recommendation}: adopt a convention in the emitter, e.g. a
\texttt{// CODEGEN-ONLY:} marker on every guard that a pinned artifact cannot
reach, and route them all to \revert{QuotientProgramInvalid} or a new
\texttt{BuildInvariantViolated}. The existing \revert{MemoryLayoutViolated} is
exactly this pattern applied to one case; generalising it is a small change with
a real payoff for triage.
\end{finding}

\subsection{The \yul{success} plumbing}

The contract mixes two idioms: \yul{success := and(success, cond)} folded to a
later check (18 sites), and \yul{success := staticcall(...)} which
\emph{overwrites} (12 sites). Every overwrite is inside an \yul{if success}
block, so no cleared flag is resurrected --- and the guard is also a gas
optimisation, since it skips the call on the failure path. The discipline is
correct as written.

\begin{finding}[Q-7]{The overwrite idiom is correct but locally unverifiable --- Observation}
\label{finding:q:success}
Checking this property requires reading all 30 sites; a single emitted
\yul{success := staticcall(...)} outside an \yul{if success} guard would silently
turn a failed earlier check into a pass, and nothing in the artifact would show
it. The property is worth stating as an invariant the generator enforces --- and
it is cheap to enforce: the emitter that renders a precompile call already knows
whether it is inside a guard.
\end{finding}

\subsection{Summary of the critique}

\begin{center}
\begin{tabularx}{\textwidth}{@{}l l X@{}}
\toprule
ID & Severity & Issue \\
\midrule
Q-1 & Medium (auditability) & 64\,\% of memory references are bare literals; three anonymous aliases \\
Q-2 & Low & Gas bounds and an error selector duplicated as literals \\
Q-3 & Medium (auditability) & Identity blocks carry no link to the constraint they encode \\
Q-4 & Low & Three header/inline claims are inaccurate or incomplete \\
Q-5 & Low & Pairing slot names invert their pairing role \\
Q-6 & Low & Codegen-only guards share the attacker-facing error taxonomy \\
Q-7 & Observation & \yul{success} overwrite discipline is correct but not locally checkable \\
\bottomrule
\end{tabularx}
\end{center}

None of these is a soundness problem. All seven are emitter changes, and five of
them (Q-1, Q-2, Q-3, Q-6, Q-7) can be enforced mechanically by a generator-side
lint rather than by review discipline --- which is the right level to fix them
at, because this artifact will be regenerated many times and reviewed by
different people each time.


%% ==== 20-robustness.tex =================================================
\section{Are these the right checks? Redundancy, ad-hoc-ness and simplification}
\label{sec:robust}

Section~\ref{sec:assess:summary} asks whether the verifier is sound.
Section~\ref{sec:quality} asks whether it is readable. This section asks a third
question that neither covers: are the checks it performs \emph{the right checks,
expressed the right way}? Is anything redundant, anything ad hoc, anything
fragile in the sense of holding by coincidence rather than by construction ---
and is there a simpler formulation that loses no rejection?

Six regions were analysed and every proposal was then handed to a second
reviewer instructed to break it: show that the rewrite loses a rejection, or
that the claimed saving is wrong. That pass did real work --- it \emph{compiled}
the proposed variants and measured them, it re-derived the algebra over random
field draws, and in two cases it built the change into the fixture and re-ran the
adversarial suites against the modified contract. Eight proposals were rejected
outright, including one that would have produced an undeployable artifact.

\begin{center}
\begin{tabular}{@{}l r@{}}
\toprule
Outcome after verification & Count \\
\midrule
Adopt (proposed and confirmed) & 23 \\
Adopt (found by the verifying pass itself) & 24 \\
Safe but marginal & 14 \\
Rejected as unsafe & 2 \\
Rejected on a wrong premise & 6 \\
\midrule
Total & 69 \\
\bottomrule
\end{tabular}
\end{center}

\subsection{The result that matters most: a rejected simplification}
\label{sec:robust:rejected}

The most instructive outcome is a proposal that looked obviously right and was
not.

\begin{finding}[RB-1]{Batching the five PCS inversions on the stack would push the contract past EIP-170 --- Rejected}
\label{finding:robust:pcsbatch}
The PCS block runs five independent Montgomery ladders, each ending in its own
\yul{scalar\_inv} --- five \texttt{modexp} calls where the contract already
contains a \yul{batch\_invert} helper that inverts an arbitrary run with one.
The algebra permits batching: the dependency graph is one-directional
($x_3$ and the four rotation points $\rightarrow$ 21 batch inputs $\rightarrow$
5 totals $\rightarrow$ 5 inverses), no set consumes another's inverse, and the
rewrite was verified to give an identical $f_{\mathrm{eval}}$ over 200 random
field draws.

The verifying pass compiled all variants rather than estimating:

\begin{center}
\begin{tabular}{@{}l r r@{}}
\toprule
Variant & Runtime bytes & $\Delta$ \\
\midrule
Baseline & 21{,}924 & --- \\
Batch the five totals on the stack & 29{,}614 & $+7{,}690$ \\
Batch the five totals through memory & 22{,}520 & $+596$ \\
\textbf{One 14-element memory batch} & \textbf{19{,}985} & $\mathbf{-1{,}939}$ \\
\bottomrule
\end{tabular}
\end{center}

(Compiled with solc 0.8.35 rather than the pinned 0.8.30, so read the deltas,
not the absolutes.) The stack variant --- the one that ``needs no memory'' and
looks cheapest --- is $5{,}038$ bytes \emph{over} the EIP-170 limit and would not
deploy. Bytecode in this artifact is driven by stack pressure inside one giant
assembly block, so a rewrite that keeps more values live costs code size
super-linearly. The right rewrite is the one nobody proposed first: batch all
fourteen distinct denominators in a single memory run and read the inverses back
from memory. That is $1{,}939$ bytes \emph{smaller} than the baseline and four
\texttt{modexp} calls cheaper, with the same rejection set --- \yul{batch\_invert}
rejects iff the product of the fourteen is zero, which is exactly when some
factor is zero, which is exactly what the five \yul{scalar\_inv} calls reject
today.

\textbf{The lesson for this codebase}: in a contract whose deployability depends
on 2{,}696 bytes of EIP-170 headroom, no code-size claim should be accepted
without a compile. Two of the twelve dimensions in this section produced
proposals whose \emph{sign} was wrong.
\end{finding}

\subsection{Adopt: measured, verified simplifications}
\label{sec:robust:adopt}

\begin{finding}[RB-2]{The accumulator codec has four mechanisms for one concept --- Adopt, $-367$ bytes and $\approx 3{,}000$ gas}
\label{finding:robust:acc}
``Is this the point at infinity'' is decided four times in
\texttt{Halo2Verifier.sol}:1010--1359: an exact four-word sentinel match
(\yul{is\_acc\_encoded\_identity}, 1065--1079), a flag prefilter and probe inside
\yul{load\_acc\_coord} (1108--1114), the \yul{if is\_id} arm in
\yul{load\_acc\_point} (1189--1210), and the \yul{decoded\_zero} rejection
(1212--1216). The probe costs a \emph{second} full seven-limb decode of every
$x$ coordinate on every proof, because a real $x$'s first word is 224 bits and
therefore always clears the \yul{>= base} prefilter.

The probe and the \yul{if is\_id} arm are provably reject-only: reaching 1189
with \yul{is\_id} set implies the sentinel check failed, which for a
packing-clean flagged $x$ implies $y$ is not the sentinel, which implies
$y \neq 0$, which makes the guard at 1201 zero. So both can go, and the
rejection set --- including the error selector --- is unchanged: the class they
used to catch now fails the field bound at 1126 before any staticcall runs, so
\yul{precompile\_failed} stays $0$ and the call site still raises
\revert{BadPointEncoding}. Measured: $21{,}912 \rightarrow 21{,}545$ runtime
bytes.

\textbf{Critical caveat, and the reason this is written down rather than just
done}: the rewrite must be applied \emph{atomically}. Deleting the
\yul{if is\_id} arm while keeping the probe is strictly worse than losing a
rejection --- \yul{load\_acc\_point} would fall through writing nothing to
\yul{dst}, leaving \mptr{ACC\_LHS\_MPTR} and \mptr{ACC\_RHS\_MPTR} holding stale
memory. The test at
\texttt{fixture\_verifier\_validates\_the\_public\_accumulator\_codec} case (c2)
covers exactly this class and must keep passing.
\end{finding}

\begin{finding}[RB-3]{A cheaper field bound with a bit-identical rejection set --- Adopt, $528$ gas per verification}
\label{finding:robust:fplo}
The Fp membership test is spelled
\yul{or(lt(hi, BLS\_P\_HI), and(eq(hi, BLS\_P\_HI), iszero(gt(lo, BLS\_P\_MINUS\_ONE\_LO))))}
at three sites. Since $p_{\mathrm{lo}}$ ends \texttt{...aaab} and
$p_{\mathrm{lo}}-1$ ends \texttt{...aaaa} --- differing by one with no borrow
into the high word, and $p_{\mathrm{lo}} \neq 0$ so there is no wraparound edge
--- \yul{iszero(gt(lo, BLS\_P\_MINUS\_ONE\_LO))} is bit-identical to
\yul{lt(lo, BLS\_P\_LO)} for every 256-bit \yul{lo}, and one opcode shorter.

This one was not argued, it was executed: the change was built into the fixture
together with three related cleanups, and all four adversarial suites were re-run
(22/22 pass) plus the exhaustive single-bit sweep --- 17{,}032 mutations, none
accepted, and the selector histogram \emph{byte-identical} to the baseline
(4 / 119 / 6840 / 7433 / 2376 / 260). That is the standard this kind of change
should be held to.
\end{finding}

\begin{finding}[RB-4]{A masking operation that is the identity function --- Adopt, $28$ bytes and $\approx 200$ gas}
\label{finding:robust:mask}
\yul{common\_uncompressed\_g1} reaches lines 783--784 only after
\yul{if shr(128, x\_hi\_word) \{ fail(...) \}}, and \yul{fail} is an
unconditional revert, not a flag. So at 783 the value is already below $2^{128}$
and \yul{and(x\_hi\_word, 0xffff...ffff)} is the identity. The masked locals feed
only the field bound; the transcript gets the raw calldata bytes via
\yul{calldatacopy}. The optimiser does not remove it, and it runs 68 times per
proof (34 points, two coordinates).
\end{finding}

\begin{finding}[RB-5]{An unlabelled ordering dependency carries the whole bound on eight attacker-controlled words --- Adopt (documentation and a build-time assertion)}
\label{finding:robust:order}
\yul{validate\_public\_accumulator} runs at line 1433, \emph{before} the instance
canonicality loop at 1511--1522 that enforces \yul{lt(inst\_be, r)} on all
nineteen public inputs. At the moment the eight accumulator words are decoded,
\yul{check\_acc\_coord\_packing} is therefore the \emph{only} bound established
on them anywhere in the program --- not one of several.

That makes its \yul{if lt(used\_bits, 256)} guard more dangerous than it looks.
The guard exists to avoid a \yul{shl} overflow when a word's limbs would occupy
$\geq 256$ bits, and it fails \emph{open}: a future render with
$\mathit{bits} \times \mathit{limbs\_per\_word} \geq 256$ emits no check for that
word at all, while the decoder silently truncates its top limb. The ordering
means nothing downstream would catch it.

The region is correct as rendered ($4 \times 56 = 224$ and $3 \times 56 = 168$,
and \mptr{BLS\_P\_MINUS\_ONE\_PACKED\_0\_WITH\_ID\_FLAG} has bit length exactly
$224$, so the bound is tight rather than lucky). The fix is not runtime code:
state the ordering dependency in the comment, and add a generation-time assertion
that $\mathit{bits} \times \mathit{limbs\_per\_word} < 256$.
\end{finding}

\subsection{Full disposition}
\label{sec:robust:table}

\begingroup
\footnotesize
\setlength{\tabcolsep}{4pt}
\begin{longtable}{@{}l l l p{0.46\textwidth}@{}}
\toprule
ID & Area & Disposition & Proposal \\
\midrule
\endfirsthead
\toprule
ID & Area & Disposition & Proposal \\
\midrule
\endhead
\texttt{P4} & PCS & Adopt & MAX\_ROTATION\_WALK\_STEPS = 4096 does not enforce the property its own comment says it enforces \\
\addlinespace
\texttt{P5} & PCS & Adopt & The 24 free rotation-table words are addressing headroom, not spare scratch — and BATCH\_INV\_SCRATCH\_MPTR is not free at PCS time \\
\addlinespace
\texttt{P6} & PCS & Adopt & The unrolled ladders buy nothing measurable in gas and cost about a third of the remaining EIP-170 headroom \\
\addlinespace
\texttt{ACC-1} & accumulator codec & Adopt & Four identity mechanisms collapse to one: the flag probe and the `if is\_id` arm are provably reject-only \\
\addlinespace
\texttt{ACC-3} & accumulator codec & Adopt & `check\_acc\_coord\_packing`'s loop is two `shr`s, and `if lt(used\_bits,256)` fails open for a future render \\
\addlinespace
\texttt{ACC-4} & accumulator codec & Adopt & The Fp-membership predicate is inlined three times; the 128-bit padding bound is spelled two ways \\
\addlinespace
\texttt{FR-1} & fragility & Adopt & Transcript-buffer bound is derived at generation time but never restated on-chain \\
\addlinespace
\texttt{FR-4} & fragility & Adopt & NUM\_INSTANCE\_CPTR assumes the proof payload is word-aligned; nothing asserts it \\
\addlinespace
\texttt{G5} & guard discipline & Adopt & The two-value `(out, precompile\_failed)` protocol pushes selector choice to the caller and uses `:=` where it needs `or` \\
\addlinespace
\texttt{FP-1} & range idioms & Adopt & One `fp\_is\_canonical(hi\_word, lo)` replaces three inline spellings AND makes both 2\^{}128 pad guards fall out for free \\
\addlinespace
\texttt{FP-3} & range idioms & Adopt & Three Fr-canonicality loops, two error idioms, and a boundary check at 1789 that provably cannot fire \\
\addlinespace
\texttt{FP-4} & range idioms & Adopt & `if lt(used\_bits, 256)` is a special case guarding a `shl` overflow trap; `iszero(shr(used\_bits, w))` is the general predicate… \\
\addlinespace
\texttt{FP-6} & range idioms & Adopt & The identity-flag 'cheap prefilter' at 1108 is not an optimisation -- it is the precondition that keeps `sub(packed,… \\
\addlinespace
\texttt{P1} & test suite & Adopt & reject\_early's gas anchor measures the EIP-7623 calldata floor, not execution depth \\
\addlinespace
\texttt{P11} & test suite & Adopt & Concrete attack classes with no case in the suite \\
\addlinespace
\texttt{P12} & test suite & Adopt & Four assertions that cannot detect the thing they are named for \\
\addlinespace
\texttt{P2} & test suite & Adopt & Delete reject\_late (one call site, 5.8\% above the floor) and use reject\_at\_pairing \\
\addlinespace
\texttt{P3} & test suite & Adopt & reject\_at\_pairing's 3/4-of-total floor admits the EIP-2537 precompile-failure band \\
\addlinespace
\texttt{P4} & test suite & Adopt & The per-proof verifying-key re-check is never executed; ERR\_VK\_MISMATCH is pinned but unreachable by any test \\
\addlinespace
\texttt{P5} & test suite & Adopt & Section 5's curve-deferral loop accepts either of two selectors at any depth \\
\addlinespace
\texttt{P7} & test suite & Adopt & encode\_coord is raw\_limbs with a shift; the packing loop is written twice, three lines apart \\
\addlinespace
\texttt{P8} & test suite & Adopt & Three tests assert only 'it reverted' where the measured outcome supports an exact selector and depth \\
\addlinespace
\texttt{P9} & test suite & Adopt & The layout pin omits ERR\_PRECOMPILE\_FAILED, cannot detect a sum-preserving section shift, and maps unparseable VK words to 0 \\
\addlinespace
\texttt{A1} & PCS & Adopt (verifier) & The measured right rewrite: one 14-element memory batch, inverses read from memory — 1,939 bytes smaller AND 4 modexps cheaper \\
\addlinespace
\texttt{A2} & PCS & Adopt (verifier) & The `m == 1` emitter special case is already dead weight — the general path renders correct code for cardinality-1 sets today \\
\addlinespace
\texttt{A3} & PCS & Adopt (verifier) & `batch\_invert`'s fail-closed contract is two separate obligations at two separate line numbers, enforced by convention at the… \\
\addlinespace
\texttt{ADD-1} & accumulator codec & Adopt (verifier) & The identity sentinel is frozen at (bits=56, n=7, limbs\_per\_word=4) while every other part of the codec is templated — false… \\
\addlinespace
\texttt{ADD-2} & accumulator codec & Adopt (verifier) & The 128-bit mask at 783-784 is dead code that the optimizer does not remove — 28 bytes and \textasciitilde{}200 gas over 34 call sites \\
\addlinespace
\texttt{ADD-3} & accumulator codec & Adopt (verifier) & The accumulator codec runs BEFORE the instance canonicality loop, so check\_acc\_coord\_packing is the only bound on those eight… \\
\addlinespace
\texttt{AV-1} & fragility & Adopt (verifier) & The VM operand clamp is derived from a cross-render constant table when it could be derived from the pinned program, for the… \\
\addlinespace
\texttt{AV-2} & fragility & Adopt (verifier) & MF-3's u8-vs-u16 split is a hand-written case analysis of one arithmetic fact that is asserted nowhere; one generation-time… \\
\addlinespace
\texttt{AV-3} & fragility & Adopt (verifier) & The runtime's largest gas bound is a bare literal at its call site; the constructor probe covers it only because both are… \\
\addlinespace
\texttt{AV-4} & fragility & Adopt (verifier) & The transcript bound is three hand-derived sums that must each mirror the emitter's absorb schedule, and one of them already… \\
\addlinespace
\texttt{V1} & guard discipline & Adopt (verifier) & Codify the rule that G3 violates: a staticcall result and its returndatasize check must be separate statements \\
\addlinespace
\texttt{V2} & guard discipline & Adopt (verifier) & `if acc\_msm\_len \{` at 1323 is dead by construction and is the accumulator twin of the dead `success` guards G1 lists \\
\addlinespace
\texttt{V3} & guard discipline & Adopt (verifier) & The 78-pair MSM gas bound exists twice — named at 295, bare literal at 3999 — and only the named copy is checked by the… \\
\addlinespace
\texttt{V4} & guard discipline & Adopt (verifier) & 'Canonical and nonzero in Fr' is spelled three different ways in two adjacent helpers \\
\addlinespace
\texttt{FP-A1} & range idioms & Adopt (verifier) & Adopt `lt(lo, BLS\_P\_LO)` after all -- it is 528 gas per verification cheaper, which FP-2 did not measure \\
\addlinespace
\texttt{FP-A2} & range idioms & Adopt (verifier) & The 128-bit truncation rule is a protocol requirement spelled as a bare literal in 8 places with no shared name \\
\addlinespace
\texttt{FP-A3} & range idioms & Adopt (verifier) & check\_acc\_coord\_packing is NOT redundant with the 1126 range check, and the source does not say why -- it is the next thing a… \\
\addlinespace
\texttt{FP-A4} & range idioms & Adopt (verifier) & scalar\_inv's two guards fuse into one shifted-origin range check -- measured, and the numbers say do not bother \\
\addlinespace
\texttt{A1} & test suite & Adopt (verifier) & Even after P1, the anchor is blind to the first 176,934 gas of execution — and no tx-level measurement can recover it \\
\addlinespace
\texttt{A2} & test suite & Adopt (verifier) & The same floor-blind anchor is copy-pasted into both sibling suites — one bug in three files, which is the real case for P10 \\
\addlinespace
\texttt{A3} & test suite & Adopt (verifier) & The proof layout is hardcoded twice inside this one file: SECTIONS and the `bounds` literal \\
\addlinespace
\texttt{A4} & test suite & Adopt (verifier) & `no\_single\_bit\_mutation\_of\_the\_fixture\_calldata\_verifies` names a universal it does not test — it covers 25\% of single-bit… \\
\addlinespace
\texttt{A5} & test suite & Adopt (verifier) & `envelope\_class` panics on any revert payload that is neither 0 nor 4 bytes, converting a taxonomy failure into a helper crash \\
\addlinespace
\texttt{A6} & test suite & Adopt (verifier) & Sections 5 and 7 are the only sweeps with no rejection-count drift guard, and section 5 hides a `continue` of the same kind P12… \\
\addlinespace
\texttt{P1} & PCS & Rejected: unsafe & Five modexps for one inversion problem: batch the per-set totals (or recompute the forward pass) to get to one \\
\addlinespace
\texttt{G3} & guard discipline & Rejected: unsafe & Nine PCS/pairing staticcalls defer to a single far-away enforcement point \\
\addlinespace
\texttt{P2} & PCS & Rejected: premise & Only 4 distinct dx values exist; the code computes 11 of them and inverts all 11 \\
\addlinespace
\texttt{ACC-6} & accumulator codec & Rejected: premise & Named pointers vs one cursor, and `if acc\_msm\_len` fails open \\
\addlinespace
\texttt{FR-2} & fragility & Rejected: premise & Constant-table index is clamped in three VM arms and unclamped in seven \\
\addlinespace
\texttt{FR-5} & fragility & Rejected: premise & Phase ordering is enforced for template sections but blind to Yul helper call sites and to reads after a region's phase \\
\addlinespace
\texttt{FP-2} & range idioms & Rejected: premise & The three Fp spellings are character-identical, and the `iszero(gt(lo, P\_MINUS\_ONE\_LO))` form is measurably better than the… \\
\addlinespace
\texttt{P6} & test suite & Rejected: premise & r is hardcoded twice and p once, none cross-checked against the artifact's own moduli \\\\
\bottomrule
\end{longtable}
\endgroup

\subsection{The same question, asked of the test suite}
\label{sec:robust:tests}

The review covered the adversarial suite as well, and found a defect in it that
is worth recording because it is the failure mode every gas-anchored replay test
is exposed to.

\begin{finding}[RB-6]{The gas anchors were measuring EIP-7623's calldata floor, not execution --- fixed}
\label{finding:robust:floor}
revm charges
$\mathit{gas\_used} = \max(\mathit{legacy\_intrinsic} + \mathit{execution},\;
\mathit{floor}_{7623})$.
For this 8{,}516-byte calldata (1{,}525 zero and 6{,}991 non-zero bytes) that is
a legacy intrinsic of $138{,}956$ against a floor of $315{,}890$. Every cheap
rejection therefore bills at exactly the floor, and the original anchor
--- $\mathit{gas} \times 2 < \mathit{accepted}$, i.e.\ $\mathit{gas} < 668{,}486$
--- was satisfied by the floor alone. A flipped selector, which executes only the
Solidity dispatcher, and a guard that ran the accumulator decoder and two G1MSMs
both report $315{,}890$. The anchor measured the calldata size and nothing else.

The suite now computes both quantities from each mutated input's own byte
composition and requires an early rejection to bill \emph{at} its floor. What
that guarantees is bounded, and the code says so rather than overclaiming:
execution below $\mathit{floor} - \mathit{legacy} = 176{,}934$ gas is
unobservable at transaction level, which is $17.3\,\%$ of an accepted run's
execution. The proof that the residual gap is real, not theoretical: the
\texttt{all-zero instances} case bills at exactly its own floor while reverting
\revert{PrecompileFailed}, which can only come from an EIP-2537 staticcall that
ran and failed.

All 13 tests still pass under the tightened anchors, which is the useful part of
the result: the guards really do fire where the suite claimed, the old
assertion just could not have detected it if they had not.
\end{finding}

Six further changes were made to the suite as a result of this review: the
section table is now derived from the verifier's own parser loop bounds rather
than written out twice; the two field moduli are matched against the artifact's
constants; \yul{envelope\_class} reports an unexpected payload through the
caller's assertion instead of panicking inside a helper; the curve-substitution
sweep now pins the exact selector and depth (all 24 cases land on
\revert{PrecompileFailed} after $\approx 1.04$M gas of execution) instead of
accepting either of two selectors at any depth; two sweeps gained
rejection-count drift guards; and the exhaustive sweep was renamed to
\texttt{no\_low\_or\_high\_bit\_flip\_...}, because it covers two bits per byte
and therefore $25\,\%$ of single-bit mutations, not all of them.

One coverage gap is documented rather than fixed: the verifier re-checks the
verifying key's \texttt{extcodesize} and \texttt{extcodehash} on \emph{every}
proof, not only at construction. Reaching that guard needs the code at the
pinned address to change after deployment, which the revm harness cannot
express, so \revert{VkMismatch} is unreachable from this file.


%% ==== 21-appendix.tex ===================================================
\appendix

\section{Appendix: reproducing and checking these artifacts}
\label{sec:appendix:repro}

\subsection{Toolchain}

\begin{center}
\begin{tabularx}{\textwidth}{@{}l X@{}}
\toprule
Component & Value \\
\midrule
Compiler & \texttt{solc 0.8.30+commit.73712a01}, pinned by the artifact's
\texttt{pragma solidity 0.8.30} (not a floating \texttt{\^{}} range) \\
Flags & \texttt{--via-ir --optimize --optimize-runs 1 --evm-version cancun
--no-cbor-metadata} \\
EVM target & Cancun opcodes (\texttt{MCOPY}) plus the Prague EIP-2537
precompiles \texttt{0x0b}--\texttt{0x0f}, and \texttt{0x05} modexp \\
Rust toolchain & \texttt{1.90.0} (\texttt{rust-toolchain.toml}) \\
Test harness & \texttt{revm 19} with \texttt{blst}, \texttt{SpecId::PRAGUE} \\
\bottomrule
\end{tabularx}
\end{center}

The \texttt{--via-ir} flag is not optional: without it the contract fails to
compile with \emph{Stack too deep} in \yul{load\_acc\_coord\_shifted}.

\subsection{Size and the deployability envelope}
\label{sec:appendix:size}

The optimiser-runs setting is load-bearing, and the artifact header says so.
Measured runtime sizes for \texttt{Halo2Verifier.sol}:

\begin{center}
\begin{tabular}{@{}l r r@{}}
\toprule
\texttt{--optimize-runs} & Runtime bytes & EIP-170 headroom \\
\midrule
1 & 21{,}880 & $+2{,}696$ \\
200 & 21{,}912 & $+2{,}664$ \\
100000 & 31{,}176 & $-6{,}600$ (undeployable) \\
\bottomrule
\end{tabular}
\end{center}

\texttt{Halo2VerifyingKey.sol} deploys to 17{,}025 bytes regardless of the
optimiser setting, because its runtime is data returned by the constructor
rather than compiled code --- \texttt{--bin-runtime} reports three bytes for it,
which is the compiled runtime section, not the deployed one.

\begin{finding}[A-1]{The header's measured sizes do not match this toolchain --- Informational}
The artifact header states ``0.8.24 at runs=1 emits 29,567 bytes and 0.8.30 at
runs=100000 emits 29,836 --- both over the EIP-170 24,576-byte limit, so neither
can be deployed''. Compiled here with the pinned \texttt{solc 0.8.30} and the
flags the crate's own \yul{compile\_solidity\_with\_runs} uses, runs=1 gives
21{,}880 bytes and runs=100000 gives 31{,}176. The conclusion the header draws
(only the pinned pair is known to be deployable) survives, but the specific
figures do not reproduce, most likely because they were taken without
\texttt{--via-ir} or without \texttt{--no-cbor-metadata}.

\textbf{Recommendation}: record the exact command line next to the numbers, or
derive them at render time, so a reader can reproduce them. A size figure a
reviewer cannot reproduce is worse than no figure.
\end{finding}

\subsection{Verifying-key payload arithmetic}
\label{sec:appendix:vkmath}

The verifying-key contract's payload is self-checking by arithmetic. Counting
the generated \yul{mstore} lines by their trailing comment:

\begin{center}
\begin{tabular}{@{}l r r l@{}}
\toprule
Region & Words & Bytes & Runtime address \\
\midrule
Header (\texttt{vk\_digest}\dots\texttt{neg\_s\_g2}) & 31 & \texttt{0x03e0} & \mptr{VK\_MPTR} $=$ \texttt{0x3680} \\
\texttt{quotient\_const} & 178 & \texttt{0x1640} & \texttt{0x3a60} $=$ \yul{q\_const\_mptr} \\
\texttt{quotient\_program} & 143 & \texttt{0x11e0} & \texttt{0x50a0} $=$ \yul{q\_program\_mptr} \\
\texttt{fixed\_comms} (27 points) & 108 & \texttt{0x0d80} & \texttt{0x6280} \\
\texttt{permutation\_comms} (18 points) & 72 & \texttt{0x0900} & \texttt{0x7000} \\
\midrule
Total payload & 532 & \texttt{0x4280} $=17{,}024$ & ends at \texttt{0x7900} $=$ \mptr{CHALLENGE\_MPTR} \\
Runtime (with \texttt{INVALID}) & --- & \texttt{0x4281} $=17{,}025$ & \\
\bottomrule
\end{tabular}
\end{center}

Every derived pointer in the verifier agrees with this table:
\texttt{0x3680} $+$ \texttt{0x03e0} $=$ \texttt{0x3a60}, \texttt{0x3680} $+$
\texttt{0x1a20} $=$ \texttt{0x50a0}, and the payload ends exactly where the
challenge slots begin. \texttt{return(runtime, 0x4281)} matches
\texttt{EXPECTED\_VK\_LENGTH}, and \texttt{EXPECTED\_VK\_PAYLOAD\_LENGTH} is one
less. The deployed codehash equals \texttt{EXPECTED\_VK\_CODEHASH}; this is
checked by
\texttt{fixture\_artifacts\_fit\_in\_a\_deployable\_contract}.

The executable bytecode is \texttt{0x11cf} $=4{,}559$ bytes, so 17 of the 4{,}576
payload bytes reserved for the program are zero tail padding. That padding is
not decorative: operand reads use whole-word \yul{mload(q\_pc)}, so an
instruction near the end reads up to 31 bytes past \yul{q\_end}
(\S\ref{sec:qvms:frame}).

\subsection{What a third party can check without the SRS}

A reviewer who does not have the structured reference string, the Moonlight
checkout, or a proving run can still check a useful amount:

\begin{enumerate}
\item \textbf{The artifacts are self-consistent.} The payload arithmetic above,
  the pointer derivations, the proof-layout tiling
  (\S\ref{sec:tp:layout}), and the memory map
  (\S\ref{sec:const:fullmap}) are all closed computations over the two files.
\item \textbf{The fixture proof verifies and nothing near it does.} The replay
  suite needs only \texttt{solc} and \texttt{revm}
  (\S\ref{sec:testing}).
\item \textbf{The error selectors are what they claim to be.} Each is
  \texttt{bytes4(keccak256("Name()"))}; all eight were recomputed
  independently for this document and match.
\item \textbf{The gas bounds match the published schedules.} \texttt{G1ADD}
  $=375$, \texttt{G1MSM} for one pair $=12{,}000$, the pairing bound
  $32{,}600\cdot 2 + 37{,}700 = 102{,}900$, and the 78-pair MSM bound
  $525{,}096$ (implying an EIP-2537 discount of 561 at $k=78$).
\end{enumerate}

What cannot be checked this way is the only thing that ultimately matters for
soundness: that the generated arithmetic is the \emph{same} arithmetic the
native Rust verifier performs. That is the job of the differential anchors in
\texttt{src/lowering/} (\texttt{transcript\_anchor.rs},
\texttt{codec\_anchor.rs}, \texttt{quotient\_numerator/native\_anchor.rs}) and
of the 247 tests under \texttt{src/}. A reviewer of this artifact who does not
also review those anchors has checked that the contract is internally coherent,
not that it verifies the right protocol.

\end{document}
